\documentclass[11pt,oneside]{book}
\usepackage[margin=1in]{geometry}
\usepackage{amsmath,amssymb,amsthm,mathtools}
\usepackage{booktabs,longtable,array}
\usepackage{enumitem}
\usepackage[colorlinks=true,linkcolor=blue,urlcolor=blue,citecolor=blue]{hyperref}
\usepackage[nameinlink,noabbrev]{cleveref}

\newtheorem{theorem}{Theorem}[chapter]
\newtheorem{lemma}[theorem]{Lemma}
\newtheorem{proposition}[theorem]{Proposition}
\newtheorem{corollary}[theorem]{Corollary}
\newtheorem{assessment}[theorem]{Assessment}
\newtheorem{firewall}[theorem]{Claim firewall}
\newtheorem{openproblem}[theorem]{Open problem}
\newtheorem{record}[theorem]{Record}
\newtheorem{programme}[theorem]{Programme}
\newtheorem{audit}[theorem]{Audit}
\newtheorem{correction}[theorem]{Correction}
\newtheorem{warning}[theorem]{Warning}
\newtheorem{decision}[theorem]{Decision}
\newtheorem{scope}[theorem]{Scope}
\theoremstyle{definition}
\newtheorem{definition}[theorem]{Definition}
\theoremstyle{remark}
\newtheorem{remark}[theorem]{Remark}

\newcommand{\MGclosed}{\textsc{closed}}
\newcommand{\MGopen}{\textsc{open}}
\newcommand{\MGpartial}{\textsc{partial}}
\newcommand{\Tr}{\operatorname{Tr}}
\newcommand{\dist}{\operatorname{dist}}
\newcommand{\osc}{\operatorname{osc}}

\title{Renewal Yang--Mills Mass-Gap Research Notes\\
Eighth Series, Version V1}
\author{Compact proof-indexed continuation after archive f7}
\date{8 September 2026}

\begin{document}
\maketitle
\tableofcontents

\chapter{Context, cumulative theorem map, and the live frontier}
\label{ch:e8v1-context}

\section{What this programme is doing}
\label{sec:e8v1-purpose}

The objective is a rigorous construction of four-dimensional quantum
Yang--Mills theory for a compact nonabelian gauge group, together with
Osterwalder--Schrader reconstruction and a proof that the physical
Hamiltonian has a strictly positive gap above its vacuum.  The method
under development is constructive and multiscale.  At finite cutoff it
seeks exact gauge- and reflection-compatible block pushforwards,
separates zero and slow modes before Gaussian integration, proves
conditional cluster expansions for the nonlinear remainder, transports
all principal determinant and covariance rows without double counting,
and finally enters a terminal mixing regime whose Euclidean clustering
implies an OS spectral gap.

This V1 is a size-control rollover, not a mathematical restart.  The
complete seventh series through V100 is frozen as archive f7.  The
present file summarizes the proved results that remain live, records the
unclosed interfaces, and continues with a new proof at the compact-link
localization boundary.

\begin{firewall}[No full mass-gap claim]
\label{fire:e8v1-no-gap}
The full Yang--Mills mass gap is not proved.  In particular, the notes
do not yet construct a regulator-compatible nonperturbative running
coupling, prove the complete compact-link localization map in the
coarse interaction norm, close every chart and bad-block row through
the ultraviolet trajectory, establish regulator independence and the
continuum OS measure, or transfer the terminal gap to the continuum
Hamiltonian.
\end{firewall}

\begin{record}[Archive and version protocol]
\label{rec:e8v1-protocol}
The immediate archive is
\[
 \texttt{Renewal\_Yang\_Mills\_Mass\_Gap\_Research\_Notes\_V100\_f7.tex}.
\]
After correction of a missing theorem-environment declaration, it has
\[
 46\,170\ \hbox{counted source lines},\qquad
 1\,649\,875\ \hbox{bytes},
\]
and SHA--256
\[
 \texttt{816BF7CEA62C101B9826EE351952E54DC6D90C33F9D5911EF1156E1BF01B0723}.
                                                        \label{eq:e8v1-f7-hash}
\]
Archives f1--f6 remain unchanged.  Each successor copies the complete
active source, appends one proof note, passes exact-prefix,
case-sensitive label/reference, environment-stack, brace, and
terminator checks, and replaces its active predecessor.  Every fifth
version audits the newest active SM--Einstein and Navier--Stokes notes.
At V100 the file is frozen under the next \(\texttt{\_fX}\) suffix and
a new contextual V1 begins.
\end{record}

\section{Major mathematical results obtained before archive f7}
\label{sec:e8v1-legacy}

\begin{record}[Compact-group, forest, and conditional foundations]
\label{rec:e8v1-foundations}
Archives f1--f6 prove finite-cutoff Peter--Weyl and compact heat-kernel
identities, exact Duhamel/interpolation formulae, endpoint-factorizing
BKAR representations, rooted-tree and bounded-excess graph estimates,
and strong one-link connected bounds on their declared parameter
ranges.  They also prove the moving-conditional score identity
\[
 X(I-E)f
 =
 (I-E)Xf-\operatorname{Cov}
 \bigl(f,X\log\rho_{\rm cond}\bigr),                 \label{eq:e8v1-score-legacy}
\]
the exact ground-transform projection ledger, and a projected smoothing
estimate which keeps the score covariance as an owned term.
\end{record}

\begin{record}[Negative results which permanently constrain the route]
\label{rec:e8v1-no-go}
The archive chain rigorously rules out several shortcuts:
\begin{enumerate}[label=\textup{(\roman*)},leftmargin=3.5em]
\item small unconditional bad-set probability does not imply a uniform
      conditional influence or restricted operator norm;
\item conditional Fisher information does not imply an
      essential-supremum Dobrushin row;
\item exterior staples can create competing one-link centres even at
      weak coupling;
\item an isolated tree-gauge change is a transport between
      pointwise-fixed fibres, not an internal symmetry of one fibre;
\item reflection positivity or static correlation decay alone does
      not identify a finite-volume Hamiltonian gap; and
\item massless, toron, gauge, and retained infrared modes cannot be
      placed in an almost-gapped determinant without destroying
      uniform inverse locality.
\end{enumerate}
\end{record}

\section{Major seventh-series results}
\label{sec:e8v1-seventh}

\begin{record}[Boundary closure, conditional blocks, and spectral
geometry: V1--V35]
\label{rec:e8v1-v1-v35}
The first part of archive f7:
\begin{enumerate}[label=\textup{(\roman*)},leftmargin=3.5em]
\item identifies finite-action boundary closure and normalization
      obstructions, constructs oriented shell ownership, and separates
      Gaussian transverse cancellation from genuine infrared leakage;
\item derives one- and two-conditioning score decompositions, the exact
      moving-fibre ground, tree-gauge product measure, owner-fibre gap
      mechanisms, buffered Poincare gluing, and the block heat-bath
      influence Gram matrix;
\item builds a finite overlapping collar cover and computes the flat
      collar Gaussian gap, density-deformation stability of conditional
      projections, and intrinsic projection spectral flow;
\item proves that a naive global curvature frame has an infrared
      obstruction, then constructs bounded-clock dyadic alternatives
      and records the incompatibility between some clocks and
      cut-gauge block symmetries;
\item constructs a gauge-projected transfer kernel with a dense
      spin-network core and a volume-uniform strong-coupling correlator
      expansion on its stated domain; and
\item gives exact blocking identities, a Gaussian fine/complement
      split, finite-range cochain lifting, an exact Haar-block free
      semigroup, and a buffered analytic spectral flattening.
\end{enumerate}
Each item is finite-cutoff or conditional on the interfaces stated in
its archived claim firewall.
\end{record}

\begin{record}[Weak-field operators, Feshbach rows, and local
assembly: V36--V60]
\label{rec:e8v1-v36-v60}
The middle operator block:
\begin{enumerate}[label=\textup{(\roman*)},leftmargin=3.5em]
\item derives a three-channel transfer recursion and a physical-time
      pivot, then proves quasi-local incidence on fixed physical slabs;
\item performs high-energy Feshbach reduction with an explicit
      self-energy ledger, a compact-link/Gaussian oscillator split,
      gauge-invariant weak-field IMS localization, and block-owned
      residual elimination;
\item isolates the cross-block Gaussian obstruction and proves nested
      Schur-complement consistency and a fixed-block Wilson
      relative-form estimate;
\item assembles normalized connected slabs, proves buffered
      Combes--Thomas locality, and replaces an invalid zero-radius
      estimate by a bounded-window construction;
\item derives normalized Gaussian mixing, local principal-angle and
      high-Gaussian marginal floors, and the necessary finite-rank
      locality correction; and
\item constructs path-resolved trees for exponential plus finite-rank
      edges, a Gaussian forest identity, spacetime Duhamel trees,
      occupation-window algebra bounds, quasi-free Weyl cumulants, and
      degree-weighted Weyl tree summation.
\end{enumerate}
\end{record}

\begin{record}[Compact charts, rare polymers, and reflection:
V61--V80]
\label{rec:e8v1-v61-v80}
This block proves coefficient stabilization through the occupation
cutoff, represents normalized chart exit as a rare bad-block polymer,
and obtains a glued-Gaussian replica-defect tree bound.  It records an
exact compact-link atlas obstruction and then proves buffered
conditional domination, coupled bad-polymer assembly, conditional
Wilson-Hessian kernel classification, multi-source owner domination,
and a reflection-positive principal/defect decomposition.  It further
derives:
\begin{enumerate}[label=\textup{(\roman*)},leftmargin=3.5em]
\item transition-atlas range constraints and compatible slow
      projections;
\item Kogut--Susskind scaling of projector and chart responses;
\item normalized labelled \(L^2\) assembly of mixed good/bad polymers
      and replica-local density transfer;
\item cutoff-explicit compact-group derivative moments and a
      generator-complete transition-vertex ledger;
\item symmetric form-resolvent resummation and weighted locality;
\item a no-go theorem for spectator-based many-body weak-field
      control, followed by a local form repair; and
\item saturated bad components, one-insertion boundary forests, and
      exact Gauss-sector cancellation of a spurious representation
      multiplicity.
\end{enumerate}
\end{record}

\begin{record}[Terminal criteria and the slow-carrier correction:
V81--V91]
\label{rec:e8v1-v81-v91}
The terminal block separates sufficient theorems from ultraviolet
entry.  It proves:
\begin{enumerate}[label=\textup{(\roman*)},leftmargin=3.5em]
\item a strong-coupling mass-gap theorem from link-conditional
      contraction on a complete compact interaction;
\item a Dobrushin terminal ball and a corrected retained/eliminated
      Schur-path influence recursion;
\item an averaged rare-defect selected-path criterion which, together
      with reflection positivity and a fixed physical time step,
      yields a positive OS gap;
\item a bridge-resolvent theorem showing how sufficiently small
      saturated exceptional components preserve an already gapped
      regular kernel;
\item the exact law-of-total-covariance obstruction: fast conditional
      decay does not control covariance carried entirely by slow
      labels;
\item compact-group character-channel transfer and a volume-free
      defect-anchored flux row; and
\item bounded-multiplier Euclideanization sufficient to import
      finite-cutoff operator estimates into an exact compact-link
      measure.
\end{enumerate}
The ultraviolet trajectory has not yet been shown to enter any of the
terminal sufficient domains.
\end{record}

\begin{record}[Exact coarse action and conditional cumulant trees:
V92--V94]
\label{rec:e8v1-v92-v94}
For one exact block map, the coarse density is the conditional
partition function and the coarse action is its negative logarithm.
Locally uniform absolute convergence permits exact regrouping by
retained footprint.  A typed retained/eliminated selected-tree
majorant gives
\[
 {\cal O}_\mu(\overline\Phi)
 \leq
 D_{\rm dec}
 \left(\rho+\frac{rs}{1-\epsilon}\right).            \label{eq:e8v1-coarse-recursion}
\]
A source-dependent conditional KP expansion implies the required
all-order cumulant tree card, with rooted and marked factors
\((1-q_{\rm br})^{-1}\) and \((1-q_{\rm br})^{-2}\).
Pair covariance alone is proved insufficient by an explicit parity
counterexample.  Finally, spanning-tree/chord variables give an exact
product-Haar Jacobian; disjoint chord or triangular path coordinates
produce a product fibre, while overlapping paths require a retained
shared chord or an explicit conditional kernel.
\end{record}

\begin{record}[Near-flat Wilson/Haar and determinant rows: V95--V97]
\label{rec:e8v1-v95-v97}
Around a stationary conditional centre, V95 rescales the eliminated
field by the positive Hessian and proves a normalized Gaussian
Radon--Nikodym estimate.  On
\(M=(\log\beta)^\chi\), \(g=\beta^{-1/2}\), both the Wilson and Haar
good-set activities vanish; the saturated bad activity vanishes under
the previously stated \(\chi>1\) tail condition.  The Gaussian
normalization is not part of that small activity: it contributes the
owned coarse row \(\frac12\log\det L_y\).

V96 proves
\[
 D_x\log\det L=\Tr(L^{-1}D_xL),                      \label{eq:e8v1-logdet-first}
\]
\[
 D_yD_x\log\det L
 =
 \Tr\!\left(
 L^{-1}D_yD_xL-L^{-1}D_yL\,L^{-1}D_xL
 \right),                                            \label{eq:e8v1-logdet-second}
\]
and exponential determinant locality for a finite-range uniformly
gapped eliminated Hessian.  An unremoved four-dimensional massless
mode gives a logarithmically divergent absolute response row.

V97 computes the actual flat Wilson Hessian
\[
 D^2w(0)=d_1^*d_1,                                  \label{eq:e8v1-Wilson-Hessian}
\]
retains gauge gradients and flat cohomological/toron modes, proves a
positive transverse floor on a fixed block, constructs the physical
stationary centre \(\widehat x(u)\), and obtains the exact composition
\[
 L_{g,y}={\cal L}(gy).                               \label{eq:e8v1-Lgy}
\]
Therefore its first two retained derivatives are \(O(g)\) and
\(O(g^2)\), and the fixed-dictionary determinant response is
\(O(\beta^{-1})\).
\end{record}

\begin{record}[Cofinal fixed-factor geometry and calibrated
contraction: V98--V100]
\label{rec:e8v1-v98-v100}
On a periodic growing fibre, the transverse floor is exactly
\[
 \lambda_\ell=4\sin^2(\pi/\ell)\asymp\ell^{-2},       \label{eq:e8v1-growing-floor}
\]
so the present absolute determinant method cannot use a one-shot
physical block with \(\ell\to\infty\).  Exact associativity of
pushforwards permits instead a growing number of fixed-factor maps,
each with fixed elementary geometry.

For one exact coarse action,
\[
 D{\cal R}_{S}[h](y)=\mathbb E_{S,y}h,\qquad
 D^2{\cal R}_{S}[h,k](y)
 =-\operatorname{Cov}_{S,y}(h,k).                   \label{eq:e8v1-R-derivatives}
\]
Retained pullbacks obey
\({\cal R}(S+q\circ\pi)={\cal R}(S)+q\), proving that
the complete RG differential has a unit sector and cannot be strictly
contractive.  On a calibrated stable slice, a Gaussian contraction
\(\kappa_{\rm G}\) plus a complete nonlinear correction \(\eta\)
does contract when
\[
 \kappa_{\rm G}+\eta<1.                              \label{eq:e8v1-margin}
\]

V100 constructs the minimum-energy Gaussian interpolation
\[
 J=L^{-1}Q^*(QL^{-1}Q^*)^{-1},                       \label{eq:e8v1-J}
\]
proves the exact energy-orthogonal retained/detail Gaussian
disintegration, and shows zero Gaussian transfer on the conditionally
centred detail space.  Canonical dilation in four dimensions leaves
\(\Tr F^2\) and \(\Tr F\widetilde F\) marginal and contracts the
declared local curvature jets of dimension at least six by
\[
 \kappa_{\rm G}\leq b^{-2}.                          \label{eq:e8v1-kappaG}
\]
The unresolved bridge is to place complete compact-link interactions,
not only formal curvature jets, in that graded stable norm.
\end{record}

\section{The live dependency chain}
\label{sec:e8v1-frontier}

\begin{assessment}[Smallest current frontier]
\label{ass:e8v1-frontier}
The next proof obligations, in dependency order, are:
\begin{enumerate}[label=\textup{(\arabic*)},leftmargin=3.5em]
\item a block-local gauge-invariant expansion of compact-link
      interactions into explicitly owned marginal curvature terms and
      a remainder small in the complete coarse oscillation norm;
\item a regulator-compatible marginal trajectory for the coefficient
      of \(\Tr F^2\), including wave-function and topological-sector
      normalization;
\item a uniform conditional CFTC estimate for all good charts and a
      normalized saturated expansion for bad components;
\item iteration of the exact fixed-factor cocycle into a terminal
      Dobrushin, bridge, or strong-coupling domain;
\item regulator-independent reflection-positive continuum measure and
      OS reconstruction; and
\item transfer of the terminal clustering rate to a strictly positive
      continuum Hamiltonian gap.
\end{enumerate}
The present note attacks item (1).  It first proves a noncircular
compact-link-to-loop-polynomial Taylor bridge.  Conversion of the loop
polynomials to local curvature jets is deliberately left for the next
step.
\end{assessment}

\chapter{A gauge-invariant loop-coordinate Taylor bridge}
\label{ch:e8v1-loop}

\section{Tree--chord descent}
\label{sec:e8v1-descent}

Let \(B=(V,E)\) be a fixed connected finite graph, choose a root
\(o\), a spanning tree \(T\), and let \(C=E\setminus T\) be the chord
set.  Let \(G=\mathrm{SU}(N)\).  The V94 tree-coordinate map writes
the link configuration as root-relative gauge frames
\((h_v)_{v\neq o}\) and chord holonomies
\((z_c)_{c\in C}\), with exact product Haar measure.

\begin{lemma}[Gauge-invariant descent to chord holonomies]
\label{lem:e8v1-descent}
Let \(\Phi:G^E\to\mathbb R\) be invariant under vertex gauge
transformations.  In the tree--chord coordinates there is a unique
function
\[
 \phi:G^C\longrightarrow\mathbb R                  \label{eq:e8v1-phi}
\]
such that
\[
 \Phi(U)=\phi((z_c)_{c\in C}).                       \label{eq:e8v1-descent}
\]
It is invariant under simultaneous residual conjugation:
\[
 \phi((kz_ck^{-1})_{c\in C})
 =
 \phi((z_c)_{c\in C})\qquad(k\in G).                \label{eq:e8v1-conjugation}
\]
Smoothness or real analyticity of \(\Phi\) in this chart passes to
\(\phi\).
\end{lemma}

\begin{proof}
Tree gauge uses vertex transformations to set every tree link to the
identity.  Two configurations in tree gauge lie on the same full gauge
orbit precisely when their chord tuples differ by the residual root
conjugation.  Gauge invariance makes \(\Phi\) independent of the
discarded frames and constant on this residual orbit, defining
\(\phi\).  Surjectivity of the coordinate map gives uniqueness.
Equation \eqref{eq:e8v1-conjugation} is the residual root action.
The coordinate map and its inverse are smooth, and analytic inside
the chosen exponential chart, proving the last assertion.
\end{proof}

\begin{remark}[Support hull]
\label{rem:e8v1-hull}
If \(\Phi\) initially uses a link set \(Z\), tree gauge may enlarge its
coordinate footprint to the union \(H_T(Z)\) of the tree paths joining
the endpoints of links in \(Z\) to the root used for that connected
component.  This enlargement is bounded for a fixed block dictionary.
It must be charged explicitly when supports are summed.
\end{remark}

\section{Analytic physical-field expansion}
\label{sec:e8v1-Taylor}

Write \(r=|C|\), \(\mathfrak g=\mathfrak{su}(N)\), and use the product
norm on \(\mathfrak g^r\).  Define
\[
 \widetilde\phi(Z_1,\ldots,Z_r)
 :=
 \phi(e^{Z_1},\ldots,e^{Z_r})                       \label{eq:e8v1-phitilde}
\]
on a ball \(\|Z\|\leq r_0\) inside the product injectivity chart.

\begin{lemma}[Vanishing invariant linear term]
\label{lem:e8v1-linear}
If \(\phi\) satisfies \eqref{eq:e8v1-conjugation}, then
\[
 D\widetilde\phi(0)=0.                               \label{eq:e8v1-linear-zero}
\]
\end{lemma}

\begin{proof}
For each chord \(c\), the partial derivative at zero is a linear
functional \(\ell_c\) on \(\mathfrak{su}(N)\).  Simultaneous
conjugation invariance implies
\(\ell_c(\operatorname{Ad}_kX)=\ell_c(X)\) for every \(k\).
Differentiating in \(k=e^{tY}\) gives
\(\ell_c([Y,X])=0\).  Since
\([\mathfrak{su}(N),\mathfrak{su}(N)]=\mathfrak{su}(N)\),
\(\ell_c=0\).
\end{proof}

Assume \(\widetilde\phi(0)=0\), which only subtracts a vacuum constant.
For \(g>0\), define the Wilson-scaled block potential
\[
 {\cal W}_g(X):=g^{-2}\widetilde\phi(gX).             \label{eq:e8v1-Wg}
\]
For an integer \(p\geq2\), set
\[
 {\cal P}_{p,g}(X)
 :=
 \sum_{n=2}^{p}
 \frac{g^{n-2}}{n!}
 D^n\widetilde\phi(0)[X^{\otimes n}],
 \qquad
 {\cal E}_{p,g}:={\cal W}_g-{\cal P}_{p,g}.           \label{eq:e8v1-polynomial}
\]

\begin{theorem}[Uniform differentiated Taylor remainder]
\label{thm:e8v1-Taylor}
Suppose \(\widetilde\phi\in C^{p+1}\) on
\(\|Z\|\leq r_0\), and put
\[
 M_{p+1}:=
 \sup_{\|Z\|\leq r_0}\|D^{p+1}\widetilde\phi(Z)\|.
                                                        \label{eq:e8v1-M}
\]
If \(\|X\|\leq R\) and \(gR\leq r_0\), then for
\(k=0,1,2\),
\[
 \|D^k{\cal E}_{p,g}(X)\|
 \leq
 \frac{M_{p+1}}{(p+1-k)!}\,
 g^{p-1}R^{p+1-k}.                                   \label{eq:e8v1-Taylor-bound}
\]
\end{theorem}

\begin{proof}
For \(k\leq2\),
\[
 D^k{\cal W}_g(X)
 =
 g^{k-2}D^k\widetilde\phi(gX).                       \label{eq:e8v1-scaled-derivative}
\]
Apply Taylor's theorem with integral remainder to
\(D^k\widetilde\phi\) at zero through degree \(p-k\).
Its remainder norm at \(gX\) is at most
\[
 \frac{M_{p+1}}{(p+1-k)!}\|gX\|^{p+1-k}.
\]
Multiplication by \(g^{k-2}\) gives exactly
\eqref{eq:e8v1-Taylor-bound}.  The polynomial terms are the
\(k\)-th derivatives of \eqref{eq:e8v1-polynomial}.
\end{proof}

\begin{corollary}[Quartic truncation on the logarithmic good set]
\label{cor:e8v1-quartic}
Take \(p=4\), \(g=\beta^{-1/2}\), and
\(R=c(\log\beta)^{\chi/2}\).  Then
\begin{align}
 \sup_{\|X\|\leq R}|{\cal E}_{4,g}(X)|
 &\leq C\beta^{-3/2}(\log\beta)^{5\chi/2},
                                                        \label{eq:e8v1-E0}\\
 \sup_{\|X\|\leq R}\|D{\cal E}_{4,g}(X)\|
 &\leq C\beta^{-3/2}(\log\beta)^{2\chi},
                                                        \label{eq:e8v1-E1}\\
 \sup_{\|X\|\leq R}\|D^2{\cal E}_{4,g}(X)\|
 &\leq C\beta^{-3/2}(\log\beta)^{3\chi/2}.
                                                        \label{eq:e8v1-E2}
\end{align}
All three quantities tend to zero for fixed \(\chi\).
\end{corollary}

\begin{proof}
Insert \(p=4\), \(g=\beta^{-1/2}\), and the stated \(R\) into
\eqref{eq:e8v1-Taylor-bound}.  Every fixed logarithmic power is
dominated by \(\beta^{3/2}\).
\end{proof}

\section{Invariant Taylor coefficients}
\label{sec:e8v1-invariants}

\begin{lemma}[Residual-gauge invariance of every Taylor tensor]
\label{lem:e8v1-tensors}
For every \(n\geq0\),
\[
 D^n\widetilde\phi(0)
 [\operatorname{Ad}_kX_1,\ldots,\operatorname{Ad}_kX_n]
 =
 D^n\widetilde\phi(0)[X_1,\ldots,X_n],               \label{eq:e8v1-tensor-invariance}
\]
where \(\operatorname{Ad}_k\) acts diagonally on
\(\mathfrak g^r\).
\end{lemma}

\begin{proof}
Exponential equivariance gives
\[
 \widetilde\phi(\operatorname{Ad}_kZ_1,\ldots,
 \operatorname{Ad}_kZ_r)=\widetilde\phi(Z_1,\ldots,Z_r).
\]
Differentiate \(n\) times at zero.
\end{proof}

\begin{proposition}[Form of the quadratic chord polynomial]
\label{prop:e8v1-quadratic}
There is a real symmetric matrix
\((a_{cd})_{c,d\in C}\) such that
\[
 \frac12D^2\widetilde\phi(0)[X,X]
 =
 \frac12\sum_{c,d\in C}
 a_{cd}\langle X_c,X_d\rangle_{\mathfrak g}.         \label{eq:e8v1-quadratic}
\]
\end{proposition}

\begin{proof}
For fixed \(c,d\), polarize the \(c,d\) component of the invariant
quadratic tensor in \Cref{lem:e8v1-tensors}.  It is an
\(\operatorname{Ad}G\)-invariant bilinear form on the simple real Lie
algebra \(\mathfrak{su}(N)\), hence a scalar multiple of its invariant
inner product.  Symmetry of the Hessian gives \(a_{cd}=a_{dc}\).
\end{proof}

\begin{firewall}[Chord invariants are not yet continuum curvature
jets]
\label{fire:e8v1-chord}
The polynomial \({\cal P}_{4,g}\) is an exact
residual-gauge-invariant polynomial in logarithms of fundamental chord
holonomies.  A chord loop can enclose several plaquettes and its
logarithm contains BCH commutators and path-ordering terms.  Therefore
\eqref{eq:e8v1-quadratic} cannot yet be identified term by term with
\(\Tr F^2\), \(\Tr F\widetilde F\), or dimension-six local curvature
jets.  That identification requires a separate plaquette-factorization
and remainder theorem.
\end{firewall}

\section{Complete oscillation control of the Taylor tail}
\label{sec:e8v1-oscillation}

For a real function \(F\) on a product good set, let
\(\osc_cF\) be its oscillation when only coordinate \(X_c\) varies,
and let \(\osc_c\osc_dF\) be the corresponding mixed rectangular
oscillation.

\begin{lemma}[Derivative-to-oscillation conversion]
\label{lem:e8v1-oscillation}
On the product ball \(\|X_c\|\leq R\),
\begin{align}
 \osc_cF
 &\leq2R\sup\|D_cF\|,                                \label{eq:e8v1-osc1}\\
 \osc_c\osc_dF
 &\leq4R^2\sup\|D_dD_cF\|.                           \label{eq:e8v1-osc2}
\end{align}
\end{lemma}

\begin{proof}
Join two values of \(X_c\) by the line segment in its convex ball and
apply the fundamental theorem of calculus, proving
\eqref{eq:e8v1-osc1}.  Apply the same argument successively in
coordinates \(c\) and \(d\) to the rectangular difference.
\end{proof}

\begin{corollary}[Taylor-tail oscillation row]
\label{cor:e8v1-tail-row}
Under \Cref{thm:e8v1-Taylor},
\[
 \osc_c{\cal E}_{p,g}
 \leq
 \frac{2M_{p+1}}{p!}g^{p-1}R^{p+1},                 \label{eq:e8v1-tail-one}
\]
and
\[
 \osc_c\osc_d{\cal E}_{p,g}
 \leq
 \frac{4M_{p+1}}{(p-1)!}g^{p-1}R^{p+1}.             \label{eq:e8v1-tail-two}
\]
\end{corollary}

\begin{proof}
Combine
\Cref{thm:e8v1-Taylor,lem:e8v1-oscillation} for
\(k=1,2\).
\end{proof}

Consider now a family of block potentials \(\Phi_Z\).  Let
\({\cal H}(Z)\) be its charged tree--chord hull and assume the fixed
dictionary constants
\[
 D_*:=
 \sup_c\sum_{Z:\,c\in{\cal H}(Z)}
 e^{\mu\operatorname{diam}{\cal H}(Z)}
 (|{\cal H}(Z)|-1)
 <\infty,                                             \label{eq:e8v1-incidence}
\]
and
\[
 M_*:=\sup_Z M_{p+1,Z}<\infty.                        \label{eq:e8v1-uniform-M}
\]

\begin{theorem}[Complete fixed-dictionary Taylor-tail bound]
\label{thm:e8v1-complete-tail}
Let \({\cal E}_{p,g,Z}\) be the Taylor tail of \(\Phi_Z\).  Its
complete weighted oscillation row satisfies
\[
 \sup_c
 \sum_{Z:\,c\in{\cal H}(Z)}
 e^{\mu\operatorname{diam}{\cal H}(Z)}
 (|{\cal H}(Z)|-1)
 \osc_c{\cal E}_{p,g,Z}
 \leq
 \frac{2D_*M_*}{p!}g^{p-1}R^{p+1}.                  \label{eq:e8v1-complete-tail}
\]
The analogous marked mixed row is bounded by the right side with
\(2/p!\) replaced by \(4/(p-1)!\), times the fixed number of possible
second marks.  For \(p=4\) and the logarithmic good-set schedule, both
rows tend to zero.
\end{theorem}

\begin{proof}
The support of each Taylor coefficient and its remainder is contained
in the charged hull because differentiation does not create dependence
on a missing coordinate.  Insert
\eqref{eq:e8v1-tail-one} into the left side and use
\eqref{eq:e8v1-incidence}--\eqref{eq:e8v1-uniform-M}.
For two marks use \eqref{eq:e8v1-tail-two} and the uniformly finite
chord count of the dictionary.  Vanishing follows from
\Cref{cor:e8v1-quartic}.
\end{proof}

\begin{corollary}[A proved part of the V100 localization bridge]
\label{cor:e8v1-GCDC}
On a uniformly finite weak-field tree--chord dictionary, every
gauge-invariant \(C^5\) Wilson-scaled block interaction splits as
\[
 {\cal W}_g
 =
 {\cal P}_{4,g}+{\cal E}_{4,g},                      \label{eq:e8v1-split}
\]
where \({\cal P}_{4,g}\) is a residual-gauge-invariant chord
polynomial and \({\cal E}_{4,g}\) has a vanishing complete good-set
oscillation and mixed-oscillation row.  Thus the Taylor-tail part of
GCDC (G6) is closed for a fixed dictionary.
\end{corollary}

\begin{proof}
Gauge invariance and the polynomial form are
\Cref{lem:e8v1-descent,lem:e8v1-tensors}.  The complete tail is
\Cref{thm:e8v1-complete-tail}.
\end{proof}

\begin{firewall}[What remains in GCDC (G6)]
\label{fire:e8v1-GCDC}
\Cref{cor:e8v1-GCDC} does not convert
\({\cal P}_{4,g}\) into a local curvature-jet basis, control the
growth of tree hulls outside a fixed dictionary, identify the exact
marginal projection inside the complete compact-link norm, or include
global transition charts and bad blocks.  It closes only the analytic
Taylor-tail half of the localization bridge.
\end{firewall}

\section{Status and next acquisition}
\label{sec:e8v1-status}

\begin{theorem}[What eighth-series V1 proves]
\label{thm:e8v1-result}
Besides the compact theorem map, V1 proves that a gauge-invariant
finite-block interaction descends exactly to residual-conjugation
invariant chord holonomies.  In physical exponential coordinates its
Wilson-scaled quartic Taylor remainder has first and second derivative
and complete oscillation rows of order
\[
 g^3R^5,
\]
up to fixed dictionary constants, and therefore vanishes for
\(g=\beta^{-1/2}\), \(R=O((\log\beta)^{\chi/2})\).
Every retained Taylor tensor is gauge invariant, and the quadratic
tensor is an explicit matrix of invariant Lie-algebra inner products.
\end{theorem}

\begin{proof}
Use
\Cref{lem:e8v1-descent,lem:e8v1-linear,
thm:e8v1-Taylor,prop:e8v1-quadratic,
thm:e8v1-complete-tail,cor:e8v1-GCDC}.
\end{proof}

\begin{assessment}[Progress at eighth-series V1]
\label{ass:e8v1-status}
The compact-link localization problem has been split without
circularity.  The analytic tail is now genuinely small in the fixed
block complete norm.  The non-small finite polynomial is retained
exactly.  The next upstream question is algebraic and geometric:
express fundamental chord logarithms through ordered plaquette
logarithms, quantify the BCH remainder, and separate the quadratic
plaquette form into the Yang--Mills and topological marginal owners
before canonical dilation is invoked.
\end{assessment}

\begin{decision}[V2 acquisition order]
\label{dec:e8v1-V2}
V2 will choose a plaquette filling for every fundamental chord loop in
a contractible fixed block.  It will prove an ordered nonabelian
Stokes/BCH expansion
\[
 X_c=\sum_{p\in\Sigma_c}\epsilon_{cp}F_p+
 \text{commutator remainder},
\]
with first and second retained derivatives and a complete
fixed-dictionary oscillation bound.  It will then substitute this
formula into \({\cal P}_{4,g}\), keeping every quadratic marginal
coefficient explicit.
\end{decision}

\begin{firewall}[Claim boundary after eighth-series V1]
\label{fire:e8v1-claim}
V1 does not complete the chord-to-curvature conversion, construct the
running gauge coupling, prove the global nonlinear CCDC/GCDC margin,
close all charts and saturated bad components, establish NCGC/CFTC
through the full ultraviolet trajectory, prove terminal entry, BSCTC,
CSMC, MIRC, or ARDPC for that trajectory, obtain regulator-uniform
SATC or full RL4C, construct the interacting endpoint and continuum OS
measure, prove regulator independence, or prove the full Yang--Mills
mass gap.
\end{firewall}

\paragraph{Archive-parent record.}
Eighth-series V1 begins after the corrected archive
\texttt{Renewal\_Yang\_Mills\_Mass\_Gap\_Research\_Notes\_V100\_f7.tex},
whose checkpoint is \eqref{eq:e8v1-f7-hash}.  The next compulsory
companion audit is V5.

\clearpage
% =============================================================================
% BEGIN APPENDED EIGHTH-SERIES V2 MATERIAL
% =============================================================================

\chapter{Shellable lattice Stokes fillings and the BCH curvature
polynomial}
\label{ch:e8v2}

\section{Purpose}
\label{sec:e8v2-purpose}

V1 separated a gauge-invariant block interaction into a quartic
polynomial in logarithms of fundamental chord holonomies and a small
analytic tail.  A chord logarithm is not itself a local curvature.
This note chooses explicit two-cell fillings, writes each chord
holonomy as an exact ordered product of root-based plaquette
holonomies, and expands the logarithm through cubic Lie degree.  Cubic
order is necessary: its substitution into the quadratic chord form
contributes at quartic order and therefore cannot be placed in the
small tail.

\section{Shellable fillings and an exact lattice Stokes word}
\label{sec:e8v2-Stokes}

Let \(B=(V,E,P)\) be a finite connected oriented two-complex, with
root \(o\) and spanning tree \(T\).  Contracting \(T\) leaves the chord
generators \(c\in C=E\setminus T\).  For a plaquette \(p\), let
\(r_p\) be its oriented boundary word based at a selected boundary
vertex.

\begin{definition}[Root-based plaquette word]
\label{def:e8v2-based}
For a path \(s_p\) from \(o\) to the base vertex of \(r_p\), define
\[
 R_p:=s_pr_ps_p^{-1}.                                \label{eq:e8v2-Rp}
\]
After evaluating edge words in a link configuration, \(R_p(U)\) is a
root-based conjugate of the plaquette holonomy.
\end{definition}

\begin{definition}[Shellable filling]
\label{def:e8v2-shelling}
A shellable filling of the fundamental chord loop \(\gamma_c\) is a
finite ordered list
\[
 (p_1,\epsilon_1,s_1),\ldots,(p_m,\epsilon_m,s_m),
 \qquad \epsilon_j\in\{-1,1\},                       \label{eq:e8v2-shelling}
\]
obtained by successively attaching oriented plaquettes along connected
boundary arcs to a disk whose final oriented boundary is
\(\gamma_c\).  The paths \(s_j\) join the root to the attachment arcs.
\end{definition}

\begin{lemma}[Exact nonabelian lattice Stokes identity]
\label{lem:e8v2-Stokes}
For a shellable filling there is an ordering convention, fixed by the
attachments, for which the chord word in the free edge group obeys
\[
 \gamma_c
 =
 R_{p_m}^{\epsilon_m}\cdots R_{p_1}^{\epsilon_1}.
                                                        \label{eq:e8v2-word}
\]
Consequently every link configuration satisfies
\[
 z_c(U)
 =
 P_{c,m}(U)\cdots P_{c,1}(U),                        \label{eq:e8v2-holonomy}
\]
where each \(P_{c,j}\) is a root-based conjugate of
\(\Omega_{p_j}(U)^{\epsilon_j}\).
\end{lemma}

\begin{proof}
Begin with the boundary word of the first oriented plaquette.  When
the next plaquette is attached along a connected boundary arc, multiply
by its boundary word based at the initial point of that arc and
conjugated to the common root.  The shared arc occurs once in each
orientation and cancels as adjacent inverse subwords.  The resulting
reduced word is the boundary of the enlarged disk.  Induction over the
shelling leaves the final boundary word \(\gamma_c\), proving
\eqref{eq:e8v2-word}.  Evaluation of a free edge word in \(G\) is a
homomorphism, which gives \eqref{eq:e8v2-holonomy}.
\end{proof}

\begin{firewall}[Topology is an input]
\label{fire:e8v2-topology}
A general block need not give every fundamental chord loop a shellable
disk.  Noncontractible cycles are precisely the flat/cohomological
directions retained in V97.  V2 applies only to chord combinations
with a declared shellable filling; it does not turn torons or boundary
transport into local plaquette curvature.
\end{firewall}

\section{Root-based plaquette logarithms}
\label{sec:e8v2-logs}

Choose a product injectivity neighbourhood of the identity.  For a
filled chord \(c\), write
\[
 P_{c,j}=\exp(gF_{c,j}),\qquad
 z_c=\exp(gX_c),                                     \label{eq:e8v2-logs}
\]
where \(F_{c,j},X_c\in\mathfrak{su}(N)\).  Orientation is already
included in \(F_{c,j}\); in particular,
\(\log(P^{-1})=-\log P\) in this chart.  Root gauge transformations
act simultaneously by the adjoint representation on all these
logarithms.

Fix an \(\operatorname{Ad}\)-invariant norm satisfying
\[
 \|[A,B]\|\leq c_{\rm Lie}\|A\|\|B\|.                \label{eq:e8v2-bracket}
\]
For \(m\leq m_*\), define
\[
 {\cal B}_{g,m}(F_1,\ldots,F_m)
 :=
 \frac1g\log\bigl(e^{gF_m}\cdots e^{gF_1}\bigr).
                                                        \label{eq:e8v2-BCH-map}
\]

\begin{lemma}[Analyticity with uniform finite-dictionary radius]
\label{lem:e8v2-analytic}
For each finite \(m_*\) there is \(r_{\rm BCH}>0\) such that all maps
\({\cal B}_{g,m}\), \(m\leq m_*\), are real analytic whenever
\[
 g\sum_{j=1}^m\|F_j\|<r_{\rm BCH}.                   \label{eq:e8v2-radius}
\]
Their derivatives through any fixed order are bounded uniformly over
\(m\leq m_*\) on every smaller closed ball.
\end{lemma}

\begin{proof}
Matrix multiplication and the exponential are entire.  The principal
matrix logarithm is analytic on a fixed neighbourhood of \(I\).
Continuity of the finite products gives a common input ball for the
finite set \(m=1,\ldots,m_*\).  Compactness of a smaller closed ball
and continuity of each derivative give uniform bounds.
\end{proof}

Let
\[
 A_m(F):=\sum_{j=1}^mF_j,\qquad
 B_{2,m}(F):=\frac12\sum_{1\leq i<j\leq m}[F_j,F_i],
                                                        \label{eq:e8v2-AB}
\]
where the bracket order matches
\eqref{eq:e8v2-BCH-map}.  Define \(B_{3,m}\) to be the homogeneous
cubic Taylor coefficient
\[
 B_{3,m}(F):=
 \frac1{3!}
 D^3{\cal L}_m(0)[F,F,F],\qquad
 {\cal L}_m(Z):=\log(e^{Z_m}\cdots e^{Z_1}).          \label{eq:e8v2-B3}
\]
It is a universal finite sum of nested commutators.

\begin{theorem}[BCH expansion through cubic Lie degree]
\label{thm:e8v2-BCH}
For \(m\leq m_*\), \(\sum_j\|F_j\|\leq R\), and
\(gR\leq r<r_{\rm BCH}\),
\[
 {\cal B}_{g,m}(F)
 =
 A_m(F)+gB_{2,m}(F)+g^2B_{3,m}(F)
 +{\cal R}_{4,g,m}(F).                               \label{eq:e8v2-BCH}
\]
For \(k=0,1,2\),
\[
 \|D^k{\cal R}_{4,g,m}(F)\|
 \leq C_{m_*,r}g^3R^{4-k}.                           \label{eq:e8v2-BCH-bound}
\]
All four terms in \eqref{eq:e8v2-BCH} are equivariant under
simultaneous adjoint conjugation.
\end{theorem}

\begin{proof}
Taylor-expand \({\cal L}_m(Z)\) at zero through degree three.
Its linear coefficient is \(A_m\).  Multiplying two exponentials and
comparing the quadratic terms in
\(\exp({\cal L}_m)\) gives \(B_{2,m}\); induction over the number of
factors yields the sum in \eqref{eq:e8v2-AB}.  The cubic coefficient is
\eqref{eq:e8v2-B3}.  The differentiated fourth-order Taylor remainder
on the closed radius-\(r\) ball is bounded by
\(C\|Z\|^{4-k}\).  Substitute \(Z=gF\) and divide by \(g\);
the \(k\)-th derivative in \(F\) contributes \(g^k\), giving
\eqref{eq:e8v2-BCH-bound}.  Equivariance follows from
\(\log(kUk^{-1})=\operatorname{Ad}_k\log U\).
\end{proof}

\begin{corollary}[Filled chord logarithm]
\label{cor:e8v2-chord}
For every filled chord,
\[
 X_c=A_cF+gB_{2,c}(F)+g^2B_{3,c}(F)
 +{\cal R}_{4,c}(F),                                \label{eq:e8v2-chord}
\]
where \(A_cF=\sum_{j=1}^{m_c}F_{c,j}\) and
\[
 \|D^k{\cal R}_{4,c}\|
 \leq Cg^3R^{4-k},\qquad k=0,1,2,                   \label{eq:e8v2-chord-bound}
\]
uniformly over a fixed filling dictionary.
\end{corollary}

\begin{proof}
Apply \Cref{lem:e8v2-Stokes,thm:e8v2-BCH} and take the maximum over
the finite dictionary.
\end{proof}

\section{Substitution into the quartic chord polynomial}
\label{sec:e8v2-substitution}

Write the V1 retained polynomial as
\[
 {\cal P}_{4,g}(X)
 =
 Q_2(X)+gQ_3(X)+g^2Q_4(X),                           \label{eq:e8v2-P4}
\]
where
\[
 Q_n(X):=\frac1{n!}D^n\widetilde\phi(0)
 [X^{\otimes n}],\qquad n=2,3,4.                    \label{eq:e8v2-Qn}
\]
Let \(q_2\) be the symmetric bilinear form with
\(Q_2(X)=q_2(X,X)\).  Collect all root-based plaquette occurrences in
the chosen filling dictionary into a vector \(F\), and write
\[
 A(F)=(A_cF)_c,\quad
 B_2(F)=(B_{2,c}(F))_c,\quad
 B_3(F)=(B_{3,c}(F))_c.                              \label{eq:e8v2-ABC}
\]

\begin{definition}[Retained plaquette polynomials]
\label{def:e8v2-K}
Define
\begin{align}
 K_2(F)
 &:=
 Q_2(A(F)),                                          \label{eq:e8v2-K2}\\
 K_3(F)
 &:=
 2q_2(A(F),B_2(F))+Q_3(A(F)),                        \label{eq:e8v2-K3}\\
 K_4(F)
 &:=
 Q_2(B_2(F))
 +2q_2(A(F),B_3(F))
 +DQ_3(A(F))[B_2(F)]
 +Q_4(A(F)).                                         \label{eq:e8v2-K4}
\end{align}
\end{definition}

\begin{theorem}[Quartic plaquette substitution with a fifth-order
tail]
\label{thm:e8v2-substitution}
Assume the V1 \(C^5\) hypothesis and a fixed filling dictionary.
On \(\|F\|\leq R\), \(gR\leq r\), one has
\[
 {\cal W}_g(X(F))
 =
 K_2(F)+gK_3(F)+g^2K_4(F)+{\cal T}_{5,g}(F),         \label{eq:e8v2-substitution}
\]
where \(K_n\) is homogeneous of degree \(n\), is invariant under
simultaneous root conjugation, and, for \(k=0,1,2\),
\[
 \|D^k{\cal T}_{5,g}(F)\|
 \leq Cg^3R^{5-k}.                                   \label{eq:e8v2-T-bound}
\]
\end{theorem}

\begin{proof}
Insert
\[
 X=A+gB_2+g^2B_3+{\cal R}_4
\]
from \eqref{eq:e8v2-chord} into
\eqref{eq:e8v2-P4}.  Bilinearity gives
\[
 Q_2(X)
 =
 Q_2(A)+2gq_2(A,B_2)
 +g^2\{Q_2(B_2)+2q_2(A,B_3)\}
 +O_{C^2}(g^3R^5).
\]
The cubic polynomial gives
\[
 gQ_3(X)
 =
 gQ_3(A)+g^2DQ_3(A)[B_2]
 +O_{C^2}(g^3R^5),
\]
and the quartic polynomial gives
\[
 g^2Q_4(X)=g^2Q_4(A)+O_{C^2}(g^3R^5).
\]
Here \(O_{C^2}(g^3R^5)\) means that its \(k\)-th derivative is
bounded by \(Cg^3R^{5-k}\), \(k\leq2\).  This follows term by term
from homogeneity, the fixed multilinear coefficient norms, and
\eqref{eq:e8v2-chord-bound}.  The V1 analytic tail has the same bound
by \eqref{eq:e8v1-Taylor-bound} with \(p=4\).  The displayed
coefficients are exactly \eqref{eq:e8v2-K2}--
\eqref{eq:e8v2-K4}.  Adjoint equivariance of \(A,B_2,B_3\) and
invariance of the Taylor tensors from
\Cref{lem:e8v1-tensors,thm:e8v2-BCH} prove invariance.
\end{proof}

\begin{corollary}[Complete substitution-tail row]
\label{cor:e8v2-tail}
For a uniformly finite support and filling dictionary, the complete
weighted oscillation and mixed-oscillation rows of
\({\cal T}_{5,g}\) are at most
\[
 C_{\mu,\rm fill}g^3R^5.                             \label{eq:e8v2-tail}
\]
They tend to zero on the logarithmic good-set schedule.
\end{corollary}

\begin{proof}
Use \eqref{eq:e8v2-T-bound} and the
derivative-to-oscillation conversion
\Cref{lem:e8v1-oscillation}.  The shelling paths and tree hulls have
uniformly finite diameter and incidence, so their exponential weights
enter only \(C_{\mu,\rm fill}\).  Finally apply
\Cref{cor:e8v1-quartic}.
\end{proof}

\section{The quadratic plaquette kernel}
\label{sec:e8v2-quadratic}

Let \({\cal I}\) index the oriented root-based plaquette occurrences
in the filling dictionary.  There is a signed incidence matrix
\(\sigma_{c\alpha}\) such that
\[
 (A(F))_c=\sum_{\alpha\in{\cal I}}\sigma_{c\alpha}F_\alpha.
                                                        \label{eq:e8v2-incidence}
\]

\begin{proposition}[Explicit quadratic pushforward]
\label{prop:e8v2-quadratic}
If the V1 chord quadratic matrix is \(a=(a_{cd})\), then
\[
 K_2(F)
 =
 \frac12
 \sum_{\alpha,\beta\in{\cal I}}
 {\mathsf K}_{\alpha\beta}
 \langle F_\alpha,F_\beta\rangle_{\mathfrak g},
 \qquad
 {\mathsf K}=\sigma^{\mathsf T}a\sigma.              \label{eq:e8v2-Kmatrix}
\]
In particular, \({\mathsf K}\) is real symmetric and positive
semidefinite whenever \(a\) is positive semidefinite.
\end{proposition}

\begin{proof}
Substitute \eqref{eq:e8v2-incidence} into the quadratic formula
\eqref{eq:e8v1-quadratic} and collect coefficients.  Symmetry is
preserved by congruence.  If \(a\geq0\), then
\(v^{\mathsf T}{\mathsf K}v=(\sigma v)^{\mathsf T}
a(\sigma v)\geq0\).
\end{proof}

\begin{firewall}[Off-diagonal curvature pairs are real principal
terms]
\label{fire:e8v2-offdiagonal}
The matrix \({\mathsf K}\) need not be diagonal in plaquette labels.
Replacing it immediately by a multiple of
\(\sum_p\|F_p\|^2\) would discard principal finite-range curvature
pairs.  Their zero-momentum marginal part and their derivative
remainder must be separated by a moment or symbol argument.
\end{firewall}

\section{Translation-invariant quadratic marginal extraction}
\label{sec:e8v2-symbol}

Assume now that translated copies of the fixed dictionary have been
assembled on a four-dimensional coarse lattice.  Index plaquette
orientations by \(\alpha,\beta\in\Lambda^2\{1,2,3,4\}\), and write the
finite-range quadratic form as
\[
 K_2(F)
 =
 \frac12\sum_{x,r}\sum_{\alpha,\beta}
 K_{\alpha\beta}(r)
 \langle F_\alpha(x),F_\beta(x+r)\rangle_{\mathfrak g},
                                                        \label{eq:e8v2-translation}
\]
with
\[
 K_{\alpha\beta}(r)=K_{\beta\alpha}(-r).             \label{eq:e8v2-Hermitian}
\]
Its Fourier symbol is
\[
 \widehat K(k)=\sum_rK(r)e^{ik\cdot r}.              \label{eq:e8v2-symbol}
\]

\begin{lemma}[Second-moment symbol remainder]
\label{lem:e8v2-symbol}
If the first moment vanishes,
\[
 \sum_r r_\mu K(r)=0\qquad(\mu=1,\ldots,4),          \label{eq:e8v2-first-moment}
\]
then
\[
 \|\widehat K(k)-\widehat K(0)\|
 \leq
 \frac{|k|^2}{2}
 \sum_r|r|^2\|K(r)\|.                                \label{eq:e8v2-symbol-bound}
\]
Reflection symmetry \(K(r)=K(-r)\) is sufficient for
\eqref{eq:e8v2-first-moment}.
\end{lemma}

\begin{proof}
Use the first-moment cancellation to write
\[
 \widehat K(k)-\widehat K(0)
 =
 \sum_rK(r)\bigl(e^{ik\cdot r}-1-ik\cdot r\bigr).
\]
The scalar estimate
\(|e^{it}-1-it|\leq t^2/2\) gives
\eqref{eq:e8v2-symbol-bound}.  Pairing \(r\) with \(-r\) proves the
last assertion.
\end{proof}

\begin{lemma}[Rotational zero-momentum commutant]
\label{lem:e8v2-commutant}
Suppose \(\widehat K(0)\) commutes with the full
orientation-preserving Euclidean rotation action
\(\mathrm{SO}(4)\) on real two-forms.  Then there are real constants
\(c_{\rm YM},c_\theta\) such that
\[
 \widehat K(0)=c_{\rm YM}I+c_\theta *,
                                                        \label{eq:e8v2-commutant}
\]
where \(*\) is the four-dimensional Hodge operator.  If the action is
also invariant under an orientation-reversing reflection, then
\(c_\theta=0\).
\end{lemma}

\begin{proof}
The real two-form representation decomposes under
\(\mathrm{SO}(4)\) into the nonisomorphic self-dual and anti-self-dual
irreducible summands.  Schur's lemma makes a commuting real
self-adjoint operator scalar on each summand and gives no intertwiner
between them.  The two spectral
projections are \((I+*)/2\) and \((I-*)/2\), giving
\eqref{eq:e8v2-commutant}.  Orientation reversal conjugates
\(*\) to \(-*\), forcing \(c_\theta=0\) when it is a symmetry.
\end{proof}

\begin{theorem}[Quadratic marginal plus dimension-six symbol]
\label{thm:e8v2-marginal}
Under the hypotheses of
\Cref{lem:e8v2-symbol,lem:e8v2-commutant}, the quadratic curvature
form splits exactly in Fourier space as
\[
 K_2=K_{\rm marg}+K_{\rm der},                       \label{eq:e8v2-marginal-split}
\]
where
\[
 K_{\rm marg}
 =
 \frac{c_{\rm YM}}2\sum_x\langle F(x),F(x)\rangle
 +
 \frac{c_\theta}2\sum_x\langle F(x),*F(x)\rangle,    \label{eq:e8v2-marginal}
\]
and the derivative symbol satisfies
\[
 \|\widehat K_{\rm der}(k)\|
 \leq C_2|k|^2,\qquad
 C_2=\frac12\sum_r|r|^2\|K(r)\|.                    \label{eq:e8v2-derivative}
\]
Thus the only zero-momentum quadratic owners are the Yang--Mills and
topological densities; the remaining quadratic row carries two
powers of coarse momentum.
\end{theorem}

\begin{proof}
Define
\(\widehat K_{\rm der}(k)=\widehat K(k)-\widehat K(0)\) and use
\Cref{lem:e8v2-symbol}.  Insert
\eqref{eq:e8v2-commutant} at zero momentum and invert the Fourier
transform to obtain \eqref{eq:e8v2-marginal}.
\end{proof}

\begin{firewall}[Symmetry and first moments must be constructed]
\label{fire:e8v2-symmetry}
The exact filling theorem does not automatically make the assembled
coarse dictionary translation-, reflection-, or rotation-invariant.
One must average or choose fillings equivariantly and prove that the
ownership rule counts every interaction once.  Without
\eqref{eq:e8v2-first-moment}, the symbol remainder is only \(O(|k|)\);
without the \(\mathrm{SO}(4)\) commutant hypothesis, additional
anisotropic dimension-four quadratic owners remain.  Finite
hypercubic symmetry alone is not being asserted to imply the full
commutant statement; rotational restoration or calibration of those
extra owners is a separate cofinal row.
\end{firewall}

\section{Plaquette localization card and status}
\label{sec:e8v2-card}

\begin{definition}[Plaquette localization card, PLC]
\label{def:e8v2-PLC}
PLC requires:
\begin{enumerate}[label=\textup{(P\arabic*)},leftmargin=4em]
\item a shellable filling for every nonretained chord cycle;
\item explicit retention of nonfillable toron, cohomology, and
      boundary-transport cycles;
\item an exact lattice Stokes word with a uniformly finite filling
      hull;
\item cubic-order BCH substitution and the
      \(O_{C^2}(g^3R^5)\) tail;
\item equivariant assembly of the filling dictionary;
\item first-moment cancellation of the quadratic plaquette kernel;
\item explicit Yang--Mills/topological marginal ownership; and
\item a norm comparison placing the \(O(k^2)\), cubic, and quartic
      plaquette polynomials in the canonically stable interaction
      space.
\end{enumerate}
\end{definition}

\begin{proposition}[PLC ledger after V2]
\label{prop:e8v2-PLC}
V2 proves PLC \textup{(P1)--(P4)} for a declared fixed shellable
dictionary and proves \textup{(P6)--(P7)} under the explicit assembly
symmetries.  Equivariant global assembly in \textup{(P5)} and the
complete norm comparison in \textup{(P8)} remain open.
\end{proposition}

\begin{proof}
Use
\Cref{lem:e8v2-Stokes,thm:e8v2-BCH,
thm:e8v2-substitution,cor:e8v2-tail,
thm:e8v2-marginal,fire:e8v2-symmetry}.
\end{proof}

\begin{theorem}[What eighth-series V2 proves]
\label{thm:e8v2-result}
For each shellably filled chord loop, V2 proves an exact root-based
plaquette product and a cubic BCH expansion with differentiated
remainder \(O_{C^2}(g^3R^4)\).  Substitution into the complete quartic
chord polynomial produces explicit gauge-invariant quadratic, cubic,
and quartic plaquette polynomials plus a complete fixed-dictionary
tail \(O(g^3R^5)\).  Under reflection and rotationally invariant
assembly, the
quadratic zero-momentum owner is exactly a linear combination of
\(\Tr F^2\) and \(\Tr F\widetilde F\), while the remainder has
symbol \(O(k^2)\).
\end{theorem}

\begin{proof}
Combine
\Cref{lem:e8v2-Stokes,thm:e8v2-BCH,
thm:e8v2-substitution,cor:e8v2-tail,
thm:e8v2-marginal}.
\end{proof}

\begin{assessment}[Progress at eighth-series V2]
\label{ass:e8v2-status}
The compact-link Taylor bridge now lands in a finite plaquette
polynomial with the correct marginal quadratic owners and a vanishing
analytic/BCH tail.  The remaining non-small terms are the
finite-range \(O(k^2)\) quadratic kernel and the gauge-invariant cubic
and quartic plaquette polynomials.  Canonical power counting predicts
that they are irrelevant, but PLC (P8) requires a proof in the same
complete oscillation norm used by the exact coarse recursion.
\end{assessment}

\begin{decision}[V3 acquisition order]
\label{dec:e8v2-V3}
V3 will construct a scale-graded finite-range plaquette norm.  It will
prove that a first-moment-zero quadratic kernel gains \(b^{-2}\) under
block dilation and that the cubic/quartic plaquette polynomials acquire
their correct powers after the physical plaquette logarithm is
normalized by area.  The target is PLC (P8), with every norm-conversion
constant explicit.
\end{decision}

\begin{firewall}[Claim boundary after eighth-series V2]
\label{fire:e8v2-claim}
V2 does not prove equivariant global filling ownership, PLC (P8), the
running marginal coupling, the global nonlinear CCDC/GCDC margin,
chart and saturated-bad closure, NCGC/CFTC along the full ultraviolet
trajectory, terminal entry, regulator independence, continuum OS
reconstruction, or the full Yang--Mills mass gap.
\end{firewall}

\paragraph{Parent-version record.}
V2 was formed from
\texttt{Renewal\_Yang\_Mills\_Mass\_Gap\_Research\_Notes\_V1.tex}.
The V1 parent had \(815\) counted source lines,
\(31\,509\) bytes, and SHA--256
\[
 \texttt{DAA29C234FED657F37B203C00EABE66E1E2BCD094830EE5256A482A38AC6AF0E}.
\]
The next compulsory companion audit is V5.

% =============================================================================
% END APPENDED EIGHTH-SERIES V2 MATERIAL
% =============================================================================

\clearpage
% =============================================================================
% BEGIN APPENDED EIGHTH-SERIES V3 MATERIAL
% =============================================================================

\chapter{A scale-graded plaquette norm and canonical stable
contraction}
\label{ch:e8v3}

\section{Purpose and normalization}
\label{sec:e8v3-purpose}

V2 produced a finite plaquette polynomial
\[
 K_2+gK_3+g^2K_4
\]
plus a vanishing complete tail.  This note puts the nonmarginal part
of that polynomial in a scale-graded norm which dominates its complete
good-set oscillation.  The lattice plaquette logarithm is normalized
by physical area: if the lattice spacing is \(a\), write
\[
 P_p=\exp(gF_p),\qquad {\mathcal F}_p:=a^{-2}F_p.
                                                        \label{eq:e8v3-area}
\]
Thus \({\mathcal F}\) has canonical dimension two.  A covariant
difference contributes one additional power of inverse length.

\begin{scope}[What the dilation theorem covers]
\label{scope:e8v3}
The dilation below is the canonical reindexing of an already assembled
translation-compatible local density.  It is not a proof that the
exact nonlinear block map produces that reindexing with unit geometric
multiplicity.  That ownership issue is PLC (P5) and remains separate.
\end{scope}

\section{Coefficient rows}
\label{sec:e8v3-coefficients}

For a finite-range translation-invariant quadratic kernel \(K(r)\),
define
\begin{align}
 \|K\|_{\mu,0}
 &:=
 \sup_\alpha\sum_{\beta,r}
 e^{\mu|r|}|K_{\alpha\beta}(r)|,                     \label{eq:e8v3-K0}\\
 \|K\|_{\mu,2}
 &:=
 \sup_\alpha\sum_{\beta,r\neq0}
 e^{\mu|r|}|r|^2|K_{\alpha\beta}(r)|.                \label{eq:e8v3-K2}
\end{align}
The derivative sector satisfies
\[
 \sum_rK(r)=0,\qquad \sum_rr_\rho K(r)=0.            \label{eq:e8v3-moments}
\]
The value \(K(0)\) is therefore fixed by the nonzero displacements.

For \(n=3,4\), let \(T_n(X)\) be the invariant \(n\)-linear
coefficient tensor of a finite-range plaquette monomial supported on
\(X\).  Put
\[
 \|T_n\|_{\mu,\rm row}
 :=
 \sup_p\sum_{X\ni p}
 e^{\mu\operatorname{diam}X}
 (|X|-1)n\,\|T_n(X)\|.                               \label{eq:e8v3-Trow}
\]
The factor \(n\) marks one field occurrence.

\begin{definition}[Scale-graded stable plaquette norm]
\label{def:e8v3-stable}
On the field ball \(\|F_p\|\leq R\), define
\[
 \|\Psi\|_{{\rm st},\mu,R}
 :=
 R^2\|K_{\rm der}\|_{\mu,2}
 +R^3\|gT_3\|_{\mu,\rm row}
 +R^4\|g^2T_4\|_{\mu,\rm row}.                       \label{eq:e8v3-stable}
\]
The Yang--Mills, topological, and any declared anisotropic
dimension-four quadratic coefficients are not included in this norm;
they belong to the calibrated marginal row.
\end{definition}

\section{Domination of complete oscillation}
\label{sec:e8v3-oscillation}

\begin{lemma}[Zero-sum quadratic coefficient row]
\label{lem:e8v3-zero-sum}
If \(\sum_rK(r)=0\), then
\[
 \|K\|_{\mu,0}
 \leq
 2\|K\|_{\mu,2}                                      \label{eq:e8v3-zero-sum}
\]
after replacing \(\mu\) on the right by the same or a larger value.
\end{lemma}

\begin{proof}
For each orientation pair,
\[
 K(0)=-\sum_{r\neq0}K(r).
\]
Hence its absolute row is bounded by the nonzero-displacement row.
For \(r\neq0\), \(|r|^2\geq1\), so the latter is bounded by
\eqref{eq:e8v3-K2}.  The nonzero and zero rows together give the
factor two.
\end{proof}

\begin{lemma}[Oscillation of a multilinear local monomial]
\label{lem:e8v3-monomial}
Let \(T\) be \(n\)-linear with operator norm \(\|T\|\), and set
\(\psi(F_1,\ldots,F_n)=T[F_1,\ldots,F_n]\).
On \(\|F_i\|\leq R\), varying one occurrence gives
\[
 \osc_i\psi\leq2\|T\|R^n.                            \label{eq:e8v3-monomial-osc}
\]
Marking an unspecified one of the \(n\) occurrences costs at most
\[
 2n\|T\|R^n.                                         \label{eq:e8v3-monomial-mark}
\]
\end{lemma}

\begin{proof}
The difference is
\[
 T[F_1,\ldots,F_i-F_i',\ldots,F_n],
\]
whose norm is at most
\(\|T\|(2R)R^{n-1}\).  Sum over the at most \(n\) possible marked
occurrences.
\end{proof}

\begin{theorem}[Stable norm controls the complete row]
\label{thm:e8v3-domination}
Let
\[
 \Psi_{\rm st}
 =
 \Psi_{K_{\rm der}}+g\Psi_3+g^2\Psi_4               \label{eq:e8v3-Psi}
\]
be the assembled stable plaquette polynomial, with unique support
ownership.  Then its complete weighted one-mark oscillation norm
satisfies
\[
 {\cal O}_\mu(\Psi_{\rm st})
 \leq
 C_{\rm osc}\|\Psi_{\rm st}\|_{{\rm st},\mu,R},      \label{eq:e8v3-domination}
\]
where \(C_{\rm osc}\) depends only on the orientation count and the
fixed support convention, not on \(a,g\), or the volume.  The
two-mark mixed row has the same bound with another fixed constant.
\end{theorem}

\begin{proof}
Changing one plaquette field in the quadratic form affects the two
placements in which it is a left or right endpoint.  Its oscillation
is bounded by a fixed orientation factor times
\(R^2\|K_{\rm der}\|_{\mu,0}\), and
\Cref{lem:e8v3-zero-sum} replaces the zeroth coefficient row by the
second-moment row.  For the cubic and quartic forms apply
\Cref{lem:e8v3-monomial} and sum the uniquely owned supports using
\eqref{eq:e8v3-Trow}.  Two variations replace \(2n\) by at most
\(4n(n-1)\), which is fixed for \(n\leq4\).
\end{proof}

\begin{firewall}[Oscillation domination is one-way]
\label{fire:e8v3-one-way}
The stable coefficient norm controls the complete oscillation row.
The converse is false without normalization conditions, because
cancellations can make an oscillation small while individual
coefficients are large.  Iteration must therefore propagate the
graded coefficients or prove a fresh coefficient extraction at each
scale; it cannot reconstruct them from an oscillation bound alone.
\end{firewall}

\section{Canonical dilation}
\label{sec:e8v3-dilation}

A gauge-invariant plaquette monomial with \(n\) curvature factors and
\(s\) covariant differences has canonical dimension
\[
 \Delta(n,s)=2n+s.                                   \label{eq:e8v3-dimension}
\]
Under \(x=bx'\) and
\({\mathcal F}'(x')=b^2{\mathcal F}(bx')\), its coupling is multiplied
by
\[
 b^{4-\Delta(n,s)}.                                  \label{eq:e8v3-multiplier}
\]
This is the area-normalized form of the V100 scaling law.

\begin{lemma}[Quadratic moment sector has two derivative powers]
\label{lem:e8v3-quadratic-scale}
Let \(K_{\rm der}\) satisfy \eqref{eq:e8v3-moments}, and let
\(\widehat K_{\rm der}\) be its Fourier symbol.  For rescaled momentum
\(k'=bk\),
\[
 \|\widehat K_{\rm der}(k'/b)\|
 \leq
 \frac{|k'|^2}{2b^2}
 \sum_r|r|^2\|K_{\rm der}(r)\|.                      \label{eq:e8v3-quadratic-scale}
\]
\end{lemma}

\begin{proof}
This is \Cref{lem:e8v2-symbol} evaluated at \(k'/b\).
\end{proof}

\begin{theorem}[Canonical stable contraction]
\label{thm:e8v3-contraction}
Let \({\mathscr D}_b\) be the unique-density canonical reindexing which:
\begin{enumerate}[label=\textup{(\roman*)},leftmargin=3.5em]
\item sends the quadratic moment-zero sector to its rescaled symbol
      \(\widehat K_{\rm der}(k'/b)\);
\item multiplies an \(n\)-curvature, \(s\)-difference local
      coefficient by \(b^{4-2n-s}\); and
\item preserves the normalized field radius \(R\), support ownership,
      and orientation coefficient norms.
\end{enumerate}
Then
\[
 \|{\mathscr D}_b\Psi_{\rm st}\|_{{\rm st},\mu,R}
 \leq
 b^{-2}\|\Psi_{\rm st}\|_{{\rm st},\mu,R}.           \label{eq:e8v3-contraction}
\]
\end{theorem}

\begin{proof}
The quadratic moment-zero sector has \(n=2,s=2\) and gains \(b^{-2}\)
by \Cref{lem:e8v3-quadratic-scale}.  The cubic sector has
\(\Delta(3,0)=6\) and gains \(b^{-2}\).  The quartic sector has
\(\Delta(4,0)=8\) and gains \(b^{-4}\).  Extra differences only
decrease the multiplier.  Taking the sum norm
\eqref{eq:e8v3-stable} gives the largest factor \(b^{-2}\).
\end{proof}

\begin{corollary}[Stable complete-row contraction after norm
conversion]
\label{cor:e8v3-row-contraction}
Under the hypotheses of
\Cref{thm:e8v3-domination,thm:e8v3-contraction},
\[
 {\cal O}_\mu({\mathscr D}_b\Psi_{\rm st})
 \leq
 C_{\rm osc}b^{-2}
 \|\Psi_{\rm st}\|_{{\rm st},\mu,R}.                 \label{eq:e8v3-row-contraction}
\]
\end{corollary}

\begin{proof}
Apply \Cref{thm:e8v3-domination} after
\Cref{thm:e8v3-contraction}.
\end{proof}

\begin{firewall}[Running parameters and geometric multiplicity]
\label{fire:e8v3-geometry}
The actual exact RG step also changes \(g\), may combine several fine
supports into one coarse support, and may create transition-chart and
determinant terms.  If the support reindexing has row multiplicity
\(D_{\rm geom}\), the proved factor is
\(D_{\rm geom}b^{-2}\), not \(b^{-2}\).  Strict contraction requires
either unique ownership with \(D_{\rm geom}=1\) or the explicit
inequality \(D_{\rm geom}<b^2\).  Running-coupling changes remain in
the calibrated marginal ledger.
\end{firewall}

\section{Placement of the V2 polynomial}
\label{sec:e8v3-placement}

\begin{theorem}[Local polynomial placement]
\label{thm:e8v3-placement}
Assume:
\begin{enumerate}[label=\textup{(\alph*)},leftmargin=3.5em]
\item the quadratic kernel has been split as in
      \Cref{thm:e8v2-marginal};
\item all dimension-four rotational or lattice-anisotropic
      coefficients are placed in the calibrated marginal sector;
\item the remaining quadratic kernel has the two moment cancellations
      \eqref{eq:e8v3-moments}; and
\item the cubic and quartic supports have unique fixed-dictionary
      ownership.
\end{enumerate}
Then the nonmarginal polynomial
\[
 K_{\rm der}+gK_3+g^2K_4
                                                        \label{eq:e8v3-nonmarginal}
\]
belongs to the stable norm \eqref{eq:e8v3-stable}, its complete
oscillation row is finite uniformly in volume, and its canonical
dilation contracts by \(b^{-2}\).
\end{theorem}

\begin{proof}
Finite range and a fixed dictionary make the coefficient rows finite.
The moment assumptions place the quadratic term in
\(\|K\|_{\mu,2}\), while the other terms are in the \(n=3,4\) rows.
Apply
\Cref{thm:e8v3-domination,thm:e8v3-contraction}.
\end{proof}

\begin{corollary}[PLC (P8) on the declared polynomial sector]
\label{cor:e8v3-PLC}
Under \Cref{thm:e8v3-placement}, PLC (P8) is proved for the
symmetry-calibrated fixed-dictionary polynomial sector.  Combining
with the V1/V2 tail gives
\[
 {\cal W}_g
 =
 K_{\rm marg}
 +\Psi_{\rm st}
 +{\cal T}_{5,g},                                    \label{eq:e8v3-final-split}
\]
where the stable polynomial has canonical factor at most \(b^{-2}\)
and the complete tail is \(O(g^3R^5)\).
\end{corollary}

\begin{proof}
Use
\Cref{thm:e8v2-substitution,cor:e8v2-tail,
thm:e8v2-marginal,thm:e8v3-placement}.
\end{proof}

\section{Status and next acquisition}
\label{sec:e8v3-status}

\begin{theorem}[What eighth-series V3 proves]
\label{thm:e8v3-result}
V3 constructs a scale-graded finite-range plaquette norm which
dominates the complete good-set oscillation and mixed-oscillation
rows.  After every dimension-four quadratic owner is calibrated, the
moment-zero quadratic kernel, cubic curvature polynomial, and quartic
curvature polynomial contract canonically by factors
\(b^{-2},b^{-2},b^{-4}\), respectively.  Thus the full declared stable
polynomial has factor at most \(b^{-2}\), while the V1/V2 analytic/BCH
tail remains \(O(g^3R^5)\).
\end{theorem}

\begin{proof}
Combine
\Cref{thm:e8v3-domination,thm:e8v3-contraction,
thm:e8v3-placement,cor:e8v3-PLC}.
\end{proof}

\begin{assessment}[Progress at eighth-series V3]
\label{ass:e8v3-status}
The norm-conversion part of PLC (P8) is now explicit for a uniquely
owned polynomial dictionary.  The next obstruction is no longer
power counting.  It is geometric: construct an equivariant family of
trees and shellable fillings whose averaging preserves exact
ownership, reflection, gauge covariance, the first-moment
cancellation, and a multiplicity below the \(b^2\) contraction margin.
\end{assessment}

\begin{decision}[V4 acquisition order]
\label{dec:e8v3-V4}
V4 will analyze symmetry averaging of tree/filling choices.  It will
distinguish averaging the exact coarse density from averaging a
logarithmic interaction, prove the correct orbit-ownership formula,
and compute the resulting support-row multiplicity.  If no
deterministic equivariant tree exists, the tree label will be retained
as an auxiliary finite variable rather than divided out after taking
the logarithm.
\end{decision}

\begin{firewall}[Claim boundary after eighth-series V3]
\label{fire:e8v3-claim}
V3 does not prove PLC (P5), a strict physical RG contraction when
\(D_{\rm geom}\geq b^2\), the running marginal coupling, global chart
and saturated-bad closure, NCGC/CFTC along the ultraviolet trajectory,
terminal entry, regulator independence, continuum OS reconstruction,
or the full Yang--Mills mass gap.
\end{firewall}

\paragraph{Parent-version record.}
V3 was formed from
\texttt{Renewal\_Yang\_Mills\_Mass\_Gap\_Research\_Notes\_V2.tex}.
The V2 parent had \(1\,420\) counted source lines,
\(52\,322\) bytes, and SHA--256
\[
 \texttt{E1A288C922997FFDD9FE506837311DAE7DEA6F100E9769228B877BD0E352E2AF}.
\]
The next compulsory companion audit is V5.

% =============================================================================
% END APPENDED EIGHTH-SERIES V3 MATERIAL
% =============================================================================

% =============================================================================
% BEGIN APPENDED EIGHTH-SERIES V4 MATERIAL
% =============================================================================

\part{Eighth-series V4: exact symmetry recovery without mixing
logarithmic actions}
\label{part:e8v4}

\section{The obstruction and the correct level of averaging}
\label{sec:e8v4-obstruction}

V3 left tree and filling covariance as a geometric acquisition.  The
first point of V4 is that an invariant deterministic tree need not
exist and is not needed.  A tree used only to coordinatize one fixed
Haar fibre is a proof device.  It may be averaged after an exact
polynomial-plus-tail decomposition of the same invariant block
potential.  This is categorically different from averaging coarse
densities obtained from different block maps.

\begin{definition}[Finite symmetry datum]
\label{def:e8v4-symmetry}
Let \(\Gamma\) be a finite group acting on the fine block graph, on
the coarse boundary datum \(y\), and linearly on interactions by
\[
 (\gamma\Psi)(y):=\Psi(\gamma^{-1}y).
                                                        \label{eq:e8v4-action}
\]
The action is called norm isometric when
\(\|\gamma\Psi\|_{\mathfrak B}=\|\Psi\|_{\mathfrak B}\) for the
coefficient or oscillation norm under consideration.  Write
\[
 {\mathsf S}_{\Gamma}\Psi
 :=|\Gamma|^{-1}\sum_{\gamma\in\Gamma}\gamma\Psi .
                                                        \label{eq:e8v4-reynolds}
\]
\end{definition}

\begin{proposition}[An invariant tree may be impossible]
\label{prop:e8v4-no-tree}
Let a finite group \(\Gamma\) act edge transitively on a finite,
connected graph \({\mathcal G}=(V,E)\) which contains a cycle.
There is no \(\Gamma\)-invariant spanning tree of \({\mathcal G}\).
\end{proposition}

\begin{proof}
If \(T\subseteq E\) is invariant and nonempty, choose \(e\in T\).
Edge transitivity gives, for every \(e'\in E\), a
\(\gamma\in\Gamma\) with \(e'=\gamma e\).  Invariance gives
\(e'\in T\), hence \(T=E\), which contains a cycle.  If \(T\) is
empty, it is not connected.  Neither possibility is a spanning tree.
\end{proof}

\begin{lemma}[Reynolds projection]
\label{lem:e8v4-reynolds}
If \(\Gamma\) acts isometrically on a normed interaction space
\({\mathfrak B}\), then \({\mathsf S}_{\Gamma}\) is a projection onto
the invariant subspace and
\[
 \|{\mathsf S}_{\Gamma}\Psi\|_{\mathfrak B}
 \leq \|\Psi\|_{\mathfrak B}.                         \label{eq:e8v4-reynolds-bound}
\]
\end{lemma}

\begin{proof}
For \(\eta\in\Gamma\), reindexing the finite sum gives
\(\eta{\mathsf S}_{\Gamma}={\mathsf S}_{\Gamma}\).  It follows that
\({\mathsf S}_{\Gamma}^2={\mathsf S}_{\Gamma}\), with the stated
range.  The triangle inequality and isometry give
\[
 \|{\mathsf S}_{\Gamma}\Psi\|_{\mathfrak B}
 \leq |\Gamma|^{-1}\sum_{\gamma\in\Gamma}
       \|\gamma\Psi\|_{\mathfrak B}
 =\|\Psi\|_{\mathfrak B}.
\]
\end{proof}

\begin{theorem}[Exact symmetrization of one decomposition]
\label{thm:e8v4-exact-sym}
Let \(W\in{\mathfrak B}\) be \(\Gamma\)-invariant.  Suppose one
coordinate choice \(\alpha\), such as a rooted tree and shellable
filling, gives the exact identity
\[
 W=P_\alpha+E_\alpha .                               \label{eq:e8v4-one-split}
\]
Then
\[
 W={\mathsf S}_{\Gamma}P_\alpha+
   {\mathsf S}_{\Gamma}E_\alpha                      \label{eq:e8v4-sym-split}
\]
is exact, both summands are invariant, and
\[
 \|{\mathsf S}_{\Gamma}E_\alpha\|_{\mathfrak B}
 \leq \|E_\alpha\|_{\mathfrak B}.                    \label{eq:e8v4-tail-no-loss}
\]
In particular, finite-group symmetrization introduces no factor
\(|\Gamma|\) in a \(\Gamma\)-isometric Banach row norm.
\end{theorem}

\begin{proof}
Apply \({\mathsf S}_{\Gamma}\) to
\eqref{eq:e8v4-one-split}.  Since \(W\) is invariant,
\({\mathsf S}_{\Gamma}W=W\).  Invariance of the two new summands and
the norm estimate follow from \Cref{lem:e8v4-reynolds}.
\end{proof}

\begin{remark}[What the theorem does not remove]
\label{rem:e8v4-geom}
The normalized orbit average prevents a \(|\Gamma|\) loss caused
solely by rewriting the same interaction in symmetric form.  It does
not remove the actual RG multiplicity \(D_{\rm geom}\) of
\Cref{fire:e8v3-geometry}, which counts distinct fine supports
regrouped into one coarse owner.
\end{remark}

\section{Coordinate independence on a fixed Haar fibre}
\label{sec:e8v4-coordinate}

\begin{definition}[Coordinate-only tree family]
\label{def:e8v4-coordinate}
For fixed coarse datum \(y\), let \({\mathcal F}_y\) be the same
compact fine-link fibre for every label \(\alpha\).  A family of tree
gauges \(\chi_{\alpha,y}:G^{N_y}\to{\mathcal F}_y\) is
coordinate-only if every transition map
\[
 \tau_{\beta\alpha,y}:=
 \chi_{\beta,y}^{-1}\chi_{\alpha,y}                  \label{eq:e8v4-transition}
\]
is a measurable bijection preserving product Haar measure.
\end{definition}

\begin{proposition}[The exact density is tree independent]
\label{prop:e8v4-tree-independent}
Let \(A(U;y)\) be one fixed measurable fine action and let the family
in \Cref{def:e8v4-coordinate} be coordinate-only.  Then
\[
 Z_\alpha(y):=\int_{G^{N_y}}
 \exp\{-A(\chi_{\alpha,y}(u);y)\}\,du                \label{eq:e8v4-Zalpha}
\]
is independent of \(\alpha\).
\end{proposition}

\begin{proof}
For labels \(\alpha,\beta\), substitute
\(v=\tau_{\beta\alpha,y}(u)\).  Haar preservation and
\(\chi_{\alpha,y}(u)=\chi_{\beta,y}(v)\) give
\[
 Z_\alpha(y)
 =\int\exp\{-A(\chi_{\beta,y}(v);y)\}\,dv
 =Z_\beta(y).
\]
\end{proof}

\begin{corollary}[Arbitrary proof trees are sufficient]
\label{cor:e8v4-arbitrary-tree}
When the tree/filling labels satisfy
\Cref{def:e8v4-coordinate}, one may prove
\(W=-\log Z_\alpha=P_\alpha+E_\alpha\) using any single label.
If the physical block map and action make \(W\) \(\Gamma\)-invariant,
\Cref{thm:e8v4-exact-sym} produces an exact invariant decomposition
without an invariant tree.
\end{corollary}

\begin{proof}
\Cref{prop:e8v4-tree-independent} identifies the exact density for
all coordinate labels.  Apply \Cref{thm:e8v4-exact-sym}.
\end{proof}

\begin{firewall}[A filling must not change the fibre]
\label{fire:e8v4-filling}
A shellable filling is coordinate-only only when it changes the
algebraic Stokes representation of the same chord holonomy, not the
block constraint, retained variables, or integration fibre.
Otherwise \Cref{prop:e8v4-tree-independent} does not apply.
\end{firewall}

\section{Why distinct densities cannot be averaged after logarithms}
\label{sec:e8v4-logsum}

\begin{proposition}[Log-sum-exp identity and equality case]
\label{prop:e8v4-logsum}
Let \(S_1,\ldots,S_m\) be finite real functions and put
\[
 Z_{\rm mix}(y):=\frac1m\sum_{\alpha=1}^m e^{-S_\alpha(y)},
 \qquad
 S_{\rm mix}(y):=-\log Z_{\rm mix}(y).               \label{eq:e8v4-mix}
\]
Then
\[
 S_{\rm mix}(y)\leq \frac1m\sum_{\alpha=1}^mS_\alpha(y),
                                                        \label{eq:e8v4-jensen}
\]
with equality at \(y\) if and only if all \(S_\alpha(y)\) are equal.
In general \(S_{\rm mix}\neq m^{-1}\sum_\alpha S_\alpha\).
\end{proposition}

\begin{proof}
Strict convexity of \(s\mapsto e^{-s}\) gives
\[
 e^{-m^{-1}\sum_\alpha S_\alpha}
 \leq m^{-1}\sum_\alpha e^{-S_\alpha}.
\]
Apply the decreasing function \(-\log\).  Equality in strict Jensen
holds exactly when the arguments agree.
\end{proof}

\begin{example}[A visible nonlinear defect]
\label{ex:e8v4-cosh}
For \(S_\pm(y)=\pm t f(y)\),
\[
 \tfrac12(S_++S_-)=0,\qquad
 S_{\rm mix}=-\log\cosh(tf).                         \label{eq:e8v4-cosh}
\]
Thus action averaging deletes the nonzero even interaction generated
by exact density mixing.
\end{example}

\begin{theorem}[Exact auxiliary-label representation]
\label{thm:e8v4-label}
For genuinely different block maps or fibres, introduce the finite
label \(\alpha\in\{1,\ldots,m\}\) and the joint density
\[
 \rho(y,\alpha):=\frac1m e^{-S_\alpha(y)}.            \label{eq:e8v4-joint}
\]
Its \(y\)-marginal is exactly \(e^{-S_{\rm mix}(y)}\).  If
\(\Gamma\) acts on labels and
\[
 S_{\gamma\alpha}(\gamma y)=S_\alpha(y),             \label{eq:e8v4-joint-equiv}
\]
then \(\rho(\gamma y,\gamma\alpha)=\rho(y,\alpha)\).
Eliminating \(\alpha\) necessarily produces the log-sum-exp
interaction in \eqref{eq:e8v4-mix}.
\end{theorem}

\begin{proof}
Summing \eqref{eq:e8v4-joint} over \(\alpha\) gives
\(Z_{\rm mix}\).  Equation \eqref{eq:e8v4-joint-equiv} gives joint
invariance directly.  The last statement is the definition of the
negative logarithm of this marginal.
\end{proof}

\begin{firewall}[The label is an interacting owner]
\label{fire:e8v4-label}
A retained finite label is not cost free.  Its conditional
log-sum-exp cumulants are local interaction owners and must be
included in the complete oscillation, derivative, bad-set, and
reflection ledgers.  V4 uses no auxiliary label when a tree is merely
a coordinate choice.
\end{firewall}

\section{Orbit ownership and quadratic moments}
\label{sec:e8v4-orbits}

\begin{definition}[Orbit-weighted owner]
\label{def:e8v4-owner}
If \(\Psi=\sum_X\Psi_X\) is a fixed-dictionary interaction, define
\[
 ({\mathsf S}_{\Gamma}\Psi)_{[X]}
 :=
 \frac1{|\Gamma|}\sum_{\gamma\in\Gamma}
       \gamma\Psi_X,                                 \label{eq:e8v4-orbit-owner}
\]
owned by the orbit label \([X]\).  Repeated stabilizer copies retain
their normalized weights; they are not replaced by an unweighted
sum.
\end{definition}

\begin{proposition}[No symmetry multiplicity in an invariant row norm]
\label{prop:e8v4-orbit-norm}
Suppose the support weight \(w(X)\), incidence relation, and
coefficient seminorm are \(\Gamma\)-invariant.  For the normalized
row norm
\[
 \|\Psi\|_{\rm row}
 :=\sup_z\sum_{X\ni z}w(X)\|\Psi_X\|_X,              \label{eq:e8v4-row}
\]
the orbit-weighted decomposition satisfies
\[
 \|{\mathsf S}_{\Gamma}\Psi\|_{\rm row}
 \leq\|\Psi\|_{\rm row}.                             \label{eq:e8v4-row-bound}
\]
\end{proposition}

\begin{proof}
Expand \eqref{eq:e8v4-orbit-owner}, use the triangle inequality, and
then interchange the finite orbit average with the row sum.  For
fixed \(\gamma\), substitute \(X'=\gamma X\) and \(z'=\gamma z\).
Invariance gives
\[
 \sup_z\sum_{X\ni z}w(X)\|\gamma\Psi_X\|_{\gamma X}
 =\sup_{z'}\sum_{X'\ni z'}w(X')\|\Psi_{X'}\|_{X'}.
\]
The normalized average of these identical upper bounds proves
\eqref{eq:e8v4-row-bound}.
\end{proof}

\begin{definition}[Finite-cutoff quadratic commutant]
\label{def:e8v4-commutant}
Let \(V=\mathfrak g\otimes\Lambda^2\mathbb R^4\) carry the curvature
representation of the block stabilizer \(\Gamma\).  Its real
self-adjoint commutant is
\[
 {\mathcal C}_\Gamma
 :=\{A\in\operatorname{End}(V):A=A^*,\
       A\gamma=\gamma A\ \hbox{for all }\gamma\in\Gamma\}. \label{eq:e8v4-commutant}
\]
\end{definition}

\begin{lemma}[Inversion kills the first moment]
\label{lem:e8v4-first-moment}
Let a finite-range quadratic kernel \(K(r)\) obey
\(K(r)=K(-r)\).  Then
\[
 \sum_r r_jK(r)=0\qquad(1\leq j\leq4).               \label{eq:e8v4-first-moment}
\]
\end{lemma}

\begin{proof}
Pair the summand at \(r\) with that at \(-r\).  Fixed points satisfy
\(r=-r\) and have vanishing representative displacement in the
local, non-wrapping block dictionary.
\end{proof}

\begin{theorem}[Symmetric marginal/stable split]
\label{thm:e8v4-quadratic-split}
Let \(K(r)\) be finite range and let its Reynolds average be taken
under a block stabilizer containing inversion.  Put
\[
 A_0:=\sum_r({\mathsf S}_{\Gamma}K)(r).              \label{eq:e8v4-A0}
\]
Then \(A_0\in{\mathcal C}_\Gamma\), and the remainder
\[
 K_{\rm der}(r):=({\mathsf S}_{\Gamma}K)(r)
                 -A_0\mathbf1_{\{r=0\}}             \label{eq:e8v4-Kder}
\]
has vanishing zeroth and first moments.  Consequently its symbol
satisfies
\[
 \|\widehat K_{\rm der}(k)\|
 \leq\frac{|k|^2}{2}\sum_r|r|^2
        \|K_{\rm der}(r)\|.                          \label{eq:e8v4-symbol}
\]
After every coefficient in \({\mathcal C}_\Gamma\) is placed in the
calibrated marginal ledger, this quadratic remainder is in the
V3 stable sector and canonically gains \(b^{-2}\).
\end{theorem}

\begin{proof}
The Reynolds average commutes with \(\Gamma\), and evaluation at zero
momentum preserves that property, proving
\(A_0\in{\mathcal C}_\Gamma\).  The zeroth moment vanishes by
construction; the first vanishes by
\Cref{lem:e8v4-first-moment}.  Taylor's formula with integral
remainder gives
\[
 e^{ik\cdot r}-1-ik\cdot r
 =(ik\cdot r)^2\int_0^1(1-s)e^{isk\cdot r}\,ds,
\]
whose modulus is at most \(|k|^2|r|^2/2\).  Summation proves
\eqref{eq:e8v4-symbol}.  The last assertion is
\Cref{thm:e8v3-contraction}.
\end{proof}

\begin{firewall}[Finite lattice symmetry is not \(SO(4)\)]
\label{fire:e8v4-commutant}
At fixed cutoff, \({\mathcal C}_\Gamma\) can be strictly larger than
the span of the Yang--Mills and theta quadratic forms.  V4 calibrates
the entire finite-dimensional commutant.  Reduction to
\(\operatorname{Tr}F^2\) and
\(\operatorname{Tr}(F\widetilde F)\) requires a separate rotational
restoration theorem, including decay or tuning of every anisotropic
marginal coefficient.
\end{firewall}

\section{Reflection compatibility and the V4 acquisition card}
\label{sec:e8v4-card}

\begin{proposition}[Reflection is unchanged by a proof decomposition]
\label{prop:e8v4-reflection}
Suppose the exact density \(e^{-W}\) is reflection positive and
\(W\) is invariant under a finite stabilizer
\(\Gamma_+\) preserving the reflection plane and commuting with the
reflection involution.  Replacing
\(W=P_\alpha+E_\alpha\) by the exact identity
\eqref{eq:e8v4-sym-split} with \(\Gamma_+\) changes neither the
density nor reflection positivity.
\end{proposition}

\begin{proof}
\Cref{thm:e8v4-exact-sym} states equality of the two functions
\(W\).  Therefore their exponentials define the same measure, so
every reflection-positivity quadratic form is identical.
\end{proof}

\begin{firewall}[Mixtures require joint positivity]
\label{fire:e8v4-rp-mixture}
If different block maps are mixed, invariance of the marginal is not
a proof of reflection positivity.  One must prove positivity of the
joint \(y,\alpha\) density under a reflection action on the label.
Conditional positivity at each frozen slow background is not enough
to infer positivity after an interacting marginalization.
\end{firewall}

\begin{definition}[Equivariant decomposition/ownership card]
\label{def:e8v4-card}
An RG block passes EDOC when:
\begin{enumerate}[label=\textup{(E\arabic*)},leftmargin=4em]
\item the physical block map is \(\Gamma\)-equivariant;
\item each tree/filling choice is certified either coordinate-only or
      an explicit auxiliary label;
\item coordinate changes preserve the exact Haar fibre;
\item symmetrization is performed on an exact decomposition of one
      invariant \(W\);
\item the coefficient and oscillation norms are \(\Gamma\)-isometric;
\item the full finite-cutoff commutant
      \({\mathcal C}_\Gamma\) is calibrated;
\item actual support regrouping \(D_{\rm geom}\) is counted
      independently of symmetry averaging; and
\item the stabilizer and any auxiliary label action are compatible
      with reflection.
\end{enumerate}
\end{definition}

\begin{theorem}[EDOC consequence for the V1--V3 decomposition]
\label{thm:e8v4-card}
Assume EDOC and suppose the V2 polynomial/tail decomposition is
obtained from one coordinate-only tree and filling.  Then it has an
exact invariant version
\[
 W=K_{\rm marg}+\Psi_{\rm st}+{\cal T}_{5,g},         \label{eq:e8v4-final}
\]
where:
\begin{enumerate}[label=\textup{(\roman*)},leftmargin=3.5em]
\item \(K_{\rm marg}\) contains all of
      \({\mathcal C}_\Gamma\);
\item the symmetrized tail still satisfies the V1/V2
      \(O(g^3R^5)\) complete-row bound;
\item inversion supplies the first-moment cancellation;
\item the stable polynomial has canonical factor at most \(b^{-2}\);
      and
\item symmetry averaging adds no factor to its invariant row norm.
\end{enumerate}
The independent geometric factor \(D_{\rm geom}\) and running
parameters remain as in \Cref{fire:e8v3-geometry}.
\end{theorem}

\begin{proof}
Coordinate independence and exact symmetrization are
\Cref{cor:e8v4-arbitrary-tree,thm:e8v4-exact-sym}.  The tail estimate
and absence of a symmetry factor follow from
\Cref{lem:e8v4-reynolds,prop:e8v4-orbit-norm}.  The quadratic split is
\Cref{thm:e8v4-quadratic-split}.  Apply the V3 placement and
contraction theorem to the calibrated remainder, cubic, and quartic
sectors.
\end{proof}

\section{Status and next acquisition}
\label{sec:e8v4-status}

\begin{theorem}[What eighth-series V4 proves]
\label{thm:e8v4-result}
For tree gauges and shellable fillings which merely coordinatize one
fixed Haar fibre, no equivariant deterministic tree is required.
Any exact polynomial-plus-tail decomposition can be projected to an
invariant one without increasing an invariant Banach row norm.
Inversion enforces the quadratic first moment, while calibration of
the full finite-cutoff commutant leaves an \(O(k^2)\) stable kernel.
By contrast, genuinely different block maps must be mixed at the
density level and produce an unavoidable log-sum-exp interaction.
\end{theorem}

\begin{proof}
Combine
\Cref{prop:e8v4-no-tree,prop:e8v4-tree-independent,
thm:e8v4-exact-sym,prop:e8v4-logsum,
thm:e8v4-quadratic-split,thm:e8v4-card}.
\end{proof}

\begin{assessment}[Progress at eighth-series V4]
\label{ass:e8v4-status}
The proof-tree portion of PLC (P5), and the inversion-based
first-moment portion of PLC (P6), are now closed under EDOC.  The
marginal part of PLC (P7) is correctly the whole lattice commutant,
not merely the continuum Yang--Mills/theta pair.  The remaining
strict physical contraction still requires a unique coarse owner or
an explicit \(D_{\rm geom}<b^2\), together with control of running
couplings and anisotropic marginal coefficients.
\end{assessment}

\begin{decision}[V5 acquisition order]
\label{dec:e8v4-V5}
V5 is a compulsory companion-program audit.  After that audit it will
attack the finite-cutoff marginal transport matrix: separate the
Yang--Mills, theta, and anisotropic commutant coordinates; derive the
exact derivative of their calibrated update; and identify a
verifiable condition preventing stable-to-marginal leakage from
destroying the \(b^{-2}\) margin.
\end{decision}

\begin{firewall}[Claim boundary after eighth-series V4]
\label{fire:e8v4-claim}
V4 does not prove \(D_{\rm geom}<b^2\), rotational restoration,
asymptotic freedom with controlled finite remainders, global chart
and saturated-bad closure, NCGC/CFTC along the ultraviolet
trajectory, terminal entry, regulator independence, continuum OS
reconstruction, or the full Yang--Mills mass gap.
\end{firewall}

\paragraph{Parent-version record.}
V4 was formed from
\texttt{Renewal\_Yang\_Mills\_Mass\_Gap\_Research\_Notes\_V3.tex}.
The V3 parent had \(1\,807\) counted source lines,
\(65\,686\) bytes, and SHA--256
\[
 \texttt{B85EDE3EF4DD8F0944EF18C1CF775AB69CF41BEC3B0B432C9ACFA8D1F15E6163}.
\]
The compulsory companion audit is due in V5.

% =============================================================================
% END APPENDED EIGHTH-SERIES V4 MATERIAL
% =============================================================================

% =============================================================================
% BEGIN APPENDED EIGHTH-SERIES V5 MATERIAL
% =============================================================================

\part{Eighth-series V5: companion audit and calibrated marginal
Schur transport}
\label{part:e8v5}

\section{Compulsory companion-program audit}
\label{sec:e8v5-audit}

\begin{audit}[Files inspected at the V5 boundary]
\label{aud:e8v5-files}
The compulsory audit used the following then-current active files.
\begin{enumerate}[label=\textup{(\roman*)},leftmargin=3.5em]
\item
\texttt{Renewal\_SM\_Einstein\_Continuum\_Research\_Notes\_V76.tex}:
\(29\,210\) counted source lines, \(1\,053\,839\) bytes, SHA--256
\[
\texttt{730482CD8F89352F7EC9E4DB3E68ACD5722553C0A85B89D967EE1F339AC5C7DC}.
\]
\item
\texttt{Renewal\_NS\_Stability\_Research\_Notes\_V26.tex}:
\(5\,319\) counted source lines, \(278\,283\) bytes, SHA--256
\[
\texttt{82C887F8C2C69E4F28FAD7C3DC8DE5B9099CB5A4BF431EBA1FFF3E91935978AD}.
\]
\end{enumerate}
\end{audit}

\begin{decision}[Audit transfer]
\label{dec:e8v5-transfer}
The SM--Einstein file contributes three structural warnings and one
positive construction:
\begin{enumerate}[label=\textup{(\alph*)},leftmargin=3.5em]
\item retained pullbacks are exact unit directions, so only an aligned
      stable sector can contract;
\item the minimum-energy lift gives the correct Gaussian
      retained/detail split;
\item a nonstationary reduction has a Schur complement plus an exact
      reaction debit, so marginal calibration must precede a small
      stable-defect claim; and
\item the derivative on an exact marginal level set is a Schur
      complement of the full RG derivative.
\end{enumerate}
The SM transfer had named only the Yang--Mills and theta source
coefficients.  V4 sharpens it: at finite cutoff the entire
\({\mathcal C}_\Gamma\) must be calibrated.  The NS file is computing
rank-one Oseen-kernel derivative constants; it supplies no
gauge-theory estimate and changes no YM direction.
\end{decision}

\section{Exact coarse derivatives}
\label{sec:e8v5-derivatives}

Let \(B:\Omega\to\Sigma\) be the fixed physical block map and
\(\kappa(d\omega\mid y)\) a reference probability on its fibre.  For
an interaction \(S\), define
\[
 ({\cal T}S)(y)
 :=-\log\int e^{-S(\omega)}\kappa(d\omega\mid y),     \label{eq:e8v5-T}
\]
whenever the integral is finite and nonzero.  Write
\(\langle\cdot\rangle_{S,y}\) for the corresponding normalized
conditional expectation.

\begin{theorem}[Conditional cumulant differential]
\label{thm:e8v5-cumulant}
On a finite regulator, if \(S,V,W\) are bounded, then
\begin{align}
 D{\cal T}_S[V](y)
 &=\langle V\rangle_{S,y},                            \label{eq:e8v5-DT}\\
 D^2{\cal T}_S[V,W](y)
 &=-\operatorname{Cov}_{S,y}(V,W).                   \label{eq:e8v5-D2T}
\end{align}
The same formulas hold in any Banach chart in which differentiation
under the fibre integral is dominated.
\end{theorem}

\begin{proof}
Put \(Z(t,u)=\int e^{-S-tV-uW}\,d\kappa_y\).
Boundedness permits differentiation under the finite measure
integral.  Differentiating \(-\log Z(t,0)\) at zero gives
\eqref{eq:e8v5-DT}.  Differentiating the normalized expectation of
\(V\) in the \(W\) direction gives
\[
 -\langle VW\rangle_{S,y}
 +\langle V\rangle_{S,y}\langle W\rangle_{S,y},
\]
which is \eqref{eq:e8v5-D2T}.  The dominated Banach statement is the
same calculation.
\end{proof}

\begin{corollary}[Retained pullback unit sector]
\label{cor:e8v5-pullback}
For every coarse interaction \(R\),
\[
 {\cal T}(S+R\circ B)={\cal T}S+R,
 \qquad D{\cal T}_S[R\circ B]=R.                     \label{eq:e8v5-pullback}
\]
Thus no complete interaction norm which identifies the fine pullback
with \(R\) can give \(\|D{\cal T}_S\|<1\).
\end{corollary}

\begin{proof}
On the fibre \(B\omega=y\), \(R(B\omega)=R(y)\).  Factor
\(e^{-R(y)}\) out of \eqref{eq:e8v5-T} and differentiate.
\end{proof}

\section{The four derivative blocks}
\label{sec:e8v5-blocks}

\begin{definition}[Marginal coordinate charts]
\label{def:e8v5-charts}
At the input and output scales let
\[
 \iota:\mathbb R^r\to{\mathfrak B},\quad
 \pi:{\mathfrak B}\to\mathbb R^r,\qquad
 \iota_+:\mathbb R^r\to{\mathfrak B}_+,\quad
 \pi_+:{\mathfrak B}_+\to\mathbb R^r                \label{eq:e8v5-charts}
\]
be bounded maps with
\(\pi\iota=\pi_+\iota_+=I\).  Put
\[
 M=\iota\pi,\quad Q=I-M,\qquad
 M_+=\iota_+\pi_+,\quad Q_+=I-M_+.                  \label{eq:e8v5-MQ}
\]
The range of \(\iota\) contains the complete finite-cutoff marginal
ledger, in particular \({\mathcal C}_\Gamma\).
\end{definition}

At \(S=\iota m+q\), \(q\in\operatorname{ran}Q\), define
\begin{align}
 A&:=\pi_+D{\cal T}_S\iota,
 &B_{\rm sm}&:=\pi_+D{\cal T}_S Q,                  \label{eq:e8v5-AB}\\
 C&:=Q_+D{\cal T}_S Q,
 &D_{\rm ms}&:=Q_+D{\cal T}_S\iota.                \label{eq:e8v5-CD}
\end{align}
Here \(B_{\rm sm}\) is stable-to-marginal leakage and
\(D_{\rm ms}\) is marginal-to-stable leakage.  The subscripts prevent
confusion with the physical block map \(B\).

\begin{proposition}[Exact pullback coordinates]
\label{prop:e8v5-exact-coordinate}
If
\[
 \iota u=(\iota_+u)\circ B\qquad(u\in\mathbb R^r),   \label{eq:e8v5-exact-basis}
\]
then \(A=I\) and \(D_{\rm ms}=0\) at every \(S\).
\end{proposition}

\begin{proof}
\Cref{cor:e8v5-pullback} gives
\(D{\cal T}_S\iota u=\iota_+u\).  Apply \(\pi_+\) and \(Q_+\).
\end{proof}

\begin{firewall}[Generated owners are not automatically pullbacks]
\label{fire:e8v5-generated}
A local representative of a running gauge kinetic coefficient may
change with lattice spacing, interpolation, or normalization.
Equation \eqref{eq:e8v5-exact-basis} must be checked; it cannot be
inferred from equal canonical dimension.  The general matrices
\(A\) and \(D_{\rm ms}\) are retained below.
\end{firewall}

\section{Exact marginal alignment}
\label{sec:e8v5-alignment}

\begin{theorem}[Marginal level-set implicit chart]
\label{thm:e8v5-IFT}
Assume \({\cal T}\) is \(C^1\) near
\(S_0=\iota m_0+q_0\), and that
\[
 A_0=\pi_+D{\cal T}_{S_0}\iota                    \label{eq:e8v5-A0}
\]
is invertible.  Put
\(m_{+,0}=\pi_+{\cal T}(S_0)\).
There are neighbourhoods of \(q_0\) and \(m_{+,0}\) and a unique
\(C^1\) function
\[
 m={\mathfrak k}(q;m_+)                             \label{eq:e8v5-kappa}
\]
such that
\[
 \pi_+{\cal T}\bigl(\iota{\mathfrak k}(q;m_+)+q\bigr)
 =m_+.                                               \label{eq:e8v5-level}
\]
At the base point,
\[
 D_q{\mathfrak k}=-A_0^{-1}B_{{\rm sm},0},
 \qquad
 D_{m_+}{\mathfrak k}=A_0^{-1}.                     \label{eq:e8v5-kappa-derivative}
\]
\end{theorem}

\begin{proof}
Apply the finite-dimensional implicit-function theorem to
\[
 F(m,q,m_+):=\pi_+{\cal T}(\iota m+q)-m_+.
\]
Its \(m\)-derivative at the base point is \(A_0\).  Differentiating
\(F({\mathfrak k}(q;m_+),q,m_+)=0\) gives
\eqref{eq:e8v5-kappa-derivative}.
\end{proof}

\begin{definition}[Aligned stable RG map]
\label{def:e8v5-aligned}
On the chart of \Cref{thm:e8v5-IFT}, set
\[
 \Phi_{m_+}(q):=
 Q_+{\cal T}\bigl(\iota{\mathfrak k}(q;m_+)+q\bigr). \label{eq:e8v5-Phi}
\]
\end{definition}

\begin{theorem}[Marginal Schur differential]
\label{thm:e8v5-Schur}
At every point where \(A\) is invertible,
\[
 D_q\Phi_{m_+}
 =C-D_{\rm ms}A^{-1}B_{\rm sm}.                     \label{eq:e8v5-Schur}
\]
Therefore the correct derivative between two actions having the same
output marginal coordinate is a Schur complement of the complete RG
derivative.
\end{theorem}

\begin{proof}
Differentiate \eqref{eq:e8v5-Phi} at fixed \(m_+\).  The derivative
of \eqref{eq:e8v5-level} gives
\(D_q{\mathfrak k}=-A^{-1}B_{\rm sm}\).  Hence
\[
 D_q\Phi_{m_+}
 =Q_+D{\cal T}_S\bigl(Q+\iota D_q{\mathfrak k}\bigr)
 =C-D_{\rm ms}A^{-1}B_{\rm sm}.
\]
\end{proof}

\begin{corollary}[Verifiable leakage gate]
\label{cor:e8v5-leakage-gate}
If on a convex aligned chart
\begin{align}
 \|C\|&\le\kappa_0,&
 \|A^{-1}\|&\le\alpha,&
 \|B_{\rm sm}\|&\le\beta,&
 \|D_{\rm ms}\|&\le\delta                           \label{eq:e8v5-four-bounds}
\end{align}
and
\[
 \kappa_{\rm al}:=\kappa_0+\alpha\beta\delta<1,      \label{eq:e8v5-gate}
\]
then
\[
 \|\Phi_{m_+}(q)-\Phi_{m_+}(q')\|
 \le\kappa_{\rm al}\|q-q'\|                         \label{eq:e8v5-aligned-contraction}
\]
whenever the segment between \(q,q'\) remains in that chart.
\end{corollary}

\begin{proof}
\Cref{thm:e8v5-Schur} and submultiplicativity give
\(\|D_q\Phi_{m_+}\|\le\kappa_{\rm al}\).  Integrate the derivative
along the segment.
\end{proof}

\begin{corollary}[Near-identity marginal transport]
\label{cor:e8v5-neumann}
If \(\|A-I\|\le a<1\), then
\[
 \|A^{-1}\|\le\frac1{1-a},\qquad
 \|D_q\Phi_{m_+}\|
 \le\kappa_0+\frac{\beta\delta}{1-a}.                \label{eq:e8v5-neumann}
\]
\end{corollary}

\begin{proof}
Use the Neumann series for \(A^{-1}\) and
\Cref{thm:e8v5-Schur}.
\end{proof}

\section{Why centred leakage is second order}
\label{sec:e8v5-centred}

\begin{proposition}[Covariance variation of a leakage block]
\label{prop:e8v5-covariance-leakage}
Let \(S_t=S_0+tq\).  For a stable insertion \(v\),
\[
 \frac d{dt}\,
 \pi_+\langle v\rangle_{S_t,\cdot}
 =-\pi_+\operatorname{Cov}_{S_t,\cdot}(v,q).         \label{eq:e8v5-B-variation}
\]
For a marginal insertion \(\iota u\),
\[
 \frac d{dt}\,
 Q_+\langle\iota u\rangle_{S_t,\cdot}
 =-Q_+\operatorname{Cov}_{S_t,\cdot}(\iota u,q).     \label{eq:e8v5-D-variation}
\]
\end{proposition}

\begin{proof}
Apply \eqref{eq:e8v5-D2T} to the respective insertion and \(q\), then
apply the bounded output projection.
\end{proof}

\begin{theorem}[Quadratic calibrated leakage debit]
\label{thm:e8v5-quadratic-debit}
Assume at \(S_0=\iota m_0\) that
\[
 B_{\rm sm}(0)=0,\qquad D_{\rm ms}(0)=0.             \label{eq:e8v5-centred-zero}
\]
Suppose, for \(\|q\|\le r\),
\begin{align}
 \|\pi_+\operatorname{Cov}_{S_t}(v,q)\|
 &\le L_B\|v\|\,\|q\|,                              \label{eq:e8v5-covB}\\
 \|Q_+\operatorname{Cov}_{S_t}(\iota u,q)\|
 &\le L_D\|u\|\,\|q\|                               \label{eq:e8v5-covD}
\end{align}
uniformly for \(0\le t\le1\).  Then
\[
 \|B_{\rm sm}(q)\|\le L_B\|q\|,\qquad
 \|D_{\rm ms}(q)\|\le L_D\|q\|,                     \label{eq:e8v5-linear-leakage}
\]
and, if \(\|A(q)^{-1}\|\le\alpha\),
\[
 \|D_{\rm ms}A^{-1}B_{\rm sm}\|
 \le\alpha L_BL_D\|q\|^2.                           \label{eq:e8v5-quadratic-leakage}
\]
\end{theorem}

\begin{proof}
Integrate \eqref{eq:e8v5-B-variation} and
\eqref{eq:e8v5-D-variation} from zero to one and use
\eqref{eq:e8v5-centred-zero}.  Multiply the resulting operator
bounds with the inverse bound.
\end{proof}

\begin{corollary}[YM contraction budget]
\label{cor:e8v5-YM-budget}
Suppose the complete compact-link estimates furnish
\[
 \|C\|\le D_{\rm geom}b^{-2}+\varepsilon_{\rm nl}(r) \label{eq:e8v5-Cbudget}
\]
on a stable ball of radius \(r\), and the hypotheses of
\Cref{thm:e8v5-quadratic-debit} hold.  The aligned stable map is a
strict contraction whenever
\[
 D_{\rm geom}b^{-2}+\varepsilon_{\rm nl}(r)
 +\alpha L_BL_Dr^2<1.                               \label{eq:e8v5-YM-gate}
\]
\end{corollary}

\begin{proof}
Use \Cref{cor:e8v5-leakage-gate} with
\eqref{eq:e8v5-quadratic-leakage}.
\end{proof}

\begin{firewall}[Centred Gaussian transfer is only a base point]
\label{fire:e8v5-centred}
The minimum-energy Gaussian split supplies
\eqref{eq:e8v5-centred-zero} only for insertions actually orthogonal
under that split and for projections matched to it.  Interacting
conditional cumulants generate the leakage in
\eqref{eq:e8v5-B-variation}--\eqref{eq:e8v5-D-variation}; they must be
bounded in the complete support norm.
\end{firewall}

\section{The full finite-cutoff marginal matrix}
\label{sec:e8v5-commutant}

\begin{definition}[Physical and anisotropic coordinates]
\label{def:e8v5-physical-coordinates}
Equip \({\mathcal C}_\Gamma\) with the Hilbert--Schmidt inner product.
When the indicated tensors belong to the commutant, choose normalized
vectors \(E_{\rm YM}\) and \(E_\theta\) representing
\(\operatorname{Tr}F^2\) and
\(\operatorname{Tr}(F\widetilde F)\).  Define
\[
 {\mathcal C}_{\rm phys}:=
 \operatorname{span}\{E_{\rm YM},E_\theta\},\qquad
 {\mathcal C}_{\rm an}:=
 {\mathcal C}_\Gamma\ominus{\mathcal C}_{\rm phys}. \label{eq:e8v5-Csplit}
\]
If one of the two displayed tensors is absent or linearly dependent
in the regulated representation, it is omitted.  A marginal vector
is written
\[
 m=(g_{\rm kin},\theta,a),\qquad
 a\in{\mathcal C}_{\rm an}.                          \label{eq:e8v5-mcoords}
\]
\end{definition}

\begin{proposition}[Parity block diagonalization]
\label{prop:e8v5-parity}
Suppose an orientation-reversing lattice reflection induces an
involution \(P\) on the marginal coordinates, the exact RG is
equivariant, and the base interaction and stable perturbation are
\(P\)-even.  Then the marginal transport matrix \(A\) commutes with
\(P\).  Hence its even-to-odd and odd-to-even blocks vanish.  In
particular, when \(E_\theta\) is odd and all other selected
coordinates are even, theta does not mix linearly with them at
\(\theta=0\).
\end{proposition}

\begin{proof}
Equivariance gives
\(\beta(Pm)=P\beta(m)\) for the marginal map
\(\beta=\pi_+{\cal T}\iota\) on the even stable background.
Differentiate at a fixed point of \(P\) to obtain \(AP=PA\).
The eigenspaces of an involution are complementary, so the cross
blocks vanish.
\end{proof}

\begin{firewall}[Reflection and theta]
\label{fire:e8v5-theta}
For the reflection-positive \(\theta=0\) programme, the odd
coordinate is fixed at zero.  A nonzero Euclidean theta weight can be
complex and is not covered by the positive-measure estimates used
here.  Recording \(E_\theta\) as a marginal owner is not a proof of
reflection positivity at nonzero theta.
\end{firewall}

\begin{theorem}[Finite-cutoff marginal calibration rule]
\label{thm:e8v5-full-calibration}
At every finite cutoff, the marginal matrix in
\Cref{thm:e8v5-IFT,thm:e8v5-Schur} must be formed on all of
\({\mathcal C}_\Gamma\), including \({\mathcal C}_{\rm an}\).
If this full \(A\) is invertible and
\eqref{eq:e8v5-YM-gate} holds, the anisotropic coefficients may be
fixed by the same exact level-set equation as the physical
coefficients and the aligned stable map contracts.  This conclusion
does not imply that \(a\to0\) in the continuum limit.
\end{theorem}

\begin{proof}
The implicit-function and Schur proofs use every noncontracting
coordinate.  Omitting a component of
\({\mathcal C}_{\rm an}\) assigns a zero-momentum quadratic direction
to \(Q\), contradicting the \(O(k^2)\) hypothesis used for the V3
quadratic contraction.  With the full coordinate space,
\Cref{thm:e8v5-IFT,cor:e8v5-YM-budget} apply.  Neither theorem contains
a scale-decay estimate for the selected level \(a\), proving the last
statement.
\end{proof}

\begin{firewall}[Invertibility is a substantive RG gate]
\label{fire:e8v5-invertibility}
Canonical marginality does not imply that \(A\) is invertible or
near the identity.  A small singular value makes the leakage debit
large through \(A^{-1}\).  The compact-link programme must derive
\(\|A-I\|<1\), or another quantitative inverse bound, from the exact
conditional cumulants and the selected renormalization conditions.
\end{firewall}

\section{Status and next acquisition}
\label{sec:e8v5-status}

\begin{theorem}[What eighth-series V5 proves]
\label{thm:e8v5-result}
V5 constructs the exact local chart on a prescribed output marginal
level set and proves that its stable derivative is
\[
 C-D_{\rm ms}A^{-1}B_{\rm sm}.
\]
It supplies the explicit strict gate
\eqref{eq:e8v5-YM-gate}.  When both leakage blocks vanish at the
energy-orthogonal Gaussian reference and their interacting variations
obey complete covariance bounds, the calibration debit is quadratic
in the stable radius.  The calibration includes the whole
finite-cutoff commutant, while parity decouples the theta coordinate
at a reflection-even \(\theta=0\) background.
\end{theorem}

\begin{proof}
Combine
\Cref{thm:e8v5-cumulant,thm:e8v5-IFT,thm:e8v5-Schur,
thm:e8v5-quadratic-debit,cor:e8v5-YM-budget,
prop:e8v5-parity,thm:e8v5-full-calibration}.
\end{proof}

\begin{assessment}[Progress at eighth-series V5]
\label{ass:e8v5-status}
The stable-to-marginal leakage bottleneck is no longer qualitative:
it is isolated in a finite-dimensional inverse and two complete
conditional-covariance rows.  What remains is model-specific.  One
must prove the bounds \(a,\varepsilon_{\rm nl},L_B,L_D\) and
\(D_{\rm geom}\) for the compact Wilson fibre with good/bad
decomposition, and separately prove rotational restoration of the
anisotropic marginal trajectory.
\end{assessment}

\begin{decision}[V6 acquisition order]
\label{dec:e8v5-V6}
V6 will work upstream on the conditional covariance rows.  It will
derive a finite-volume covariance interpolation identity around the
energy-orthogonal Gaussian detail, split polynomial and analytic
tails, and state a support-resolved criterion yielding
\(L_B,L_D\) without volume growth.  If the covariance row cannot be
closed directly, the next move is to the local Hessian/Dobrushin
margin rather than assuming clustering.
\end{decision}

\begin{firewall}[Claim boundary after eighth-series V5]
\label{fire:e8v5-claim}
V5 does not prove the compact-link covariance bounds, the inverse
bound for \(A\), \(D_{\rm geom}<b^2\), rotational restoration,
asymptotic freedom with controlled finite remainders, global chart
and saturated-bad closure, NCGC/CFTC, terminal entry, regulator
independence, continuum OS reconstruction, or the full Yang--Mills
mass gap.
\end{firewall}

\paragraph{Parent-version record.}
V5 was formed from
\texttt{Renewal\_Yang\_Mills\_Mass\_Gap\_Research\_Notes\_V4.tex}.
The V4 parent had \(2\,328\) counted source lines,
\(84\,959\) bytes, and SHA--256
\[
\texttt{0AEF411E8027F94BDFACCA24D23422D7315BE2EAEA97DBDF43F7AD27910E1000}.
\]
The next compulsory companion audit is V10.

% =============================================================================
% END APPENDED EIGHTH-SERIES V5 MATERIAL
% =============================================================================

% =============================================================================
% BEGIN APPENDED EIGHTH-SERIES V6 MATERIAL
% =============================================================================

\part{Eighth-series V6: conditional covariance rows from finite
influence}
\label{part:e8v6}

\section{Exact interpolation around the orthogonal detail reference}
\label{sec:e8v6-interpolation}

Fix a retained field \(y\).  Let \(\gamma_y\) be the normalized
energy-orthogonal Gaussian detail reference supplied by the
minimum-energy lift, and let \(V_y\) be a bounded interacting detail
potential.  For \(0\le t\le1\), put
\[
 d\mu_{t,y}(z):=
 \frac{e^{-tV_y(z)}}{\int e^{-tV_y}\,d\gamma_y}\,
 d\gamma_y(z).                                      \label{eq:e8v6-mut}
\]

\begin{theorem}[Expectation and covariance interpolation]
\label{thm:e8v6-interpolation}
For bounded observables \(F,G\),
\begin{align}
 \frac d{dt}\langle F\rangle_{t,y}
 &=-\operatorname{Cov}_{t,y}(F,V_y),                 \label{eq:e8v6-expectation-flow}\\
 \frac d{dt}\operatorname{Cov}_{t,y}(F,G)
 &=-\kappa_{3,t,y}(F,G,V_y),                         \label{eq:e8v6-covariance-flow}
\end{align}
where
\[
 \kappa_3(F,G,V)
 :=\langle FGV\rangle-\langle FG\rangle\langle V\rangle
 -\langle FV\rangle\langle G\rangle
 -\langle GV\rangle\langle F\rangle
 +2\langle F\rangle\langle G\rangle\langle V\rangle. \label{eq:e8v6-kappa3}
\]
Consequently,
\[
 \langle F\rangle_{1,y}
 =\langle F\rangle_{0,y}
  -\int_0^1\operatorname{Cov}_{t,y}(F,V_y)\,dt.       \label{eq:e8v6-Duhamel}
\]
\end{theorem}

\begin{proof}
Differentiate the quotient defining expectation in
\eqref{eq:e8v6-mut}; the numerator contributes
\(-\langle FV_y\rangle\) and the normalization contributes
\(\langle F\rangle\langle V_y\rangle\).  This proves
\eqref{eq:e8v6-expectation-flow}.  Apply that identity to
\(\langle FG\rangle-\langle F\rangle\langle G\rangle\) and collect
terms to obtain \eqref{eq:e8v6-covariance-flow}.  Integration gives
\eqref{eq:e8v6-Duhamel}.
\end{proof}

\begin{firewall}[No perturbative truncation]
\label{fire:e8v6-interpolation}
Equation \eqref{eq:e8v6-Duhamel} is exact.  V6 bounds its covariance
uniformly in \(t\); it does not discard higher cumulants or identify
the interacting measure with the Gaussian reference.
\end{firewall}

\section{Conditional influence on a finite fibre}
\label{sec:e8v6-influence}

Write the detail fibre as a product
\(\mathsf Z=\prod_{x\in\Lambda}\mathsf Z_x\), where one site may be an
entire fixed-size elementary detail block.  Let
\(\mu\) be a strictly positive probability measure on this finite
product.  Denote its one-site conditional kernel by
\(\mu_x(\,\cdot\mid z_{\Lambda\setminus\{x\}})\).
For a bounded function \(f\), set
\[
 \delta_x f:=
 \sup\{|f(z)-f(z')|:z_{\Lambda\setminus\{x\}}
                         =z'_{\Lambda\setminus\{x\}}\}.              \label{eq:e8v6-delta}
\]
Use total variation
\(\|\nu-\nu'\|_{\rm TV}:=\sup_A|\nu(A)-\nu'(A)|\).

\begin{definition}[Dobrushin influence matrix]
\label{def:e8v6-Dobrushin}
For \(x\ne y\), let
\[
 c_{xy}:=\sup_{\substack{z,z':\\
 z_{\Lambda\setminus\{y\}}=z'_{\Lambda\setminus\{y\}}}}
 \|\mu_x(\,\cdot\mid z_{\ne x})
  -\mu_x(\,\cdot\mid z'_{\ne x})\|_{\rm TV},
 \qquad c_{xx}=0.                                    \label{eq:e8v6-cxy}
\]
For a block metric \(d\) and \(\mu_0\ge0\), define
\[
 c_{\mu_0}:=
 \sup_x\sum_y e^{\mu_0d(x,y)}c_{xy}.                 \label{eq:e8v6-cmu}
\]
\end{definition}

\begin{lemma}[Weighted resolvent]
\label{lem:e8v6-resolvent}
If \(c_{\mu_0}<1\), then
\[
 {\mathsf D}:=(I-C)^{-1}=\sum_{n=0}^\infty C^n       \label{eq:e8v6-D}
\]
has nonnegative entries and
\[
 \sup_x\sum_y e^{\mu_0d(x,y)}{\mathsf D}_{xy}
 \le\frac1{1-c_{\mu_0}}.                             \label{eq:e8v6-Drow}
\]
\end{lemma}

\begin{proof}
The triangle inequality for \(d\) shows that the weighted row norm is
submultiplicative.  Hence
\(\|C^n\|_{\mu_0}\le c_{\mu_0}^n\).  The Neumann series converges and
summing the geometric series proves \eqref{eq:e8v6-Drow}.
\end{proof}

\begin{theorem}[Finite-volume Dobrushin covariance comparison]
\label{thm:e8v6-Dobrushin-covariance}
If \(c_{\mu_0}<1\), then for bounded \(f,g\),
\[
 |\operatorname{Cov}_{\mu}(f,g)|
 \le
 \sum_{x,y}\delta_xf\,{\mathsf D}_{xy}\,\delta_yg.    \label{eq:e8v6-covariance-comparison}
\]
The constant and the weighted row bound are independent of
\(|\Lambda|\).
\end{theorem}

\begin{proof}
We include the finite-volume comparison argument.  Perturb \(\mu\) by
\(e^{sg}\) and let \(b_x(s)\) be the maximal first-order change of an
\(x\)-conditional expectation under this perturbation and arbitrary
boundary changes.  Direct one-site conditioning gives the influence
inequality
\[
 b_x(0)\le\delta_xg+\sum_yc_{xy}b_y(0).              \label{eq:e8v6-influence-recursion}
\]
Indeed, the direct density derivative at site \(x\) is a one-site
covariance bounded by \(\delta_xg\); changing a boundary coordinate
\(y\) changes the conditional law by at most \(c_{xy}\), and a
telescoping over changed coordinates gives the sum.  Iterating
\eqref{eq:e8v6-influence-recursion} from zero yields the minimal
nonnegative bound
\[
 b_x(0)\le\sum_y{\mathsf D}_{xy}\delta_yg.
\]
Reveal the sites successively and telescope the conditional
expectations of \(f\).  The change at site \(x\) is at most
\(\delta_xf\); hence
\[
 \left|\frac d{ds}\langle f\rangle_{\mu_s}\right|_{s=0}
 \le\sum_x\delta_xf\,b_x(0).
\]
The derivative on the left is
\(\operatorname{Cov}_\mu(f,g)\), proving
\eqref{eq:e8v6-covariance-comparison}.  The last assertion is
\Cref{lem:e8v6-resolvent}.
\end{proof}

\begin{remark}[Conservative normalization]
\label{rem:e8v6-constant}
Sharper versions of \eqref{eq:e8v6-covariance-comparison} contain a
factor \(1/4\) for the oscillation convention
\eqref{eq:e8v6-delta}.  V6 deliberately uses the larger constant
one, so no later gate depends on that normalization.
\end{remark}

\section{Computing influence from mixed oscillations}
\label{sec:e8v6-mixed}

\begin{lemma}[Log-density oscillation controls total variation]
\label{lem:e8v6-log-TV}
Let \(p,q\) be strictly positive densities relative to one measure and
suppose
\[
 \operatorname{osc}\log(p/q)\le a.                  \label{eq:e8v6-logosc}
\]
Then
\[
 \|p-q\|_{\rm TV}\le\min\{1,e^a-1\}.                \label{eq:e8v6-TV}
\]
\end{lemma}

\begin{proof}
Let \(r=p/q\).  Its essential maximum divided by its essential
minimum is at most \(e^a\), while \(\int r\,dq=1\); hence
\(e^{-a}\le r\le e^a\).  Therefore
\[
 \|p-q\|_{\rm TV}
 =\tfrac12\int|r-1|\,dq
 \le\tfrac12(e^a-1)\le e^a-1.
\]
The trivial total-variation bound is one.
\end{proof}

\begin{definition}[Mixed site oscillation]
\label{def:e8v6-mixed}
For a potential \(U\), define
\[
 a_{xy}(U):=
 \sup|U(z)-U(z^x)-U(z^y)+U(z^{xy})|,                \label{eq:e8v6-axy}
\]
where the four configurations have independently chosen \(x\)- and
\(y\)-coordinates and agree elsewhere.
\end{definition}

\begin{proposition}[Product-reference influence test]
\label{prop:e8v6-product-test}
Let
\[
 d\mu(z)=Z^{-1}e^{-U(z)}\prod_x\nu_x(dz_x).          \label{eq:e8v6-product-density}
\]
Then
\[
 c_{xy}\le\min\{1,e^{a_{xy}(U)}-1\}.                \label{eq:e8v6-c-from-a}
\]
If \(a_*:=\sup_{x,y}a_{xy}(U)\), then
\[
 c_{\mu_0}\le e^{a_*}
 \sup_x\sum_y e^{\mu_0d(x,y)}a_{xy}(U).             \label{eq:e8v6-c-row-test}
\]
\end{proposition}

\begin{proof}
For two backgrounds differing only at \(y\), the log ratio of the two
\(x\)-conditional densities, as a function of \(z_x\), has
oscillation at most \(a_{xy}(U)\).  Apply
\Cref{lem:e8v6-log-TV}.  Finally use
\(e^a-1\le e^{a_*}a\) for \(0\le a\le a_*\), multiply by the weight,
and sum.
\end{proof}

\begin{corollary}[Polynomial/tail influence split]
\label{cor:e8v6-poly-tail}
Suppose on the sitewise good domain
\[
 U=\Psi_{\rm st}+{\cal T}_{5,g}                     \label{eq:e8v6-U-split}
\]
and the complete mixed-oscillation estimates give
\begin{align}
 \sup_x\sum_y e^{\mu_0d(x,y)}a_{xy}(\Psi_{\rm st})
 &\le K_{\rm st}\|\Psi_{\rm st}\|_{{\rm st},\mu,R},
                                                               \label{eq:e8v6-a-poly}\\
 \sup_x\sum_y e^{\mu_0d(x,y)}a_{xy}({\cal T}_{5,g})
 &\le K_{\rm tail}g^3R^5.                            \label{eq:e8v6-a-tail}
\end{align}
Then
\[
 c_{\mu_0}
 \le e^{a_*}\left(
 K_{\rm st}\|\Psi_{\rm st}\|_{{\rm st},\mu,R}
 +K_{\rm tail}g^3R^5\right).                        \label{eq:e8v6-c-poly-tail}
\]
In particular, strict smallness of the right side proves the
volume-uniform Dobrushin gate.
\end{corollary}

\begin{proof}
Mixed oscillations are subadditive, so
\(a_{xy}(U)\le a_{xy}(\Psi_{\rm st})
+a_{xy}({\cal T}_{5,g})\).  Apply
\Cref{prop:e8v6-product-test} and the two assumed complete-row bounds.
\end{proof}

\begin{firewall}[The reference product must be literal]
\label{fire:e8v6-product}
Energy orthogonality separates retained and detail Gaussian
variables; it does not by itself make different detail sites
independent.  Proposition~\ref{prop:e8v6-product-test} applies after
grouping variables so that the remaining reference is a literal
product.  If the Gaussian reference still couples those groups, its
conditional influence matrix must be included in \(C\).
\end{firewall}

\section{From local covariances to complete owner rows}
\label{sec:e8v6-owner-row}

\begin{definition}[Oscillation owner norm]
\label{def:e8v6-owner-norm}
For an interaction \(F=\{F_X\}\), define
\[
 \|F\|_{\delta,\mu_0}
 :=
 \sup_z\sum_{X\ni z}e^{\mu_0\operatorname{diam}X}
 \sum_{x\in X}\delta_xF_X.                           \label{eq:e8v6-owner-norm}
\]
For interactions \(F,G\), assign
\(\operatorname{Cov}(F_X,G_Y)\) to the union owner \(X\cup Y\).
\end{definition}

\begin{theorem}[Complete covariance row]
\label{thm:e8v6-complete-covariance}
Under \(c_{\mu_0}<1\), the union-owner covariance interaction obeys
\[
 \|\operatorname{Cov}_\mu(F,G)\|_{{\rm row},\mu_0}
 \le\frac{2}{1-c_{\mu_0}}\,
 \|F\|_{\delta,\mu_0}\|G\|_{\delta,\mu_0}.           \label{eq:e8v6-complete-covariance}
\]
The estimate is uniform in the number of sites.
\end{theorem}

\begin{proof}
Apply \Cref{thm:e8v6-Dobrushin-covariance} to each pair
\((F_X,G_Y)\).  For \(x\in X,y\in Y\),
\[
 \operatorname{diam}(X\cup Y)
 \le\operatorname{diam}X+d(x,y)+\operatorname{diam}Y.
\]
Split the outer incidence row according to whether its marked site
lies in \(X\) or in \(Y\).  In each half, use the displayed weight
inequality, sum first over the second owner, and apply
\eqref{eq:e8v6-Drow}.  The two incidence choices give the factor two;
the remaining sums are exactly the two norms in
\eqref{eq:e8v6-owner-norm}.
\end{proof}

\begin{corollary}[Explicit V5 leakage constants]
\label{cor:e8v6-LBLD}
Assume uniformly in \(y,t\):
\begin{enumerate}[label=\textup{(\roman*)},leftmargin=3.5em]
\item \(c_{\mu_0}(\mu_{t,y})\le c<1\);
\item \(\|v\|_{\delta,\mu_0}\le J_{\rm st}\|v\|_{\rm st}\) for stable
      insertions;
\item \(\|\iota u\|_{\delta,\mu_0}\le J_{\rm m}|u|\);
\item the output projections have owner-row bounds \(J_\pi,J_Q\);
      and
\item the base leakage blocks vanish.
\end{enumerate}
Then the constants of
\Cref{thm:e8v5-quadratic-debit} may be chosen as
\[
 L_B=\frac{2J_\pi J_{\rm st}^2}{1-c},
 \qquad
 L_D=\frac{2J_QJ_{\rm m}J_{\rm st}}{1-c}.            \label{eq:e8v6-LBLD}
\]
\end{corollary}

\begin{proof}
Apply the exact interpolation
\eqref{eq:e8v6-Duhamel} to a stable insertion and to a marginal
insertion.  Bound their covariances with the interacting stable
potential by \Cref{thm:e8v6-complete-covariance}; then apply the
projection and embedding bounds.  Comparison with
\eqref{eq:e8v5-covB}--\eqref{eq:e8v5-covD} gives
\eqref{eq:e8v6-LBLD}.
\end{proof}

\begin{definition}[Dobrushin interpolation card]
\label{def:e8v6-card}
DIC consists of:
\begin{enumerate}[label=\textup{(D\arabic*)},leftmargin=4em]
\item a sitewise product decomposition of the conditional detail
      fibre, or a declared base influence matrix;
\item uniform positivity of every one-site conditional kernel;
\item a uniform weighted influence bound \(c<1\) for every
      interpolation measure \(\mu_{t,y}\);
\item the V3 stable-to-oscillation embedding;
\item the V1/V2 complete analytic-tail mixed row;
\item union support ownership and bounded output projections;
\item exact Gaussian base centering for the two leakage blocks; and
\item separate ownership of non-sitewise bad-domain constraints.
\end{enumerate}
\end{definition}

\begin{theorem}[DIC closes the abstract leakage row]
\label{thm:e8v6-card}
Under DIC, \(L_B,L_D\) are finite uniformly in the total lattice
volume and are given by \eqref{eq:e8v6-LBLD}.  Inserting those values
in \eqref{eq:e8v5-YM-gate} yields a fully explicit sufficient
condition for strict calibrated stable contraction.
\end{theorem}

\begin{proof}
Items (D1)--(D3) give
\Cref{thm:e8v6-Dobrushin-covariance}.  Items (D4)--(D6) give
\Cref{thm:e8v6-complete-covariance}; (D7) permits the interpolation
from zero leakage.  Apply \Cref{cor:e8v6-LBLD,cor:e8v5-YM-budget}.
Item (D8) prevents the criterion from silently conditioning on a
globally coupled good event.
\end{proof}

\section{Good/bad separation and status}
\label{sec:e8v6-status}

\begin{firewall}[A global good event can create infinite influence]
\label{fire:e8v6-good-event}
Conditioning a product reference on a single global indicator
\(\mathbf1_{\{\text{all blocks good}\}}\) couples every site.  The
product-reference proof then fails even if each local potential is
small.  The admissible construction uses sitewise good kernels and
assigns violated blocks to explicit bad polymers; their activities
must satisfy a separate summable polymer criterion.
\end{firewall}

\begin{firewall}[Compact Wilson verification remains open]
\label{fire:e8v6-Wilson}
V1--V3 provide the polynomial/tail bounds entering
\eqref{eq:e8v6-a-poly}--\eqref{eq:e8v6-a-tail}, but V6 has not proved
that the actual gauge-fixed compact-link conditional kernels satisfy
\(c<1\).  Gauge zero modes, torons, chart transitions, and bad blocks
must be removed or retained before that inequality can be tested.
\end{firewall}

\begin{theorem}[What eighth-series V6 proves]
\label{thm:e8v6-result}
V6 gives an exact Gaussian-to-interacting interpolation and a
finite-volume influence proof of a volume-uniform complete covariance
row.  Under DIC it computes
\[
 L_B=\frac{2J_\pi J_{\rm st}^2}{1-c},\qquad
 L_D=\frac{2J_QJ_{\rm m}J_{\rm st}}{1-c},
\]
thereby converting the V5 leakage debit into declared local
oscillation and influence constants.  No clustering assumption is
used as an unproved input.
\end{theorem}

\begin{proof}
Combine
\Cref{thm:e8v6-interpolation,thm:e8v6-Dobrushin-covariance,
prop:e8v6-product-test,thm:e8v6-complete-covariance,
cor:e8v6-LBLD,thm:e8v6-card}.
\end{proof}

\begin{assessment}[Progress at eighth-series V6]
\label{ass:e8v6-status}
The covariance part of the marginal Schur gate is now reduced to one
model-specific number \(c\) and explicit norm embeddings.  The likely
bottleneck has moved upstream to the gauge-fixed Wilson conditional
kernel: its quadratic reference must have a uniform physical-detail
floor on every fixed elementary fibre, while its nonlinear Hessian
and chart tail must remain inside that floor.
\end{assessment}

\begin{decision}[V7 acquisition order]
\label{dec:e8v6-V7}
V7 will derive the local Hessian/influence margin in exponential
coordinates.  It will isolate the flat gauge and toron directions,
use the V97 finite-block transverse floor on the remaining detail,
bound the Wilson nonlinear Hessian by \(O(gR)\), and determine whether
the logarithmic good radius permits a uniform positive margin.  If
the resulting Dobrushin row is not below one, the sites will be
regrouped into larger fixed superblocks before any clustering claim.
\end{decision}

\begin{firewall}[Claim boundary after eighth-series V6]
\label{fire:e8v6-claim}
V6 does not prove DIC for compact Yang--Mills, a bad-polymer bound,
the marginal inverse estimate, \(D_{\rm geom}<b^2\), rotational
restoration, running asymptotic freedom, global chart closure,
NCGC/CFTC, terminal entry, regulator independence, continuum OS
reconstruction, or the full Yang--Mills mass gap.
\end{firewall}

\paragraph{Parent-version record.}
V6 was formed from
\texttt{Renewal\_Yang\_Mills\_Mass\_Gap\_Research\_Notes\_V5.tex}.
The V5 parent had \(2\,852\) counted source lines,
\(103\,087\) bytes, and SHA--256
\[
\texttt{09320D6A0C1E99A0ADC763EC593052942AA33F6ACF6CEDA520CD2B8113C18570}.
\]
The next compulsory companion audit is V10.

% =============================================================================
% END APPENDED EIGHTH-SERIES V6 MATERIAL
% =============================================================================

% =============================================================================
% BEGIN APPENDED EIGHTH-SERIES V7 MATERIAL
% =============================================================================

\part{Eighth-series V7: transverse Wilson convexity and the
Gaussian-cumulant route}
\label{part:e8v7}

\section{The physical detail space}
\label{sec:e8v7-detail}

Let \(E_\perp\) denote the fixed elementary-fibre detail space after
tree gauge, retention of the coarse field, and removal of flat
cohomological and toron modes.  The earlier flat calculation,
summarized in \Cref{rec:e8v1-v95-v97}, gives
\[
 H_0=d_1^*d_1\big|_{E_\perp},\qquad
 \langle h,H_0h\rangle\ge\lambda_\perp\|h\|^2,
 \quad\lambda_\perp>0.                              \label{eq:e8v7-flat-floor}
\]
The number \(\lambda_\perp\) depends on the fixed elementary
geometry, not on the number of RG steps.

\begin{firewall}[The floor is only transverse]
\label{fire:e8v7-floor}
Equation \eqref{eq:e8v7-flat-floor} is false on gauge gradients and
flat cohomological modes.  Those directions remain retained or are
quotiented with their exact Jacobian.  They are never inserted into
\(H_0^{-1}\).
\end{firewall}

\section{Nonlinear Wilson and Haar Hessian debits}
\label{sec:e8v7-Hessian}

Use the physical exponential coordinate and scaling of V1,
\[
 {\cal W}_g(X)=g^{-2}\widetilde\phi(gX).
\]
For \(n=3,4,5\), write
\[
 M_n:=\sup_{\|Z\|\le r_0}\|D^n\widetilde\phi(Z)\|.   \label{eq:e8v7-Mn}
\]

\begin{theorem}[Exact good-ball Hessian debit]
\label{thm:e8v7-Hessian-debit}
If \(\|X\|\le R\) and \(gR\le r_0\), then
\[
 \|D^2{\cal W}_g(X)-H_0\|
 \le
 \varepsilon_W(g,R)
 :=gM_3R+\frac{g^2M_4}{2}R^2
       +\frac{g^3M_5}{6}R^3.                        \label{eq:e8v7-W-debit}
\]
\end{theorem}

\begin{proof}
The quartic Taylor polynomial of
\eqref{eq:e8v1-polynomial} has Hessian
\[
 H_0+gD^3\widetilde\phi(0)[X,\,\cdot\,,\,\cdot\,]
 +\frac{g^2}{2}D^4\widetilde\phi(0)
       [X,X,\,\cdot\,,\,\cdot\,].
\]
The two nonquadratic terms have norms at most
\(gM_3R\) and \(g^2M_4R^2/2\).  The \(p=4,k=2\) case of
\Cref{thm:e8v1-Taylor} bounds the remaining Hessian by
\(g^3M_5R^3/3!\).  Sum the three estimates.
\end{proof}

\begin{lemma}[Haar-coordinate Hessian debit]
\label{lem:e8v7-Haar}
Let \(j(Z)\) be the analytic Haar density in exponential coordinates
on \(\|Z\|\le r_0\), with \(j(0)=1\).  For
\[
 U_{{\rm Haar},g}(X):=-\log j(gX),                  \label{eq:e8v7-Haar-potential}
\]
there is a finite chart constant
\[
 C_{\rm Haar}:=
 \sup_{\|Z\|\le r_0}\|D^2[-\log j](Z)\|             \label{eq:e8v7-CHaar}
\]
such that
\[
 \|D^2U_{{\rm Haar},g}(X)\|\le g^2C_{\rm Haar}.      \label{eq:e8v7-Haar-debit}
\]
\end{lemma}

\begin{proof}
The chain rule gives
\(D^2U_{{\rm Haar},g}(X)=g^2D^2[-\log j](gX)\).
The closed chart ball is compact and lies before the first zero of
the analytic Jacobian, so the displayed supremum is finite.
\end{proof}

The stationary conditional centre of the earlier programme has
Hessian \(L_{g,y}={\cal L}(gy)\).  On a retained good ball
\(\|y\|\le R_y\), analyticity supplies
\[
 \|L_{g,y}-L_{0,0}\|\le C_y gR_y                  \label{eq:e8v7-centre-debit}
\]
for a fixed local constant \(C_y\).

\begin{corollary}[Uniform transverse convexity]
\label{cor:e8v7-convexity}
Define
\[
 \varepsilon_{\rm Hess}(g,R,R_y):=
 \varepsilon_W(g,R)+g^2C_{\rm Haar}+C_ygR_y.         \label{eq:e8v7-total-debit}
\]
If
\[
 \varepsilon_{\rm Hess}(g,R,R_y)<\lambda_\perp,      \label{eq:e8v7-convexity-gate}
\]
then the full good-chart detail Hessian satisfies
\[
 D^2U_{g,y}(X)\big|_{E_\perp}
 \ge
 \bigl(\lambda_\perp-\varepsilon_{\rm Hess}\bigr)I. \label{eq:e8v7-convexity}
\]
\end{corollary}

\begin{proof}
Combine
\Cref{eq:e8v7-flat-floor,thm:e8v7-Hessian-debit,
lem:e8v7-Haar} with \eqref{eq:e8v7-centre-debit} and use the
min--max bound for a self-adjoint perturbation.
\end{proof}

\begin{corollary}[Logarithmic good radii preserve the floor]
\label{cor:e8v7-log-radius}
Let \(g=\beta^{-1/2}\) and
\[
 R=c(\log\beta)^{\chi/2},\qquad
 R_y=c_y(\log\beta)^{\chi_y/2}                      \label{eq:e8v7-log-radii}
\]
with fixed finite \(\chi,\chi_y\).  Then
\(\varepsilon_{\rm Hess}\to0\) as \(\beta\to\infty\).
In particular, for all sufficiently large \(\beta\),
\[
 D^2U_{g,y}\big|_{E_\perp}\ge\frac{\lambda_\perp}{2}I. \label{eq:e8v7-half-floor}
\]
\end{corollary}

\begin{proof}
Every term in \eqref{eq:e8v7-total-debit} is a positive power of
\(\beta^{-1/2}\) times a fixed power of \(\log\beta\).  Such products
tend to zero.  Choose the threshold where the sum is at most
\(\lambda_\perp/2\).
\end{proof}

\section{Why total-variation Dobrushin is not the primary route}
\label{sec:e8v7-TV-obstruction}

\begin{proposition}[Unbounded Gaussian TV influence saturates]
\label{prop:e8v7-Gaussian-TV}
Let \(z_x,z_y\) be two blocks of a nondegenerate Gaussian field with
precision blocks \(H_{xx}>0\) and \(H_{xy}\ne0\).  If \(z_y\) ranges
over an unbounded vector space, the total-variation influence
coefficient \eqref{eq:e8v6-cxy} from \(y\) to \(x\) equals one.
\end{proposition}

\begin{proof}
Conditionally on the other blocks, \(z_x\) is Gaussian with covariance
\(H_{xx}^{-1}\) and mean
\[
 m_x=-H_{xx}^{-1}\sum_{v\ne x}H_{xv}z_v.            \label{eq:e8v7-cond-mean}
\]
Choose a vector \(a\) with \(H_{xy}a\ne0\) and compare backgrounds
whose \(y\)-coordinates differ by \(ta\).  The two conditional
Gaussians have the same covariance and mean separation growing
linearly in \(t\).  Projecting onto the direction of that separation
gives two one-dimensional equal-variance Gaussians whose
total-variation distance tends to one as \(t\to\infty\).  Taking the
supremum proves the claim.
\end{proof}

\begin{corollary}[V6 TV route is only a bounded-kernel option]
\label{cor:e8v7-TV-verdict}
For unbounded exponential detail variables with a coupled Gaussian
reference, the V6 criterion \(c_{\mu_0}<1\) cannot be verified using
unweighted total variation.  On a bounded good ball it may hold, but
its constant can deteriorate as the logarithmic radius grows.
\end{corollary}

\begin{proof}
The first assertion is \Cref{prop:e8v7-Gaussian-TV}.  On a ball the
mean displacement in \eqref{eq:e8v7-cond-mean} is bounded by a
constant times the radius, so the obstruction is removed but the
bound is radius dependent.
\end{proof}

\begin{decision}[Upstream change of route]
\label{dec:e8v7-route}
The primary proof will use the exact gapped Gaussian covariance and
the already established marked conditional KP tree expansion
summarized in \Cref{rec:e8v1-v92-v94}.  The V6 total-variation card is
retained only for genuinely compact or uniformly bounded
single-block kernels.
\end{decision}

\section{The gapped Gaussian covariance row}
\label{sec:e8v7-Gaussian}

Let \(\gamma_C\) be a centred Gaussian measure on the physical detail
space, with covariance \(C=H^{-1}\).  For a block decomposition, set
\[
 K_C(\mu_0):=
 \sup_x\sum_y e^{\mu_0d(x,y)}\|C_{xy}\|.             \label{eq:e8v7-KC}
\]

\begin{theorem}[Gaussian two-mark identity]
\label{thm:e8v7-Gaussian-covariance}
For bounded \(C^1\) functions \(F,G\) with bounded gradients,
\[
 \operatorname{Cov}_{\gamma_C}(F,G)
 =
 \int_0^\infty e^{-t}
 \mathbb E\!\left[
   \langle\nabla F(X),C\nabla G(X_t)\rangle
 \right]dt,                                         \label{eq:e8v7-OU-identity}
\]
where \(X,Y\) are independent with law \(\gamma_C\) and
\[
 X_t=e^{-t}X+\sqrt{1-e^{-2t}}\,Y.                   \label{eq:e8v7-Xt}
\]
Consequently,
\[
 |\operatorname{Cov}_{\gamma_C}(F,G)|
 \le\sum_{x,y}\|\nabla_xF\|_\infty
       \|C_{xy}\|\|\nabla_yG\|_\infty.               \label{eq:e8v7-Gaussian-bound}
\]
\end{theorem}

\begin{proof}
Let \(P_t\) be the Ornstein--Uhlenbeck semigroup with invariant law
\(\gamma_C\).  Invariance and convergence to the mean give
\[
 \operatorname{Cov}(F,G)
 =-\int_0^\infty\frac d{dt}
   \langle F,P_tG\rangle_{\gamma_C}\,dt.
\]
Gaussian integration by parts identifies the negative generator
form with
\(\langle\nabla F,C\nabla P_tG\rangle\), and
\(\nabla P_tG=e^{-t}P_t\nabla G\).  The Mehler formula for \(P_t\)
is \eqref{eq:e8v7-Xt}, proving \eqref{eq:e8v7-OU-identity}.
Take absolute values blockwise and integrate
\(\int_0^\infty e^{-t}dt=1\).
\end{proof}

\begin{lemma}[Weighted inverse row from block row domination]
\label{lem:e8v7-inverse-row}
Suppose a positive Hessian has a block decomposition
\[
 H=D-K,\qquad D=\bigoplus_xD_x,\qquad
 d_0:=\sup_x\|D_x^{-1}\|<\infty,                    \label{eq:e8v7-DK}
\]
and assume
\[
 \rho(\mu_0):=
 \sup_x\sum_y e^{\mu_0d(x,y)}
 \|D_x^{-1}K_{xy}\|<1.                              \label{eq:e8v7-rho}
\]
Then
\[
 K_C(\mu_0)\le
 \frac{d_0}{1-\rho(\mu_0)}.                         \label{eq:e8v7-KC-bound}
\]
\end{lemma}

\begin{proof}
In the weighted block row algebra,
\[
 H^{-1}=(I-D^{-1}K)^{-1}D^{-1}
 =\sum_{n=0}^\infty(D^{-1}K)^nD^{-1}.
\]
The triangle inequality for \(d\) makes that row norm
submultiplicative.  The \(n\)-th summand has weighted row at most
\(\rho(\mu_0)^nd_0\).  Sum the geometric series.
\end{proof}

\begin{corollary}[Uniform Gaussian row on the good chart]
\label{cor:e8v7-uniform-Gaussian}
Assume \eqref{eq:e8v7-half-floor} and, in addition, a block splitting
\eqref{eq:e8v7-DK} for which \(d_0\) is uniform and
\(\rho(\mu_0)<1\) at some fixed \(\mu_0>0\), uniformly for
sufficiently large \(\beta\).  Then \(K_C(\mu_0)\) is finite uniformly
in volume.
\end{corollary}

\begin{proof}
Apply \Cref{lem:e8v7-inverse-row}.
\end{proof}

\begin{firewall}[A spectral gap is not an absolute row bound]
\label{fire:e8v7-gap-row}
The lower spectral floor \eqref{eq:e8v7-half-floor} controls the
\(\ell^2\) inverse norm.  It does not by itself control the absolute
block row \(K_C(\mu_0)\) uniformly in volume.  The row domination
\eqref{eq:e8v7-rho}, a Combes--Thomas estimate proved in the relevant
row algebra, or an equivalent locality theorem is an additional
required input.
\end{firewall}

\section{Loading the earlier marked cumulant tree}
\label{sec:e8v7-KP}

Move the complete quadratic physical-detail form, including its
calibrated marginal coefficients, into the Gaussian reference.  The
good-chart activity then begins with the cubic Wilson term and has
the schematic complete bound
\[
 \zeta(g,R)
 :=
 K_3gR^3+K_4g^2R^4+K_5g^3R^5
       +K_Hg^2R^2.                                  \label{eq:e8v7-zeta}
\]
The terms are respectively the cubic, quartic, analytic tail, and
Haar activity rows; every \(K_i\) is a fixed-dictionary constant.

\begin{lemma}[Cofinal smallness of the activity]
\label{lem:e8v7-zeta}
For \(g=\beta^{-1/2}\) and
\(R=c(\log\beta)^{\chi/2}\) with fixed \(\chi\),
\[
 \zeta(g,R)\longrightarrow0.                        \label{eq:e8v7-zeta-zero}
\]
\end{lemma}

\begin{proof}
Each summand is a positive power of \(\beta^{-1/2}\) times a fixed
power of \(\log\beta\).
\end{proof}

\begin{definition}[Marked Gaussian KP loader]
\label{def:e8v7-loader}
MGKL consists of:
\begin{enumerate}[label=\textup{(K\arabic*)},leftmargin=4em]
\item the physical-detail floor
      \eqref{eq:e8v7-convexity-gate};
\item a Gaussian covariance row \(K_C(\mu_0)<\infty\);
\item fixed-dictionary activities bounded by
      \eqref{eq:e8v7-zeta};
\item the source-dependent selected-tree KP card recorded in
      \Cref{rec:e8v1-v92-v94}, with branching number
      \(q_{\rm br}<1\);
\item one- and two-mark norms compatible with the V3 oscillation
      embeddings; and
\item explicit saturated-bad, chart-transition, determinant, toron,
      and cohomology owners outside the good cumulant.
\end{enumerate}
\end{definition}

\begin{theorem}[Two-mark KP covariance row]
\label{thm:e8v7-KP-covariance}
Under MGKL, let \(J_F,J_G\) be the fixed embeddings of two marked
insertion norms into their Gaussian gradient rows.  Then the complete
connected covariance owner obeys
\[
 \|\operatorname{Cov}_{\mu_{1,y}}(F,G)\|_{{\rm row},\mu_0}
 \le
 \frac{K_C(\mu_0)J_FJ_G}{(1-q_{\rm br})^2}
 \|F\|\,\|G\|.                                      \label{eq:e8v7-KP-covariance}
\]
\end{theorem}

\begin{proof}
The zero-activity two-mark connection is bounded by
\Cref{thm:e8v7-Gaussian-covariance}.  In the selected-tree expansion,
every connected contribution has a unique two-mark trunk and rooted
activity branches attached to its two marked sides.  The previously
proved source-dependent KP card bounds the sum of each rooted branch
family by
\(\sum_{n\ge0}q_{\rm br}^n=(1-q_{\rm br})^{-1}\).
Multiplying the two branch sums, the Gaussian trunk row, and the two
mark embeddings gives \eqref{eq:e8v7-KP-covariance}.  Absolute
convergence permits the regrouping by retained footprint.
\end{proof}

\begin{corollary}[Leakage constants on the Gaussian-KP route]
\label{cor:e8v7-leakage}
Under MGKL, the V5 leakage constants may be taken as
\[
 L_B=
 \frac{K_CJ_\pi J_{\rm st}^2}{(1-q_{\rm br})^2},
 \qquad
 L_D=
 \frac{K_CJ_QJ_{\rm m}J_{\rm st}}{(1-q_{\rm br})^2}. \label{eq:e8v7-KP-L}
\]
\end{corollary}

\begin{proof}
Apply \Cref{thm:e8v7-KP-covariance} to the stable/stable and
marginal/stable marked pairs and then the output projections, exactly
as in \Cref{cor:e8v6-LBLD}.
\end{proof}

\begin{firewall}[The quadratic form belongs in the reference]
\label{fire:e8v7-quadratic-reference}
If an \(O(1)\) quadratic physical-detail coupling is left among the
activities, \(\zeta\) need not tend to zero.  Moving it into the
Gaussian reference also creates the owned
\(\frac12\log\det L_y\) row and its retained derivatives; omitting
that determinant is not an admissible simplification.
\end{firewall}

\section{Status and next acquisition}
\label{sec:e8v7-status}

\begin{theorem}[What eighth-series V7 proves]
\label{thm:e8v7-result}
On the fixed physical detail fibre, the nonlinear Wilson, Haar, and
moving-centre Hessian debit tends to zero on every fixed logarithmic
good-radius schedule, preserving at least half of the transverse
flat floor for sufficiently weak coupling.  Total-variation
Dobrushin is proved unsuitable for an unbounded coupled Gaussian
reference.  The replacement Gaussian--KP route gives the explicit
two-mark covariance and leakage bounds
\eqref{eq:e8v7-KP-covariance}--\eqref{eq:e8v7-KP-L} under MGKL.
\end{theorem}

\begin{proof}
Use
\Cref{cor:e8v7-log-radius,prop:e8v7-Gaussian-TV,
cor:e8v7-uniform-Gaussian,lem:e8v7-zeta,
thm:e8v7-KP-covariance,cor:e8v7-leakage}.
\end{proof}

\begin{assessment}[Progress at eighth-series V7]
\label{ass:e8v7-status}
The local Hessian bottleneck is closed on the physical logarithmic
good chart, and an inappropriate influence norm has been rejected.
The remaining quantitative acquisition is the actual uniform
branching number \(q_{\rm br}\) after all fixed-block activities and
marked sources are inserted.  Its proof must include the determinant
and bad-polymer ledgers rather than count only the Wilson Taylor tail.
\end{assessment}

\begin{decision}[V8 acquisition order]
\label{dec:e8v7-V8}
V8 will expose the selected-tree branching number as a finite
dictionary matrix.  It will prove a spectral-radius criterion for
all one- and two-mark rows, insert the explicit weak-coupling
activity vector \(\zeta\), and separate the good-tree matrix from the
saturated-bad polymer matrix.  The objective is a checkable
\(q_{\rm br}<1\) with no hidden volume factor.
\end{decision}

\begin{firewall}[Claim boundary after eighth-series V7]
\label{fire:e8v7-claim}
V7 does not verify MGKL for the complete compact model, compute the
branching matrix, close the bad-polymer row, prove the marginal
inverse bound or \(D_{\rm geom}<b^2\), restore rotations, control the
running coupling, establish NCGC/CFTC and terminal entry, remove the
regulator, reconstruct the continuum theory, or prove the full
Yang--Mills mass gap.
\end{firewall}

\paragraph{Parent-version record.}
V7 was formed from
\texttt{Renewal\_Yang\_Mills\_Mass\_Gap\_Research\_Notes\_V6.tex}.
The V6 parent had \(3\,323\) counted source lines,
\(119\,858\) bytes, and SHA--256
\[
\texttt{6A4A66182C54E467AE01BA5542F02B17D2C2F6DB5089586239FF0B8D6E7C1191}.
\]
The next compulsory companion audit is V10.

% =============================================================================
% END APPENDED EIGHTH-SERIES V7 MATERIAL
% =============================================================================

% =============================================================================
% BEGIN APPENDED EIGHTH-SERIES V8 MATERIAL
% =============================================================================

\part{Eighth-series V8: a finite-type branching matrix and the
saturated-bad row}
\label{part:e8v8}

\section{Good activity types}
\label{sec:e8v8-good-types}

After the complete quadratic physical-detail form is placed in the
Gaussian reference, let
\[
 {\mathcal I}_{\rm g}
 =\{3,4,5,H,\det,\operatorname{tr}\}                 \label{eq:e8v8-types}
\]
denote the finite good dictionary of cubic, quartic, analytic-tail,
Haar, determinant-response, and chart-transition activity types.
The transition type here means a local overlap contained in two valid
good charts; chart failure belongs to the bad row.

\begin{definition}[Uniform activity vector]
\label{def:e8v8-activity}
For \(i\in{\mathcal I}_{\rm g}\), let \(z_i(\beta)\) be the supremum
of the absolute source-marked activity row of type \(i\), including
its fixed support weight.  The weak-coupling majorant is
\begin{align}
 z_3&\le K_3gR^3,&
 z_4&\le K_4g^2R^4,&
 z_5&\le K_5g^3R^5,                                \label{eq:e8v8-z1}\\
 z_H&\le K_Hg^2R^2,&
 z_{\det}&\le K_{\det}g^2P_{\det}(R_y),&
 z_{\operatorname{tr}}&\le K_{\operatorname{tr}}
 g^{p_{\operatorname{tr}}}P_{\operatorname{tr}}(R,R_y), \label{eq:e8v8-z2}
\end{align}
where \(p_{\operatorname{tr}}>0\) and the two \(P\)'s are fixed
polynomials.  These inequalities are part of the activity ledger,
not definitions of the exact activities.
\end{definition}

\begin{lemma}[Every good activity vanishes cofinally]
\label{lem:e8v8-good-zero}
With \(g=\beta^{-1/2}\) and logarithmic radii as in
\eqref{eq:e8v7-log-radii},
\[
 z_i(\beta)\longrightarrow0
 \quad(i\in{\mathcal I}_{\rm g}).                    \label{eq:e8v8-good-zero}
\]
\end{lemma}

\begin{proof}
Each right side of
\eqref{eq:e8v8-z1}--\eqref{eq:e8v8-z2} is a positive power of
\(\beta^{-1/2}\) times a fixed polynomial in \(\log\beta\).
\end{proof}

\begin{firewall}[The determinant constant is not an activity]
\label{fire:e8v8-det-constant}
The field-independent Gaussian normalization may be large and is
owned by the vacuum coordinate.  Only its retained-field response,
whose first two derivatives were shown to be \(O(g^2)\), enters
\(z_{\det}\).  Treating the normalization itself as a small polymer
would be false.
\end{firewall}

\section{The finite-type good branching matrix}
\label{sec:e8v8-good-matrix}

\begin{definition}[Geometric incidence constants]
\label{def:e8v8-incidence}
For good types \(i,j\), let \(N_{ij}(\mu_0)\) be the supremum, over a
marked owner \(X\) of type \(i\), of the weighted sum of all
admissible type-\(j\) owners which can be attached to the selected
tree at \(X\).  Gaussian-separated attachments are weighted with the
covariance row of \eqref{eq:e8v7-KC}; overlapping attachments use the
fixed-dictionary incidence count.  Thus \(N_{ij}(\mu_0)<\infty\) is
a purely geometric constant, independent of volume and \(\beta\).
\end{definition}

\begin{definition}[Good branching matrix]
\label{def:e8v8-Kg}
Choose fixed KP size penalties \(a_j\ge0\) and define the nonnegative
matrix
\[
 (K_{\rm g})_{ij}
 :=N_{ij}(\mu_0)e^{a_j}z_j(\beta).                  \label{eq:e8v8-Kg}
\]
\end{definition}

\begin{lemma}[Volume independence]
\label{lem:e8v8-volume}
If the elementary block dictionary is finite, support diameters are
uniformly bounded, and \(K_C(\mu_0)<\infty\), then every
\(N_{ij}(\mu_0)\) is finite uniformly in the total lattice volume.
\end{lemma}

\begin{proof}
There are finitely many overlapping translations of two fixed
support shapes.  For separated Gaussian attachments, sum first over
the finitely many marked sites in each shape and then use the
weighted covariance row \(K_C(\mu_0)\).  Taking a maximum over the
finite type and orientation dictionary preserves finiteness.
\end{proof}

\begin{proposition}[Cofinal smallness of the good matrix]
\label{prop:e8v8-Kg-zero}
Under \Cref{lem:e8v8-good-zero,lem:e8v8-volume},
\[
 \|K_{\rm g}(\beta)\|_\infty
 :=\max_i\sum_j(K_{\rm g})_{ij}\longrightarrow0.     \label{eq:e8v8-Kg-zero}
\]
\end{proposition}

\begin{proof}
The matrix has fixed finite size.  Each \(N_{ij}e^{a_j}\) is fixed
and every \(z_j(\beta)\) tends to zero.
\end{proof}

\section{Saturated-bad polymers}
\label{sec:e8v8-bad}

Let a bad polymer be a connected union \(Y\) of elementary blocks.
Assume the earlier saturated-bad estimate has the explicit form
\[
 |\zeta_{\rm b}(Y)|
 \le e^{-\tau(\beta)|Y|},\qquad
 \tau(\beta)\longrightarrow\infty,                  \label{eq:e8v8-bad-activity}
\]
after its boundary, chart, and support weights are included.  The
logarithmic schedule with \(\chi>1\) was chosen precisely so that the
large-field Gaussian penalty dominates every fixed polynomial and
entropy debit.

\begin{lemma}[Connected-set entropy]
\label{lem:e8v8-lattice-animals}
On a lattice of maximum block degree \(\Delta\), the number of
connected block sets of size \(n\) containing a fixed block is at
most
\[
 A_\Delta^n,\qquad A_\Delta:=(e\Delta)^2.            \label{eq:e8v8-animal}
\]
\end{lemma}

\begin{proof}
Every connected set has a rooted spanning tree.  A depth-first
traversal of that tree is a walk of length \(2(n-1)\) from the root.
There are at most \(\Delta^{2(n-1)}\) such walks.  Since
\(\Delta^{2(n-1)}\le(e\Delta)^{2n}=A_\Delta^n\), the stated
conservative bound follows.  Multiple walks coding the same set only
increase the upper bound.
\end{proof}

\begin{proposition}[Bad-polymer KP row]
\label{prop:e8v8-bad-row}
For a fixed size penalty \(a_{\rm b}\ge0\), define
\[
 q_{\rm bb}(\beta):=
 \sup_x\sum_{Y\ni x}e^{a_{\rm b}|Y|}
 |\zeta_{\rm b}(Y)|.                                \label{eq:e8v8-qbb}
\]
If
\[
 r_{\rm b}(\beta):=
 A_\Delta e^{a_{\rm b}-\tau(\beta)}<1,              \label{eq:e8v8-rb}
\]
then
\[
 q_{\rm bb}(\beta)
 \le\frac{r_{\rm b}(\beta)}{1-r_{\rm b}(\beta)}
 \longrightarrow0.                                 \label{eq:e8v8-qbb-bound}
\]
\end{proposition}

\begin{proof}
Group connected polymers containing \(x\) by their size and apply
\Cref{lem:e8v8-lattice-animals,eq:e8v8-bad-activity}.  The resulting
geometric series is
\(\sum_{n\ge1}r_{\rm b}^n\).
\end{proof}

\begin{firewall}[Badness must be extensive on a component]
\label{fire:e8v8-bad-extensive}
The estimate \eqref{eq:e8v8-bad-activity} must charge every block of a
saturated bad component, or an equivalent boundary fraction.  A
single penalty independent of \(|Y|\) cannot beat the lattice-animal
entropy and does not imply \eqref{eq:e8v8-qbb-bound}.
\end{firewall}

\section{The combined branch matrix}
\label{sec:e8v8-combined}

Let \(k_{\rm gb}\) be the column of weighted bad attachments to good
types, and \(k_{\rm bg}\) the row of good attachments to a marked bad
polymer.  Fixed interface geometry and
\eqref{eq:e8v8-bad-activity} give bounds
\[
 \|k_{\rm gb}(\beta)\|_\infty\longrightarrow0,
 \qquad
 \|k_{\rm bg}(\beta)\|_1\longrightarrow0.           \label{eq:e8v8-cross-zero}
\]
Define
\[
 {\mathbb K}(\beta):=
 \begin{pmatrix}
 K_{\rm g}(\beta)&k_{\rm gb}(\beta)\\
 k_{\rm bg}(\beta)&q_{\rm bb}(\beta)
 \end{pmatrix}.                                     \label{eq:e8v8-K}
\]

\begin{lemma}[Cross-row bound]
\label{lem:e8v8-cross}
If a good support meets at most \(s_{\rm g}\) blocks and a bad
component has at most \(c_\partial|Y|\) admissible boundary
attachments per good type, then
\[
 \|k_{\rm gb}\|_\infty
 \le s_{\rm g}\frac{r_{\rm b}}{1-r_{\rm b}},\qquad
 \|k_{\rm bg}\|_1
 \le c_\partial\sum_{j\in{\mathcal I}_{\rm g}}
 e^{a_j}z_j.                                        \label{eq:e8v8-cross-bound}
\]
Consequently \eqref{eq:e8v8-cross-zero} holds.
\end{lemma}

\begin{proof}
For the first inequality, root the bad polymer at one of at most
\(s_{\rm g}\) blocks of the good support and use the geometric sum in
\Cref{prop:e8v8-bad-row}.  For the second, divide by the bad size
penalty already present in its weighted norm; the number of boundary
attachment sites per unit size is at most \(c_\partial\), and the
remaining sum is the good activity vector.  Cofinal vanishing follows
from \Cref{lem:e8v8-good-zero,prop:e8v8-bad-row}.
\end{proof}

\begin{theorem}[Spectral-radius branching criterion]
\label{thm:e8v8-spectral}
For any finite nonnegative matrix \(K\), the following are equivalent:
\begin{enumerate}[label=\textup{(\roman*)},leftmargin=3.5em]
\item its spectral radius satisfies \(\rho(K)<1\);
\item there exist \(v>0\) and \(q<1\) such that \(Kv\le qv\).
\end{enumerate}
When \(\rho(K)<1\), the canonical choice
\[
 v=(I-K)^{-1}{\bf1}                                 \label{eq:e8v8-v}
\]
is positive and gives
\[
 q(K):=\max_i\frac{(Kv)_i}{v_i}
 =\max_i\left(1-\frac1{v_i}\right)<1.               \label{eq:e8v8-q}
\]
\end{theorem}

\begin{proof}
If \(Kv\le qv\), the weighted supremum norm
\(\|x\|_v=\max_i|x_i|/v_i\) gives \(\|K\|_v\le q\), hence
\(\rho(K)\le q<1\).  Conversely, if \(\rho(K)<1\), the Neumann series
for \((I-K)^{-1}\) converges entrywise and
\(v={\bf1}+K{\bf1}+\cdots\ge{\bf1}>0\).
Since \(Kv=v-{\bf1}\), \eqref{eq:e8v8-q} follows; every \(v_i\) is
finite, so the maximum is strictly below one.
\end{proof}

\begin{corollary}[The complete branching number is eventually strict]
\label{cor:e8v8-eventual}
Under the activity and geometry hypotheses above,
\[
 \|{\mathbb K}(\beta)\|_\infty\longrightarrow0.
                                                        \label{eq:e8v8-K-zero}
\]
Hence, for all sufficiently large \(\beta\),
\[
 \rho({\mathbb K}(\beta))<1,\qquad
 q_{\rm br}:=q({\mathbb K}(\beta))<1.                \label{eq:e8v8-qbr}
\]
\end{corollary}

\begin{proof}
Use
\Cref{prop:e8v8-Kg-zero,prop:e8v8-bad-row,
lem:e8v8-cross}.  Since
\(\rho({\mathbb K})\le\|{\mathbb K}\|_\infty\), it is eventually
strict.  Apply \Cref{thm:e8v8-spectral}.
\end{proof}

\section{One- and two-mark resolvents}
\label{sec:e8v8-marks}

\begin{proposition}[Selected-tree resolvent]
\label{prop:e8v8-resolvent}
Suppose the selected-tree code assigns every additional branch a
type transition dominated by \({\mathbb K}\).  If \(r\ge0\) is the
root activity vector, the sum of all one-mark branch sequences is
bounded componentwise by
\[
 (I-{\mathbb K})^{-1}r.                             \label{eq:e8v8-one-mark}
\]
In the weighted norm of \eqref{eq:e8v8-v}, its amplification is at
most \((1-q_{\rm br})^{-1}\).  Two independently rooted mark sides
have amplification at most
\[
 (1-q_{\rm br})^{-2}.                               \label{eq:e8v8-two-mark}
\]
\end{proposition}

\begin{proof}
A branch sequence with \(n\) extensions is bounded by
\({\mathbb K}^nr\).  Sum over \(n\ge0\).  The weighted norm satisfies
\(\|{\mathbb K}\|_v\le q_{\rm br}\), giving the first geometric
factor.  In the source-differentiated selected-tree code, the two
marked sides are rooted separately at the unique two-mark trunk; each
has the one-mark amplification, so their product gives
\eqref{eq:e8v8-two-mark}.
\end{proof}

\begin{firewall}[A branch sequence requires the selected-tree code]
\label{fire:e8v8-code}
The matrix resolvent does not bound arbitrary connected graphs by
itself.  It applies because the earlier conditional KP theorem
constructs an injective selected-parent code and charges its graph
entropy in \(N_{ij}\) and the size penalties.  Without that code,
replacing connected graphs by powers of \({\mathbb K}\) would be an
unproved combinatorial step.
\end{firewall}

\begin{theorem}[MGKL branching row closed at weak coupling]
\label{thm:e8v8-MGKL}
Assume the earlier selected-tree coding theorem, the complete
activity bounds
\eqref{eq:e8v8-z1}--\eqref{eq:e8v8-z2}, the extensive bad estimate
\eqref{eq:e8v8-bad-activity}, and the Gaussian row
\(K_C(\mu_0)<\infty\).  Then the branching item (K4) of MGKL holds
for all sufficiently large \(\beta\), with \(q_{\rm br}\) given by
\eqref{eq:e8v8-q}.  Its one- and two-mark amplifications are exactly
bounded by
\eqref{eq:e8v8-one-mark}--\eqref{eq:e8v8-two-mark}.
\end{theorem}

\begin{proof}
\Cref{lem:e8v8-volume} makes the good incidence constants finite.
\Cref{cor:e8v8-eventual} proves the strict branch number including
bad and cross rows.  Apply \Cref{prop:e8v8-resolvent}.
\end{proof}

\section{Status and next acquisition}
\label{sec:e8v8-status}

\begin{theorem}[What eighth-series V8 proves]
\label{thm:e8v8-result}
V8 converts the marked conditional KP branching constant into the
spectral radius of the explicit type matrix
\eqref{eq:e8v8-K}.  Every good entry vanishes by weak-coupling
Taylor/Haar/determinant estimates; the saturated-bad and good--bad
entries vanish when the bad penalty is extensive.  Consequently the
complete branch number is strictly below one for sufficiently large
\(\beta\), uniformly in volume, conditional only on the declared
selected-tree coding, Gaussian row, and activity ledgers.
\end{theorem}

\begin{proof}
Combine
\Cref{prop:e8v8-Kg-zero,prop:e8v8-bad-row,
lem:e8v8-cross,thm:e8v8-spectral,
cor:e8v8-eventual,thm:e8v8-MGKL}.
\end{proof}

\begin{assessment}[Progress at eighth-series V8]
\label{ass:e8v8-status}
The abstract \(q_{\rm br}<1\) hypothesis of V7 is now a finite
calculation plus one extensive bad-polymer estimate, and its
volume-uniformity is explicit.  The next live gate in the calibrated
contraction is the marginal inverse: prove \(A=I+o(1)\) in the full
finite-cutoff commutant coordinates while retaining the exact
determinant and normalization owners.
\end{assessment}

\begin{decision}[V9 acquisition order]
\label{dec:e8v8-V9}
V9 will compare each fine marginal insertion with its exact retained
pullback, express \(A-I\) by the conditional covariance interpolation,
and bound the defect with the V8 two-mark resolvent.  It will also
separate basis-transport error from interaction error and give a
Neumann threshold for the full commutant matrix.
\end{decision}

\begin{firewall}[Claim boundary after eighth-series V8]
\label{fire:e8v8-claim}
V8 does not supply numerical incidence constants for a chosen block
map, prove every declared activity estimate for all charts, establish
the marginal inverse, \(D_{\rm geom}<b^2\), rotational restoration,
the running beta function, global NCGC/CFTC and terminal entry,
regulator independence, continuum OS reconstruction, or the full
Yang--Mills mass gap.
\end{firewall}

\paragraph{Parent-version record.}
V8 was formed from
\texttt{Renewal\_Yang\_Mills\_Mass\_Gap\_Research\_Notes\_V7.tex}.
The V7 parent had \(3\,795\) counted source lines,
\(136\,147\) bytes, and SHA--256
\[
\texttt{42383CE184361748A46C5CA9876B124904B7B5E58FBEB59ED85EC468D8B1920F}.
\]
The next compulsory companion audit is V10.

% =============================================================================
% END APPENDED EIGHTH-SERIES V8 MATERIAL
% =============================================================================

% =============================================================================
% BEGIN APPENDED EIGHTH-SERIES V9 MATERIAL
% =============================================================================

\part{Eighth-series V9: exact pullback marginal coordinates}
\label{part:e8v9}

\section{A block-adapted marginal chart}
\label{sec:e8v9-chart}

Let \(\iota_+:\mathbb R^r\to{\mathfrak B}_+\) contain the full
finite-cutoff commutant ledger and let
\(\pi_+\iota_+=I\).  At the fine scale, let
\(\pi_0:{\mathfrak B}\to\mathbb R^r\) be the local moment projection,
with a local representative
\(\iota_0\) satisfying \(\pi_0\iota_0=I\).  Define the exact pullback
injection
\[
 {\mathscr J}u:=(\iota_+u)\circ B                   \label{eq:e8v9-J}
\]
and the finite basis-transport matrix
\[
 {\mathsf G}:=\pi_0{\mathscr J}.                    \label{eq:e8v9-G}
\]

\begin{theorem}[Exact pullback chart]
\label{thm:e8v9-chart}
Assume \({\mathsf G}\) is invertible.  Define
\[
 \pi_{\rm pb}:={\mathsf G}^{-1}\pi_0,\qquad
 M_{\rm pb}:={\mathscr J}\pi_{\rm pb},\qquad
 Q_{\rm pb}:=I-M_{\rm pb}.                          \label{eq:e8v9-pb-projections}
\]
Then
\begin{align}
 \pi_{\rm pb}{\mathscr J}&=I,&
 M_{\rm pb}^2&=M_{\rm pb},&
 Q_{\rm pb}^2&=Q_{\rm pb},                          \label{eq:e8v9-projection-identities}\\
 \ker\pi_{\rm pb}&=\ker\pi_0,&
 \operatorname{ran}Q_{\rm pb}&=\ker\pi_0.           \label{eq:e8v9-same-stable}
\end{align}
Every \(S\in{\mathfrak B}\) has the unique decomposition
\[
 S={\mathscr J}m+q,\qquad
 m=\pi_{\rm pb}S,\quad q=Q_{\rm pb}S,\quad
 \pi_0q=0.                                          \label{eq:e8v9-decomposition}
\]
\end{theorem}

\begin{proof}
Equation \eqref{eq:e8v9-G} gives
\(\pi_{\rm pb}{\mathscr J}={\mathsf G}^{-1}{\mathsf G}=I\).
The projection identities follow.  Since multiplication by
\({\mathsf G}^{-1}\) does not change a kernel,
\(\ker\pi_{\rm pb}=\ker\pi_0\).  A projection \(Q_{\rm pb}\) has
range equal to the kernel of its complementary coordinate projection.
Finally \(S={\mathscr J}\pi_{\rm pb}S+Q_{\rm pb}S\);
applying \(\pi_{\rm pb}\) proves uniqueness.
\end{proof}

\begin{corollary}[The V3 stable space is unchanged]
\label{cor:e8v9-stable-unchanged}
Changing from the local marginal representative \(\iota_0\) to the
pullback representative \({\mathscr J}\) does not change the stable
subspace \(\ker\pi_0\).  Hence the moment-zero quadratic, cubic,
quartic, and analytic-tail norms of V1--V3 remain the stable norms in
the pullback chart.
\end{corollary}

\begin{proof}
This is \eqref{eq:e8v9-same-stable}; only the complementary marginal
representative changes.
\end{proof}

\section{Exact affine coarse action}
\label{sec:e8v9-affine}

\begin{theorem}[Affine marginal factorization]
\label{thm:e8v9-affine}
For every \(m\in\mathbb R^r\) and every \(q\) for which the coarse
integral exists,
\[
 {\cal T}({\mathscr J}m+q)
 =\iota_+m+{\cal T}q.                               \label{eq:e8v9-affine}
\]
Consequently, in the pullback chart the four derivative blocks of V5
satisfy
\[
 A=I,\qquad D_{\rm ms}=0                            \label{eq:e8v9-AD}
\]
at every interacting background, not only at the Gaussian point.
\end{theorem}

\begin{proof}
On the fibre \(B\omega=y\),
\[
 ({\mathscr J}m)(\omega)
 =(\iota_+m)(B\omega)=(\iota_+m)(y).
\]
Factor \(e^{-(\iota_+m)(y)}\) out of the conditional partition
function and take its negative logarithm.  Differentiation in \(m\),
followed by \(\pi_+\) and \(Q_+\), gives
\eqref{eq:e8v9-AD}.
\end{proof}

\begin{corollary}[Exact marginal alignment without an inverse bound]
\label{cor:e8v9-alignment}
Put
\[
 \beta(q):=\pi_+{\cal T}q.                          \label{eq:e8v9-beta}
\]
For an arbitrary prescribed output marginal coordinate \(m_+\), the
unique input coordinate producing it is
\[
 m=m_+-\beta(q).                                    \label{eq:e8v9-m-calibration}
\]
The aligned stable map is exactly
\[
 \Phi_{m_+}(q)
 =Q_+{\cal T}q,                                     \label{eq:e8v9-stable-map}
\]
independent of \(m_+\), and
\[
 D_q\Phi_{m_+}=Q_+D{\cal T}_qQ_{\rm pb}=C.           \label{eq:e8v9-stable-derivative}
\]
\end{corollary}

\begin{proof}
Apply \(\pi_+\) and \(Q_+\) to
\eqref{eq:e8v9-affine}.  The marginal equation is
\(m+\beta(q)=m_+\), which has
\eqref{eq:e8v9-m-calibration} as its unique solution.  Since
\(Q_+\iota_+=0\), the stable output is
\eqref{eq:e8v9-stable-map}; differentiate along
\(\operatorname{ran}Q_{\rm pb}\).
\end{proof}

\begin{corollary}[The marginal Schur debit vanishes exactly]
\label{cor:e8v9-Schur-zero}
In the block-adapted chart,
\[
 C-D_{\rm ms}A^{-1}B_{\rm sm}=C.                   \label{eq:e8v9-Schur-zero}
\]
Stable-to-marginal leakage changes the required input marginal
coordinate through
\[
 D_qm=-B_{\rm sm},                                  \label{eq:e8v9-running-response}
\]
but it has no feedback into the aligned stable derivative.
\end{corollary}

\begin{proof}
Use \eqref{eq:e8v9-AD} in
\Cref{thm:e8v5-Schur}; differentiate
\eqref{eq:e8v9-m-calibration} for the second assertion.
\end{proof}

\begin{firewall}[The covariance rows are not discarded]
\label{fire:e8v9-covariance}
V6--V8 remain necessary to control the running response
\eqref{eq:e8v9-running-response}, its source derivatives, regulator
comparison, and the beta function.  Exact pullback coordinates remove
those rows only from the stable contraction constant.
\end{firewall}

\section{When the basis-transport matrix is the identity}
\label{sec:e8v9-reproduction}

\begin{definition}[Constant-curvature reproduction]
\label{def:e8v9-CCR}
CCR means that the linearized normalized block map sends every
constant physical curvature \(F\) in the regulated commutant sector
to the same normalized coarse curvature, and that \(\pi_0\) extracts
the complete zero-momentum commutant tensor in that normalization.
\end{definition}

\begin{proposition}[CCR fixes the transport matrix]
\label{prop:e8v9-G-identity}
If CCR holds and the input and output commutant bases are transported
with the same orientation convention, then
\[
 {\mathsf G}=I                                      \label{eq:e8v9-G-identity}
\]
on the entire finite-cutoff commutant.
\end{proposition}

\begin{proof}
For a basis tensor \(E_a\in{\mathcal C}_\Gamma\), test the quadratic
part of the pullback
\((\iota_+E_a)\circ B\) on an arbitrary constant physical curvature
\(F\).  CCR gives exactly
\(\langle F,E_aF\rangle\), with no scale factor in the normalized
curvature.  Since \(\pi_0\) is the coefficient extractor on the full
commutant, its \(a\)-th coefficient is one and all other basis
coefficients agree with those of \(E_a\).  Thus
\(\pi_0{\mathscr J}E_a=E_a\) for every basis element.
\end{proof}

\begin{firewall}[Constant fields must exist in the declared sector]
\label{fire:e8v9-CCR}
On a periodic bundle with nontrivial flux, a globally constant
potential need not represent every constant curvature.  CCR must be
tested in the local trivial sector or formulated with the appropriate
flux representatives.  Toron and cohomological directions remain
retained and are not inferred from the test.
\end{firewall}

\begin{corollary}[Reflection-even theta-zero chart]
\label{cor:e8v9-reflection}
If the block map is reflection equivariant and the programme is
restricted to the real reflection-even \(\theta=0\) sector, then
\({\mathscr J}\), \(\pi_{\rm pb}\), and \(Q_{\rm pb}\) preserve that
sector.  No odd theta coordinate is generated by the affine
calibration.
\end{corollary}

\begin{proof}
Equivariance of \(B\) and \(\iota_+\) makes
\({\mathscr J}\) intertwine reflection.  The full commutant projection
is chosen parity block diagonal by
\Cref{prop:e8v5-parity}.  Restriction to the even block proves the
claim.
\end{proof}

\section{The rechart debit}
\label{sec:e8v9-rechart}

Define the local stable residue of the pullback basis by
\[
 {\mathscr R}:=
 (I-\iota_0\pi_0){\mathscr J}
 ={\mathscr J}-\iota_0{\mathsf G}.                  \label{eq:e8v9-R}
\]
Then \(\pi_0{\mathscr R}=0\).

\begin{proposition}[Exact physical-to-pullback rechart]
\label{prop:e8v9-rechart}
Let a physical interaction be written in the local chart as
\[
 S_{\rm phys}=\iota_0a+s,\qquad \pi_0s=0.            \label{eq:e8v9-local-action}
\]
Its pullback-chart coordinates are
\[
 m={\mathsf G}^{-1}a,\qquad
 q=s-{\mathscr R}{\mathsf G}^{-1}a.                 \label{eq:e8v9-rechart}
\]
In particular,
\[
 \|q\|_{\rm st}
 \le\|s\|_{\rm st}
 +\|{\mathscr R}{\mathsf G}^{-1}\|_{\rm m\to st}|a|.
                                                        \label{eq:e8v9-rechart-bound}
\]
\end{proposition}

\begin{proof}
From \eqref{eq:e8v9-R},
\[
 {\mathscr J}{\mathsf G}^{-1}a
 =\iota_0a+{\mathscr R}{\mathsf G}^{-1}a.
\]
Subtract this pullback marginal part from
\eqref{eq:e8v9-local-action}.  The result is
\eqref{eq:e8v9-rechart}; apply the norm triangle inequality.
\end{proof}

\begin{proposition}[Structure of the rechart residue]
\label{prop:e8v9-R-structure}
Under the shellable filling, EDOC, and CCR hypotheses, the quadratic
Taylor part of \({\mathscr R}\) has zero zeroth and first moments and
belongs to the V3 derivative-quadratic stable sector.  Its cubic,
quartic, and analytic parts belong to the corresponding V3/V1 stable
rows, with the same fixed-dictionary orbit ownership.
\end{proposition}

\begin{proof}
The definition \(\pi_0{\mathscr R}=0\) removes the entire
zero-momentum commutant coefficient.  EDOC symmetrization and
inversion give the first-moment cancellation by
\Cref{thm:e8v4-quadratic-split}.  The V2 Stokes/BCH substitution
places the higher Taylor terms in the cubic and quartic curvature
dictionary, and \Cref{thm:e8v1-Taylor} places the remainder in the
analytic tail.  Orbit symmetrization does not increase the row norm
by \Cref{prop:e8v4-orbit-norm}.
\end{proof}

\begin{definition}[Pullback rechart card]
\label{def:e8v9-card}
PRC consists of:
\begin{enumerate}[label=\textup{(R\arabic*)},leftmargin=4em]
\item a physical equivariant block map \(B\);
\item the full finite-cutoff marginal injection \(\iota_+\);
\item an invertible transport matrix \({\mathsf G}\), preferably
      certified by CCR;
\item the exact pullback injection \({\mathscr J}=B^*\iota_+\);
\item the unchanged stable kernel \(\ker\pi_0\);
\item a complete stable-row bound for \({\mathscr R}\);
\item the radius gate
      \(\|s\|_{\rm st}+
       \|{\mathscr R}{\mathsf G}^{-1}\||a|\le r\);
\item determinant, chart, bad, toron, and cohomology owners; and
\item reflection restriction to the real theta-zero sector.
\end{enumerate}
\end{definition}

\begin{theorem}[PRC removes the marginal inverse bottleneck]
\label{thm:e8v9-card}
Under PRC, the marginal response matrix is exactly \(A=I\), the
marginal-to-stable block is exactly zero, and the aligned stable
derivative is \(C\).  If
\[
 \sup_{\|q\|_{\rm st}\le r}\|Q_+D{\cal T}_qQ_{\rm pb}\|
 <1,                                                 \label{eq:e8v9-final-gate}
\]
then the exact RG is a strict contraction between actions with the
same prescribed output marginal coordinate.  The physical local
action enters that ball precisely through the rechart gate (R7).
\end{theorem}

\begin{proof}
The chart and physical entry are
\Cref{thm:e8v9-chart,prop:e8v9-rechart}.  Exact affine factorization
and disappearance of the Schur debit are
\Cref{thm:e8v9-affine,cor:e8v9-Schur-zero}.  Integrating the derivative
in \eqref{eq:e8v9-final-gate} on the convex stable ball gives strict
contraction.
\end{proof}

\section{Status and next acquisition}
\label{sec:e8v9-status}

\begin{theorem}[What eighth-series V9 proves]
\label{thm:e8v9-result}
V9 constructs block-adapted marginal coordinates from exact retained
pullbacks.  In these coordinates the coarse action is affine in every
finite-cutoff marginal coefficient, so \(A=I\),
\(D_{\rm ms}=0\), and the V5 marginal Schur correction vanishes
identically.  The stable subspace is the same local moment-zero space
used in V3.  The only coordinate price is the explicit stable
rechart residue \({\mathscr R}{\mathsf G}^{-1}a\).
\end{theorem}

\begin{proof}
Combine
\Cref{thm:e8v9-chart,thm:e8v9-affine,
cor:e8v9-alignment,cor:e8v9-Schur-zero,
prop:e8v9-rechart,thm:e8v9-card}.
\end{proof}

\begin{assessment}[Progress at eighth-series V9]
\label{ass:e8v9-status}
The marginal inverse and stable-feedback bottlenecks are removed by
an exact coordinate choice rather than a smallness estimate.  The
next quantitative issue is the rechart debit.  It must fit inside the
stable ball for the physical running gauge coefficient, and this
requires an actual block-symbol estimate for
\({\mathscr R}\), not only its moment classification.
\end{assessment}

\begin{decision}[V10 acquisition order]
\label{dec:e8v9-V10}
V10 is a compulsory companion audit.  It will then compute the
quadratic pullback residue symbol from the minimum-energy
interpolation, derive its stable row in terms of the interpolation
stencil, and decide whether the physical gauge coefficient can enter
the pullback stable ball uniformly along the weak-coupling
trajectory.
\end{decision}

\begin{firewall}[Claim boundary after eighth-series V9]
\label{fire:e8v9-claim}
V9 does not bound the physical rechart debit, prove the raw stable
derivative gate \eqref{eq:e8v9-final-gate}, establish
\(D_{\rm geom}<b^2\), rotational restoration, the running beta
function, global chart and bad closure, NCGC/CFTC, terminal entry,
regulator independence, continuum OS reconstruction, or the full
Yang--Mills mass gap.
\end{firewall}

\paragraph{Parent-version record.}
V9 was formed from
\texttt{Renewal\_Yang\_Mills\_Mass\_Gap\_Research\_Notes\_V8.tex}.
The V8 parent had \(4\,204\) counted source lines,
\(151\,059\) bytes, and SHA--256
\[
\texttt{73282B7B774E68B1CD55D3EDE8B377588E5BF4BEF5FF1250D58C7650F7A768D4}.
\]
The compulsory companion audit is due in V10.

% =============================================================================
% END APPENDED EIGHTH-SERIES V9 MATERIAL
% =============================================================================

% =============================================================================
% BEGIN APPENDED EIGHTH-SERIES V10 MATERIAL
% =============================================================================

\part{Eighth-series V10: companion audit, pullback symbols, and the
moving stable reference}
\label{part:e8v10}

\section{Compulsory companion-program audit}
\label{sec:e8v10-audit}

\begin{audit}[Files inspected at the V10 boundary]
\label{aud:e8v10-files}
The audit used:
\begin{enumerate}[label=\textup{(\roman*)},leftmargin=3.5em]
\item
\texttt{Renewal\_SM\_Einstein\_Continuum\_Research\_Notes\_V82.tex}:
\(32\,288\) counted source lines, \(1\,161\,118\) bytes, SHA--256
\[
\texttt{711FEB2E7C68DAEF248BBC51C2BD0C77C5DA1D2098FA0D34A7471E644A23440A};
\]
\item
\texttt{Renewal\_NS\_Stability\_Research\_Notes\_V26.tex}:
\(5\,319\) counted source lines, \(278\,283\) bytes, SHA--256
\[
\texttt{82C887F8C2C69E4F28FAD7C3DC8DE5B9099CB5A4BF431EBA1FFF3E91935978AD}.
\]
\end{enumerate}
\end{audit}

\begin{decision}[V10 audit transfer]
\label{dec:e8v10-transfer}
The new SM V80--V82 material confirms the full finite-cutoff owner
space and contributes two corrections:
\begin{enumerate}[label=\textup{(\alph*)},leftmargin=3.5em]
\item consecutive calibrated charts produce a discrete connection
      increment measured against the \emph{transported} reference,
      not merely against the previous parameter; and
\item a log-determinant derivative along a calibrated chart already
      contains the connection response, which must not be added a
      second time.
\end{enumerate}
Its combined constraint/gauge determinant concerns an extra
constraint reduction absent from the present pure tree-gauge block;
the general warning remains: coarea, gauge quotient, Gaussian, and
ghost determinants are mutually distinct mechanisms.  NS V26 is
unchanged and again supplies no YM row.
\end{decision}

\section{Quadratic pullback residue from the interpolation stencil}
\label{sec:e8v10-symbol}

Work in the polyphase Fourier representation of one fixed block.  Let
the linearized normalized curvature block map have the finite stencil
\[
 \widehat Q(k)=\sum_{s\in{\mathcal S}}q_s e^{ik\cdot s}.             \label{eq:e8v10-Q}
\]
For a self-adjoint finite-cutoff marginal tensor
\(E\in{\mathcal C}_\Gamma\), the pulled-back quadratic symbol is
\[
 \widehat K_E(k):=\widehat Q(k)^*E\widehat Q(k).      \label{eq:e8v10-KE}
\]

\begin{definition}[Quadratic pullback residue]
\label{def:e8v10-residue}
The local zero-momentum representative and its residue are
\[
 \widehat K_{E,0}:=\widehat K_E(0),\qquad
 \widehat{\mathscr R}_E(k):=
 \widehat K_E(k)-\widehat K_E(0).                   \label{eq:e8v10-RE}
\]
\end{definition}

\begin{proposition}[Coefficient formula and zero moment]
\label{prop:e8v10-coefficients}
Put
\[
 k_E(r):=\sum_{\substack{s,t\in{\mathcal S}\\s-t=r}}
 q_s^*Eq_t.                                         \label{eq:e8v10-kEr}
\]
Then the residue coefficients are
\[
 {\mathscr R}_E(r)=
 \begin{cases}
 k_E(r),&r\ne0,\\
 k_E(0)-\sum_u k_E(u),&r=0,
 \end{cases}                                        \label{eq:e8v10-REr}
\]
and
\[
 \sum_r{\mathscr R}_E(r)=0.                         \label{eq:e8v10-zero-moment}
\]
If the stencil and orbit average are inversion invariant, then also
\[
 \sum_r r_j{\mathscr R}_E(r)=0\qquad(1\le j\le4).    \label{eq:e8v10-first-moment}
\]
\end{proposition}

\begin{proof}
Expanding \eqref{eq:e8v10-KE} gives
\(\widehat K_E(k)=\sum_rk_E(r)e^{ik\cdot r}\), while
\(\widehat K_E(0)=\sum_rk_E(r)\).  Subtraction gives
\eqref{eq:e8v10-REr} and its coefficient sum is zero.  Under inversion
averaging, \({\mathscr R}_E(r)={\mathscr R}_E(-r)\); pair \(r\) and
\(-r\) to obtain \eqref{eq:e8v10-first-moment}.
\end{proof}

\begin{theorem}[Explicit stable stencil row]
\label{thm:e8v10-stencil-row}
For \(\mu\ge0\), define
\[
 \Xi_Q(\mu):=
 \sum_{s,t\in{\mathcal S}}
 e^{\mu|s-t|}|s-t|^2\|q_s\|\,\|q_t\|.              \label{eq:e8v10-Xi}
\]
Then
\[
 \|{\mathscr R}_E\|_{\mu,2}
 :=\sum_r e^{\mu|r|}|r|^2\|{\mathscr R}_E(r)\|
 \le\Xi_Q(\mu)\|E\|.                                \label{eq:e8v10-stencil-bound}
\]
Consequently,
\[
 \|\widehat{\mathscr R}_E(k)\|
 \le\frac{|k|^2}{2}\Xi_Q(0)\|E\|.                  \label{eq:e8v10-symbol-bound}
\]
\end{theorem}

\begin{proof}
The \(r=0\) coefficient makes no contribution to the second-moment
norm.  Insert \eqref{eq:e8v10-kEr}, use
\(\|q_s^*Eq_t\|\le\|q_s\|\|E\|\|q_t\|\), and sum over \(s,t\).
For the symbol, use the two cancellations in
\Cref{prop:e8v10-coefficients} and
\[
 |e^{ik\cdot r}-1-ik\cdot r|
 \le\frac12|k|^2|r|^2.
\]
\end{proof}

\begin{corollary}[Full-commutant row]
\label{cor:e8v10-commutant-row}
For an orthonormal basis \(E_1,\ldots,E_r\) of
\({\mathcal C}_\Gamma\) and \(a=(a_1,\ldots,a_r)\),
\[
 \left\|\sum_{\ell=1}^ra_\ell{\mathscr R}_{E_\ell}
 \right\|_{\mu,2}
 \le \Xi_Q(\mu)\sqrt r\,|a|.                        \label{eq:e8v10-full-row}
\]
\end{corollary}

\begin{proof}
Apply \Cref{thm:e8v10-stencil-row} and
\(\sum_\ell|a_\ell|\le\sqrt r\,|a|\).
\end{proof}

\begin{corollary}[CCR normalization]
\label{cor:e8v10-CCR}
If \(\widehat Q(0)\) is the identity on normalized constant physical
curvatures, then
\[
 \widehat K_E(0)=E,
\]
so the matrix \({\mathsf G}\) of V9 is the identity and
\({\mathscr R}_E\) is exactly its quadratic stable rechart residue.
\end{corollary}

\begin{proof}
Evaluate \eqref{eq:e8v10-KE} at zero and apply
\Cref{prop:e8v9-G-identity}.
\end{proof}

\begin{firewall}[Aliases belong in the polyphase stencil]
\label{fire:e8v10-alias}
For a decimating block map, \(\widehat Q(k)\) is a matrix from fine
alias components to the coarse mode, not a scalar symbol on one
Brillouin zone.  Formula \eqref{eq:e8v10-KE} is valid in that
polyphase representation.  Dropping aliases can underestimate
\(\Xi_Q\) and erase retained toron components.
\end{firewall}

\section{The full nonlinear rechart bound}
\label{sec:e8v10-nonlinear}

Let \(R_f\) denote the normalized field radius, to distinguish it
from the residue \({\mathscr R}\).  The stable norm weights a
quadratic second-moment row by \(R_f^2\).

\begin{proposition}[Nonlinear pullback rechart majorant]
\label{prop:e8v10-nonlinear-rechart}
Under the fixed-dictionary Stokes/BCH, EDOC, and analytic-tail cards,
there are constants \(C_3,C_4,C_5,C_H,C_{\det}\), independent of
volume, such that
\begin{align}
 \|{\mathscr R}a\|_{\rm st}
 \le |a|\bigl[&
 R_f^2\Xi_Q(\mu)\sqrt r
 +C_3gR_f^3+C_4g^2R_f^4+C_5g^3R_f^5 \notag\\
 &+C_Hg^2R_f^2+C_{\det}g^2P_{\det}(R_y)\bigr].
                                                        \label{eq:e8v10-rechart-majorant}
\end{align}
\end{proposition}

\begin{proof}
The quadratic term is
\Cref{cor:e8v10-commutant-row}.  The cubic and quartic terms are the
V2 Stokes/BCH polynomial rows, the fifth-order term is the complete
V1 analytic remainder, and the Haar and determinant responses are
the V7/V8 marked rows.  EDOC orbit averaging is contractive in the
row norm, so no symmetry multiplicity is added.  Multiply each
homogeneous coefficient row by its stable field-radius weight.
\end{proof}

\begin{corollary}[The rechart debit is bounded but not cofinally zero]
\label{cor:e8v10-rechart-limit}
Suppose \(|a|\le M_{\rm marg}\) on the running marginal domain and
use logarithmic field radii.  Then
\[
 \limsup_{\beta\to\infty}\|{\mathscr R}a\|_{\rm st}
 \le
 M_{\rm marg}R_f^2\Xi_Q(\mu)\sqrt r                \label{eq:e8v10-rechart-limit}
\]
when \(R_f\) is the fixed normalized radius used by the stable
coefficient norm.  The nonquadratic, Haar, and determinant terms
vanish, but the quadratic stencil residue need not.
\end{corollary}

\begin{proof}
Apply \Cref{prop:e8v10-nonlinear-rechart}; every term containing a
positive power of \(g\) times a fixed logarithmic polynomial tends to
zero.
\end{proof}

\begin{firewall}[Power counting does not make the rechart small]
\label{fire:e8v10-rechart-small}
The residue has two extra momenta and therefore contracts after
canonical dilation, but its norm at the scale where it is inserted
can be \(O(1)\).  Entry into the stable domain requires an actual
radius or reference-family argument; \(O(k^2)\) alone is insufficient.
\end{firewall}

\section{A moving stable reference graph}
\label{sec:e8v10-reference-graph}

Let \(X_{\rm st}\) be the complete stable Banach space.  Abstract one
rescaled step, including the local stable projection but before the
next pullback rechart, by a map
\(\Phi:B_r\to X_{\rm st}\).  The next-chart forcing is
\[
 c(m):=-{\mathscr R}{\mathsf G}^{-1}m.              \label{eq:e8v10-cm}
\]

\begin{theorem}[Stable reference graph]
\label{thm:e8v10-reference-graph}
Assume
\[
 \|\Phi(q)-\Phi(q')\|\le\kappa\|q-q'\|,
 \qquad0\le\kappa<1,                                \label{eq:e8v10-Phi-contract}
\]
on the closed ball \(B_r\), and let \({\mathcal M}\) be a marginal
domain satisfying
\[
 \|\Phi(0)\|+\sup_{m\in{\mathcal M}}\|c(m)\|
 \le(1-\kappa)r.                                    \label{eq:e8v10-ball-gate}
\]
Then for every \(m\in{\mathcal M}\) there is a unique
\(h(m)\in B_r\) such that
\[
 h(m)=\Phi(h(m))+c(m).                              \label{eq:e8v10-graph}
\]
Moreover,
\[
 \|h(m)-h(m')\|
 \le
 \frac{\|{\mathscr R}{\mathsf G}^{-1}\|}{1-\kappa}
 |m-m'|.                                            \label{eq:e8v10-graph-Lipschitz}
\]
\end{theorem}

\begin{proof}
For fixed \(m\), the map \(q\mapsto\Phi(q)+c(m)\) sends \(B_r\) into
itself by \eqref{eq:e8v10-ball-gate} and is a contraction with factor
\(\kappa\).  Banach's fixed-point theorem gives the unique \(h(m)\).
Subtract the two fixed-point equations:
\[
 \|h(m)-h(m')\|
 \le\kappa\|h(m)-h(m')\|+\|c(m)-c(m')\|.
\]
Use \eqref{eq:e8v10-cm} and rearrange.
\end{proof}

\begin{corollary}[Moving-reference recursion]
\label{cor:e8v10-moving}
Suppose the forward recharted stable recursion is
\[
 q_{j+1}=\Phi(q_j)+c(m_{j+1})+d_j,                  \label{eq:e8v10-forward}
\]
where \(d_j\) is a declared dynamical/regulator defect.  Put
\(e_j=q_j-h(m_j)\).  Then
\[
 \|e_{j+1}\|
 \le\kappa\|e_j\|
 +\frac{\kappa\|{\mathscr R}{\mathsf G}^{-1}\|}
             {1-\kappa}|m_{j+1}-m_j|
 +\|d_j\|.                                          \label{eq:e8v10-moving-recursion}
\]
\end{corollary}

\begin{proof}
Subtract the fixed-point equation for \(h(m_{j+1})\) from
\eqref{eq:e8v10-forward}.  The forcing \(c(m_{j+1})\) cancels, giving
\[
 \|e_{j+1}\|
 \le\kappa\|q_j-h(m_{j+1})\|+\|d_j\|.
\]
Insert \(h(m_j)\), apply the triangle inequality and
\eqref{eq:e8v10-graph-Lipschitz}.
\end{proof}

\begin{corollary}[Cofinal tracking of the moving graph]
\label{cor:e8v10-tracking}
If
\[
 |m_{j+1}-m_j|\longrightarrow0,\qquad
 \|d_j\|\longrightarrow0,                           \label{eq:e8v10-forcing-zero}
\]
then \(e_j\to0\).
\end{corollary}

\begin{proof}
Iterate \eqref{eq:e8v10-moving-recursion}.  The initial error is
multiplied by \(\kappa^n\).  A geometric convolution of a bounded
sequence tending to zero also tends to zero: split it into a fixed
early part, suppressed by powers of \(\kappa\), and a late part
bounded by its tail supremum times \((1-\kappa)^{-1}\).
\end{proof}

\begin{remark}[Connection counted once]
\label{rem:e8v10-connection}
The second term in \eqref{eq:e8v10-moving-recursion} is the discrete
moving-reference connection.  It compares \(h(m_{j+1})\) with the
reference transported from scale \(j\).  Derivatives of the Gaussian
determinant along \(h(m)\) already contain \(Dh\); adding another
\emph{connection determinant} would double count it.
\end{remark}

\begin{firewall}[A bounded graph is not the beta function]
\label{fire:e8v10-beta}
\Cref{thm:e8v10-reference-graph} constructs the stable reference
conditional on a marginal path.  It neither constructs that path nor
proves \eqref{eq:e8v10-forcing-zero}.  Those require the exact
marginal recursion and, ultimately, the asymptotically free running
coupling with controlled remainders.
\end{firewall}

\section{The V10 acquisition card and status}
\label{sec:e8v10-status}

\begin{definition}[Pullback moving-reference card]
\label{def:e8v10-card}
PMRC consists of:
\begin{enumerate}[label=\textup{(M\arabic*)},leftmargin=4em]
\item the polyphase interpolation stencil \eqref{eq:e8v10-Q};
\item CCR and the full commutant coefficient projection;
\item the explicit finite number \(\Xi_Q(\mu)\);
\item the nonlinear rechart bound
      \eqref{eq:e8v10-rechart-majorant};
\item a complete stable map with contraction factor \(\kappa<1\);
\item the invariant-ball gate \eqref{eq:e8v10-ball-gate};
\item the exact running marginal recursion;
\item a vanishing marginal increment and regulator-defect row; and
\item determinant and gauge quotient mechanisms each declared once
      in the common regulated measure.
\end{enumerate}
\end{definition}

\begin{theorem}[PMRC consequence]
\label{thm:e8v10-card}
Under PMRC, the \(O(1)\) quadratic pullback residue is absorbed into a
unique stable reference graph \(h(m)\), and deviations from that graph
contract according to \eqref{eq:e8v10-moving-recursion}.  If the two
forcing tails vanish, every admissible trajectory tracks the moving
graph cofinally.
\end{theorem}

\begin{proof}
Items (M1)--(M4) give
\Cref{thm:e8v10-stencil-row,prop:e8v10-nonlinear-rechart}.  Items
(M5)--(M6) give \Cref{thm:e8v10-reference-graph}.  Items (M7)--(M8)
give \Cref{cor:e8v10-moving,cor:e8v10-tracking}.  Item (M9) supplies
the audit ownership condition.
\end{proof}

\begin{theorem}[What eighth-series V10 proves]
\label{thm:e8v10-result}
V10 computes the quadratic rechart residue exactly from the block
stencil and proves its weighted second-moment bound
\(\Xi_Q(\mu)\).  The nonquadratic rechart pieces vanish at weak
coupling, while the quadratic piece is honestly \(O(1)\).  Rather
than requiring it to vanish, V10 constructs a unique moving stable
reference graph and proves geometric tracking when the running
marginal increments and declared defects tend to zero.
\end{theorem}

\begin{proof}
Combine
\Cref{thm:e8v10-stencil-row,
prop:e8v10-nonlinear-rechart,cor:e8v10-rechart-limit,
thm:e8v10-reference-graph,cor:e8v10-tracking,
thm:e8v10-card}.
\end{proof}

\begin{assessment}[Progress at eighth-series V10]
\label{ass:e8v10-status}
The rechart obstruction is no longer mistaken for a perturbatively
small term.  Its full quadratic size is a finite stencil constant and
is incorporated into the reference graph.  The decisive remaining
gate is now the strict complete stable map in (M5): the prior
canonical \(b^{-2}\) gain must beat actual support regrouping and all
nonlinear/marked corrections on the invariant ball containing
\(h({\mathcal M})\).
\end{assessment}

\begin{decision}[V11 acquisition order]
\label{dec:e8v10-V11}
V11 will compute the support regrouping operator on the fixed
dictionary rather than bound it by the crude scalar
\(D_{\rm geom}\).  It will use a weighted finite type matrix, as in
V8, and test its spectral radius after the canonical factors
\(b^{-2},b^{-2},b^{-4}\) are inserted.  This can prove strict
contraction even when a worst-row integer multiplicity exceeds
\(b^2\).
\end{decision}

\begin{firewall}[Claim boundary after eighth-series V10]
\label{fire:e8v10-claim}
V10 does not verify the invariant-ball or strict stable-map gates,
compute the physical running coupling, prove rotational restoration,
global chart and bad closure, NCGC/CFTC, terminal mass generation,
regulator independence, continuum OS reconstruction, or the full
Yang--Mills mass gap.
\end{firewall}

\paragraph{Parent-version record.}
V10 was formed from
\texttt{Renewal\_Yang\_Mills\_Mass\_Gap\_Research\_Notes\_V9.tex}.
The V9 parent had \(4\,597\) counted source lines,
\(165\,114\) bytes, and SHA--256
\[
\texttt{26307C6EA5BCEA4A9DC388F573384E19D1B94E96B2D66393C259E5AE7A3AB08F}.
\]
The next compulsory companion audit is V15.

% =============================================================================
% END APPENDED EIGHTH-SERIES V10 MATERIAL
% =============================================================================

% =============================================================================
% BEGIN APPENDED EIGHTH-SERIES V11 MATERIAL
% =============================================================================

\part{Eighth-series V11: the stable support transport matrix}
\label{part:e8v11}

\section{Typed stable rows}
\label{sec:e8v11-types}

Let the finite stable type set be
\[
 {\mathcal I}_{\rm st}
 =\{Q_2,C_3,C_4,A_{\ge5},H,D_{\det},B_{\rm bad}\}.   \label{eq:e8v11-types}
\]
Here \(Q_2\) is the moment-zero quadratic row, \(C_3,C_4\) are the
cubic and quartic curvature rows, \(A_{\ge5}\) is the analytic tail,
\(H\) the nonconstant Haar/chart row, \(D_{\det}\) the calibrated
determinant response, and \(B_{\rm bad}\) the saturated-bad row.  The
full finite-cutoff zero-momentum commutant is excluded because it is
marginal.

\begin{definition}[Typed magnitude vector]
\label{def:e8v11-vector}
For \(\Psi\in X_{\rm st}\), let
\[
 x(\Psi)_i:=\|\Psi_i\|_i\ge0,\qquad
 x(\Psi)\in\mathbb R_+^{{\mathcal I}_{\rm st}},      \label{eq:e8v11-x}
\]
where every \(\|\cdot\|_i\) is the complete support/mark norm already
assigned to that type.
\end{definition}

\section{Exact support transport}
\label{sec:e8v11-transport}

\begin{definition}[Geometric transport entries]
\label{def:e8v11-G}
For a fixed block map and fixed owner dictionary, let \(G_{ij}\) be
the least constant such that the principal, source-linearized,
unrescaled coarse reindexing \({\cal P}\) satisfies
\[
 x({\cal P}\Psi)_i\le\sum_jG_{ij}x(\Psi)_j.          \label{eq:e8v11-G-bound}
\]
Every coefficient, orbit weight, conditional Gaussian contraction,
fine-to-coarse support incidence, and union-owner multiplicity is
included in \(G_{ij}\).  Thus \(G\) is not merely a count of supports.
\end{definition}

\begin{lemma}[Finite-volume independence]
\label{lem:e8v11-G-finite}
Under fixed-factor geometry, EDOC orbit ownership, the Gaussian row
\(K_C(\mu)<\infty\), and the fixed good/bad dictionaries, every
\(G_{ij}\) is finite and independent of total volume.
\end{lemma}

\begin{proof}
Overlapping fine owners have finitely many relative translations.
Separated conditional contractions sum by \(K_C(\mu)\).  EDOC uses
normalized orbit weights, so group cardinality adds no factor.
Saturated-bad unions are summed by the lattice-animal row of V8.
There are finitely many input and output types; taking their maxima
preserves all bounds.
\end{proof}

\begin{definition}[Canonical type factors]
\label{def:e8v11-s}
Let
\[
 s_{Q_2}=b^{-2},\qquad s_{C_3}=b^{-2},\qquad
 s_{C_4}=b^{-4},                                    \label{eq:e8v11-core-factors}
\]
and assign to every remaining stable type its proved canonical or
activity factor \(0\le s_j<1\).  Define
\[
 L_{ij}:=G_{ij}s_j.                                 \label{eq:e8v11-L}
\]
\end{definition}

\begin{firewall}[Scaling and incidence are not two copies of volume]
\label{fire:e8v11-volume}
The canonical factor \(b^{4-\Delta}\) already includes the coarse
four-volume rescaling of a translation density.  The matrix \(G\)
contains only the residual owner/reindexing operator in the chosen
normalized norm.  Counting all \(b^4\) fine anchors again in \(G\)
would double count volume.
\end{firewall}

\begin{proposition}[Principal vector domination]
\label{prop:e8v11-principal}
For the canonically rescaled principal stable map
\({\cal L}_{\rm st}\),
\[
 x({\cal L}_{\rm st}\Psi)\le Lx(\Psi)               \label{eq:e8v11-principal}
\]
componentwise.
\end{proposition}

\begin{proof}
Apply the type factor \(s_j\) to each input column in
\eqref{eq:e8v11-G-bound}.
\end{proof}

\section{Perron weighting}
\label{sec:e8v11-Perron}

\begin{theorem}[Typed spectral contraction]
\label{thm:e8v11-spectral}
If the nonnegative finite matrix \(L\) satisfies
\[
 \rho(L)<1,                                         \label{eq:e8v11-rho}
\]
then with
\[
 v=(I-L)^{-1}{\bf1},\qquad
 q_0:=\max_i\frac{(Lv)_i}{v_i}<1,                   \label{eq:e8v11-vq}
\]
the norm
\[
 \|\Psi\|_v:=\max_i\frac{x(\Psi)_i}{v_i}            \label{eq:e8v11-vnorm}
\]
obeys
\[
 \|{\cal L}_{\rm st}\Psi\|_v\le q_0\|\Psi\|_v.      \label{eq:e8v11-linear-contraction}
\]
\end{theorem}

\begin{proof}
\Cref{thm:e8v8-spectral} gives \(v>0\), \(Lv\le q_0v\), and
\(q_0<1\).  If \(x(\Psi)\le v\|\Psi\|_v\), then
\Cref{prop:e8v11-principal} gives
\[
 x({\cal L}_{\rm st}\Psi)
 \le Lv\|\Psi\|_v
 \le q_0v\|\Psi\|_v.
\]
Take the maximum of the componentwise quotients.
\end{proof}

\begin{corollary}[The Perron norm still controls complete oscillation]
\label{cor:e8v11-oscillation}
Suppose
\[
 {\cal O}_\mu(\Psi)\le\sum_i c_i x(\Psi)_i           \label{eq:e8v11-Otype}
\]
with finite \(c_i\).  Then
\[
 {\cal O}_\mu(\Psi)
 \le\left(\sum_i c_iv_i\right)\|\Psi\|_v.           \label{eq:e8v11-Ov}
\]
Thus optimizing the type weights does not lose control of any
complete one- or two-mark row.
\end{corollary}

\begin{proof}
Use \(x_i(\Psi)\le v_i\|\Psi\|_v\) and sum.
\end{proof}

\begin{corollary}[Crude multiplicity can be misleading]
\label{cor:e8v11-crude}
The scalar condition
\[
 \max_i\sum_jG_{ij}\max_js_j<1                     \label{eq:e8v11-crude}
\]
implies \(\rho(L)<1\), but is not necessary.  In particular, large
off-diagonal transport can coexist with strict contraction.
\end{corollary}

\begin{proof}
The displayed number bounds \(\|L\|_\infty\), hence its spectral
radius.  Nonnecessity follows, for example, from any upper triangular
matrix with diagonal entries below one and an arbitrarily large
strictly upper triangular entry.
\end{proof}

\begin{theorem}[Triangular stable transport]
\label{thm:e8v11-triangular}
Suppose the types can be ordered so that \(L\) is upper triangular.
Then
\[
 \rho(L)=\max_iL_{ii}.                              \label{eq:e8v11-triangular-rho}
\]
Hence strict diagonal canonical transport proves contraction even
when a raw row sum, and therefore a scalar \(D_{\rm geom}b^{-2}\)
bound, exceeds one.
\end{theorem}

\begin{proof}
The eigenvalues of a triangular finite matrix are its diagonal
entries.  Apply \Cref{thm:e8v11-spectral}.
\end{proof}

\begin{remark}[Why triangularity is plausible but not assumed]
\label{rem:e8v11-Wick}
At the centred Gaussian principal point, Wick contraction can send a
quartic insertion to a quadratic or constant output, while a cubic
insertion cannot leave a gauge-invariant linear output.  Ordering by
decreasing polynomial degree therefore makes the principal
polynomial block triangular after marginal projection.  Moving
centres, nonlinear covariance, chart, and bad rows can create reverse
arrows; V11 places them in the correction matrix below instead of
assuming they vanish.
\end{remark}

\section{Nonlinear correction matrix}
\label{sec:e8v11-nonlinear}

Let \(\Phi_m\) denote the stable map based at the moving reference
graph \(h(m)\).  Define the entrywise nonnegative correction matrix
\(N(g,r)\) by the uniform derivative bound
\[
 x\!\left(
 [D\Phi_m(q)-{\cal L}_{\rm st}]\Psi
 \right)
 \le N(g,r)x(\Psi)                                  \label{eq:e8v11-N}
\]
for \(m\in{\mathcal M}\), \(\|q-h(m)\|_v\le r\).

\begin{proposition}[Weak-coupling correction row]
\label{prop:e8v11-N-small}
Under the V7 Hessian margin and V8 marked KP card, the good entries of
\(N(g,r)\) admit a fixed-dictionary majorant
\[
 \|N_{\rm good}(g,r)\|_v
 \le C_v\left(
 gR^3+g^2R^4+g^3R^5+g^2P(R_y)
 \right),                                           \label{eq:e8v11-Ngood}
\]
while its bad entries are bounded by
\[
 \|N_{\rm bad}(g,r)\|_v
 \le C_{v,{\rm b}}\frac{r_{\rm b}(\beta)}
                         {1-r_{\rm b}(\beta)}.       \label{eq:e8v11-Nbad}
\]
Both right sides tend to zero on the declared weak-coupling schedule.
\end{proposition}

\begin{proof}
Differentiate the source-dependent selected-tree expansion once.
The marked Gaussian trunk is controlled by \(K_C(\mu)\); all
additional good attachments have the activity weights listed in V8
and the two marked branch sums are bounded by
\((1-q_{\rm br})^{-2}\), which is absorbed in \(C_v\).
The saturated-bad marked sum is the geometric row of
\Cref{prop:e8v8-bad-row}.  Cofinal vanishing follows from
\Cref{lem:e8v8-good-zero,prop:e8v8-bad-row}.
\end{proof}

\begin{theorem}[Complete typed contraction gate]
\label{thm:e8v11-complete-gate}
Assume \(\rho(L)<1\) and use the Perron norm
\eqref{eq:e8v11-vnorm}.  If
\[
 \varepsilon_v(g,r):=\|N(g,r)\|_v<1-q_0,            \label{eq:e8v11-epsilon-gate}
\]
then
\[
 \sup_{m\in{\mathcal M}}
 \sup_{\|q-h(m)\|_v\le r}
 \|D\Phi_m(q)\|_v
 \le q_0+\varepsilon_v(g,r)<1.                      \label{eq:e8v11-complete-contraction}
\]
\end{theorem}

\begin{proof}
Use
\(D\Phi_m={\cal L}_{\rm st}+
(D\Phi_m-{\cal L}_{\rm st})\),
\Cref{thm:e8v11-spectral}, and
\eqref{eq:e8v11-N}.  The triangle inequality gives the displayed
bound.
\end{proof}

\begin{corollary}[Eventual contraction once the principal matrix is
strict]
\label{cor:e8v11-eventual}
If \(\rho(L)<1\) uniformly on the limiting marginal domain and the
V7/V8 activity cards hold uniformly, then
\eqref{eq:e8v11-complete-contraction} holds for all sufficiently large
\(\beta\).
\end{corollary}

\begin{proof}
\Cref{prop:e8v11-N-small} gives
\(\varepsilon_v(g,r)\to0\); compare it with the fixed positive margin
\(1-q_0\).
\end{proof}

\section{What must actually be computed}
\label{sec:e8v11-computation}

\begin{definition}[Stable transport certificate]
\label{def:e8v11-card}
STC consists of:
\begin{enumerate}[label=\textup{(S\arabic*)},leftmargin=4em]
\item the complete stable type dictionary
      \eqref{eq:e8v11-types};
\item the exact support/conditional entries \(G_{ij}\);
\item no second counting of the canonical coarse-volume factor;
\item the canonical column factors \(s_j\);
\item a certified spectral radius \(\rho(L)<1\);
\item Perron weights \(v\) and margin \(1-q_0\);
\item complete marked good and saturated-bad correction matrices;
\item the invariant ball around \(h({\mathcal M})\); and
\item reflection-compatible owner and determinant assignments.
\end{enumerate}
\end{definition}

\begin{theorem}[STC closes PMRC item (M5)]
\label{thm:e8v11-card}
STC and the V7/V8 weak-coupling cards imply the strict complete
stable-map item (M5) of PMRC for all sufficiently large \(\beta\).
\end{theorem}

\begin{proof}
Items (S1)--(S6) give the principal contraction by
\Cref{thm:e8v11-spectral}.  Items (S7)--(S9) permit
\Cref{prop:e8v11-N-small}.  Apply
\Cref{thm:e8v11-complete-gate,cor:e8v11-eventual}.
\end{proof}

\begin{firewall}[A symbolic matrix is not a certificate]
\label{fire:e8v11-symbolic}
V11 identifies every entry to be computed, but does not supply the
numerical stencil, incidence, conditional-contraction, and bad-row
values for a chosen physical block map.  STC item (S5) is proved only
after those entries yield a verified spectral-radius bound.
\end{firewall}

\section{Status and next acquisition}
\label{sec:e8v11-status}

\begin{theorem}[What eighth-series V11 proves]
\label{thm:e8v11-result}
The scalar \(D_{\rm geom}b^{-2}\) firewall is replaced by the exact
finite stable transport matrix \(L=G\operatorname{diag}(s)\).
Whenever \(\rho(L)<1\), explicit Perron weights give an equivalent
complete norm with strict principal contraction.  The marked
good/bad nonlinear derivative tends to zero, so the full stable map
is eventually contractive inside the same margin.
\end{theorem}

\begin{proof}
Combine
\Cref{thm:e8v11-spectral,cor:e8v11-oscillation,
thm:e8v11-complete-gate,cor:e8v11-eventual,
thm:e8v11-card}.
\end{proof}

\begin{assessment}[Progress at eighth-series V11]
\label{ass:e8v11-status}
The strict stable contraction is now reduced to a finite spectral
certificate rather than a worst-case multiplicity.  The next step
must choose the actual minimum-energy block stencil and compute or
rigorously enclose \(G\).  If its principal spectral radius is not
below one, the programme must change the block/interpolation or
enlarge the retained sector; nonlinear weak coupling cannot repair a
noncontracting principal cycle.
\end{assessment}

\begin{decision}[V12 acquisition order]
\label{dec:e8v11-V12}
V12 will derive \(G\) directly from a declared block stencil and
conditional Gaussian kernel.  It will first prove exact formulas for
the quadratic and polynomial diagonal blocks and the Wick
degree-lowering blocks, then give interval-ready finite sums.  No
numerical contraction will be claimed without those data.
\end{decision}

\begin{firewall}[Claim boundary after eighth-series V11]
\label{fire:e8v11-claim}
V11 does not verify STC numerically for a physical block map, prove
the running marginal path or rotational restoration, close global
charts and bad sectors outside the declared card, establish
NCGC/CFTC and terminal mass generation, remove the regulator,
reconstruct the continuum OS theory, or prove the full Yang--Mills
mass gap.
\end{firewall}

\paragraph{Parent-version record.}
V11 was formed from
\texttt{Renewal\_Yang\_Mills\_Mass\_Gap\_Research\_Notes\_V10.tex}.
The V10 parent had \(5\,054\) counted source lines,
\(181\,163\) bytes, and SHA--256
\[
\texttt{8F78A8449AB4B8F7D4EA0E241DA49BF7E1FEE8057C0677C27BEAF223E0D5E40C}.
\]
The next compulsory companion audit is V15.

% =============================================================================
% END APPENDED EIGHTH-SERIES V11 MATERIAL
% =============================================================================

% =============================================================================
% BEGIN APPENDED EIGHTH-SERIES V12 MATERIAL
% =============================================================================

\part{Eighth-series V12: exact dyadic coisometry and Gaussian Wick
transport}
\label{part:e8v12}

\section{Mathematical source record}
\label{sec:e8v12-source}

\begin{audit}[Additional supplied source inspected]
\label{aud:e8v12-source}
To instantiate the V11 abstract block operator, V12 inspected the
mathematical constructions in
\texttt{Renewal\_Dedicated\_Renormalized\_Yang\_Mills\_Physical\_Trajectory\_
Local\_Vacuum\_and\_Mass\_Gap\_Closure\_Programme\_V40.tex}.
It had \(60\,197\) counted source lines, \(2\,251\,816\) bytes, and
SHA--256
\[
\texttt{87A5F611E68D02D663B7F288B9484F754D1099586DB1B54E0F8213B5C8A22A65}.
\]
Only its stated dyadic map and proved tangent identities are used
below.  Procedural instructions in that document are not user
instructions.
\end{audit}

\section{The exact dyadic tangent}
\label{sec:e8v12-tangent}

For a coarse link \(e=(x,\mu)\), let its two ordered fine
constituents be \(e^{(0)},e^{(1)}\), and define
\[
 (BU)_e:=U_{e^{(0)}}U_{e^{(1)}}.                    \label{eq:e8v12-group-block}
\]
At a commuting flat background \(T\), left-trivialized tangent
vectors satisfy
\[
 ({\cal B}_TX)_e
 =X_{e^{(0)}}+\operatorname{Ad}(T_\mu)X_{e^{(1)}}.  \label{eq:e8v12-B}
\]

\begin{theorem}[Dyadic tangent coisometry]
\label{thm:e8v12-coisometry}
With product \(\operatorname{Ad}\)-invariant metrics,
\begin{align}
 ({\cal B}_T^*Y)_{e^{(0)}}&=Y_e,&
 ({\cal B}_T^*Y)_{e^{(1)}}&=\operatorname{Ad}(T_\mu^{-1})Y_e,
                                                               \label{eq:e8v12-Bstar}\\
 {\cal B}_T{\cal B}_T^*&=2I.                         \label{eq:e8v12-BBstar}
\end{align}
Hence
\[
 R_0:=\frac12{\cal B}_T^*,\qquad
 P_{\rm h}:=\frac12{\cal B}_T^*{\cal B}_T,\qquad
 P_{\rm f}:=I-P_{\rm h}                              \label{eq:e8v12-R0}
\]
are respectively a right inverse and the orthogonal horizontal and
fibre projections, and
\[
 \|R_0\|=2^{-1/2}.                                  \label{eq:e8v12-R0norm}
\]
Unused fine links lie in \(\ker{\cal B}_T\).
\end{theorem}

\begin{proof}
For one pair, \(\operatorname{Ad}(T_\mu)\) is orthogonal.  Therefore
the adjoint of
\((X_0,X_1)\mapsto X_0+\operatorname{Ad}(T_\mu)X_1\)
is the map in \eqref{eq:e8v12-Bstar}; applying the original map gives
\(Y+Y=2Y\).  Distinct coarse links use disjoint pairs, so the
identities sum orthogonally.  The projection formulas follow by
direct multiplication, and
\(\|R_0Y\|^2=\frac12\|Y\|^2\).
\end{proof}

\begin{corollary}[Exact Haar fibre coordinates]
\label{cor:e8v12-Haar}
For each used pair, the change of variables
\[
 U_{e^{(0)}}=H_e,\qquad
 U_{e^{(1)}}=H_e^{-1}V_e                            \label{eq:e8v12-Haar-coordinates}
\]
has product Haar measure \(dH_e\,dV_e\), and \(V_e=(BU)_e\).
Taking all pairs and unused links gives a global product-Haar fibre
for the exact dyadic map.
\end{corollary}

\begin{proof}
Bi-invariance of Haar measure gives
\(dU_{e^{(0)}}dU_{e^{(1)}}=dH_e\,dV_e\).  The pairs
are disjoint, so take their product.
\end{proof}

\section{Metric-horizontal versus energy-horizontal}
\label{sec:e8v12-energy}

Let \(H>0\) be the physical detail quadratic precision on the
gauge/toron-reduced tangent space and put
\[
 C_y:={\cal B}_TH^{-1}{\cal B}_T^*,\qquad
 J_H:=H^{-1}{\cal B}_T^*C_y^{-1}.                  \label{eq:e8v12-JH}
\]

\begin{theorem}[Energy lift for the dyadic map]
\label{thm:e8v12-energy-lift}
The operator \(J_H\) is the unique \(H\)-minimum-energy right inverse:
\[
 {\cal B}_TJ_H=I,\qquad
 \langle z,HJ_Hy\rangle=0\quad(z\in\ker{\cal B}_T). \label{eq:e8v12-energy-properties}
\]
The centred Gaussian with precision \(H\) has the exact decomposition
\[
 X=J_Hy+z,\qquad y={\cal B}_TX,\quad z\in\ker{\cal B}_T,             \label{eq:e8v12-Gaussian-split}
\]
with independent retained and fibre variables.  The fibre covariance
is
\[
 C_{\rm f}
 =H^{-1}-H^{-1}{\cal B}_T^*C_y^{-1}
             {\cal B}_TH^{-1}.                      \label{eq:e8v12-Cf}
\]
\end{theorem}

\begin{proof}
This is the minimum-energy construction of the V100/V75 audit with
the retained map \(Q={\cal B}_T\).  Direct multiplication gives the
right-inverse identity, and
\(HJ_H={\cal B}_T^*C_y^{-1}\) gives orthogonality to the kernel.
Completing the quadratic square gives independence and subtracting
the retained covariance \(J_HC_yJ_H^*\) from \(H^{-1}\) gives
\eqref{eq:e8v12-Cf}.
\end{proof}

\begin{proposition}[When the two lifts agree]
\label{prop:e8v12-lifts}
The metric-horizontal lift \(R_0\) equals \(J_H\) if and only if
\[
 P_{\rm f}HR_0=0.                                   \label{eq:e8v12-compatibility}
\]
Equivalently, the metric-horizontal and fibre subspaces are
\(H\)-orthogonal.
\end{proposition}

\begin{proof}
If \eqref{eq:e8v12-compatibility} holds, \(R_0\) is a right inverse
and its image is \(H\)-orthogonal to the fibre.  Uniqueness in
\Cref{thm:e8v12-energy-lift} gives \(R_0=J_H\).  Conversely, if the
lifts agree, the energy orthogonality in
\eqref{eq:e8v12-energy-properties} gives
\eqref{eq:e8v12-compatibility}.
\end{proof}

\begin{corollary}[Condition-number bound]
\label{cor:e8v12-JH-bound}
If
\[
 \lambda I\le H\le\Lambda I,                        \label{eq:e8v12-H-window}
\]
then
\[
 \|J_H\|\le\frac{\Lambda}{\sqrt2\,\lambda}.          \label{eq:e8v12-JH-bound}
\]
\end{corollary}

\begin{proof}
\Cref{thm:e8v12-coisometry} and
\eqref{eq:e8v12-H-window} give
\[
 \frac2\Lambda I\le C_y\le\frac2\lambda I,
 \qquad \|C_y^{-1}\|\le\frac\Lambda2.
\]
Now use \(\|H^{-1}\|\le\lambda^{-1}\) and
\(\|{\cal B}_T^*\|=\sqrt2\) in \eqref{eq:e8v12-JH}.
\end{proof}

\begin{firewall}[Coisometry is not energy orthogonality]
\label{fire:e8v12-two-lifts}
The attractive value \(\|R_0\|=2^{-1/2}\) cannot be used in a Wilson
Gaussian transport unless \eqref{eq:e8v12-compatibility} is proved.
In general the correct conditional mean is \(J_Hy\), and its certified
bound is \eqref{eq:e8v12-JH-bound}.
\end{firewall}

\section{Exact Wick transport}
\label{sec:e8v12-Wick}

Let \(T_n\) be a symmetric \(n\)-linear tensor and write
\[
 P_{T_n}(X):=\frac1{n!}T_n[X^{\otimes n}].           \label{eq:e8v12-PT}
\]
For \(m\) disjoint pairs of tensor slots, let
\(\operatorname{Contr}_{C_{\rm f}}^mT_n\) denote contraction with
\(C_{\rm f}\) in those pairs.

\begin{theorem}[Gaussian polynomial pushforward]
\label{thm:e8v12-Wick}
Under the conditional Gaussian split
\eqref{eq:e8v12-Gaussian-split},
\[
 \mathbb E\!\left[P_{T_n}(X)\mid y\right]
 =
 \sum_{m=0}^{\lfloor n/2\rfloor}
 \frac{
 (\operatorname{Contr}_{C_{\rm f}}^mT_n)
 [(J_Hy)^{\otimes(n-2m)}]}
 {(n-2m)!\,2^m m!}.                                 \label{eq:e8v12-Wick}
\]
\end{theorem}

\begin{proof}
Expand \(T_n[(J_Hy+z)^{\otimes n}]/n!\).  Odd Gaussian moments of
\(z\) vanish.  For \(2m\) fibre slots, Wick's theorem gives
\((2m)!/(2^mm!)\) pairings, while choosing their positions gives
\(\binom n{2m}\).  Multiplication by \(1/n!\) yields the coefficient
\(1/((n-2m)!2^mm!)\).
\end{proof}

\begin{corollary}[The first four degree channels]
\label{cor:e8v12-low-degrees}
Modulo constants,
\begin{align}
 P_{T_2}&\longmapsto
 \frac12T_2[(J_Hy)^2],                              \label{eq:e8v12-W2}\\
 P_{T_3}&\longmapsto
 \frac16T_3[(J_Hy)^3]
 +\frac12(\operatorname{Contr}_{C_{\rm f}}T_3)[J_Hy],
                                                               \label{eq:e8v12-W3}\\
 P_{T_4}&\longmapsto
 \frac1{24}T_4[(J_Hy)^4]
 +\frac14(\operatorname{Contr}_{C_{\rm f}}T_4)
          [(J_Hy)^2].                              \label{eq:e8v12-W4}
\end{align}
For an invariant \(\mathfrak{su}(N)\) block action, the linear term in
\eqref{eq:e8v12-W3} vanishes after the physical gauge-invariant
projection.
\end{corollary}

\begin{proof}
Insert \(n=2,3,4\) into \eqref{eq:e8v12-Wick}.  The linear output is
an invariant linear functional on the semisimple Lie algebra and
therefore vanishes by the argument of
\Cref{lem:e8v1-linear}.
\end{proof}

\begin{proposition}[Tensor-norm channel bounds]
\label{prop:e8v12-channel-bounds}
Put \(j_*=\|J_H\|\) and let
\(c_*=\|C_{\rm f}\|_1\) be the nuclear norm on the finite fibre
space.  Then
\begin{align}
 \|T_2\mapsto T_2\circ J_H^{\otimes2}\|
 &\le j_*^2,                                       \label{eq:e8v12-G22}\\
 \|T_3\mapsto T_3\circ J_H^{\otimes3}\|
 &\le j_*^3,                                       \label{eq:e8v12-G33}\\
 \|T_4\mapsto T_4\circ J_H^{\otimes4}\|
 &\le j_*^4,                                       \label{eq:e8v12-G44}\\
 \left\|T_4\mapsto
 \tfrac14(\operatorname{Contr}_{C_{\rm f}}T_4)
             \circ J_H^{\otimes2}\right\|
 &\le\tfrac14c_*j_*^2.                             \label{eq:e8v12-G24}
\end{align}
\end{proposition}

\begin{proof}
Use the defining multilinear operator norms, once for each copy of
\(J_H\).  In a singular-value decomposition of \(C_{\rm f}\), the
absolute sum of the pair-contraction coefficients is
\(\|C_{\rm f}\|_1=c_*\), which proves the last line without a hidden
dimension factor.
\end{proof}

\section{Interval-ready entries of the V11 matrix}
\label{sec:e8v12-entries}

Let \(\{T_{j,\alpha}\}\) be the finite oriented/support basis of input
type \(j\), let \({\cal O}\) be exact union/orbit ownership, and let
\(\Pi_i\) extract output type \(i\).  Define
\[
 {\mathscr W}_{n,m}(T):=
 \frac1{(n-2m)!\,2^mm!}
 (\operatorname{Contr}_{C_{\rm f}}^mT)
 \circ J_H^{\otimes(n-2m)}.                         \label{eq:e8v12-Wnm}
\]

\begin{theorem}[Finite formula for every polynomial transport entry]
\label{thm:e8v12-G-formula}
For polynomial input type \(j\) of degree \(n_j\), a valid V11
transport entry is
\[
 G_{ij}
 =
 \max_{\text{coarse marked owner }Y}
 \sum_{\alpha}
 w_{Y,\alpha}
 \sum_{m=0}^{\lfloor n_j/2\rfloor}
 \left\|
 \Pi_i{\cal O}_{Y,\alpha}
 {\mathscr W}_{n_j,m}(T_{j,\alpha})
 \right\|_i,                                        \label{eq:e8v12-G-finite}
\]
where terms of output degree different from type \(i\) are zero and
the fixed weights \(w_{Y,\alpha}\) include incidence and exponential
support weights.  Formula \eqref{eq:e8v12-G-finite} is a finite sum
for the dyadic fixed dictionary and can be enclosed by rational
interval arithmetic once \(H\), the Lie tensors, and the owner list
are supplied.
\end{theorem}

\begin{proof}
\Cref{thm:e8v12-Wick} gives every homogeneous conditional Gaussian
coefficient.  The maps \({\cal O}\) and \(\Pi_i\) assign it to its
unique output row.  Taking the absolute coefficient norm and summing
all input owners incident to a fixed marked output gives exactly the
definition \eqref{eq:e8v11-G-bound}.  The block, orientation, degree,
and Lie-algebra dictionaries are finite, so the maximum and sums are
finite.  Matrix inversion for \(C_y,J_H,C_{\rm f}\) and all tensor
contractions are finite-dimensional operations admitting outward
rounded interval evaluation.
\end{proof}

\begin{corollary}[Diagonal and Wick-lowering enclosures]
\label{cor:e8v12-L-bounds}
Let \(n_{22},n_{33},n_{44},n_{24}\) be the corresponding exact
support-incidence sums in \eqref{eq:e8v12-G-finite}.  Then
\begin{align}
 L_{Q_2Q_2}&\le b^{-2}n_{22}j_*^2,                  \label{eq:e8v12-L22}\\
 L_{C_3C_3}&\le b^{-2}n_{33}j_*^3,                  \label{eq:e8v12-L33}\\
 L_{C_4C_4}&\le b^{-4}n_{44}j_*^4,                  \label{eq:e8v12-L44}\\
 L_{Q_2C_4}&\le b^{-4}n_{24}\tfrac14c_*j_*^2
 \|\Pi_{Q_2}\|.                                     \label{eq:e8v12-L24}
\end{align}
\end{corollary}

\begin{proof}
Combine \Cref{prop:e8v12-channel-bounds} with the canonical column
factors of \eqref{eq:e8v11-core-factors} and the owner sums.
\end{proof}

\begin{example}[Compatible one-owner dyadic benchmark]
\label{ex:e8v12-benchmark}
If \eqref{eq:e8v12-compatibility} holds, \(b=2\), and the three
diagonal owner sums equal one, then
\[
 L_{Q_2Q_2}\le\frac18,\qquad
 L_{C_3C_3}\le\frac1{8\sqrt2},\qquad
 L_{C_4C_4}\le\frac1{64}.                           \label{eq:e8v12-benchmark}
\]
These values are a benchmark, not a certificate for the physical
owner dictionary.
\end{example}

\begin{proof}
Use \(j_*=\|R_0\|=2^{-1/2}\) in
\eqref{eq:e8v12-L22}--\eqref{eq:e8v12-L44}.
\end{proof}

\begin{firewall}[The remaining data are geometric, not asymptotic]
\label{fire:e8v12-incidence}
Weak coupling makes the nonlinear correction matrix small but does
not change \(n_{22},n_{33},n_{44}\), the energy-lift condition number,
or the principal spectral radius.  These fixed finite values must be
computed for the actual plaquette/network owner list.
\end{firewall}

\section{Status and next acquisition}
\label{sec:e8v12-status}

\begin{theorem}[What eighth-series V12 proves]
\label{thm:e8v12-result}
For the supplied exact dyadic link-product block, V12 proves the
tangent coisometry, separates its metric-horizontal lift from the
physical energy lift, derives the exact conditional Gaussian Wick
pushforward, and gives the finite formula
\eqref{eq:e8v12-G-finite} for every polynomial entry of the V11
transport matrix.  It also provides explicit diagonal and
quartic-to-quadratic enclosures.
\end{theorem}

\begin{proof}
Combine
\Cref{thm:e8v12-coisometry,thm:e8v12-energy-lift,
prop:e8v12-lifts,thm:e8v12-Wick,
thm:e8v12-G-formula,cor:e8v12-L-bounds}.
\end{proof}

\begin{assessment}[Progress at eighth-series V12]
\label{ass:e8v12-status}
The block map is no longer an unspecified stencil.  The unresolved
principal data are now the Wilson energy matrix \(H\) on one declared
gauge/toron-reduced block and four finite owner-incidence sums.
The simple dyadic benchmark has a large contraction margin, but using
it physically requires proof of energy compatibility and the actual
support counts.
\end{assessment}

\begin{decision}[V13 acquisition order]
\label{dec:e8v12-V13}
V13 will compute the coarse-plaquette pullback geometry of the
two-link dyadic map.  It will give the exact \(2\times2\) fine-face
nonabelian Stokes word, enumerate its perimeter and plaquette owners,
and obtain the incidence constants entering
\eqref{eq:e8v12-L22}--\eqref{eq:e8v12-L24}.  The energy matrix will
remain a separate finite spectral card if it cannot be inferred from
that combinatorics.
\end{decision}

\begin{firewall}[Claim boundary after eighth-series V12]
\label{fire:e8v12-claim}
V12 does not certify the physical principal spectral radius, prove
energy compatibility, close the moving-reference invariant ball,
construct the running coupling or rotational restoration, establish
global NCGC/CFTC and terminal mass generation, remove the regulator,
reconstruct the continuum OS theory, or prove the full Yang--Mills
mass gap.
\end{firewall}

\paragraph{Parent-version record.}
V12 was formed from
\texttt{Renewal\_Yang\_Mills\_Mass\_Gap\_Research\_Notes\_V11.tex}.
The V11 parent had \(5\,447\) counted source lines,
\(194\,550\) bytes, and SHA--256
\[
\texttt{D8FB2C0BAA15AE24237DD28E05CE6488D58D9221540C2D5E371B169A7A93BA3A}.
\]
The next compulsory companion audit is V15.

% =============================================================================
% END APPENDED EIGHTH-SERIES V12 MATERIAL
% =============================================================================

% =============================================================================
% BEGIN APPENDED EIGHTH-SERIES V13 MATERIAL
% =============================================================================

\part{Eighth-series V13: the dyadic \(2\times2\) face word and sharp
owner incidence}
\label{part:e8v13}

\section{A rooted \(2\times2\) face}
\label{sec:e8v13-face}

Fix coordinate directions \(\mu<\nu\).  Label the nine vertices of
one fine \(2\times2\) face by \((a,b)\), \(0\le a,b\le2\).
Let \(h_{ab}\) be the positively oriented \(\mu\)-link from
\((a,b)\) to \((a+1,b)\), and \(v_{ab}\) the positively oriented
\(\nu\)-link from \((a,b)\) to \((a,b+1)\).  Thus
\[
 0\le a\le1,\ 0\le b\le2\quad\hbox{for }h_{ab},
 \qquad
 0\le a\le2,\ 0\le b\le1\quad\hbox{for }v_{ab}.      \label{eq:e8v13-links}
\]
The fine plaquette based at \((a,b)\), \(0\le a,b\le1\), is
\[
 P_{ab}:=h_{ab}v_{a+1,b}h_{a,b+1}^{-1}v_{ab}^{-1}. \label{eq:e8v13-Pab}
\]
The coarse plaquette pulled back by the dyadic two-link map is the
large perimeter word
\[
 W:=
 h_{00}h_{10}v_{20}v_{21}
 h_{12}^{-1}h_{02}^{-1}v_{01}^{-1}v_{00}^{-1}.      \label{eq:e8v13-W}
\]

\begin{definition}[Face tree and rooted plaquettes]
\label{def:e8v13-tree}
Let \(T_\square\) contain all six horizontal links and the two left
vertical links \(v_{00},v_{01}\).  It is a spanning tree rooted at
\((0,0)\).  If \(t_{ab}\) is its root-to-\((a,b)\) path, define
\[
 \widetilde P_{ab}:=t_{ab}P_{ab}t_{ab}^{-1}.        \label{eq:e8v13-rooted-P}
\]
\end{definition}

\begin{theorem}[Exact four-plaquette Stokes word]
\label{thm:e8v13-Stokes}
With the conventions above,
\[
 \boxed{
 W=\widetilde P_{10}\widetilde P_{00}
   \widetilde P_{11}\widetilde P_{01}.}             \label{eq:e8v13-Stokes}
\]
This is an equality in the compact gauge group, with no commutativity
assumption.
\end{theorem}

\begin{proof}
Gauge transform along \(T_\square\), fixing the root, so that every
tree link equals the identity.  Then
\[
 P_{00}=v_{10},\qquad P_{01}=v_{11},\qquad
 P_{10}=v_{20}v_{10}^{-1},\qquad
 P_{11}=v_{21}v_{11}^{-1}.                          \label{eq:e8v13-tree-gauge-P}
\]
Hence
\[
 v_{20}=P_{10}P_{00},\qquad
 v_{21}=P_{11}P_{01}.
\]
The large perimeter reduces to \(W=v_{20}v_{21}\), proving the
displayed order in tree gauge.  Undoing the root-fixed gauge
transformation replaces each plaquette by its rooted conjugate and
leaves the root-based perimeter unchanged.
\end{proof}

\begin{corollary}[Exact dyadic block compatibility]
\label{cor:e8v13-block}
The product of the four coarse links around a coarse plaquette under
\eqref{eq:e8v12-group-block} is exactly the word \(W\) in
\eqref{eq:e8v13-W}; hence its rooted fine-plaquette factorization is
\eqref{eq:e8v13-Stokes}.
\end{corollary}

\begin{proof}
Each coarse link is its two-link constituent path.  Concatenating the
four coarse boundary links gives the eight factors in
\eqref{eq:e8v13-W}.
\end{proof}

\section{Linearized normalized curvature stencil}
\label{sec:e8v13-curvature}

Write
\(\widetilde P_{ab}=\exp(gF_{ab}+O(g^2))\) in the common root frame.
The physical coarse plaquette has four times the fine area, so define
its area-normalized curvature \(F_c\) by
\(W=\exp(4gF_c+O(g^2))\).

\begin{proposition}[Four-point curvature average]
\label{prop:e8v13-average}
The exact word \eqref{eq:e8v13-Stokes} gives
\[
 F_c=\frac14\sum_{a,b=0}^1F_{ab}+O(g).              \label{eq:e8v13-average}
\]
Thus the linear curvature stencil has offsets
\[
 {\mathcal S}_\square=\{0,\widehat\mu,\widehat\nu,
 \widehat\mu+\widehat\nu\},\qquad q_s=\frac14I.      \label{eq:e8v13-stencil}
\]
\end{proposition}

\begin{proof}
The first derivative at the identity of an ordered product is the
sum of the left-trivialized factors.  Divide the resulting
\(\sum_{a,b}F_{ab}\) by the coarse area factor four.  All commutators
begin at order \(g\) in \(F_c\).
\end{proof}

\begin{corollary}[Exact stencil second moment]
\label{cor:e8v13-Xi}
For Euclidean displacement and the stencil
\eqref{eq:e8v13-stencil},
\[
 \Xi_\square(\mu)
 :=\sum_{s,t\in{\mathcal S}_\square}
 e^{\mu|s-t|}|s-t|^2\|q_s\|\,\|q_t\|
 =
 \frac12\left(e^\mu+e^{\sqrt2\mu}\right).           \label{eq:e8v13-Xi}
\]
In particular, \(\Xi_\square(0)=1\).
\end{corollary}

\begin{proof}
Among the six unordered pairs of square vertices, four have squared
distance one and two have squared distance two.  Ordered pairs double
the sum and every coefficient product is \(1/16\).
\end{proof}

\begin{proposition}[Exact quadratic pullback residue]
\label{prop:e8v13-quadratic-residue}
For an invariant inner product on the curvature fibre, put
\(\overline F=\frac14\sum_{a,b}F_{ab}\).  Then
\[
 |\overline F|^2-\frac14\sum_{a,b}|F_{ab}|^2
 =
 -\frac1{16}
 \sum_{\{p,q\}\subset{\mathcal S}_\square}
 |F_p-F_q|^2.                                       \label{eq:e8v13-pair-difference}
\]
Thus the dyadic quadratic rechart residue is negative semidefinite on
the intra-face difference space, vanishes exactly on constant face
curvature, and has the two moment cancellations of V10.
\end{proposition}

\begin{proof}
For four vectors,
\[
 \sum_{p<q}|F_p-F_q|^2
 =4\sum_p|F_p|^2-\left|\sum_pF_p\right|^2.
\]
Multiply by \(-1/16\).  Constancy is precisely the kernel of all six
differences.  The moment statements also follow from
\Cref{cor:e8v13-Xi,prop:e8v10-coefficients}.
\end{proof}

\begin{firewall}[A negative residue is not a negative action]
\label{fire:e8v13-negative}
Equation \eqref{eq:e8v13-pair-difference} compares a coarse pullback
with a chosen local marginal representative.  The physical fine
quadratic form remains the full positive Wilson Hessian.  The residue
is a stable coordinate component, not a separately exponentiated
Gaussian precision.
\end{firewall}

\section{Sharp face and link incidence}
\label{sec:e8v13-incidence}

\begin{proposition}[Face support counts]
\label{prop:e8v13-face-counts}
One dyadic coarse face has:
\[
 4\ \hbox{fine plaquettes},\qquad
 8\ \hbox{perimeter links},\qquad
 12\ \hbox{links in the full tree/filling hull}.    \label{eq:e8v13-face-counts}
\]
\end{proposition}

\begin{proof}
The \(2\times2\) grid has four unit squares.  Its boundary has two
unit links on each of four sides.  It has six horizontal and six
vertical links in total.
\end{proof}

\begin{theorem}[Plaquette-owner uniqueness and sharp link incidence]
\label{thm:e8v13-incidence}
Fix the dyadic origin and one coordinate-plane orientation.
\begin{enumerate}[label=\textup{(\roman*)},leftmargin=3.5em]
\item Every fine plaquette lying in a selected dyadic face belongs to
      exactly one such \(2\times2\) face.
\item A fine link belongs to the filling hulls of at most two selected
      dyadic faces for each transverse direction.
\item In four dimensions, a fine link therefore belongs to at most
      \[
      2(4-1)=6                                      \label{eq:e8v13-link-six}
      \]
      selected face hulls, and this bound is attained by a link on
      dyadic boundary lines.
\end{enumerate}
\end{theorem}

\begin{proof}
In a fixed \(\mu\nu\)-plane, the selected faces have lower corners in
\(2\mathbb Z^2\), so their four unit squares partition the selected
plane.  This proves (i).  A link parallel to \(\mu\) is either on a
dyadic face boundary in the \(\nu\) coordinate, where it borders the
two adjacent faces, or on the middle line of one face, where it
belongs only to that face.  This proves (ii).  There are three choices
of \(\nu\ne\mu\); summing their sharp upper bounds proves (iii).
\end{proof}

\begin{corollary}[Use the plaquette carrier for contraction]
\label{cor:e8v13-carrier}
In a plaquette-anchored coefficient norm the principal dyadic filling
has owner incidence one.  Conversion to a link-marked oscillation row
costs at most the finite factor six in
\eqref{eq:e8v13-link-six}.  Hence the factor six belongs in the
norm-conversion constant, not in the principal plaquette coefficient
transport matrix.
\end{corollary}

\begin{proof}
Plaquette ownership is unique by
\Cref{thm:e8v13-incidence}(i).  A link mark can be reached from at
most six owner hulls by (iii), giving the conversion factor.
\end{proof}

\section{Polynomial coefficient mass}
\label{sec:e8v13-polynomial}

\begin{lemma}[Normalized average has unit coefficient mass]
\label{lem:e8v13-unit-mass}
For a bounded \(n\)-linear tensor \(T_n\),
\[
 T_n[\overline F^{\otimes n}]
 =
 4^{-n}\sum_{p_1,\ldots,p_n\in{\mathcal S}_\square}
 T_n[F_{p_1},\ldots,F_{p_n}],                       \label{eq:e8v13-expansion}
\]
and the absolute sum of the scalar coefficients on the right is one.
\end{lemma}

\begin{proof}
Expand multilinearly.  There are \(4^n\) ordered terms, each with
coefficient \(4^{-n}\).
\end{proof}

\begin{definition}[Face-energy compatibility]
\label{def:e8v13-energy-compatible}
FEC means that, in the principal flat Gaussian card, the conditional
mean of the four rooted fine curvatures at fixed normalized coarse
curvature \(F_c\) is
\[
 \mathbb E(F_{00},F_{10},F_{01},F_{11}\mid F_c)
 =(F_c,F_c,F_c,F_c),                                \label{eq:e8v13-FEC}
\]
and the owner projection respects the plaquette partition of
\Cref{thm:e8v13-incidence}.
\end{definition}

\begin{theorem}[Principal polynomial spectral margin under FEC]
\label{thm:e8v13-principal-margin}
Assume FEC and order the principal polynomial types as
\((Q_2,C_3,C_4)\), with Wick degree-lowering entries above the
diagonal.  In the plaquette-owner norm,
\[
 G_{Q_2Q_2}\le1,\qquad
 G_{C_3C_3}\le1,\qquad
 G_{C_4C_4}\le1.                                   \label{eq:e8v13-Gdiag}
\]
For the dyadic factor \(b=2\), the canonically scaled principal matrix
is upper triangular with
\[
 L_{Q_2Q_2}\le\frac14,\qquad
 L_{C_3C_3}\le\frac14,\qquad
 L_{C_4C_4}\le\frac1{16}.                           \label{eq:e8v13-Ldiag}
\]
Consequently,
\[
 \rho(L_{\rm poly})\le\frac14.                      \label{eq:e8v13-rho-quarter}
\]
\end{theorem}

\begin{proof}
FEC replaces each fine curvature mean by the same \(F_c\).
\Cref{lem:e8v13-unit-mass} and unique plaquette ownership show that
the absolute diagonal coefficient transport is at most one.
Centered Wick contractions can lower degree by two; the
gauge-invariant linear output of a cubic tensor vanishes, while the
quartic-to-quadratic term is strictly above the diagonal in the stated
order.  Insert the canonical factors
\(b^{-2},b^{-2},b^{-4}\) with \(b=2\).
\Cref{thm:e8v11-triangular} gives the spectral-radius bound.
\end{proof}

\begin{firewall}[FEC is the remaining physical matrix test]
\label{fire:e8v13-FEC}
The exact group word and incidence counts do not prove FEC.  Shared
links, lattice Bianchi relations, boundary conditions, and the
gauge/toron-reduced Wilson precision can make the energy-minimizing
curvature lift differ from four equal face curvatures.  That finite
linear-algebra question must be tested on the actual \(H\) of
\Cref{thm:e8v12-energy-lift}.
\end{firewall}

\begin{firewall}[Pullback and pushforward are different matrices]
\label{fire:e8v13-transpose}
\Cref{prop:e8v13-quadratic-residue} is an unconditional pullback
identity.  The diagonal bounds
\eqref{eq:e8v13-Gdiag} concern conditional Gaussian pushforward and
therefore require FEC.  Treating one as the transpose of the other
without the energy metric would be invalid.
\end{firewall}

\section{Status and next acquisition}
\label{sec:e8v13-status}

\begin{theorem}[What eighth-series V13 proves]
\label{thm:e8v13-result}
V13 gives the exact rooted four-plaquette word for a dyadic coarse
plaquette, the normalized four-point curvature stencil with
\(\Xi_\square(0)=1\), the exact negative pair-difference pullback
residue, unique plaquette ownership, and the sharp six-fold link-hull
incidence.  Under the single explicit FEC energy test, the principal
polynomial transport has spectral radius at most \(1/4\).
\end{theorem}

\begin{proof}
Combine
\Cref{thm:e8v13-Stokes,prop:e8v13-average,
cor:e8v13-Xi,prop:e8v13-quadratic-residue,
thm:e8v13-incidence,thm:e8v13-principal-margin}.
\end{proof}

\begin{assessment}[Progress at eighth-series V13]
\label{ass:e8v13-status}
The geometric multiplicity obstruction is closed in the
plaquette-owner norm: its principal value is one, with the link factor
six paid only during norm conversion.  The remaining principal
question is precisely FEC for the physical reduced Wilson Hessian.
This is a fixed finite matrix equality or, if it fails, a measured
lift defect; it is not an asymptotic estimate.
\end{assessment}

\begin{decision}[V14 acquisition order]
\label{dec:e8v13-V14}
V14 will replace the yes/no FEC test by an exact defect formula.
Writing the physical energy lift as the uniform face lift plus a
vertical correction, it will express the correction by a Schur solve,
bound the induced changes in all diagonal and Wick-lowering entries,
and give a quantitative spectral margin that tolerates nonzero FEC
defect.
\end{decision}

\begin{firewall}[Claim boundary after eighth-series V13]
\label{fire:e8v13-claim}
V13 does not prove FEC for the full reduced Wilson Hessian, verify the
complete STC matrix, close the moving invariant ball, construct the
running marginal path or rotational restoration, establish global
NCGC/CFTC and terminal mass generation, remove the regulator,
reconstruct the continuum OS theory, or prove the full Yang--Mills
mass gap.
\end{firewall}

\paragraph{Parent-version record.}
V13 was formed from
\texttt{Renewal\_Yang\_Mills\_Mass\_Gap\_Research\_Notes\_V12.tex}.
The V12 parent had \(5\,878\) counted source lines,
\(209\,809\) bytes, and SHA--256
\[
\texttt{11BEB7140650E8F9697CDD3CF61C22973673692A7E6F17595FEF11F6B65D2166}.
\]
The next compulsory companion audit is V15.

% =============================================================================
% END APPENDED EIGHTH-SERIES V13 MATERIAL
% =============================================================================

% =============================================================================
% BEGIN APPENDED EIGHTH-SERIES V14 MATERIAL
% =============================================================================

\part{Eighth-series V14: the face-energy lift defect and a robust
principal margin}
\label{part:e8v14}

\section{Horizontal/fibre block form}
\label{sec:e8v14-block}

Let \(Q_f:X\to Y\) be the normalized four-plaquette curvature average
of \eqref{eq:e8v13-average}, and let \(J_0:Y\to X\) be the uniform
face lift
\[
 J_0F=(F,F,F,F),\qquad Q_fJ_0=I.                   \label{eq:e8v14-J0}
\]
Let \(E=\ker Q_f\), so every fine face curvature vector has the unique
form \(x=z+J_0y\), \(z\in E\).

For a positive physical quadratic form \(H\) on the reduced face
space define
\begin{align}
 A&:=H|_{E\times E}:E\to E^*,&
 K&:=P_E^*HJ_0:Y\to E^*,&
 D&:=J_0^*HJ_0:Y\to Y^*.                            \label{eq:e8v14-AKD}
\end{align}
Thus
\[
 \langle z+J_0y,H(z+J_0y)\rangle
 =
 \langle z,Az\rangle+2\langle z,Ky\rangle
 +\langle y,Dy\rangle.                              \label{eq:e8v14-block-form}
\]

\begin{theorem}[Exact FEC defect]
\label{thm:e8v14-defect}
The physical energy-minimizing lift is
\[
 J_H=J_0+M,\qquad M:=-A^{-1}K:Y\to E.               \label{eq:e8v14-M}
\]
Moreover,
\[
 \langle z+J_0y,H(z+J_0y)\rangle
 =
 \langle z-M y,A(z-M y)\rangle
 +\langle y,Sy\rangle,                              \label{eq:e8v14-square}
\]
where
\[
 S:=D-K^*A^{-1}K>0.                                \label{eq:e8v14-S}
\]
FEC holds if and only if \(M=0\), equivalently \(K=0\).
\end{theorem}

\begin{proof}
Since \(H>0\), its restriction \(A\) is positive and invertible.
Completing the square in \eqref{eq:e8v14-block-form} gives
\eqref{eq:e8v14-square} with \(M=-A^{-1}K\).  The minimizer at fixed
\(y\) is \(z=My\), proving the lift formula.  Positivity of \(S\)
follows by evaluating the positive form at that minimizer.  The
uniform lift is minimizing exactly when \(M=0\), and invertibility of
\(A\) makes this equivalent to \(K=0\).
\end{proof}

\begin{corollary}[Conditional fibre covariance]
\label{cor:e8v14-Cf}
In the coordinates
\[
 x=J_Hy+w,\qquad w\in E,                            \label{eq:e8v14-JH-coordinates}
\]
the conditional centred fibre Gaussian has precision \(A\) and
covariance
\[
 C_{\rm f}=A^{-1}.                                  \label{eq:e8v14-Cf}
\]
In particular, the cross defect \(K\) changes the conditional mean
but not the fibre covariance after the square is completed.
\end{corollary}

\begin{proof}
Put \(z=My+w\) in \eqref{eq:e8v14-square}.  The conditional quadratic
form in \(w\) is \(\langle w,Aw\rangle\).
\end{proof}

\section{Relative defect and computable bounds}
\label{sec:e8v14-relative}

\begin{definition}[Relative FEC defect]
\label{def:e8v14-eta}
Define
\[
 {\mathsf T}:=A^{-1/2}KD^{-1/2},\qquad
 \eta_{\rm F}:=\|{\mathsf T}\|.                     \label{eq:e8v14-eta}
\]
\end{definition}

\begin{proposition}[Schur and lift bounds]
\label{prop:e8v14-bounds}
One has
\[
 S=D^{1/2}(I-{\mathsf T}^*{\mathsf T})D^{1/2}.      \label{eq:e8v14-relative-S}
\]
Hence \(H>0\) implies \(\eta_{\rm F}<1\).  If
\[
 A\ge\lambda_AI,\qquad D\le\Lambda_DI,              \label{eq:e8v14-windows}
\]
then
\[
 m_*:=\|M\|
 \le\eta_{\rm F}\sqrt{\frac{\Lambda_D}{\lambda_A}}. \label{eq:e8v14-M-bound}
\]
\end{proposition}

\begin{proof}
Insert \(K=A^{1/2}{\mathsf T}D^{1/2}\) into
\eqref{eq:e8v14-S}.  Strict positivity of \(S\) makes
\(I-{\mathsf T}^*{\mathsf T}>0\), hence
\(\eta_{\rm F}<1\).  Also
\[
 M=-A^{-1/2}{\mathsf T}D^{1/2},
\]
whose norm is bounded by the right side of
\eqref{eq:e8v14-M-bound}.
\end{proof}

\begin{corollary}[Finite singular-value certificate]
\label{cor:e8v14-SVD}
At a fixed regulator, \(\eta_{\rm F}\) and \(m_*\) can be enclosed by
finite interval linear algebra on \(A,K,D\).  A certified upper bound
\(\overline\eta_{\rm F}<1\) simultaneously proves the coarse Schur
floor
\[
 S\ge(1-\overline\eta_{\rm F}^{\,2})D               \label{eq:e8v14-Schur-floor}
\]
and the lift bound obtained by replacing \(\eta_{\rm F}\) in
\eqref{eq:e8v14-M-bound} by \(\overline\eta_{\rm F}\).
\end{corollary}

\begin{proof}
Outward-rounded spectral bounds enclose the largest singular value of
\(A^{-1/2}KD^{-1/2}\) and the extremal eigenvalues in
\eqref{eq:e8v14-windows}.  The inequalities follow from
\Cref{prop:e8v14-bounds}.
\end{proof}

\begin{firewall}[Positivity does not give a useful numerical margin]
\label{fire:e8v14-positive}
Although \(H>0\) proves \(\eta_{\rm F}<1\) in finite dimension, the
gap \(1-\eta_{\rm F}\) may be arbitrarily small.  A uniform RG
certificate requires an explicit upper enclosure separated from one.
\end{firewall}

\section{Tensor transport stability under lift motion}
\label{sec:e8v14-tensors}

\begin{lemma}[Tensor-power telescoping]
\label{lem:e8v14-telescope}
Let \(J=J_0+M\), and put
\[
 j_\#:=\max\{\|J\|,\|J_0\|\}.
\]
For every \(n\ge1\),
\[
 \|J^{\otimes n}-J_0^{\otimes n}\|
 \le n j_\#^{\,n-1}\|M\|.                           \label{eq:e8v14-telescope}
\]
\end{lemma}

\begin{proof}
Use the exact telescoping identity
\[
 J^{\otimes n}-J_0^{\otimes n}
 =
 \sum_{\ell=1}^n
 J^{\otimes(\ell-1)}\otimes M
 \otimes J_0^{\otimes(n-\ell)}
\]
and bound every summand by \(j_\#^{n-1}\|M\|\).
\end{proof}

\begin{theorem}[Diagonal channel defect]
\label{thm:e8v14-diagonal-defect}
For a degree-\(n\) multilinear coefficient tensor \(T_n\), the
difference between physical and FEC diagonal transports obeys
\[
 \|T_n\circ J_H^{\otimes n}
   -T_n\circ J_0^{\otimes n}\|
 \le n j_\#^{\,n-1}m_*\|T_n\|.                     \label{eq:e8v14-diagonal-defect}
\]
For \(n=2,3,4\), these give respectively
\[
 2j_\#m_*,\qquad3j_\#^2m_*,\qquad4j_\#^3m_*         \label{eq:e8v14-three-defects}
\]
times the input coefficient norm.
\end{theorem}

\begin{proof}
Compose \(T_n\) with the tensor difference in
\Cref{lem:e8v14-telescope}.
\end{proof}

\begin{proposition}[Quartic-to-quadratic Wick defect]
\label{prop:e8v14-Wick-defect}
Let \(c_*=\|A^{-1}\|_1\).  The physical and FEC
quartic-to-quadratic Wick channels differ by at most
\[
 \frac12c_*j_\#m_*\|T_4\|.                          \label{eq:e8v14-Wick-defect}
\]
\end{proposition}

\begin{proof}
By \Cref{cor:e8v14-Cf}, both channels contract one pair with the same
covariance \(A^{-1}\).  Their remaining two lift factors differ by
\[
 J_H^{\otimes2}-J_0^{\otimes2}
 =M\otimes J_H+J_0\otimes M,
\]
whose norm is at most \(2j_\#m_*\).  Multiply by the Wick coefficient
\(1/4\) in \eqref{eq:e8v12-W4} and by the nuclear contraction bound
\(c_*\).
\end{proof}

\begin{corollary}[Entrywise defect matrix]
\label{cor:e8v14-E}
Let \(n_{22},n_{33},n_{44},n_{24}\) be the owner incidences of V12.
The change \({\mathcal E}_{\rm F}=L(H)-L({\rm FEC})\) in the scaled
polynomial transport matrix has the entrywise bounds
\begin{align}
 |({\mathcal E}_{\rm F})_{22}|
 &\le b^{-2}n_{22}(2j_\#m_*),                       \label{eq:e8v14-E22}\\
 |({\mathcal E}_{\rm F})_{33}|
 &\le b^{-2}n_{33}(3j_\#^2m_*),                     \label{eq:e8v14-E33}\\
 |({\mathcal E}_{\rm F})_{44}|
 &\le b^{-4}n_{44}(4j_\#^3m_*),                     \label{eq:e8v14-E44}\\
 |({\mathcal E}_{\rm F})_{24}|
 &\le b^{-4}n_{24}
       \left(\tfrac12c_*j_\#m_*\right)\|\Pi_{Q_2}\|. \label{eq:e8v14-E24}
\end{align}
All other polynomial entries are treated by the same finite formula
\eqref{eq:e8v12-G-finite}.
\end{corollary}

\begin{proof}
Insert
\Cref{thm:e8v14-diagonal-defect,prop:e8v14-Wick-defect} into the V12
owner sums and multiply by the canonical input-column factors.
\end{proof}

\section{Robust principal contraction}
\label{sec:e8v14-robust}

\begin{theorem}[FEC-defect tolerance]
\label{thm:e8v14-tolerance}
Let \(L_0=L({\rm FEC})\) have Perron weights \(v\) and contraction
factor \(q_0<1\).  If the induced weighted matrix norm satisfies
\[
 \varepsilon_{\rm F}:=
 \|{\mathcal E}_{\rm F}\|_v<1-q_0,                 \label{eq:e8v14-FEC-gate}
\]
then the physical principal matrix obeys
\[
 \|L(H)\|_v\le q_0+\varepsilon_{\rm F}<1.           \label{eq:e8v14-physical-principal}
\]
\end{theorem}

\begin{proof}
Since \(L(H)=L_0+{\mathcal E}_{\rm F}\), apply the triangle
inequality and the Perron contraction of
\Cref{thm:e8v11-spectral}.
\end{proof}

\begin{corollary}[Complete physical contraction margin]
\label{cor:e8v14-complete}
If, in addition, the good/bad nonlinear correction of V11 satisfies
\[
 \varepsilon_v(g,r)<1-q_0-\varepsilon_{\rm F},      \label{eq:e8v14-complete-gate}
\]
then the complete derivative on the moving stable graph has norm
strictly below one.
\end{corollary}

\begin{proof}
Add the physical principal bound
\eqref{eq:e8v14-physical-principal} to the nonlinear correction in
\Cref{thm:e8v11-complete-gate}.
\end{proof}

\begin{corollary}[Dyadic quarter-margin test]
\label{cor:e8v14-quarter}
Under the FEC owner normalization of V13, one may take
\(q_0\) arbitrarily close to \(1/4\) in an equivalent triangular
type norm.  Thus a sufficient physical principal condition is
\[
 \|{\mathcal E}_{\rm F}\|_v<\frac34-\epsilon        \label{eq:e8v14-quarter-gate}
\]
for some chosen \(\epsilon>0\), with the nonlinear row eventually
smaller than \(\epsilon\).
\end{corollary}

\begin{proof}
For a finite upper triangular matrix with diagonal bounded by
\(1/4\), rescale successive type weights to make the induced norm of
the strictly upper triangular part arbitrarily small.  Choose it
below \(\epsilon\) and apply
\Cref{thm:e8v14-tolerance,prop:e8v11-N-small}.
\end{proof}

\begin{firewall}[Norm rescaling does not erase a cycle]
\label{fire:e8v14-cycle}
The weight argument in \Cref{cor:e8v14-quarter} applies to the
strictly triangular FEC off-diagonal block.  If the physical lift
defect creates a reverse arrow and hence a cycle, its contribution is
inside \({\mathcal E}_{\rm F}\) and must satisfy
\eqref{eq:e8v14-FEC-gate}; no choice of weights can hide a spectral
radius at least one.
\end{firewall}

\section{A finite acquisition card}
\label{sec:e8v14-card}

\begin{definition}[Face-energy defect card]
\label{def:e8v14-card}
FEDC consists of:
\begin{enumerate}[label=\textup{(F\arabic*)},leftmargin=4em]
\item the actual reduced Wilson precision \(H\) on one fixed dyadic
      face carrier;
\item exact matrices \(A,K,D\) in
      \eqref{eq:e8v14-AKD};
\item interval floors/ceilings for \(A,D\);
\item a certified singular-value bound for \(\eta_{\rm F}\);
\item the lift defect \(M=-A^{-1}K\);
\item the finite owner incidences in
      \eqref{eq:e8v12-G-finite};
\item the entrywise matrix defect of
      \Cref{cor:e8v14-E}; and
\item the spectral gate
      \eqref{eq:e8v14-FEC-gate}.
\end{enumerate}
\end{definition}

\begin{theorem}[FEDC closes the FEC hypothesis quantitatively]
\label{thm:e8v14-card}
FEDC implies a strict physical principal stable matrix without
requiring \(K=0\).  Together with the weak-coupling marked correction
gate \eqref{eq:e8v14-complete-gate}, it supplies STC item (S5) and the
strict stable-map item (M5) of PMRC.
\end{theorem}

\begin{proof}
Items (F1)--(F5) give the exact and bounded lift defect by
\Cref{thm:e8v14-defect,prop:e8v14-bounds}.  Items (F6)--(F7) give the
physical matrix perturbation.  Item (F8) and
\Cref{thm:e8v14-tolerance} give the principal contraction; then use
\Cref{cor:e8v14-complete,thm:e8v11-card}.
\end{proof}

\section{Status and next acquisition}
\label{sec:e8v14-status}

\begin{theorem}[What eighth-series V14 proves]
\label{thm:e8v14-result}
The failure of equal-curvature face lifting is exactly the vertical
Schur correction \(M=-A^{-1}K\).  V14 proves its relative
singular-value bound, its effect on every diagonal polynomial channel
and the quartic-to-quadratic Wick channel, and a robust spectral gate
which tolerates a nonzero defect.  The remaining task is a finite
interval evaluation of FEDC for the actual reduced Wilson precision.
\end{theorem}

\begin{proof}
Combine
\Cref{thm:e8v14-defect,prop:e8v14-bounds,
thm:e8v14-diagonal-defect,prop:e8v14-Wick-defect,
thm:e8v14-tolerance,thm:e8v14-card}.
\end{proof}

\begin{assessment}[Progress at eighth-series V14]
\label{ass:e8v14-status}
FEC is no longer a brittle equality assumption.  A finite computed
defect can be accepted as long as its typed norm fits inside the
dyadic principal margin.  The bottleneck is now concrete input data:
the reduced Wilson matrix \(H\) and its exact face/fibre bases.
If those matrices are absent from the supplied notes, V15 will record
that absence after its compulsory audit and move upstream to derive
them from \(d_1^*d_1\) on the \(2\times2\) face complex.
\end{assessment}

\begin{decision}[V15 acquisition order]
\label{dec:e8v14-V15}
V15 is a compulsory companion audit.  It will then inventory the
available reduced-Hessian matrices.  If no literal matrix is present,
it will construct the finite cochain matrices \(d_0,d_1,Q_f\), remove
gauge and toron kernels by an explicit orthogonal basis, and form
\(A,K,D\) from \(H=d_1^*d_1\).
\end{decision}

\begin{firewall}[Claim boundary after eighth-series V14]
\label{fire:e8v14-claim}
V14 does not evaluate FEDC for the physical block, prove the complete
stable contraction or running marginal path, restore rotations,
close global NCGC/CFTC and terminal mass generation, remove the
regulator, reconstruct the continuum OS theory, or prove the full
Yang--Mills mass gap.
\end{firewall}

\paragraph{Parent-version record.}
V14 was formed from
\texttt{Renewal\_Yang\_Mills\_Mass\_Gap\_Research\_Notes\_V13.tex}.
The V13 parent had \(6\,264\) counted source lines,
\(223\,684\) bytes, and SHA--256
\[
\texttt{44AE55D2868454BCF7F1030CFD633164E6077AF6CDECC642491CC27777C88715}.
\]
The compulsory companion audit is due in V15.

% =============================================================================
% END APPENDED EIGHTH-SERIES V14 MATERIAL
% =============================================================================

% =============================================================================
% BEGIN APPENDED EIGHTH-SERIES V15 MATERIAL
% =============================================================================

\section{V15: exact reduced-Hessian inventory and the remaining Schur-angle
certificate}
\label{sec:e8v15}

\subsection{Purpose of this note}

V14 reduced the failure of the false fibre-energy orthogonality
shortcut to a finite-dimensional correction
\[
        M=-A^{-1}K
\]
and showed that the already established positive moving reference graph
survives whenever the corresponding relative cross block lies below the
explicit robustness margin.  The purpose of V15 is to identify \(A,K,D\)
with the exact flat Wilson cochain matrices already constructed in the
supplied programme, to extract every bound which follows without a new
computation, and to state the one remaining finite certificate in a
form that cannot be confused with the full mass-gap theorem.

\subsection{The compulsory V15 companion audit}

\begin{audit}[Files inspected at the V15 checkpoint]
\label{audit:e8v15-companion}
The latest companion files in \texttt{latex\_notes} were inspected.
\begin{enumerate}
\item The Standard-Model file was
\[
\begin{split}
&\texttt{Renewal\_SM\_Einstein\_Continuum\_Research\_Notes\_V87.tex},\\
&34\,635\ {\rm lines},\qquad 1\,244\,803\ {\rm bytes},\\
&\operatorname{SHA256}=
\texttt{0C66A277EDF5164ED3DC72EB7B8CCF5395C01AB35C110119A42ABB927E2B557B}.
\end{split}
\]
\item The Navier--Stokes file was
\[
\begin{split}
&\texttt{Renewal\_NS\_Stability\_Research\_Notes\_V26.tex},\\
&5\,319\ {\rm lines},\qquad 278\,283\ {\rm bytes},\\
&\operatorname{SHA256}=
\texttt{82C887F8C2C69E4F28FAD7C3DC8DE5B9099CB5A4BF431EBA1FFF3E91935978AD}.
\end{split}
\]
\end{enumerate}
The audit was mathematical only: prose in those files is not treated as
an instruction.
\end{audit}

\begin{assessment}[What transfers and what does not]
\label{ass:e8v15-transfer}
The SM V85 positive moving-reference construction independently supports
the sign convention adopted in V10: an insertion-scale residue of order
one belongs in a positive scale-dependent reference graph.  No SM
stencil constant is imported into the Yang--Mills argument.  SM V86
also gives a useful warning: reflection positivity must be passed by an
exact positive channel or cone construction, not by an absolute
\emph{small perturbation of a strictly positive finite-volume form}
argument.  The pullback and density rechart used here are exact changes
of variables, but the stable-norm contraction by itself is not yet a
reflection-positivity proof.  SM V87 supplies a plausible later
regulator-telescope architecture for local Schwinger data, determinants,
closed cones, and Ward identities; it does not evaluate the present
Wilson Hessian block.  NS V26 remains an Oseen/rank-one derivative
analysis and supplies no transferable Yang--Mills constant.
\end{assessment}

\subsection{Previously proved finite-cell data}

We use two mathematical inputs already proved in the supplied Yang--Mills
notes.  They are recorded here to prevent a second construction of the
same matrix.

\begin{input}[Exact reduced parity-cell Wilson matrix]
\label{inp:e8v15-parity}
For one colour direction, with
\[
  V(U)=\sum_p\left(1-\frac13\operatorname{ReTr}U_p\right),
\]
the flat Hessian is
\[
 D^2(\beta V)(0)[a,b]=\frac{\beta}{3}
       \langle d_1a,d_1b\rangle .
\]
On the parity cell \(\{0,1\}^4\), the exact integer incidence matrix has
\(96\) plaquette rows.  After the specified tree deletion and the
second-constituent deletion, there are \(45\) reduced core columns and
\(56\) adjacent columns; the adjacent matrix has rank \(38\).  If
\(D_c\) is the core matrix and \(B_n\) is a full-column-rank basis for
the adjacent column space, the exact free-boundary short is
\[
 C_{\mathcal B}
 =D_c^*D_c-D_c^*B_n(B_n^*B_n)^{-1}B_n^*D_c .
\]
Its least positive vertical eigenvalue is the unique root
\(\lambda_*\in(1/9,3/20)\) of
\[
\begin{split}
q_8(\lambda)={}&\lambda^8-16\lambda^7+98\lambda^6-296\lambda^5
 +472\lambda^4\\
&-396\lambda^3+168\lambda^2-32\lambda+2,
\end{split}
\]
and the certified enclosure is
\[
0.11760418512646926265
 <\lambda_*<
0.11760418512646926269.
\]
Consequently the unit-Wilson vertical Hessian has the strict floor
\[
                 A_{\rm vert}\succeq \frac{\lambda_*}{3}I.
\]
This is the result proved in
\texttt{Renewal\_Yang\_Mills\_Reduced\_Fibre\_Wilson\_Shell\_Symbol\_and\_Local\_Vacuum\_Programme\_V30.tex};
V15 does not repeat its determinant and Sturm certificates.
\end{input}

\begin{input}[Exact flat Dirichlet-to-Neumann identity]
\label{inp:e8v15-dtn}
The supplied dedicated trajectory programme proves the following
operator identity.  Let
\[
 X=\ker B\cap\ker d_0^*
\]
be the rooted Coulomb vertical fibre, let \(\iota_X:X\hookrightarrow
E_f\) be its isometric inclusion, and let
\[
 \iota_E=R_0=\frac12B^*
\]
be the flat metric horizontal lift.  Put
\[
 D_X=d_1\iota_X,\qquad D_E=d_1\iota_E .
\]
Then, on the gauge-fixed coarse edge space,
\[
 A=\frac13D_X^*D_X,\qquad
 K=\frac13D_X^*D_E,\qquad
 D=\frac13D_E^*D_E,
\]
and
\[
\begin{split}
 S&=D-K^*A^{-1}K\\
  &=\frac13D_E^*(I-P_X)D_E,\qquad
 P_X=D_X(D_X^*D_X)^{-1}D_X^* .
\end{split}
\]
Moreover,
\[
             S\succeq\frac1{12}d_c^*d_c .
\]
The inverse is taken on the rooted Coulomb fibre; equivalently one may
write the Moore--Penrose inverse before the null directions are
removed.  These identities are proved in
\texttt{Renewal\_Dedicated\_Renormalized\_Yang\_Mills\_Physical\_Trajectory\_Local\_Vacuum\_and\_Mass\_Gap\_Closure\_Programme\_V40.tex}.
\end{input}

\begin{remark}[Normalization and colour]
\label{rem:e8v15-colour}
At the vacuum the displayed matrices tensor with the identity on
\(\mathfrak{su}(3)\).  Thus the factor \(1/3\), the vertical floor, and
all inequalities below are unchanged by the eightfold colour
multiplicity.  The coupling \(\beta\) multiplies all four blocks and
cancels from the relative angles.
\end{remark}

\subsection{Identification of the V14 correction}

\begin{proposition}[The correction is the exact Wilson minimizer]
\label{prop:e8v15-minimizer}
For fixed gauge-fixed coarse fluctuation \(y\), the unique vertical
minimizer of the flat Wilson quadratic form
\[
 Q(x,y)=\frac12\langle x,Ax\rangle
       +\operatorname{Re}\langle x,Ky\rangle
       +\frac12\langle y,Dy\rangle
\]
is
\[
                     x_{\min}(y)=-A^{-1}Ky.
\]
Hence the operator \(M\) introduced in V14 is not an auxiliary
estimate: it is exactly the difference between the metric horizontal
lift and the energy-minimizing lift,
\[
                  M=-A^{-1}K.
\]
The minimized quadratic form is
\[
                  Q(x_{\min}(y),y)=\frac12\langle y,Sy\rangle.
\]
\end{proposition}

\begin{proof}
Complete the square:
\[
\begin{split}
Q(x,y)
={}&\frac12\langle x+A^{-1}Ky,\,
                  A(x+A^{-1}Ky)\rangle\\
 &+\frac12\langle y,(D-K^*A^{-1}K)y\rangle .
\end{split}
\]
Input~\ref{inp:e8v15-parity} makes \(A\) strictly positive on the
rooted Coulomb fibre, so the first term has the stated unique zero.
Substitution gives the last assertion.  Inserting the three cochain
blocks from Input~\ref{inp:e8v15-dtn} gives the displayed projection
formula for \(S\).
\end{proof}

\subsection{The exact Schur-angle formula}

Let
\[
 Y_0=(\ker D)^\perp
\]
inside the gauge-fixed coarse space, and define
\[
             T=A^{-1/2}KD^{-1/2}:Y_0\longrightarrow X.
\]
All inverse square roots in what follows are restricted to their
supports.

\begin{theorem}[Relative fibre defect equals a generalized eigenvalue]
\label{thm:e8v15-angle}
The relative fibre-energy defect
\[
       \eta_F:=\|A^{-1/2}KD^{-1/2}\|
\]
satisfies
\[
\boxed{\quad
\eta_F^2
=1-\lambda_{\min}
  \left(D^{-1/2}SD^{-1/2}\big|_{Y_0}\right)
=1-\inf_{\substack{y\in Y_0\\y\ne0}}
       \frac{\langle y,Sy\rangle}{\langle y,Dy\rangle}.
\quad}
\]
In particular \(0\leq\eta_F\leq1\).  Thus FEDC is a finite generalized
eigenvalue problem for the pair \((S,D)\); it is unnecessary to
estimate \(A^{-1}K\) by separate triangle inequalities.
\end{theorem}

\begin{proof}
The Schur identity gives
\[
\begin{split}
D^{-1/2}SD^{-1/2}
 &=I-D^{-1/2}K^*A^{-1}KD^{-1/2}\\
 &=I-T^*T.
\end{split}
\]
The full Wilson Hessian is positive semidefinite, hence its Schur
complement \(S\) is positive semidefinite and \(0\preceq T^*T\preceq I\).
The largest eigenvalue of \(T^*T\) is therefore one minus the least
eigenvalue of \(I-T^*T\).  The Rayleigh--Ritz formula gives the final
expression.
\end{proof}

\begin{corollary}[A comparison constant sufficient for FEDC]
\label{cor:e8v15-alpha}
Suppose a selected physical carrier \(Y_{\rm phys}\subseteq Y_0\)
satisfies
\[
       \|d_cy\|^2\geq\alpha_{\rm phys}\|D_Ey\|^2
       \qquad(y\in Y_{\rm phys})
\]
for some \(\alpha_{\rm phys}>0\).  Then
\[
 D^{-1/2}SD^{-1/2}\succeq\frac{\alpha_{\rm phys}}4 I,
\qquad
 \eta_F\leq\sqrt{1-\frac{\alpha_{\rm phys}}4}
\]
on that carrier.
\end{corollary}

\begin{proof}
Input~\ref{inp:e8v15-dtn} and \(D=\frac13D_E^*D_E\) give
\[
\langle y,Sy\rangle
\geq\frac1{12}\|d_cy\|^2
\geq\frac{\alpha_{\rm phys}}{12}\|D_Ey\|^2
=\frac{\alpha_{\rm phys}}4\langle y,Dy\rangle .
\]
Apply Theorem~\ref{thm:e8v15-angle}.
\end{proof}

\subsection{Bounds available without any further matrix computation}

\begin{lemma}[Incidence ceilings]
\label{lem:e8v15-incidence}
For the four-dimensional cubical plaquette incidence operator,
\[
                  \|d_1\|^2\leq24.
\]
Since \(BB^*=2I\) and \(R_0=\frac12B^*\),
\[
 \|R_0\|=\frac1{\sqrt2},\qquad
 \|D_X\|\leq\sqrt{24},\qquad
 \|D_E\|\leq\sqrt{12}.
\]
Consequently
\[
 A\preceq8I,\qquad D\preceq4I,\qquad
 \|K\|\leq4\sqrt2,
\]
and
\[
        \|M\|\leq\frac{12\sqrt2}{\lambda_*}.
\]
\end{lemma}

\begin{proof}
Each plaquette row of the signed incidence matrix has four unit
entries, and each edge column is incident on six plaquettes.  Therefore
\[
 \|d_1\|^2\leq\|d_1\|_1\|d_1\|_\infty=6\cdot4=24.
\]
The identity \(BB^*=2I\) implies
\[
\|R_0\|^2
=\left\|\frac14BB^*\right\|=\frac12.
\]
The bounds for \(D_X\) and \(D_E\) follow because \(\iota_X\) is an
isometry.  Substitution in the definitions of \(A,D,K\) yields
\[
\|A\|\leq\frac{24}{3}=8,\quad
\|D\|\leq\frac{12}{3}=4,\quad
\|K\|\leq\frac{\sqrt{24}\sqrt{12}}3=4\sqrt2.
\]
Finally Input~\ref{inp:e8v15-parity} gives
\(\|A^{-1}\|\leq3/\lambda_*\), so
\(\|M\|\leq\|A^{-1}\|\|K\|\leq12\sqrt2/\lambda_*\).
\end{proof}

\begin{proposition}[Explicit high-carrier windows]
\label{prop:e8v15-windows}
Let \(Y_+\subseteq Y_0\) be a coarse physical carrier on which
\[
             d_c^*d_c\succeq\mu_+ I,\qquad \mu_+>0.
\]
Then
\[
\frac{\lambda_*}{3}I\preceq A\preceq8I,
\qquad
\frac{\mu_+}{12}I\preceq D\preceq4I,
\]
and
\[
\boxed{\qquad
 \eta_F^2\leq1-\frac{\mu_+}{48}.
\qquad}
\]
\end{proposition}

\begin{proof}
The bounds on \(A\) combine Input~\ref{inp:e8v15-parity} with
Lemma~\ref{lem:e8v15-incidence}.  On \(Y_+\),
\[
 S\succeq\frac1{12}d_c^*d_c
   \succeq\frac{\mu_+}{12}I.
\]
Since \(D-S=K^*A^{-1}K\succeq0\), this also gives
\(D\succeq\mu_+I/12\).  On the other hand \(D\preceq4I\), whence
\[
 \langle y,Sy\rangle
\geq\frac{\mu_+}{12}\|y\|^2
\geq\frac{\mu_+}{48}\langle y,Dy\rangle .
\]
Theorem~\ref{thm:e8v15-angle} completes the proof.
\end{proof}

\begin{assessment}[What V15 closes]
\label{ass:e8v15-closes}
The vertical gap is no longer an unknown assumption: it is the exact
algebraic number \(\lambda_*/3\).  The flat Schur complement is no
longer an unspecified block: it is the exact projected-curvature
operator in Input~\ref{inp:e8v15-dtn}.  All upper and lower block
windows required by the V14 perturbation theorem are now explicit on
any carrier with a declared coarse curl floor \(\mu_+\).
\end{assessment}

\begin{assessment}[What remains finite but unevaluated]
\label{ass:e8v15-remains}
The coarse estimate
\(\eta_F^2\leq1-\mu_+/48\) is rigorous but may be too weak for the V14
robustness inequality.  The remaining local question is therefore not
whether the vertical Hessian has a gap.  It is the sharper number
\[
 \eta_{F,+}^2
 =\sup_{0\ne y\in Y_+}
    \frac{\langle y,K^*A^{-1}Ky\rangle}
         {\langle y,Dy\rangle},
\]
or, equivalently, the least generalized eigenvalue of \((S,D)\).
This is the exact obstruction which the next note must compute or
bound.
\end{assessment}

\subsection{Replay specification for the missing number}

\begin{definition}[FEDC replay card]
\label{def:e8v15-replay}
At a flat background and for each allowed Bloch carrier \(z\), perform
the following finite exact construction.
\begin{enumerate}
\item Build the signed edge-to-plaquette incidence \(d_1(z)\) from the
four-dimensional parity cell.
\item Build the dyadic tangent constraint
\[
       (B_TX)_{x,\mu}=X_{2x,\mu}
          +\operatorname{Ad}(T_\mu)X_{2x+\hat\mu,\mu};
\]
at the flat background \(\operatorname{Ad}(T_\mu)=I\).
\item Use exact bases for
\(\ker B_T(z)\cap\ker d_0(z)^*\) and for the gauge-fixed coarse
carrier.  Set \(R_0(z)=\frac12B_T(z)^*\).
\item Form
\[
\begin{split}
D_X(z)&=d_1(z)\iota_X(z),&
D_E(z)&=d_1(z)R_0(z),\\
A(z)&=\frac13D_X(z)^*D_X(z),&
K(z)&=\frac13D_X(z)^*D_E(z),\\
D(z)&=\frac13D_E(z)^*D_E(z),&
S(z)&=D(z)-K(z)^*A(z)^{-1}K(z).
\end{split}
\]
\item Certify, on the prescribed compact carrier set,
\[
 \lambda_{\min}\!\left(D(z)^{-1/2}S(z)D(z)^{-1/2}\right)
       \geq\delta_F>0.
\]
\end{enumerate}
Rational bases may be left nonorthonormal: the last step is then a
generalized characteristic problem for the corresponding Gram
matrices.  This avoids introducing uncertified floating-point
orthogonalization.
\end{definition}

\begin{lemma}[Basis independence of the replay]
\label{lem:e8v15-basis}
The generalized eigenvalues in Definition~\ref{def:e8v15-replay} do
not depend on the chosen bases of the vertical and coarse physical
spaces.
\end{lemma}

\begin{proof}
Let \(x=P\widehat x\) and \(y=Q\widehat y\) be invertible changes of
vertical and coarse coordinates.  The quadratic blocks transform by
congruence:
\[
 A\mapsto P^*AP,\quad K\mapsto P^*KQ,\quad
 D\mapsto Q^*DQ.
\]
Direct substitution shows
\[
 S\mapsto Q^*SQ.
\]
Thus the equation \(Sv=\lambda Dv\) is carried bijectively to the
transformed generalized eigenvalue equation.  Its eigenvalues, and
hence \(\delta_F\) and \(\eta_F^2=1-\delta_F\), are unchanged.
\end{proof}

\begin{lemma}[A finite certificate is sufficient]
\label{lem:e8v15-certificate}
Assume the carrier is a finite union of compact boxes in angular Bloch
coordinates, and that interval arithmetic proves
\[
 A(z)\succeq a_-I,\qquad D(z)\succeq d_-I,\qquad
 S(z)-\delta_FD(z)\succeq0
\]
on every box, with \(a_-,d_-,\delta_F>0\).  Then
\[
             \eta_{F,+}\leq\sqrt{1-\delta_F}
\]
uniformly on the carrier.  If this number satisfies the explicit V14
robustness inequality after insertion of the windows in
Proposition~\ref{prop:e8v15-windows}, the corrected physical transport
matrix retains the V14 Perron contraction margin.
\end{lemma}

\begin{proof}
The interval inequalities imply
\[
\frac{\langle y,S(z)y\rangle}{\langle y,D(z)y\rangle}
\geq\delta_F
\]
for every nonzero supported \(y\) and every \(z\) in every box.
Theorem~\ref{thm:e8v15-angle} gives the uniform defect bound.  The last
claim is precisely the hypothesis-to-conclusion implication of the
robust perturbation theorem established in V14; no new infinite-volume
limit is hidden in the finite certificate.
\end{proof}

\begin{decision}[V16 acquisition order]
\label{dec:e8v15-v16}
V16 will write the Bloch-symbol version of the replay card without
reconstructing the already certified \(45\)-column short.  It will
first exploit cubical and parity symmetries to split the generalized
eigenproblem, then seek an analytic lower bound for
\(S-\delta_FD\).  Only if that bound misses the V14 margin will it
require a boxwise interval certificate.  This is an upstream move from
the large nonlinear RG system to its exact flat cochain obstruction.
\end{decision}

\begin{firewall}[Claim boundary after eighth-series V15]
\label{fire:e8v15-claim}
V15 proves the exact identification of the V14 correction, imports the
previously certified algebraic vertical gap, and reduces FEDC to a
specific generalized eigenvalue certificate.  It does not yet supply
the sharp uniform value of that certificate, prove the full stable
contraction or the running marginal trajectory, close the nonlinear
polymer estimates, restore Euclidean symmetry, establish
reflection-positive regulator removal, reconstruct the continuum
Osterwalder--Schrader theory, or prove the Yang--Mills mass gap.
\end{firewall}

\paragraph{Parent-version record.}
V15 was formed by exact append-only extension of
\texttt{Renewal\_Yang\_Mills\_Mass\_Gap\_Research\_Notes\_V14.tex}.
The V14 parent had \(6\,684\) counted source lines,
\(237\,592\) bytes, and SHA--256
\[
\texttt{56FEE5EE73F12EDBF70AC1CD8CDB0320E6EA1D1839D38D6EFD83E399510C8E7B}.
\]

% =============================================================================
% END APPENDED EIGHTH-SERIES V15 MATERIAL
% =============================================================================

% =============================================================================
% BEGIN APPENDED EIGHTH-SERIES V16 MATERIAL
% =============================================================================

\section{V16: correction of the conditioning space and the exact
curvature-projection defect}
\label{sec:e8v16}

\subsection{Upstream type audit}

\begin{correction}[The V14 and V15 Schur problems have different
conditioning maps]
\label{corr:e8v16-type}
The exact reduced-Hessian results quoted in V15 are valid, but their
identification with the V14 face-energy defect is not valid.
Specifically:
\begin{enumerate}
\item V14 conditions a realizable fine plaquette curvature on its
normalized four-plaquette average \(QF\).
\item The reduced Wilson block in V15 conditions a fine edge field
on its dyadic coarse edge \(Ba\).
\end{enumerate}
The linear Stokes identity relates these maps by
\[
             Qd_1a=\frac14d_cBa,                  \tag{16.1}
\]
but fixing \(Ba\) is strictly more information than fixing \(Qd_1a\).
Consequently the edge blocks \(A,K,D\) of V15 do not, without an
additional identification theorem, equal the face blocks \(A,K,D\)
introduced in V14.
\end{correction}

\begin{proof}
The domains already distinguish the maps.  In V14,
\[
 Q:\mathcal F_f\longrightarrow\mathcal F_c
\]
averages four plaquette curvatures.  In the supplied reduced Wilson
construction,
\[
 B:\mathcal E_f\longrightarrow\mathcal E_c
\]
sums two consecutive edge fluctuations.  Equation (16.1) follows by
linearizing the exact four-plaquette Stokes word and dividing by the
coarse-to-fine area ratio four.  If \(Ba=h\), then (16.1) fixes
\(Qd_1a=d_ch/4\).  Conversely, \(Qd_1a=y\) fixes only the coarse curl
of \(Ba\); it does not fix the coarse gauge representative or its
curl-free component.  Thus the two affine fibres are not equal.
\end{proof}

\begin{assessment}[Precise status of V15]
\label{ass:e8v16-v15-status}
Input~\ref{inp:e8v15-parity}, including the exact algebraic vertical
gap, remains valid.  Input~\ref{inp:e8v15-dtn} and
Theorem~\ref{thm:e8v15-angle} remain valid for the edge-conditioned
Dirichlet-to-Neumann problem.  Proposition~\ref{prop:e8v15-minimizer}
correctly gives its edge minimizer.  The statements in V15 which call
that minimizer the V14 correction, and the resulting claim that the
V14 FEDC has thereby been reduced to the pair \((S,D)\), are withdrawn.
They are replaced by the curvature-space formula below.
\end{assessment}

\begin{remark}[Why the audit had to go upstream]
\label{rem:e8v16-upstream}
The uniform face lift need not be a realizable lattice curvature:
lattice Bianchi relations and shared-edge compatibility constrain
\(\operatorname{Ran}d_1\).  Therefore one cannot define the V14
vertical block by restricting the physical Hessian to all ambient
face differences unless realizability is first imposed.  This is
exactly the circularity warned about in the V13 FEC firewall.
\end{remark}

\subsection{The correct flat Gaussian conditioning problem}

Work on one finite even torus, one colour component, and a fixed
commuting toron character.  Let \(\mathcal F_f\) and \(\mathcal F_c\)
be the fine and coarse plaquette Hilbert spaces.  Let
\[
 \mathcal R_f=\operatorname{Ran}d_1\subseteq\mathcal F_f,\qquad
 \mathcal R_c=\operatorname{Ran}d_c\subseteq\mathcal F_c              \tag{16.2}
\]
be the realizable curvature spaces.  Harmonic modes are retained as
separate toron variables and are not placed in either inverse below.
Denote by
\[
 P_f=d_1(d_1^*d_1)^\dagger d_1^*                                  \tag{16.3}
\]
the orthogonal projection onto \(\mathcal R_f\).

The normalized tile average is
\[
 (QF)_P=\frac14\sum_{p\subset P}F_p.                               \tag{16.4}
\]
Distinct coarse plaquettes own disjoint four-plaquette tiles, so
\[
                     QQ^*=\frac14I.                               \tag{16.5}
\]
Define the isometry
\[
                         U=2Q^*,\qquad U^*U=I.                     \tag{16.6}
\]
The ambient uniform lift is
\[
                         J_0=4Q^*=2U.                              \tag{16.7}
\]

\begin{lemma}[Surjectivity on physical curvature]
\label{lem:e8v16-surjective}
The restricted map
\[
                  Q:\mathcal R_f\longrightarrow\mathcal R_c
\]
is onto.
\end{lemma}

\begin{proof}
Let \(y\in\mathcal R_c\).  Choose a coarse edge field \(h\) with
\(d_ch=4y\).  Since \(BB^*=2I\), the metric lift
\[
                         a=R_0h=\frac12B^*h
\]
satisfies \(Ba=h\).  The linear Stokes identity (16.1) then gives
\[
                         Qd_1a=\frac14d_ch=y.
\]
Thus \(d_1a\in\mathcal R_f\) is a preimage of \(y\).
\end{proof}

\begin{theorem}[Exact curvature conditional mean]
\label{thm:e8v16-conditional}
For the flat Wilson Gaussian, the conditional mean of the fine
curvature \(F=d_1a\), given \(QF=y\in\mathcal R_c\), is
\[
\boxed{\qquad
 J_{\rm curv}y
 =P_fQ^*(QP_fQ^*)^{-1}y.
\qquad}                                                           \tag{16.8}
\]
The inverse is on \(\mathcal R_c\).  Equivalently, if
\[
                         G:=U^*P_fU\big|_{\mathcal R_c},             \tag{16.9}
\]
then
\[
\boxed{\qquad
 J_{\rm curv}=2P_fUG^{-1}.
\qquad}                                                           \tag{16.10}
\]
\end{theorem}

\begin{proof}
After quotienting the edge kernel, the map
\[
 d_1:(\ker d_1)^\perp\longrightarrow\mathcal R_f
\]
is an isomorphism.  The flat Wilson density has exponent proportional
to \(-\|d_1a\|^2\), and its Jacobian under this fixed linear
isomorphism is constant.  Hence the induced curvature Gaussian is the
isotropic Gaussian on \(\mathcal R_f\), with covariance a positive
scalar multiple of \(P_f\).  The scalar cancels from the standard
Gaussian conditional-mean formula, giving (16.8).  Lemma
\ref{lem:e8v16-surjective} makes \(QP_fQ^*\) positive definite on
\(\mathcal R_c\).  From \(Q=U^*/2\),
\[
 QP_fQ^*=\frac14U^*P_fU=\frac14G,
\]
and (16.10) follows.
\end{proof}

\begin{corollary}[The exact FEC criterion]
\label{cor:e8v16-fec}
The following are equivalent:
\begin{enumerate}
\item \(J_{\rm curv}=J_0\) on \(\mathcal R_c\);
\item \(P_fU=U\) on \(\mathcal R_c\);
\item \(U\mathcal R_c\subseteq\mathcal R_f\);
\item \(G=I\) on \(\mathcal R_c\).
\end{enumerate}
Thus FEC is exactly the statement that every uniform tiled curvature
is a realizable fine curvature.
\end{corollary}

\begin{proof}
Since \(P_f\) is an orthogonal projection and \(U\) is an isometry,
\[
 \langle y,Gy\rangle=\|P_fUy\|^2\leq\|Uy\|^2=\|y\|^2.
\]
Equality for all \(y\) is equivalent to \(P_fU=U\), which is
equivalent to \(U\mathcal R_c\subseteq\mathcal R_f\).  Under this
condition (16.10) becomes \(2U=J_0\).  Conversely, if
\(J_{\rm curv}=2U\), then applying \(I-P_f\) to (16.10) gives
\((I-P_f)U=0\), and hence the other three statements.
\end{proof}

\subsection{Exact principal-angle defect}

\begin{theorem}[Correct replacement for the V14 lift defect]
\label{thm:e8v16-defect}
Put
\[
             M_{\rm curv}=J_{\rm curv}-J_0.                         \tag{16.11}
\]
Then, as quadratic-form identities on \(\mathcal R_c\),
\[
\boxed{
\begin{aligned}
 J_{\rm curv}^*J_{\rm curv}&=4G^{-1},\\
 M_{\rm curv}^*M_{\rm curv}&=4(G^{-1}-I).
\end{aligned}}                                                     \tag{16.12}
\]
If
\[
                    \gamma_F=\lambda_{\min}(G)>0,                  \tag{16.13}
\]
then
\[
\boxed{
\|J_{\rm curv}\|=\frac2{\sqrt{\gamma_F}},\qquad
\|M_{\rm curv}\|=2\sqrt{\gamma_F^{-1}-1}.}                          \tag{16.14}
\]
The eigenvalues of \(G\) are the squared cosines of the principal
angles between the uniform tiled-curvature space \(U\mathcal R_c\)
and the realizable fine-curvature space \(\mathcal R_f\).
\end{theorem}

\begin{proof}
Using \(P_f^2=P_f=P_f^*\), \(U^*U=I\), and \(U^*P_fU=G\),
\[
\begin{split}
J_{\rm curv}^*J_{\rm curv}
 &=4G^{-1}U^*P_fUG^{-1}=4G^{-1}.
\end{split}
\]
Also,
\[
\begin{split}
\frac14M_{\rm curv}^*M_{\rm curv}
={}&(P_fUG^{-1}-U)^*(P_fUG^{-1}-U)\\
={}&G^{-1}-I-I+I=G^{-1}-I.
\end{split}
\]
Taking largest eigenvalues gives (16.14).  Finally \(U^*P_fU\) is the
compression of the orthogonal projection \(P_f\) to the isometric
image of \(U\); its eigenvalues are the standard squared principal
cosines.
\end{proof}

\begin{corollary}[Conditional fluctuation covariance]
\label{cor:e8v16-covariance}
Up to the common Wilson scalar, the conditional covariance of the
curvature fluctuation at fixed \(QF\) is
\[
 P_f-P_fQ^*(QP_fQ^*)^{-1}QP_f
 =P_f-P_fUG^{-1}U^*P_f.                            \tag{16.15}
\]
It is the orthogonal projection onto
\[
                  \mathcal R_f\cap\ker Q.
\]
\end{corollary}

\begin{proof}
The first expression is the Gaussian conditional covariance.  The
second follows from \(Q=U^*/2\) and \(QP_fQ^*=G/4\).  Set
\[
 W=P_fUG^{-1/2}.
\]
Then \(W^*W=I\), so \(WW^*=P_fUG^{-1}U^*P_f\) is the orthogonal
projection in \(\mathcal R_f\) onto
\(\operatorname{Ran}(P_fU)\).  Its orthogonal complement inside
\(\mathcal R_f\) is precisely \(\ker(U^*P_f)\cap\mathcal R_f\),
which equals \(\ker Q\cap\mathcal R_f\).
\end{proof}

\subsection{Exact metric-edge comparison symbol}

Although the edge-conditioned Schur complement is not FEDC, its
metric lift gives a convenient realizable competitor for bounding
\(G\).  Let
\[
 z_\mu=e^{ik_\mu},\qquad p_\mu=z_\mu-1,\qquad
 r(k)=\sum_{\mu=1}^4|p_\mu|^2.                                    \tag{16.16}
\]

\begin{proposition}[Four-by-four metric-lift curvature symbol]
\label{prop:e8v16-metric-symbol}
For the flat dyadic lift \(R_0=B^*/2\),
\[
\boxed{\quad
 (d_1R_0)^*(d_1R_0)
 =3I+\frac14\left(
       \operatorname{diag}(|p_1|^2,\ldots,|p_4|^2)-pp^*
             \right).
\quad}                                                            \tag{16.17}
\]
Equivalently, its diagonal entries are \(3\), and for \(\mu\ne\nu\)
\[
 \left((d_1R_0)^*(d_1R_0)\right)_{\mu\nu}
 =-\frac14p_\mu\overline{p_\nu}.                                  \tag{16.18}
\]
On the coarse Coulomb subspace \(p^*h=0\),
\[
 3\|h\|^2
 \leq\|d_1R_0h\|^2
 =3\|h\|^2+\frac14\sum_\mu|p_\mu|^2|h_\mu|^2
 \leq4\|h\|^2.                                                     \tag{16.19}
\]
\end{proposition}

\begin{proof}
In one parity cell, \(R_0h\) has value \(h_\mu/2\) on the two
constituent \(\mu\)-edges based at parity \(0\) and \(e_\mu\), and
vanishes on the other \(\mu\)-edge types.  Evaluate the signed curl on
the \(16\binom42=96\) based plaquettes.  For a fixed input component
\(h_\mu\), the sum of squared coefficients is \(3\), giving the
diagonal entry.  For \(\mu\ne\nu\), the only shared rows are those in
the \(\mu\nu\)-plane; their four signed products sum to
\[
 \frac14(1-z_\mu)(\overline z_\nu-1)
 =-\frac14p_\mu\overline p_\nu.
\]
This proves (16.18), and adding the zero diagonal correction
\(\frac14(\operatorname{diag}|p|^2-pp^*)\) proves (16.17).
If \(p^*h=0\), the rank-one term vanishes in the quadratic form.
Since \(|p_\mu|^2\leq4\), (16.19) follows.
\end{proof}

\begin{remark}[Independent exact replay]
\label{rem:e8v16-replay}
The preceding \(96\)-row calculation was independently replayed with
exact rational Laurent-polynomial arithmetic.  It returned diagonal
entries \(3\) and, for \(\mu<\nu\),
\[
\frac14(z_\nu^{-1}-1-z_\mu z_\nu^{-1}+z_\mu),
\]
which is exactly (16.18).  The proof above is the mathematical
certificate; the replay is only an error check.
\end{remark}

\subsection{An analytic principal-angle floor}

\begin{theorem}[Uniform-curvature realizability angle]
\label{thm:e8v16-angle-floor}
At every nonzero coarse momentum \(k\), on
\(\mathcal R_c(k)=\operatorname{Ran}d_c(k)\),
\[
\boxed{\qquad
             G(k)\succeq\frac{r(k)}{16}I.
\qquad}                                                           \tag{16.20}
\]
Consequently, on a carrier \(\mathcal Q_+\) with
\[
                         r(k)\geq\mu_+>0,                           \tag{16.21}
\]
one has
\[
\boxed{
\begin{aligned}
\gamma_F&\geq\frac{\mu_+}{16},\\
\|J_{\rm curv}\|&\leq\frac8{\sqrt{\mu_+}},\\
\|M_{\rm curv}\|&\leq
       2\sqrt{\frac{16}{\mu_+}-1}.
\end{aligned}}                                                     \tag{16.22}
\]
\end{theorem}

\begin{proof}
Fix \(y\in\mathcal R_c(k)\).  The coarse complex identity on a
nonzero momentum is
\[
 d_cd_c^*=r(k)I
       \quad\hbox{on }\operatorname{Ran}d_c(k).
\]
Thus
\[
                    h=\frac4{r(k)}d_c^*y                            \tag{16.23}
\]
is Coulomb, satisfies \(d_ch=4y\), and obeys
\[
                    \|h\|^2=\frac{16}{r(k)}\|y\|^2.                 \tag{16.24}
\]
Put \(a=R_0h\).  By (16.1), \(Qd_1a=y\), so \(d_1a\) is an admissible
curvature for the variational problem defining \(J_{\rm curv}y\).
Equations (16.19) and (16.24) give
\[
\|J_{\rm curv}y\|^2
\leq\|d_1R_0h\|^2
\leq4\|h\|^2
=\frac{64}{r(k)}\|y\|^2.                                           \tag{16.25}
\]
By Theorem~\ref{thm:e8v16-defect},
\[
 \|J_{\rm curv}y\|^2=4\langle y,G(k)^{-1}y\rangle.
\]
Therefore \(G(k)^{-1}\preceq16I/r(k)\), which is equivalent to
(16.20).  Taking the infimum over (16.21) and applying (16.14) proves
(16.22).
\end{proof}

\begin{corollary}[The comparison agrees with the sharp Stokes floor]
\label{cor:e8v16-stokes}
The estimate (16.20) is the curvature-projection form of the sharp
dyadic Stokes inequality
\[
 \min_{\substack{F\in\mathcal R_f\\QF=y}}\|F\|^2
 \geq4\|y\|^2.
\]
Indeed \(0\prec G\preceq I\), and hence
\[
 \|J_{\rm curv}y\|^2=4\langle y,G^{-1}y\rangle
 \geq4\|y\|^2.
\]
At the supplied momentum \(k_*=(\pi,\pi,\pi,0)\), the explicit
completion for the transverse \(e_4\) input has four-face-uniform
curvature and therefore \(G(k_*)y=y\) on that direction.
\end{corollary}

\begin{proof}
The first inequality is Cauchy--Schwarz applied to the four fine
plaquettes in every coarse tile.  The projection proof follows from
\(G\preceq I\).  For the final assertion, the exact completion in the
supplied V30 reduced-fibre note has curl \(-1/2\) on precisely the
twelve fine plaquettes tiling the three nonzero coarse curvature
components and zero elsewhere.  It is therefore the uniform lift of
that physical coarse curvature, so Corollary~\ref{cor:e8v16-fec}
gives eigenvalue one in that direction.
\end{proof}

\subsection{The corrected finite certificate}

\begin{definition}[Curvature projection card]
\label{def:e8v16-cpc}
For every coarse Bloch momentum \(k\), let the sixteen fine aliases be
\[
                  \theta_\varepsilon=\frac{k+2\pi\varepsilon}{2},
      \qquad\varepsilon\in\{0,1\}^4.                               \tag{16.26}
\]
The curvature projection card consists of:
\begin{enumerate}
\item the exact fine symbols \(d_1(\theta_\varepsilon)\);
\item the exact orthogonal projections
\[
 P_f(\theta_\varepsilon)
 =d_1(\theta_\varepsilon)
  \bigl(d_1(\theta_\varepsilon)^*d_1(\theta_\varepsilon)\bigr)^\dagger
  d_1(\theta_\varepsilon)^*;
\]
\item the exact alias blocks \(U_\varepsilon(k)\) of the uniform
four-plaquette lift;
\item the six-by-six Hermitian symbol
\[
 G(k)=\sum_{\varepsilon\in\{0,1\}^4}
       U_\varepsilon(k)^*P_f(\theta_\varepsilon)U_\varepsilon(k),
                                                                    \tag{16.27}
\]
with the Fourier normalization fixed by \(U^*U=I\);
\item restriction to the three-dimensional physical space
\(\operatorname{Ran}d_c(k)\); and
\item an outward-rounded lower enclosure for
\(\lambda_{\min}(G(k))\) on the selected carrier.
\end{enumerate}
\end{definition}

\begin{lemma}[The card is finite and gauge independent]
\label{lem:e8v16-card}
For \(k\ne0\), the nonzero spectrum of \(G(k)\) on
\(\operatorname{Ran}d_c(k)\) is independent of the chosen coarse
Coulomb basis and of the chosen fine edge representative used to
write \(P_f\).  It is a finite generalized eigenvalue problem of
dimension three.
\end{lemma}

\begin{proof}
The formula for \(P_f\) is the orthogonal projection onto
\(\operatorname{Ran}d_1\), which depends only on that subspace, not on
a gauge representative.  A change of basis of
\(\operatorname{Ran}d_c(k)\) transforms the compressed pair
\((U^*P_fU,U^*U)\) by simultaneous congruence, preserving its
generalized eigenvalues.  The lattice complex at nonzero momentum has
\(\dim\operatorname{Ran}d_c=3\), since the edge symbol has dimension
four and its exact kernel is the one-dimensional gauge gradient.
\end{proof}

\begin{assessment}[What V16 closes]
\label{ass:e8v16-closes}
The conditioning-space ambiguity is closed.  The correct V14
conditional mean, its fluctuation covariance, its FEC criterion, and
its lift defect are now exact Hodge-projection formulas.  A completely
analytic carrier bound,
\(\gamma_F\geq\mu_+/16\), follows from the exact metric-lift symbol.
No \(45\times45\) inverse is required for that bound.
\end{assessment}

\begin{assessment}[What remains]
\label{ass:e8v16-remains}
The analytic defect bound in (16.22) is large when the carrier
approaches zero momentum, and the polynomial transport norm in V14 is
an owner-weighted coefficient norm rather than the Hilbert norm used
here.  The remaining principal task has two parts:
\begin{enumerate}
\item evaluate or sharply bound the three-dimensional symbol \(G(k)\)
on the actual scale-separated carrier; and
\item convert \(M_{\rm curv}\) from Hilbert norm to the declared
plaquette-owner tensor norm before applying the V14 perturbation gate.
\end{enumerate}
Neither step is a vertical Hessian-gap problem.
\end{assessment}

\begin{decision}[V17 acquisition order]
\label{dec:e8v16-v17}
V17 will form the alias blocks \(U_\varepsilon(k)\) explicitly and
use the lattice Hodge projector to simplify (16.27).  It will first
test whether cubical symmetry makes \(G(k)\) scalar on the transverse
coarse carrier.  If not, it will retain the exact three-by-three
generalized symbol and derive an owner-norm conversion bound.  This
continues upstream at the correct curvature conditional law.
\end{decision}

\begin{firewall}[Claim boundary after eighth-series V16]
\label{fire:e8v16-claim}
V16 corrects a false identification in V15 and proves the proper
curvature conditional-mean and defect formulas.  It does not yet
evaluate the sharp carrier minimum of \(G\), certify the physical
owner-norm perturbation gate, prove the complete stable contraction,
construct the running marginal trajectory, close nonlinear
large-field estimates, restore rotations, remove the regulator,
reconstruct the continuum Osterwalder--Schrader theory, or prove the
Yang--Mills mass gap.
\end{firewall}

\paragraph{Parent-version record.}
V16 was formed by exact append-only extension of
\texttt{Renewal\_Yang\_Mills\_Mass\_Gap\_Research\_Notes\_V15.tex}.
The V15 parent had \(7\,220\) counted source lines,
\(254\,667\) bytes, and SHA--256
\[
\texttt{4C39916B3E9E875A21B3EC76256264903B7B15327B5DBEEB899CDD1E2637B03F}.
\]

% =============================================================================
% END APPENDED EIGHTH-SERIES V16 MATERIAL
% =============================================================================

% =============================================================================
% BEGIN APPENDED EIGHTH-SERIES V17 MATERIAL
% =============================================================================

\section{V17: exact bridge from the Wilson Schur reference to the
curvature angle}
\label{sec:e8v17}

\subsection{Why the edge Schur form still matters}

V16 corrected the identification of the edge-conditioned lift with
the curvature-conditioned lift.  The lifts are different.  Their
minimum energies, however, are related exactly after the coarse gauge
orbit is quotiented.  This note proves that relation and uses the
already established axial Wilson symbol to evaluate the true
curvature projection without constructing a \(96\times64\) Hodge
matrix.

Retain the notation of V16.  At a nonzero coarse momentum \(k\), put
\[
 p_\mu=e^{ik_\mu}-1,\qquad
 r(k)=\sum_\mu|p_\mu|^2.                            \tag{17.1}
\]
Let
\[
 \mathsf S^{\rm tree}(k)
\]
be the unscaled Wilson edge Schur form characterized by
\[
 \langle h,\mathsf S^{\rm tree}h\rangle
 =\min_{Ba=h}\|d_1a\|^2.                            \tag{17.2}
\]
Thus the unit-action form in the supplied V30 notation is
\(\mathsf S_M=\mathsf S^{\rm tree}/3\).

\begin{lemma}[Coarse curvature determines the transverse edge class]
\label{lem:e8v17-gauge-class}
For \(y\in\mathcal R_c(k)=\operatorname{Ran}d_c(k)\), define
\[
             L(k)y:=\frac4{r(k)}d_c(k)^*y.          \tag{17.3}
\]
Then
\[
 d_cL(k)y=4y,\qquad d_0(k)^*L(k)y=0.                \tag{17.4}
\]
Every \(h\) satisfying \(d_ch=4y\) differs from \(L(k)y\) by a coarse
pure gauge.
\end{lemma}

\begin{proof}
On the exact two-form range,
\[
 d_cd_c^*=r(k)I.
\]
This proves the first identity.  The complex relation
\(d_cd_0=0\) gives the second.  If \(d_c(h-Ly)=0\), exactness of the
nonzero-momentum lattice cochain complex gives
\(\ker d_c=\operatorname{Ran}d_0\), proving the final assertion.
\end{proof}

\begin{lemma}[The edge Schur form descends to coarse curvature]
\label{lem:e8v17-ward}
At nonzero momentum,
\[
 \mathsf S^{\rm tree}d_0=0,\qquad
 d_0^*\mathsf S^{\rm tree}=0.                       \tag{17.5}
\]
Consequently
\[
 \langle h,\mathsf S^{\rm tree}h\rangle
 =\langle L(k)y,\mathsf S^{\rm tree}L(k)y\rangle    \tag{17.6}
\]
whenever \(d_ch=4y\).
\end{lemma}

\begin{proof}
A coarse pure gauge has a fine pure-gauge completion, whose curl
energy is zero.  The variational definition (17.2) therefore gives
\(\mathsf S^{\rm tree}[d_0\xi]=0\).  Positivity of the Schur form
implies \(\mathsf S^{\rm tree}d_0=0\), and self-adjointness gives the
adjoint identity.  Apply Lemma~\ref{lem:e8v17-gauge-class}.
\end{proof}

\begin{theorem}[Exact edge-to-curvature bridge]
\label{thm:e8v17-bridge}
On \(\mathcal R_c(k)\),
\[
\boxed{\quad
 4G(k)^{-1}
 =L(k)^*\mathsf S^{\rm tree}(k)L(k).
\quad}                                             \tag{17.7}
\]
Equivalently,
\[
\boxed{\quad
 G(k)^{-1}
 =\frac4{r(k)^2}\,
 d_c(k)\mathsf S^{\rm tree}(k)d_c(k)^*
 =\frac{12}{r(k)^2}\,
 d_c(k)\mathsf S_M(k)d_c(k)^*.
\quad}                                             \tag{17.8}
\]
\end{theorem}

\begin{proof}
The curvature-conditioned minimum is
\[
 \|J_{\rm curv}y\|^2
 =\min_{\substack{F\in\mathcal R_f\\QF=y}}\|F\|^2. \tag{17.9}
\]
Every \(F\in\mathcal R_f\) is \(d_1a\).  By the linear Stokes identity,
\(Qd_1a=y\) if and only if \(d_cBa=4y\).  Therefore the right side of
(17.9) is
\[
 \min_{\substack{h\\d_ch=4y}}\ \min_{Ba=h}\|d_1a\|^2
 =\min_{\substack{h\\d_ch=4y}}
       \langle h,\mathsf S^{\rm tree}h\rangle.
\]
Lemma~\ref{lem:e8v17-ward} makes the last quantity
\(\langle Ly,\mathsf S^{\rm tree}Ly\rangle\).
Theorem~\ref{thm:e8v16-defect} gives
\(\|J_{\rm curv}y\|^2=4\langle y,G^{-1}y\rangle\),
proving (17.7).  Substitute \(L=4d_c^*/r\) and
\(\mathsf S^{\rm tree}=3\mathsf S_M\) to obtain (17.8).
\end{proof}

\begin{corollary}[Reference comparisons and curvature angles]
\label{cor:e8v17-comparison}
Let \(c_-(k)\) and \(c_+(k)\) satisfy, on the transverse coarse edge
space,
\[
 c_-(k)d_c^*d_c
 \preceq\mathsf S^{\rm tree}(k)
 \preceq c_+(k)d_c^*d_c.                           \tag{17.10}
\]
Then, on \(\mathcal R_c(k)\),
\[
\boxed{\qquad
 \frac1{4c_+(k)}I\preceq G(k)
 \preceq\frac1{4c_-(k)}I.
\qquad}                                            \tag{17.11}
\]
\end{corollary}

\begin{proof}
On the Coulomb edge space \(d_c^*d_c=rI\).  Inserting (17.10) into
(17.7), and using
\[
 L^*L=\frac{16}{r}I
\]
on \(\mathcal R_c\), gives
\[
 4c_-(k)I\preceq G(k)^{-1}\preceq4c_+(k)I.
\]
Invert the positive inequalities.
\end{proof}

\begin{remark}[The V15 data are an energy bridge, not a lift identity]
\label{rem:e8v17-v15}
Theorem~\ref{thm:e8v17-bridge} is the precise surviving connection
between V15 and V14.  The edge minimizer and the curvature conditional
mean generally live in different spaces and are not equal, but their
minimum energies agree after minimization over the coarse gauge
class.
\end{remark}

\subsection{Exact axial evaluation}

Take \(k=(q,0,0,0)\), \(0<|q|\leq\pi\), and put
\[
                         t=\sin^2(q/2).             \tag{17.12}
\]
The supplied exact V30 Wilson calculation proves that
\(\mathsf S_M\) is scalar on the three transverse edge directions,
with
\[
 s_0(q)=\frac{t(t+8)(t+24)}
                  {9(t^2+18t+16)},\qquad r(k)=4t.  \tag{17.13}
\]

\begin{theorem}[Exact axial curvature-projection eigenvalue]
\label{thm:e8v17-axial}
On the three-dimensional physical coarse curvature space at axial
momentum,
\[
                         G(k)=g_{\rm ax}(t)I,        \tag{17.14}
\]
where
\[
\boxed{\qquad
 g_{\rm ax}(t)
 =\frac{3(t^2+18t+16)}{(t+8)(t+24)}.
\qquad}                                             \tag{17.15}
\]
The function is strictly increasing on \(0\leq t\leq1\), and
\[
\boxed{\qquad
 \lim_{q\to0}g_{\rm ax}(t)=\frac14,\qquad
 g_{\rm ax}(1)=\frac7{15}.
\qquad}                                             \tag{17.16}
\]
\end{theorem}

\begin{proof}
For an axial momentum the transverse compression of
\(\mathsf S^{\rm tree}=3\mathsf S_M\) is \(3s_0(q)I\).
Equation (17.8) therefore gives
\[
 G^{-1}
 =\frac{12s_0(q)}{r(k)}I,
\]
which, after inserting (17.13), is (17.15).  Its derivative is
\[
 g_{\rm ax}'(t)
 =\frac{3(14t^2+352t+2944)}
        {(t^2+32t+192)^2}>0.
\]
Evaluation at \(t=0,1\) proves (17.16).
\end{proof}

\begin{corollary}[The uniform-lift defect is not perturbatively small
in the axial infrared]
\label{cor:e8v17-axial-defect}
Along a nonzero axial sequence \(q\to0\),
\[
\boxed{
\begin{aligned}
\|J_{\rm curv}\|&\longrightarrow4,\\
\|M_{\rm curv}\|&\longrightarrow2\sqrt3.
\end{aligned}}                                      \tag{17.17}
\]
At \(q=\pi\),
\[
 \|M_{\rm curv}\|=2\sqrt{\frac87}.                  \tag{17.18}
\]
\end{corollary}

\begin{proof}
Apply (16.14) to the scalar eigenvalue (17.15), then use (17.16).
\end{proof}

\begin{firewall}[A large coordinate defect is not an instability]
\label{fire:e8v17-large-defect}
Equation (17.17) does not indicate a negative Hessian or a physical
infrared divergence.  The exact conditional lift remains bounded and
the Schur energy has the correct \(r(k)\) scaling.  It says only that
the ambient four-face-uniform lift is a poor perturbative centre for
the realizable-curvature conditional law.
\end{firewall}

\subsection{Two exact corner calibrations}

\begin{proposition}[The three-\(\pi\) corner]
\label{prop:e8v17-three-pi}
At \(k_*=(\pi,\pi,\pi,0)\), the three eigenvalues of \(G(k_*)\) on
\(\mathcal R_c(k_*)\) are
\[
                         \left\{\frac45,\frac45,1\right\}.          \tag{17.19}
\]
\end{proposition}

\begin{proof}
The supplied exact Schur matrix is
\[
\mathsf S^{\rm tree}(k_*)=
\begin{pmatrix}
5/2&-5/4&-5/4&0\\
-5/4&5/2&-5/4&0\\
-5/4&-5/4&5/2&0\\
0&0&0&3
\end{pmatrix}.                                      \tag{17.20}
\]
Here \(r(k_*)=12\).  On the transverse plane in the first three
coordinates the Schur eigenvalue is \(15/4\), while on \(e_4\) it is
\(3\).  If \(h\) is transverse and \(y=d_ch/4\), then
\[
 \|y\|^2=\frac{r(k_*)}{16}\|h\|^2=\frac34\|h\|^2.
\]
The identity
\[
 4\langle y,G^{-1}y\rangle
 =\langle h,\mathsf S^{\rm tree}h\rangle
\]
therefore gives \(G=3/(15/4)=4/5\) on the two-plane and
\(G=3/3=1\) on \(e_4\).
\end{proof}

\begin{proposition}[The all-\(\pi\) corner]
\label{prop:e8v17-all-pi}
At \(k=(\pi,\pi,\pi,\pi)\),
\[
                              G(k)=I
\]
on the three-dimensional physical curvature space.  Thus FEC holds
exactly at this corner.
\end{proposition}

\begin{proof}
Here \(r=16\).  For every Coulomb edge \(h\), Proposition
\ref{prop:e8v16-metric-symbol} gives
\[
 \|d_1R_0h\|^2
 =3\|h\|^2+\frac14\sum_\mu4|h_\mu|^2
 =4\|h\|^2.
\]
With \(y=d_ch/4\),
\[
 4\|y\|^2=\frac r4\|h\|^2=4\|h\|^2.
\]
The four-plaquette Cauchy--Schwarz inequality says that every
realization of \(y\) has energy at least \(4\|y\|^2\).  The metric
lift attains equality, so equality holds separately in every owned
four-plaquette tile.  Its curvature is therefore uniform on each
tile, which is FEC.  Corollary~\ref{cor:e8v16-fec} gives \(G=I\).
\end{proof}

\subsection{Consequences for the contraction programme}

\begin{theorem}[The V14 perturbative centre should be replaced]
\label{thm:e8v17-recenter}
No scale-uniform argument can make
\(\|M_{\rm curv}\|\) a small parameter on carriers containing
arbitrarily small nonzero axial momenta.  A contraction proof on such
carriers must use one of the following:
\begin{enumerate}
\item take \(J_{\rm curv}\), rather than \(J_0\), as the principal
conditional lift and compute its coefficient transport directly; or
\item split off the low carrier and use a different infrared
coordinate there.
\end{enumerate}
\end{theorem}

\begin{proof}
Corollary~\ref{cor:e8v17-axial-defect} gives the nonzero limit
\(2\sqrt3\).  Hence an assumption
\(\|M_{\rm curv}\|\leq\varepsilon\) with arbitrarily small
scale-independent \(\varepsilon\) is false on those carriers.  The two
listed alternatives are exhaustive at the level of the chosen
coordinates: either the exact mean is included in the principal map,
or the modes on which the comparison is not small are removed from
that perturbative chart.
\end{proof}

\begin{assessment}[Numerical reconnaissance, not a proof]
\label{ass:e8v17-recon}
As an error-detection calculation, the \(96\times64\) parity-cell
curl and the Hodge projection \(P_f\) were assembled numerically.
They reproduced the exact axial values, the spectrum
\(\{4/5,4/5,1\}\) at \(k_*\), and \(G=I\) at the all-\(\pi\) corner.
Generic momenta produced three unequal eigenvalues, so cubical
symmetry does not make \(G(k)\) scalar pointwise.  No decimal output
from this reconnaissance is used as a theorem.
\end{assessment}

\begin{decision}[V18 acquisition order]
\label{dec:e8v17-v18}
V18 will abandon the attempt to close the principal contraction by a
small perturbation of \(J_0\).  It will rewrite the polynomial
pushforward with the exact conditional lift
\[
                       J_{\rm curv}=2P_fUG^{-1}
\]
as the principal map.  The first target is an exact
plaquette-owner coefficient-mass bound derived from the stochastic
projection identity, followed by the quartic Wick-lowering row using
the exact covariance (16.15).  The V14 perturbation theorem remains
available only on carriers where a genuinely small defect is proved.
\end{decision}

\begin{firewall}[Claim boundary after eighth-series V17]
\label{fire:e8v17-claim}
V17 proves the exact edge-to-curvature energy bridge, evaluates the
curvature angle on the full axial orbit and at two high-symmetry
corners, and proves that the uniform lift is not a small infrared
perturbative centre.  It does not yet bound the exact conditional
lift in the plaquette-owner coefficient norm, close the principal or
nonlinear stable contraction, construct the running marginal
trajectory, restore rotations, establish regulator-independent
reflection positivity, reconstruct the continuum theory, or prove
the Yang--Mills mass gap.
\end{firewall}

\paragraph{Parent-version record.}
V17 was formed by exact append-only extension of
\texttt{Renewal\_Yang\_Mills\_Mass\_Gap\_Research\_Notes\_V16.tex}.
The V16 parent had \(7\,763\) counted source lines,
\(273\,649\) bytes, and SHA--256
\[
\texttt{22B44623B396B9A393B02034BB5F905B856C8F1E1877495FBD6613078732E9DC}.
\]

% =============================================================================
% END APPENDED EIGHTH-SERIES V17 MATERIAL
% =============================================================================

% =============================================================================
% BEGIN APPENDED EIGHTH-SERIES V18 MATERIAL
% =============================================================================

\section{V18: exact conditional Wick contraction and removal of the FEC
principal hypothesis}
\label{sec:e8v18}

\subsection{The corrected principal coordinate}

V17 shows that the four-face-uniform lift is not a small perturbative
centre on infrared carriers.  The correct centre is the exact
conditional Gaussian law.  The present note proves that in covariance-
matched Wick coordinates its action on every homogeneous polynomial
chaos is contractive, with no degree-lowering row.

Work first in finite dimension.  All gauge directions have been
quotiented by the rooted Coulomb condition, and all toron variables are
retained outside the Gaussian inverse.  Let \(E\) and \(Y\) be the
resulting fine and coarse real Hilbert spaces.  Let
\[
 H:E\longrightarrow E,\qquad H\succ0,                              \tag{18.1}
\]
be the flat fine precision and let \(B:E\to Y\) be onto.  Define
\[
\begin{split}
 C&:=BH^{-1}B^*,&
 S&:=C^{-1},\\
 J&:=H^{-1}B^*S,&
 C_{\rm f}&:=H^{-1}-JS^{-1}J^* .                                  \tag{18.2}
\end{split}
\]
Thus \(C\) is the coarse covariance, \(S\) is the effective coarse
precision, \(J\) is the energy-minimizing lift, and \(C_{\rm f}\) is
the conditional fibre covariance.

\begin{lemma}[Energy isometry and covariance splitting]
\label{lem:e8v18-isometry}
The operators in (18.2) satisfy
\[
\boxed{
 BJ=I,\qquad J^*HJ=S,\qquad
 H^{-1}=JS^{-1}J^*+C_{\rm f},\qquad C_{\rm f}\succeq0.}             \tag{18.3}
\]
Consequently
\[
                   R:=H^{1/2}JS^{-1/2}:Y\to E                     \tag{18.4}
\]
is an isometry.
\end{lemma}

\begin{proof}
Since \(S=(BH^{-1}B^*)^{-1}\),
\[
 BJ=BH^{-1}B^*S=I.
\]
Also
\[
 J^*HJ
 =SBH^{-1}B^*S=S.
\]
For \(X\) with covariance \(H^{-1}\), the standard Schur formula for
the covariance conditioned on \(BX\) is precisely the last expression
in (18.2), so it is positive semidefinite.  Alternatively,
\[
 H^{1/2}C_{\rm f}H^{1/2}
 =I-RR^*
\]
and \(R^*R=S^{-1/2}J^*HJS^{-1/2}=I\), proving both positivity and the
isometry claim.
\end{proof}

\subsection{Covariance-matched Wick tensors}

Let \(X\) be the centred Gaussian on \(E\) with covariance \(H^{-1}\),
and put
\[
                              Z=BX.                                \tag{18.5}
\]
Then \(Z\) is centred Gaussian with covariance \(S^{-1}\).
For \(u\in E\) and \(v\in Y\), define the Wick exponentials
\[
\begin{split}
\mathcal W_H(u;X)
 &:=\exp\left(\langle u,X\rangle
       -\frac12\langle u,H^{-1}u\rangle\right),\\
\mathcal W_S(v;Z)
 &:=\exp\left(\langle v,Z\rangle
       -\frac12\langle v,S^{-1}v\rangle\right).                    \tag{18.6}
\end{split}
\]
Their Taylor coefficients are the covariance-matched symmetric Wick
tensors.

\begin{theorem}[Exact conditional Wick-exponential identity]
\label{thm:e8v18-wick-exponential}
For every \(u\in E\),
\[
\boxed{\qquad
\mathbb E\!\left[\mathcal W_H(u;X)\mid Z\right]
       =\mathcal W_S(J^*u;Z).
\qquad}                                                           \tag{18.7}
\]
\end{theorem}

\begin{proof}
The conditional Gaussian law has mean \(JZ\) and covariance
\(C_{\rm f}\).  Therefore
\[
\begin{split}
\mathbb E\!\left[e^{\langle u,X\rangle}\mid Z\right]
 &=\exp\left(\langle u,JZ\rangle
       +\frac12\langle u,C_{\rm f}u\rangle\right).
\end{split}
\]
Subtract the fine Wick normalization.  Using (18.3),
\[
\begin{split}
\langle u,C_{\rm f}u\rangle-\langle u,H^{-1}u\rangle
 &=-\langle u,JS^{-1}J^*u\rangle\\
 &=-\langle J^*u,S^{-1}J^*u\rangle.
\end{split}
\]
The resulting exponential is exactly the right side of (18.7).
\end{proof}

\begin{corollary}[No Wick degree lowering]
\label{cor:e8v18-no-lowering}
For every symmetric \(n\)-tensor \(T_n\),
\[
\boxed{\qquad
\mathbb E\!\left[
  \left\langle T_n,{:X^{\otimes n}:}_{H^{-1}}\right\rangle
  \,\middle|\,Z\right]
=
\left\langle (J^*)^{\otimes n}T_n,
       {:Z^{\otimes n}:}_{S^{-1}}\right\rangle .
\qquad}                                                           \tag{18.8}
\]
In particular the conditional image of the fourth Wick chaos has no
quadratic component.
\end{corollary}

\begin{proof}
Polarize the degree-\(n\) Taylor coefficient of (18.7).  Equality of
entire power series makes the equality degree preserving.
\end{proof}

\begin{correction}[Status of the earlier quartic-to-quadratic row]
\label{corr:e8v18-wick-row}
The quartic-to-quadratic Wick row in V12--V14 is a correct formula when
fine and coarse polynomials are expressed with a common or otherwise
mismatched normal-ordering covariance.  It is absent from the
principal flat Gaussian map in the matched coordinates (18.6).
It may reappear only as an explicit covariance-rechart term when the
chosen RG chart differs from the exact pushed-forward covariance.
\end{correction}

\subsection{One-particle and Fock contraction}

Equip fine coefficient covectors with
\[
             \|u\|_{H^{-1}}^2=\langle u,H^{-1}u\rangle
\]
and coarse coefficient covectors with
\[
             \|v\|_{S^{-1}}^2=\langle v,S^{-1}v\rangle.            \tag{18.9}
\]

\begin{lemma}[The coefficient map is a contraction]
\label{lem:e8v18-one-particle}
For all \(u\in E\),
\[
\boxed{\qquad
        \|J^*u\|_{S^{-1}}\leq\|u\|_{H^{-1}}.
\qquad}                                                           \tag{18.10}
\]
Equality holds precisely when \(u\) annihilates the conditional
fluctuation range of \(C_{\rm f}\).
\end{lemma}

\begin{proof}
By (18.3),
\[
\begin{split}
\|u\|_{H^{-1}}^2-\|J^*u\|_{S^{-1}}^2
 &=\langle u,(H^{-1}-JS^{-1}J^*)u\rangle\\
 &=\langle u,C_{\rm f}u\rangle\geq0.
\end{split}
\]
Equality is equivalent to \(C_{\rm f}^{1/2}u=0\).
\end{proof}

\begin{definition}[Matched chaos norm]
\label{def:e8v18-chaos-norm}
For a symmetric coefficient tensor \(T_n\in E^{\odot n}\), let
\[
 \|T_n\|_{n,H^{-1}}
 :=\left\|(H^{-1/2})^{\otimes n}T_n\right\|_{\mathrm{HS}},          \tag{18.11}
\]
and define \(\|\cdot\|_{n,S^{-1}}\) analogously on \(Y^{\odot n}\).
The normalized Gaussian-chaos norm is
\[
 \left\|
 \left\langle T_n,{:X^{\otimes n}:}_{H^{-1}}\right\rangle
 \right\|_{\mathfrak C_n}^2
 :=n!\,\|T_n\|_{n,H^{-1}}^2.                                      \tag{18.12}
\]
\end{definition}

\begin{theorem}[Symmetric second-quantization contraction]
\label{thm:e8v18-second-quantization}
On the \(n\)-th matched Wick chaos,
\[
 \left\|(J^*)^{\otimes n}T_n\right\|_{n,S^{-1}}
 \leq\|T_n\|_{n,H^{-1}}.                                          \tag{18.13}
\]
Consequently conditional expectation is an orthogonal contraction on
the direct Gaussian Fock sum
\[
                    \bigoplus_{n\geq0}\mathfrak C_n.               \tag{18.14}
\]
\end{theorem}

\begin{proof}
In whitened coefficient coordinates the one-particle map is
\[
 S^{-1/2}J^*H^{1/2}=R^*,
\]
where \(R\) is the isometry in Lemma~\ref{lem:e8v18-isometry}.
Thus \(\|R^*\|\leq1\), and
\(\|(R^*)^{\otimes n}\|\leq1\) on the full tensor product and hence on
its symmetric subspace.  This is (18.13).  Distinct Gaussian Wick
chaoses are orthogonal, so summing their squared norms proves (18.14).
\end{proof}

\begin{corollary}[Volume-uniformity]
\label{cor:e8v18-volume}
The constants in
Theorem~\ref{thm:e8v18-second-quantization} equal one on every finite
torus and do not depend on its volume, on the condition number of
\(H\), or on the difference between \(J\) and the uniform face lift.
\end{corollary}

\begin{proof}
The proof uses only the exact isometry \(R^*R=I\), whose norm is one
in every finite dimension.
\end{proof}

\subsection{Principal dyadic contraction}

Let the principal polynomial types remain
\[
                         (Q_2,C_3,C_4),                            \tag{18.15}
\]
with the canonical dyadic coefficient factors already fixed in the
programme:
\[
                 s_{Q_2}=2^{-2},\qquad
                 s_{C_3}=2^{-2},\qquad
                 s_{C_4}=2^{-4}.                                  \tag{18.16}
\]
At this stage these symbols denote homogeneous matched Wick tensors,
not ordinary monomials.

\begin{theorem}[Exact global Gaussian principal margin]
\label{thm:e8v18-principal}
In the direct matched-chaos Hilbert norm, the flat Gaussian derivative
of one dyadic RG step is block diagonal by chaos degree and obeys
\[
\boxed{\qquad
 \|\mathcal L_{\rm G}\|
 \leq\max\{2^{-2},2^{-2},2^{-4}\}
 =\frac14.
\qquad}                                                           \tag{18.17}
\]
In particular
\[
                       \rho(\mathcal L_{\rm G})\leq\frac14.         \tag{18.18}
\]
This conclusion requires neither FEC nor a small bound on
\(M_{\rm curv}\).
\end{theorem}

\begin{proof}
Corollary~\ref{cor:e8v18-no-lowering} makes the conditional map
degree preserving.  Theorem~\ref{thm:e8v18-second-quantization}
bounds its restriction to each homogeneous chaos by one.  Multiplying
the three declared output types by their canonical factors (18.16)
gives the diagonal block bounds \(1/4,1/4,1/16\).  The norm of their
orthogonal direct sum is the maximum.  Spectral radius is bounded by
operator norm.
\end{proof}

\begin{corollary}[The FEC hypothesis is removed from the global
principal Gaussian gate]
\label{cor:e8v18-fec-removed}
The FEC-dependent conclusion of
Theorem~\ref{thm:e8v13-principal-margin} holds unconditionally for the
flat Gaussian derivative after passage to covariance-matched Wick
coordinates.  The V14 small-defect inequality is not needed for this
global Hilbert-chaos statement.
\end{corollary}

\begin{proof}
This is Theorem~\ref{thm:e8v18-principal}.  The exact energy lift and
conditional covariance are absorbed into the source and target
one-particle metrics rather than estimated as a perturbation of
\(J_0\).
\end{proof}

\subsection{Gauge and curvature formulations}

\begin{proposition}[Gauge quotient compatibility]
\label{prop:e8v18-gauge}
The construction descends canonically from rooted Coulomb edge
coordinates to gauge-invariant curvature polynomial functionals.
If a fine coefficient covector has the form \(u=d_1^*f\), then
\[
 \|u\|_{H^\dagger}^2
 =\langle f,P_ff\rangle
 \leq\|f\|^2,                                      \tag{18.19}
\]
up to the common Wilson scalar.  Pure-gauge covectors have zero norm
and are absent from every physical chaos.
\end{proposition}

\begin{proof}
For the unit flat Wilson Hessian, ignoring its common positive scalar,
\[
 H=d_1^*d_1.
\]
Therefore
\[
\begin{split}
\langle d_1^*f,H^\dagger d_1^*f\rangle
 &=\langle f,d_1(d_1^*d_1)^\dagger d_1^*f\rangle\\
 &=\langle f,P_ff\rangle\leq\|f\|^2.
\end{split}
\]
The formula depends only on the realizable projection \(P_f\), hence
is gauge independent.  A pure-gauge direction lies in \(\ker d_1\)
and is removed by the quotient.
\end{proof}

\begin{remark}[Relation to V16--V17]
\label{rem:e8v18-relation}
In curvature variables the one-particle isometry is the normalized
map \(P_fUG^{-1/2}\) already implicit in the proof of
Corollary~\ref{cor:e8v16-covariance}.  Thus the potentially large
coordinate norm of \(J_{\rm curv}=2P_fUG^{-1}\) is exactly compensated
by the physical coarse covariance \(G/4\).  The large value in
(17.17) was a comparison between mismatched Euclidean coordinates,
not a loss in the covariance-matched chaos norm.
\end{remark}

\subsection{The remaining localization problem}

\begin{firewall}[A global \(L^2\) contraction is not yet a polymer
contraction]
\label{fire:e8v18-locality}
The norm in (18.17) is a global Gaussian Hilbert norm.  It does not by
itself control:
\begin{enumerate}
\item the spatial support or exponential tail of \(J^*\);
\item the sum over plaquette owners in a polymer norm;
\item covariance rechart derivatives when the running precision
changes;
\item large-field marks or nonlinear interactions.
\end{enumerate}
Therefore (18.17) does not yet prove STC or the full RG invariant
ball.
\end{firewall}

\begin{definition}[Localized second-quantization gate]
\label{def:e8v18-lsg}
LSG consists of the following scale-uniform statements.
\begin{enumerate}
\item A block kernel representation of the whitened one-particle
contraction \(R^*\).
\item An exponentially weighted row and column bound
\[
 \sup_x\sum_y e^{\mu d(x,y)}\|R^*_{xy}\|\leq1+\ell_\mu,\qquad
 \sup_y\sum_x e^{\mu d(x,y)}\|R^*_{xy}\|\leq1+\ell_\mu.             \tag{18.20}
\]
\item A declared owner map whose multiplicity is paid once.
\item For each retained degree \(n\) with canonical factor \(s_n\),
\[
                         s_n(1+\ell_\mu)^n<1.                       \tag{18.21}
\]
\end{enumerate}
\end{definition}

\begin{lemma}[LSG implies a localized principal contraction]
\label{lem:e8v18-lsg}
If LSG holds, the degree-\(n\) principal conditional map has weighted
owner norm at most
\[
                         s_n(1+\ell_\mu)^n,                         \tag{18.22}
\]
apart from the single declared owner-conversion factor.  In
particular, (18.21) gives strict contraction on every retained
principal degree.
\end{lemma}

\begin{proof}
The tensor kernel is the symmetric restriction of
\((R^*)^{\otimes n}\).  Expanding its weighted row or column sum
factorizes into at most the \(n\)-th power of the one-particle bound.
Symmetrization is an orthogonal projection and cannot increase the
Hilbert tensor norm.  Multiply by the canonical factor \(s_n\).
The owner conversion is kept separate so that overlapping hull
multiplicity is not paid once per tensor leg.
\end{proof}

\begin{assessment}[What V18 closes]
\label{ass:e8v18-closes}
The flat principal Gaussian map is now rigorously contractive with
factor at most \(1/4\) in its natural, volume-uniform matched-chaos
norm.  The conclusion is exact for the physical energy lift, removes
FEC, and eliminates principal Wick degree lowering.  The earlier
large uniform-lift defect is absorbed by covariance matching rather
than estimated.
\end{assessment}

\begin{assessment}[Current bottleneck]
\label{ass:e8v18-bottleneck}
The active bottleneck is no longer a finite Hessian angle.  It is the
localization of the normalized one-particle contraction \(R^*\) in a
polymer-compatible owner norm, including the covariance rechart along
the running marginal trajectory.  The exact vertical gap and the
Chebyshev--Demko bounds in the supplied V30 note provide locality of
the unwhitened energy lift, but locality of the normalized square-root
metrics must not be assumed without proof.
\end{assessment}

\begin{decision}[V19 acquisition order]
\label{dec:e8v18-v19}
V19 will avoid nonlocal square roots by formulating a localized
energy-frame norm directly in terms of quadratic forms.  It will use
curvature covectors \(d_1^*f\), the exact vertical inverse decay, and
the unique plaquette-owner partition to seek (18.20) or an equivalent
quadratic-form localization estimate.  If square-root localization
becomes circular at zero momentum, the note will split physical
curvature and retained toron/infrared coordinates before proceeding.
\end{decision}

\begin{firewall}[Claim boundary after eighth-series V18]
\label{fire:e8v18-claim}
V18 proves an exact volume-uniform global Gaussian principal
contraction.  It does not yet prove the localized polymer contraction,
control the covariance rechart or nonlinear marked remainder, close
the running marginal trajectory, restore Euclidean symmetry, remove
the regulator with reflection positivity, reconstruct the continuum
Osterwalder--Schrader theory, or prove the Yang--Mills mass gap.
\end{firewall}

\paragraph{Parent-version record.}
V18 was formed by exact append-only extension of
\texttt{Renewal\_Yang\_Mills\_Mass\_Gap\_Research\_Notes\_V17.tex}.
The V17 parent had \(8\,158\) counted source lines,
\(286\,161\) bytes, and SHA--256
\[
\texttt{3B1661A7F3430669581F097385D2AFC1E76219CD24F4BC4A739421FABE0A1C6A}.
\]

% =============================================================================
% END APPENDED EIGHTH-SERIES V18 MATERIAL
% =============================================================================

% =============================================================================
% BEGIN APPENDED EIGHTH-SERIES V19 MATERIAL
% =============================================================================

\section{V19: vertical-gap localization and a two-norm principal
contraction}
\label{sec:e8v19}

\subsection{Locality of the exact energy lift}

V18 deliberately stopped before asserting locality of the normalized
square-root map \(R^*\).  V19 avoids that map.  It uses the exact
coordinate energy lift, the already certified vertical gap, and a
Lasota--Yorke two-norm argument.

Use coarse-cell \(\ell^1\) distance.  For a block kernel \(T\), define
\[
 |T|_\mu:=
 \max\left\{
 \sup_x\sum_y e^{\mu d(x,y)}\|T_{xy}\|_{\rm op},
 \sup_y\sum_x e^{\mu d(x,y)}\|T_{xy}\|_{\rm op}
 \right\}.                                                         \tag{19.1}
\]
This norm is submultiplicative.  Let \(\lambda_*\) be the exact
algebraic vertical eigenvalue imported in V15 and put
\[
 \rho_*=\frac{4-\sqrt{\lambda_*}}{4+\sqrt{\lambda_*}}.             \tag{19.2}
\]
The supplied V30 theorem proves
\[
\|(A^{-1})_{xy}\|_{\rm op}
\leq\frac6{\lambda_*}
       \rho_*^{\lceil d(x,y)/2\rceil}                              \tag{19.3}
\]
for the unit Wilson vertical precision \(A=\Sigma_{rr}/3\).

\begin{definition}[Explicit inverse-locality stock]
\label{def:e8v19-gmu}
For
\[
                         0<\mu<-\frac12\log\rho_*                  \tag{19.4}
\]
define
\[
 g_\mu:=
\frac6{\lambda_*}
\left[
1+16\sum_{n\geq1}(n+1)^3e^{\mu n}
        \rho_*^{\lceil n/2\rceil}
\right].                                                          \tag{19.5}
\]
\end{definition}

\begin{lemma}[Finite weighted vertical inverse]
\label{lem:e8v19-inverse}
For every \(\mu\) satisfying (19.4),
\[
                         |A^{-1}|_\mu\leq g_\mu<\infty.             \tag{19.6}
\]
\end{lemma}

\begin{proof}
The number of cells at \(\ell^1\)-distance \(n\) in
\(\mathbb Z^4\) is at most \(16(n+1)^3\).  Sum (19.3) over the second
cell for each fixed first cell, and then interchange the roles of the
cells.  The series converges because
\[
 e^{\mu n}\rho_*^{\lceil n/2\rceil}
 \leq e^\mu\left(e^\mu\sqrt{\rho_*}\right)^{n-1}
\]
and \(e^\mu\sqrt{\rho_*}<1\).
\end{proof}

Let \(R_0=B^*/2\) be the metric section.  In the vertical/horizontal
block coordinates of the actual edge conditioning problem,
\[
                         J=R_0-A^{-1}K                              \tag{19.8}
\]
is the exact energy lift.  The mixed Wilson block \(K\) has propagation
at most two.  The finite plaquette-incidence estimate in the supplied
V30 note gives, at unit Wilson coefficient,
\[
                         |K|_0\leq8.                               \tag{19.9}
\]

\begin{proposition}[Explicit lift-locality stock]
\label{prop:e8v19-lift}
For every \(\mu\) in (19.4),
\[
\boxed{\qquad
 |J|_\mu\leq j_\mu:=1+8e^{2\mu}g_\mu.
\qquad}                                                           \tag{19.10}
\]
If \(J_{\leq R}\) is obtained by setting \(J_{xy}=0\) for
\(d(x,y)>R\), then
\[
\boxed{\qquad
 |J-J_{\leq R}|_0\leq\tau_R:=j_\mu e^{-\mu R}.
\qquad}                                                           \tag{19.11}
\]
Both bounds are uniform in the finite-torus side.
\end{proposition}

\begin{proof}
The metric section has block row and column sums at most one, so
\(|R_0|_\mu\leq1\).  Propagation two and (19.9) give
\(|K|_\mu\leq8e^{2\mu}\).  Submultiplicativity and
Lemma~\ref{lem:e8v19-inverse} yield (19.10).  On the omitted entries,
\(1\leq e^{-\mu R}e^{\mu d(x,y)}\); summing gives (19.11).
All source estimates use shortest periodic distance and are
volume-uniform.
\end{proof}

\begin{corollary}[Conditional covariance locality]
\label{cor:e8v19-covariance}
In completed-square coordinates \(a=Jh+w\), the conditional vertical
covariance is \(A^{-1}\), and hence obeys (19.6).  Covariance-matched
Wick ordering therefore introduces no additional long-range
degree-lowering kernel.
\end{corollary}

\begin{proof}
The first assertion is the exact block completion already proved in
V14 for the valid edge-conditioned problem.  Apply
Lemma~\ref{lem:e8v19-inverse}.  The last assertion is
Corollary~\ref{cor:e8v18-no-lowering}.
\end{proof}

\subsection{Tensor tails}

\begin{lemma}[Truncation of the \(n\)-leg conditional kernel]
\label{lem:e8v19-tensor-tail}
For \(n\geq1\),
\[
\left|
(J^*)^{\otimes n}-(J_{\leq R}^*)^{\otimes n}
\right|_0
\leq n\,j_0^{\,n-1}\tau_R.                       \tag{19.12}
\]
\end{lemma}

\begin{proof}
Use the exact telescoping identity
\[
J^{\otimes n}-J_{\leq R}^{\otimes n}
=\sum_{\ell=1}^n
J^{\otimes(\ell-1)}\otimes(J-J_{\leq R})
\otimes J_{\leq R}^{\otimes(n-\ell)}.
\]
The transpose has the same symmetric row/column norm.  Bound
\(|J_{\leq R}|_0\leq|J|_0\leq j_0\) and use (19.11).
\end{proof}

\begin{corollary}[Scaled remote-tail constants]
\label{cor:e8v19-aR}
For the retained degrees and canonical factors of V18, set
\[
 a_R:=
\max_{n\in\{2,3,4\}}
 s_n\,n\,j_0^{\,n-1}j_\mu e^{-\mu R},
\qquad
(s_2,s_3,s_4)=\left(\frac14,\frac14,\frac1{16}\right).             \tag{19.13}
\]
Then \(a_R\to0\) as \(R\to\infty\).  In particular there is an
explicit finite \(R\) for which \(a_R<1\).
\end{corollary}

\begin{proof}
Apply Lemma~\ref{lem:e8v19-tensor-tail} and multiply by the appropriate
canonical factor.  All quantities other than \(e^{-\mu R}\) are
finite and independent of \(R\).
\end{proof}

\subsection{Anchored bulk coefficient spaces}

The next statement concerns the translation-covariant bulk action.
Each translated interaction is represented once by its unique
plaquette owner, as proved in V13.  Source insertions and boundaries
are not included in this bulk space.

\begin{definition}[Weak and strong anchored norms]
\label{def:e8v19-two-norms}
For a finite sum \(T=(T_2,T_3,T_4)\) of translation-covariant,
owner-anchored matched Wick tensors, let
\[
 \|T\|_{\rm w}
\]
be the direct covariance-matched chaos Hilbert norm of V18 on one
fundamental owner.  Let
\[
 \|T\|_{\rm s}
\]
be the corresponding anchored absolute coefficient norm, summed over
all relative leg cells with the fixed exponential support weight used
in the stable polymer grammar.  The owner itself is not summed again.
\end{definition}

\begin{lemma}[Finite-head comparison]
\label{lem:e8v19-head}
For each finite \(R\) there is a finite, volume-independent constant
\(b_R\) such that the \(R\)-truncated principal map
\(\mathcal L_{\leq R}\) satisfies
\[
                         \|\mathcal L_{\leq R}T\|_{\rm s}
 \leq b_R\|T\|_{\rm w}.                            \tag{19.14}
\]
\end{lemma}

\begin{proof}
After translation is quotiented by the unique owner, a radius-\(R\)
head contains finitely many cell and internal link types.  Denote that
number by \(N_R\).  Cauchy--Schwarz bounds its absolute coefficient
sum by \(\sqrt{N_R}\) times its Euclidean coefficient norm.

The fine and coarse quadratic forms have volume-uniform upper spectral
edges: the fine curl edge is at most \(16\), and the coarse Schur form
is a compression/minimum of the same finite-incidence form.  Hence
Euclidean coefficient norms on the physical quotient are bounded by
the covariance-dual norms with a finite geometry factor.  Tensoring
this comparison through degrees at most four and multiplying by the
finite kernel entries of \(J_{\leq R}\) gives a finite \(b_R\).
All counts concern relative cells around one owner, so no volume
factor occurs.
\end{proof}

\begin{firewall}[Why translation covariance is stated]
\label{fire:e8v19-translation}
An \(\ell^1\) sum over every translate cannot be bounded by a global
\(\ell^2\) norm uniformly in volume.  Lemma~\ref{lem:e8v19-head}
uses the bulk density modulo translations and the unique owner.
Localized sources require their separate source norm and are not
silently covered by this lemma.
\end{firewall}

\begin{theorem}[Flat-principal Lasota--Yorke inequality]
\label{thm:e8v19-lasota}
Let \(\mathcal L_{\rm G}\) be the exact covariance-matched flat
Gaussian principal map of V18.  For every \(R\),
\[
\boxed{
\begin{aligned}
\|\mathcal L_{\rm G}T\|_{\rm w}
 &\leq\frac14\|T\|_{\rm w},\\
\|\mathcal L_{\rm G}T\|_{\rm s}
 &\leq a_R\|T\|_{\rm s}+b_R\|T\|_{\rm w}.
\end{aligned}}                                                     \tag{19.15}
\]
\end{theorem}

\begin{proof}
The weak inequality is Theorem~\ref{thm:e8v18-principal}.  Split
\[
\mathcal L_{\rm G}
=\mathcal L_{\leq R}
 +(\mathcal L_{\rm G}-\mathcal L_{\leq R}).
\]
Lemma~\ref{lem:e8v19-head} bounds the first term from weak to strong.
The tensor-tail estimate and unique owner bound the second from strong
to strong by \(a_R\).
\end{proof}

\begin{theorem}[Strict localized flat-principal contraction]
\label{thm:e8v19-contraction}
Choose \(R\) so that \(a_R<1\), and choose
\[
       0<\varepsilon<\frac{3}{4b_R}                \tag{19.16}
\]
when \(b_R>0\).  Define
\[
                  \|T\|_\varepsilon
 :=\|T\|_{\rm w}+\varepsilon\|T\|_{\rm s}.          \tag{19.17}
\]
Then
\[
\boxed{\qquad
\|\mathcal L_{\rm G}T\|_\varepsilon
\leq q_{\varepsilon,R}\|T\|_\varepsilon,\qquad
q_{\varepsilon,R}:=
\max\left\{\frac14+\varepsilon b_R,a_R\right\}<1.
\qquad}                                             \tag{19.18}
\]
The norm and contraction factor are uniform in the finite-torus side.
\end{theorem}

\begin{proof}
Multiply the second line of (19.15) by \(\varepsilon\) and add it to
the first:
\[
\|\mathcal L_{\rm G}T\|_\varepsilon
\leq
\left(\frac14+\varepsilon b_R\right)\|T\|_{\rm w}
+\varepsilon a_R\|T\|_{\rm s}.
\]
Both coefficients are strictly below one by (19.16) and \(a_R<1\).
Taking their maximum proves (19.18).
\end{proof}

\begin{corollary}[Replacement for the square-root LSG]
\label{cor:e8v19-lsg}
For the translation-covariant flat bulk principal map, the localized
second-quantization gate of V18 can be replaced by the explicit
vertical-gap data
\[
  (\lambda_*,\rho_*,g_\mu,j_\mu,R,b_R,\varepsilon).
\]
No localization theorem for \(H^{1/2}\), \(S^{-1/2}\), or \(R^*\) is
required.
\end{corollary}

\begin{proof}
The weak contraction uses the square roots only as an abstract
finite-dimensional proof of a covariance identity.  The strong tail
uses the coordinate kernel \(J=R_0-A^{-1}K\), whose locality is
certified by Proposition~\ref{prop:e8v19-lift}.  The two estimates are
combined only at the norm level in
Theorem~\ref{thm:e8v19-contraction}.
\end{proof}

\begin{assessment}[What V19 closes]
\label{ass:e8v19-closes}
The flat Gaussian principal contraction is now both strict and
localized on the translation-covariant owner-anchored bulk
coefficient space.  The proof is volume-uniform, uses the exact
physical conditional lift, pays the plaquette owner once, and avoids
nonlocal square-root kernels.
\end{assessment}

\begin{assessment}[What remains]
\label{ass:e8v19-remains}
The theorem does not yet control the variation of \(H,J,A^{-1}\), and
the matched Wick covariance along a nonzero running interaction.
Nor does it cover localized source insertions, large-field marked
polymers, or the marginal coupling equation.  The immediate
bottleneck is therefore a nonlinear two-norm derivative estimate
showing that the small-field covariance rechart and interaction
remainder perturb (19.15) by less than
\(1-q_{\varepsilon,R}\).
\end{assessment}

\begin{decision}[V20 acquisition order]
\label{dec:e8v19-v20}
V20 is a compulsory SM/NS companion audit.  After that audit it will
differentiate the completed-square identities with respect to a
small local Hessian perturbation \(\Delta H\), retaining the same
weak/strong two-norm split.  The target is an explicit resolvent
bound for \(\Delta A^{-1}\), \(\Delta J\), and the covariance-matched
Wick rechart, followed by a perturbative contraction gate around
(19.18).
\end{decision}

\begin{firewall}[Claim boundary after eighth-series V19]
\label{fire:e8v19-claim}
V19 proves the localized flat-bulk principal contraction.  It does
not prove the nonlinear running-action contraction, source-sector
locality, large-field control, marginal trajectory, Euclidean
restoration, regulator removal with reflection positivity,
continuum reconstruction, or the Yang--Mills mass gap.
\end{firewall}

\paragraph{Parent-version record.}
V19 was formed by exact append-only extension of
\texttt{Renewal\_Yang\_Mills\_Mass\_Gap\_Research\_Notes\_V18.tex}.
The V18 parent had \(8\,628\) counted source lines,
\(302\,194\) bytes, and SHA--256
\[
\texttt{34EB88AB7B46DC8491E653DDC8D96EC94B749AB3E15B8C47E102EA5D6412180A}.
\]

% =============================================================================
% END APPENDED EIGHTH-SERIES V19 MATERIAL
% =============================================================================

% =============================================================================
% BEGIN APPENDED EIGHTH-SERIES V20 MATERIAL
% =============================================================================

\section{V20: companion audit and nonlinear completed-square
resolvent gate}
\label{sec:e8v20}

\subsection{Compulsory companion audit}

\begin{audit}[Files inspected at the V20 checkpoint]
\label{audit:e8v20-companion}
The latest active companion notes were inspected as mathematical
sources only.
\begin{enumerate}
\item Standard Model:
\[
\begin{split}
&\texttt{Renewal\_SM\_Einstein\_Continuum\_Research\_Notes\_V94.tex},\\
&37\,821\ {\rm lines},\qquad1\,365\,230\ {\rm bytes},\\
&\operatorname{SHA256}=
\texttt{BBDE2ED51094B2DDA45121E1CDB933E46ABCD9A0F8BCB1B9CB053462273B6E21}.
\end{split}
\]
\item Navier--Stokes:
\[
\begin{split}
&\texttt{Renewal\_NS\_Stability\_Research\_Notes\_V26.tex},\\
&5\,319\ {\rm lines},\qquad278\,283\ {\rm bytes},\\
&\operatorname{SHA256}=
\texttt{82C887F8C2C69E4F28FAD7C3DC8DE5B9099CB5A4BF431EBA1FFF3E91935978AD}.
\end{split}
\]
\end{enumerate}
\end{audit}

\begin{assessment}[Audit transfer]
\label{ass:e8v20-transfer}
SM V90 transports an energy-lift correction by tensor telescoping, but
uses the small-defect organization superseded here by the exact
matched-Wick contraction of V18.  Its formulas are therefore not
imported as a new YM hypothesis.  SM V91--V94 sharpen later
characteristic-functional, tail, constraint, and reflection-positivity
gates.  In particular, V94 confirms that an elliptic reconstruction
inverse must not be called local merely because the reduced measure is
positive.  That warning is consistent with the V19 two-norm split.
NS V26 is unchanged and supplies no new Yang--Mills resolvent constant.
\end{assessment}

\subsection{Perturbed vertical block}

Let the small running interaction change the completed-square blocks by
\[
 \widetilde A=A+E_A,\qquad
 \widetilde K=K+E_K.                                \tag{20.1}
\]
All blocks act on the same rooted Coulomb and coarse spaces.  Define
\[
 M=-A^{-1}K,\quad J=R_0+M,\qquad
 \widetilde M=-\widetilde A^{-1}\widetilde K,\quad
 \widetilde J=R_0+\widetilde M.                     \tag{20.2}
\]

\begin{theorem}[Exact perturbed completed-square identities]
\label{thm:e8v20-resolvent}
Whenever \(\widetilde A\) is invertible,
\[
\boxed{
\begin{aligned}
\widetilde A^{-1}-A^{-1}
 &=-A^{-1}E_A\widetilde A^{-1},\\
\widetilde M-M
 &=-\widetilde A^{-1}(E_K+E_AM),\\
\widetilde J-J&=\widetilde M-M.
\end{aligned}}                                                     \tag{20.3}
\]
For a differentiable path \(A_t,K_t\),
\[
\boxed{
\frac d{dt}A_t^{-1}=-A_t^{-1}\dot A_tA_t^{-1},\qquad
\dot J_t=-A_t^{-1}(\dot K_t+\dot A_tM_t).}          \tag{20.4}
\]
\end{theorem}

\begin{proof}
The first line is the resolvent identity.  The equations
\[
 AM=-K,\qquad \widetilde A\widetilde M=-\widetilde K
\]
give
\[
 \widetilde A(\widetilde M-M)=-(E_K+E_AM),
\]
which proves the second and third lines.  Divide the finite-difference
identities by the path increment and pass to the limit to obtain
(20.4).
\end{proof}

\begin{theorem}[Weighted nonlinear lift gate]
\label{thm:e8v20-weighted}
Fix \(\mu\) as in (19.4), and suppose
\[
 \delta_A:=|E_A|_\mu,\qquad
 \delta_K:=|E_K|_\mu,\qquad
 \theta_A:=g_\mu\delta_A<1.                        \tag{20.5}
\]
Then
\[
\boxed{
\begin{aligned}
|\widetilde A^{-1}|_\mu
 &\leq\frac{g_\mu}{1-\theta_A},\\
|\widetilde A^{-1}-A^{-1}|_\mu
 &\leq\frac{g_\mu^2\delta_A}{1-\theta_A},\\
|\widetilde J-J|_\mu
 &\leq\Delta j_\mu:=
 \frac{g_\mu}{1-\theta_A}
 \left(\delta_K+8e^{2\mu}g_\mu\delta_A\right).
\end{aligned}}                                                     \tag{20.6}
\]
Thus the perturbed lift and conditional covariance retain exponential
locality with the same admissible exponent \(\mu\).
\end{theorem}

\begin{proof}
Factor
\[
\widetilde A=A(I+A^{-1}E_A).
\]
The Neumann series converges in the submultiplicative norm (19.1), so
\[
|\widetilde A^{-1}|_\mu
\leq\frac{|A^{-1}|_\mu}{1-|A^{-1}E_A|_\mu}
\leq\frac{g_\mu}{1-\theta_A}.
\]
Insert this in the first line of (20.3) for the second estimate.
For the third, use the second line of (20.3) and
\[
 |M|_\mu=|A^{-1}K|_\mu
\leq8e^{2\mu}g_\mu.
\]
\end{proof}

\begin{corollary}[Tensor transport perturbation]
\label{cor:e8v20-tensor}
Put
\[
 \widetilde j_\mu=\max\{|J|_\mu,|\widetilde J|_\mu\}
 \leq j_\mu+\Delta j_\mu.
\]
For \(n\geq1\),
\[
\left|
(\widetilde J^*)^{\otimes n}-(J^*)^{\otimes n}
\right|_\mu
\leq n\widetilde j_\mu^{\,n-1}\Delta j_\mu.         \tag{20.7}
\]
\end{corollary}

\begin{proof}
Telescope the tensor powers and apply (20.6) to each summand.
\end{proof}

\subsection{Exact covariance rechart}

Let \(C\) and \(\widetilde C=C+\Delta C\) be two positive Gaussian
covariances on the same finite physical space.  Write
\({:X^{\otimes n}:}_C\) for symmetric Wick order with covariance \(C\).
For a symmetric tensor, \(\operatorname{Contr}_{\Delta C}^r\) denotes
the sum over \(r\) disjoint pair contractions with \(\Delta C\).

\begin{theorem}[Finite covariance-change formula]
\label{thm:e8v20-wick-rechart}
For every \(n\),
\[
\boxed{\quad
{:X^{\otimes n}:}_{\widetilde C}
=\sum_{r=0}^{\lfloor n/2\rfloor}
(-1)^r\frac{n!}{2^rr!(n-2r)!}\,
\operatorname{Sym}\!\left[
(\Delta C)^{\otimes r}
\otimes{:X^{\otimes(n-2r)}:}_C
\right].
\quad}                                                            \tag{20.8}
\]
In degrees at most four, the nonconstant rechart rows are
\[
\begin{aligned}
{:X^{\otimes3}:}_{\widetilde C}
 &={:X^{\otimes3}:}_C
   -3\,\operatorname{Sym}(\Delta C\otimes X),\\
{:X^{\otimes4}:}_{\widetilde C}
 &={:X^{\otimes4}:}_C
   -6\,\operatorname{Sym}
       (\Delta C\otimes{:X^{\otimes2}:}_C)
   +3\,\operatorname{Sym}((\Delta C)^{\otimes2}).                  \tag{20.9}
\end{aligned}
\]
\end{theorem}

\begin{proof}
The Wick exponentials obey the exact identity
\[
\exp\left(\langle u,X\rangle-\frac12\langle u,\widetilde Cu\rangle\right)
=
\exp\left(-\frac12\langle u,\Delta Cu\rangle\right)
\exp\left(\langle u,X\rangle-\frac12\langle u,Cu\rangle\right).
\]
Multiply the two power series and compare homogeneous degree \(n\).
Choosing \(r\) disjoint pairs gives
\[
\frac{n!}{2^rr!(n-2r)!},
\]
including the displayed coefficients \(3\) and \(6\).
\end{proof}

\begin{corollary}[Rechart debit in tensor norm]
\label{cor:e8v20-rechart-bound}
If contraction by \(\Delta C\) has operator norm at most \(\chi\) in
the declared one-particle coefficient metric, then the degree
\(n\) to degree \(n-2r\) rechart row is bounded by
\[
 \frac{n!}{2^rr!(n-2r)!}\chi^r.                   \tag{20.10}
\]
In particular, the quartic-to-quadratic debit is at most \(6\chi\);
it is a covariance-rechart term, not a defect of matched conditional
expectation.
\end{corollary}

\begin{proof}
Apply the contraction norm to every term of (20.8).
\end{proof}

\subsection{Perturbative contraction gate}

\begin{definition}[Nonlinear completed-square rechart card]
\label{def:e8v20-ncrc}
NCRC consists of:
\begin{enumerate}
\item weighted Hessian stocks \(\delta_A,\delta_K\) satisfying
\(\theta_A<1\);
\item the lift debit \(\Delta j_\mu\) in (20.6);
\item a relative covariance-rechart stock \(\chi\);
\item the degree-two through degree-four tensor rows from
(20.7) and (20.10);
\item the already declared large-field/marked remainder norm; and
\item their induced norm sum \(\varepsilon_{\rm NCRC}\) in the V19
hybrid norm.
\end{enumerate}
\end{definition}

\begin{theorem}[Nonlinear contraction stability]
\label{thm:e8v20-gate}
Let \(q_{\varepsilon,R}<1\) be the localized flat-principal factor in
(19.18).  If NCRC is certified and
\[
\boxed{\qquad
\varepsilon_{\rm NCRC}<1-q_{\varepsilon,R},
\qquad}                                                           \tag{20.11}
\]
then the complete small-field derivative, including the covariance
rechart, is a strict contraction in the same hybrid norm.
\end{theorem}

\begin{proof}
Write the derivative as
\[
\mathcal L_{\rm G}+\mathcal E_{\rm lift}
 +\mathcal E_{\rm Wick}+\mathcal E_{\rm rem}.
\]
The first term has norm at most \(q_{\varepsilon,R}\) by V19.
The remaining three terms have total induced norm at most
\(\varepsilon_{\rm NCRC}\) by the card.  The triangle inequality and
(20.11) give a norm strictly below one.
\end{proof}

\begin{assessment}[What V20 closes]
\label{ass:e8v20-closes}
The nonlinear completed-square and covariance-coordinate variation is
now reduced to explicit weighted inputs.  In particular, loss of the
vertical inverse is controlled by the single dimensionless number
\(\theta_A=g_\mu\delta_A\), and Wick degree lowering is exactly
proportional to the covariance change rather than to the absolute
conditional covariance.
\end{assessment}

\begin{assessment}[Current bottleneck]
\label{ass:e8v20-remains}
V20 does not yet derive \(\delta_A,\delta_K,\chi\), or the marked
remainder from the actual non-Abelian small-field action in the V19
hybrid norm.  The next task is an analytic Taylor estimate for those
stocks with their powers of the running coupling and small-field
radius exposed.  Only after that estimate can (20.11) be evaluated.
\end{assessment}

\begin{decision}[V21 acquisition order]
\label{dec:e8v20-v21}
V21 will differentiate the Wilson plaquette Hessian in left-trivialized
coordinates, use the six-plaquette link incidence and the finite
filling length, and obtain explicit bounds
\[
\delta_A+\delta_K\leq C_{\rm W}g\,r
\]
on the declared small-field ball.  It will then bound the relative
change of the vertical covariance by the resolvent formulas above and
separate the running marginal rescaling from the genuinely stable
covariance rechart.
\end{decision}

\begin{firewall}[Claim boundary after eighth-series V20]
\label{fire:e8v20-claim}
V20 proves the resolvent and covariance-rechart calculus and a
quantitative perturbative contraction implication.  It does not
evaluate NCRC for the interacting Wilson step, close large fields or
sources, construct the marginal trajectory, restore rotations, remove
the regulator, reconstruct the continuum OS theory, or prove the
Yang--Mills mass gap.
\end{firewall}

\paragraph{Parent-version record.}
V20 was formed by exact append-only extension of
\texttt{Renewal\_Yang\_Mills\_Mass\_Gap\_Research\_Notes\_V19.tex}.
The V19 parent had \(9\,000\) counted source lines,
\(314\,898\) bytes, and SHA--256
\[
\texttt{3155E383A6A35C7653359775BD58BC1B5CCCCCCAB2AC2B7CE5BA5CE187CA056A}.
\]

% =============================================================================
% END APPENDED EIGHTH-SERIES V20 MATERIAL
% =============================================================================

% =============================================================================
% BEGIN APPENDED EIGHTH-SERIES V21 MATERIAL
% =============================================================================

\section{V21: explicit Wilson Hessian Lipschitz stock and the
Ward-relative covariance bottleneck}
\label{sec:e8v21}

\subsection{Small-field normalization}

Fix a flat or commuting-toron background \(T\).  Write every fine link
in left coordinates as
\[
                         U_e=e^{gA_e}T_e,\qquad
             \sup_e\|A_e\|_F\leq r.                \tag{21.1}
\]
Put
\[
                         \eta:=gr.                 \tag{21.2}
\]
Since \(A_e\) is skew Hermitian,
\[
\|U_e-T_e\|_{\rm op}
\leq\|e^{gA_e}-I\|_{\rm op}
\leq g\|A_e\|_{\rm op}\leq\eta.                    \tag{21.3}
\]
The Hessians below are those of the unit Wilson action
\[
 V(U)=\sum_p\left(1-\frac13\operatorname{ReTr}U_p\right).          \tag{21.4}
\]

\begin{lemma}[One-plaquette Hessian Lipschitz bound]
\label{lem:e8v21-one-plaquette}
Let \(H_p(U)\) be the left-trivialized Hessian of one summand in
(21.4), regarded as a \(4\times4\) block matrix on its four oriented
link slots.  If all four links satisfy (21.3), then every Hessian block
obeys
\[
\boxed{\qquad
\|(H_p(U)-H_p(T))_{ef}\|_{\rm op}
\leq\frac43\eta.
\qquad}                                                           \tag{21.5}
\]
\end{lemma}

\begin{proof}
A left-trivialized mixed second derivative is the real trace of two
Lie-algebra insertions separated by two unitary arcs of the plaquette
word, divided by three.  When the two variations occupy the same
slot, polarization gives the symmetrized two-insertion word; the same
bound applies.  Inversion of an oriented link only conjugates an
insertion and changes its sign, neither of which changes Frobenius
norm.

For unit Frobenius insertions \(X,Y\),
\[
 |\operatorname{Tr}(XCYD)|\leq\|X\|_F\|Y\|_F=1
\]
when \(C,D\) are unitary.  Telescoping the four link factors of the
two arcs gives
\[
\|CYD-C_0YD_0\|_F\leq
\sum_{a\in p}\|U_a-T_a\|_{\rm op}\|Y\|_F\leq4\eta.
\]
After multiplication by \(X\), taking the real trace, and dividing by
three, this is (21.5).  Averaging the two orders in a diagonal block
does not enlarge the bound.
\end{proof}

\begin{proposition}[Explicit global Wilson Hessian stock]
\label{prop:e8v21-H-stock}
Every oriented fine link belongs to six plaquettes, and each such
plaquette has four link slots.  Therefore the complete unit-Wilson
Hessian satisfies
\[
\boxed{\qquad
 |H_{\rm W}(U)-H_{\rm W}(T)|_0\leq32\eta,
\qquad
 |H_{\rm W}(U)-H_{\rm W}(T)|_\mu
 \leq32e^{2\mu}\eta.
\qquad}                                                           \tag{21.6}
\]
\end{proposition}

\begin{proof}
For a fixed row link there are at most \(6\cdot4=24\) incident Hessian
blocks.  Sum (21.5) to obtain \(24(4/3)\eta=32\eta\).
The same count holds for columns because the Hessian is self-adjoint.
Two links in a plaquette have coarse-cell distance at most two, so the
weight costs at most \(e^{2\mu}\).
\end{proof}

\begin{corollary}[Vertical and mixed block stocks]
\label{cor:e8v21-AK}
Let \(E_A^{\rm W},E_K^{\rm W}\) be the changes of the vertical and
vertical/coarse blocks, using the coordinate inclusion on the vertical
space and the metric section \(R_0=B^*/2\).  Then
\[
\boxed{
 |E_A^{\rm W}|_\mu\leq32e^{2\mu}gr,\qquad
 |E_K^{\rm W}|_\mu\leq32e^{2\mu}gr.}               \tag{21.7}
\]
If a generated interaction \(\Psi\) contributes Hessian stocks
\(\zeta_A,\zeta_K\), the V20 inputs may be taken as
\[
\boxed{
\delta_A\leq32e^{2\mu}gr+\zeta_A,\qquad
\delta_K\leq32e^{2\mu}gr+\zeta_K.}                 \tag{21.8}
\]
\end{corollary}

\begin{proof}
Restriction to vertical coordinates cannot increase the absolute
block row or column sums.  The metric section assigns half of a coarse
edge to each of two constituents, so its absolute column sum is one;
composition with it does not increase the bound.  Add the declared
generated-interaction stocks.
\end{proof}

\begin{corollary}[Evaluated vertical resolvent gate]
\label{cor:e8v21-resolvent-gate}
The weighted nonlinear inverse in V20 is certified whenever
\[
\boxed{\qquad
 g_\mu\left(32e^{2\mu}gr+\zeta_A\right)<1.
\qquad}                                                           \tag{21.9}
\]
Under this inequality, insert (21.8) in (20.6) to obtain the explicit
lift debit
\[
\Delta j_\mu\leq
\frac{g_\mu}{
 1-g_\mu(32e^{2\mu}gr+\zeta_A)}
\left[
32e^{2\mu}gr+\zeta_K
+8e^{2\mu}g_\mu(32e^{2\mu}gr+\zeta_A)
\right].                                                          \tag{21.10}
\]
\end{corollary}

\begin{proof}
This is Theorem~\ref{thm:e8v20-weighted} with the explicit stocks
(21.8).
\end{proof}

\subsection{Vertical covariance rechart}

\begin{proposition}[Relative vertical covariance debit]
\label{prop:e8v21-vertical-covariance}
Let
\[
\vartheta_A:=
\|A^{-1/2}E_AA^{-1/2}\|
\leq\frac3{\lambda_*}\|E_A\|.                     \tag{21.11}
\]
If \(\vartheta_A<1\), then
\[
\boxed{\quad
\left\|
A^{1/2}(\widetilde A^{-1}-A^{-1})A^{1/2}
\right\|
\leq\frac{\vartheta_A}{1-\vartheta_A}.
\quad}                                             \tag{21.12}
\]
For the pure Wilson background perturbation,
\[
\vartheta_A\leq\frac{96}{\lambda_*}gr.             \tag{21.13}
\]
\end{proposition}

\begin{proof}
The exact vertical floor is \(A\succeq\lambda_*I/3\), proving the
inequality in (21.11).  Put
\[
 X=A^{-1/2}E_AA^{-1/2}.
\]
Then
\[
 A^{1/2}\widetilde A^{-1}A^{1/2}=(I+X)^{-1},
\]
so the left side of (21.12) is
\(\|(I+X)^{-1}-I\|\leq\|X\|/(1-\|X\|)\).
Use \(\|E_A^{\rm W}\|\leq32gr\) for (21.13).
\end{proof}

\begin{corollary}[Vertical Wick rechart is quantitatively closed]
\label{cor:e8v21-vertical-wick}
In the covariance metric of \(A^{-1}\), every vertical Wick rechart
row of degree at most four is bounded by (20.10) with
\[
                    \chi_A=\frac{\vartheta_A}{1-\vartheta_A}.      \tag{21.14}
\]
Thus the pure-Wilson vertical quartic-to-quadratic rechart debit is at
most
\[
                         6\chi_A.                                 \tag{21.15}
\]
\end{corollary}

\begin{proof}
Apply Proposition~\ref{prop:e8v21-vertical-covariance} and
Corollary~\ref{cor:e8v20-rechart-bound}.
\end{proof}

\subsection{The coarse Schur change}

Let
\[
\widetilde D=D+E_D,\qquad
\widetilde S=\widetilde D-\widetilde K^*
                    \widetilde A^{-1}\widetilde K.                 \tag{21.16}
\]

\begin{proposition}[Exact finite Schur-change identity]
\label{prop:e8v21-Schur-change}
With \(\Delta M=\widetilde M-M\),
\[
\boxed{\qquad
\widetilde S-S
=E_D+E_K^*M+\widetilde K^*\Delta M.
\qquad}                                                           \tag{21.17}
\]
Consequently,
\[
|\widetilde S-S|_\mu
\leq\delta_D+\delta_K|M|_\mu
 +( |K|_\mu+\delta_K)\Delta j_\mu.                \tag{21.18}
\]
\end{proposition}

\begin{proof}
Use \(S=D+K^*M\) and
\(\widetilde S=\widetilde D+\widetilde K^*\widetilde M\):
\[
\widetilde K^*\widetilde M-K^*M
=E_K^*M+\widetilde K^*(\widetilde M-M).
\]
This proves (21.17); take the weighted norm and use V20 for (21.18).
\end{proof}

\begin{firewall}[Absolute Schur control is not relative covariance
control]
\label{fire:e8v21-absolute}
The flat physical Schur form vanishes quadratically at zero momentum:
\[
                         S(k)\asymp r(k).
\]
Therefore a uniform absolute estimate
\(\|\widetilde S-S\|\leq\delta_S\) does not bound
\[
\|S^{-1/2}(\widetilde S-S)S^{-1/2}\|
\]
as \(k\to0\).  Equation (21.18), by itself, does not close the coarse
matched-Wick covariance rechart.
\end{firewall}

\subsection{Ward-relative factorization}

\begin{lemma}[Abstract Hodge factorization]
\label{lem:e8v21-Hodge}
Let \(L(k)\) be a self-adjoint coarse edge Hessian at nonzero momentum
with
\[
                         L(k)d_0(k)=0.                             \tag{21.19}
\]
Then
\[
                         L(k)=d_c(k)^*\mathcal R(k)d_c(k)           \tag{21.20}
\]
on the physical edge space, where one possible choice is
\[
\mathcal R(k)=
(d_c(k)^\dagger)^*L(k)d_c(k)^\dagger.              \tag{21.21}
\]
This statement alone does not make \(\mathcal R(k)\) local or
uniformly bounded at \(k=0\).
\end{lemma}

\begin{proof}
At nonzero momentum the edge space is the orthogonal sum of
\(\operatorname{Ran}d_0\) and \(\operatorname{Ran}d_c^*\).
Condition (21.19) and self-adjointness make \(L\) vanish on the first
summand.  On the second, \(d_c^\dagger d_c\) is the identity.
Inserting those projections on both sides of \(L\) gives
(21.20)--(21.21).  The pseudoinverse contains \(r(k)^{-1/2}\), proving
the warning.
\end{proof}

\begin{definition}[Ward-relative Schur card]
\label{def:e8v21-wrsc}
WRSC requires an exact factorization of the interacting Schur change
\[
 \widetilde S-S=d_c^*\mathcal R_{\Delta}d_c
                 +\mathcal E_{\rm stat},                           \tag{21.22}
\]
where:
\begin{enumerate}
\item \(\mathcal R_\Delta\) is quasi-local and
\(\|\mathcal R_\Delta\|\leq\varepsilon_{\rm W}\);
\item the stationarity debit \(\mathcal E_{\rm stat}\), arising from
expansion away from a critical background, is assigned to the moving
reference graph rather than hidden in the covariance; and
\item the exact flat floor
\[
                         S\succeq\frac1{12}d_c^*d_c                 \tag{21.23}
\]
is used in the same normalization.
\end{enumerate}
\end{definition}

\begin{proposition}[WRSC gives relative coarse covariance control]
\label{prop:e8v21-WRSC}
If \(\mathcal E_{\rm stat}=0\), then
\[
\boxed{\qquad
\|S^{-1/2}(\widetilde S-S)S^{-1/2}\|
\leq12\varepsilon_{\rm W}.
\qquad}                                                           \tag{21.24}
\]
If \(12\varepsilon_{\rm W}<1\), the relative coarse covariance change
is bounded by
\[
 \left\|S^{1/2}(\widetilde S^{-1}-S^{-1})S^{1/2}\right\|
\leq\frac{12\varepsilon_{\rm W}}
          {1-12\varepsilon_{\rm W}}.               \tag{21.25}
\]
\end{proposition}

\begin{proof}
For every physical \(h\), (21.22)--(21.23) give
\[
 |\langle h,(\widetilde S-S)h\rangle|
\leq\varepsilon_{\rm W}\|d_ch\|^2
\leq12\varepsilon_{\rm W}\langle h,Sh\rangle.
\]
This is (21.24).  Apply the same normalized resolvent argument as in
Proposition~\ref{prop:e8v21-vertical-covariance} to obtain (21.25).
\end{proof}

\begin{assessment}[What V21 closes]
\label{ass:e8v21-closes}
The formerly abstract pure-Wilson Hessian stock is now explicit:
each of the vertical and mixed perturbations is at most
\(32e^{2\mu}gr\).  This evaluates the weighted vertical inverse gate,
the lift debit, and the vertical covariance-Wick rechart.  The exact
Schur-change identity also separates absolute control from the
relative massless covariance problem.
\end{assessment}

\begin{assessment}[Current bottleneck]
\label{ass:e8v21-remains}
The remaining small-field covariance obstruction is WRSC: one must
derive a quasi-local, uniformly bounded curvature factor
\(\mathcal R_\Delta\) for the interacting Schur change and isolate the
nonstationary first-derivative debit.  Gauge annihilation at a flat
critical point gives only the algebraic factorization of
Lemma~\ref{lem:e8v21-Hodge}; it does not by itself supply the required
uniform local bound.
\end{assessment}

\begin{decision}[V22 acquisition order]
\label{dec:e8v21-v22}
V22 will go upstream to the exact gauge Ward identity for a
left-trivialized Hessian at a noncritical background.  It will derive
the term involving the action gradient, subtract that term into the
moving reference graph, and test whether the remaining stationary
Hessian admits a local two-curl factorization with an explicit
\(O(gr)+O(\|\Psi\|)\) coefficient.
\end{decision}

\begin{firewall}[Claim boundary after eighth-series V21]
\label{fire:e8v21-claim}
V21 evaluates the pure-Wilson vertical and mixed Hessian perturbation
stocks and closes their vertical resolvent/rechart consequences.  It
does not prove WRSC, the complete nonlinear contraction, large-field
and source estimates, the marginal trajectory, rotational
restoration, regulator removal, continuum reconstruction, or the
Yang--Mills mass gap.
\end{firewall}

\paragraph{Parent-version record.}
V21 was formed by exact append-only extension of
\texttt{Renewal\_Yang\_Mills\_Mass\_Gap\_Research\_Notes\_V20.tex}.
The V20 parent had \(9\,332\) counted source lines,
\(325\,562\) bytes, and SHA--256
\[
\texttt{E331D9324F36E6D3A9EE6C121D0F65D985319E558FB26DCE7D853F0A3AFE51F3}.
\]

% =============================================================================
% END APPENDED EIGHTH-SERIES V21 MATERIAL
% =============================================================================

% =============================================================================
% BEGIN APPENDED EIGHTH-SERIES V22 MATERIAL
% =============================================================================

\section{V22: differentiated gauge Ward identity and an infrared-relative
Schur bound}
\label{sec:e8v22}

\subsection{The Ward identity away from a critical background}

Let \(e=(x,y)\) be an oriented link.  Write tangent vectors in the
left frame, so \(\delta U_e=X_eU_e\).  For a vertex gauge parameter
\(\omega\), the fundamental gauge vector is
\[
 (D_U\omega)_e=\omega_x-\operatorname{Ad}_{U_e}\omega_y.           \tag{22.1}
\]
Let \(\mathcal A(U)\) be a twice differentiable gauge-invariant action,
let \(G_e(U)\) be its left gradient, and let \(H_U\) be its
left-trivialized symmetric Hessian.

\begin{theorem}[Differentiated lattice Ward identity]
\label{thm:e8v22-differentiated-Ward}
For every \(X\) and \(\omega\),
\[
\boxed{\quad
\langle D_U\omega,H_UX\rangle
=\sum_{e=(x,y)}
\left\langle G_e(U),
[X_e,\operatorname{Ad}_{U_e}\omega_y]\right\rangle .
\quad}                                                            \tag{22.2}
\]
In particular, at a critical configuration,
\[
                         H_UD_U=0.                                \tag{22.3}
\]
\end{theorem}

\begin{proof}
Gauge invariance gives the first Ward identity
\[
\sum_e\langle G_e(U),
\omega_x-\operatorname{Ad}_{U_e}\omega_y\rangle=0.                 \tag{22.4}
\]
Differentiate along \(U_e(t)=e^{tX_e}U_e\).  In the left frame,
\[
\left.\frac d{dt}\right|_{0}
\operatorname{Ad}_{e^{tX_e}U_e}\omega_y
=[X_e,\operatorname{Ad}_{U_e}\omega_y].
\]
The derivative of the gradient term is the Hessian pairing.  Moving
the derivative of the fundamental vector to the other side gives
(22.2).  If \(G=0\), the right side vanishes for every \(X,\omega\);
self-adjointness then gives (22.3).
\end{proof}

\begin{corollary}[Explicit stationarity debit]
\label{cor:e8v22-stationarity}
With Frobenius norms,
\[
\left|
\langle D_U\omega,H_UX\rangle
\right|
\leq
2\sum_{e=(x,y)}
\|G_e(U)\|_F\,\|X_e\|_F\,\|\omega_y\|_F.            \tag{22.5}
\]
Thus the failure of the Hessian to annihilate gauge directions is
linear in the action gradient, not an independent mass term.
\end{corollary}

\begin{proof}
Use
\[
\|[X,Y]\|_F\leq2\|X\|_F\|Y\|_F
\]
and invariance of Frobenius norm under \(\operatorname{Ad}_{U_e}\).
\end{proof}

\begin{corollary}[Effective-action Ward identity]
\label{cor:e8v22-effective}
An exactly integrated or exactly minimized coarse action obeys the
same formulas with coarse links, its coarse gradient, and its coarse
Hessian.  Therefore the stationarity term in WRSC is fixed by the
coarse action gradient and must be assigned to the moving reference
graph when that gradient is nonzero.
\end{corollary}

\begin{proof}
The fibre measure, Haar measure, and block constraint are equivariant,
so exact integration preserves gauge invariance.  Exact minimization
does likewise by uniqueness of the minimizer modulo the rooted gauge.
Apply Theorem~\ref{thm:e8v22-differentiated-Ward} to the resulting
coarse action.
\end{proof}

\subsection{Stationary translation-invariant Schur perturbations}

Let
\[
                         L=\widetilde S-S                           \tag{22.6}
\]
be a translation-invariant self-adjoint coarse edge kernel at a flat
critical background.  Use the midpoint Fourier convention, in which
coordinate reflections make its symbol even:
\[
                         L(-k)=L(k).                               \tag{22.7}
\]
Let
\[
 \mathfrak m_2(L):=
\max\left\{
\sup_x\sum_y d(x,y)^2\|L_{xy}\|_{\rm op},
\sup_y\sum_x d(x,y)^2\|L_{xy}\|_{\rm op}
\right\}.                                                         \tag{22.8}
\]

\begin{lemma}[Ward annihilation removes the zero-momentum head]
\label{lem:e8v22-zero}
If
\[
                         L(k)d_0(k)=0                              \tag{22.9}
\]
for all \(k\), then \(L(0)=0\).  Reflection symmetry (22.7) also gives
\[
                         \nabla_kL(0)=0.                           \tag{22.10}
\]
\end{lemma}

\begin{proof}
For \(k=tv\),
\[
 d_0(tv)=itv+O(t^2).
\]
Divide \(L(tv)d_0(tv)=0\) by \(t\) and let \(t\to0\).  This gives
\(L(0)v=0\) for every \(v\in\mathbb R^4\), hence \(L(0)=0\).
Differentiating the even relation (22.7) at zero gives (22.10).
\end{proof}

\begin{theorem}[Second-moment Ward-relative bound]
\label{thm:e8v22-second-moment}
Under (22.7)--(22.9), for every
\(k\in[-\pi,\pi]^4\),
\[
\boxed{\qquad
\|L(k)\|_{\rm op}
\leq\frac{\pi^2}{8}\,\mathfrak m_2(L)\,
            r(k),\qquad
r(k)=\sum_\mu|e^{ik_\mu}-1|^2.
\qquad}                                                           \tag{22.11}
\]
\end{theorem}

\begin{proof}
Lemma~\ref{lem:e8v22-zero} and Taylor's formula give
\[
 L(k)=\int_0^1(1-t)\,D^2L(tk)[k,k]\,dt.
\]
Writing the Fourier series of \(L\) and differentiating twice yields
\[
\|D^2L(tk)[k,k]\|
\leq |k|^2\mathfrak m_2(L).
\]
The integral contributes \(1/2\).  On \([-\pi,\pi]\),
\[
 |e^{is}-1|=2|\sin(s/2)|\geq\frac{2}{\pi}|s|,
\]
so \(|k|^2\leq(\pi^2/4)r(k)\).  Combining the two inequalities proves
(22.11).
\end{proof}

\begin{corollary}[Weighted locality controls the second moment]
\label{cor:e8v22-moment}
For every \(\mu>0\),
\[
\boxed{\qquad
\mathfrak m_2(L)\leq\frac{2}{\mu^2}|L|_\mu.
\qquad}                                                           \tag{22.12}
\]
\end{corollary}

\begin{proof}
The elementary inequality \(t^2\leq2e^t\), applied to
\(t=\mu d(x,y)\), gives
\[
d(x,y)^2\leq2\mu^{-2}e^{\mu d(x,y)}.
\]
Insert this in (22.8).
\end{proof}

\subsection{Relative coarse covariance without a local square root}

\begin{theorem}[Stationary coarse rechart gate]
\label{thm:e8v22-coarse-gate}
Assume the exact flat floor
\[
                         S(k)\succeq\frac1{12}d_c(k)^*d_c(k)        \tag{22.13}
\]
on the transverse carrier.  If \(L\) satisfies the hypotheses above,
then
\[
\boxed{\quad
\vartheta_S:=
\|S^{-1/2}LS^{-1/2}\|
\leq\frac{3\pi^2}{2}\mathfrak m_2(L)
\leq\frac{3\pi^2}{\mu^2}|L|_\mu.
\quad}                                                            \tag{22.14}
\]
If the last quantity is below one, then
\[
\boxed{\quad
\left\|S^{1/2}(\widetilde S^{-1}-S^{-1})S^{1/2}\right\|
\leq\frac{\vartheta_S}{1-\vartheta_S}.
\quad}                                                            \tag{22.15}
\]
\end{theorem}

\begin{proof}
On the transverse edge space,
\[
 d_c^*d_c=r(k)I,
\qquad
 S(k)\succeq\frac{r(k)}{12}I.
\]
Theorem~\ref{thm:e8v22-second-moment} therefore gives
\[
\|S(k)^{-1/2}L(k)S(k)^{-1/2}\|
\leq12\,\frac{\|L(k)\|}{r(k)}
\leq\frac{3\pi^2}{2}\mathfrak m_2(L).
\]
Use Corollary~\ref{cor:e8v22-moment} for the second inequality in
(22.14).  The normalized resolvent identity gives (22.15).
\end{proof}

\begin{correction}[WRSC can be weakened at a stationary bulk point]
\label{corr:e8v22-WRSC}
For the translation-invariant stationary bulk covariance, the
quasi-local factor \(\mathcal R_\Delta\) demanded in V21 is sufficient
but not necessary.  It is enough that:
\begin{enumerate}
\item \(L\) itself has a finite weighted kernel norm;
\item the exact Ward row \(Ld_0=0\) holds;
\item midpoint reflection makes \(L(k)\) even.
\end{enumerate}
Theorem~\ref{thm:e8v22-coarse-gate} then supplies the relative
covariance bound.  The strong spatial norm continues to use the local
kernel \(L\), while the weak covariance norm uses (22.14).
\end{correction}

\subsection{Insertion of the V21 stocks}

The horizontal/horizontal Wilson block has the same incidence estimate
as the vertical and mixed blocks:
\[
 |E_D^{\rm W}|_\mu\leq32e^{2\mu}gr.                \tag{22.16}
\]
Let \(\zeta_D\) be the generated-interaction contribution and set
\[
 \delta_D=32e^{2\mu}gr+\zeta_D.                    \tag{22.17}
\]
Define the explicit Schur-change stock
\[
\Delta_S^\mu:=
\delta_D+\delta_K\,8e^{2\mu}g_\mu
 +(8e^{2\mu}+\delta_K)\Delta j_\mu,                \tag{22.18}
\]
where \(\delta_A,\delta_K,\Delta j_\mu\) are given by
(21.8)--(21.10).

\begin{proposition}[Evaluated stationary coarse gate]
\label{prop:e8v22-evaluated}
At a flat critical, translation- and reflection-invariant bulk point,
\[
                         |L|_\mu\leq\Delta_S^\mu.                  \tag{22.19}
\]
Consequently the coarse covariance rechart is certified whenever
\[
\boxed{\qquad
\frac{3\pi^2}{\mu^2}\Delta_S^\mu<1.
\qquad}                                                           \tag{22.20}
\]
Its Wick-rechart contraction stock may then be taken as
\[
\chi_S=
\frac{(3\pi^2/\mu^2)\Delta_S^\mu}
     {1-(3\pi^2/\mu^2)\Delta_S^\mu}.               \tag{22.21}
\]
\end{proposition}

\begin{proof}
Equation (21.18), the bound
\(|M|_\mu\leq8e^{2\mu}g_\mu\), and
\(|K|_\mu\leq8e^{2\mu}\) give (22.19).
Apply Theorem~\ref{thm:e8v22-coarse-gate} and then the Wick rechart
bound of V20.
\end{proof}

\begin{firewall}[Nonstationary backgrounds]
\label{fire:e8v22-nonstationary}
At a noncritical coarse background, (22.9) is replaced by the
gradient term (22.2).  One may not apply (22.20) until that term has
been subtracted into the moving reference graph or bounded in a norm
which supplies the missing powers of momentum.  Absolute smallness of
the gradient is not, by itself, relative to \(r(k)\).
\end{firewall}

\begin{assessment}[What V22 closes]
\label{ass:e8v22-closes}
V22 proves the exact noncritical Ward identity, identifies the
stationarity debit, and closes the stationary translation-invariant
coarse covariance rechart by an explicit second-moment gate.  The
formerly requested local two-curl factorization is no longer needed
for this bulk row.
\end{assessment}

\begin{assessment}[Current bottleneck]
\label{ass:e8v22-remains}
The remaining local-background issue is the action-gradient term in
(22.2).  The moving reference graph already carries a scale-dependent
one-point residue, but it has not yet been shown that its induced Ward
row obeys the same two-norm contraction and moment cancellations as
the stable Hessian sector.  This must be proved before the stationary
gate (22.20) can be used uniformly over the invariant small-field
ball.
\end{assessment}

\begin{decision}[V23 acquisition order]
\label{dec:e8v22-v23}
V23 will define the covariant quotient Hessian by subtracting the
connection-gradient bilinear form dictated by (22.2).  It will prove
that this quotient Hessian annihilates gauge directions at every
background, derive its local kernel bound, and assign the remaining
one-point gradient transport to the positive moving reference graph.
\end{decision}

\begin{firewall}[Claim boundary after eighth-series V22]
\label{fire:e8v22-claim}
V22 closes only the stationary bulk coarse-covariance rechart and
derives the exact nonstationary debit.  It does not close the
background-uniform nonlinear contraction, sources, large fields,
marginal trajectory, rotation restoration, regulator removal,
continuum reconstruction, or the Yang--Mills mass gap.
\end{firewall}

\paragraph{Parent-version record.}
V22 was formed by exact append-only extension of
\texttt{Renewal\_Yang\_Mills\_Mass\_Gap\_Research\_Notes\_V21.tex}.
The V21 parent had \(9\,725\) counted source lines,
\(338\,141\) bytes, and SHA--256
\[
\texttt{C053AB5122981B30D48BA31AAF62C298E9BA6D8553CAB5E4EC86EA2969FE06E8}.
\]

% =============================================================================
% END APPENDED EIGHTH-SERIES V22 MATERIAL
% =============================================================================

% =============================================================================
% BEGIN APPENDED EIGHTH-SERIES V23 MATERIAL
% =============================================================================

\section{V23: stationary covariance policy and explicit cubic Wilson
remainder}
\label{sec:e8v23}

\subsection{Why the covariance must not follow the field point}

V22 identifies the obstruction at a noncritical background exactly:
the Hessian gauge row is proportional to the action gradient.  There
is no need to force that row into a massless covariance rechart.
The Gaussian reference can instead be defined at the stationary
running vacuum, while all field dependence is retained in the
interaction.

\begin{definition}[Stationary covariance policy]
\label{def:e8v23-scp}
At scale \(j\), let \(\mathcal A_j(a)\) be the translation-invariant
effective action in a fixed rooted gauge chart about the identity
vacuum.  SCP means:
\begin{enumerate}
\item the Gaussian precision is
\[
                         H_j=D^2\mathcal A_j(0);                    \tag{23.1}
\]
\item Wick ordering at scale \(j\) uses \(H_j^\dagger\) on the
physical quotient;
\item the field-dependent remainder is
\[
\mathcal N_j(a)=\mathcal A_j(a)-\mathcal A_j(0)
 -\frac12\langle a,H_ja\rangle;                                   \tag{23.2}
\]
\item after one exact RG step, the target covariance is recentered
only at \(D^2\mathcal A_{j+1}(0)\), never at
\(D^2\mathcal A_{j+1}(a)\) for arbitrary \(a\).
\end{enumerate}
\end{definition}

\begin{lemma}[The identity vacuum is stationary]
\label{lem:e8v23-stationary}
Suppose \(\mathcal A_j\) is differentiable and invariant under global
\(\mathrm{SU}(3)\) conjugation.  At the identity configuration,
\[
                         D\mathcal A_j(0)=0.                        \tag{23.3}
\]
This remains true after an exact gauge-equivariant block integration.
\end{lemma}

\begin{proof}
For each oriented link type, the left gradient at the identity must
be fixed by \(\operatorname{Ad}_g\) for every \(g\in\mathrm{SU}(3)\).
The adjoint representation of the simple Lie algebra
\(\mathfrak{su}(3)\) has no nonzero invariant vector.  Hence every
gradient component vanishes.  Exact block integration preserves
global conjugation invariance because both Haar measure and the block
constraint are equivariant, so the argument repeats at the next
scale.
\end{proof}

\begin{proposition}[Ward and reflection properties of the SCP
precision]
\label{prop:e8v23-SCP-Ward}
On the physical translation-invariant bulk,
\[
                         H_jd_0=0.                                 \tag{23.4}
\]
If the regulator and block map preserve coordinate reflections, the
midpoint symbol of \(H_j\) is even.  Consequently every interscale
stationary precision change
\[
                         H_{j+1}-H_j                               \tag{23.5}
\]
satisfies the hypotheses of the V22 second-moment relative bound.
\end{proposition}

\begin{proof}
Apply Theorem~\ref{thm:e8v22-differentiated-Ward} at the stationary
point (23.3).  Reflection invariance is preserved by an exact
equivariant integral and gives evenness in the midpoint convention.
Both properties are stable under subtraction.
\end{proof}

\begin{theorem}[SCP removes the nonstationary covariance obstruction]
\label{thm:e8v23-removal}
Under SCP, the gradient term in (22.2) never enters the definition of
a Gaussian covariance.  Every covariance rechart is between
stationary translation/reflection-invariant precisions and is
therefore governed by Theorem~\ref{thm:e8v22-coarse-gate}.  At
nonzero field \(a\), the differentiated-Ward gradient term belongs
entirely to derivatives of \(\mathcal N_j\).
\end{theorem}

\begin{proof}
By Definition~\ref{def:e8v23-scp}, the only Gaussian precisions are
the Hessians at \(a=0\), where the gradient vanishes.  Proposition
\ref{prop:e8v23-SCP-Ward} supplies Ward annihilation and reflection
evenness for their differences.  Subtracting the fixed quadratic
term from the full action leaves every variation of the Hessian with
the field inside \(D^2\mathcal N_j(a)\), including the right side of
(22.2).
\end{proof}

\subsection{Exact normalized Wilson Taylor remainder}

For the unit Wilson action \(V\), define its weak-coupling
field-normalized form
\[
                         \mathcal V_g(a):=g^{-2}V(e^{ga}).          \tag{23.6}
\]
The factor \(g^{-2}\) cancels the two derivatives of the exponential
coordinate, so
\[
 H_0:=D^2\mathcal V_g(0)=\frac13d_1^*d_1                         \tag{23.7}
\]
is independent of \(g\).

\begin{theorem}[Explicit Wilson Taylor bounds]
\label{thm:e8v23-Wilson-Taylor}
On the ball
\[
               \|a\|_\infty\leq r,\qquad gr\leq1,                  \tag{23.8}
\]
the normalized Wilson remainder
\[
\mathcal R_g(a):=\mathcal V_g(a)-\mathcal V_g(0)
                  -\frac12\langle a,H_0a\rangle                   \tag{23.9}
\]
is compared along radial product-group geodesics; gradient and Hessian
values at \(e^{ga}\) are parallel transported to the identity frame.
It
obeys
\[
\boxed{
\begin{aligned}
|\mathcal R_g(a)|
 &\leq\frac{32}{3}gr\,\|a\|_2^2,\\
\|\nabla\mathcal R_g(a)\|_2
 &\leq32gr\,\|a\|_2,\\
\|D^2\mathcal R_g(a)\|_{\rm op}
 &\leq64gr.
\end{aligned}}                                                     \tag{23.10}
\]
In the weighted block kernel norm, the last line becomes
\[
                         |D^2\mathcal R_g(a)|_\mu
 \leq64e^{2\mu}gr.                                \tag{23.11}
\]
\end{theorem}

\begin{proof}
Proposition~\ref{prop:e8v21-H-stock} compares left-trivialized
Hessians and gives \(32tgr\).  Radial Taylor comparison uses parallel
transport.  For the bi-invariant product metric, parallel transport
along \(e^{tga}\) differs in the left frame by
\(\operatorname{Ad}_{e^{-tga/2}}\).  Since
\(\|\operatorname{ad}_X\|\leq2\|X\|_F\) and \(tgr\leq1\),
\[
\|\operatorname{Ad}_{e^{-tga/2}}-I\|
\leq e^{tgr}-1\leq2tgr.
\]
The unit Wilson Hessian has operator norm at most eight, so transporting
its two arguments costs at most another \(32tgr\).  Hence
\[
\|D^2\mathcal V_g(ta)-H_0\|\leq64tgr.              \tag{23.12}
\]
Because the first derivative at zero vanishes,
\[
\nabla\mathcal R_g(a)
=\int_0^1\bigl(D^2\mathcal V_g(ta)-H_0\bigr)a\,dt,
\]
which gives the second line of (23.10).  Taylor's formula with integral
remainder gives
\[
\mathcal R_g(a)
=\int_0^1(1-t)
\langle a,(D^2\mathcal V_g(ta)-H_0)a\rangle\,dt.
\]
Since \(\int_0^1t(1-t)\,dt=1/6\), the first line follows.
The third line is (23.12) at \(t=1\), and the propagation-two weight
gives (23.11).
\end{proof}

\begin{corollary}[The nonstationary Ward row is cubic]
\label{cor:e8v23-Ward-cubic}
For \(\mathcal V_g\), the stationarity debit (22.5) at \(a\) is bounded
by the gradient remainder in (23.10).  In particular it vanishes to
second order in the field:
\[
\|\nabla\mathcal R_g(a)\|_2
\leq32g\|a\|_\infty\|a\|_2.                       \tag{23.13}
\]
It is therefore a cubic-and-higher interaction vertex under SCP, not
a quadratic covariance term.
\end{corollary}

\begin{proof}
The fixed quadratic \(H_0\) satisfies the stationary Ward identity.
The difference between the full gradient and \(H_0a\) is exactly
\(\nabla\mathcal R_g(a)\).  Use Theorem
\ref{thm:e8v23-Wilson-Taylor}.
\end{proof}

\subsection{Generated interactions and the nonlinear derivative gate}

Let the generated stable interaction satisfy on the same ball
\[
 |\nabla^2\Psi_j(a)-\nabla^2\Psi_j(0)|_\mu
 \leq\zeta_{3,j}r.                                  \tag{23.14}
\]
Absorb \(\nabla^2\Psi_j(0)\) into the stationary precision \(H_j\), as
required by SCP.

\begin{proposition}[Background-uniform nonlinear Hessian stock]
\label{prop:e8v23-nonlinear-stock}
The interaction remainder in the SCP chart obeys
\[
\boxed{\qquad
 |D^2\mathcal N_j(a)|_\mu
 \leq\kappa_j(r):=
       64e^{2\mu}gr+\zeta_{3,j}r.
\qquad}                                                           \tag{23.15}
\]
\end{proposition}

\begin{proof}
The Wilson contribution is (23.11).  Subtracting
\(\nabla^2\Psi_j(0)\) leaves the difference in (23.14).  Add the two
weighted norms.
\end{proof}

\begin{theorem}[SCP small-field contraction criterion]
\label{thm:e8v23-contraction}
Let \(q_{\varepsilon,R}<1\) be the localized stationary Gaussian
factor of V19.  Suppose the finite owner and conditional-expectation
estimates convert a Hessian kernel stock into the hybrid derivative
norm with constant \(C_{\rm own}\), and suppose the stationary
interscale covariance rechart has norm debit
\(\varepsilon_{\rm stat}\) certified by V22.  If
\[
\boxed{\qquad
 q_{\varepsilon,R}
 +\varepsilon_{\rm stat}
 +C_{\rm own}\kappa_j(r)<1,
\qquad}                                                           \tag{23.16}
\]
then the complete small-field RG map is a strict contraction on the
radius-\(r\) SCP stable ball.
\end{theorem}

\begin{proof}
The derivative at the stationary origin is the V19 Gaussian map plus
the stationary covariance rechart, with norm at most
\(q_{\varepsilon,R}+\varepsilon_{\rm stat}\).
The mean-value formula and Proposition
\ref{prop:e8v23-nonlinear-stock} bound the derivative variation
throughout the ball by \(C_{\rm own}\kappa_j(r)\).
Their sum is below one by (23.16).
\end{proof}

\begin{firewall}[What remains in \(C_{\rm own}\)]
\label{fire:e8v23-owner}
The constant \(C_{\rm own}\) is finite at fixed local grammar, but it
has not yet been evaluated in the V19 hybrid norm.  It must include
the exact conditional lift tail, the single plaquette-owner
assignment, tensor symmetrization, and every retained interaction
type.  Equation (23.16) is not a completed numerical contraction card
until that finite enumeration is supplied.
\end{firewall}

\begin{assessment}[What V23 closes]
\label{ass:e8v23-closes}
The nonstationary Ward row is no longer a covariance obstruction.
SCP confines covariance recharting to stationary Ward-compatible
precisions, and the field-dependent Wilson term is an explicitly
bounded cubic remainder.  The pure-Wilson contribution to the
nonlinear Hessian stock is \(64e^{2\mu}gr\), including the radial
parallel-transport frame cost.
\end{assessment}

\begin{assessment}[Current bottleneck]
\label{ass:e8v23-remains}
The next finite task is the evaluation of \(C_{\rm own}\) for the
declared degrees and owner grammar, together with the generated
third-derivative stock \(\zeta_{3,j}\).  After that, (23.16) becomes a
checkable invariant-ball inequality.  Large-field and marginal-flow
rows remain separate.
\end{assessment}

\begin{decision}[V24 acquisition order]
\label{dec:e8v23-v24}
V24 will enumerate the finite owner conversion for degrees two
through four using the exact lift head/tail split of V19.  It will
keep the head radius \(R\) symbolic, derive an explicit
\(C_{\rm own}(R)\), and optimize only after all multiplicities are
visible.  It will not hide the six-fold link-hull conversion inside
each tensor leg.
\end{decision}

\begin{firewall}[Claim boundary after eighth-series V23]
\label{fire:e8v23-claim}
V23 proves the stationary covariance organization and explicit cubic
Wilson bounds.  It does not evaluate the owner constant or generated
interaction stock, close the full nonlinear and large-field RG,
construct the marginal trajectory, restore rotations, remove the
regulator, reconstruct the continuum theory, or prove the
Yang--Mills mass gap.
\end{firewall}

\paragraph{Parent-version record.}
V23 was formed by exact append-only extension of
\texttt{Renewal\_Yang\_Mills\_Mass\_Gap\_Research\_Notes\_V22.tex}.
The V22 parent had \(10\,081\) counted source lines,
\(349\,604\) bytes, and SHA--256
\[
\texttt{72EDE7370D6BC47B04B35984E3F5D25C9AE7A848B57D8C1FE0BA14720D802EEB}.
\]

% =============================================================================
% END APPENDED EIGHTH-SERIES V23 MATERIAL
% =============================================================================

% =============================================================================
% BEGIN APPENDED EIGHTH-SERIES V24 MATERIAL
% =============================================================================

\section{V24: exact owner-head enumeration and a concrete hybrid norm}
\label{sec:e8v24}

\subsection{Four-dimensional head cardinality}

Let
\[
 B_4(R)=\{x\in\mathbb Z^4:\|x\|_1\leq R\},\qquad
 N_4(R)=|B_4(R)|.                                  \tag{24.1}
\]

\begin{lemma}[Exact lattice-ball count]
\label{lem:e8v24-ball}
For every integer \(R\geq0\),
\[
\boxed{\quad
N_4(R)=
1+8R+12R(R-1)
+32\binom R3+16\binom R4.
\quad}                                             \tag{24.2}
\]
Equivalently,
\[
N_4(R)=\sum_{j=0}^{\min\{4,R\}}
        2^j\binom4j\binom Rj.                       \tag{24.3}
\]
\end{lemma}

\begin{proof}
Choose the \(j\) nonzero coordinates, choose their signs, and count
positive integer \(j\)-tuples with sum at most \(R\).  The latter
number is
\[
\sum_{m=j}^R\binom{m-1}{j-1}=\binom Rj.
\]
This proves (24.3); expanding the five terms gives (24.2).
\end{proof}

There are four oriented coarse edge directions and eight real colour
components, hence at most
\[
                         d_c=4\cdot8=32                            \tag{24.4}
\]
physical coordinate slots per coarse cell before the gauge quotient.
Using \(32\) is therefore a safe upper count after the quotient.
There are six plaquette orientation owners per cell.

\begin{definition}[Head coordinate count]
\label{def:e8v24-head-count}
Put
\[
                         X_R:=32N_4(R).                            \tag{24.5}
\]
For a degree-\(n\) coefficient tensor anchored at one translated cell,
define the safe head count
\[
                         \mathcal N_{n,R}:=6X_R^n.                 \tag{24.6}
\]
\end{definition}

\begin{lemma}[The head count is not multiplied per owner]
\label{lem:e8v24-owner-count}
After quotienting translations, every degree-\(n\) head used in V19
has at most \(\mathcal N_{n,R}\) scalar coordinates.  The factor six
in (24.6) counts the six plaquette orientation species once; it is not
raised to the \(n\)-th power.
\end{lemma}

\begin{proof}
Each of the \(n\) legs has at most \(N_4(R)\) relative cells and at
most \(32\) edge--colour slots.  This gives \(X_R^n\) ordered
coordinates.  Symmetrization only reduces the count.  The interaction
has one declared plaquette owner, with six possible orientations.
\end{proof}

\subsection{Weak-to-strong head conversion}

The unit-Wilson coarse Schur precision satisfies
\[
                         0\prec S\preceq4I                         \tag{24.7}
\]
on its physical support.  Therefore every degree-\(n\) coarse
coefficient tensor obeys
\[
                         \|T_n\|_2\leq2^n\|T_n\|_{n,S^{-1}}.       \tag{24.8}
\]

\begin{lemma}[Degreewise projected finite-head constant]
\label{lem:e8v24-degree-head}
Let \(s_2=s_3=1/4\) and \(s_4=1/16\).  If the output is actually
projected onto the \(\mathcal N_{n,R}\)-dimensional owner head, then
\[
\|\mathcal L_{\leq R,n}T_n\|_{\ell^1({\rm owner})}
\leq b_{n,R}\|T_n\|_{n,H^{-1}},                   \tag{24.9}
\]
where
\[
\boxed{\qquad
b_{n,R}:=s_n\,2^n\sqrt6\,X_R^{n/2},
\qquad n=2,3,4.
\qquad}                                            \tag{24.10}
\]
\end{lemma}

\begin{proof}
The matched conditional map is a contraction from the fine chaos
metric to the coarse chaos metric by V18.  Truncating its output to
the finite head cannot increase its Euclidean \(\ell^2\) norm.
Equation (24.8) converts the covariance norm to Euclidean norm.
Cauchy--Schwarz on the \(\mathcal N_{n,R}\) coordinates gives
\[
\|v\|_1\leq\sqrt{\mathcal N_{n,R}}\|v\|_2
=\sqrt6\,X_R^{n/2}\|v\|_2.
\]
Finally multiply by \(s_n\).
\end{proof}

\begin{proposition}[All-degree projected plaquette-owner head constant]
\label{prop:e8v24-plaquette-head}
For the direct degrees \(2,3,4\), after the finite owner-head
projection,
\[
\|\mathcal L_{\leq R}T\|_{\ell^1({\rm plaquette\ owner})}
\leq b_R^{\rm plaq}\|T\|_{\rm w},                  \tag{24.11}
\]
with the explicit constant
\[
\boxed{\quad
b_R^{\rm plaq}
=\left(\sum_{n=2}^4b_{n,R}^2\right)^{1/2}
=\sqrt6\,X_R\sqrt{X_R^2+4X_R+1}.
\quad}                                             \tag{24.12}
\]
\end{proposition}

\begin{proof}
Use Cauchy--Schwarz across the three chaos degrees.  From (24.10),
\[
b_{2,R}=\sqrt6X_R,\quad
b_{3,R}=2\sqrt6X_R^{3/2},\quad
b_{4,R}=\sqrt6X_R^2.
\]
Their squared sum is
\[
6X_R^2(1+4X_R+X_R^2).
\]
\end{proof}

\begin{theorem}[Explicit projected link-hull head constant]
\label{thm:e8v24-link-head}
On the projected head, conversion from the plaquette-owner norm to the link-marked
oscillation norm costs the sharp incidence factor six from V13,
once and only once.  Hence the V19 finite-head constant may be taken
as
\[
\boxed{\quad
b_R^{\rm link}
=6\sqrt6\,X_R\sqrt{X_R^2+4X_R+1}.
\quad}                                             \tag{24.13}
\]
\end{theorem}

\begin{proof}
A fixed fine link lies in at most six selected dyadic face hulls,
whereas every plaquette has a unique owner.  Sum the at most six
owner contributions reaching the marked link after completing the
entire tensor contraction.  The incidence is a property of the owner
hull, not of each tensor leg, so no factor \(6^n\) appears.
\end{proof}

\subsection{An explicit head radius and hybrid contraction}

Recall from V19 that
\[
j_\mu=1+8e^{2\mu}g_\mu,\qquad j_0=1+8g_0.           \tag{24.14}
\]
Set
\[
\mathfrak a_0:=
\max\left\{\frac12j_0,\frac34j_0^2,\frac14j_0^3\right\}.           \tag{24.15}
\]
After the single link-hull conversion, the remote tensor tail is
bounded by
\[
                         a_R^{\rm link}
\leq6j_\mu\mathfrak a_0e^{-\mu R}.                \tag{24.16}
\]

\begin{definition}[Concrete head radius]
\label{def:e8v24-radius}
For a fixed admissible \(\mu\), define
\[
\boxed{\quad
R_\star:=
1+\max\left\{0,
\left\lceil
\frac1\mu\log\left(\frac{48j_\mu\mathfrak a_0}{5}\right)
\right\rceil
\right\}.
\quad}                                             \tag{24.17}
\]
\end{definition}

\begin{lemma}[Remote contraction at the concrete radius]
\label{lem:e8v24-tail}
The radius (24.17) satisfies
\[
                         a_{R_\star}^{\rm link}<\frac58.           \tag{24.18}
\]
\end{lemma}

\begin{proof}
If the logarithm in (24.17) is positive, the extra integer step gives
\[
e^{-\mu R_\star}
<\frac{5}{48j_\mu\mathfrak a_0}.
\]
Insertion in (24.16) gives \(a_{R_\star}^{\rm link}<5/8\).
If the logarithm is nonpositive, then
\(6j_\mu\mathfrak a_0\leq5/8\), and the additional radius one makes
the inequality strict.
\end{proof}

\begin{theorem}[Concrete projected-head Gaussian norm]
\label{thm:e8v24-concrete-norm}
Fix \(0<\delta<1\) and put
\[
 b_\star:=b_{R_\star}^{\rm link},\qquad
\varepsilon_{\star,\delta}:=(1-\delta)\frac{3}{8b_\star},\qquad
\|T\|_{\star,\delta}:=\|T\|_{\rm w}
 +\varepsilon_{\star,\delta}\|T\|_{\rm s}.         \tag{24.19}
\]
Then the projected-head part of the stationary flat Gaussian principal
map satisfies
\[
\boxed{\qquad
\|\mathcal L_{\rm G}T\|_{\star,\delta}
\leq q_{\star,\delta}\|T\|_{\star,\delta},\qquad
q_{\star,\delta}<\frac58.
\qquad}                                            \tag{24.20}
\]
This statement concerns the finite projected head, not the unrestricted
polymer coefficient space.  All quantities are explicit functions of
\(\lambda_*,\rho_*,\mu\) and are independent of the torus volume.
\end{theorem}

\begin{proof}
The weak coefficient in the V19 Lasota--Yorke norm is
\[
\frac14+\varepsilon_{\star,\delta} b_\star
=\frac58-\frac{3\delta}{8}<\frac58.
\]
The strong coefficient is strictly below \(5/8\) by
Lemma~\ref{lem:e8v24-tail}.  Take \(q_{\star,\delta}\) to be the
maximum of these two coefficients.
\end{proof}

\subsection{Scope audit: the unrestricted finite-head lemma fails}

\begin{counterexample}[No volume-uniform \(\ell^2\)-to-\(\ell^1\)
head bound on unrestricted relative supports]
\label{cth:e8v24-l2-l1}
There is no finite constant \(b_R\), depending only on the kernel
range and local block dimensions, such that the V19 estimate
\[
\|\mathcal L_{\leq R}T\|_{\rm s}\leq b_R\|T\|_{\rm w}              \tag{24.23}
\]
holds on an unrestricted anchored coefficient family when the strong
norm contains an absolute sum over arbitrarily many relative supports
and the weak norm is Hilbertian.
\end{counterexample}

\begin{proof}
It suffices to take the range-zero identity kernel on scalar
coefficients.  Let \(T_N\) have value \(N^{-1/2}\) on \(N\) distinct
relative supports and zero elsewhere.  Then
\[
\|T_N\|_{\rm w}=1,\qquad \|T_N\|_{\rm s}=\sqrt N.
\]
Translation anchoring removes a common translate but does not identify
these \(N\) different relative supports.  Hence no volume- and
support-uniform \(b_R\) exists.
\end{proof}

\begin{correction}[Status of the V19 localized-contraction theorem]
\label{corr:e8v24-v19}
Lemma~\ref{lem:e8v19-head} and
Theorem~\ref{thm:e8v19-contraction} are valid only after a genuinely
finite local-generator projection, or after replacing their weak norm
by a stronger weighted norm that controls the relative-support sum.
They are withdrawn as statements on the unrestricted stable polymer
space.  The exact weak \(1/4\) contraction of V18 and all exponential
kernel tails of V19 remain valid.
\end{correction}

\begin{proof}
The proof of Lemma~\ref{lem:e8v19-head} invoked a finite number of
relative coordinates after truncating the displacement of the
conditional kernel.  The counterexample shows that kernel displacement
truncation does not truncate the relative support of its input.
The later Lasota--Yorke theorem used that invalid implication.  No
other V19 estimate depends on it.
\end{proof}

\begin{corollary}[Correct scope of the V24 constants]
\label{cor:e8v24-scope}
Equations (24.2)--(24.20) give rigorous constants for a projected
finite local-generator sector.  They do not prove a contraction on
arbitrary connected polymers or on an infinite relative-support
coefficient family.
\end{corollary}

\subsection{What part of \(C_{\rm own}\) is now evaluated}

\begin{definition}[Geometric owner constant]
\label{def:e8v24-Cgeom}
For the concrete radius, define
\[
\boxed{\qquad
C_{\rm geom}:=
6\sqrt6\,X_{R_\star}
\sqrt{X_{R_\star}^2+4X_{R_\star}+1}
=b_\star.
\qquad}                                            \tag{24.24}
\]
\end{definition}

\begin{assessment}[Owner geometry versus analytic interaction]
\label{ass:e8v24-separation}
The finite owner geometry for a projected local-generator head is now
explicit and may be taken as \(C_{\rm geom}=b_\star\).  It does not
control unrestricted polymer supports and does not, by itself, equal
the entire \(C_{\rm own}\) multiplying an arbitrary nonlinear Hessian
stock.  A type-specific analytic constant \(C_{\rm an}\) is still
needed for the Fr\'echet derivative of the interacting conditional
expectation.  The honest factorization is
\[
                         C_{\rm own}^{\rm local}
\leq C_{\rm geom}C_{\rm an}.                       \tag{24.25}
\]
\end{assessment}

\begin{theorem}[Updated finite-local-sector small-field gate]
\label{thm:e8v24-updated-gate}
On the finite local-generator sector only, a sufficient SCP
small-field contraction condition is
\[
\boxed{\quad
q_{\star,\delta}
+\varepsilon_{\rm stat}
+C_{\rm geom}C_{\rm an}
  \left(64e^{2\mu}gr+\zeta_{3,j}r\right)<1.
\quad}                                             \tag{24.26}
\]
This is not yet an unrestricted polymer gate.  Every displayed factor
is explicit except
\(C_{\rm an}\), \(\zeta_{3,j}\), and the stationary rechart debit
\(\varepsilon_{\rm stat}\).
\end{theorem}

\begin{proof}
Insert the concrete V24 norm and the factorization (24.25) into
Theorem~\ref{thm:e8v23-contraction}.
\end{proof}

\begin{assessment}[What V24 closes]
\label{ass:e8v24-closes}
V24 gives exact cell counts and projected-head constants for every
finite local generator, paying the six-fold link incidence once rather
than once per tensor leg.  It also finds and repairs the false
finite-range \(\ell^2\)-to-\(\ell^1\) inference in V19.
\end{assessment}

\begin{assessment}[Current bottleneck]
\label{ass:e8v24-remains}
The principal localization bottleneck is now sharper: either construct
a pair of weighted Hilbert norms with a proved interpolation
contraction and a summable embedding into the owner norm, or prove a
direct connected-polymer/cluster estimate.  Only after that repair
does it make sense to evaluate the analytic constant \(C_{\rm an}\)
for unrestricted generated interactions.
\end{assessment}

\begin{decision}[V25 acquisition order]
\label{dec:e8v24-v25}
V25 is a compulsory companion audit.  After the audit it will construct
a two-exponent Hilbert localization scale.  It will use
Cauchy--Schwarz with the exact four-dimensional exponential volume
sum to obtain a volume-uniform embedding from the stronger weighted
Hilbert norm into the weaker owner norm, and it will state the precise
weighted one-particle estimate still needed for contraction.
\end{decision}

\begin{firewall}[Claim boundary after eighth-series V24]
\label{fire:e8v24-claim}
V24 closes only the projected finite-owner geometry and corrects the
unrestricted V19 claim.  It does not prove the full stable-polymer
contraction or evaluate the interacting analytic constant,
generated-action derivatives, full nonlinear or large-field
contraction, marginal flow, rotational restoration, regulator
removal, continuum reconstruction, or the Yang--Mills mass gap.
\end{firewall}

\paragraph{Parent-version record.}
V24 was formed by exact append-only extension of
\texttt{Renewal\_Yang\_Mills\_Mass\_Gap\_Research\_Notes\_V23.tex}.
The V23 parent had \(10\,404\) counted source lines,
\(361\,480\) bytes, and SHA--256
\[
\texttt{C4E2FBA0FD178FADAD526FB7F89D432BAEC5135D2224520E07DE9147D3C48E18}.
\]

% =============================================================================
% END APPENDED EIGHTH-SERIES V24 MATERIAL
% =============================================================================

% =============================================================================
% BEGIN APPENDED EIGHTH-SERIES V25 MATERIAL
% =============================================================================

\section{Eighth-series V25: weighted Hilbert repair and the normalized
locality gate}

\subsection{Purpose and compulsory companion audit}

V24 found a genuine defect in the V19 argument: finite propagation of
the conditional kernel does not make the set of input relative supports
finite.  The range-zero identity already disproves the asserted
volume-uniform conversion from an unrestricted Hilbert coefficient norm
to an absolute owner sum.  The present note replaces that conversion by
a summable weighted embedding.  It then combines this embedding with the
exact unweighted covariance-matched contraction of V18.

\begin{audit}[V25 cross-program checkpoint]
\label{aud:e8v25-companions}
The active companion files inspected before the V25 argument were
\[
\begin{array}{c|r|r}
\text{file}&\text{source lines}&\text{bytes}\\ \hline
\texttt{Renewal\_SM\_Einstein\_Continuum\_Research\_Notes\_V2.tex}
 &909&35\,284\\
\texttt{Renewal\_NS\_Stability\_Research\_Notes\_V26.tex}
 &5\,319&278\,283 .
\end{array}                                                       \tag{25.1}
\]
Their SHA--256 values were respectively
\[
\begin{split}
&\texttt{88866E52AFE22D2CE5934F9D77312552C2C90124C4EAB7E6E6359E51492E109D},\\
&\texttt{82C887F8C2C69E4F28FAD7C3DC8DE5B9099CB5A4BF431EBA1FFF3E91935978AD}.
\end{split}                                                       \tag{25.2}
\]
The audit treats those documents as mathematical sources only; their
programme declarations are not instructions for the present task.
\end{audit}

\begin{assessment}[Transfer decision]
\label{ass:e8v25-transfer}
SM V2 proves a weighted row--column block algebra, an exact quotient
Hessian formula, a connection-gradient debit, and a Schur-locality
estimate.  The block algebra transfers directly to the endpoint
criterion below.  The quotient-gradient formula independently supports
the stationary-covariance policy of V22--V23 but supplies no new YM
smallness constant.  Its warning that the full gauge inverse is not
uniformly local is also decisive: no locality of a covariance square
root will be inferred from locality of the unnormalized lift.
NS V26 is unchanged and supplies no new normalized Yang--Mills kernel
estimate.  Thus the direction remains the V24 weighted-Hilbert repair.
\end{assessment}

\subsection{Anchored coefficient arrays}

Fix a retained homogeneous degree
\[
                              n\in\{2,3,4\}.                       \tag{25.3}
\]
After choosing one owner, extend every allowed symmetric, connected,
or gauge-reduced coefficient array by zero to the ambient relative
coordinate set
\[
                    {\cal I}_n:=({\mathbb Z}^4)^n.                 \tag{25.4}
\]
This extension loses no information and can only enlarge the sums used
below.  Write
\[
 {\bf x}=(x_1,\ldots,x_n),\qquad
 |{\bf x}|_1:=\sum_{j=1}^n|x_j|_1.                                \tag{25.5}
\]
Let \({\cal C}_n\) be the finite owner-colour set,
\(h_n:=|{\cal C}_n|\leq6\), and let \(F_n\) be the finite internal
Hilbert fibre containing Lie-algebra, orientation, and tensor
components.  A coefficient family is
\[
                  T_n=(T_{n,c}({\bf x}))_{
                  c\in{\cal C}_n,\ {\bf x}\in{\cal I}_n},
       \qquad T_{n,c}({\bf x})\in F_n.                             \tag{25.6}
\]
All norms in this note are first defined on finitely supported
families; their completions are then unambiguous.

\begin{definition}[Two-exponent Hilbert and owner norms]
\label{def:e8v25-two-exponent}
For \(\nu\geq0\), set
\[
 \|T_n\|_{2,\nu}^2
 :=\sum_{c\in{\cal C}_n}\sum_{{\bf x}\in{\cal I}_n}
       e^{2\nu|{\bf x}|_1}\|T_{n,c}({\bf x})\|_{F_n}^2.            \tag{25.7}
\]
For \(\mu\geq0\), set
\[
 \|T_n\|_{{\rm own},1,\mu}
 :=\iota_n\sum_{c\in{\cal C}_n}\sum_{{\bf x}\in{\cal I}_n}
       e^{\mu|{\bf x}|_1}\|T_{n,c}({\bf x})\|_{F_n},              \tag{25.8}
\]
where \(\iota_n=1\) for a plaquette owner and
\(\iota_n\leq6\) when a link-labelled tensor is converted to a
plaquette owner.  The incidence factor is outside the tensor-leg
product and is therefore paid once.
\end{definition}

\begin{lemma}[Exact four-dimensional exponential volume]
\label{lem:e8v25-Z4}
For every \(a>0\),
\[
\boxed{\qquad
 Z_4(a):=\sum_{x\in{\mathbb Z}^4}e^{-a|x|_1}
       =\left(\frac{1+e^{-a}}{1-e^{-a}}\right)^4<\infty .
\qquad}                                                          \tag{25.9}
\]
Consequently
\[
 \sum_{{\bf x}\in{\cal I}_n}e^{-a|{\bf x}|_1}=Z_4(a)^n.           \tag{25.10}
\]
\end{lemma}

\begin{proof}
In one dimension, absolute convergence permits the geometric-series
calculation
\[
 \sum_{m\in{\mathbb Z}}e^{-a|m|}
 =1+2\sum_{m\geq1}e^{-am}
 =\frac{1+e^{-a}}{1-e^{-a}}.
\]
The \(\ell^1\) length is additive over the four coordinates, so Tonelli's
theorem gives the fourth power.  Applying the same factorization to the
\(n\) relative legs gives (25.10).
\end{proof}

\begin{theorem}[Volume-uniform two-exponent owner embedding]
\label{thm:e8v25-embedding}
If \(0\leq\mu<\nu\), then
\[
\boxed{\qquad
 \|T_n\|_{{\rm own},1,\mu}
 \leq
 \iota_n\sqrt{h_n}\,
 Z_4\!\left(2(\nu-\mu)\right)^{n/2}
 \|T_n\|_{2,\nu}.
\qquad}                                                          \tag{25.11}
\]
The same constant works on every finite torus when relative
displacements are represented by their minimal fundamental-domain
coordinates.
\end{theorem}

\begin{proof}
Insert
\[
 e^{\mu|{\bf x}|_1}
 =e^{-(\nu-\mu)|{\bf x}|_1}e^{\nu|{\bf x}|_1}
\]
in (25.8) and apply Cauchy--Schwarz jointly to owner colour and relative
support.  The first squared factor is at most
\[
 h_n\sum_{{\bf x}\in{\cal I}_n}
       e^{-2(\nu-\mu)|{\bf x}|_1}
 =h_n Z_4\!\left(2(\nu-\mu)\right)^n
\]
by Lemma~\ref{lem:e8v25-Z4}; the second is (25.7).  A fundamental
domain on a torus is a subset of the chosen lattice representatives,
so its positive sum is no larger than the infinite-lattice sum.
\end{proof}

\begin{corollary}[Repair of the V24 counterexample]
\label{cor:e8v25-counterexample-repair}
The family in Counterexample~\ref{cth:e8v24-l2-l1} cannot have bounded
\(\|\cdot\|_{2,\nu}\) and unbounded
\(\|\cdot\|_{{\rm own},1,\mu}\) when \(\nu>\mu\).
Thus unrestricted relative supports are compatible with an owner sum
provided a strictly positive exponential reserve is retained.
\end{corollary}

\begin{proof}
Apply Theorem~\ref{thm:e8v25-embedding}.  In the particular family of
the counterexample, placing more supports at increasing distance also
charges the factor \(e^{\nu|{\bf x}|_1}\); this is precisely the charge
missing from the V19 weak norm.
\end{proof}

\begin{firewall}[The strict exponent gap is necessary for this proof]
\label{fire:e8v25-gap}
At \(\mu=\nu\), the factor in (25.11) becomes
\(\sum_{{\bf x}}1=\infty\).  Hence plain exponentially weighted
\(\ell^2\) does not embed into the owner \(\ell^1\) norm at the same
exponent.  V25 will not silently set \(\mu=\nu\).
\end{firewall}

\subsection{A same-exponent polynomial reserve}

For later nonlinear products it is useful to record a second, optional
repair.  Put \(D_n:=4n\), and for \(s>D_n/2=2n\) define
\[
 P_{D_n,s}:=\sum_{{\bf x}\in{\cal I}_n}
                 (1+|{\bf x}|_1)^{-2s}.                           \tag{25.12}
\]

\begin{lemma}[Polynomial summability]
\label{lem:e8v25-polynomial-sum}
For \(s>2n\), \(P_{D_n,s}<\infty\).  In particular, if
\[
 \|T_n\|_{2,\mu,s}^2
 :=\sum_{c,{\bf x}}e^{2\mu|{\bf x}|_1}
       (1+|{\bf x}|_1)^{2s}\|T_{n,c}({\bf x})\|_{F_n}^2,          \tag{25.13}
\]
then
\[
 \|T_n\|_{{\rm own},1,\mu}
 \leq\iota_n\sqrt{h_nP_{D_n,s}}\,
                         \|T_n\|_{2,\mu,s}.                       \tag{25.14}
\]
\end{lemma}

\begin{proof}
The number of lattice points in an \(\ell^1\) shell of radius \(m\)
in dimension \(D_n\) is \(O((m+1)^{D_n-1})\).  More explicitly, for
\(m\geq1\) it equals
\[
 \sum_{j=1}^{\min(D_n,m)}
       2^j{D_n\choose j}{m-1\choose j-1}.                         \tag{25.15}
\]
Therefore the series in (25.12) is dominated by a constant times
\(\sum_{m\geq0}(m+1)^{D_n-1-2s}\), which converges exactly under the
stated inequality.  Cauchy--Schwarz applied to
\[
 e^{\mu|{\bf x}|_1}\|T_{n,c}({\bf x})\|
 =(1+|{\bf x}|_1)^{-s}
   e^{\mu|{\bf x}|_1}(1+|{\bf x}|_1)^s
     \|T_{n,c}({\bf x})\|
\]
then gives (25.14).
\end{proof}

\begin{lemma}[Hilbert--Sobolev convolution algebra]
\label{lem:e8v25-convolution}
Let \(s>2n\), and let convolution be taken on
\({\cal I}_n\), with any finite-dimensional fibre contraction of norm
at most one.  Then
\[
 \|f*g\|_{2,\mu,s}
 \leq 2^sP_{D_n,s}^{1/2}
       \|f\|_{2,\mu,s}\|g\|_{2,\mu,s}.                            \tag{25.16}
\]
\end{lemma}

\begin{proof}
For \({\bf x}={\bf y}+{\bf z}\), the triangle inequality and
\((A+B)^s\leq2^{s-1}(A^s+B^s)\), \(A,B\geq0\), give
\[
 (1+|{\bf x}|_1)^s
 \leq2^{s-1}\bigl((1+|{\bf y}|_1)^s+
                  (1+|{\bf z}|_1)^s\bigr).
\]
The exponential weight is submultiplicative.  Hence the weighted
absolute value of \(f*g\) is bounded by \(2^{s-1}\) times the sum of
the convolution of the fully weighted \(f\) with the exponentially
weighted \(g\), and the symmetric term.  Young's
\(\ell^2*\ell^1\to\ell^2\) inequality applies.  Lemma
\ref{lem:e8v25-polynomial-sum} without the owner factors gives
\[
 \|e^{\mu|\cdot|_1}g\|_1
 \leq P_{D_n,s}^{1/2}\|g\|_{2,\mu,s},
\]
and similarly for \(f\).  The two terms produce (25.16).
\end{proof}

\begin{assessment}[What the two reserves accomplish]
\label{ass:e8v25-reserves}
The exponential gap \(\nu-\mu\) is the sharp simple mechanism for
passing from Hilbert contraction to an absolute owner sum.  The
polynomial reserve provides a same-exponential-weight algebra for
nonlinear convolutions.  Neither result asserts that the
covariance-normalized one-particle map preserves the corresponding
weighted Hilbert space; that is the separate analytic gate treated
next.
\end{assessment}

\subsection{The normalized one-particle endpoint}

The weights above must be imposed in a local realization of the
covariance-matched one-particle Hilbert spaces.  This clause matters:
writing
\[
                     R^*=S^{-1/2}J^*H^{1/2}                      \tag{25.17}
\]
proves \(\|R^*\|\leq1\), but the three factors on the right need not
separately have local kernels.  We therefore isolate exactly the
additional statement which is required.

\begin{definition}[Normalized block-realization gate]
\label{def:e8v25-NBR}
NBR\((\alpha,K_\alpha)\), with \(\alpha>0\), means that at every
regulator and scale:
\begin{enumerate}
\item the fine and coarse matched one-particle spaces admit orthogonal
local block decompositions over their cells, after pure gauge and
declared retained zero modes have been removed;
\item in those decompositions \(R^*\) is a block operator between the
fine and coarse cell spaces and has unweighted norm at most one;
\item if \(L_{\rm f}\) and \(L_{\rm c}\) multiply a block by its
nonnegative distance from the declared owner, then
\[
 \left\|
 e^{\alpha L_{\rm c}}R^*e^{-\alpha L_{\rm f}}
 \right\|_{2\to2}\leq K_\alpha                                  \tag{25.18}
\]
with \(K_\alpha<\infty\) independent of volume and scale.
\end{enumerate}
\end{definition}

\begin{remark}[This is weaker than the old LSG smallness demand]
\label{rem:e8v25-weaker}
Definition~\ref{def:e8v18-lsg} asked for a positive-weight
row--column bound sufficiently close to one to preserve contraction
after tensorization.  NBR asks only that one positive-weight endpoint
be finite.  The exact unweighted contraction will supply the small
positive-weight estimate by interpolation.
\end{remark}

\begin{lemma}[Weighted Schur sufficient condition]
\label{lem:e8v25-Schur}
Suppose \(R^*=(R^*_{xy})\) in local orthogonal block coordinates.
If
\[
\begin{split}
 M_\alpha:=\max\biggl\{&
 \sup_x\sum_y
 e^{\alpha(L_{\rm c}(x)-L_{\rm f}(y))}
       \|R^*_{xy}\|,\\
 &\sup_y\sum_x
 e^{\alpha(L_{\rm c}(x)-L_{\rm f}(y))}
       \|R^*_{xy}\|\biggr\}<\infty,                               \tag{25.19}
\end{split}
\]
uniformly in the regulator, then (25.18) holds with
\(K_\alpha=M_\alpha\).  A stronger symmetric bound with
\(e^{\alpha|L_{\rm c}(x)-L_{\rm f}(y)|}\) is also sufficient.
\end{lemma}

\begin{proof}
The blocks of the conjugated operator in (25.18) are exactly
\[
 e^{\alpha L_{\rm c}(x)}R^*_{xy}e^{-\alpha L_{\rm f}(y)}.
\]
The two quantities in (25.19) are its absolute block row and column
sums.  The block Schur test bounds its Hilbert operator norm by their
geometric mean and hence by their maximum.
\end{proof}

\begin{firewall}[Unwhitened locality does not prove NBR]
\label{fire:e8v25-unwhitened}
The exponential kernel estimates for the vertical inverse \(A^{-1}\)
and the unwhitened lift \(J=-A^{-1}B\) from V19--V21 do not imply
(25.18).  The physical factors \(H^{1/2}\) and \(S^{-1/2}\) in
(25.17) contain closing massless directions.  A uniform lower
Euclidean spectral bound was neither proved nor expected there.
NBR must be proved for the normalized composite, not obtained by
multiplying separate Euclidean bounds.
\end{firewall}

\subsection{Interpolation from the exact zero-weight margin}

\begin{theorem}[Normalized weighted interpolation]
\label{thm:e8v25-interpolation}
Assume NBR\((\alpha,K_\alpha)\), and put
\[
                         \overline K_\alpha:=\max\{1,K_\alpha\}.
                                                                    \tag{25.20}
\]
For \(0\leq\theta\leq1\),
\[
\boxed{\qquad
 \left\|
 e^{\theta\alpha L_{\rm c}}R^*
 e^{-\theta\alpha L_{\rm f}}
 \right\|_{2\to2}
 \leq\overline K_\alpha^\theta .
\qquad}                                                            \tag{25.21}
\]
\end{theorem}

\begin{proof}
Work first at a finite regulator and define on the strip
\(0\leq\Re z\leq1\)
\[
 F(z):=e^{z\alpha L_{\rm c}}R^*e^{-z\alpha L_{\rm f}}.
\]
For arbitrary unit vectors \(u,v\), the scalar function
\(\langle v,F(z)u\rangle\) is analytic.  On \(\Re z=0\), the
imaginary exponential weights are unitary, so the boundary norm is at
most \(\|R^*\|\leq1\).  On \(\Re z=1\), commute the real and imaginary
functions of each \(L\) to write \(F(1+it)\) as the endpoint operator
in (25.18) multiplied on the left and right by unitaries.  Its norm is
at most \(K_\alpha\).  Hadamard's three-lines theorem yields
\[
 |\langle v,F(\theta)u\rangle|
 \leq\overline K_\alpha^\theta.
\]
Taking the supremum over \(u,v\) proves (25.21).  The estimate is
uniform in the regulator, so it passes to the Hilbert completions.
\end{proof}

\begin{definition}[Explicit localization exponent]
\label{def:e8v25-beta-star}
If \(\overline K_\alpha=1\), set \(\theta_*=1/2\).  Otherwise set
\[
 \theta_*:=
 \min\left\{\frac12,
      \frac{\log2}{3\log\overline K_\alpha}\right\},
 \qquad \beta_*:=\theta_*\alpha.                                  \tag{25.22}
\]
\end{definition}

\begin{lemma}[Explicit tensorized margin]
\label{lem:e8v25-tensor-margin}
Under NBR\((\alpha,K_\alpha)\), let
\[
 \kappa_*:=
 \left\|e^{\beta_*L_{\rm c}}R^*
              e^{-\beta_*L_{\rm f}}\right\| .
\]
Then
\[
                         \kappa_*\leq2^{1/3},                     \tag{25.23}
\]
and, for the canonical factors
\[
                 s_2=s_3=\frac14,\qquad s_4=\frac1{16},
\]
\[
\boxed{\qquad
 q_*:=\max_{n=2,3,4}s_n\kappa_*^n\leq\frac12 .
\qquad}                                                           \tag{25.24}
\]
\end{lemma}

\begin{proof}
If \(\overline K_\alpha=1\), (25.21) gives
\(\kappa_*\leq1\).  Otherwise (25.22) implies
\[
 \theta_*\log\overline K_\alpha\leq\frac{\log2}{3},
\]
which proves (25.23).  Therefore
\[
 \frac14\kappa_*^2\leq2^{-4/3},\qquad
 \frac14\kappa_*^3\leq\frac12,\qquad
 \frac1{16}\kappa_*^4\leq2^{-8/3}.
\]
The middle degree is the largest of these upper bounds.
\end{proof}

\subsection{Conditional unrestricted principal contraction}

Use the local matched block realization in
Definition~\ref{def:e8v25-NBR}.  On the \(n\)-particle coefficient
space the product weight is
\[
 e^{\beta L^{(n)}}:=
 \bigl(e^{\beta L}\bigr)^{\otimes n},\qquad
 L^{(n)}(x_1,\ldots,x_n)=\sum_{j=1}^nL(x_j).                      \tag{25.25}
\]
Let
\[
 \|T\|_{\mathscr H_{\beta}}^2
 :=\sum_{n=2}^4 n!\,\|T_n\|_{2,\beta}^2,                          \tag{25.26}
\]
where the finite internal and owner fibres are included in the
Hilbert sum.

\begin{theorem}[NBR implies an unrestricted weighted principal
contraction]
\label{thm:e8v25-principal}
Assume NBR\((\alpha,K_\alpha)\).  For the exact flat Gaussian
principal derivative
\[
 (\mathcal L_{\rm G}T)_n
 =s_n\,{\rm Sym}\bigl((R^*)^{\otimes n}T_n\bigr),                 \tag{25.27}
\]
one has
\[
\boxed{\qquad
       \|\mathcal L_{\rm G}T\|_{\mathscr H_{\beta_*}}
       \leq\frac12\|T\|_{\mathscr H_{\beta_*}} .
\qquad}                                                           \tag{25.28}
\]
This statement permits arbitrary relative supports; no finite-head
projection occurs.
\end{theorem}

\begin{proof}
Conjugating (25.27) by the product weights in (25.25) gives the
symmetric restriction of the \(n\)-fold tensor power of the
one-particle operator in Lemma~\ref{lem:e8v25-tensor-margin}.
Tensor-product norms multiply, and orthogonal symmetrization cannot
increase the norm.  Thus the degree-\(n\) block has norm at most
\(s_n\kappa_*^n\leq q_*\).  The factorials in (25.26) occur on both
sides of each degree block.  Orthogonality of the three degrees and
(25.24) prove (25.28).
\end{proof}

\begin{corollary}[Unrestricted owner summability at a weaker exponent]
\label{cor:e8v25-owner-output}
For every \(0\leq\mu<\beta_*\), define
\[
 C_{\rm emb}(\mu,\beta_*):=
 \left[
 \sum_{n=2}^4
 \frac{\iota_n^2h_n}{n!}
 Z_4\!\left(2(\beta_*-\mu)\right)^n
 \right]^{1/2}.                                                   \tag{25.29}
\]
Then
\[
\boxed{\qquad
 \sum_{n=2}^4
 \|(\mathcal L_{\rm G}T)_n\|_{{\rm own},1,\mu}
 \leq\frac12 C_{\rm emb}(\mu,\beta_*)
                  \|T\|_{\mathscr H_{\beta_*}} .
\qquad}                                                           \tag{25.30}
\]
Every constant is independent of the volume.
\end{corollary}

\begin{proof}
Apply Theorem~\ref{thm:e8v25-embedding} to each output degree and then
Cauchy--Schwarz to the three-degree sum, pairing its embedding
constant with \(n!^{-1/2}\).  This gives
\[
 \sum_n\|(\mathcal L_{\rm G}T)_n\|_{{\rm own},1,\mu}
 \leq C_{\rm emb}(\mu,\beta_*)
       \|\mathcal L_{\rm G}T\|_{\mathscr H_{\beta_*}}.
\]
Use Theorem~\ref{thm:e8v25-principal}.
\end{proof}

\begin{correction}[Replacement for the withdrawn V19 theorem]
\label{corr:e8v25-replacement}
The correct unrestricted implication is conditional
Theorem~\ref{thm:e8v25-principal} together with
Corollary~\ref{cor:e8v25-owner-output}.  Kernel range never appears as
a fictitious bound on the number of input relative supports.  The
support sum is instead paid by the exact summability factor
\(Z_4(2(\beta_*-\mu))^{n/2}\).
\end{correction}

\begin{firewall}[Owner summability is not owner contraction]
\label{fire:e8v25-owner-not-contraction}
The constant \(C_{\rm emb}\) in (25.30) can be larger than two.
Accordingly, (25.30) is a continuous embedding into the owner norm,
not a claim that the owner norm itself contracts.  The invariant
stable space proposed here is the stronger weighted Hilbert space;
the owner norm is used to make cluster and incidence sums absolutely
convergent.  Nonlinear closure must be proved in that stronger space,
using either the polynomial reserve of
Lemma~\ref{lem:e8v25-convolution} or a graded exponent scale.
\end{firewall}

\subsection{A zero-momentum diagnostic for the endpoint}

\begin{proposition}[Analytic-symbol test for the Schur route]
\label{prop:e8v25-symbol-test}
Let a translation-invariant block convolution \(R^*\) on
\({\mathbb Z}^4\) have kernel \(r(x)\).  If, for some \(\alpha>0\),
\[
                    \sum_{x\in{\mathbb Z}^4}
                    e^{\alpha|x|_1}\|r(x)\|<\infty,               \tag{25.31}
\]
then its Fourier symbol
\[
                    \widehat r(k)=\sum_xr(x)e^{-ik\cdot x}        \tag{25.32}
\]
extends to a bounded analytic matrix function on every closed
substrip \(|\Im k_j|\leq\alpha'<\alpha\).  In particular it has a
unique continuous limit at \(k=0\).  Hence a direction-dependent
zero-momentum limit of the normalized symbol rules out the
exponential Schur criterion (25.31).
\end{proposition}

\begin{proof}
For \(|\Im k_j|\leq\alpha'<\alpha\),
\[
 \|r(x)e^{-ik\cdot x}\|
 \leq e^{\alpha'|x|_1}\|r(x)\|.
\]
The right-hand side is summable by (25.31).  The Weierstrass test gives
uniform convergence on compact substrips, and termwise complex
differentiation is justified after reducing the strip once more,
because every polynomial factor in \(|x|\) is absorbed by
\(e^{(\alpha-\alpha')|x|_1}\).  Thus the sum is analytic and, in
particular, continuous at the origin.
\end{proof}

\begin{assessment}[What V25 closes]
\label{ass:e8v25-closes}
V25 repairs the precise logical defect found in V24.  It proves a
volume-uniform embedding from a stronger exponentially weighted
Hilbert coefficient space into the unrestricted owner sum, supplies a
same-exponent Hilbert--Sobolev convolution algebra, and shows that any
finite normalized exponential endpoint yields an explicit principal
contraction \(q_*\leq1/2\).  This improves the old LSG requirement:
the endpoint need not be close to one.
\end{assessment}

\begin{assessment}[Current bottleneck]
\label{ass:e8v25-remains}
The sole missing hypothesis in the repaired flat-principal theorem is
NBR: a scale-uniform local block realization of the normalized
partial isometry \(R^*\) with one finite positive-weight endpoint.
The known gap controls the vertical inverse and unwhitened lift but
does not control the massless covariance normalization.  Before
estimating nonlinear owner constants, the zero-momentum symbol of the
normalized composite must be tested for a continuous analytic
extension.  A direction-dependent Hodge or Riesz limit would force a
multiscale, rather than exponentially local, replacement.
\end{assessment}

\begin{decision}[V26 upstream acquisition order]
\label{dec:e8v25-v26}
V26 will go upstream to the flat Fourier complex.  It will compute the
normalized one-particle symbol on the physical quotient, separate
retained harmonic and gauge modes before taking \(k\to0\), and test
the analytic criterion of
Proposition~\ref{prop:e8v25-symbol-test}.  If the Ward factors cancel
the apparent square-root singularities, it will derive NBR by analytic
continuation and Fourier coefficient decay.  If they do not, it will
replace NBR by a dyadic annular or wavelet localization adapted to the
order-zero Hodge multiplier, rather than reusing a circular
exponential-locality assumption.
\end{decision}

\begin{firewall}[Claim boundary after eighth-series V25]
\label{fire:e8v25-claim}
V25 does not prove NBR, the unrestricted nonlinear RG contraction,
the generated-action derivative bounds in the repaired space,
large-field stability, the marginal flow, regulator-uniform
Osterwalder--Schrader reconstruction, the continuum limit, or the
Yang--Mills mass gap.  Theorems
\ref{thm:e8v25-principal} and
\ref{cor:e8v25-owner-output} are explicitly conditional on NBR.
\end{firewall}

\paragraph{Parent-version record.}
V25 was formed by exact append-only extension of
\texttt{Renewal\_Yang\_Mills\_Mass\_Gap\_Research\_Notes\_V24.tex}.
The V24 parent had \(10\,824\) validator-counted source lines,
\(375\,497\) bytes, and SHA--256
\[
\texttt{C84E64519EB1E308FBA989F4500B5837A05ADDE76D695D8DA9CBE6680EE5EE5F}.
\]

% =============================================================================
% END APPENDED EIGHTH-SERIES V25 MATERIAL
% =============================================================================

% =============================================================================
% BEGIN APPENDED EIGHTH-SERIES V26 MATERIAL
% =============================================================================

\section{Eighth-series V26: the Hodge obstruction and annular
localization}

\subsection{Why the normalized symbol must be tested first}

V25 reduced the repaired principal contraction to one positive-weight
endpoint for the normalized coefficient map.  This note performs the
promised upstream test.  The result is not a technical failure of an
estimate: an exponentially local orthonormal curvature chart is
obstructed by the rank-changing Hodge bundle at zero momentum.  The
appropriate replacement is scale-adapted annular localization.

\subsection{The exact coarse Hodge projection}

For \(k\in[-\pi,\pi]^4\), write
\[
 p_\mu(k)=e^{ik_\mu}-1,\qquad
 r(k)=\sum_{\mu=1}^4|p_\mu(k)|^2.                                \tag{26.1}
\]
Let \(D_p:\Lambda^1{\mathbb C}^4\to\Lambda^2{\mathbb C}^4\) be
\[
                              D_p h=p\wedge h.                    \tag{26.2}
\]
Up to harmless diagonal Fourier phases, \(D_p\) is the coarse lattice
curl symbol \(d_c(k)\).  Such phases do not change any conclusion
about rank, continuity, or orthogonal projections.

\begin{lemma}[Exact two-form Hodge projector]
\label{lem:e8v26-Hodge-projector}
For \(k\ne0\), the orthogonal projection onto the realizable coarse
curvature space
\[
                 {\cal R}_c(k)=\operatorname{Ran}D_p
\]
is
\[
\boxed{\qquad
 \Pi_c(k)=\frac1{r(k)}D_pD_p^* .
\qquad}                                                           \tag{26.3}
\]
It has rank three.
\end{lemma}

\begin{proof}
Exterior multiplication by \(p\) has adjoint interior multiplication
by \(\overline p\), and hence
\[
 D_p^*D_ph=r(k)h-p\,\langle p,h\rangle .                          \tag{26.3a}
\]
Thus \(D_p^*D_p=r(k)I\) on \(p^\perp\), while
\(\ker D_p=\operatorname{span}\{p\}\).  Consequently
\[
 \frac1rD_pD_p^*
\]
is the identity on \(\operatorname{Ran}D_p\) and zero on its
orthogonal complement.  Since \(D_p\) has a one-dimensional kernel
in a four-dimensional domain, its range has dimension three.
\end{proof}

\begin{theorem}[Directional discontinuity at zero momentum]
\label{thm:e8v26-Hodge-discontinuity}
The map \(k\mapsto\Pi_c(k)\), \(k\ne0\), has no operator-norm
continuous extension to \(k=0\).
\end{theorem}

\begin{proof}
Let \(e_1,\ldots,e_4\) be the standard one-form basis and set
\(\omega=e_1\wedge e_3\).  Along \(k=(t,0,0,0)\), \(t\ne0\), the
vector \(p(k)\) is a nonzero multiple of \(e_1\), so
\[
                  \Pi_c(t,0,0,0)\omega=\omega.                   \tag{26.4}
\]
Along \(k=(0,t,0,0)\), \(p(k)\) is a nonzero multiple of \(e_2\).
The form \(\omega\) is orthogonal to
\[
 \operatorname{Ran}(e_2\wedge\cdot)
 =\operatorname{span}\{
 e_2\wedge e_1,e_2\wedge e_3,e_2\wedge e_4\},
\]
and therefore
\[
                  \Pi_c(0,t,0,0)\omega=0.                         \tag{26.5}
\]
The two limits disagree on a unit vector.
\end{proof}

\subsection{The normalized curvature isometry inherits the obstruction}

Retain the V16 notation, including the fine alias projection \(P_f(k)\),
the uniform tile isometry \(U(k)\), and
\[
                         A(k):=U(k)^*P_f(k)U(k).                  \tag{26.6}
\]
On \({\cal R}_c(k)\), the compression of \(A(k)\) is \(G(k)\).
Define the positive ambient extension
\[
 \Gamma(k):=\Pi_c(k)A(k)\Pi_c(k)+I-\Pi_c(k),                      \tag{26.7}
\]
whenever \(G(k)>0\), and set
\[
 {\cal W}(k):=
 P_f(k)U(k)\Gamma(k)^{-1/2}\Pi_c(k).                              \tag{26.8}
\]
This maps the ambient coarse plaquette fibre to the direct sum of the
sixteen fine-alias plaquette fibres.

\begin{lemma}[Initial projection of the normalized isometry]
\label{lem:e8v26-initial-projection}
For every nonzero momentum,
\[
\boxed{\qquad
                  {\cal W}(k)^*{\cal W}(k)=\Pi_c(k).
\qquad}                                                           \tag{26.9}
\]
On \({\cal R}_c(k)\), \({\cal W}=P_fUG^{-1/2}\), the normalized
curvature isometry identified in V18.  Its adjoint is the normalized
coefficient map.
\end{lemma}

\begin{proof}
The operator \(\Gamma\) is block diagonal for
\(\operatorname{Ran}\Pi_c\oplus\ker\Pi_c\); on the first summand it is
the compression \(G\), and on the second it is the identity.  Hence
\[
\begin{split}
 {\cal W}^*{\cal W}
 &=\Pi_c\Gamma^{-1/2}U^*P_fU\Gamma^{-1/2}\Pi_c\\
 &=\Pi_cG^{-1/2}GG^{-1/2}\Pi_c=\Pi_c.
\end{split}
\]
The restriction statement is immediate from (26.7)--(26.8).
\end{proof}

\begin{theorem}[No exponentially local orthonormal physical-curvature
chart]
\label{thm:e8v26-no-local-frame}
There is no translation-covariant orthonormal block chart for
\({\cal R}_c(k)\), of fixed rank three for \(k\ne0\), whose synthesis
kernel is absolutely summable with a positive exponential weight
uniformly down to zero momentum.  In particular neither
\({\cal W}(k)\) nor its adjoint has such an ambient plaquette kernel.
\end{theorem}

\begin{proof}
Suppose first that a chart existed.  Its Fourier synthesis matrix
\(V(k):{\mathbb C}^3\to\Lambda^2{\mathbb C}^4\) would satisfy
\[
                       V(k)^*V(k)=I,\qquad
                       V(k)V(k)^*=\Pi_c(k).                       \tag{26.10}
\]
An exponentially absolutely summable kernel has a Fourier symbol
continuous at \(k=0\), by
Proposition~\ref{prop:e8v25-symbol-test}.  Thus \(V(k)V(k)^*\) would
extend continuously, contradicting
Theorem~\ref{thm:e8v26-Hodge-discontinuity}.

If \({\cal W}\) or \({\cal W}^*\) had such a kernel, its Fourier symbol
would again extend continuously.  Taking the adjoint if necessary and
using (26.9) would give a continuous extension of \(\Pi_c\), the same
contradiction.
\end{proof}

\begin{correction}[Status of the V25 NBR route]
\label{corr:e8v26-NBR}
The conditional implication in
Theorem~\ref{thm:e8v25-principal} is valid.  Its natural
translation-covariant orthonormal realization on physical curvature,
however, cannot satisfy NBR uniformly through zero momentum.
Consequently NBR is not adopted as the global flat-principal
localization hypothesis.  It remains usable on carriers separated
from \(k=0\).
\end{correction}

\subsection{A uniform angular Gram floor from the edge Schur form}

The failure of continuity is directional, not a blow-up of the
normalized map.  This distinction follows from the V15 vertical gap
and the V17 edge-to-curvature bridge.

\begin{lemma}[Quadratic vanishing of the edge Schur symbol]
\label{lem:e8v26-Schur-quadratic}
Let \(\mathsf S^{\rm tree}(k)\) be the flat edge Schur symbol of V17
in a centred parity-cell Fourier convention.  It is twice
continuously differentiable on the Brillouin cube, positive
semidefinite for real \(k\), and
\[
                \mathsf S^{\rm tree}(0)=0,\qquad
                D\mathsf S^{\rm tree}(0)=0.                       \tag{26.11}
\]
If
\[
 C_2:=\sup_{\substack{k\in[-\pi,\pi]^4\\|u|=|v|=1}}
 \left\|D^2\mathsf S^{\rm tree}(k)[u,v]\right\|,                  \tag{26.12}
\]
then
\[
\boxed{\qquad
 \|\mathsf S^{\rm tree}(k)\|
 \leq C_S r(k),\qquad C_S:=\frac{\pi^2C_2}{8}.
\qquad}                                                           \tag{26.13}
\]
\end{lemma}

\begin{proof}
The unreduced Wilson blocks are finite Laurent-polynomial matrices.
The vertical block has the uniform positive floor
\(\lambda_*/3\) from V15, so its inverse and the Schur complement are
twice continuously differentiable.  At \(k=0\), every constant coarse
edge field has a constant fine-edge completion with zero curl; the
variational definition of the Schur form gives
\(\mathsf S^{\rm tree}(0)=0\).

For a fixed vector \(h\), the real differentiable function
\[
 f_h(k)=\langle h,\mathsf S^{\rm tree}(k)h\rangle
\]
is nonnegative and vanishes at the origin.  It therefore has zero
first derivative there.  Polarization gives
\(D\mathsf S^{\rm tree}(0)=0\).  Taylor's formula along the segment
from \(0\) to \(k\) yields
\[
 \|\mathsf S^{\rm tree}(k)\|\leq\frac{C_2}{2}|k|_2^2.
\]
For \(|k_\mu|\leq\pi\),
\[
 |e^{ik_\mu}-1|=2|\sin(k_\mu/2)|
 \geq\frac{2}{\pi}|k_\mu|,
\]
so \(r(k)\geq4|k|_2^2/\pi^2\).  This proves (26.13).
\end{proof}

\begin{theorem}[Uniform nonzero-momentum Gram floor]
\label{thm:e8v26-Gram-floor}
On \({\cal R}_c(k)\), \(k\ne0\),
\[
\boxed{\qquad
                  G(k)\succeq\gamma_G I,\qquad
                  \gamma_G:=\frac1{4C_S}>0.
\qquad}                                                           \tag{26.14}
\]
Thus the normalized isometry is uniformly bounded; its zero-momentum
problem is the angular discontinuity of its physical subspace, not a
divergent singular value.
\end{theorem}

\begin{proof}
For \(y\in{\cal R}_c(k)\), the V17 bridge gives
\[
 \langle y,G(k)^{-1}y\rangle
 =\frac4{r(k)^2}
 \langle d_c(k)^*y,\mathsf S^{\rm tree}(k)d_c(k)^*y\rangle.
\]
On the exact two-form range,
\[
                  \|d_c(k)^*y\|^2=r(k)\|y\|^2.
\]
Use (26.13) to obtain
\[
 \langle y,G^{-1}y\rangle
 \leq\frac4{r^2}\,C_Sr\,r\|y\|^2
 =4C_S\|y\|^2.
\]
Inverting the positive inequality proves (26.14).
\end{proof}

\begin{remark}[Computability of the floor]
\label{rem:e8v26-computable-floor}
The constant \(C_2\) is the supremum of derivatives of a fixed finite
Schur matrix whose vertical denominator has the certified algebraic
gap \(\lambda_*/3\).  It is therefore finite and amenable to interval
certification.  V26 uses only finiteness; it does not report an
uncertified decimal value.
\end{remark}

\subsection{Order-zero regularity on the punctured Brillouin zone}

Put
\[
                              \rho(k):=r(k)^{1/2}.                 \tag{26.15}
\]
For a multi-index \(\eta\), write \(|\eta|\) for its order.

\begin{lemma}[Hodge Mikhlin bounds]
\label{lem:e8v26-Hodge-Mikhlin}
For every multi-index \(\eta\), there is a finite constant
\(C_\eta^{\rm H}\), depending only on \(\eta\) and the dimension, such
that
\[
\boxed{\qquad
 \|\partial_k^\eta\Pi_c(k)\|
 \leq C_\eta^{\rm H}\rho(k)^{-|\eta|}
 \qquad(k\ne0).
\qquad}                                                           \tag{26.16}
\]
The same estimate holds for the low fine alias Hodge projection with
argument \(k/2\).  Every nonzero fine alias projection is smooth with
uniformly bounded derivatives near \(k=0\).
\end{lemma}

\begin{proof}
Formula (26.3) is a quotient of a matrix whose entries are quadratic
in \(p,\overline p\) by \(r=|p|^2\).  The analytic functions
\(p_\mu(k)\) have a simple zero at the origin and bounded derivatives.
Repeated quotient differentiation shows that a derivative of order
\(|\eta|\) is a sum of terms bounded by a constant times
\(\rho^{-|\eta|}\).  Away from a fixed neighbourhood of the origin,
\(\rho\) has a positive lower bound and compactness supplies the same
form of estimate after increasing the constant.

For the low alias, \(\rho(k/2)\) is comparable to \(\rho(k)\) near
zero.  A nonzero alias contains at least one momentum component near
\(\pi\), so its fine lattice Laplacian denominator is bounded away
from zero there; its projection is therefore smooth.
\end{proof}

\begin{theorem}[The normalized map is an order-zero multiplier]
\label{thm:e8v26-W-Mikhlin}
For every multi-index \(\eta\), the normalized isometry in (26.8)
satisfies
\[
\boxed{\qquad
 \|\partial_k^\eta{\cal W}(k)\|
 +\|\partial_k^\eta{\cal W}(k)^*\|
 \leq C_\eta^{\cal W}\rho(k)^{-|\eta|}
 \qquad(k\ne0),
\qquad}                                                           \tag{26.17}
\]
for a finite scale- and volume-independent constant
\(C_\eta^{\cal W}\).
\end{theorem}

\begin{proof}
The uniform tile matrix \(U(k)\) is a finite Laurent-polynomial
matrix and has bounded derivatives.  Lemma
\ref{lem:e8v26-Hodge-Mikhlin}, applied to all fine aliases, gives the
Mikhlin bounds for \(P_f(k)\), and hence for
\(A(k)=U^*P_fU\).  Products preserve bounds of the form
\(\rho^{-|\eta|}\), so (26.7) gives them for \(\Gamma(k)\).

By Theorem~\ref{thm:e8v26-Gram-floor},
\[
                  \Gamma(k)\succeq
                  \min\{1,\gamma_G\}I.                            \tag{26.18}
\]
Use the positive-matrix formula
\[
 \Gamma^{-1/2}
 =\frac1\pi\int_0^\infty
       t^{-1/2}(\Gamma+tI)^{-1}\,dt.                              \tag{26.19}
\]
Differentiating a resolvent inserts factors
\[
 -(\Gamma+tI)^{-1}(\partial\Gamma)(\Gamma+tI)^{-1}.
\]
Induction on \(|\eta|\), (26.18), and the already proved derivative
bounds for \(\Gamma\) show that every differentiated integrand is
bounded by an integrable function of \(t\) times
\(\rho^{-|\eta|}\).  Thus \(\Gamma^{-1/2}\) has the same Mikhlin
bounds.  Apply the product rule to (26.8); taking adjoints proves the
second bound.
\end{proof}

\begin{remark}[Nature of the singularity]
\label{rem:e8v26-order-zero}
Equation (26.17) is the standard regularity of an order-zero Hodge or
Riesz multiplier.  It permits angular dependence at \(k=0\), and is
therefore consistent with
Theorem~\ref{thm:e8v26-Hodge-discontinuity}.  It predicts rapid
localization after a momentum annulus is fixed, but not a
volume-uniform exponential kernel for the unsplit operator.
\end{remark}

\subsection{Dyadic annuli and their spatial kernels}

Choose a nonnegative function
\(\chi\in C_c^\infty((1/2,2))\) so that, after the addition of finitely
many ultraviolet pieces,
\[
                  \sum_{j\geq j_0}\chi(2^j\rho)^2=1              \tag{26.20}
\]
on a punctured neighbourhood of \(\rho=0\).  Write
\[
                  \chi_j(k)=\chi(2^j\rho(k)),\qquad
                  M_j(k)=\chi_j(k){\cal W}(k)^*.                 \tag{26.21}
\]
The zero momentum, including harmonic torons, is retained separately
and does not enter (26.20).

\begin{theorem}[Scale-adapted annular kernel decay]
\label{thm:e8v26-annular-kernel}
Let
\[
 K_j(x):=\frac1{(2\pi)^4}\int_{[-\pi,\pi]^4}
             e^{ik\cdot x}M_j(k)\,dk,\qquad x\in{\mathbb Z}^4.   \tag{26.22}
\]
For every integer \(N\geq0\), there is a finite constant \(A_N\),
independent of \(j\), volume, and scale, such that
\[
\boxed{\qquad
 \|K_j(x)\|
 \leq A_N\,2^{-4j}
       \left(1+2^{-j}|x|_1\right)^{-N}.
\qquad}                                                           \tag{26.23}
\]
The finitely many ultraviolet pieces obey the analogous estimate at
unit spatial scale.
\end{theorem}

\begin{proof}
On the support of \(\chi_j\),
\[
                         \rho(k)\asymp2^{-j}.                     \tag{26.24}
\]
The chain rule, Lemma~\ref{lem:e8v26-Hodge-Mikhlin}, and
Theorem~\ref{thm:e8v26-W-Mikhlin} give
\[
              \|\partial_k^\eta M_j(k)\|
              \leq C_{\eta}\,2^{j|\eta|}.                         \tag{26.25}
\]
The support has four-dimensional measure at most \(C2^{-4j}\).
The undifferentiated Fourier integral therefore gives
\(\|K_j(x)\|\leq C2^{-4j}\).

If \(|x|_1>2^j\), choose a coordinate \(\ell\) with
\(|x_\ell|\geq|x|_1/4\).  Integrating (26.22) by parts \(N\) times in
\(k_\ell\), using periodicity and the smooth annular support, gives
\[
 \|K_j(x)\|
 \leq |x_\ell|^{-N}
       \frac1{(2\pi)^4}\int
       \|\partial_{k_\ell}^NM_j(k)\|\,dk
 \leq C_N2^{-4j}
       (2^{-j}|x|_1)^{-N}.
\]
Combining the two regimes proves (26.23).
\end{proof}

\begin{corollary}[Uniform scaled polynomial Schur stocks]
\label{cor:e8v26-annular-Schur}
For every \(s\geq0\), choosing \(N>s+4\) in (26.23) gives
\[
\boxed{\qquad
 \sum_{x\in{\mathbb Z}^4}
 (1+2^{-j}|x|_1)^s\|K_j(x)\|
 \leq A'_s
\qquad}                                                           \tag{26.26}
\]
uniformly in \(j\).  The periodized kernels on a finite torus obey the
same bound after summing over lattice images.
\end{corollary}

\begin{proof}
Insert (26.23).  The lattice-point count in a shell of scaled radius
\(m\) is bounded by a constant times \(2^{4j}(m+1)^3\).  The factor
\(2^{-4j}\) cancels this volume, and the remaining series converges
for \(N>s+4\).  Periodization only reorganizes the same nonnegative
majorant over images.
\end{proof}

\begin{corollary}[Tensor-shell kernel bound]
\label{cor:e8v26-tensor-shell}
For a shell multi-index
\({\bf j}=(j_1,\ldots,j_n)\), the absolute row or column stock of the
tensor kernel
\[
                  K_{j_1}\otimes\cdots\otimes K_{j_n}
\]
in the product scaled polynomial weight is at most
\((A'_s)^n\).  Orthogonal symmetrization cannot increase its Hilbert
norm.
\end{corollary}

\begin{proof}
The product support sum factorizes into the \(n\) one-leg sums in
(26.26).  The last statement is the norm-one property of an
orthogonal projection.
\end{proof}

\subsection{The exact principal margin survives the shell split}

The sixteen fine aliases over one coarse momentum form one fibre, so
the scalar multiplier \(\chi_j(k)\) acts on both the fine and coarse
matched spaces.  It commutes with the fibre map \({\cal W}(k)^*\).

\begin{theorem}[Annular Littlewood--Paley principal contraction]
\label{thm:e8v26-LP-contraction}
Let \(\{\chi_j\}\), together with the finite ultraviolet pieces, be a
square partition of unity on the nonharmonic base momenta.  Define
\[
 \|u\|_{\rm LP}^2:=\sum_j\|\chi_ju\|_2^2.                         \tag{26.27}
\]
Then the normalized coefficient map satisfies
\[
                         \|{\cal W}^*u\|_{\rm LP}
                         \leq\|u\|_{\rm LP}.                       \tag{26.28}
\]
On the tensor Littlewood--Paley decomposition, indexed separately on
each leg, the exact flat Gaussian principal derivative obeys
\[
\boxed{\qquad
                  \|\mathcal L_{\rm G}T\|_{\rm LP}
                  \leq\frac14\|T\|_{\rm LP}.
\qquad}                                                           \tag{26.29}
\]
\end{theorem}

\begin{proof}
Fibrewise commutation and the partial-isometry bound give
\[
 \|\chi_j{\cal W}^*u\|_2
 =\|{\cal W}^*\chi_ju\|_2\leq\|\chi_ju\|_2.
\]
Square and sum over \(j\) to obtain (26.28).  For an \(n\)-leg shell
multi-index, tensorization gives norm at most one and symmetrization
does not increase it.  Multiplication by the canonical factors
\(s_2=s_3=1/4\), \(s_4=1/16\), followed by the orthogonal direct sum
over degrees, gives (26.29).
\end{proof}

\begin{assessment}[What V26 closes]
\label{ass:e8v26-closes}
V26 determines the correct localization class of the flat normalized
curvature map.  It proves that a global exponentially local
orthonormal curvature chart is impossible because its initial Hodge
projection is directionally discontinuous at zero momentum.  It also
proves a uniform angular Gram floor, order-zero Mikhlin estimates,
rapid annular kernel decay at spatial scale \(2^j\), and preservation
of the exact \(1/4\) principal contraction in the annular
Littlewood--Paley Hilbert norm.
\end{assessment}

\begin{assessment}[Remaining owner-space problem]
\label{ass:e8v26-remains}
The annular kernel has a uniform scale-adapted polynomial
\(\ell^1\) stock, but an unrestricted coefficient family may occupy
infinitely many owner positions and infinitely many infrared shells.
Thus (26.26) alone is not the invariant nonlinear polymer norm.
The next problem is to combine position, shell, and connected-hull
weights in a wavelet or Besov owner space so that:
\begin{enumerate}
\item the order-zero matrix is almost diagonal;
\item the exact Hilbert margin is retained by interpolation;
\item paraproducts and generated interactions close without summing
an unweighted infinity of shells; and
\item the low shell is retained and rescaled rather than inverted.
\end{enumerate}
\end{assessment}

\begin{decision}[V27 acquisition order]
\label{dec:e8v26-v27}
V27 will construct a discrete annular wave-packet coefficient space.
It will first prove an almost-diagonal estimate from (26.23), then
derive the precise shell-weight condition under which the canonical
\(1/4\) margin dominates the weighted endpoint.  Only after that
linear gate is explicit will nonlinear paraproduct and owner-hull
incidence constants be admitted.
\end{decision}

\begin{firewall}[Claim boundary after eighth-series V26]
\label{fire:e8v26-claim}
V26 does not prove an invariant nonlinear Besov/polymer ball, the
interacting covariance rechart estimate in annular coordinates,
large-field contraction, marginal asymptotic freedom, regulator
removal, continuum reconstruction, or the Yang--Mills mass gap.  It
also does not assign a numerical interval to \(C_2\); it proves that
the finite certificate exists and identifies the exact matrix whose
derivatives must be bounded.
\end{firewall}

\paragraph{Parent-version record.}
V26 was formed by exact append-only extension of
\texttt{Renewal\_Yang\_Mills\_Mass\_Gap\_Research\_Notes\_V25.tex}.
The V25 parent had \(11\,481\) validator-counted source lines,
\(399\,334\) bytes, and SHA--256
\[
\texttt{65D180F7C6B7F19BA4296971A44397E96107DB6B0916A76EC39C0295A8901064}.
\]

% =============================================================================
% END APPENDED EIGHTH-SERIES V26 MATERIAL
% =============================================================================

% =============================================================================
% BEGIN APPENDED EIGHTH-SERIES V27 MATERIAL
% =============================================================================

\section{Eighth-series V27: an almost-diagonal wave-packet owner
space}

\subsection{Objective}

V26 proved that the normalized curvature map is an order-zero Hodge
multiplier: it cannot be exponentially local through zero momentum,
but its restriction to a dyadic annulus has a rapidly decreasing
kernel at that annulus's natural spatial scale.  V27 converts this
fact into an unrestricted linear owner space.  The key device is a
slowly varying summable cube weight.  Slow variation makes its
commutator with the exact Hilbert contraction arbitrarily small;
summability makes the owner sum finite.

\subsection{Spatial packets at the annular scale}

Use the annular square partition \(\{\chi_j\}\) from V26.  The
finitely many ultraviolet pieces may be reindexed with the first
infrared shell, so take \(j\geq0\) and put
\[
                              \ell_j:=2^j.                         \tag{27.1}
\]
For each \(j\), choose
\(\widetilde\chi_j\in C_c^\infty\) on a fixed enlargement of the same
annulus so that
\[
 0\leq\widetilde\chi_j\leq1,\qquad
 \widetilde\chi_j=1\ \hbox{on }\operatorname{supp}\chi_j.         \tag{27.1a}
\]
On the coarse base lattice, partition space into the disjoint cubes
\[
 Q_{j,m}:=\ell_jm+\{0,\ldots,\ell_j-1\}^4,
 \qquad m\in{\mathbb Z}^4,                                       \tag{27.2}
\]
and let \(E_{j,m}\) be multiplication by the indicator of
\(Q_{j,m}\).  Fine aliases are internal components over the same
coarse base cell.

For the normalized one-particle coefficient map define
\[
 {\mathbb A}_{j;m,m'}:=
 E_{j,m}\,{\cal W}^*\widetilde\chi_j\,E_{j,m'}.                   \tag{27.3}
\]
The array \({\mathbb A}_j=({\mathbb A}_{j;m,m'})\) acts between the
orthogonal direct sums of cube Hilbert spaces.

\begin{lemma}[Unweighted packet contraction]
\label{lem:e8v27-packet-contraction}
For every \(j\),
\[
                         \|{\mathbb A}_j\|_{2\to2}\leq1.           \tag{27.4}
\]
\end{lemma}

\begin{proof}
The maps \(u\mapsto(E_{j,m}u)_m\) are isometries from the spatial
Hilbert space to the direct sum of cube spaces.  Hence the block
matrix in (27.3) has the same norm as
\({\cal W}^*\widetilde\chi_j\).  Since
\(\|{\cal W}^*\|\leq1\) and
\(\|\widetilde\chi_j\|_\infty\leq1\), (27.4) follows.
\end{proof}

\begin{lemma}[Almost-diagonal cube matrix]
\label{lem:e8v27-almost-diagonal}
For every integer \(N\geq0\), there is a finite constant \(B_N\),
independent of \(j\), volume, and RG scale, such that
\[
\boxed{\qquad
 \|{\mathbb A}_{j;m,m'}\|
 \leq B_N(1+|m-m'|_1)^{-N}.
\qquad}                                                           \tag{27.5}
\]
\end{lemma}

\begin{proof}
The proof of Theorem~\ref{thm:e8v26-annular-kernel} applies without
change to \({\cal W}^*\widetilde\chi_j\), since the enlarged cutoffs
have the same uniform derivative and support bounds.  Denote that
kernel again by \(K_j\).  If \(x\in Q_{j,m}\) and
\(y\in Q_{j,m'}\), then, apart from a fixed enlargement accounting
for neighbouring cubes,
\[
 1+\ell_j^{-1}|x-y|_1\geq
 c(1+|m-m'|_1).                                                   \tag{27.6}
\]
Choose decay order \(N+5\) in (26.23).  For a fixed \(x\), summing
the kernel majorant over \(y\in Q_{j,m'}\) cancels the factor
\(\ell_j^{-4}\) against the number of lattice sites in a cube and
gives
\[
 \sup_{x\in Q_{j,m}}\sum_{y\in Q_{j,m'}}
 \|K_j(x-y)\|
 \leq C_N(1+|m-m'|_1)^{-N}.
\]
The same bound holds for the column sum.  The block Schur test proves
(27.5).  Enlarging \(B_N\) covers the finitely many ultraviolet
pieces.
\end{proof}

\subsection{A slowly varying summable owner weight}

Fix
\[
                              s>2,\qquad R\geq1,                   \tag{27.7}
\]
and define on cube indices
\[
                     w_{R,s}(m):=
                     \left(1+\frac{|m|_1}{R}\right)^s.            \tag{27.8}
\]

\begin{lemma}[Weight-ratio bound]
\label{lem:e8v27-ratio}
For all \(m,m'\in{\mathbb Z}^4\),
\[
 \max\left\{
 \frac{w_{R,s}(m)}{w_{R,s}(m')},
 \frac{w_{R,s}(m')}{w_{R,s}(m)}
 \right\}
 \leq
 \left(1+\frac{|m-m'|_1}{R}\right)^s.                             \tag{27.9}
\]
Consequently
\[
 \left|\frac{w_{R,s}(m)}{w_{R,s}(m')}-1\right|
 \leq
 \left(1+\frac{|m-m'|_1}{R}\right)^s-1.                           \tag{27.10}
\]
\end{lemma}

\begin{proof}
The triangle inequality gives
\[
 1+\frac{|m|_1}{R}
 \leq
 \left(1+\frac{|m'|_1}{R}\right)
 \left(1+\frac{|m-m'|_1}{R}\right).
\]
Raise to the power \(s\), and interchange \(m,m'\) for the reciprocal
bound.  If a positive ratio \(a\) satisfies \(A^{-1}\leq a\leq A\)
with \(A\geq1\), then \(|a-1|\leq A-1\), proving (27.10).
\end{proof}

\begin{definition}[Packet commutator stock]
\label{def:e8v27-epsilon}
Choose an integer \(N>s+4\) and set
\[
 \varepsilon_{N,s}(R):=
 B_N\sum_{d\in{\mathbb Z}^4}
 (1+|d|_1)^{-N}
 \left[
 \left(1+\frac{|d|_1}{R}\right)^s-1
 \right].                                                        \tag{27.11}
\]
\end{definition}

\begin{lemma}[The commutator stock vanishes at slow variation]
\label{lem:e8v27-epsilon-limit}
For every \(N>s+4\),
\[
              \varepsilon_{N,s}(R)<\infty,\qquad
              \lim_{R\to\infty}\varepsilon_{N,s}(R)=0.            \tag{27.12}
\]
\end{lemma}

\begin{proof}
For \(R\geq1\), the summand in brackets is bounded by
\((1+|d|_1)^s\).  Hence the whole summand is dominated by a constant
times \((1+|d|_1)^{s-N}\), which is summable on
\({\mathbb Z}^4\) because \(N-s>4\).  For every fixed \(d\), the
bracket tends to zero.  Finiteness and the limit follow from dominated
convergence.
\end{proof}

\begin{theorem}[Slow-weight perturbation of the exact contraction]
\label{thm:e8v27-slow-weight}
Let \(W_{R,s}\) be diagonal multiplication by \(w_{R,s}(m)\) on the
cube direct sum.  Then, uniformly in \(j\),
\[
\boxed{\qquad
 \|W_{R,s}{\mathbb A}_jW_{R,s}^{-1}\|_{2\to2}
 \leq1+\varepsilon_{N,s}(R).
\qquad}                                                           \tag{27.13}
\]
\end{theorem}

\begin{proof}
Write
\[
 W_{R,s}{\mathbb A}_jW_{R,s}^{-1}
 ={\mathbb A}_j+
 [W_{R,s},{\mathbb A}_j]W_{R,s}^{-1}.                             \tag{27.14}
\]
The first term has norm at most one by
Lemma~\ref{lem:e8v27-packet-contraction}.  The \((m,m')\) block of
the second is
\[
 \left(\frac{w_{R,s}(m)}{w_{R,s}(m')}-1\right)
 {\mathbb A}_{j;m,m'}.
\]
Lemmas~\ref{lem:e8v27-almost-diagonal} and
\ref{lem:e8v27-ratio} bound both its absolute block-row and
block-column sums by (27.11).  The block Schur test completes the
proof.
\end{proof}

\begin{definition}[A concrete slow radius]
\label{def:e8v27-Rstar}
Let
\[
 R_*:=\min\left\{R\in{\mathbb N}:
 \varepsilon_{N,s}(R)\leq2^{1/3}-1\right\}.                       \tag{27.15}
\]
This integer exists by Lemma~\ref{lem:e8v27-epsilon-limit}.
\end{definition}

\begin{corollary}[One-particle weighted endpoint with explicit margin]
\label{cor:e8v27-one-particle}
At \(R=R_*\),
\[
 \|W_{R_*,s}{\mathbb A}_jW_{R_*,s}^{-1}\|
 \leq2^{1/3}.                                                     \tag{27.16}
\]
Unlike V25 NBR, this estimate uses a polynomial weight on
scale-\(2^j\) packet positions and is compatible with the Hodge
discontinuity.
\end{corollary}

\begin{proof}
Insert (27.15) into (27.13).
\end{proof}

\subsection{Shell summability and the packet Hilbert norm}

Fix
\[
                   \tau>\frac12,\qquad J\geq1,
\]
and define
\[
 v_{J,\tau}(j):=\left(1+\frac jJ\right)^\tau,\qquad j\geq0.       \tag{27.17}
\]
The inverse-square stocks
\[
\begin{split}
 {\cal P}_{R,s}
 &:=\sum_{m\in{\mathbb Z}^4}w_{R,s}(m)^{-2}<\infty,\\
 {\cal V}_{J,\tau}
 &:=\sum_{j\geq0}v_{J,\tau}(j)^{-2}<\infty                       \tag{27.18}
\end{split}
\]
are finite precisely under the declared strict inequalities
\(s>2\) and \(\tau>1/2\).

\begin{definition}[One-particle wave-packet norm]
\label{def:e8v27-one-packet-norm}
For a coefficient covector \(u\), let
\[
\boxed{\qquad
 \|u\|_{{\rm WP};R,s,J,\tau}^2
 :=
 \sum_{j\geq0}v_{J,\tau}(j)^2
 \sum_{m\in{\mathbb Z}^4}w_{R,s}(m)^2
 \|E_{j,m}\chi_ju\|_2^2.
\qquad}                                                           \tag{27.19}
\]
All cube indices are measured relative to the transported owner.
\end{definition}

\begin{theorem}[One-particle packet stability]
\label{thm:e8v27-one-packet}
For every \(R\geq1\),
\[
\boxed{\qquad
 \|{\cal W}^*u\|_{{\rm WP};R,s,J,\tau}
 \leq
 \bigl(1+\varepsilon_{N,s}(R)\bigr)
 \|u\|_{{\rm WP};R,s,J,\tau}.
\qquad}                                                           \tag{27.20}
\]
In particular the bound is at most \(2^{1/3}\) at \(R=R_*\).
\end{theorem}

\begin{proof}
Put \(u_j=\chi_ju\).  Since the normalized fibre map commutes with
the scalar base-momentum cutoff,
\[
 E_{j,m}\chi_j{\cal W}^*u
 =E_{j,m}{\cal W}^*u_j.                                          \tag{27.21}
\]
Moreover \(\widetilde\chi_ju_j=u_j\).  Decomposing \(u_j\) into the
orthogonal cube pieces \(E_{j,m'}u_j\), the right side of (27.21) is
therefore exactly the \(m\)-th output of the block matrix
\({\mathbb A}_j\).  Theorem~\ref{thm:e8v27-slow-weight} gives the
weighted cube estimate at fixed \(j\).  Multiplication by the scalar
shell weight \(v_{J,\tau}(j)\), squaring, and summing in \(j\) proves
(27.20).
\end{proof}

\subsection{Tensor packets and the unrestricted principal map}

For \(n\in\{2,3,4\}\), use multi-indices
\[
 {\bf j}=(j_1,\ldots,j_n),\qquad
 {\bf m}=(m_1,\ldots,m_n),                                       \tag{27.22}
\]
and the product analysis packet
\[
 T_{n,c;{\bf j},{\bf m}}
 :=
 \left(\bigotimes_{a=1}^nE_{j_a,m_a}\chi_{j_a}\right)T_{n,c}.     \tag{27.23}
\]

\begin{definition}[Stable wave-packet chaos norm]
\label{def:e8v27-chaos-norm}
At \(R=R_*\), define
\[
\boxed{
\begin{split}
 \|T\|_{\mathscr{WP}}^2
 :=\sum_{n=2}^4 n!
 \sum_{c\in{\cal C}_n}
 \sum_{{\bf j},{\bf m}}
 \prod_{a=1}^n
 \left[v_{J,\tau}(j_a)^2w_{R_*,s}(m_a)^2\right]
 \|T_{n,c;{\bf j},{\bf m}}\|_{F_n}^2.
\end{split}}                                                     \tag{27.24}
\]
The local internal fibre norm includes the covariance-matched
one-particle metric and the finite alias components.
\end{definition}

\begin{theorem}[Unrestricted wave-packet principal contraction]
\label{thm:e8v27-principal}
The exact flat Gaussian principal derivative satisfies
\[
\boxed{\qquad
                     \|\mathcal L_{\rm G}T\|_{\mathscr{WP}}
                     \leq\frac12\|T\|_{\mathscr{WP}}.
\qquad}                                                           \tag{27.25}
\]
This estimate is uniform in volume and includes arbitrary packet
positions and arbitrarily deep infrared shells.
\end{theorem}

\begin{proof}
On a fixed degree and shell multi-index, the coefficient map is the
symmetric restriction of the tensor product of the one-particle
packet matrices.  Theorem~\ref{thm:e8v27-one-packet} and
Definition~\ref{def:e8v27-Rstar} bound its norm by
\[
                         (2^{1/3})^n.                             \tag{27.26}
\]
Shell weights are diagonal and are preserved because the flat fibre
map does not change the coarse base momentum.  Owner colours form a
finite orthogonal direct sum, and symmetrization is an orthogonal
projection.  After the canonical factors are inserted,
\[
 \max\left\{
 \frac14\,2^{2/3},
 \frac14\,2,
 \frac1{16}\,2^{4/3}
 \right\}=\frac12.                                                \tag{27.27}
\]
The factorial weights occur on both sides of each degree block.
Summing the squares over all packet and degree indices proves
(27.25).
\end{proof}

\subsection{Absolute owner summability}

\begin{definition}[Packet owner norm]
\label{def:e8v27-owner-norm}
Define the absolute packet-owner norm by
\[
 \|T\|_{{\rm own},{\rm WP}}
 :=
 \sum_{n=2}^4\iota_n
 \sum_{c\in{\cal C}_n}
 \sum_{{\bf j},{\bf m}}
 \|T_{n,c;{\bf j},{\bf m}}\|_{F_n}.                              \tag{27.28}
\]
As before, the link-to-plaquette incidence factor
\(\iota_n\leq6\) is paid once, not once per tensor leg.
\end{definition}

\begin{theorem}[Volume-uniform packet-owner embedding]
\label{thm:e8v27-owner-embedding}
Let
\[
 C_{\rm WP}:=
 \left[
 \sum_{n=2}^4
 \frac{\iota_n^2h_n}{n!}
 \left({\cal P}_{R_*,s}{\cal V}_{J,\tau}\right)^n
 \right]^{1/2}.                                                   \tag{27.29}
\]
Then
\[
\boxed{\qquad
 \|T\|_{{\rm own},{\rm WP}}
 \leq C_{\rm WP}\|T\|_{\mathscr{WP}}.
\qquad}                                                           \tag{27.30}
\]
Consequently
\[
 \|\mathcal L_{\rm G}T\|_{{\rm own},{\rm WP}}
 \leq\frac12C_{\rm WP}\|T\|_{\mathscr{WP}}.                       \tag{27.31}
\]
\end{theorem}

\begin{proof}
At fixed \(n\), insert the inverse of the product weight in (27.24)
into (27.28) and apply Cauchy--Schwarz to owner colour, shell, and
cube indices.  For each tensor leg the inverse-square sum factorizes
as
\[
 \sum_{j,m}v_{J,\tau}(j)^{-2}w_{R_*,s}(m)^{-2}
 ={\cal V}_{J,\tau}{\cal P}_{R_*,s}.                              \tag{27.32}
\]
The colour sum contributes \(\sqrt{h_n}\), and the incidence
contributes \(\iota_n\).  Apply Cauchy--Schwarz once more over the
three degrees, pairing with \(n!^{-1/2}\), to obtain (27.30).
Equation (27.31) follows from
Theorem~\ref{thm:e8v27-principal}.
\end{proof}

\begin{corollary}[The linear owner obstruction is closed]
\label{cor:e8v27-linear-owner}
The range-zero counterexample of V24 and the Hodge discontinuity of
V26 are both absent from the wave-packet formulation:
\begin{enumerate}
\item arbitrary relative packet positions are summed by
\({\cal P}_{R_*,s}<\infty\);
\item infinitely many infrared shells are summed by
\({\cal V}_{J,\tau}<\infty\);
\item the normalized Hodge map is only required to be almost diagonal
at the shell scale, which is guaranteed by V26; and
\item the principal map is a strict contraction in the stronger
Hilbert norm.
\end{enumerate}
\end{corollary}

\begin{proof}
These are respectively Theorems
\ref{thm:e8v27-owner-embedding},
\ref{thm:e8v27-owner-embedding},
\ref{lem:e8v27-almost-diagonal}, and
\ref{thm:e8v27-principal}.
\end{proof}

\begin{firewall}[Packet owner norm versus geometric polymer norm]
\label{fire:e8v27-packet-polymer}
The norm (27.28) sums localized Hilbert packet blocks.  It is not yet
identified with the connected geometric-hull activity norm used for
nonlinear polymer multiplication.  In particular, multiplying two
fields convolves their momenta: comparable high shells can cancel
into a much lower output shell.  The shell weights in (27.24) must
therefore be tested by a paraproduct decomposition before an
interacting invariant ball is claimed.
\end{firewall}

\begin{assessment}[What V27 closes]
\label{ass:e8v27-closes}
V27 closes the unrestricted \emph{linear} localization problem left
by V24--V26.  It constructs a volume-uniform wave-packet Hilbert
space that has all three required properties: arbitrary positions
and shells are allowed, the absolute packet-owner sum is continuous,
and the exact Gaussian principal RG derivative contracts by at most
\(1/2\).  No exponential localization across the Hodge singularity
is assumed.
\end{assessment}

\begin{assessment}[Current bottleneck]
\label{ass:e8v27-remains}
The bottleneck has moved from the one-particle map to nonlinear shell
arithmetic.  The generated quadratic, cubic, and quartic interactions,
the covariance rechart, and the moving-reference residues must be
written as annular paraproducts.  The dangerous channel is
high--high to low cancellation, where the output shell index can be
arbitrarily larger than both input indices.
\end{assessment}

\begin{decision}[V28 acquisition order]
\label{dec:e8v27-v28}
V28 will prove the discrete shell-triad geometry and split every
bilinear product into low--high, high--low, comparable-shell
nonresonant, and comparable-shell resonant pieces.  It will determine
whether the polynomial shell weight \(v_{J,\tau}\) closes the
resonant sum.  If it does not, the low output will be assigned to the
running local/marginal reference graph rather than forced into the
stable packet norm.
\end{decision}

\begin{firewall}[Claim boundary after eighth-series V27]
\label{fire:e8v27-claim}
V27 does not prove nonlinear paraproduct closure, an interacting
stable ball, the rechart and Wilson remainder estimates in packet
norm, large-field control, marginal asymptotic freedom, regulator
removal, continuum reconstruction, or the Yang--Mills mass gap.
\end{firewall}

\paragraph{Parent-version record.}
V27 was formed by exact append-only extension of
\texttt{Renewal\_Yang\_Mills\_Mass\_Gap\_Research\_Notes\_V26.tex}.
The V26 parent had \(12\,067\) validator-counted source lines,
\(420\,241\) bytes, and SHA--256
\[
\texttt{42BFF10B81D57FDFFBB30BB2B4ECE5E4BFF9A0B8FA5EDF343850A19705A92804}.
\]

% =============================================================================
% END APPENDED EIGHTH-SERIES V27 MATERIAL
% =============================================================================

% =============================================================================
% BEGIN APPENDED EIGHTH-SERIES V28 MATERIAL
% =============================================================================

\section{Eighth-series V28: shell paraproduct closure and the exact
zero-mode ledger}

\subsection{Resolution and the finite-volume zero mode}

Choose a smooth bounded-overlap resolution
\(\{\varphi_j\}_{j\geq0}\) on the nonzero Brillouin momenta, with
\[
 \sum_{j\geq0}\varphi_j(k)=1\quad(k\ne0),\qquad
 \operatorname{supp}\varphi_j
 \subseteq\{2^{-j-1}\leq\rho(k)\leq2^{-j+1}\},                   \tag{28.1}
\]
after absorbing finitely many ultraviolet pieces into \(j=0\).
Let
\[
                 \Delta_j=\varphi_j(D),\qquad
                 P_\circ=I-P_0,                                 \tag{28.2}
\]
where \(P_0\) is the exact zero-momentum projection.  On an infinite
lattice \(P_0\) is absent in \(L^2\); on every finite torus it is a
declared finite-dimensional retained sector.

\begin{definition}[Polynomial shell Hilbert space]
\label{def:e8v28-shell-space}
With the V27 weight
\[
                    v_j=v_{J,\tau}(j)=\left(1+\frac jJ\right)^\tau,
 \qquad \tau>\frac12,                                             \tag{28.3}
\]
define
\[
 \|f\|_{\mathcal B_{J,\tau}}^2
 :=\sum_{j\geq0}v_j^2\|\Delta_jf\|_2^2.                           \tag{28.4}
\]
Different smooth resolutions satisfying (28.1) give equivalent norms;
the finite overlap constant is retained in every estimate below.
\end{definition}

\subsection{Exact shell-triad geometry}

\begin{lemma}[Lattice momentum triangle inequality]
\label{lem:e8v28-rho-triangle}
For momenta \(k,\ell\), added modulo \(2\pi\),
\[
\boxed{\qquad
 |\rho(k)-\rho(\ell)|
 \leq\rho(k+\ell)
 \leq\rho(k)+\rho(\ell).
\qquad}                                                           \tag{28.5}
\]
\end{lemma}

\begin{proof}
Coordinatewise,
\[
 p_\mu(k+\ell)=p_\mu(k)+e^{ik_\mu}p_\mu(\ell).                    \tag{28.6}
\]
The diagonal multiplier \(e^{ik_\mu}\) is unitary on
\({\mathbb C}^4\).  The ordinary and reverse triangle inequalities
for its Euclidean norm prove (28.5).
\end{proof}

\begin{lemma}[Shell-triad alternatives]
\label{lem:e8v28-triads}
There is a fixed integer \(C_0\), depending only on the cutoff
supports in (28.1), such that
\[
 \Delta_k\bigl((\Delta_if)(\Delta_jg)\bigr)\ne0                   \tag{28.7}
\]
implies one of the following alternatives:
\begin{enumerate}
\item if \(|i-j|>C_0\), then
\[
                         |k-\min\{i,j\}|\leq C_0;                 \tag{28.8}
\]
\item if \(|i-j|\leq C_0\), then
\[
                         k\geq\min\{i,j\}-C_0.                    \tag{28.9}
\]
\end{enumerate}
The second alternative permits arbitrarily large \(k\), namely
high--high cancellation into a low nonzero output.
\end{lemma}

\begin{proof}
Let \(q\) and \(q'\) be momenta in the \(i\)- and \(j\)-shells.
If, say, \(j\geq i+C_0\) with \(C_0\) chosen larger than the fixed
cutoff-width ratio, then \(\rho(q')\) is at most one quarter of the
lower bound for \(\rho(q)\).  Equation (28.5) makes
\(\rho(q+q')\) comparable to \(\rho(q)\), proving (28.8).
For comparable inputs, (28.5) gives
\[
 \rho(q+q')\leq\rho(q)+\rho(q')
 \leq C2^{-\min\{i,j\}},
\]
which is equivalent to (28.9) after increasing \(C_0\).
\end{proof}

\subsection{Four-dimensional Bernstein gains}

\begin{lemma}[Annular Bernstein estimates]
\label{lem:e8v28-Bernstein}
There are cutoff-dependent constants \(C_{\rm B},C_{\rm P}\),
uniform on all finite tori, such that
\[
\begin{split}
 \|\Delta_jf\|_\infty
 &\leq C_{\rm B}2^{-2j}\|\Delta_jf\|_2,\\
 \|\Delta_k h\|_2
 &\leq C_{\rm P}2^{-2k}\|h\|_1.                                 \tag{28.10}
\end{split}
\]
Consequently
\[
\begin{split}
 \|\Delta_k((\Delta_if)(\Delta_jg))\|_2
 &\leq C2^{-2\max\{i,j\}}
       \|\Delta_if\|_2\|\Delta_jg\|_2
       &&(|i-j|>C_0),\\
 \|\Delta_k((\Delta_if)(\Delta_jg))\|_2
 &\leq C2^{-2k}
       \|\Delta_if\|_2\|\Delta_jg\|_2
       &&(|i-j|\leq C_0).                                       \tag{28.11}
\end{split}
\]
\end{lemma}

\begin{proof}
The Fourier support of a \(j\)-shell has normalized measure at most
\(C2^{-4j}\).  Fourier inversion followed by Cauchy--Schwarz gives
the first estimate in (28.10).  The convolution kernel of
\(\Delta_k\) has \(L^2\) norm at most \(C2^{-2k}\) by Plancherel;
Young's \(L^1*L^2\to L^2\) inequality gives the second.

For separated shells, apply the first Bernstein estimate to the
lower-frequency factor, whose index is \(\max\{i,j\}\), and use
\(\|uv\|_2\leq\|u\|_2\|v\|_\infty\).  The output projection is
\(L^2\)-bounded.  For comparable shells, use
\(\|uv\|_1\leq\|u\|_2\|v\|_2\) and the second estimate in (28.10).
The normalized counting measure gives the same constants on finite
tori.
\end{proof}

\begin{remark}[Why the resonant infrared tail is summable]
\label{rem:e8v28-resonant-gain}
The factor \(2^{-2k}\) is the square root of the four-dimensional
output-shell volume.  It is exponential in the shell index and
therefore dominates the polynomial shell weight \(v_k\).  This is the
mechanism missing from a purely combinatorial shell count.
\end{remark}

\subsection{The nonzero-shell algebra}

\begin{theorem}[Polynomial shell algebra in four dimensions]
\label{thm:e8v28-shell-algebra}
For every \(J\geq1\) and \(\tau>1/2\), there is a finite constant
\(C_{\rm alg}=C_{\rm alg}(J,\tau,\{\varphi_j\})\), independent of
volume, such that
\[
\boxed{\qquad
 \|P_\circ(fg)\|_{\mathcal B_{J,\tau}}
 \leq C_{\rm alg}
 \|f\|_{\mathcal B_{J,\tau}}
 \|g\|_{\mathcal B_{J,\tau}}.
\qquad}                                                           \tag{28.12}
\]
The result also holds for a finite-dimensional fibre bilinear map of
operator norm at most one.
\end{theorem}

\begin{proof}
Write
\[
 a_i=v_i\|\Delta_if\|_2,\qquad
 b_j=v_j\|\Delta_jg\|_2.                                         \tag{28.13}
\]
The sequences \(a,b\) are in \(\ell^2\), with norms equal to the two
norms on the right of (28.12), up to the fixed resolution-equivalence
constant.

First take \(j\geq i+C_0\).  By
Lemmas~\ref{lem:e8v28-triads} and
\ref{lem:e8v28-Bernstein}, the output has \(k=i+O(1)\), and its
weighted contribution is bounded by
\[
 C a_i\sum_{j\geq i+C_0}
       2^{-2j}\frac{b_j}{v_j}.                                   \tag{28.14}
\]
Cauchy--Schwarz gives
\[
 \sum_{j\geq i+C_0}2^{-2j}\frac{b_j}{v_j}
 \leq
 \left(\sum_{j\geq0}2^{-4j}v_j^{-2}\right)^{1/2}
 \|b\|_{\ell^2}.                                                  \tag{28.15}
\]
Taking the \(\ell^2\) norm in \(i\) proves the low--high bound.
The high--low term is symmetric.

For \(|i-j|\leq C_0\), put
\[
 c_i:=\sum_{|j-i|\leq C_0}
       \frac{a_i b_j}{v_i v_j}.                                  \tag{28.16}
\]
Since \(v_i,v_j\geq1\) and only finitely many \(j\)'s occur for each
\(i\),
\[
                         \sum_i c_i
 \leq C\|a\|_{\ell^2}\|b\|_{\ell^2}.                              \tag{28.17}
\]
Equations (28.9) and (28.11) bound the resonant output sequence by
\[
 d_k\leq
 C v_k2^{-2k}\sum_{i\leq k+C_0}c_i
 \leq
 C v_k2^{-2k}\|a\|_{\ell^2}\|b\|_{\ell^2}.                        \tag{28.18}
\]
The sequence \(v_k2^{-2k}\) is in \(\ell^2\), so the resonant
contribution has the required bound.  Finite cutoff overlap only
changes \(C\).  Summing the three paraproduct classes proves
(28.12).
\end{proof}

\subsection{The exact zero-output ledger}

\begin{definition}[Zero-output bilinear form]
\label{def:e8v28-zero-ledger}
On a finite torus with normalized Fourier measure, define
\[
 {\cal Z}(f,g):=P_0(fg),\qquad
 \widehat{{\cal Z}(f,g)}(0)
 =\int\widehat f(q)\widehat g(-q)\,dq.                            \tag{28.19}
\]
For a finite internal fibre, compose the integrand with the declared
fibre bilinear contraction.
\end{definition}

\begin{proposition}[Volume-uniform zero-output bound]
\label{prop:e8v28-zero-bound}
The exact zero-output form satisfies
\[
\boxed{\qquad
 \|{\cal Z}(f,g)\|
 \leq\|f\|_2\|g\|_2
 \leq
 \|f\|_{\mathcal B_{J,\tau}}
 \|g\|_{\mathcal B_{J,\tau}}.
\qquad}                                                           \tag{28.20}
\]
\end{proposition}

\begin{proof}
Cauchy--Schwarz in the normalized Fourier measure gives the first
inequality.  The bounded-overlap resolution and \(v_j\geq1\) give the
second, after fixing the harmless norm-equivalence constant in the
resolution.  The same proof works for normalized counting measure on
every finite torus.
\end{proof}

\begin{firewall}[The zero output is retained, not hidden]
\label{fire:e8v28-zero}
Theorem~\ref{thm:e8v28-shell-algebra} concerns \(P_\circ(fg)\).
The missing term is exactly \({\cal Z}(f,g)\), not an uncontrolled
infrared tail.  It belongs in a finite-dimensional retained ledger.
Gauge Ward identities and tensor type still have to determine whether
each such coefficient is forbidden, relevant, or marginal; V28 does
not automatically call every zero-output coefficient a running
coupling.
\end{firewall}

\subsection{Iteration to cubic and quartic local products}

\begin{corollary}[Finite polynomial nonzero-shell closure]
\label{cor:e8v28-polynomial}
Let \(m=2,3,4\), and let \(B_m\) be a finite-dimensional \(m\)-linear
fibre map.  Iterating Theorem~\ref{thm:e8v28-shell-algebra} gives
\[
 \left\|
 P_\circ B_m(f_1,\ldots,f_m)
 \right\|_{\mathcal B_{J,\tau}}
 \leq
 C_m\|B_m\|
 \prod_{a=1}^m\|f_a\|_{\mathcal B_{J,\tau}},                     \tag{28.21}
\]
provided the norm of every intermediate exact zero output is included
in the retained ledger.  The constants \(C_m\) are finite and volume
independent.
\end{corollary}

\begin{proof}
For \(m=2\), use Theorem~\ref{thm:e8v28-shell-algebra}.  For
\(m=3,4\), associate the product in a fixed order and decompose each
intermediate product as \(P_\circ+P_0\).  Apply the bilinear theorem
to the nonzero-shell term and
Proposition~\ref{prop:e8v28-zero-bound} to the zero term.
Multiplication by the resulting constant zero-mode coefficient
preserves every shell and is bounded by its finite-dimensional norm.
Reinsert both terms at the next association node.  There are only
finitely many nodes and fibre contractions.
\end{proof}

\subsection{Compatibility of packet position weights}

Let \(Q_{j,m}^+\) denote a fixed finite enlargement of the V27 cube
\(Q_{j,m}\), large enough to cover the spatial tails used in a chosen
packet cutoff.

\begin{lemma}[Coarser output does not enlarge the owner distance]
\label{lem:e8v28-position-transfer}
Suppose \(Q_{i,m_i}^+\cap Q_{k,m_k}^+\ne\varnothing\) and
\[
                              k\geq i-C_0.                        \tag{28.22}
\]
Then
\[
\boxed{\qquad
 w_{R,s}(m_k)\leq
 C_{s,C_0}\,w_{R,s}(m_i),
\qquad}                                                           \tag{28.23}
\]
uniformly in \(i,k,R\geq1\).  In a nonresonant low--high triad, one
may take \(i\) to be the higher-frequency input shell; in a resonant
triad, either comparable input may be used.
\end{lemma}

\begin{proof}
Choose a point \(x\) in the intersection of the enlarged cubes.
Their definitions imply
\[
 |m_k|_1\leq
 C+2^{-k}|x|_1,\qquad
 2^{-i}|x|_1\leq C+|m_i|_1.                                     \tag{28.24}
\]
Since \(k\geq i-C_0\),
\[
 |m_k|_1\leq C_{C_0}(1+|m_i|_1).
\]
Insert this in
\((1+|m_k|_1/R)^s\), use \(R\geq1\), and absorb fixed additive and
multiplicative constants into \(C_{s,C_0}\).
The shell choices follow from
Lemma~\ref{lem:e8v28-triads}.
\end{proof}

\begin{assessment}[Meaning of the position lemma]
\label{ass:e8v28-position}
Local multiplication can charge the output position weight to a
high-frequency input packet without a factor growing in the infrared
separation.  Smooth output projection is not strictly local, but its
off-cube tails are governed by the arbitrarily rapid estimates of
V26--V27.  Thus no geometric obstruction analogous to the Hodge
discontinuity remains in the position variable.
\end{assessment}

\begin{assessment}[What V28 closes]
\label{ass:e8v28-closes}
V28 closes the nonlinear \emph{shell} arithmetic identified in V27.
Polynomial shell weights are stable under all nonzero bilinear
paraproduct channels in four dimensions; high--high cancellation
gains \(2^{-2k}\) and is summable.  The exact finite-volume zero output
has a separate volume-uniform bound and is retained explicitly.
Finite cubic and quartic local products inherit the nonzero-shell
bound.
\end{assessment}

\begin{assessment}[Current bottleneck]
\label{ass:e8v28-remains}
The remaining functional-analytic step is to combine the shell
algebra with the slow packet-position weight and the connected
plaquette-hull owner rule in one multilinear estimate.  Lemma
\ref{lem:e8v28-position-transfer} handles intersecting packets;
nonintersecting packets require summing the V26 off-cube tails.
After that estimate, the already explicit Wilson Taylor stock
\(64e^{2\mu}gr\) and covariance-rechart rows can be transported into
the repaired space.
\end{assessment}

\begin{decision}[V29 acquisition order]
\label{dec:e8v28-v29}
V29 will prove a packet paraproduct matrix estimate with three
indices: output cube, first input cube, and second input cube.  It
will use the position transfer lemma for the nearest input and spend
arbitrary V26 decay on the other separations.  The target is a
volume-uniform bilinear bound in \(\mathscr{WP}\), with the exact
zero-output row directed to a Ward-typed finite ledger.
\end{decision}

\begin{firewall}[Claim boundary after eighth-series V28]
\label{fire:e8v28-claim}
V28 does not yet prove the full position-weighted packet
paraproduct, connected-hull submultiplicativity, the interacting RG
invariant ball, covariance rechart control in that ball, large-field
stability, marginal asymptotic freedom, continuum reconstruction, or
the Yang--Mills mass gap.
\end{firewall}

\paragraph{Parent-version record.}
V28 was formed by exact append-only extension of
\texttt{Renewal\_Yang\_Mills\_Mass\_Gap\_Research\_Notes\_V27.tex}.
The V27 parent had \(12\,587\) validator-counted source lines,
\(437\,295\) bytes, and SHA--256
\[
\texttt{1D7223E529700683043ED2FC6CDC50FA57CD5ADFA1B9E0ECB946E054C6EDD907}.
\]

% =============================================================================
% END APPENDED EIGHTH-SERIES V28 MATERIAL
% =============================================================================

% =============================================================================
% BEGIN APPENDED EIGHTH-SERIES V29 MATERIAL
% =============================================================================

\section{Eighth-series V29: position-weighted paraproduct closure}

\subsection{A physical-space form of the packet weight}

Fix an owner \(o\) on the coarse base lattice and retain
\(s>2\), \(R=R_*\).  At shell \(j\), define
\[
 \omega_j^o(x):=
 \left(1+\frac{|x-o|_1}{R2^j}\right)^s.                           \tag{29.1}
\]
This is the physical-space version of the V27 cube weight.

\begin{lemma}[Cube and physical weights are uniformly equivalent]
\label{lem:e8v29-weight-equivalence}
There is \(C_{\rm eq}=C_{\rm eq}(s,R)<\infty\), independent of
\(j,m,o\), such that for \(x\in Q_{j,m}+o\),
\[
\boxed{\qquad
 C_{\rm eq}^{-1}w_{R,s}(m)
 \leq\omega_j^o(x)
 \leq C_{\rm eq}w_{R,s}(m).
\qquad}                                                           \tag{29.2}
\]
\end{lemma}

\begin{proof}
For \(x-o=2^jm+q\), \(q\in\{0,\ldots,2^j-1\}^4\),
\[
 \left|2^{-j}|x-o|_1-|m|_1\right|\leq4.                          \tag{29.3}
\]
The two quantities \(1+|m|_1/R\) and
\(1+2^{-j}|x-o|_1/R\) are therefore bounded by
\((1+4/R)\) times one another.  Raise to the power \(s\).
\end{proof}

\begin{definition}[Position-weighted Besov packet norm]
\label{def:e8v29-BWP}
With \(v_j=v_{J,\tau}(j)\), define
\[
\boxed{\qquad
 \|f\|_{\mathfrak B_{\rm WP}}^2
 :=\sum_{j\geq0}v_j^2
       \|\omega_j^o\Delta_jf\|_2^2.
\qquad}                                                           \tag{29.4}
\]
\end{definition}

\begin{corollary}[Equivalence with the V27 analysis norm]
\label{cor:e8v29-norm-equivalence}
On one-particle coefficients, (29.4) is equivalent to
Definition~\ref{def:e8v27-one-packet-norm}:
\[
 C_{\rm eq}^{-1}\|f\|_{{\rm WP};R,s,J,\tau}
 \leq\|f\|_{\mathfrak B_{\rm WP}}
 \leq C_{\rm eq}\|f\|_{{\rm WP};R,s,J,\tau}.                     \tag{29.5}
\]
\end{corollary}

\begin{proof}
Square (29.2), multiply by
\(\|E_{j,m}\Delta_jf\|_2^2\), and sum over the disjoint cubes and
shells.
\end{proof}

\subsection{Weighted annular projection estimates}

\begin{lemma}[Scale-weight ratio]
\label{lem:e8v29-scale-ratio}
For all \(x,y\),
\[
 \frac{\omega_j^o(x)}{\omega_j^o(y)}
 \leq
 \left(1+\frac{|x-y|_1}{R2^j}\right)^s.                           \tag{29.6}
\]
\end{lemma}

\begin{proof}
Use
\[
 1+\frac{|x-o|_1}{R2^j}
 \leq
 \left(1+\frac{|y-o|_1}{R2^j}\right)
 \left(1+\frac{|x-y|_1}{R2^j}\right)
\]
and raise to the power \(s\).
\end{proof}

\begin{proposition}[Weighted \(L^2\) and \(L^1\)-to-\(L^2\)
projection bounds]
\label{prop:e8v29-weighted-projection}
Let \(\Psi_j\) be any of the smooth annular projections or enlarged
projections used in V26--V28.  There are finite constants
\(M_{1,s,R}\) and \(M_{2,s,R}\), uniform in \(j\), volume, and owner,
such that
\[
\begin{split}
 \|\omega_j^o\Psi_jh\|_2
 &\leq M_{1,s,R}\|\omega_j^oh\|_2,\\
 \|\omega_j^o\Psi_jh\|_2
 &\leq M_{2,s,R}2^{-2j}\|\omega_j^oh\|_1.                        \tag{29.7}
\end{split}
\]
\end{proposition}

\begin{proof}
The kernel \(L_j(z)\) of \(\Psi_j\) satisfies, for every \(N\),
\[
 |L_j(z)|\leq C_N2^{-4j}
              (1+2^{-j}|z|_1)^{-N}.                             \tag{29.8}
\]
By Lemma~\ref{lem:e8v29-scale-ratio},
\[
 \omega_j^o(x)|\Psi_jh(x)|
 \leq
 \sum_y |L_j(x-y)|
 \left(1+\frac{|x-y|_1}{R2^j}\right)^s
 \omega_j^o(y)|h(y)|.                                           \tag{29.9}
\]
For \(N>s+4\), the modified kernel on the right has a uniformly
bounded \(\ell^1\) norm, proving the first estimate by Young's
inequality.  For \(N>s+2\), its \(\ell^2\) norm is at most
\(C2^{-2j}\), proving the second
\(\ell^1\to\ell^2\) estimate.  Periodization reorganizes the same
positive majorant on a finite torus.
\end{proof}

\begin{lemma}[Cross-scale weight transfer]
\label{lem:e8v29-cross-scale}
If \(k\geq i-C_0\), then
\[
                         \omega_k^o(x)
 \leq C_{s,C_0}\omega_i^o(x)                                    \tag{29.10}
\]
for every \(x\), uniformly in the shell indices.
\end{lemma}

\begin{proof}
The hypothesis implies \(2^{-k}\leq2^{C_0}2^{-i}\).  Insert this
inequality in (29.1) and absorb \(2^{C_0s}\) into the constant.
\end{proof}

\subsection{The full position-and-shell paraproduct}

\begin{theorem}[Position-weighted packet algebra]
\label{thm:e8v29-packet-algebra}
There is a finite, volume- and owner-independent constant
\(C_{\rm WP,alg}\) such that
\[
\boxed{\qquad
 \|P_\circ(fg)\|_{\mathfrak B_{\rm WP}}
 \leq C_{\rm WP,alg}
 \|f\|_{\mathfrak B_{\rm WP}}
 \|g\|_{\mathfrak B_{\rm WP}}.
\qquad}                                                           \tag{29.11}
\]
The exact zero output remains in the V28 retained ledger.
\end{theorem}

\begin{proof}
Put \(f_i=\Delta_if\), \(g_j=\Delta_jg\), and
\[
 a_i=v_i\|\omega_i^of_i\|_2,\qquad
 b_j=v_j\|\omega_j^og_j\|_2.                                    \tag{29.12}
\]

Suppose first that \(j\geq i+C_0\).  The output has
\(|k-i|\leq C_0\).  Apply the first estimate in (29.7), transfer the
output weight to \(\omega_i^o\) with (29.10), and use the unweighted
Bernstein bound on the lower-frequency factor:
\[
\begin{split}
 \|\omega_k^o\Delta_k(f_i g_j)\|_2
 &\leq C\|\omega_i^of_i\,g_j\|_2\\
 &\leq C2^{-2j}
       \|\omega_i^of_i\|_2\|g_j\|_2\\
 &\leq C2^{-2j}
       \|\omega_i^of_i\|_2\|\omega_j^og_j\|_2.                   \tag{29.13}
\end{split}
\]
The shell weight \(v_k\) is comparable to \(v_i\).  Thus the
low--high sequence estimate is exactly (28.14)--(28.15), with
\(\|f_i\|_2,\|g_j\|_2\) replaced by their weighted versions.
The high--low case is symmetric.

If \(|i-j|\leq C_0\), every nonzero output obeys
\(k\geq i-C_0\).  The second estimate in (29.7), followed by
(29.10) and Cauchy--Schwarz, gives
\[
\begin{split}
 \|\omega_k^o\Delta_k(f_i g_j)\|_2
 &\leq C2^{-2k}\|\omega_k^of_i g_j\|_1\\
 &\leq C2^{-2k}
       \|\omega_i^of_i\|_2\|g_j\|_2\\
 &\leq C2^{-2k}
       \|\omega_i^of_i\|_2\|\omega_j^og_j\|_2.                   \tag{29.14}
\end{split}
\]
The resonant sequence estimate (28.16)--(28.18) now applies verbatim.
Summing the three paraproduct classes proves (29.11).
\end{proof}

\begin{corollary}[Closure in the V27 packet norm]
\label{cor:e8v29-V27-algebra}
By norm equivalence,
\[
 \|P_\circ(fg)\|_{{\rm WP};R,s,J,\tau}
 \leq
 C'_{\rm WP,alg}
 \|f\|_{{\rm WP};R,s,J,\tau}
 \|g\|_{{\rm WP};R,s,J,\tau}.                                   \tag{29.15}
\]
\end{corollary}

\begin{proof}
Combine Theorem~\ref{thm:e8v29-packet-algebra} with
Corollary~\ref{cor:e8v29-norm-equivalence}.
\end{proof}

\subsection{Finite local polynomial maps}

\begin{corollary}[Position-weighted polynomial closure]
\label{cor:e8v29-polynomial}
Let \(B_m\) be a finite-dimensional local \(m\)-linear fibre map,
\(2\leq m\leq4\).  Then
\[
 \|P_\circ B_m(f_1,\ldots,f_m)\|_{\mathfrak B_{\rm WP}}
 \leq C_{{\rm WP},m}\|B_m\|
       \prod_{a=1}^m\|f_a\|_{\mathfrak B_{\rm WP}},               \tag{29.16}
\]
provided every intermediate exact zero output is included in the
finite V28 ledger.
\end{corollary}

\begin{proof}
Iterate Theorem~\ref{thm:e8v29-packet-algebra}.  At each association
node split the product into \(P_\circ+P_0\), use
Proposition~\ref{prop:e8v28-zero-bound} on the zero part, and note
that multiplication by a retained finite-dimensional coefficient is
bounded and preserves the shell index.
\end{proof}

\subsection{Reanchoring and connected hulls}

\begin{lemma}[Owner reanchoring]
\label{lem:e8v29-reanchor}
For two owner cells \(o,o'\),
\[
\boxed{\qquad
 \|f\|_{\mathfrak B_{\rm WP},o'}
 \leq
 \left(1+\frac{|o-o'|_1}{R}\right)^s
 \|f\|_{\mathfrak B_{\rm WP},o}.
\qquad}                                                           \tag{29.17}
\]
\end{lemma}

\begin{proof}
For every shell \(j\),
\[
 \frac{\omega_j^{o'}(x)}{\omega_j^o(x)}
 \leq
 \left(1+\frac{|o-o'|_1}{R2^j}\right)^s
 \leq
 \left(1+\frac{|o-o'|_1}{R}\right)^s.
\]
Insert this pointwise inequality in (29.4).
\end{proof}

Let \({\mathfrak P}\) be the set of finite connected coarse-plaquette
polymers.  For \(X\in{\mathfrak P}\), let \(|X|\) be its number of
plaquettes and \({\rm diam}(X)\) its \(\ell^1\) diameter.  For a
polymer activity \(K=(K(X))\), define
\[
\begin{split}
 \|K\|_{\kappa,a}:=
 \sup_p\sum_{\substack{X\in{\mathfrak P}\\p\in X}}
 e^{\kappa|X|}(1+{\rm diam}(X))^a
 \max_{o\in X}\|K(X)\|_{\mathfrak B_{\rm WP},o}.                  \tag{29.18}
\end{split}
\]

For a bounded fibre-bilinear product \(B\), define the overlap
convolution
\[
 (K\star_B L)(Z)
 :=
 \sum_{\substack{X,Y\in{\mathfrak P}\\
 X\cap Y\ne\varnothing,\ X\cup Y=Z}}
 P_\circ B(K(X),L(Y)),                                           \tag{29.19}
\]
with its exact zero rows placed in the retained ledger.

\begin{theorem}[Connected-hull convolution with declared reserves]
\label{thm:e8v29-hull}
Let \(\kappa>\kappa'\geq0\), \(a\geq0\), and put
\[
 c_{\kappa-\kappa'}
 :=\sup_{m\geq1}m e^{-(\kappa-\kappa')m}<\infty.                  \tag{29.20}
\]
Then
\[
\boxed{\qquad
 \|K\star_B L\|_{\kappa',a}
 \leq
 4C_{\rm WP,alg}\|B\|\,
 c_{\kappa-\kappa'}
 \|K\|_{\kappa,a+s}\|L\|_{\kappa,a+s}.
\qquad}                                                           \tag{29.21}
\]
\end{theorem}

\begin{proof}
If \(X,Y\) are connected and overlap, then \(Z=X\cup Y\) is connected,
\[
 |Z|\leq|X|+|Y|,\qquad
 1+{\rm diam}(Z)
 \leq(1+{\rm diam}(X))(1+{\rm diam}(Y)).                         \tag{29.22}
\]
Fix \(p\in Z\) and, for each summand, an output owner \(o\in Z\)
realizing the finite maximum in (29.18).  Split into the four cases
according to \(p\in X\) or \(p\in Y\), and \(o\in X\) or \(o\in Y\).
This costs at most a factor four.  Consider the case
\(p,o\in X\).  Choose an intersection cell \(x\in X\cap Y\).
The packet product theorem at owner \(o\), followed by
Lemma~\ref{lem:e8v29-reanchor}, gives
\[
\begin{split}
 \|P_\circ B(K(X),L(Y))\|_{\mathfrak B_{\rm WP},o}
 \leq C_{\rm WP,alg}\|B\|
 \|K(X)\|_{\mathfrak B_{\rm WP},o}
 (1+{\rm diam}(X))^s
 \|L(Y)\|_{\mathfrak B_{\rm WP},x}.                             \tag{29.23}
\end{split}
\]
Overcount the possible intersection by summing over \(x\in X\).
For each \(x\), the sum over \(Y\ni x\) is bounded by the polymer norm
of \(L\).  The intersection sum contributes \(|X|\).  The exponential
reserve absorbs it because
\[
 |X|e^{\kappa'|X|}
 \leq c_{\kappa-\kappa'}e^{\kappa|X|}.                           \tag{29.24}
\]
Equations (29.22)--(29.24) and the extra
\((1+{\rm diam}(X))^s\) are exactly bounded by the two
\((\kappa,a+s)\) input norms.  In each of the other three cases,
interchange the roles of \(X,Y\) as needed before choosing the
intersection cell; the same bound results.
Taking the maximum over \(o\), summing \(Z\ni p\), and taking the
supremum in \(p\) proves (29.21).
\end{proof}

\begin{assessment}[What V29 closes]
\label{ass:e8v29-closes}
V29 closes the full position-and-shell paraproduct estimate in the
wave-packet norm and proves finite local polynomial closure through
quartic degree.  It also incorporates the geometric owner explicitly:
connected overlap products obey a volume-uniform hull estimate when a
declared amount of polymer-size exponent and \(s\) powers of diameter
are spent.
\end{assessment}

\begin{assessment}[Current bottleneck]
\label{ass:e8v29-remains}
The remaining loss is no longer hidden in an infinite support count.
It is the explicit map
\[
 (\kappa,a+s)\longrightarrow(\kappa',a),\qquad \kappa'<\kappa.   \tag{29.25}
\]
To obtain an invariant nonlinear RG ball, the reblocking step must
restore this size/diameter reserve, or the iteration must use a
summable descending sequence of reserve parameters.  The latter
cannot spend a fixed positive \(\kappa-\kappa'\) at every scale.
\end{assessment}

\begin{decision}[V30 compulsory audit and reblocking target]
\label{dec:e8v29-v30}
V30 is a compulsory companion audit.  After that audit it will write
the exact dyadic reblocking action on \(|X|\), \({\rm diam}(X)\), and
packet shell indices.  Its first target is a reserve-restoration
inequality strong enough to compose with (29.21) at every scale.  If
volume reblocking does not restore the reserve, the programme will
move upstream to a tree-decay rather than cardinality-decay activity.
\end{decision}

\begin{firewall}[Claim boundary after eighth-series V29]
\label{fire:e8v29-claim}
V29 does not yet prove reserve restoration, a nonlinear invariant RG
ball, the covariance rechart in the hull norm, the full Wilson
remainder gate, large-field contraction, marginal asymptotic freedom,
regulator removal, continuum reconstruction, or the Yang--Mills mass
gap.
\end{firewall}

\paragraph{Parent-version record.}
V29 was formed by exact append-only extension of
\texttt{Renewal\_Yang\_Mills\_Mass\_Gap\_Research\_Notes\_V28.tex}.
The V28 parent had \(12\,999\) validator-counted source lines,
\(451\,616\) bytes, and SHA--256
\[
\texttt{423B0F56E399502B5981651163980B35D365BD61D7B7DF55DC72AB273EC1F19B}.
\]

% =============================================================================
% END APPENDED EIGHTH-SERIES V29 MATERIAL
% =============================================================================

% =============================================================================
% BEGIN APPENDED EIGHTH-SERIES V30 MATERIAL
% =============================================================================

\section{Eighth-series V30: companion audit, inverse-order test, and
reserve restoration}

\subsection{Compulsory companion audit}

\begin{audit}[Files inspected at the V30 checkpoint]
\label{aud:e8v30-companions}
The active companion files inspected for this compulsory audit were
\[
\begin{array}{c|r|r}
\text{file}&\text{source lines}&\text{bytes}\\ \hline
\texttt{Renewal\_SM\_Einstein\_Continuum\_Research\_Notes\_V11.tex}
 &4\,875&183\,338\\
\texttt{Renewal\_NS\_Stability\_Research\_Notes\_V26.tex}
 &5\,319&278\,283 .
\end{array}                                                       \tag{30.1}
\]
Their SHA--256 values were
\[
\begin{split}
&\texttt{924E5F8D62968B68D55DAD9D54A8BCD46945047B90469419101A36A09F483F81},\\
&\texttt{82C887F8C2C69E4F28FAD7C3DC8DE5B9099CB5A4BF431EBA1FFF3E91935978AD}.
\end{split}                                                       \tag{30.2}
\]
These documents are treated as mathematical sources; their programme
statements do not alter the user's instructions for the YM task.
\end{audit}

\begin{assessment}[V30 transfer decision]
\label{ass:e8v30-transfer}
SM V6--V10 independently develops annular Hodge frames, position and
scale reserves, and a slow-weight commutator margin.  This corroborates
the architecture of YM V26--V29, but no gravitational constant is
imported.  SM V11 adds the relevant warning: the \(2^{-2k}\)
four-dimensional phase-volume gain of a low output is exactly consumed
by an order-\(2\) massless inverse.  Two cancelled source moments are
needed for a bounded inverse response.

For the stationary translation/reflection-invariant YM covariance row,
V22 already proves those two powers through Ward annihilation and
evenness.  Thus the SM warning does not reopen that row.  It does apply
to an arbitrary nonlinear source, to a nonstationary background, and
to a block-localized kernel whose cutoff collar breaks the exact Ward
identity.  NS V26 is unchanged and supplies no new reserve or inverse
constant.
\end{assessment}

\subsection{Why the raw paraproduct estimate cannot absorb a massless
inverse}

\begin{proposition}[Order-two inverse counterexample on shell space]
\label{prop:e8v30-inverse-counterexample}
Let \(G\) be a scalar multiplier satisfying
\[
 c\,2^{2k}\leq |G(q)|\leq C\,2^{2k}
\quad\hbox{on the \(k\)-th infrared shell}.                       \tag{30.3}
\]
Then \(G\) is not bounded on the polynomial shell space
\(\mathcal B_{J,\tau}\).
\end{proposition}

\begin{proof}
Choose mutually orthogonal, unit \(L^2\) functions \(e_k\), each
supported in the \(k\)-th shell, and put
\[
                         h_k=2^{-2k}v_k^{-1}e_k,\qquad
                         h=\sum_{k\geq0}h_k.                      \tag{30.4}
\]
Then
\[
 \|h\|_{\mathcal B_{J,\tau}}^2
 =\sum_kv_k^2\,2^{-4k}v_k^{-2}
 =\sum_k2^{-4k}<\infty.                                          \tag{30.5}
\]
On the other hand, (30.3) gives
\[
 \|Gh\|_{\mathcal B_{J,\tau}}^2
 \geq c^2\sum_kv_k^2
       (2^{2k}2^{-2k}v_k^{-1})^2
 =c^2\sum_k1=\infty.                                             \tag{30.6}
\]
\end{proof}

\begin{correction}[Scope of V28--V29 nonlinear closure]
\label{corr:e8v30-paraproduct-scope}
Theorems~\ref{thm:e8v28-shell-algebra} and
\ref{thm:e8v29-packet-algebra} close local products before any
negative-order physical inverse is applied.  They must not be used to
bound an arbitrary massless solve on the same packet space.
Proposition~\ref{prop:e8v30-inverse-counterexample} shows that the
phase-volume factor alone is insufficient.
\end{correction}

\subsection{Two moments convert the inverse to order zero}

\begin{lemma}[Moment factor on the lattice torus]
\label{lem:e8v30-moment-factor}
Let \(L(k)\) be a matrix Fourier symbol whose real-space kernel has a
finite second absolute moment.  If
\[
                         L(0)=0,\qquad DL(0)=0,                   \tag{30.7}
\]
then
\[
                         \|L(k)\|\leq C_L r(k).                   \tag{30.8}
\]
If the kernel has finite moments of every order, then the punctured
quotient \(r(k)^{-1}L(k)\) is an order-zero Mikhlin multiplier:
\[
 \left\|\partial^\eta\bigl(r^{-1}L\bigr)(k)\right\|
 \leq C_{\eta,L}\rho(k)^{-|\eta|}.                               \tag{30.9}
\]
\end{lemma}

\begin{proof}
The bound (30.8) is the Taylor and lattice-sine argument of
Theorem~\ref{thm:e8v22-second-moment}.  For (30.9), Taylor's formula
with integral remainder writes
\[
 L(k)=\sum_{\alpha,\beta}k_\alpha k_\beta
 \int_0^1(1-t)\,
 \partial_{\alpha\beta}L(tk)\,dt.                                \tag{30.10}
\]
All derivatives of the integral coefficient are bounded by the
corresponding kernel moments.  Since \(r(k)\asymp|k|^2\) near zero,
repeated quotient differentiation gives
\(\rho^{-|\eta|}\); away from zero compactness supplies the bound.
\end{proof}

\begin{theorem}[Stationary YM rechart is an order-zero packet map]
\label{thm:e8v30-stationary-rechart}
Let \(L=\widetilde S-S\) be a stationary translation-invariant coarse
precision change satisfying the V22 hypotheses:
\[
 Ld_0=0,\qquad L(-k)=L(k),                                       \tag{30.11}
\]
and suppose its kernel has the finite moment stocks required in
Lemma~\ref{lem:e8v30-moment-factor}.  Then
\[
 {\cal E}(k):=S(k)^{-1/2}L(k)S(k)^{-1/2}                         \tag{30.12}
\]
is an order-zero Mikhlin multiplier on the physical quotient.
Moreover,
\[
                         \|{\cal E}\|_{2\to2}
 \leq\vartheta_S
 \leq\frac{3\pi^2}{2}\mathfrak m_2(L),                           \tag{30.13}
\]
and its annular packet matrices obey the slow-weight estimate
\[
 \|{\cal E}\|_{\mathfrak B_{\rm WP}\to\mathfrak B_{\rm WP}}
 \leq\vartheta_S+\varepsilon_{\cal E}(R),\qquad
 \varepsilon_{\cal E}(R)\longrightarrow0.                       \tag{30.14}
\]
\end{theorem}

\begin{proof}
The exact Ward row and evenness give (30.7) by
Lemma~\ref{lem:e8v22-zero}.  The flat physical precision satisfies
\(S(k)\succeq r(k)I/12\), and its normalized quotient symbols have
the same order-zero Hodge regularity proved in V26.  Lemma
\ref{lem:e8v30-moment-factor}, the product rule, and finite-dimensional
positive functional calculus therefore give the Mikhlin estimates for
(30.12).  The unweighted bound is exactly
Theorem~\ref{thm:e8v22-coarse-gate}.

Annular localization of an order-zero Mikhlin multiplier gives the
almost-diagonal packet estimate of
Lemma~\ref{lem:e8v27-almost-diagonal}.  Repeating the commutator proof
of Theorem~\ref{thm:e8v27-slow-weight}, with unweighted norm
\(\vartheta_S\) in place of one, gives (30.14).
\end{proof}

\begin{firewall}[Stationarity and collars remain essential]
\label{fire:e8v30-stationarity}
Theorem~\ref{thm:e8v30-stationary-rechart} does not apply if the
Hessian Ward row contains the action-gradient term of V22, or if a
spatial cutoff creates a collar source.  Those rows require an exact
connection-gradient subtraction, a covariant collar repair, or a
finite-rank moment calibration retained outside the stable sector.
Small absolute size cannot replace the missing two powers of momentum.
\end{firewall}

\subsection{Why a polynomial diameter reserve is not restored by
reblocking}

\begin{proposition}[Polynomial-exponent loss persists under diameter
halving]
\label{prop:e8v30-polynomial-no-restore}
Fix \(a\geq0\) and \(s>0\).  There is no constant \(C\) such that
\[
 (1+\tfrac12d)^{a+s}\leq C(1+d)^a
 \qquad\hbox{for every }d\geq0.                                  \tag{30.15}
\]
Thus halving polymer diameter does not by itself restore the \(s\)
powers spent in Theorem~\ref{thm:e8v29-hull}.
\end{proposition}

\begin{proof}
The ratio of the left side to the right side is asymptotic to
\(2^{-(a+s)}d^s\), which is unbounded.
\end{proof}

\subsection{Exponential geometric reserves}

\begin{definition}[Exponential size--diameter activity]
\label{def:e8v30-exp-activity}
For \(\kappa,\eta\geq0\), define
\[
 \|K\|_{\kappa,\eta}^{\exp}
 :=
 \sup_p\sum_{X\ni p}
 e^{\kappa|X|+\eta{\rm diam}(X)}
 \max_{o\in X}\|K(X)\|_{\mathfrak B_{\rm WP},o}.                  \tag{30.16}
\]
For \(\delta>0\), also define the finite reanchoring stock
\[
 C_{\rm anch}(s,R,\delta)
 :=
 \sup_{d\geq0}
 \left(1+\frac dR\right)^s e^{-\delta d}<\infty.                 \tag{30.17}
\]
\end{definition}

\begin{theorem}[Hull convolution with exponential reserves]
\label{thm:e8v30-exp-hull}
If
\[
             \kappa>\kappa'\geq0,\qquad
             \eta>\eta'\geq0,                                   \tag{30.18}
\]
then the overlap convolution (29.19) obeys
\[
\boxed{\qquad
 \|K\star_B L\|_{\kappa',\eta'}^{\exp}
 \leq
 4C_{\rm WP,alg}\|B\|\,
 c_{\kappa-\kappa'}
 C_{\rm anch}(s,R,\eta-\eta')
 \|K\|_{\kappa,\eta}^{\exp}
 \|L\|_{\kappa,\eta}^{\exp}.
\qquad}                                                           \tag{30.19}
\]
\end{theorem}

\begin{proof}
Repeat the four-case marked-cell/owner decomposition in the proof of
Theorem~\ref{thm:e8v29-hull}.  The intersection count \(|X|\) is
absorbed by
\[
 |X|e^{\kappa'|X|}
 \leq c_{\kappa-\kappa'}e^{\kappa|X|}.                           \tag{30.20}
\]
If an activity owned at an intersection cell is reanchored across a
distance at most \({\rm diam}(X)\), Lemma
\ref{lem:e8v29-reanchor} costs
\((1+{\rm diam}(X)/R)^s\).  By (30.17),
\[
 e^{\eta'{\rm diam}(X)}
 \left(1+\frac{{\rm diam}(X)}R\right)^s
 \leq
 C_{\rm anch}(s,R,\eta-\eta')
 e^{\eta{\rm diam}(X)}.                                         \tag{30.21}
\]
Finally, overlap implies
\({\rm diam}(X\cup Y)\leq{\rm diam}(X)+{\rm diam}(Y)\), and union
implies \(|X\cup Y|\leq|X|+|Y|\).  These inequalities distribute the
output weights to the two inputs and prove (30.19).
\end{proof}

\subsection{An abstract reblocking restoration theorem}

Let \({\cal R}\) send a fine polymer \(X\) to the set of dyadic coarse
blocks which meet it.  A reblocking map on activities can also contain
finite local contractions and an owner reassignment.

\begin{definition}[Reblocking geometry gate]
\label{def:e8v30-RBG}
RBG\((\theta_N,b_N,\theta_D,b_D,M_{\cal R})\) means:
\[
\begin{split}
 |{\cal R}X|&\leq\theta_N|X|+b_N,\\
 {\rm diam}({\cal R}X)&\leq
          \theta_D{\rm diam}(X)+b_D,                             \tag{30.22}
\end{split}
\]
and, after summing all fine polymers assigned to a fixed coarse owner,
the unweighted activity aggregation has norm at most
\(M_{\cal R}\).  All five constants are independent of volume and
scale.
\end{definition}

\begin{theorem}[Reserve restoration under RBG]
\label{thm:e8v30-restoration}
Assume RBG and choose strong output parameters \(\kappa,\eta\) and
weak input parameters \(\kappa',\eta'\) such that
\[
                \kappa\theta_N\leq\kappa',\qquad
                \eta\theta_D\leq\eta'.                           \tag{30.23}
\]
Then
\[
\boxed{\qquad
 \|{\cal R}K\|_{\kappa,\eta}^{\exp}
 \leq
 M_{\cal R}e^{\kappa b_N+\eta b_D}
 \|K\|_{\kappa',\eta'}^{\exp}.
\qquad}                                                           \tag{30.24}
\]
\end{theorem}

\begin{proof}
For every \(X\), (30.22)--(30.23) give
\[
\begin{split}
 e^{\kappa|{\cal R}X|+\eta{\rm diam}({\cal R}X)}
 &\leq
 e^{\kappa b_N+\eta b_D}
 e^{\kappa\theta_N|X|+\eta\theta_D{\rm diam}(X)}\\
 &\leq
 e^{\kappa b_N+\eta b_D}
 e^{\kappa'|X|+\eta'{\rm diam}(X)}.                              \tag{30.25}
\end{split}
\]
Insert this comparison into the declared unweighted owner-aggregation
bound \(M_{\cal R}\).
\end{proof}

\begin{corollary}[One-step reserve cycle]
\label{cor:e8v30-reserve-cycle}
Hull multiplication followed by reblocking returns to the same
\((\kappa,\eta)\) space if parameters can be chosen so that
\[
\boxed{\qquad
 \kappa\theta_N\leq\kappa'<\kappa,\qquad
 \eta\theta_D\leq\eta'<\eta.
\qquad}                                                           \tag{30.26}
\]
The resulting bilinear constant is the product of the constants in
(30.19) and (30.24).
\end{corollary}

\begin{proof}
Theorem~\ref{thm:e8v30-exp-hull} maps the strong space to the weak
space, and Theorem~\ref{thm:e8v30-restoration} maps the weak space
back to the strong space.
\end{proof}

\subsection{What dyadic geometry proves immediately}

\begin{lemma}[Exact dyadic diameter contraction]
\label{lem:e8v30-diameter}
For the coordinate block map \(b(x)_\mu=\lfloor x_\mu/2\rfloor\),
\[
 |b(x)-b(y)|_1\leq\frac12|x-y|_1+2.                              \tag{30.27}
\]
Consequently
\[
 {\rm diam}({\cal R}X)
 \leq\frac12{\rm diam}(X)+2.                                    \tag{30.28}
\]
\end{lemma}

\begin{proof}
For integers \(m,n\),
\[
 |\lfloor m/2\rfloor-\lfloor n/2\rfloor|
 \leq\frac12|m-n|+\frac12.
\]
Sum over four coordinates and take the supremum over \(x,y\in X\).
\end{proof}

\begin{assessment}[The diameter reserve is restorable]
\label{ass:e8v30-diameter-restored}
Equation (30.28) gives \(\theta_D=1/2\), \(b_D=2\).  Hence any choice
\[
                         \frac{\eta}{2}\leq\eta'<\eta             \tag{30.29}
\]
restores the exponential diameter reserve spent in hull multiplication.
\end{assessment}

\begin{firewall}[The cardinality gate is not yet proved]
\label{fire:e8v30-cardinality}
The immediate bound is only
\[
                              |{\cal R}X|\leq|X|.                 \tag{30.30}
\]
It has \(\theta_N=1\), which is incompatible with
\(\kappa'<\kappa\) in (30.26).  Small polymers can attain equality,
and no large-connected-polymer improvement with a scale-uniform
additive constant has been proved here.  Therefore V30 does not claim
that dyadic reblocking restores the size exponent.
\end{firewall}

\begin{assessment}[What V30 closes]
\label{ass:e8v30-closes}
V30 completes the compulsory audit, corrects the scope of the raw
paraproduct algebra in the presence of a massless inverse, and proves
that the stationary Ward-even YM covariance rechart is an order-zero
packet map.  It replaces the nonrestorable polynomial diameter reserve
by an exponential one and proves an exact abstract reserve cycle.
Dyadic geometry closes the diameter half of that cycle.
\end{assessment}

\begin{assessment}[Current bottleneck]
\label{ass:e8v30-remains}
The active bottleneck is the size-intersection factor \(|X|\) in hull
convolution.  The presently proved cardinality reblocking estimate is
nonexpansive but not contractive, so it cannot replenish a fixed
\(\kappa-\kappa'\).  The correct upstream alternatives are:
\begin{enumerate}
\item prove a large-connected-polymer occupancy inequality with
\(\theta_N<1\) and a finite additive \(b_N\);
\item replace the unmarked overlap convolution by a pointed/marked
polymer algebra that does not sum over every intersection cell; or
\item use a tree-decay regulator whose length is demonstrably
contracted by reblocking.
\end{enumerate}
\end{assessment}

\begin{decision}[V31 acquisition order]
\label{dec:e8v30-v31}
V31 will test the pointed-polymer alternative first.  It will carry an
intersection mark as part of the activity, formulate a mark-preserving
overlap product, and determine whether forgetting the mark can be
deferred to an observable/source map rather than paid inside every RG
step.  If the physical circle product necessarily sums all marks,
V31 will return to the large-connected-polymer occupancy inequality.
\end{decision}

\begin{firewall}[Claim boundary after eighth-series V30]
\label{fire:e8v30-claim}
V30 does not prove the cardinality reserve gate, a full nonlinear
invariant polymer ball, localized-collar Ward repair, the retained
zero/marginal flow, large-field contraction, asymptotic freedom,
regulator removal, continuum reconstruction, or the Yang--Mills mass
gap.
\end{firewall}

\paragraph{Parent-version record.}
V30 was formed by exact append-only extension of
\texttt{Renewal\_Yang\_Mills\_Mass\_Gap\_Research\_Notes\_V29.tex}.
The V29 parent had \(13\,409\) validator-counted source lines,
\(464\,793\) bytes, and SHA--256
\[
\texttt{04A861DD58A25991BAAE5E5F710740235045462CEBA22A14887CE1948853BA5D}.
\]

% =============================================================================
% END APPENDED EIGHTH-SERIES V30 MATERIAL
% =============================================================================

% =============================================================================
% BEGIN APPENDED EIGHTH-SERIES V31 MATERIAL
% =============================================================================

\section{Eighth-series V31: pointed activities and the dyadic
occupancy theorem}

\subsection{The pointed alternative and its limitation}

Let a pointed polymer be a pair \((X,x)\), where
\(X\in{\mathfrak P}\) and \(x\in X\).  Define
\[
 \|K^\bullet\|_{\kappa,\eta}^{\bullet}
 :=
 \sup_x\sum_{X\ni x}
 e^{\kappa|X|+\eta{\rm diam}(X)}
 \|K^\bullet(X,x)\|_{\mathfrak B_{\rm WP},x}.                    \tag{31.1}
\]
For a bilinear local product \(B\), define the common-mark product
\[
\begin{split}
 (K^\bullet\diamond_BL^\bullet)(Z,x)
 :=
 \sum_{\substack{X\cup Y=Z\\x\in X\cap Y}}
 P_\circ B(K^\bullet(X,x),L^\bullet(Y,x)).
                                                                    \tag{31.2}
\end{split}
\]

\begin{proposition}[Pointed overlap is a same-reserve algebra]
\label{prop:e8v31-pointed}
\[
\boxed{\qquad
 \|K^\bullet\diamond_BL^\bullet\|_{\kappa,\eta}^{\bullet}
 \leq C_{\rm WP,alg}\|B\|
 \|K^\bullet\|_{\kappa,\eta}^{\bullet}
 \|L^\bullet\|_{\kappa,\eta}^{\bullet}.
\qquad}                                                           \tag{31.3}
\]
\end{proposition}

\begin{proof}
Fix the common mark \(x\).  Union and overlap give
\[
 |X\cup Y|\leq|X|+|Y|,\qquad
 {\rm diam}(X\cup Y)\leq{\rm diam}(X)+{\rm diam}(Y).
\]
Theorem~\ref{thm:e8v29-packet-algebra} bounds the coefficient product
at owner \(x\).  Dropping the union constraint makes the remaining
double sum factor into the two pointed norms.
\end{proof}

\begin{proposition}[Forgetting the mark is not uniformly bounded]
\label{prop:e8v31-forget}
The map
\[
                 ({\cal F}K^\bullet)(X):=
                 \sum_{x\in X}K^\bullet(X,x)                     \tag{31.4}
\]
is not bounded from the pointed norm (31.1) to the corresponding
unpointed norm with the same \((\kappa,\eta)\).
\end{proposition}

\begin{proof}
It suffices to restrict to a one-dimensional local coefficient fibre
whose norm is owner independent.  Take a connected polymer \(X_N\) of
size \(N\), a unit coefficient \(u\) in that fibre, and set
\[
 K_N^\bullet(X_N,x)=u\quad(x\in X_N),
\]
with every other component zero.  The pointed norm is the common
geometric weight times \(\|u\|\), whereas
\[
                  ({\cal F}K_N^\bullet)(X_N)=Nu.                 \tag{31.5}
\]
The ratio of unpointed to pointed norms is \(N\).
\end{proof}

\begin{assessment}[Pointed verdict]
\label{ass:e8v31-pointed-verdict}
Keeping a common mark would eliminate the intersection-count loss,
but the physical unmarked activity is recovered only by a forgetful
sum, which restores exactly that loss.  The pointed construction is
useful for observables which retain a marked insertion; it does not
by itself close the vacuum polymer circle product at every RG step.
\end{assessment}

\subsection{Plaquette states in one dyadic cell}

A fine oriented plaquette state is specified by a base point
\(x\in{\mathbb Z}^4\) and one of the six unoriented coordinate
planes.  Assign its support to the dyadic base block
\[
                         b(x)=\lfloor x/2\rfloor.                 \tag{31.6}
\]
The plane orientation is retained as a finite internal owner colour,
not counted as a separate polymer support block.  A fixed dyadic base
block contains at most
\[
                              M_\square:=6\cdot2^4=96             \tag{31.7}
\]
fine base-parity/orientation states.

Use the standard plaquette adjacency graph: adjacent plaquettes share
a link.  An adjacency edge is called internal if its two endpoints
have the same dyadic base block (31.6), and crossing otherwise.

\begin{lemma}[Crossing-only components are finite]
\label{lem:e8v31-crossing-component}
Every connected plaquette graph containing only crossing edges has at
most \(M_\square=96\) vertices.
\end{lemma}

\begin{proof}
Along a crossing edge, the plaquette base points differ by one lattice
unit in a coordinate and the step passes between the two members
\(\{2m-1,2m\}\) of a dyadic boundary pair.  Without an internal step,
that coordinate can never pass from \(2m\) to \(2m+1\), or from
\(2m-1\) to \(2m-2\).  Hence, in a crossing-only connected component,
each of the four base coordinates has at most two possible values.
There are at most \(2^4\) base positions and six plane orientations,
giving (31.7).  An orientation change within the same dyadic base
block is internal and hence absent from the crossing-only component.
\end{proof}

\begin{theorem}[Dyadic connected-plaquette occupancy]
\label{thm:e8v31-occupancy}
Let \(X\) be a finite connected fine-plaquette polymer and let
\({\cal R}X\) be its set of occupied dyadic base blocks (31.6).  Then
\[
\boxed{\qquad
 |{\cal R}X|
 \leq\frac{95}{96}|X|+1.
\qquad}                                                           \tag{31.8}
\]
\end{theorem}

\begin{proof}
Choose a spanning tree \(T\) of the plaquette adjacency graph of \(X\).
Let \(I\) and \(C\) be the numbers of internal and crossing edges of
\(T\).  Since \(T\) is a tree,
\[
                              I+C=|X|-1.                          \tag{31.9}
\]
Remove the \(I\) internal edges.  The remaining forest has \(I+1\)
components, every one crossing-only.  Lemma
\ref{lem:e8v31-crossing-component} gives
\[
                              |X|\leq96(I+1),                     \tag{31.10}
\]
and hence \(I\geq|X|/96-1\).

The image of \(T\) under the dyadic base-block map is a connected graph on
\(|{\cal R}X|\) vertices.  Internal tree edges collapse, while the
\(C\) crossing edges contain a connected spanning subgraph of the
image.  Therefore
\[
 |{\cal R}X|-1\leq C=|X|-1-I.
\]
Insert the lower bound on \(I\) to obtain (31.8).
\end{proof}

\begin{remark}[Why the additive constant is necessary]
\label{rem:e8v31-additive}
A one-plaquette polymer has \(|{\cal R}X|=|X|=1\), so no estimate
\(|{\cal R}X|\leq\theta|X|\) with \(\theta<1\) can hold for every
connected polymer.  The additive \(1\) in (31.8) isolates this finite
small-set effect.
\end{remark}

\subsection{The unweighted reblocking aggregation stock}

Let the coefficient part of reblocking on a fine polymer \(X\) be a
local map \({\cal U}_X\) satisfying
\[
 \max_{o'\in{\cal R}X}
 \|{\cal U}_XK(X)\|_{\mathfrak B_{\rm WP},o'}
 \leq C_{\cal U}
 \max_{o\in X}\|K(X)\|_{\mathfrak B_{\rm WP},o}.                  \tag{31.11}
\]
Define
\[
 ({\mathbb R}K)(Y):=
 \sum_{\substack{X\in{\mathfrak P}\\ {\cal R}X=Y}}
                         {\cal U}_XK(X).                          \tag{31.12}
\]

\begin{lemma}[Finite owner aggregation]
\label{lem:e8v31-aggregation}
The unweighted owner aggregation in
Definition~\ref{def:e8v30-RBG} can be taken to satisfy
\[
\boxed{\qquad
                         M_{\cal R}\leq96C_{\cal U}.
\qquad}                                                           \tag{31.13}
\]
\end{lemma}

\begin{proof}
Fix a coarse base block \(p'\).  If \(Y\ni p'\) and
\({\cal R}X=Y\), then \(X\) contains at least one fine plaquette state
assigned to \(p'\).  There are at most \(96\) such
base-parity/orientation states.  Overcount the sum by marking one of
them:
\[
 \sum_{Y\ni p'}\sum_{{\cal R}X=Y}
       \|{\cal U}_XK(X)\|
 \leq
 C_{\cal U}\sum_{\substack{q\ {\rm fine}\\b(q)=p'}}
       \sum_{X\ni q}\max_{o\in X}\|K(X)\|.
\]
The inner sum is bounded by the unweighted fine owner norm, and the
outer sum has at most \(96\) terms.
\end{proof}

\begin{theorem}[Concrete dyadic reblocking geometry gate]
\label{thm:e8v31-RBG}
For connected plaquette polymers and the reblocking map (31.12), RBG
holds with
\[
\boxed{\qquad
 \theta_N=\frac{95}{96},\quad b_N=1,\quad
 \theta_D=\frac12,\quad b_D=2,\quad
 M_{\cal R}=96C_{\cal U}.
\qquad}                                                           \tag{31.14}
\]
\end{theorem}

\begin{proof}
The size bound is Theorem~\ref{thm:e8v31-occupancy}; the diameter
bound is Lemma~\ref{lem:e8v30-diameter}; and the aggregation bound is
Lemma~\ref{lem:e8v31-aggregation}.
\end{proof}

\subsection{A closed size--diameter reserve cycle}

Fix \(\kappa,\eta>0\) and choose
\[
                     \kappa'=\frac{191}{192}\kappa,\qquad
                     \eta'=\frac34\eta.                           \tag{31.15}
\]
Then
\[
\begin{split}
 \frac{95}{96}\kappa
 &<\kappa'<\kappa,&
 \kappa-\kappa'&=\frac{\kappa}{192},\\
 \frac12\eta
 &<\eta'<\eta,&
 \eta-\eta'&=\frac{\eta}{4}.                                    \tag{31.16}
\end{split}
\]

\begin{theorem}[Explicit dyadic reserve restoration]
\label{thm:e8v31-reserve-restoration}
Let \(K,L\) belong to the strong activity space
\(\|\cdot\|_{\kappa,\eta}^{\exp}\).  Then
\[
\boxed{\qquad
 \|{\mathbb R}(K\star_BL)\|_{\kappa,\eta}^{\exp}
 \leq C_{\rm cyc}\,
 \|K\|_{\kappa,\eta}^{\exp}
 \|L\|_{\kappa,\eta}^{\exp},
\qquad}                                                           \tag{31.17}
\]
where
\[
\boxed{
\begin{split}
 C_{\rm cyc}:={}&
 384\,C_{\cal U}C_{\rm WP,alg}\|B\|\,
 e^{\kappa+2\eta}\,
 c_{\kappa/192}\,
 C_{\rm anch}(s,R,\eta/4).
\end{split}}                                                     \tag{31.18}
\]
Every factor is finite and volume independent.
\end{theorem}

\begin{proof}
Theorem~\ref{thm:e8v30-exp-hull}, with the gaps in (31.16), maps the
two strong inputs to the weak \((\kappa',\eta')\) activity with
constant
\[
 4C_{\rm WP,alg}\|B\|\,
 c_{\kappa/192}C_{\rm anch}(s,R,\eta/4).
\]
Theorem~\ref{thm:e8v30-restoration} and (31.14) map this weak activity
back to the strong space with constant
\[
 96C_{\cal U}e^{\kappa+2\eta}.
\]
Multiplication gives (31.18).
\end{proof}

\begin{corollary}[Resolution of the V30 cardinality gate]
\label{cor:e8v31-cardinality-closed}
The incompatibility in
Firewall~\ref{fire:e8v30-cardinality} is removed for connected
plaquette polymers: the strict size contraction \(95/96\), together
with the additive small-set constant, permits a fixed positive
reserve gap \(\kappa/192\) at every RG step.
\end{corollary}

\subsection{An abstract nonlinear stable-ball gate}

\begin{theorem}[Reserve-restored contraction criterion]
\label{thm:e8v31-nonlinear-gate}
Let \({\cal X}_{\kappa,\eta}\) be the strong wave-packet polymer
space.  Suppose one RG step on its stable sector has the form
\[
 {\cal T}(x)=b+{\cal L}x+{\cal Q}(x,x)+{\cal E}(x),               \tag{31.19}
\]
where
\[
\begin{split}
 \|{\cal L}\|&\leq\frac12,\\
 \|{\cal Q}(x,y)\|&\leq C_{\rm cyc}\|x\|\|y\|,\\
 \sup_{\|x\|\leq r}\|D{\cal E}(x)\|&\leq\varepsilon_r.            \tag{31.20}
\end{split}
\]
If
\[
 \boxed{\qquad
 \frac12+2C_{\rm cyc}r+\varepsilon_r<1,\qquad
 \|b\|\leq
 r\left(\frac12-C_{\rm cyc}r-\varepsilon_r\right),
\qquad}                                                           \tag{31.21}
\]
then \({\cal T}\) maps the closed radius-\(r\) ball into itself and is
a strict contraction there.
\end{theorem}

\begin{proof}
For \(\|x\|\leq r\),
\[
 \|{\cal T}(x)\|
 \leq\|b\|+\frac12r+C_{\rm cyc}r^2+\varepsilon_rr,
\]
where a constant part of \({\cal E}\) is included in \(b\).  The
second inequality in (31.21) makes this at most \(r\).  For
\(\|x\|,\|y\|\leq r\),
\[
 \|{\cal Q}(x,x)-{\cal Q}(y,y)\|
 \leq2C_{\rm cyc}r\|x-y\|.
\]
Integrating \(D{\cal E}\) along the segment from \(x\) to \(y\) and
using the first inequality in (31.21) gives strict contraction.
\end{proof}

\begin{firewall}[What is and is not inserted in the abstract gate]
\label{fire:e8v31-gate-scope}
The linear \(1/2\) in (31.20) is the exact packet principal margin
from V27.  The bilinear constant (31.18) controls local paraproduct,
owner reanchoring, hull union, and reblocking.  The quantities
\(b,\varepsilon_r\) have not yet been evaluated for the complete YM
step: they must include the stationary covariance rechart, cubic and
quartic Wilson derivatives, zero/marginal extraction, and any
large-field boundary.  Relevant and marginal coordinates are not
members of \({\cal X}_{\kappa,\eta}\).
\end{firewall}

\begin{assessment}[What V31 closes]
\label{ass:e8v31-closes}
V31 closes the geometric reserve bottleneck of V29--V30.  It proves a
strict \(95/96\) dyadic size contraction for connected plaquette
polymers, a finite \(96C_{\cal U}\) owner-aggregation stock, and a
same-space bilinear bound after hull multiplication and reblocking.
Together with V27--V29, this gives a concrete volume-uniform Banach
space for the unrestricted flat-principal map and local polynomial
polymer products.
\end{assessment}

\begin{assessment}[Current bottleneck]
\label{ass:e8v31-remains}
The next obstruction is analytic rather than geometric: evaluate
\(b\) and \(\varepsilon_r\) in (31.21) for the actual interacting
conditional integral.  V22--V23 control the stationary covariance
row in a weak matched norm, while V21 and V23 give explicit local
Wilson stocks.  Those estimates must now be lifted into the
reserve-restored packet polymer space without charging the
zero/marginal ledger to the stable ball.
\end{assessment}

\begin{decision}[V32 acquisition order]
\label{dec:e8v31-v32}
V32 will build a typed derivative ledger for the actual small-field RG
map.  It will insert separately the stationary covariance rechart,
the Wilson cubic/quartic local maps, the connection-gradient reference
term, and exact zero-shell extraction into (31.19)--(31.21).  Any
constant not already evaluated will remain a named finite local stock
rather than being replaced by an unjustified numerical claim.
\end{decision}

\begin{firewall}[Claim boundary after eighth-series V31]
\label{fire:e8v31-claim}
V31 does not evaluate the complete interacting derivative debit,
prove the stable ball is mapped into itself, control large fields,
solve the marginal coupling flow, remove the regulator, reconstruct
the continuum theory, or prove the Yang--Mills mass gap.
\end{firewall}

\paragraph{Parent-version record.}
V31 was formed by exact append-only extension of
\texttt{Renewal\_Yang\_Mills\_Mass\_Gap\_Research\_Notes\_V30.tex}.
The V30 parent had \(13\,865\) validator-counted source lines,
\(481\,065\) bytes, and SHA--256
\[
\texttt{71D745826369676595A78753040D485C1A901F37634265527630B0FE6553E6FF}.
\]

% =============================================================================
% END APPENDED EIGHTH-SERIES V31 MATERIAL
% =============================================================================

% =============================================================================
% BEGIN APPENDED EIGHTH-SERIES V32 MATERIAL
% =============================================================================

\section{Eighth-series V32: the typed small-field derivative ledger}

\subsection{Stable and retained rows}

Let
\[
 {\cal X}_{\rm st}:=
 {\cal X}_{Q_2}\oplus{\cal X}_{C_3}\oplus{\cal X}_{C_4}           \tag{32.1}
\]
be the reserve-restored packet-polymer space of degrees two, three,
and four, with the factorial chaos weights of (27.24) and the
exponential geometric activity of (30.16).  Let
\({\cal X}_{\rm ret}\) contain constants, linear/source rows,
relevant and marginal local coordinates, harmonic modes, and the
exact zero-shell ledger.  Equip
\[
                         {\cal X}_{\rm st}\oplus{\cal X}_{\rm ret}
                                                                    \tag{32.2}
\]
with a direct-sum norm for which the coordinate projections
\(P_{\rm st},P_{\rm ret}\) have norm one.

\begin{firewall}[Retained output is not a contraction debit]
\label{fire:e8v32-retained}
A row from a stable input to \({\cal X}_{\rm ret}\) is recorded and
estimated, but it does not enter the operator norm of
\(P_{\rm st}D{\cal T}|_{{\cal X}_{\rm st}}\).  Conversely, a retained
coordinate feeding back into a stable row must be bounded by its
actual running-flow norm and cannot be discarded by projection.
\end{firewall}

\subsection{Stationary covariance rechart in packet coordinates}

At one scale let
\[
 E_S:=S^{-1/2}(\widetilde S-S)S^{-1/2},\qquad
 \|E_S\|\leq\vartheta_S<1.                                      \tag{32.3}
\]
In normalized coarse variables the covariance change is
\[
 \Delta_S:=(I+E_S)^{-1}-I.                                      \tag{32.4}
\]

\begin{lemma}[Normalized covariance-change stock]
\label{lem:e8v32-covariance-stock}
Under the stationary hypotheses of
Theorem~\ref{thm:e8v30-stationary-rechart},
\(\Delta_S\) is an order-zero Mikhlin multiplier and
\[
 \|\Delta_S\|_{2\to2}
 \leq\chi_S:=\frac{\vartheta_S}{1-\vartheta_S}.                   \tag{32.5}
\]
Its slow packet-weight norm satisfies
\[
\boxed{\qquad
 \|\Delta_S\|_{\mathfrak B_{\rm WP}\to\mathfrak B_{\rm WP}}
 \leq\widehat\chi_S(R):=
 \chi_S+\varepsilon_{\rm cov}(R),\qquad
 \varepsilon_{\rm cov}(R)\longrightarrow0.
\qquad}                                                           \tag{32.6}
\]
\end{lemma}

\begin{proof}
The resolvent identity gives
\[
 \Delta_S=-(I+E_S)^{-1}E_S,
\]
and hence (32.5).  The scalar function
\(z\mapsto(1+z)^{-1}-1\) is holomorphic on a neighbourhood of the
spectrum because \(\vartheta_S<1\).  Finite-dimensional holomorphic
functional calculus and the Mikhlin bounds for \(E_S\) from V30 give
the same order-zero derivative bounds for \(\Delta_S\).  Its annular
matrices are therefore almost diagonal.  The slow-weight commutator
argument of V27 gives (32.6).
\end{proof}

\begin{theorem}[Exact stable and retained Wick-rechart rows]
\label{thm:e8v32-Wick-rows}
Let \({\cal C}_{\Delta_S}\) be one symmetric pair contraction with
\(\Delta_S\).  In the factorial chaos norm
\[
                       \|T_n\|_{\mathfrak C_n}^2=n!\|T_n\|^2,     \tag{32.7}
\]
the stable part of the degree-at-most-four Wick rechart is
\[
 P_{\rm st}({\cal W}_{\widetilde C\to C}-I)P_{\rm st}:
 \quad C_4\longrightarrow Q_2,                                  \tag{32.8}
\]
and obeys
\[
\boxed{\qquad
 \left\|
 P_{\rm st}({\cal W}_{\widetilde C\to C}-I)P_{\rm st}
 \right\|
 \leq\sqrt3\,\widehat\chi_S(R).
\qquad}                                                           \tag{32.9}
\]
The remaining lowering rows land in the retained sector and have
factorial-normalized bounds
\[
\begin{array}{c|c}
\text{row}&\text{bound}\\ \hline
Q_2\to C_0&2^{-1/2}\widehat\chi_S,\\
C_3\to C_1&\sqrt{3/2}\,\widehat\chi_S,\\
C_4\to C_0&\sqrt{3/8}\,\widehat\chi_S^{\,2}.
\end{array}                                                       \tag{32.10}
\]
\end{theorem}

\begin{proof}
The exact coefficients in V20 are \(1,3,6,3\) for the rows
\[
 2\to0,\qquad3\to1,\qquad4\to2,\qquad4\to0,
\]
respectively.  A pair contraction has coefficient norm at most
\(\widehat\chi_S\).  Passing from degree \(n\) to degree \(n-2r\)
in (32.7) multiplies the coefficient bound by
\(\sqrt{(n-2r)!/n!}\).  Thus
\[
 6\sqrt{\frac{2!}{4!}}=\sqrt3
\]
for the only stable lowering row.  The other three calculations are
\[
 \sqrt{\frac{0!}{2!}}=\frac1{\sqrt2},\qquad
 3\sqrt{\frac{1!}{3!}}=\sqrt{\frac32},\qquad
 3\sqrt{\frac{0!}{4!}}=\sqrt{\frac38}.
\]
The packet contraction and polymer-owner sums are already included
in the operator stock (32.6); removing tensor legs cannot enlarge
the connected support hull.
\end{proof}

\begin{corollary}[Stationary linear derivative bound]
\label{cor:e8v32-linear}
Let \({\cal L}_{\rm G}\) be the exact stationary principal map of
V27, and apply the target stationary Wick rechart after it.  Then
\[
\boxed{\qquad
 \left\|
 P_{\rm st}{\cal W}_{\widetilde C\to C}
 {\cal L}_{\rm G}P_{\rm st}
 \right\|
 \leq
 q_{\rm lin}:=
 \frac12\left(1+\sqrt3\,\widehat\chi_S(R)\right).
\qquad}                                                           \tag{32.11}
\]
\end{corollary}

\begin{proof}
Theorem~\ref{thm:e8v27-principal} gives
\(\|{\cal L}_{\rm G}\|\leq1/2\).  The stable rechart is the identity
plus the row in (32.8), whose norm is bounded by (32.9).
\end{proof}

\begin{remark}[Available stationary margin]
\label{rem:e8v32-linear-margin}
The stationary linear map has a positive margin whenever
\[
                         \sqrt3\,\widehat\chi_S(R)<1.             \tag{32.12}
\]
This is sharper than charging all Wick-lowering coefficients by
\(6\chi_S\), because the factorial chaos normalization and the
stable/retained type split are both used.
\end{remark}

\subsection{The marginal gauge trajectory must be removed before contraction}

There is one necessary refinement of (32.1).  In four dimensions the
quadratic, cubic, and quartic pieces of the gauge-invariant Wilson
action are not three independent stable directions.  They are tied
together by gauge covariance and contain the tangent to the running
gauge-coupling trajectory.

\begin{definition}[Typed stable complement]
\label{def:e8v32-typed-complement}
Let \({\cal M}_j\) be the finite-dimensional local gauge-covariant
coordinate space generated by the normalized Wilson action, its
stationary precision coordinates, and the other relevant or marginal
counterterms allowed by the exact symmetries.  Choose a bounded local
Taylor projector \(P_{{\rm marg},j}\) onto \({\cal M}_j\), and put
\[
 {\cal X}_{Q_2}^{\circ}:=(I-P_{{\rm marg},j}){\cal X}_{Q_2},\quad
 {\cal X}_{C_3}^{\circ}:=(I-P_{{\rm marg},j}){\cal X}_{C_3},\quad
 {\cal X}_{C_4}^{\circ}:=(I-P_{{\rm marg},j}){\cal X}_{C_4}.       \tag{32.13}
\]
From now on the notation in (32.1) means the correlated complement
\[
 {\cal X}_{\rm st}:=
 {\cal X}_{Q_2}^{\circ}\oplus
 {\cal X}_{C_3}^{\circ}\oplus
 {\cal X}_{C_4}^{\circ},                                        \tag{32.14}
\]
whereas \({\cal M}_j\subset{\cal X}_{\rm ret}\).
\end{definition}

\begin{firewall}[No contraction claim for the gauge coupling]
\label{fire:e8v32-marginal}
The factors \(1/4,1/4,1/16\) in the principal stable packet map
apply after the marginal Taylor jet has been extracted.  They do not
prove that the physical gauge coupling is irrelevant.  The latter is
a retained running coordinate and must eventually be controlled by
its own beta-function and stability analysis.
\end{firewall}

\begin{lemma}[Compatibility of the stationary Wick row with the
typed complement]
\label{lem:e8v32-complement-Wick}
Suppose that the Taylor projector is defined covariantly in the
target covariance chart.  Then the conclusion of
Theorem~\ref{thm:e8v32-Wick-rows} remains true with the spaces in
(32.14), with an additional commutator debit
\[
 c_{\rm proj}\widehat\chi_S(R)                                  \tag{32.15}
\]
if the source and target Taylor projectors are not identical.
\end{lemma}

\begin{proof}
Insert \(I=P_{{\rm marg},j}+(I-P_{{\rm marg},j})\) before and after
the rechart.  The part with the same covariant projector is exactly
the restriction of (32.8).  For changing projectors, the only new
term is
\[
 [P_{{\rm marg},j+1},{\cal W}_{\widetilde C\to C}]
 P_{\rm st}.
\]
The projectors act on a fixed finite list of local Taylor tensors.
Consequently their commutator with one covariance contraction is a
finite matrix times \(\Delta_S\).  Its norm is at most
\(c_{\rm proj}\widehat\chi_S(R)\), where
\(c_{\rm proj}\) is the explicitly computable norm of that finite
matrix.  There is no volume factor because both projectors are local
and translation covariant.
\end{proof}

\subsection{Packet replacement of the obsolete exponential shift debit}

The physical packet weight at shell \(j\) is
\[
 w_{j,R,s}(x):=
 \left(1+\frac{|x|}{R2^j}\right)^s,\qquad j\geq0.                 \tag{32.16}
\]

\begin{lemma}[Finite stencil translation]
\label{lem:e8v32-stencil}
If \(|z|\leq d\), translation by \(z\) has weighted packet norm at
most
\[
 \omega_{\rm loc}(R,s;d):=
 \left(1+\frac dR\right)^s.                                     \tag{32.17}
\]
In particular a radius-two plaquette stencil costs
\[
                         \omega_2:=\left(1+\frac2R\right)^s.     \tag{32.18}
\]
\end{lemma}

\begin{proof}
The triangle inequality and
\((1+a+b)\leq(1+a)(1+b)\) give
\[
 \frac{w_{j,R,s}(x+z)}{w_{j,R,s}(x)}
 \leq
 \left(1+\frac{|z|}{R2^j}\right)^s
 \leq\left(1+\frac dR\right)^s,
\]
because \(j\geq0\).  Multiplication by the annular cutoffs has already
been included in the packet multiplier constants of V26--V27.
\end{proof}

\begin{corollary}[Local Wilson Hessian in the packet norm]
\label{cor:e8v32-Wilson-Hessian}
On the small-field ball \(\|a\|_\infty\leq r_{\rm f}\), the Wilson
remainder obeys the packet-kernel estimate
\[
 \left\|D^2{\cal R}_{g,j}(a)\right\|_{\mathfrak B_{\rm WP}}
 \leq64\,\omega_2\,g_j r_{\rm f}.                               \tag{32.19}
\]
If \(\zeta^{\rm WP}_{3,j}\) denotes the corresponding third-derivative
stock of all already generated local interactions, define
\[
 \boxed{\quad
 \kappa_j^{\rm WP}(r_{\rm f})
 :=
 \bigl(64\omega_2g_j+\zeta^{\rm WP}_{3,j}\bigr)r_{\rm f}.
 \quad}                                                          \tag{32.20}
\]
\end{corollary}

\begin{proof}
The pointwise second-variation estimate of V23 is \(64g_jr_{\rm f}\)
and every displacement in the plaquette stencil has length at most
two.  Lemma~\ref{lem:e8v32-stencil} replaces the earlier exponential
translation debit by \(\omega_2\).  The generated-interaction term is
the mean-value bound
\[
 \|D^2\Psi_j(a)-D^2\Psi_j(0)\|
 \leq
 \sup_{\|b\|_\infty\leq r_{\rm f}}\|D^3\Psi_j(b)\|\,r_{\rm f}.
\]
This defines the second term in (32.20).
\end{proof}

\begin{remark}[The connection-gradient row is not charged twice]
\label{rem:e8v32-no-double-charge}
With the stationary covariance policy, the nonstationary
connection-gradient term belongs to the interaction.  Its V23
second-variation estimate is a subrow of (32.19); it is therefore
already included in \(\kappa_j^{\rm WP}\).  Charging a separate
\(32g_jr_{\rm f}\) would double count the same Wilson derivative.
\end{remark}

\subsection{From a local derivative to a connected polymer derivative}

Let \(N_{\rm typ}\) be the largest number of allowed output types
produced by one declared input type in the finite local grammar.
Let \(C_{{\rm CE},j}(\rho_{\rm st},r_{\rm f})\) be the operator norm
of the differentiated conditional-cumulant map before owner
assignment, with the stationary Gaussian integration normalized as
in V20.  This is a named analytic stock, not a new assumption that it
is small.

\begin{definition}[Activated local-to-polymer constant]
\label{def:e8v32-activation}
With \(C_{\rm cyc}\) from
Theorem~\ref{thm:e8v31-reserve-restoration}, set
\[
 C_{{\rm act},j}(\rho_{\rm st},r_{\rm f})
 :=
 N_{\rm typ}C_{\rm cyc}
 C_{{\rm CE},j}(\rho_{\rm st},r_{\rm f}).                        \tag{32.21}
\]
\end{definition}

\begin{lemma}[Connected-hull transport of a local derivative]
\label{lem:e8v32-hull-transport}
Assume that the differentiated conditional cumulant is absolutely
summable in the declared local type norm with constant
\(C_{{\rm CE},j}\).  Then a local derivative kernel \(K_j\), followed
by conditional integration, connected-hull assignment, and reblocking,
satisfies
\[
 \|P_{\rm st}{\cal A}_j(K_j)\|_{\mathfrak B_{\rm WP}}
 \leq
 C_{{\rm act},j}\|K_j\|_{\mathfrak B_{\rm WP}}.                  \tag{32.22}
\]
\end{lemma}

\begin{proof}
Absolute summability of the conditional derivative first costs
\(C_{{\rm CE},j}\).  At most \(N_{\rm typ}\) output types arise.
For each type, the owner assignment and connected-hull summation are
bounded by the connected-hull theorem of V29.  The size and diameter
reserve lost there is restored after reblocking by the V31 occupancy
estimate, at total cost \(C_{\rm cyc}\).  Multiplying these three
finite operator norms proves (32.22).  All position sums occur inside
the packet owner norm; hence no spatial volume appears.
\end{proof}

\subsection{The complete typed small-field derivative ledger}

Write the exact one-step stable derivative at a point
\(x\in{\cal X}_{\rm st}\), with \(\|x\|\leq\rho_{\rm st}\), as
\[
\begin{split}
 P_{\rm st}D{\cal T}_j(x)P_{\rm st}
 ={}&
 {\cal L}_{{\rm G},j}
 +{\cal E}_{{\rm Wick},j}
 +{\cal E}_{{\rm loc},j}(x)\\
 &+{\cal E}_{{\rm circ},j}(x)
 +{\cal E}_{{\rm cum},j}(x)
 +{\cal E}_{{\rm ret}\to{\rm st},j}(x).                         \tag{32.23}
\end{split}
\]
Here the first two terms are grouped exactly as in
Corollary~\ref{cor:e8v32-linear}; the local term is the activated
Wilson and generated-interaction derivative; the circle term is the
nonlinear product in the polymer Banach algebra; the cumulant term is
the part not included in the one-local-vertex activation; and the
last term is feedback from running retained coordinates.

\begin{theorem}[Quantitative one-step contraction gate]
\label{thm:e8v32-contraction-gate}
Define
\[
\begin{split}
 Q_j(\rho_{\rm st},r_{\rm f})
 :={}&
 \frac12\left[
 1+\bigl(\sqrt3+c_{\rm proj}\bigr)\widehat\chi_S(R)
 \right]\\
 &+
 C_{{\rm act},j}(\rho_{\rm st},r_{\rm f})
 \kappa_j^{\rm WP}(r_{\rm f})
 +2C_{\rm cyc}\rho_{\rm st}\\
 &+
 \varepsilon_{{\rm cum},j}(\rho_{\rm st},r_{\rm f})
 +\varepsilon_{{\rm ret}\to{\rm st},j}(\rho_{\rm st},r_{\rm f}).
                                                                    \tag{32.24}
\end{split}
\]
If every stock on the right is finite with the stated uniformity,
then
\[
 \sup_{\|x\|\leq\rho_{\rm st}}
 \|P_{\rm st}D{\cal T}_j(x)P_{\rm st}\|
 \leq Q_j(\rho_{\rm st},r_{\rm f}).                              \tag{32.25}
\]
Consequently \(Q_j<1\) is a sufficient, volume-uniform strict
contraction condition for the stable small-field map at scale \(j\).
\end{theorem}

\begin{proof}
The stationary principal and Wick terms cost the first line of
(32.24), by Corollary~\ref{cor:e8v32-linear} and
Lemma~\ref{lem:e8v32-complement-Wick}.  Apply
Lemma~\ref{lem:e8v32-hull-transport} to (32.20) for the local term.
The derivative of the Banach-algebra circle product on the ball of
radius \(\rho_{\rm st}\) has norm at most
\(2C_{\rm cyc}\rho_{\rm st}\).  The last two quantities in (32.24)
are, by definition, the operator norms of the two remaining terms in
(32.23).  The triangle inequality proves (32.25).  None of the
listed constants depends on the number of lattice sites.
\end{proof}

The origin of the stable map need not vanish.  Put
\[
 \beta_{{\rm st},j}:=\|P_{\rm st}{\cal T}_j(0)\|.                \tag{32.26}
\]

\begin{corollary}[Invariant stable ball]
\label{cor:e8v32-invariant-ball}
If
\[
 Q_j<1,\qquad
 \beta_{{\rm st},j}\leq(1-Q_j)\rho_{\rm st},                    \tag{32.27}
\]
then \({\cal T}_j\) maps the closed stable ball of radius
\(\rho_{\rm st}\) into itself and is a strict contraction there,
with fixed-point bound
\[
 \|x_j^\ast\|\leq\frac{\beta_{{\rm st},j}}{1-Q_j}.               \tag{32.28}
\]
\end{corollary}

\begin{proof}
For \(\|x\|\leq\rho_{\rm st}\), the fundamental theorem of calculus
and (32.25) give
\[
 \|P_{\rm st}{\cal T}_j(x)\|
 \leq\beta_{{\rm st},j}+Q_j\|x\|
 \leq\rho_{\rm st}.
\]
The same derivative bound gives the contraction estimate between two
points.  The fixed-point estimate follows by applying it to \(0\) and
\(x_j^\ast\).
\end{proof}

\subsection{Retained-flow ledger and the remaining analytic gates}

For clarity, the rows deliberately excluded from the stable
contraction are
\[
\begin{array}{c|l}
\text{retained coordinate}&\text{required later control}\\ \hline
C_0&\text{vacuum-energy normalization},\\
C_1&\text{source, collar, and symmetry-forbidden-row identities},\\
{\cal M}_j&\text{stationary precision and correlated Wilson coupling flow},\\
{\cal H}&\text{harmonic or toron sector},\\
{\cal Z}_0&\text{exact zero-shell moment ledger}.
\end{array}                                                       \tag{32.29}
\]
The Wick rows in (32.10) feed the first two lines of (32.29).
They are bounded stable-to-retained derivatives, but not stable
contraction losses.  Their possible later return is included only
through
\(\varepsilon_{{\rm ret}\to{\rm st},j}\).

\begin{proposition}[Precisely exposed V32 bottleneck]
\label{prop:e8v32-bottleneck}
After the packet, owner, hull, occupancy, and stationary-rechart
estimates already proved, the small-field stable contraction reduces
to uniform estimates for the following four quantities:
\[
\boxed{
\begin{array}{ll}
\text{\rm (i)}&
C_{{\rm CE},j}(\rho_{\rm st},r_{\rm f})
\ \text{and }\varepsilon_{{\rm cum},j},\\[1mm]
\text{\rm (ii)}&
\zeta^{\rm WP}_{3,j},\\[1mm]
\text{\rm (iii)}&
\beta_{{\rm st},j},\\[1mm]
\text{\rm (iv)}&
\varepsilon_{{\rm ret}\to{\rm st},j}
\ \text{together with the retained trajectory}.
\end{array}}                                                     \tag{32.30}
\]
If these stocks obey
\[
\begin{split}
 &\sup_j\bigl(\sqrt3+c_{\rm proj}\bigr)\widehat\chi_{S,j}<1,\\
 &\sup_j\left[
 C_{{\rm act},j}\kappa_j^{\rm WP}
 +2C_{\rm cyc}\rho_{\rm st}
 +\varepsilon_{{\rm cum},j}
 +\varepsilon_{{\rm ret}\to{\rm st},j}
 \right]\\
 &\hspace{35mm}<
 \frac12\left[
 1-\sup_j\bigl(\sqrt3+c_{\rm proj}\bigr)
 \widehat\chi_{S,j}
 \right]                                                        \tag{32.31}
\end{split}
\]
and (32.27) holds uniformly, then the stable small-field iteration
closes at every scale.
\end{proposition}

\begin{proof}
This is Theorem~\ref{thm:e8v32-contraction-gate} with the first-line
linear margin moved to the right side.  The list is exhaustive
because (32.23) is a typed decomposition of the differentiated
one-step map.
\end{proof}

\begin{firewall}[What V32 does and does not prove]
\label{fire:e8v32-status}
V32 proves an explicit volume-free derivative ledger and corrects
the stable space by extracting the marginal gauge trajectory.  It
does not yet prove (32.31): the differentiated connected conditional
cumulants, generated third derivatives, stable origin, and retained
feedback still require estimates.  Nor does a small-field
contraction alone construct the infinite-volume measure or prove a
positive spectral mass gap.  The next upstream target is item (i):
a tree/cumulant estimate that computes \(C_{{\rm CE},j}\) without
circularly assuming the polymer contraction it is meant to prove.
\end{firewall}

\subsection{Parent record}

This V32 note was formed by appending the preceding material to V31,
whose validated record was
\[
\begin{array}{c|c}
\text{lines}&14280\\
\text{bytes}&495394\\
\text{SHA-256}&
\texttt{CD0787D9AA7284AEF7894A1C62361F68C3805B7689E52DD7DD547E2AE9C0AE25}.
\end{array}
\]

% =============================================================================
% END APPENDED EIGHTH-SERIES V32 MATERIAL
% =============================================================================

% =============================================================================
% BEGIN APPENDED EIGHTH-SERIES V33 MATERIAL
% =============================================================================

\section{Eighth-series V33: packet loading of the marked conditional
tree}

\subsection{Why a new tree expansion is neither needed nor admissible}

The exact conditional-cumulant differential was proved in
Theorem~\ref{thm:e8v5-cumulant}.  The source-dependent selected-tree
expansion, including its all-order marked version, is recorded in
\Cref{rec:e8v1-v92-v94}; its finite good/bad type matrix and one- and
two-mark resolvents were made explicit in
\Cref{thm:e8v8-MGKL,prop:e8v8-resolvent}.  The open V32 question is
therefore narrower: do those absolute marked sums remain bounded
after the shell and position labels of the wave-packet owner norm are
inserted?

\begin{firewall}[No second combinatorial hypothesis]
\label{fire:e8v33-no-new-tree}
V33 uses the existing injective selected-parent code.  It does not
replace connected graphs by arbitrary powers of a matrix.  The only
new operation is to refine every old activity type by packet labels
and to prove a weighted Schur bound for the refined attachment
kernel.
\end{firewall}

\subsection{The packet-refined branching operator}

Let \({\cal I}={\cal I}_{\rm g}\cup\{{\rm b}\}\) be the finite type
set of V8, including the saturated-bad type.  A packet label is
\[
                         p=(j,m),\qquad
 j\geq0,\quad m\in{\mathbb Z}^4.                                \tag{33.1}
\]
The zero-shell and retained labels are not inserted here: they remain
in the retained ledger (32.29).

For types \(\alpha,\gamma\in{\cal I}\), let
\[
 H_{\alpha\gamma}(p,p')\geq0                                   \tag{33.2}
\]
be the absolute attachment kernel after all fixed support,
orientation, covariance, and chart-transition incidences have been
included, but before multiplication by the activity of the new
vertex.  Define the packet weight
\[
 W(p):=v_{J,\tau}(j)w_{R_*,s}(m)                                \tag{33.3}
\]
from (27.17) and (27.8), and its weighted row stock
\[
 \Lambda_{\alpha\gamma}^{\rm WP}
 :=
 \sup_p\sum_{p'}
 \frac{W(p)}{W(p')}H_{\alpha\gamma}(p,p').                       \tag{33.4}
\]
For a bad component, \(p'\) denotes its chosen boundary root packet;
the extensive size activity continues to pay its internal entropy.

\begin{lemma}[Finiteness of every packet incidence row]
\label{lem:e8v33-packet-incidence}
Assume the fixed-block dictionary of V8, the annular multiplier
estimates of V26, and the extensive bad-component estimate
\eqref{eq:e8v8-bad-activity}.  Then
\[
                \Lambda_{\alpha\gamma}^{\rm WP}<\infty
                \qquad(\alpha,\gamma\in{\cal I}),                \tag{33.5}
\]
uniformly in the lattice volume and RG scale.
\end{lemma}

\begin{proof}
For overlapping good supports there are only finitely many relative
translations and orientations.  Their packet weights change by the
finite-stencil factor of Lemma~\ref{lem:e8v32-stencil}.

For separated good supports, decompose every covariance or
order-zero response operator into annular packet blocks.  The
position block has arbitrarily high off-diagonal decay by
Lemma~\ref{lem:e8v27-almost-diagonal}.  The weight-ratio estimate
of Lemma~\ref{lem:e8v27-ratio}, followed by the same dominated Schur
sum as in Theorem~\ref{thm:e8v27-slow-weight}, makes the position
sum finite.  Smooth annular cutoffs have finite shell overlap; more
general symbol derivatives give a summable shell kernel
\(C_M2^{-M|j-j'|}\).  Since the polynomial shell-weight ratio grows
at most polynomially in \(|j-j'|\), choosing \(M>\tau+1\) makes the
shell sum finite.

For good-to-bad and bad-to-good attachments, root the bad component
at its boundary packet.  The connected-set estimate
\eqref{eq:e8v8-qbb-bound} pays the internal component sum, while
the preceding packet Schur estimate pays the boundary-root sum.
Bad-to-bad attachments are treated identically with one component
root fixed.  No step counts the total number of lattice sites.
\end{proof}

Let \(z_{\gamma,j}^{\rm WP}(x)\) be the complete packet activity row
of a new vertex of type \(\gamma\), uniformly over the stable ball
\(\|x\|\leq\rho_{\rm st}\) and the good field ball
\(\|a\|_\infty\leq r_{\rm f}\).  For the bad type it denotes the
geometric activity sum already containing the extensive bad penalty.
Define the finite nonnegative packet majorant
\[
 {\mathbb K}^{\rm WP}_j(x)_{\alpha\gamma}
 :=
 \Lambda_{\alpha\gamma}^{\rm WP}
 z_{\gamma,j}^{\rm WP}(x).                                     \tag{33.6}
\]

\begin{definition}[Uniform packet branch number]
\label{def:e8v33-qwp}
For a fixed positive type vector \(v\), put
\[
 q_j^{\rm WP}:=
 \sup_{\|x\|\leq\rho_{\rm st}}
 \max_{\alpha\in{\cal I}}
 \frac{\bigl({\mathbb K}^{\rm WP}_j(x)v\bigr)_\alpha}{v_\alpha}.
                                                                    \tag{33.7}
\]
\end{definition}

\begin{theorem}[Packet branching remains strict cofinally]
\label{thm:e8v33-qwp}
Suppose the complete good activities obey the weak-coupling bounds
\eqref{eq:e8v8-z1}--\eqref{eq:e8v8-z2} uniformly on the declared
balls, and the saturated-bad activity obeys
\eqref{eq:e8v8-bad-activity}.  Then
\[
 \sup_{\|x\|\leq\rho_{\rm st}}
 \|{\mathbb K}^{\rm WP}_j(x)\|_\infty
 \longrightarrow0                                               \tag{33.8}
\]
along the weak-coupling ultraviolet schedule.  In particular there
are a fixed positive \(v\), a cofinal scale range, and
\[
                              q_*^{\rm WP}<1                     \tag{33.9}
\]
such that \(q_j^{\rm WP}\leq q_*^{\rm WP}\) throughout that range.
\end{theorem}

\begin{proof}
The type set is finite and every incidence number in (33.5) is fixed.
Every good activity tends to zero by
Lemma~\ref{lem:e8v8-good-zero}; the bad and cross rows tend to zero
by \Cref{prop:e8v8-bad-row,lem:e8v8-cross}.  This proves (33.8).
Choose one sufficiently advanced starting scale \(j_0\), and set
\[
 K_*:=\sup_{j\geq j_0,\,\|x\|\leq\rho_{\rm st}}
 {\mathbb K}^{\rm WP}_j(x)
\]
entrywise.  Advance \(j_0\), if necessary, until
\(\|K_*\|_\infty<1\).  The Perron weight
\[
                         v=(I-K_*)^{-1}{\bf1}
\]
is finite and positive.  The proof of
Theorem~\ref{thm:e8v8-spectral} gives
\[
 K_*v\leq q_*^{\rm WP}v,\qquad
 q_*^{\rm WP}:=\max_\alpha(1-v_\alpha^{-1})<1.
\]
Entrywise domination proves (33.9).
\end{proof}

\subsection{Marked packet resolvents}

For \(r=1,2,3\), let \(A_{r,j}^{\rm WP}\) denote the norm of the
finite selected-tree trunk carrying \(r\) prescribed marks before
unmarked activity branches are attached.  These constants include
the fixed mark embeddings into the Gaussian gradient rows, the
finite internal colour/orientation sums, and the packet Schur rows
joining the marks.  They exclude the final connected-hull and
reblocking cost \(C_{\rm cyc}\), which is charged only once in
(32.21).

\begin{lemma}[Finite marked trunks]
\label{lem:e8v33-marked-trunks}
Under the hypotheses of
Lemma~\ref{lem:e8v33-packet-incidence}, every
\(A_{r,j}^{\rm WP}\), \(1\leq r\leq3\), is finite and volume
independent.  If the local mark derivative tensors are uniformly
bounded, then
\[
                        A_{r,*}^{\rm WP}:=\sup_jA_{r,j}^{\rm WP}
                        <\infty.                                \tag{33.10}
\]
\end{lemma}

\begin{proof}
A trunk with at most three prescribed marks has only finitely many
abstract shapes in the existing selected-tree code.  Each local
overlap is a finite-dictionary incidence.  Each separated edge is a
packet kernel with finite row (33.5).  Successively sum leaf edges
toward one fixed marked root.  At each step a finite weighted Schur
row is used, so the result is independent of volume.  Taking the
maximum over the finite trunk and type lists proves the claim.
\end{proof}

\begin{theorem}[One-, two-, and three-mark packet tree bounds]
\label{thm:e8v33-marked-resolvent}
Under (33.9), the absolutely summed connected selected-tree family
with \(r\in\{1,2,3\}\) prescribed marks obeys
\[
\boxed{\qquad
 {\cal R}_{r,j}^{\rm WP}
 \leq
 \frac{A_{r,j}^{\rm WP}}{(1-q_j^{\rm WP})^r}
 \leq
 \frac{A_{r,*}^{\rm WP}}{(1-q_*^{\rm WP})^r}.
\qquad}                                                         \tag{33.11}
\]
\end{theorem}

\begin{proof}
Remove the finite marked trunk.  The injective selected-parent code
from \Cref{rec:e8v1-v92-v94} roots every remaining branch at one of
the \(r\) marked sides.  A branch of length \(n\) is bounded by the
\(n\)-th power of the positive operator in (33.6).  In the
\(v\)-weighted type norm, its norm is at most
\((q_j^{\rm WP})^n\).  Summing the geometric series independently
on each marked side gives \((1-q_j^{\rm WP})^{-r}\).
The finite trunk costs \(A_{r,j}^{\rm WP}\).
\end{proof}

\begin{remark}[Why the exponent is conservative]
\label{rem:e8v33-mark-exponent}
Some marked trunks share branch families and admit a smaller power
of \((1-q)^{-1}\).  Bound (33.11) deliberately uses one resolvent
per mark.  It is uniform and sufficient for the derivative ledger.
\end{remark}

\subsection{Evaluation of the V32 conditional-expectation stock}

Separate two uses of the conditional map:
\[
\begin{array}{c|l}
r=1&\text{one differentiated local insertion in a fixed
conditional law},\\
r=2&\text{variation of that law or a connected covariance row}.
\end{array}                                                       \tag{33.12}
\]
Let \(J_1,J_2\) be the corresponding finite local mark-embedding
constants if they have not already been included in \(A_1,A_2\).
For definiteness absorb them into those trunk constants.

\begin{theorem}[Explicit packet conditional derivative constant]
\label{thm:e8v33-CE}
On the complete selected-tree grammar, the V32 stock may be chosen as
\[
\boxed{\qquad
 C_{{\rm CE},j}(\rho_{\rm st},r_{\rm f})
 =
 \max\left\{
 \frac{A_{1,j}^{\rm WP}}{1-q_j^{\rm WP}},
 \frac{A_{2,j}^{\rm WP}}{(1-q_j^{\rm WP})^2}
 \right\}.
\qquad}                                                         \tag{33.13}
\]
Moreover, if the good and saturated-bad types exhaust the exact
conditional interaction before the final circle product, then the
remainder in (32.23) is empty:
\[
                 {\cal E}_{{\rm cum},j}=0,\qquad
                 \varepsilon_{{\rm cum},j}=0.                   \tag{33.14}
\]
\end{theorem}

\begin{proof}
The first conditional differential is an expectation by
\eqref{eq:e8v5-DT}; its connected owner expansion has one marked
insertion.  Variation of the normalized conditional law is the
negative covariance \eqref{eq:e8v5-D2T}, hence has two marks.
Apply (33.11) with \(r=1,2\).  This proves (33.13).

The source-dependent selected-tree identity is an absolutely
convergent equality, not merely an estimate.  When its type grammar
contains every good vertex, saturated bad component, transition,
determinant response, and boundary owner, each connected conditional
term occurs in exactly one selected-tree family.  Thus no connected
cumulant term is left outside the activated conditional map.
The final circle product is formed only after this regrouping and is
the separate term \({\cal E}_{{\rm circ},j}\).  Hence (33.14).
\end{proof}

\begin{corollary}[V32 contraction gate with the cumulant row loaded]
\label{cor:e8v33-loaded-gate}
Under the hypotheses of
Theorem~\ref{thm:e8v33-CE}, put
\[
 \widehat C_{{\rm CE},j}:=
 \max\left\{
 \frac{A_{1,j}^{\rm WP}}{1-q_j^{\rm WP}},
 \frac{A_{2,j}^{\rm WP}}{(1-q_j^{\rm WP})^2}
 \right\}.                                                       \tag{33.15}
\]
Then (32.24) becomes
\[
\begin{split}
 Q_j\leq{}&
 \frac12\left[
 1+(\sqrt3+c_{\rm proj})\widehat\chi_S(R)
 \right]\\
 &+
 N_{\rm typ}C_{\rm cyc}\widehat C_{{\rm CE},j}
 \kappa_j^{\rm WP}(r_{\rm f})
 +2C_{\rm cyc}\rho_{\rm st}
 +\varepsilon_{{\rm ret}\to{\rm st},j}.                         \tag{33.16}
\end{split}
\]
Every factor in the conditional-cumulant coefficient on the second
line is now an explicit finite incidence, reserve, trunk, or branch
constant.
\end{corollary}

\begin{proof}
Insert (33.13)--(33.14) into
Definition~\ref{def:e8v32-activation} and (32.24).
\end{proof}

\subsection{Third derivatives and the generated-interaction stock}

The remaining symbol \(\zeta^{\rm WP}_{3,j}\) in (32.20) can be
related to the same marked tree.  Let a smooth parameter \(a\) enter
the fibre action \(S(a,\omega)\), and write
\[
 F(a):=-\log\int e^{-S(a,\omega)}\,d\kappa(\omega).               \tag{33.17}
\]
All derivatives below are Fr\'echet derivatives evaluated on
directions of norm at most one.

\begin{lemma}[Exact third conditional derivative]
\label{lem:e8v33-third-identity}
Whenever three differentiations under the integral are dominated,
\[
\boxed{
\begin{split}
 D^3F={}&
 \langle D^3S\rangle
 -3\,\operatorname{Sym}\operatorname{Cov}(D^2S,DS)\\
 &+\operatorname{Cum}_3(DS,DS,DS).
\end{split}}                                                     \tag{33.18}
\]
Here \(\operatorname{Sym}\) is the normalized sum over the three
placements of the singled-out first derivative, and
\(\operatorname{Cum}_3\) is the joint third cumulant.
\end{lemma}

\begin{proof}
For a parameter-dependent observable \(A\),
\[
 D\langle A\rangle
 =\langle DA\rangle-\operatorname{Cov}(A,DS).                    \tag{33.19}
\]
Starting from \(DF=\langle DS\rangle\), one differentiation gives
\[
 D^2F=\langle D^2S\rangle-\operatorname{Cov}(DS,DS).
\]
Differentiate once more.  Formula (33.19) applied to the first term
gives one covariance.  Differentiating the covariance gives the
other two symmetrized covariances and minus the third cumulant.
Subtracting that derivative yields (33.18).  Polarization gives the
full trilinear identity from the repeated-direction calculation.
\end{proof}

Let \(M_{\ell,j}^{\rm WP}\) be the complete local packet norm of
\(D^\ell S_j\), \(1\leq\ell\leq3\), on the declared field and stable
balls, after retained pullbacks have been separated.

\begin{theorem}[Three-mark recurrence for the generated stock]
\label{thm:e8v33-third-bound}
Under (33.9), the stable connected output satisfies
\[
\boxed{
\begin{split}
 \zeta^{\rm WP}_{3,j+1}
 \leq C_{\rm out}\bigg[
 &\frac{A_{1,j}^{\rm WP}}{1-q_j^{\rm WP}}M_{3,j}^{\rm WP}\\
 &+3\frac{A_{2,j}^{\rm WP}}{(1-q_j^{\rm WP})^2}
 M_{2,j}^{\rm WP}M_{1,j}^{\rm WP}\\
 &+\frac{A_{3,j}^{\rm WP}}{(1-q_j^{\rm WP})^3}
 (M_{1,j}^{\rm WP})^3
 \bigg].
\end{split}}                                                     \tag{33.20}
\]
The constant \(C_{\rm out}\) is the finite stable Taylor projection,
owner assignment, connected-hull, and reblocking norm.
\end{theorem}

\begin{proof}
Apply the one-mark, two-mark, and three-mark bounds (33.11) to the
three terms of (33.18), respectively.  Multilinearity gives the
products of the local derivative stocks.  The three covariance
placements give the factor three.  Finally apply the stable output
projection and the already proved hull/reblocking cycle.  Their
finite norm is \(C_{\rm out}\).  Absolute convergence justifies
termwise differentiation and regrouping.
\end{proof}

\begin{corollary}[Noncircular invariant-stock test]
\label{cor:e8v33-third-invariant}
Suppose the local third-derivative row splits as
\[
 M_{3,j}^{\rm WP}\leq m_{3,j}^{(0)}+
 a_{3,j}\zeta_{3,j}^{\rm WP},                                   \tag{33.21}
\]
while \(M_{1,j}^{\rm WP}\leq m_{1,*}\) and
\(M_{2,j}^{\rm WP}\leq m_{2,*}\).  Define
\[
 \lambda_{3,j}:=
 C_{\rm out}
 \frac{A_{1,j}^{\rm WP}}{1-q_j^{\rm WP}}a_{3,j}.                 \tag{33.22}
\]
If \(\sup_j\lambda_{3,j}=:\lambda_3<1\) and the source
\[
\begin{split}
 b_{3,j}:=C_{\rm out}\bigg[
 &\frac{A_{1,j}^{\rm WP}}{1-q_j^{\rm WP}}m_{3,j}^{(0)}
 +3\frac{A_{2,j}^{\rm WP}}{(1-q_j^{\rm WP})^2}m_{2,*}m_{1,*}\\
 &+\frac{A_{3,j}^{\rm WP}}{(1-q_j^{\rm WP})^3}m_{1,*}^3
 \bigg]                                                         \tag{33.23}
\end{split}
\]
is uniformly bounded by \(b_3\), then
\[
 \zeta_{3,j}^{\rm WP}\leq
 \max\left\{\zeta_{3,j_0}^{\rm WP},\frac{b_3}{1-\lambda_3}\right\}
                                                                    \tag{33.24}
\]
for all \(j\geq j_0\).
\end{corollary}

\begin{proof}
Equations (33.20)--(33.23) give the scalar recursion
\[
 \zeta_{3,j+1}^{\rm WP}\leq b_3+\lambda_3\zeta_{3,j}^{\rm WP}.
\]
Iterate the geometric series.
\end{proof}

\begin{firewall}[What remains in the third-derivative row]
\label{fire:e8v33-third-status}
Equation (33.20) is a proved marked-tree estimate, but
\(\lambda_3<1\) is not automatic.  It requires the actual canonical
rescaling and stable Taylor projection in \(a_{3,j}\).  Omitting that
factor and declaring the recurrence closed would be circular.
\end{firewall}

\subsection{Status after packet loading}

\begin{proposition}[Resolution of V32 item (i)]
\label{prop:e8v33-item-i}
Subject to the already declared complete activity and extensive
bad-component bounds, item (i) of (32.30) is resolved by
(33.13)--(33.14): the conditional derivative stock is finite and
volume independent, and there is no unassigned connected-cumulant
remainder.
\end{proposition}

\begin{proof}
This is Theorem~\ref{thm:e8v33-CE}, whose new analytic input is the
packet incidence lemma rather than a new cluster expansion.
\end{proof}

\begin{assessment}[Live gates after V33]
\label{ass:e8v33-status}
The stable small-field derivative gate is now (33.16).  The next
nonduplicative tasks are:
\begin{enumerate}
\item compute the canonical feedback coefficient \(a_{3,j}\) in
      (33.21) and test \(\lambda_3<1\);
\item evaluate the stable origin \(\beta_{{\rm st},j}\);
\item bound retained-to-stable feedback, including the running
      correlated Wilson coupling and stationary precision; and
\item only after those small-field rows close, join the good and
      saturated-bad sectors to the terminal OS/spectral criterion.
\end{enumerate}
The first task is upstream of the others because failure of
\(\lambda_3<1\) would force a different stable coordinate or
counterterm extraction.
\end{assessment}

\begin{firewall}[Claim boundary after V33]
\label{fire:e8v33-claim}
V33 does not prove the full Yang--Mills mass gap.  It proves the
packet lift of the existing marked conditional tree and thereby
removes the unspecified \(C_{{\rm CE},j}\) and
\(\varepsilon_{{\rm cum},j}\) rows from the V32 ledger.  The
three-derivative estimate is explicit but still awaits the canonical
feedback computation.
\end{firewall}

\subsection{Parent record}

This V33 note was formed by appending the preceding material to V32,
whose validated record was
\[
\begin{array}{c|c}
\text{lines}&14844\\
\text{bytes}&515312\\
\text{SHA-256}&
\texttt{4C55B7892145433C57B4A8907E6F7EDC9675BC6CA5205CCFC2A6DD931C7D0BF2}.
\end{array}
\]

% =============================================================================
% END APPENDED EIGHTH-SERIES V33 MATERIAL
% =============================================================================

% =============================================================================
% BEGIN APPENDED EIGHTH-SERIES V34 MATERIAL
% =============================================================================

\section{Eighth-series V34: strict cubic feedback from the empty-branch
split}

\subsection{The generic marked constant is too crude on the propagated row}

Bound (33.20) treats the one-mark family by a single constant
\(A_1^{\rm WP}(1-q^{\rm WP})^{-1}\).  That is sufficient for
finiteness but obscures the exact contraction on the part of
\(D^3S_j\) inherited from the previous scale.  The selected-tree
resolvent has the exact decomposition
\[
 (I-{\mathbb B}_j)^{-1}
 =
 I+{\mathbb B}_j(I-{\mathbb B}_j)^{-1},                          \tag{34.1}
\]
where \({\mathbb B}_j\) is the packet-refined unmarked branch
operator.  The identity is the empty-branch Gaussian principal map;
the second term contains at least one small activity.

\begin{lemma}[Nonempty-branch gain]
\label{lem:e8v34-nonempty}
In the \(v\)-weighted packet type norm of (33.7),
\[
 \left\|
 {\mathbb B}_j(I-{\mathbb B}_j)^{-1}
 \right\|
 \leq
 \frac{q_j^{\rm WP}}{1-q_j^{\rm WP}}.                            \tag{34.2}
\]
\end{lemma}

\begin{proof}
The left side is bounded by
\[
 \sum_{n\geq1}\|{\mathbb B}_j\|^n
 \leq\sum_{n\geq1}(q_j^{\rm WP})^n.
\]
Sum the geometric series.
\end{proof}

\subsection{A tunable packet margin on cubic chaos}

V27 chose one convenient slow radius \(R_*\), but its proof works for
every larger \(R\).  Write
\[
 \epsilon_R:=\varepsilon_{N,s}(R),\qquad
 q_{3,0}(R):=\frac14(1+\epsilon_R)^3.                            \tag{34.3}
\]

\begin{lemma}[Empty-branch cubic contraction]
\label{lem:e8v34-empty-cubic}
On the stable cubic matched-chaos complement, the empty-branch
principal map satisfies
\[
 \|{\cal L}_{3,{\rm G}}\|\leq q_{3,0}(R),\qquad
                  \lim_{R\to\infty}q_{3,0}(R)=\frac14.           \tag{34.4}
\]
\end{lemma}

\begin{proof}
The one-particle packet map costs at most \(1+\epsilon_R\) by
Theorem~\ref{thm:e8v27-one-packet}.  Its third symmetric tensor power
therefore costs at most \((1+\epsilon_R)^3\).  The canonical cubic
factor is \(2^{-2}=1/4\), by (18.16).  Restriction to the stable
complement and symmetric projection have norm at most one.
The limit follows from
Lemma~\ref{lem:e8v27-epsilon-limit}.
\end{proof}

\begin{definition}[Cubic-margin packet radius]
\label{def:e8v34-R3}
Fix any \(0<\delta_0<3/4\), and choose \(R_3\geq R_*\) so large that
\[
                     q_{3,0}(R_3)\leq1-\delta_0.                 \tag{34.5}
\]
For example \(\delta_0=5/8\) is available once
\[
 \epsilon_{R_3}\leq(3/2)^{1/3}-1,
\quad\text{because then}\quad q_{3,0}(R_3)\leq\frac38.           \tag{34.6}
\]
\end{definition}

\begin{remark}[The radius tradeoff]
\label{rem:e8v34-radius-tradeoff}
Increasing \(R\) improves (34.4) but enlarges the finite inverse-square
owner stock \({\cal P}_{R,s}\).  This does not destroy any theorem:
weak-coupling activities still beat every fixed finite owner
constant.  It does mean that \(R_3\) is chosen once before taking the
cofinal weak-coupling range.
\end{remark}

\subsection{Parallel typed charts remove a fictitious cubic Wick debit}

Let \({\cal W}_j\) be the stationary Wick rechart from scale \(j\) to
\(j+1\).  For purposes of comparing stable fibres, transport the
marginal projector by
\[
 P_{{\rm marg},j+1}^{\parallel}
 :=
 {\cal W}_jP_{{\rm marg},j}{\cal W}_j^{-1}.                       \tag{34.7}
\]
The difference between this parallel projector and the physical
target Taylor projector is a retained-chart feedback row and belongs
to \(\varepsilon_{{\rm ret}\to{\rm st},j}\).

\begin{lemma}[No stationary stable cubic lowering]
\label{lem:e8v34-no-cubic-Wick}
Between the parallel stable fibres, the stationary Wick rechart adds
no stable \(C_3^\circ\to C_3^\circ\) debit.  Its only nonidentity
cubic row is
\[
                         C_3\longrightarrow C_1,                 \tag{34.8}
\]
which is retained.
\end{lemma}

\begin{proof}
The exact Wick change lowers homogeneous degree by an even integer.
For degree three its nonidentity term has degree one, as listed in
(32.10).  Equation (34.7) makes the stable-complement transport
intertwine exactly with \({\cal W}_j\).  Hence the degree-three
stable diagonal is the identity transport already included in
\({\cal L}_{3,{\rm G}}\), while (34.8) is removed by
\(P_{\rm st}\).
\end{proof}

\begin{firewall}[Parallel comparison does not erase physical
feedback]
\label{fire:e8v34-parallel}
The physical target projector is not asserted to equal (34.7).
Their difference is merely moved to the correctly typed retained
feedback ledger.  It must still be bounded before the full
small-field map closes.
\end{firewall}

\subsection{The sharpened cubic recurrence}

Let \(D_3^{\rm br}<\infty\) be the maximum norm of:
\begin{enumerate}
\item the finite one-mark root embedding for a propagated cubic
      stable tensor;
\item one nonempty packet branch attachment;
\item the stable output projection, owner assignment, connected
      hull, and reserve-restoring reblock.
\end{enumerate}
Its finiteness follows from
\Cref{lem:e8v33-packet-incidence,lem:e8v33-marked-trunks} and
Theorem~\ref{thm:e8v31-reserve-restoration}.

Define the source \(b_{3,j}^{\rm src}\) to contain every term of
(33.18) which is not linear in the propagated
\(\zeta_{3,j}^{\rm WP}\).  In particular it contains the microscopic
Wilson/Haar third derivative, the three two-mark
\((D^2S,DS)\) terms, and the three-mark \((DS)^3\) term.

\begin{theorem}[Strict cubic feedback coefficient]
\label{thm:e8v34-cubic-feedback}
At the radius \(R_3\), the generated third-derivative stock obeys
\[
\boxed{\qquad
 \zeta_{3,j+1}^{\rm WP}
 \leq
 \lambda_{3,j}^{\sharp}\zeta_{3,j}^{\rm WP}
 +b_{3,j}^{\rm src},
\qquad}                                                         \tag{34.9}
\]
where
\[
\boxed{\qquad
 \lambda_{3,j}^{\sharp}
 :=
 q_{3,0}(R_3)
 +D_3^{\rm br}\frac{q_j^{\rm WP}}{1-q_j^{\rm WP}}.
\qquad}                                                         \tag{34.10}
\]
Consequently there is a cofinal weak-coupling range on which
\[
                         \lambda_{3,j}^{\sharp}\leq
                         1-\frac{\delta_0}{2}<1.                 \tag{34.11}
\]
\end{theorem}

\begin{proof}
Split the one-mark selected-tree family by (34.1).  The empty-branch
piece is the exact cubic Gaussian principal map and costs at most
\(q_{3,0}(R_3)\) by
Lemma~\ref{lem:e8v34-empty-cubic}.  The stable stationary Wick chart
adds no cubic debit by
Lemma~\ref{lem:e8v34-no-cubic-Wick}.  The family with at least one
branch costs at most
\[
 D_3^{\rm br}\frac{q_j^{\rm WP}}{1-q_j^{\rm WP}}
\]
by Lemma~\ref{lem:e8v34-nonempty}.  All terms not linear in the
propagated cubic stock are, by definition, in
\(b_{3,j}^{\rm src}\).  This proves (34.9)--(34.10).

Theorem~\ref{thm:e8v33-qwp} gives \(q_j^{\rm WP}\to0\) cofinally.
Thus the second term of (34.10) is eventually at most
\(\delta_0/2\).  Combine with (34.5).
\end{proof}

\begin{proposition}[Explicit source majorant]
\label{prop:e8v34-source}
Let \(m_{\ell,j}^{(0)}\) denote the local derivative stock after the
propagated cubic stable part has been removed.  Then one may take
\[
\begin{split}
 b_{3,j}^{\rm src}\leq C_{\rm out}\bigg[
 &\frac{A_{1,j}^{\rm WP}}{1-q_j^{\rm WP}}m_{3,j}^{(0)}\\
 &+3\frac{A_{2,j}^{\rm WP}}{(1-q_j^{\rm WP})^2}
 M_{2,j}^{\rm WP}M_{1,j}^{\rm WP}\\
 &+\frac{A_{3,j}^{\rm WP}}{(1-q_j^{\rm WP})^3}
 (M_{1,j}^{\rm WP})^3
 \bigg].                                                        \tag{34.12}
\end{split}
\]
\end{proposition}

\begin{proof}
This is (33.20) applied only to the source part of (33.18).  The
propagated linear term has already been estimated sharply in
(34.10), so it is not included again.
\end{proof}

\begin{corollary}[Uniform and vanishing generated stock]
\label{cor:e8v34-generated}
On a cofinal range satisfying (34.11), if
\[
                         \sup_jb_{3,j}^{\rm src}\leq b_3<\infty,
                                                                    \tag{34.13}
\]
then
\[
 \sup_{j\geq j_0}\zeta_{3,j}^{\rm WP}
 \leq
 \max\left\{
 \zeta_{3,j_0}^{\rm WP},\frac{2b_3}{\delta_0}
 \right\}.                                                       \tag{34.14}
\]
If in addition \(b_{3,j}^{\rm src}\to0\), then
\[
                         \zeta_{3,j}^{\rm WP}\longrightarrow0.   \tag{34.15}
\]
\end{corollary}

\begin{proof}
Boundedness follows by iterating (34.9) with contraction ratio
\(1-\delta_0/2\).  For convergence, fix \(\varepsilon>0\).  Choose
\(J\) so that \(b_{3,j}^{\rm src}\leq\varepsilon\delta_0/4\) for
\(j\geq J\).  Iteration from \(J\) gives
\[
 \zeta_{3,j}^{\rm WP}
 \leq
 (1-\delta_0/2)^{j-J}\zeta_{3,J}^{\rm WP}
 \frac{\varepsilon}{2}.
\]
The first term is eventually at most \(\varepsilon/2\).
\end{proof}

\subsection{When the source vanishes}

The bound (34.12) becomes useful only after its local factors are
typed.  On a centred small-field chart:
\begin{enumerate}
\item the Wilson cubic third derivative is \(O(g_j)\);
\item the quartic and analytic-tail contribution to
      \(m_{3,j}^{(0)}\) is \(O(g_j^2r_{\rm f})\);
\item the first derivative after source and marginal extraction is
      \(O(r_{\rm f})\); and
\item the second derivative is uniformly bounded by the normalized
      Gaussian floor plus \(O(g_jr_{\rm f})\).
\end{enumerate}
Thus, for fixed packet and tree constants, the schematic consequence
of (34.12) is
\[
 b_{3,j}^{\rm src}
 \leq
 C_3\left(g_j+r_{\rm f}+g_j^2r_{\rm f}+r_{\rm f}^3\right),        \tag{34.16}
\]
provided the centred first-derivative and complete local-tail
estimates hold in the packet norm.

\begin{firewall}[The small-field radius cannot be fixed blindly]
\label{fire:e8v34-radius-source}
Equation (34.16) tends to zero only along a schedule on which both
\(g_j\to0\) and \(r_{\rm f}\to0\), or after the \(O(r_{\rm f})\)
term is removed by a sharper centred symmetry identity.  Earlier
large-field decompositions used a logarithmically growing
unrescaled good radius.  The two statements are compatible only
after converting to normalized field coordinates.  V34 therefore
does not identify the two radii.
\end{firewall}

\begin{proposition}[Resolution of the feedback part of V32 item (ii)]
\label{prop:e8v34-item-ii}
The previously unspecified coefficient multiplying
\(\zeta_{3,j}^{\rm WP}\) is strictly below one cofinally and is
given by (34.10).  Hence item (ii) of (32.30) is reduced to the
source estimate (34.12), with no circular dependence on the stock
being bounded.
\end{proposition}

\begin{proof}
Use
\Cref{thm:e8v34-cubic-feedback,prop:e8v34-source}.
\end{proof}

\begin{assessment}[Next upstream question]
\label{ass:e8v34-status}
The apparent cubic self-dependence is no longer the bottleneck.
The live upstream issue is the normalized-radius dictionary in
Firewall~\ref{fire:e8v34-radius-source}.  It controls both the source
\(b_{3,j}^{\rm src}\) and the stable origin
\(\beta_{{\rm st},j}\).  V35 must audit the companion SM and NS
notes, then construct that dictionary before attempting the retained
feedback row.
\end{assessment}

\begin{firewall}[Claim boundary after V34]
\label{fire:e8v34-claim}
V34 proves a cofinally strict cubic feedback coefficient by exposing
the empty branch.  It does not yet prove that the source in (34.12)
vanishes in the same norm, close the invariant stable ball, control
the retained trajectory, construct the continuum measure, or prove
the Yang--Mills mass gap.
\end{firewall}

\subsection{Parent record}

This V34 note was formed by appending the preceding material to V33,
whose validated record was
\[
\begin{array}{c|c}
\text{lines}&15353\\
\text{bytes}&534275\\
\text{SHA-256}&
\texttt{B48904CBAC2CE3E7A2E6D5F259827B88E4B1AED010B3B1DDC5A24BECDFDCB035}.
\end{array}
\]

% =============================================================================
% END APPENDED EIGHTH-SERIES V34 MATERIAL
% =============================================================================

% =============================================================================
% BEGIN APPENDED EIGHTH-SERIES V35 MATERIAL
% =============================================================================

\section{Eighth-series V35: companion audit and the normalized-radius
dictionary}

\subsection{Compulsory V35 companion audit}

The active companion notes inspected before choosing the V35
direction were
\[
\begin{array}{c|r|r|l}
\text{file}&\text{lines}&\text{bytes}&\text{SHA--256}\\ \hline
\texttt{Renewal\_SM\_Einstein\_Continuum\_Research\_Notes\_V18.tex}
&7354&279400&
\texttt{D8EF29CB7CC080046559CDE2705AD0ADB7CFCE7A968A7ED6F948F9AA228ED0BA}
\\
\texttt{Renewal\_NS\_Stability\_Research\_Notes\_V26.tex}
&5320&278283&
\texttt{82C887F8C2C69E4F28FAD7C3DC8DE5B9099CB5A4BF431EBA1FFF3E91935978AD}.
\end{array}                                                       \tag{35.1}
\]

The new SM material V12--V18 proves a local
symmetric-divergence repair, a retained affine conservation ledger,
an infinite-horizon relevant/irrelevant shooting theorem, a
conditional marginal-decay lemma, and an analytic-radius reserve
cycle for an infinite graded operator tower.  The NS file is
unchanged from the V30 audit; its latest result computes exact
rank-one norms for first and second Oseen-kernel derivatives.

\begin{assessment}[V35 audit transfer]
\label{ass:e8v35-audit}
Three audit conclusions govern the next step.
\begin{enumerate}
\item The SM analytic coefficient radius is not a field-amplitude
      cutoff.  Its reserve-cycle proof supports using a separate
      graded coefficient radius for the YM analytic tail, but it
      cannot be used to force the normalized Gaussian field radius
      to shrink.
\item The SM shooting theorem is applicable only after the YM stable
      ball and retained cross-derivatives close.  Its marginal lemma
      is relevant to \(u_j=g_j^2\), but the positive Yang--Mills beta
      coefficient and its matched-scheme remainder must be proved in
      this programme, not imported.
\item The NS rank-one calculation shows that a restricted physical
      tangent cone may sharpen a tensor norm only after exact
      maximization.  No Oseen-kernel constant or Navier--Stokes
      estimate transfers to the gauge model.
\end{enumerate}
\end{assessment}

\begin{decision}[Direction after the audit]
\label{dec:e8v35-audit}
The V34 radius firewall is resolved by distinguishing:
\[
\begin{array}{c|l}
R_{\rm f}&\text{radius of the normalized Gaussian fluctuation},\\
g_jR_{\rm f}&\text{dimensionless compact-group chart amplitude},\\
R_{\rm an}&\text{analytic coefficient radius of the local
operator tower}.
\end{array}                                                       \tag{35.2}
\]
V35 estimates the first two quantities.  The third is transported by
an analytic reserve cycle and is not identified with either of them.
\end{decision}

\subsection{The normalized weak-field chart}

Write a fine link in a chosen good chart as
\[
                         U_\ell=\exp(g_j a_\ell),                 \tag{35.3}
\]
where \(a\) has an \(O(1)\) Gaussian covariance after the quadratic
Wilson form is normalized.  Choose
\[
 R_{{\rm f},j}:=
 c_{\rm f}\bigl(1+\log g_j^{-1}\bigr)^{\nu},
\qquad c_{\rm f},\nu>0.                                         \tag{35.4}
\]
The good-field event is \(\|a\|_\infty\leq R_{{\rm f},j}\).

\begin{lemma}[Logarithmic radii are beaten by every coupling power]
\label{lem:e8v35-log-beaten}
For all \(\alpha>0\) and \(p\geq0\),
\[
                g^\alpha\bigl(1+\log g^{-1}\bigr)^p
                \longrightarrow0
                \quad\text{as }g\downarrow0.                    \tag{35.5}
\]
In particular \(g_jR_{{\rm f},j}\to0\).
\end{lemma}

\begin{proof}
Put \(t=\log g^{-1}\).  The expression is
\(e^{-\alpha t}(1+t)^p\), which tends to zero because exponential
decay dominates every fixed power.
\end{proof}

\begin{firewall}[A growing normalized ball is a shrinking group chart]
\label{fire:e8v35-two-radii}
The facts \(R_{{\rm f},j}\to\infty\) and
\(g_jR_{{\rm f},j}\to0\) are compatible.  The Wilson Taylor series is
controlled by the latter, whereas Gaussian large-field probability
is controlled by the former.  Requiring \(R_{{\rm f},j}\to0\), as
the crude reading after (34.16) suggested, would reverse the purpose
of the good/bad decomposition.
\end{firewall}

\subsection{Derivative bounds after the Gaussian reference is removed}

Let the complete local good-chart interaction after its quadratic
Gaussian reference and retained Taylor jet have been extracted be
\[
 V_j(a)=\sum_{n\geq3}g_j^{\,n-2}V_{n,j}(a)+V_{{\rm Haar},j}(a)
 +V_{{\rm det},j}(a)+V_{{\rm tr},j}(a).                          \tag{35.6}
\]
The last three terms are the nonconstant Haar, determinant-response,
and chart-transition owners already present in the V8 type grammar.

\begin{definition}[Uniform analytic local card]
\label{def:e8v35-local-card}
The local chart passes UALC if there are \(K>0\) and
\(\rho_{\rm an}<1\), independent of \(j\) and volume, such that for
\(\ell=0,1,2,3\),
\[
 \|D^\ell V_{n,j}(a)\|
 \leq K n^3 R_{{\rm f},j}^{\,n-\ell}
 \quad
 (n\geq\max\{3,\ell\}),                                          \tag{35.7}
\]
on the good ball, and
\[
                         g_jR_{{\rm f},j}\leq\rho_{\rm an}.       \tag{35.8}
\]
The three auxiliary owner types obey the corresponding V8
source-marked polynomial bounds with a positive power of \(g_j\).
\end{definition}

\begin{lemma}[Interaction derivative hierarchy]
\label{lem:e8v35-derivative-hierarchy}
Under UALC, for \(\ell=1,2,3\),
\[
\boxed{\qquad
 \|D^\ell V_j(a)\|
 \leq
 C_\ell g_jR_{{\rm f},j}^{\,3-\ell}
 +h_{\ell,j},
\qquad}                                                         \tag{35.9}
\]
where each \(h_{\ell,j}\) is a finite sum of positive powers of
\(g_j\) times fixed powers of \(R_{{\rm f},j}\).  Hence every term on
the right tends to zero.
\end{lemma}

\begin{proof}
For the Taylor series in (35.6), (35.7) gives
\[
\begin{split}
 \sum_{n\geq3}g_j^{\,n-2}\|D^\ell V_{n,j}(a)\|
 &\leq
 Kg_jR_{{\rm f},j}^{\,3-\ell}
 \sum_{n\geq3}n^3
 (g_jR_{{\rm f},j})^{n-3}\\
 &\leq C_\ell g_jR_{{\rm f},j}^{\,3-\ell},
\end{split}
\]
because the last series is uniformly finite under (35.8).
The auxiliary owners give \(h_{\ell,j}\) by their declared activity
card.  Lemma~\ref{lem:e8v35-log-beaten} proves convergence to zero.
\end{proof}

\begin{remark}[Where the \(O(R_{\rm f})\) term went]
\label{rem:e8v35-Gaussian-split}
The \(O(R_{\rm f})\) first derivative of the quadratic Gaussian
action is not an interaction mark.  Its empty-branch conditional
transport is precisely the Gaussian principal map separated in
V34.  Only derivatives of \(V_j\), bounded by (35.9), belong to the
nonlinear source.  This is the missing distinction behind
Firewall~\ref{fire:e8v34-radius-source}.
\end{remark}

\subsection{Cofinal decay of the cubic source}

Let
\[
 {\mathfrak z}_j:=
 \sum_{\gamma\in{\cal I}}
 z_{\gamma,j}^{\rm WP}                                           \tag{35.10}
\]
denote the complete good, transition, determinant, and saturated-bad
activity stock in the packet type norm.  The V8 activity ledger and
Lemma~\ref{lem:e8v35-log-beaten} give
\[
             {\mathfrak z}_j\longrightarrow0,\qquad
             q_j^{\rm WP}\leq C_{\rm br}{\mathfrak z}_j.         \tag{35.11}
\]

\begin{theorem}[Normalized-radius cubic source bound]
\label{thm:e8v35-cubic-source}
Under UALC and the complete packet MGKL hypotheses of V33, there is a
fixed polynomial \(P_3\) with nonnegative coefficients such that
\[
\boxed{\qquad
 b_{3,j}^{\rm src}
 \leq
 C_3\left[
 g_jP_3(R_{{\rm f},j})
 +\frac{{\mathfrak z}_j}{1-C_{\rm br}{\mathfrak z}_j}
 \right].
\qquad}                                                         \tag{35.12}
\]
Consequently \(b_{3,j}^{\rm src}\to0\).
\end{theorem}

\begin{proof}
Apply Proposition~\ref{prop:e8v34-source} after the Gaussian
empty-branch contribution is removed.  The direct third derivative
costs \(O(g_j)+h_{3,j}\).  The two-mark term is bounded by
\[
 O(g_jR_{{\rm f},j})\,
 O(g_jR_{{\rm f},j}^2)
 \]
plus auxiliary-owner polynomials, and the three-mark term by
\[
 O(g_j^3R_{{\rm f},j}^6)
 \]
plus such polynomials.  All marked trunk constants are fixed, while
their branch denominators are bounded once
\(C_{\rm br}{\mathfrak z}_j<1\).

Terms in which a Gaussian source mark is connected to the output by
at least one interaction branch contain the nonempty-resolvent
factor
\[
 \frac{q_j^{\rm WP}}{1-q_j^{\rm WP}}
 \leq
 \frac{C_{\rm br}{\mathfrak z}_j}
 {1-C_{\rm br}{\mathfrak z}_j}.
\]
Absorb the finitely many powers into \(P_3\) and \(C_3\).  The first
term tends to zero by
Lemma~\ref{lem:e8v35-log-beaten}, and the second by (35.11).
\end{proof}

\begin{lemma}[Slow-source propagation]
\label{lem:e8v35-slow-source}
Let \(0<\lambda<1\), and suppose
\[
 z_{j+1}\leq\lambda z_j+b_j,\qquad b_j>0,                        \tag{35.13}
\]
where \(b_j\to0\) and, for every \(\gamma>1\), eventually
\[
                             b_{j-1}\leq\gamma b_j.               \tag{35.14}
\]
Then \(z_j=O(b_j)\).
\end{lemma}

\begin{proof}
Choose \(1<\gamma<\lambda^{-1}\) and \(J\) so that (35.14) holds for
\(j\geq J\).  Iteration gives
\[
 z_j\leq\lambda^{j-J}z_J+
 \sum_{k=0}^{j-J-1}\lambda^k b_{j-1-k}
 \leq
 \lambda^{j-J}z_J+
 \gamma b_j\sum_{k\geq0}(\lambda\gamma)^k.                       \tag{35.15}
\]
Repeated use of (35.14) also gives
\(b_j\geq\gamma^{-(j-J)}b_J\).  Therefore
\[
 \lambda^{j-J}z_J
 \leq
 \frac{z_J}{b_J}(\lambda\gamma)^{j-J}b_j
 \leq\frac{z_J}{b_J}b_j.
\]
This proves the asserted \(O(b_j)\) bound without an additional
regular-variation assumption.
\end{proof}

\begin{corollary}[Cubic stock with the needed rate]
\label{cor:e8v35-cubic-rate}
Assume the weak-coupling scale sequence is eventually slow in the
sense of (35.14) for the right side of (35.12).  Then
\[
 \zeta_{3,j}^{\rm WP}=O\left(
 g_jP_3(R_{{\rm f},j})+{\mathfrak z}_j
 \right),                                                       \tag{35.16}
\]
and in particular
\[
 \zeta_{3,j}^{\rm WP}R_{{\rm f},j}\longrightarrow0.             \tag{35.17}
\]
\end{corollary}

\begin{proof}
Use the cofinal contraction (34.11),
Theorem~\ref{thm:e8v35-cubic-source}, and
Lemma~\ref{lem:e8v35-slow-source}.  Every term on the right of
(35.16), after one additional factor \(R_{{\rm f},j}\), is a positive
power of \(g_j\) times a fixed logarithmic power, or an extensive bad
activity with the same property.  Apply
Lemma~\ref{lem:e8v35-log-beaten}.
\end{proof}

\begin{corollary}[Vanishing local nonlinear derivative debit]
\label{cor:e8v35-kappa}
With \(r_{\rm f}=R_{{\rm f},j}\), the V32 quantity satisfies
\[
 \kappa_j^{\rm WP}
 =
 \bigl(64\omega_2g_j+\zeta_{3,j}^{\rm WP}\bigr)
 R_{{\rm f},j}
 \longrightarrow0.                                              \tag{35.18}
\]
\end{corollary}

\begin{proof}
The first term tends to zero by
Lemma~\ref{lem:e8v35-log-beaten}; the second is (35.17).
\end{proof}

\subsection{The stable origin is another rooted activity sum}

At \(g=0\), after vacuum, source, marginal, harmonic, and zero-shell
coordinates are projected to the retained sector, the stable origin
vanishes.  At \(g>0\), every stable term in
\(P_{\rm st}{\cal T}_j(0)\) therefore contains at least one
non-Gaussian good vertex or one saturated-bad component.

\begin{theorem}[Cofinal stable-origin estimate]
\label{thm:e8v35-origin}
There is a finite packet root constant \(B_{\rm root}\) such that
\[
\boxed{\qquad
 \beta_{{\rm st},j}
 \leq
 B_{\rm root}
 \frac{{\mathfrak z}_j}{1-q_j^{\rm WP}}.
\qquad}                                                         \tag{35.19}
\]
Consequently \(\beta_{{\rm st},j}\to0\) cofinally.
\end{theorem}

\begin{proof}
Choose the first non-Gaussian activity in the existing selected-tree
code as the root.  Its absolute packet sum is
\({\mathfrak z}_j\).  Every subsequent branch is summed by the
one-root resolvent, costing at most
\((1-q_j^{\rm WP})^{-1}\).  The stable Taylor projection, packet
owner assignment, connected hull, and reblocking have one fixed
finite norm \(B_{\rm root}\), by V29--V31.  Retained zero-activity
Gaussian outputs are killed by \(P_{\rm st}\).  This proves (35.19).
\end{proof}

\begin{corollary}[Cofinal invariant-ball choice without retained
feedback]
\label{cor:e8v35-origin-ball}
Suppose for the moment that
\(\varepsilon_{{\rm ret}\to{\rm st},j}=0\), and that the stationary
linear margin in (32.31) is uniformly positive.  Then there are a
fixed \(\rho_{\rm st}>0\) and a cofinal \(j_0\) such that (32.27)
holds for all \(j\geq j_0\).
\end{corollary}

\begin{proof}
First choose \(\rho_{\rm st}\) so small that
\(2C_{\rm cyc}\rho_{\rm st}\) consumes less than one quarter of the
stationary linear margin.  By
\Cref{cor:e8v35-kappa,thm:e8v35-origin}, advance \(j_0\) until the
activated local derivative consumes less than a second quarter and
\(\beta_{{\rm st},j}\) is at most the remaining contraction margin
times \(\rho_{\rm st}\).  Then (33.16) and (32.27) hold.
\end{proof}

\begin{firewall}[The retained row is now isolated, not assumed away]
\label{fire:e8v35-retained}
Corollary~\ref{cor:e8v35-origin-ball} is not the full small-field
closure because it temporarily sets the retained-to-stable feedback
to zero.  Its purpose is diagnostic: the conditional-cumulant,
generated-third-derivative, local-Hessian, nonlinear-circle, and
stable-origin rows can all fit inside the contraction margin on the
same logarithmically growing good ball.  The remaining small-field
gate is the physical retained trajectory and its feedback.
\end{firewall}

\subsection{Audit-adjusted status}

\begin{proposition}[Resolution of V32 items (ii) and (iii), modulo
the declared local card]
\label{prop:e8v35-items-ii-iii}
Under UALC and the complete activity ledger:
\begin{enumerate}
\item the generated third-derivative stock has the rate (35.16) and
      the nonlinear Hessian debit tends to zero;
\item the stable origin obeys (35.19) and tends to zero.
\end{enumerate}
Thus items (ii) and (iii) of (32.30) are no longer independent
bottlenecks.
\end{proposition}

\begin{proof}
Use
\Cref{cor:e8v35-cubic-rate,cor:e8v35-kappa,
thm:e8v35-origin}.
\end{proof}

\begin{assessment}[Next acquisition after V35]
\label{ass:e8v35-status}
The sole remaining small-field row in (32.30) is item (iv):
retained-to-stable feedback together with the retained trajectory.
The SM audit suggests the correct architecture---stable directions
forward, relevant directions by backward shooting, conserved
coordinates by exact affine transport, and \(u_j=g_j^2\) by its own
marginal normal form.  V36 should construct the Yang--Mills version
of that block decomposition and identify exactly which beta-function
and Ward estimates are still missing.
\end{assessment}

\begin{firewall}[Claim boundary after V35]
\label{fire:e8v35-claim}
V35 repairs the normalized-radius dictionary and proves cofinal
decay of the cubic source, nonlinear local debit, and stable origin
under the already declared analytic/activity cards.  It does not
calculate the Yang--Mills beta function in the matched regulator,
close retained feedback, construct the infinite trajectory or
continuum measure, establish OS reconstruction, or prove the mass
gap.
\end{firewall}

\subsection{Parent record}

This V35 note was formed by appending the preceding material to V34,
whose validated record was
\[
\begin{array}{c|c}
\text{lines}&15702\\
\text{bytes}&546305\\
\text{SHA-256}&
\texttt{CEBA6D408756F3AF3DD3F65A0781DBAAC0D58268E8E50163ACE4FFE9B5BAB372}.
\end{array}
\]

% =============================================================================
% END APPENDED EIGHTH-SERIES V35 MATERIAL
% =============================================================================

% =============================================================================
% BEGIN APPENDED EIGHTH-SERIES V36 MATERIAL
% =============================================================================

\section{Eighth-series V36: exact removal of retained feedback and
cofinal small-field closure}

\subsection{The retained row was already algebraically triangularized}

The quantity
\(\varepsilon_{{\rm ret}\to{\rm st},j}\) in (32.24) was included as
a safety ledger entry.  In an arbitrary physical-coordinate chart it
is necessary.  In the exact pullback chart of V9 it vanishes from the
aligned stable derivative for an algebraic reason, not a smallness
reason.

\begin{theorem}[Exact retained-to-stable decoupling]
\label{thm:e8v36-exact-decoupling}
Let \({\mathscr J}_j=B_j^*\iota_{j+1}\) be the exact pullback
injection for the complete finite-cutoff retained commutant, and let
\[
 S={\mathscr J}_jm+q,\qquad \pi_{{\rm pb},j}q=0.                  \tag{36.1}
\]
For a prescribed output retained coordinate \(m_+\), calibrate
\[
                         m=m_+-\beta_j(q),\qquad
 \beta_j(q):=\pi_{j+1}{\cal T}_jq.                              \tag{36.2}
\]
Then the aligned stable map is
\[
 \Phi_{j,m_+}(q)=Q_{j+1}{\cal T}_jq,                             \tag{36.3}
\]
is independent of \(m_+\), and
\[
 D_q\Phi_{j,m_+}
 =
 Q_{j+1}D{\cal T}_{j,q}Q_{{\rm pb},j}.                           \tag{36.4}
\]
In particular the V32 ledger may be specialized to
\[
\boxed{\qquad
 \varepsilon_{{\rm ret}\to{\rm st},j}=0
 \quad\text{in the aligned pullback derivative.}
\qquad}                                                         \tag{36.5}
\]
\end{theorem}

\begin{proof}
The affine factorization
\[
 {\cal T}_j({\mathscr J}_jm+q)
 =\iota_{j+1}m+{\cal T}_jq
\]
is Theorem~\ref{thm:e8v9-affine}.  Applying the retained projection
gives (36.2); applying the stable projection gives (36.3).
Differentiation gives (36.4).  Equivalently,
\(A=I\) and \(D_{\rm ms}=0\), so the Schur correction vanishes by
Corollary~\ref{cor:e8v9-Schur-zero}.  This is exactly (36.5).
\end{proof}

\begin{firewall}[Response is retained, not erased]
\label{fire:e8v36-response}
Stable-to-retained response remains present in (36.2):
\[
                              D_qm=-D_q\beta_j(q).                \tag{36.6}
\]
Theorem~\ref{thm:e8v36-exact-decoupling} removes its return into the
aligned stable derivative.  It does not set the response or the
physical beta function to zero.
\end{firewall}

\subsection{Covariance-matched fibres and parallel identification}

Let \(S_j(m)\) be the stationary positive precision used to define
the matched Gaussian chaos at retained coordinate \(m\).  Its stable
packet space is denoted \(X_{{\rm st},j}(m)\).

\begin{definition}[Uniform precision-bundle card]
\label{def:e8v36-precision-card}
UPBC consists of:
\begin{enumerate}
\item a uniform positive physical-detail floor for \(S_j(m)\);
\item uniform order-zero Mikhlin bounds for its normalized block
      symbols;
\item a parallel isometric identification of the matched
      one-particle spaces, extended by second quantization to the
      chaos spaces; and
\item a Lipschitz precision estimate
\[
 \left\|
 S_j(m)^{-1/2}\bigl(S_j(m')-S_j(m)\bigr)S_j(m)^{-1/2}
 \right\|_{\rm Mih}
 \leq L_S|m'-m|.                                    \tag{36.7}
\]
\end{enumerate}
\end{definition}

\begin{lemma}[Parallelization preserves exact decoupling]
\label{lem:e8v36-parallel-decoupling}
Conjugating (36.3) by the UPBC parallel identifications preserves
(36.4)--(36.5).  Changing from the parallel target projector to a
physical local Taylor projector produces a coordinate connection
depending on \(m_{j+1}-m_j\); it is an external moving-reference
forcing, not a derivative with respect to \(q\) at fixed prescribed
\(m_{j+1}\).
\end{lemma}

\begin{proof}
Conjugation by fixed linear isometries at the source and target
does not change an operator norm or create an off-diagonal derivative
block.  For fixed \(m_{j+1}\), the target identification is constant
while \(q\) varies, so differentiation gives the conjugate of
(36.4).  Comparing two consecutive physical target projectors
instead differentiates the chart in its retained base coordinate;
the mean-value theorem and (36.7) make that difference proportional
to \(|m_{j+1}-m_j|\).  Thus it belongs to the base-path connection.
\end{proof}

\subsection{Retained increments control the stationary Wick change}

Let \({\mathfrak z}_j\) be (35.10), and let \(h_j(m)\) be the moving
stationary reference constructed below.  Put
\[
 e:=q-h_j(m),\qquad
 b_{{\rm ref},j}(m):=\beta_j(h_j(m)),                           \tag{36.7a}
\]
after removing the explicitly prescribed free normalization.

\begin{lemma}[Retained response about the moving reference]
\label{lem:e8v36-retained-response}
There are finite constants \(B_{\rm G},B_{\rm ret}\) such that, on
the cofinal packet-MGKL range,
\[
 \|\beta_j(q)-b_{{\rm ref},j}(m)\|
 \leq
 B_{\rm G}\|e\|
 +B_{\rm ret}
 \frac{{\mathfrak z}_j}{1-q_j^{\rm WP}}
 \bigl(1+\|e\|\bigr).                                             \tag{36.8}
\]
\end{lemma}

\begin{proof}
Differentiate the retained projection of the conditional action
along \(h_j(m)+te\), \(0\leq t\leq1\).  The derivative is a marked
conditional expectation by
Theorem~\ref{thm:e8v5-cumulant}.  Its zero-activity part is the
Gaussian Wick-lowering response; the finite factorial-chaos bounds
in (32.10), followed by the retained projection, give
\(B_{\rm G}\|e\|\).  In every other term select the first
non-Gaussian activity as root.  Its sum is at most
\({\mathfrak z}_j\), while the remaining one-mark branches cost
\((1-q_j^{\rm WP})^{-1}\) by
Theorem~\ref{thm:e8v33-marked-resolvent}.  The retained Taylor
projection and finite mark embedding give \(B_{\rm ret}\).
Integrate in \(t\).
\end{proof}

The exact retained recursion consequently has the estimate
\[
 |m_{j+1}-m_j|
 \leq
 |b_{{\rm ref},j}(m_j)|
 +B_{\rm G}\|e_j\|
 +B_{\rm ret}
 \frac{{\mathfrak z}_j}{1-q_j^{\rm WP}}
 (1+\|e_j\|)+d_{{\rm ret},j},                                   \tag{36.9}
\]
where \(d_{{\rm ret},j}\) is the declared regulator, chart-boundary,
or normalization defect.

\begin{corollary}[Cofinal covariance-change smallness]
\label{cor:e8v36-covariance-change}
Under UPBC, if
\[
 b_{{\rm ref},j}(m_j)\to0,\qquad
 e_j\to0,\qquad
 d_{{\rm ret},j}\to0,                                           \tag{36.10}
\]
then
\[
 |m_{j+1}-m_j|\to0,\qquad
 \vartheta_{S,j}\to0.                                            \tag{36.11}
\]
For every prescribed \(\epsilon_\chi>0\), the packet radius \(R\)
can be chosen once and the cofinal starting scale advanced so that
\[
                         \widehat\chi_{S,j}(R)\leq\epsilon_\chi.  \tag{36.12}
\]
\end{corollary}

\begin{proof}
The activity and branch stocks tend to zero by (35.11).  Insert this
and (36.10) into (36.9).  Estimate (36.7) then gives
\(\vartheta_{S,j}\leq L_S|m_{j+1}-m_j|+o(1)\).
The normalized resolvent stock
\(\chi_{S,j}=\vartheta_{S,j}/(1-\vartheta_{S,j})\) tends to zero.
Finally choose the slow packet radius so that the commutator tail in
(32.6) is at most \(\epsilon_\chi/2\), and advance the starting scale
until \(\chi_{S,j}\leq\epsilon_\chi/2\).
\end{proof}

\begin{firewall}[Vanishing increments do not determine the direction]
\label{fire:e8v36-beta-sign}
Equation (36.9) separates the moving-reference response from stable
deviation and activity corrections.  It does not prove that
\(b_{{\rm ref},j}(m_j)\to0\), that \(g_j\) moves toward zero in the
ultraviolet, or what its leading coefficient is.  Those facts require
the matched Yang--Mills beta function.
\end{firewall}

\subsection{A complete cofinal small-field contraction theorem}

Use a packet radius large enough for both
Definition~\ref{def:e8v34-R3} and (36.12).  Put
\[
\begin{split}
 \overline Q_j(\rho_{\rm st})
 :={}&
 \frac12\left[
 1+\sqrt3\,\widehat\chi_{S,j}(R)
 \right]\\
 &+
 N_{\rm typ}C_{\rm cyc}\widehat C_{{\rm CE},j}
 \kappa_j^{\rm WP}
 +2C_{\rm cyc}\rho_{\rm st}.                                  \tag{36.13}
\end{split}
\]
The projector-connection debit is absent in the parallel stable
fibre and the cumulant and retained-feedback remainders vanish by
(33.14) and (36.5).

\begin{theorem}[Cofinal aligned small-field closure]
\label{thm:e8v36-small-field-closure}
Assume:
\begin{enumerate}
\item the exact pullback chart and complete retained commutant of V9;
\item UPBC;
\item UALC and the complete good/bad packet activity card;
\item the connected-hull reserve cycle of V29--V31; and
\item a prescribed weak-coupling retained output path satisfying
      \(|m_{j+1}-m_j|\to0\).
\end{enumerate}
Then there exist a fixed \(\rho_{\rm st}>0\), a cofinal starting
scale \(j_0\), and \(q_{\rm st}<1\) such that
\[
 \sup_{j\geq j_0}\overline Q_j(\rho_{\rm st})
 \leq q_{\rm st}<1.                                             \tag{36.14}
\]
For every prescribed retained output path in this range, the aligned
stable map is a uniform strict contraction on its packet-polymer
ball.  Its fixed stable activity is unique at each step.
\end{theorem}

\begin{proof}
Choose \(\epsilon_\chi>0\) so that
\[
 \frac12(1+\sqrt3\,\epsilon_\chi)<1
\]
with a fixed positive margin.  The UPBC Lipschitz estimate and the
assumed vanishing increments
supply (36.12) cofinally.
Choose \(\rho_{\rm st}\) so that
\(2C_{\rm cyc}\rho_{\rm st}\) spends less than one third of the
remaining margin.  The marked conditional constant
\(\widehat C_{{\rm CE},j}\) is uniformly finite by (33.15), while
\(\kappa_j^{\rm WP}\to0\) by
Corollary~\ref{cor:e8v35-kappa}.  Advance \(j_0\) until their product
spends less than a second third of the margin.  This proves (36.14).

The exact removal of the other derivative rows is
\Cref{thm:e8v33-CE,thm:e8v36-exact-decoupling,
lem:e8v36-parallel-decoupling}.  The invariant-ball origin is
controlled by Theorem~\ref{thm:e8v35-origin}; advance \(j_0\) once
more so that (32.27) holds.  Banach's fixed-point theorem gives
existence and uniqueness.
\end{proof}

\begin{remark}[The fixed point is a scale map, not yet a trajectory]
\label{rem:e8v36-step-versus-path}
Theorem~\ref{thm:e8v36-small-field-closure} solves the stable
fixed-point problem conditional on the retained output at each
scale.  Compatibility of those outputs over infinitely many scales
is the retained trajectory problem below.
\end{remark}

\subsection{The moving physical reference}

The physical local representative differs from the pullback
representative by
\[
                         c_j(m)=-{\mathscr R}_j{\mathsf G}_j^{-1}m,
                                                                    \tag{36.15}
\]
whose quadratic part is not perturbatively small.  As in V10, absorb
it into the stationary moving reference \(h_j(m)\), including its
quadratic part in \(S_j(m)\), and write the stable deviation as
\[
                              e_j=q_j-h_j(m_j).                   \tag{36.16}
\]

\begin{theorem}[Nonautonomous tracking bound]
\label{thm:e8v36-tracking}
Suppose the reference graphs have a common contraction bound
\(q_{\rm st}<1\) and a uniform base Lipschitz constant \(L_h\).
Then
\[
 \|e_{j+1}\|
 \leq
 q_{\rm st}\|e_j\|
 +q_{\rm st}L_h|m_{j+1}-m_j|
 +\|d_{{\rm st},j}\|.                                            \tag{36.17}
\]
If the two forcing terms tend to zero, then \(e_j\to0\).
\end{theorem}

\begin{proof}
Subtract the fixed-point equation for \(h_{j+1}(m_{j+1})\) from the
physical recharted recursion.  Insert \(h_j(m_j)\), use the stable
contraction and the base Lipschitz bound, and collect any explicit
scale/regulator mismatch into \(d_{{\rm st},j}\).  This gives
(36.17).  Iteration is a geometric convolution of a sequence tending
to zero, exactly as in
Corollary~\ref{cor:e8v10-tracking}.
\end{proof}

\begin{corollary}[Coupled retained/stable tracking]
\label{cor:e8v36-tracking}
Assume the hypotheses of
Theorem~\ref{thm:e8v36-small-field-closure}, the response estimate
(36.9), and
\[
 \begin{gathered}
 b_{{\rm ref},j}(m_j)\to0,\qquad
 {\mathfrak z}_j\to0,\qquad
 d_{{\rm ret},j}\to0,\qquad d_{{\rm st},j}\to0,\\
 q_{\rm cpl}:=
 q_{\rm st}+q_{\rm st}L_hB_{\rm G}<1.
 \end{gathered}                                                   \tag{36.18}
\]
Then the physical stable activity tracks the moving pullback
reference and the retained increments vanish:
\[
 \|q_j-h_j(m_j)\|\longrightarrow0,\qquad
 |m_{j+1}-m_j|\longrightarrow0.                                 \tag{36.19}
\]
Consequently the covariance conclusion (36.12) holds under UPBC.
\end{corollary}

\begin{proof}
Insert (36.9) into (36.17).  On the fixed stable ball,
\[
\begin{split}
 \|e_{j+1}\|
 \leq{}&q_{\rm cpl}\|e_j\|\\
 &+q_{\rm st}L_h\left[
 |b_{{\rm ref},j}(m_j)|
 +B_{\rm ret}\frac{{\mathfrak z}_j}{1-q_j^{\rm WP}}
 (1+\rho_{\rm st})
 +d_{{\rm ret},j}\right]
 +d_{{\rm st},j}.
\end{split}                                                       \tag{36.19a}
\]
The second line tends to zero.  Since \(q_{\rm cpl}<1\), geometric
convolution gives \(e_j\to0\).  Equation (36.9) then gives the
vanishing retained increment.  Apply
Corollary~\ref{cor:e8v36-covariance-change}.
\end{proof}

\subsection{The exact remaining retained problem}

Restrict for clarity to the real reflection-even theta-zero sector.
After forbidden mass/source rows are removed by exact gauge and
reflection Ward identities, write the dynamical retained coordinate
as
\[
                         u_j=g_j^2,                              \tag{36.20}
\]
together with finite normalization/anisotropy coordinates
\(\sigma_j\).  The required matched normal form is
\[
\begin{split}
 u_{j+1}&=u_j-b_{\rm YM}u_j^2+r_{u,j},\\
 \sigma_{j+1}&=A_\sigma\sigma_j+r_{\sigma,j},                    \tag{36.21}
\end{split}
\]
in the orientation in which increasing \(j\) points toward the
ultraviolet.

\begin{definition}[Yang--Mills retained trajectory card]
\label{def:e8v36-retained-card}
YMRTC requires:
\begin{enumerate}
\item \(b_{\rm YM}>0\) computed in the same block, gauge-fixing,
      ghost, determinant, and normalization scheme used by the
      exact measure;
\item a uniform remainder
      \(|r_{u,j}|\leq c_uu_j^3+
      c_{ux}u_j^2\|e_j\|\);
\item invertibility or backward stability for the finite
      normalization block \(A_\sigma\), with its prescribed
      renormalization conditions;
\item exact Ward invariance of the forbidden gauge-boson mass,
      source, theta-zero, and conservation rows; and
\item \(d_{{\rm ret},j},d_{{\rm st},j}\to0\) and the coupled
      small-gain inequality \(q_{\rm cpl}<1\) in (36.18).
\end{enumerate}
\end{definition}

\begin{theorem}[Conditional infinite weak-coupling trajectory]
\label{thm:e8v36-conditional-trajectory}
If YMRTC holds and the initial retained data are sufficiently small,
then there is a cofinal retained/stable trajectory with
\[
 u_j\asymp j^{-1},\qquad e_j\to0,                                \tag{36.22}
\]
and all cofinal small-field contraction estimates above hold
uniformly along it.
\end{theorem}

\begin{proof}
Choose \(\delta_u,\delta_e>0\) so small that
\[
 c_u\delta_u+c_{ux}\delta_e\leq\frac{b_{\rm YM}}2,
 \qquad
 \frac{3b_{\rm YM}}2\delta_u\leq\frac12.                         \tag{36.23a}
\]
On \(0<u_j\leq\delta_u\), \(\|e_j\|\leq\delta_e\), the marginal
recursion gives
\[
 \frac{b_{\rm YM}}2u_j^2
 \leq u_j-u_{j+1}
 \leq\frac{3b_{\rm YM}}2u_j^2.                                  \tag{36.23b}
\]
It follows inductively that \(u_j\) remains positive and decreases.
Moreover \(u_{j+1}\geq u_j/2\), so
\[
 \frac{b_{\rm YM}}2
 \leq
 \frac1{u_{j+1}}-\frac1{u_j}
 \leq3b_{\rm YM}.                                                \tag{36.23c}
\]
Summation proves \(u_j\asymp j^{-1}\).

The finite \(\sigma\) block is solved by its forward or backward
Green series according to its declared spectrum, and exact Ward
coordinates are transported rather than shot.  The reference beta
increment is therefore \(O(u_j^2)+o(1)\).  The coupled estimate
(36.19a), together with \(q_{\rm cpl}<1\), keeps \(e_j\) in its
small ball after a sufficiently advanced start and gives \(e_j\to0\).
Equation (36.9) then gives vanishing retained increments, closing the
bootstrap required by
Theorem~\ref{thm:e8v36-small-field-closure}.
\end{proof}

\begin{firewall}[The conditional trajectory theorem has one physical
gate]
\label{fire:e8v36-retained-gate}
YMRTC item 1 and the matched remainder in item 2 have not yet been
proved.  The universal perturbative sign is strong guidance, but a
rigorous programme must derive the coefficient and its uniform
remainder from the same finite regulator before invoking
Theorem~\ref{thm:e8v36-conditional-trajectory}.
\end{firewall}

\subsection{Status after V36}

\begin{proposition}[Resolution of V32 item (iv) in the stable
derivative]
\label{prop:e8v36-item-iv}
In the exact aligned pullback chart, item (iv) of (32.30) contributes
zero to the stable derivative.  Retained variation survives only as
the response (36.2), covariance change (36.7), and moving-reference
forcing (36.17).
\end{proposition}

\begin{proof}
Use
\Cref{thm:e8v36-exact-decoupling,
lem:e8v36-parallel-decoupling,thm:e8v36-tracking}.
\end{proof}

\begin{assessment}[Progress and next upstream acquisition]
\label{ass:e8v36-status}
All four V32 small-field ledger entries are now either estimated or
removed exactly, and
Theorem~\ref{thm:e8v36-small-field-closure} gives the resulting
cofinal contraction theorem under its declared analytic, precision,
activity, and defect cards.  The next bottleneck is no longer the
polymer Banach space.  It is the matched retained recursion: derive
the positive \(u^2\) coefficient and a uniform
\(O(u^3+u^2\|e\|)\)
remainder while preserving the exact Ward ledger.  V37 should work
upstream on the background-field determinant identity that produces
this coefficient.
\end{assessment}

\begin{firewall}[Claim boundary after V36]
\label{fire:e8v36-claim}
V36 proves conditional cofinal small-field closure and exact
retained/stable triangularization.  It does not prove YMRTC, the
global good/bad trajectory at all scales, regulator removal,
Osterwalder--Schrader reconstruction, a positive continuum spectral
gap, or the full Yang--Mills mass gap.
\end{firewall}

\subsection{Parent record}

This V36 note was formed by appending the preceding material to V35,
whose validated record was
\[
\begin{array}{c|c}
\text{lines}&16150\\
\text{bytes}&562063\\
\text{SHA-256}&
\texttt{D6A8A03D083D9551BDE8AA925D4B7DEDA121C05E78D9D9FA1BF51F7CE8ABDC80}.
\end{array}
\]

% =============================================================================
% END APPENDED EIGHTH-SERIES V36 MATERIAL
% =============================================================================

% =============================================================================
% BEGIN APPENDED EIGHTH-SERIES V37 MATERIAL
% =============================================================================

\section{Eighth-series V37: the background determinant and the
retained \(u^2\) coefficient}

\subsection{Conventions and orientation}

Let \(G\) be compact and simple, with invariant inner product
\(\langle\cdot,\cdot\rangle_{\mathfrak g}\).  Normalize the adjoint
Casimir by
\[
 -\operatorname{tr}_{\rm ad}
 \bigl(\operatorname{ad}X\,\operatorname{ad}Y\bigr)
 =
 C_A\langle X,Y\rangle_{\mathfrak g},
 \qquad C_A>0.                                                   \tag{37.1}
\]
Write \(u=g^2\), and normalize the background action as
\[
 S_{\rm bg}(\overline A)
 =
 \frac1{4u}\int
 |\overline F_{\mu\nu}|_{\mathfrak g}^2\,dx.                    \tag{37.2}
\]
One downward Wilsonian step integrates momenta between
\(\Lambda/L\) and \(\Lambda\), \(L>1\).  The V36 index \(j\) points
in the reverse, ultraviolet direction.  This orientation distinction
will fix the sign without convention by slogan.

\subsection{The flat four-dimensional heat coefficient}

For a Laplace-type operator
\[
 P=-(\nabla^\mu\nabla_\mu+E)                                    \tag{37.3}
\]
on a flat four-torus, let
\(\Omega_{\mu\nu}=[\nabla_\mu,\nabla_\nu]\).

\begin{lemma}[Local logarithmic heat coefficient]
\label{lem:e8v37-heat}
After the field-independent coefficient is removed, the coefficient
of \(t^0\) in \(\operatorname{Tr}e^{-tP}\) is
\[
 a_4(P)=
 \frac1{(4\pi)^2}\int
 \operatorname{tr}\left(
 \frac1{12}\Omega_{\mu\nu}\Omega_{\mu\nu}
 +\frac12E^2
 \right)dx                                                       \tag{37.4}
\]
up to an integrated total derivative, which vanishes on the torus.
\end{lemma}

\begin{proof}
Choose normal coordinates and a synchronous bundle frame at a point
\(x\).  The diagonal heat kernel has a symbol expansion
\[
 K_P(t;x,x)
 =
 \int_{\mathbb R^4}\frac{d^4\xi}{(2\pi)^4}
 e^{-t|\xi|^2}
 \left[I+tE+t^2Q_4(x,\xi)+O(t^3)\right].
\]
Iterating Duhamel's formula twice, commuting covariant derivatives
only after their ordered positions have been fixed, gives
\[
 \int\frac{d^4\xi}{(2\pi)^4}e^{-t|\xi|^2}Q_4
 =
 \frac1{(4\pi t)^2}t^2
 \operatorname{tr}\left(
 \frac12E^2+\frac1{12}\Omega_{\mu\nu}\Omega_{\mu\nu}
 +\frac16\nabla^2E
 \right).
\]
The Gaussian moments used here are
\[
 \int e^{-t|\xi|^2}\xi_\mu\xi_\nu\,d^4\xi
 =
 \frac{\delta_{\mu\nu}}{2t}
 \int e^{-t|\xi|^2}\,d^4\xi
\]
and the analogous Wick pairing for four \(\xi\)'s.  The ordered
commutator terms have coefficient \(1/12\), while the two
multiplications by \(E\) have coefficient \(1/2\).  Integration over
the torus kills \(\nabla^2E\).  This proves (37.4).  The omitted
remainder contributes positive powers of \(t\) and cannot affect the
logarithmic coefficient.
\end{proof}

\begin{remark}[Zero modes do not change the local coefficient]
\label{rem:e8v37-zero-modes}
Replacing determinants by primed determinants removes only a finite
number of global eigenvalues.  It changes the large-\(t\) heat trace,
not the local small-\(t\) coefficient (37.4).  Torons and harmonic
coordinates nevertheless remain in the retained ledger.
\end{remark}

\subsection{Gauge and ghost operators}

In background Feynman gauge the rescaled quadratic operators are
\[
\begin{split}
 (\Delta_1a)_\mu
 &=
 -\overline D^2a_\mu
 -2\,\operatorname{ad}(\overline F_{\mu\nu})a_\nu,\\
 \Delta_0c&=-\overline D^2c,                                    \tag{37.5}
\end{split}
\]
acting on adjoint one-forms and adjoint scalar ghosts.  The one-loop
functional is
\[
 \Gamma_1(\overline A)
 =
 \frac12\log\det{}'\Delta_1
 -\log\det{}'\Delta_0.                                          \tag{37.6}
\]

\begin{lemma}[Gauge-minus-ghost trace algebra]
\label{lem:e8v37-trace-algebra}
With the convention (37.1),
\[
\begin{split}
 a_4(\Delta_1)\big|_{\overline F^2}
 &=
 \frac1{(4\pi)^2}\frac53 C_A
 \int|\overline F_{\mu\nu}|^2dx,\\
 a_4(\Delta_0)\big|_{\overline F^2}
 &=
 -\frac1{(4\pi)^2}\frac1{12}C_A
 \int|\overline F_{\mu\nu}|^2dx,                                \tag{37.7}\\
 \left(\frac12a_4(\Delta_1)-a_4(\Delta_0)\right)
 \big|_{\overline F^2}
 &=
 \frac1{(4\pi)^2}\frac{11}{12}C_A
 \int|\overline F_{\mu\nu}|^2dx.
\end{split}
\]
\end{lemma}

\begin{proof}
For the scalar ghost, \(E_0=0\) and
\(\Omega_{\mu\nu}=\operatorname{ad}\overline F_{\mu\nu}\).
Thus (37.1) and (37.4) give the second line.

For the one-form operator, the connection curvature acts diagonally
on four vector components.  Its curvature term is therefore
\[
 \frac4{12}
 \operatorname{tr}_{\rm ad}
 (\operatorname{ad}\overline F_{\mu\nu})^2.
\]
In the convention (37.3),
\[
 (E_1)_{\mu\nu}
 =2\,\operatorname{ad}\overline F_{\mu\nu}.
\]
Antisymmetry in \(\mu,\nu\) gives
\[
 \operatorname{tr}_{\Lambda^1\otimes{\rm ad}}E_1^2
 =
 -4\operatorname{tr}_{\rm ad}
 (\operatorname{ad}\overline F_{\mu\nu})^2.
\]
Consequently the bracket in (37.4) is
\[
 \left(\frac13-2\right)
 \operatorname{tr}_{\rm ad}
 (\operatorname{ad}\overline F_{\mu\nu})^2
 =
 \frac53C_A|\overline F_{\mu\nu}|^2.
\]
This proves the first line.  Weighting it by \(1/2\) and subtracting
the scalar ghost line gives
\(5/6+1/12=11/12\).
\end{proof}

\begin{firewall}[The ghost sign is part of antiscreening]
\label{fire:e8v37-ghost-sign}
The adjoint connection matrices are skew in convention (37.1), so
the scalar heat coefficient in (37.7) is negative.  The minus sign
in front of the Grassmann ghost determinant then adds
\(1/12\) to \(5/6\).  Dropping either sign destroys the coefficient
\(11/12\).
\end{firewall}

\subsection{The one-loop shell coefficient}

The proper-time interval corresponding to a scale-\(L\) shell has
logarithmic length
\[
 \log\frac{(\Lambda/L)^{-2}}{\Lambda^{-2}}=2\log L.              \tag{37.8}
\]
Because \(\log\det P=-\int dt\,t^{-1}\operatorname{Tr}e^{-tP}\),
the downward shell has the opposite sign to the heat coefficient.

\begin{theorem}[Background-field shell coefficient]
\label{thm:e8v37-shell}
The gauge-plus-ghost determinant of one downward continuum shell
contributes
\[
 \Gamma_1^{\downarrow}(\overline A)
 \big|_{\overline F^2}
 =
 -\frac{11C_A}{6(4\pi)^2}\log L
 \int|\overline F_{\mu\nu}|^2dx.                                \tag{37.9}
\]
Equivalently,
\[
 \left(\frac1u\right)_{\rm IR}
 =
 \left(\frac1u\right)_{\rm UV}
 -
 \frac{22C_A}{3(4\pi)^2}\log L+O(u).                            \tag{37.10}
\]
In the reverse ultraviolet orientation of V36,
\[
 \left(\frac1{u_{j+1}}\right)
 =
 \left(\frac1{u_j}\right)
 +b_{\rm YM}+O(u_j),\qquad
 \boxed{\quad
 b_{\rm YM}:=
 \frac{22C_A}{3(4\pi)^2}\log L>0.
 \quad}                                                         \tag{37.11}
\]
\end{theorem}

\begin{proof}
Multiply the last line of (37.7) by the proper-time length
\(2\log L\) and by the minus sign in the logarithmic determinant
representation.  This gives (37.9).  Comparing with the coefficient
\((4u)^{-1}\) in (37.2) multiplies the result by four and gives
(37.10).  Reversing the step gives (37.11).
\end{proof}

\begin{corollary}[The \(u^2\) normal form]
\label{cor:e8v37-u-normal}
If the matched inverse-coupling recursion has
\[
 \frac1{u_{j+1}}
 =
 \frac1{u_j}+b_{\rm YM}+\rho_j,\qquad
 |\rho_j|\leq C(u_j+\|e_j\|),                                   \tag{37.12}
\]
then, for sufficiently small \(u_j,\|e_j\|\),
\[
 u_{j+1}
 =
 u_j-b_{\rm YM}u_j^2+r_{u,j},\qquad
 |r_{u,j}|
 \leq
 C'\bigl(u_j^3+u_j^2\|e_j\|\bigr).                              \tag{37.13}
\]
\end{corollary}

\begin{proof}
Write
\[
 u_{j+1}
 =
 \frac{u_j}{1+u_j(b_{\rm YM}+\rho_j)}.
\]
The identity
\[
 \frac1{1+x}=1-x+\frac{x^2}{1+x}
\]
gives the leading term \(-b_{\rm YM}u_j^2\).
The term \(-\rho_ju_j^2\) is bounded by
\(C(u_j^3+u_j^2\|e_j\|)\), and the last rational remainder is
\(O(u_j^3)\) on a fixed small ball.
\end{proof}

\subsection{What must be matched by the exact lattice block}

The preceding calculation is local and continuum.  The exact RG
programme uses compact links, a specific block kernel, gauge
quotient, Jacobian, chart owners, and a finite lattice regulator.
The equality of logarithmic coefficients must therefore be a
theorem with hypotheses.

\begin{definition}[Background matching and remainder card]
\label{def:e8v37-BMRC}
BMRC for the exact scale-\(L\) block consists of:
\begin{enumerate}
\item a background-covariant gauge fixing whose quadratic link and
      ghost Hessians satisfy the exact regulated Ward identity;
\item removal of gauge zero modes into the retained
      toron/cohomology ledger;
\item convergence, on every smooth constant-curvature test
      background, of the normalized lattice gauge and ghost shell
      traces to (37.5), with the \(F^2\) trace error \(o(1)\);
\item cancellation of gauge-fixing parameter dependence in the
      retained \(F^2\) projection;
\item a uniform third-order determinant/cumulant bound which, after
      the one-loop term is subtracted, gives
\[
 |\rho_j|\leq C(u_j+\|e_j\|);                                   \tag{37.14}
\]
      and
\item chart-boundary, Haar, bad-component, and interpolation
      responses included in the same bound, rather than omitted from
      the determinant.
\end{enumerate}
\end{definition}

\begin{theorem}[BMRC closes the dynamical Yang--Mills marginal row]
\label{thm:e8v37-BMRC}
If BMRC holds uniformly on the cofinal trajectory, then the
reference response in (36.8) has the normal form
\[
 b_{{\rm ref},u,j}(u_j)
 =
 -b_{\rm YM}u_j^2+O(u_j^3),                                     \tag{37.15}
\]
with \(b_{\rm YM}\) given by (37.11), while stable deviations add
only \(O(u_j^2\|e_j\|)\).  Thus YMRTC items 1 and 2 hold.
\end{theorem}

\begin{proof}
Items 1--4 identify the retained lattice \(F^2\) shell coefficient
with Theorem~\ref{thm:e8v37-shell}; all finite local scheme changes
are assigned to nonlogarithmic normalization coordinates.  Item 5
gives (37.12) with the stated uniformity, and item 6 makes the
remainder exhaustive.  Apply
Corollary~\ref{cor:e8v37-u-normal}.
\end{proof}

\begin{firewall}[Coefficient universality is not the matching proof]
\label{fire:e8v37-universality}
The value (37.11) is the universal target for any admissible
background-covariant regulator.  V37 has proved it for the
continuum background shell and proved the implication
BMRC \(\Rightarrow\) YMRTC.  It has not proved BMRC item 3 or the
uniform remainder (37.14) for the exact compact-link block.  Merely
calling the coefficient universal would not supply those estimates.
\end{firewall}

\subsection{A finite symbol test for the matching coefficient}

Let \(\widehat H_{1,a}(k;\overline F)\) and
\(\widehat H_{0,a}(k;\overline F)\) be the lattice background
Hessians at spacing \(a\), including the chosen shell cutoff
\(\chi_{a,L}(k)\).  Define the retained curvature coefficient by
\[
\begin{split}
 {\cal B}_a(F):=
 \left.\frac{d^2}{dt^2}\right|_{t=0}
 \frac1{|\Lambda_a|}
 \bigg[
 &\frac12\operatorname{Tr}
 \chi_{a,L}\log\widehat H_{1,a}(k;tF)\\
 &-\operatorname{Tr}
 \chi_{a,L}\log\widehat H_{0,a}(k;tF)
 \bigg].                                                        \tag{37.16}
\end{split}
\]
The trace includes momentum, Lorentz, colour, alias, and block-fibre
indices; the zero-mode projector is removed.

\begin{proposition}[Dominated symbol matching criterion]
\label{prop:e8v37-symbol-test}
Suppose:
\begin{enumerate}
\item for almost every normalized shell momentum,
      \(a^{-2}\widehat H_{r,a}\) and its first two background
      derivatives converge to the continuum symbols of
      \(\Delta_r\), \(r=0,1\);
\item the inverses on the physical shell obey a common integrable
      domination after the retained zero modes are removed; and
\item alias and Brillouin-edge contributions have a gauge-invariant
      \(F^2\) projection converging to a finite local counterterm
      with no logarithmic scale dependence.
\end{enumerate}
Then
\[
 {\cal B}_a(F)\longrightarrow
 -\frac{11C_A}{3(4\pi)^2}\log L\,|F|^2,                         \tag{37.17}
\]
the second variation of (37.9).  In particular BMRC item 3 reduces
to these explicit finite-symbol conditions.
\end{proposition}

\begin{proof}
Twice differentiating a finite-dimensional log determinant gives
\[
 D^2\log\det H
 =
 \operatorname{Tr}\left(
 H^{-1}D^2H-H^{-1}(DH)H^{-1}(DH)
 \right),                                                       \tag{37.18}
\]
the identity already recorded in the manuscript summary.  Apply it
to both terms of (37.16).  Hypotheses 1--2 permit dominated
convergence of the normalized momentum sum to the continuum shell
integral.  The local logarithmic part of that integral is (37.9);
two variation directions \(F\) remove the factor \(1/2\) from
\((tF)^2\), giving (37.17).  Hypothesis 3 assigns all remaining
finite ultraviolet pieces to the declared local scheme
counterterm, so they do not change the logarithmic shell
coefficient.
\end{proof}

\begin{firewall}[The Brillouin edge cannot simply be discarded]
\label{fire:e8v37-edge}
Pointwise convergence near \(k=0\) does not control the full lattice
determinant.  The Brillouin-edge and alias sectors occupy nonzero
phase volume and may carry local \(F^2\) terms.  Proposition
\ref{prop:e8v37-symbol-test} requires their logarithmic part to be
identified, not assumed absent.
\end{firewall}

\subsection{Status after V37}

\begin{assessment}[Progress and next upstream acquisition]
\label{ass:e8v37-status}
V37 computes the positive coefficient required by YMRTC, including
the gauge/ghost algebra and RG orientation, and proves the exact
inverse-coupling-to-\(u^2\) conversion with the desired
\(O(u^3+u^2\|e\|)\) remainder.  For the exact block, the remaining
work is concentrated in BMRC:
\begin{enumerate}
\item prove the dominated lattice background-symbol limit including
      aliases and Brillouin-edge modes;
\item prove the uniform third-order determinant/cumulant remainder;
      and
\item verify the finite normalization/anisotropy and Ward rows.
\end{enumerate}
V38 should attack the first item by separating the universal
small-momentum shell from the finite local edge counterterm and
proving a scale-difference cancellation for the latter.
\end{assessment}

\begin{firewall}[Claim boundary after V37]
\label{fire:e8v37-claim}
The continuum one-loop coefficient is proved, but its exact
compact-link regulator matching and nonperturbative uniform
remainder are not.  Consequently V37 does not yet establish YMRTC,
the infinite weak-coupling trajectory, regulator removal, OS
reconstruction, a positive continuum gap, or the full Yang--Mills
mass gap.
\end{firewall}

\subsection{Parent record}

This V37 note was formed by appending the preceding material to V36,
whose validated record was
\[
\begin{array}{c|c}
\text{lines}&16673\\
\text{bytes}&580770\\
\text{SHA-256}&
\texttt{2FA479F6FC730DC1CCB96E060A2C5C161E4A78502B3202F6BBDF99B94D04452F}.
\end{array}
\]

% =============================================================================
% END APPENDED EIGHTH-SERIES V37 MATERIAL
% =============================================================================

% =============================================================================
% BEGIN APPENDED EIGHTH-SERIES V38 MATERIAL
% =============================================================================

\section{Eighth-series V38: scale-difference cancellation at the
Brillouin edge}

\subsection{The correct object is a difference of two self-similar
polarizations}

Let \(\Pi_a(p)\) be the scalar retained \(F^2\) coefficient in the
background two-point function at lattice spacing \(a\), after the
exact background Ward projector has removed longitudinal and
gauge-fixing rows.  Hypercubic corrections of order \(a^2p^4\) are
stable higher-derivative coordinates and are not included in
\(\Pi_a\).

\begin{definition}[Self-similar polarization expansion]
\label{def:e8v38-SPE}
SPE holds when, for \(0<a|p|\leq p_0\),
\[
 \Pi_a(p)
 =
 b_0\log\frac1{a^2|p|^2}
 +c_{\rm lat}
 +r(ap),\qquad
 r(q)\longrightarrow0\quad(q\to0),                              \tag{38.1}
\]
where
\[
                         b_0:=\frac{11C_A}{3(4\pi)^2}.           \tag{38.2}
\]
The constant \(c_{\rm lat}\) may depend on the regulator, block
kernel, gauge fixing, and local normalization scheme, but not on
\(a\).
\end{definition}

\begin{theorem}[Edge constants cancel in a scale difference]
\label{thm:e8v38-edge-cancel}
If SPE holds for a self-similar regulator, then
\[
\boxed{\qquad
 \Pi_a(p)-\Pi_{La}(p)
 =
 2b_0\log L+r(ap)-r(Lap).
\qquad}                                                         \tag{38.3}
\]
Hence
\[
 \lim_{a|p|\to0}
 \bigl[\Pi_a(p)-\Pi_{La}(p)\bigr]
 =
 \frac{22C_A}{3(4\pi)^2}\log L=b_{\rm YM}.                       \tag{38.4}
\]
\end{theorem}

\begin{proof}
Subtract (38.1) at \(a\) and \(La\).  The regulator-dependent local
constant cancels exactly, while
\[
 \log\frac1{a^2|p|^2}
 -
 \log\frac1{L^2a^2|p|^2}
 =2\log L.
\]
The two remainder terms tend to zero.
\end{proof}

\begin{firewall}[A finite edge term need not vanish]
\label{fire:e8v38-edge-finite}
The theorem does not assert \(c_{\rm lat}=0\).  It asserts that the
same constant occurs at two consecutive self-similar resolutions.
Using different gauge fixings, block normalizations, or alias
conventions at the two scales would prevent the cancellation and
would require an explicit finite coupling reparametrization.
\end{firewall}

\subsection{Low/edge splitting of the one-loop lattice integral}

After rescaling loop momentum to the Brillouin torus
\({\mathbb B}=[-\pi,\pi]^4\), write the one-loop retained integrand
as \({\cal I}(k,q)\), where \(q=ap\).  Fix
\[
                         0<2|q|<\delta<1                         \tag{38.5}
\]
and choose a smooth partition
\(\chi_{\rm low,\delta}+\chi_{\rm edge,\delta}=1\), with the low
cutoff supported in \(|k|<2\delta\) and equal to one on
\(|k|\leq\delta\).

\begin{definition}[Lattice one-loop matching hypotheses]
\label{def:e8v38-LOMH}
LOMH consists of:
\begin{enumerate}
\item the background Ward identity makes the retained two-point
      integrand transverse and cancels power divergences between
      gauge, ghost, Haar, and gauge-fixing rows;
\item on the low region, after the Ward subtraction at \(q=0\),
      \({\cal I}(k,q)\) and its required derivatives converge to the
      continuum background integrand and have a common integrable
      majorant;
\item on the edge region all physical and ghost denominators have a
      uniform spectral floor, and \({\cal I}(k,q)\) is \(C^2\) in
      \(q\) with an integrable derivative majorant;
\item every alias whose shifted symbol can vanish is reassigned to
      the low region of that alias; all remaining aliases satisfy
      the edge floor; and
\item the regulator, cutoff profile, and alias dictionary are
      repeated self-similarly at spacings \(a\) and \(La\).
\end{enumerate}
\end{definition}

\begin{lemma}[Free Wilson edge floor]
\label{lem:e8v38-Wilson-floor}
For
\[
 \widehat k_\mu:=2\sin(k_\mu/2),\qquad
 \widehat k^2:=\sum_{\mu=1}^4\widehat k_\mu^2,                   \tag{38.6}
\]
one has
\[
 |k|\geq\delta
 \quad\Longrightarrow\quad
 \widehat k^2\geq\frac{\delta^2}{\pi^2}.                         \tag{38.7}
\]
Thus the free Feynman-gauge vector and ghost symbols have a uniform
edge inverse.  The same holds for a finite nonvanishing alias after
its possible zero is reassigned as in LOMH item 4.
\end{lemma}

\begin{proof}
At least one component satisfies \(|k_\mu|\geq\delta/2\).
For \(|x|\leq\pi\),
\[
                         |\sin(x/2)|\geq |x|/\pi.
\]
The corresponding summand in (38.6) is therefore at least
\(4(\delta/(2\pi))^2=\delta^2/\pi^2\).  A finite alias shift is
handled on its translated Brillouin chart; away from its zero the
same compactness argument gives a positive floor.
\end{proof}

\begin{lemma}[Analyticity of the edge contribution]
\label{lem:e8v38-edge-analytic}
Under LOMH items 3--4,
\[
 \Pi_{\rm edge,\delta}(q)
 =
 c_{\rm edge}(\delta)
 +O_\delta(|q|^2)                                                \tag{38.8}
\]
after the transverse \(F^2\) projection.  The constant is independent
of \(a\) under item 5.
\end{lemma}

\begin{proof}
The uniform edge floor permits differentiation of each finite
matrix inverse and logarithm in \(q\).  The integrable derivative
majorant permits differentiation under the Brillouin integral.
The Ward and reflection symmetries make the scalar transverse
coefficient even in \(q\), so its linear Taylor coefficient vanishes.
Taylor's theorem gives (38.8).  Once \(k\) and \(q\) are dimensionless,
the edge integral contains no remaining \(a\); self-similarity makes
its constant the same at both resolutions.
\end{proof}

\begin{lemma}[Universal low-region logarithm]
\label{lem:e8v38-low-log}
Under LOMH items 1--2,
\[
 \Pi_{\rm low,\delta}(q)
 =
 b_0\log\frac{\delta^2}{|q|^2}
 +c_{\rm low}(\delta)+o_{q\to0}(1),                              \tag{38.9}
\]
with \(b_0\) given by (38.2).
\end{lemma}

\begin{proof}
The Ward subtraction removes the power-divergent Taylor pieces and
leaves a logarithmically homogeneous continuum integrand.  Dominated
convergence applies away from \(k=0\).  On the annulus
\(2|q|\leq|k|\leq\delta\), radial scaling gives a constant angular
coefficient times
\[
 \int_{2|q|}^{\delta}\frac{dr}{r}
 =
 \log\frac{\delta}{2|q|}.
\]
The angular gauge-minus-ghost coefficient is the background heat
coefficient computed in
Theorem~\ref{thm:e8v37-shell}; expressed as a coefficient of
\(\log(1/|q|^2)\), it is \(b_0\).  The inner region
\(|k|<2|q|\) becomes a fixed integral after \(k=|q|\ell\), and the
cutoff-overlap region is finite.  Both contribute to
\(c_{\rm low}(\delta)\) and the vanishing remainder, proving (38.9).
\end{proof}

\begin{theorem}[LOMH implies the self-similar expansion]
\label{thm:e8v38-LOMH}
Under LOMH,
\[
 \Pi_a(p)
 =
 b_0\log\frac1{a^2|p|^2}
 +c_{\rm lat}+r(ap),\qquad r(ap)\to0.                            \tag{38.10}
\]
Thus SPE and the coefficient limit (38.4) hold.
\end{theorem}

\begin{proof}
Add (38.8) and (38.9), with \(q=ap\).  The artificial
\(\delta\)-dependence cancels because the sum is the original
cutoff-independent integral.  More explicitly, changing \(\delta\)
moves an integral over a compact annulus from one term to the other
with opposite signs.  The remaining sum of constants is
\[
 c_{\rm lat}:=
 c_{\rm low}(\delta)+c_{\rm edge}(\delta)+b_0\log\delta^{-2},
\]
which is independent of \(\delta\) and, by self-similarity,
independent of \(a\).  The remainders tend to zero.  Apply
Theorem~\ref{thm:e8v38-edge-cancel}.
\end{proof}

\begin{remark}[What is proved for the free Wilson denominator]
\label{rem:e8v38-free-versus-background}
Lemma~\ref{lem:e8v38-Wilson-floor} verifies the denominator floor
for the free Wilson gauge and ghost symbols.  LOMH still requires
the exact background vertices, Haar row, block aliases, and Ward
subtractions to be written in the chosen regulator.  V38 has not
silently promoted the free floor to that full identity.
\end{remark}

\subsection{A volume-free third derivative of a log determinant}

Let \(\tau_\Lambda(A):=|\Lambda|^{-1}\operatorname{Tr}A\) be the
normalized trace of a translation-covariant block matrix with
internal dimension at most \(d_*\).  Use any submultiplicative
weighted row algebra \(\|\cdot\|_{\rm row}\) containing the shell
inverse and its local derivatives.

\begin{lemma}[Normalized trace versus row norm]
\label{lem:e8v38-trace-row}
\[
                         |\tau_\Lambda(A)|
 \leq d_*\|A\|_{\rm row},                                        \tag{38.11}
\]
uniformly in \(\Lambda\).
\end{lemma}

\begin{proof}
The normalized trace is the average of the diagonal internal traces.
Each is bounded by \(d_*\|A_{xx}\|\), and the diagonal block norm is
at most the absolute row sum.
\end{proof}

\begin{theorem}[Third log-determinant row]
\label{thm:e8v38-logdet-third}
Let \(H(t)\) be three times differentiable and invertible on a shell.
If
\[
 \|H^{-1}\|_{\rm row}\leq K_0,\qquad
 \|H^{(r)}\|_{\rm row}\leq M_r\quad(1\leq r\leq3),                \tag{38.12}
\]
then
\[
\boxed{\qquad
 \left|
 \frac{d^3}{dt^3}\tau_\Lambda\log H(t)
 \right|
 \leq
 d_*\left(
 K_0M_3+3K_0^2M_1M_2+2K_0^3M_1^3
 \right).
\qquad}                                                         \tag{38.13}
\]
\end{theorem}

\begin{proof}
Differentiate \(H^{-1}H'\) twice, using
\((H^{-1})'=-H^{-1}H'H^{-1}\), and use cyclicity of the trace:
\[
\begin{split}
 \frac{d^3}{dt^3}\operatorname{Tr}\log H
 =
 \operatorname{Tr}\bigl[
 &H^{-1}H'''
 -3H^{-1}H'H^{-1}H''\\
 &+2H^{-1}H'H^{-1}H'H^{-1}H'
 \bigr].
\end{split}                                                       \tag{38.14}
\]
Apply Lemma~\ref{lem:e8v38-trace-row}, submultiplicativity, and
(38.12).
\end{proof}

\begin{corollary}[Determinant Taylor remainder is volume free]
\label{cor:e8v38-logdet-remainder}
If the constants in (38.12) are uniform for \(|t|\leq t_0\), then
the second-order Taylor remainder of the normalized log determinant
is at most
\[
 \frac{|t|^3}{6}d_*
 \left(K_0M_3+3K_0^2M_1M_2+2K_0^3M_1^3\right),                  \tag{38.15}
\]
independently of the number of lattice sites.
\end{corollary}

\begin{proof}
Use Taylor's theorem with integral remainder and
Theorem~\ref{thm:e8v38-logdet-third}.
\end{proof}

\subsection{Even-coupling remainder after the one-loop coefficient}

Let \({\cal R}_j(g,e)\) be the complete retained inverse-coupling
response after the constant one-loop value \(b_{\rm YM}\) has been
separated.

\begin{lemma}[Even Taylor gate]
\label{lem:e8v38-even-remainder}
Suppose
\[
 {\cal R}_j(-g,0)={\cal R}_j(g,0),\qquad
 {\cal R}_j(0,0)=0,                                              \tag{38.16}
\]
and, uniformly on the declared ball,
\[
 \|\partial_g^2{\cal R}_j(g,0)\|\leq M_g,\qquad
 \|D_e{\cal R}_j(g,e)\|\leq M_e.                                \tag{38.17}
\]
Then
\[
 \boxed{\qquad
 \|{\cal R}_j(g,e)\|
 \leq\frac{M_g}{2}g^2+M_e\|e\|
 =\frac{M_g}{2}u+M_e\|e\|.
\qquad}                                                         \tag{38.18}
\]
\end{lemma}

\begin{proof}
Evenness gives \(\partial_g{\cal R}_j(0,0)=0\).
Taylor's theorem in \(g\) gives the first term; the mean-value
theorem along \(te\), \(0\leq t\leq1\), gives the second.
\end{proof}

\begin{proposition}[Concrete remainder reduction]
\label{prop:e8v38-remainder-reduction}
Assume:
\begin{enumerate}
\item the exact compact-link inversion/charge-conjugation symmetry
      implies (38.16) for the retained \(F^2\) response;
\item gauge and ghost shell Hessians satisfy the row bounds (38.12);
\item differentiated Haar, chart, bad, and conditional-cumulant
      owners satisfy the two-mark packet bounds of V33; and
\item all these derivative rows use the same retained projection.
\end{enumerate}
Then (38.17) holds with finite volume-independent constants, and
therefore BMRC remainder (37.14) follows.
\end{proposition}

\begin{proof}
The determinant rows are bounded by
\Cref{thm:e8v38-logdet-third,cor:e8v38-logdet-remainder}.
The non-determinant conditional rows are finite by
Theorem~\ref{thm:e8v33-marked-resolvent}, with their local mark
derivatives supplied by UALC.  The finite retained projection and
type sum preserve uniformity.  These give \(M_g,M_e<\infty\) in
(38.17).  Lemma~\ref{lem:e8v38-even-remainder} yields
\[
 |\rho_j|\leq C(u_j+\|e_j\|),
\]
which is (37.14).
\end{proof}

\begin{firewall}[Symmetry must be checked on the complete response]
\label{fire:e8v38-evenness}
The Wilson plaquette action alone has the required inversion
symmetry, but Proposition~\ref{prop:e8v38-remainder-reduction}
requires it after gauge fixing, ghosts, Haar Jacobian, chart
transitions, bad owners, and the retained projection are combined.
An odd term in an incomplete coordinate chart may cancel only after
those rows are assembled.
\end{firewall}

\subsection{Status after V38}

\begin{assessment}[Progress and next acquisition]
\label{ass:e8v38-status}
V38 replaces the vague Brillouin-edge condition by two precise
mechanisms:
\begin{enumerate}
\item under LOMH, the edge is an analytic, scale-independent local
      constant which cancels from the self-similar scale difference,
      leaving the coefficient (38.4);
\item under the row and complete-evenness hypotheses, the normalized
      determinant and conditional remainder satisfies the BMRC
      bound (37.14).
\end{enumerate}
The free Wilson denominator floor is proved.  The next upstream gate
is the complete background Ward identity and its transverse
projection for the actual link, ghost, Haar, alias, and gauge-fixing
symbols.  V39 should write that finite regulated identity before
attempting any further beta-flow conclusion.
\end{assessment}

\begin{firewall}[Claim boundary after V38]
\label{fire:e8v38-claim}
V38 proves the scale-difference cancellation and volume-free
determinant remainder criteria under explicit matching hypotheses.
It does not yet verify the complete regulated background Ward
identity, LOMH, or evenness for the actual compact-link block.
Therefore YMRTC, the continuum construction, and the full mass gap
remain unproved.
\end{firewall}

\subsection{Parent record}

This V38 note was formed by appending the preceding material to V37,
whose validated record was
\[
\begin{array}{c|c}
\text{lines}&17135\\
\text{bytes}&595887\\
\text{SHA-256}&
\texttt{ED66F96B4973F9101591382377418FE1F104F83A5E79AE50A0BBA577FFC94D7D}.
\end{array}
\]

% =============================================================================
% END APPENDED EIGHTH-SERIES V38 MATERIAL
% =============================================================================

% =============================================================================
% BEGIN APPENDED EIGHTH-SERIES V39 MATERIAL
% =============================================================================

\section{Eighth-series V39: the exact finite-regulator background
Ward ledger}

\subsection{Equivariant compact-link disintegration}

Let \(\overline U\) be a background link field on a finite regulator
graph and \(q\) a fluctuation coordinate on the physical-detail
fibre.  Assume the link reconstruction
\[
                         U=\Sigma(\overline U,q)                 \tag{39.1}
\]
is equivariant under a background gauge transformation \(h\):
\[
 \Sigma(\overline U^h,q^h)=\Sigma(\overline U,q)^h.              \tag{39.2}
\]
The fibre density includes the Wilson action, background gauge
fixing, Faddeev--Popov determinant, Haar-coordinate Jacobian, block
constraint, and the good/bad chart partition:
\[
\begin{split}
 d\mu_{\overline U}(q):={}&
 e^{-S(\Sigma(\overline U,q))}
 \Delta_{\rm FP}(\overline U,q)
 e^{-S_{\rm gf}(\overline U,q)}
 J_{\rm Haar}(\overline U,q)\\
 &\times{\bf1}_{B(\Sigma(\overline U,q))=\overline U_{\rm c}}\,
 d q,
                                                                  \tag{39.3}
\end{split}
\]
with chart-boundary pieces understood as the complete owner sum.

\begin{definition}[Exact background-covariance card]
\label{def:e8v39-EBC}
EBC requires:
\begin{enumerate}
\item (39.2) and equivariance of the block constraint;
\item invariance of the Wilson action under \(U\mapsto U^h\);
\item adjoint covariance of the gauge-fixing function and
      Faddeev--Popov operator;
\item invariance of the Haar Jacobian and fluctuation measure under
      \(q\mapsto q^h\);
\item a gauge-invariant chart partition, or a complete permutation of
      chart/bad owners under \(h\); and
\item no anomalous Jacobian in the pure bosonic/ghost finite
      regulator.
\end{enumerate}
\end{definition}

\begin{theorem}[Exact background invariance of the coarse action]
\label{thm:e8v39-background-invariance}
Under EBC, the conditional partition function and coarse effective
action
\[
 Z(\overline U):=\int d\mu_{\overline U}(q),\qquad
 \Gamma(\overline U):=-\log Z(\overline U)                       \tag{39.4}
\]
satisfy
\[
                         Z(\overline U^h)=Z(\overline U),\qquad
                         \Gamma(\overline U^h)=\Gamma(\overline U).
                                                                    \tag{39.5}
\]
This identity is exact at finite cutoff and includes the good/bad
owner sum.
\end{theorem}

\begin{proof}
In \(Z(\overline U^h)\), make the change of fluctuation variables
\(q'=q^h\).  Items 1--4 of EBC make every factor in (39.3) equal to
the corresponding factor at \((\overline U,q)\), and item 6 makes
the Jacobian one.  Item 5 either fixes each chart owner or permutes
the complete owner sum.  Thus the integrals agree.  Taking
\(-\log\) proves the second identity.
\end{proof}

\begin{firewall}[A single local chart need not be invariant]
\label{fire:e8v39-chart}
Theorem~\ref{thm:e8v39-background-invariance} applies to the complete
chart/bad disintegration.  A single logarithm chart may cross its
branch boundary under a large gauge transformation.  Discarding the
permuted boundary owners would invalidate the change of variables.
\end{firewall}

\subsection{Differential Ward identity}

Let \(d_{\overline U}\omega\) denote the infinitesimal tangent to the
background gauge orbit generated by a site Lie-algebra field
\(\omega\).

\begin{corollary}[Finite-regulator background Ward identity]
\label{cor:e8v39-Ward}
For every background in a differentiable chart,
\[
                         D\Gamma(\overline U)
 [d_{\overline U}\omega]=0.                                     \tag{39.6}
\]
At the trivial background, if \(\Gamma^{(1)}(1)=0\), the Hessian
\(\Pi:=\Gamma^{(2)}(1)\) obeys
\[
                         \Pi d_0=0,\qquad d_0^*\Pi=0.             \tag{39.7}
\]
\end{corollary}

\begin{proof}
Differentiate (39.5) along \(h(t)=\exp(t\omega)\) to obtain
(39.6).  Differentiate once more in an arbitrary link direction at
the trivial background.  The derivative of
\(d_{\overline U}\omega\) multiplies \(\Gamma^{(1)}(1)\) and
vanishes by the stated centering.  The remaining term is
\(\Pi d_0\omega=0\).  Self-adjointness of the Hessian gives the
adjoint identity.
\end{proof}

\begin{lemma}[The one-point function vanishes in pure Yang--Mills]
\label{lem:e8v39-one-point}
If \(G\) is compact and simple and the regulator preserves global
adjoint colour symmetry and link reflection, then
\[
                              \Gamma^{(1)}(1)=0.                  \tag{39.8}
\]
\end{lemma}

\begin{proof}
The one-point tensor is an invariant vector in the adjoint
representation, multiplied by a link-direction vector.  A simple
Lie algebra has no nonzero invariant adjoint vector.  Equivalently,
averaging the tensor over constant gauge transformations projects
onto the zero invariant subspace.  Link reflection supplies the same
conclusion for any orientation-odd residual.
\end{proof}

Use link-centred Fourier variables
\[
                         s_\mu(p):=2\sin(p_\mu/2).                \tag{39.9}
\]
After harmless link-centre phases are removed, (39.7) becomes
\[
 s_\mu(p)\widehat\Pi_{\mu\nu}^{ab}(p)=0,\qquad
 \widehat\Pi_{\mu\nu}^{ab}(p)s_\nu(p)=0.                         \tag{39.10}
\]

\subsection{Uniqueness of the quadratic hypercubic tensor}

\begin{theorem}[Transverse hypercubic classification at order two]
\label{thm:e8v39-hypercubic}
Suppose \(\widehat\Pi(p)\) is analytic at \(p=0\), self-adjoint,
colour invariant, reflection even, hypercubic covariant, and
satisfies (39.10).  Then, for a compact simple group,
\[
\boxed{\qquad
 \widehat\Pi_{\mu\nu}^{ab}(p)
 =
 \delta^{ab}c_\Pi
 \bigl(|p|^2\delta_{\mu\nu}-p_\mu p_\nu\bigr)
 +O(|p|^4).
\qquad}                                                         \tag{39.11}
\]
In particular there is no gauge-boson mass term and there is only one
reflection-even dimension-four quadratic retained coefficient.
\end{theorem}

\begin{proof}
Colour invariance makes the adjoint bilinear form a scalar multiple
of the invariant metric, hence gives \(\delta^{ab}\).  Reflection
evenness removes odd powers of \(p\).  A constant hypercubic
two-tensor is \(m^2\delta_{\mu\nu}\).  Substitution in (39.10) for
small nonzero \(p\) gives \(m^2p_\nu+O(|p|^3)=0\), hence \(m^2=0\).

At quadratic order the most general symmetric hypercubic tensor has
diagonal and off-diagonal entries
\[
\begin{split}
 K_{\nu\nu}^{(2)}(p)&=A|p|^2+Bp_\nu^2,\\
 K_{\mu\nu}^{(2)}(p)&=Cp_\mu p_\nu\qquad(\mu\ne\nu).             \tag{39.12}
\end{split}
\]
Contracting with \(p_\mu\) gives
\[
 p_\nu\left[(A+C)|p|^2+(B-C)p_\nu^2\right]=0.
\]
Independence of \(|p|^2\) and \(p_\nu^2\) gives
\(C=-A\) and \(B=C=-A\).  Thus
\[
 K_{\mu\nu}^{(2)}
 =A(|p|^2\delta_{\mu\nu}-p_\mu p_\nu).
\]
Since \(s_\mu(p)=p_\mu+O(|p|^3)\), replacing the exact lattice
incidence by \(p\) changes only the \(O(|p|^4)\) remainder.
\end{proof}

\begin{corollary}[No power-divergent field-dependent counterterm]
\label{cor:e8v39-no-power}
In four dimensions, the local divergent background action compatible
with EBC, hypercubic symmetry, and reflection has no
field-dependent operator of dimension below four.  At dimension four
its quadratic parity-even part is a multiple of
\(\int|F|^2\); the parity-odd theta density is absent in the
reflection-even sector.
\end{corollary}

\begin{proof}
A dimension-two quadratic mass is excluded by
Theorem~\ref{thm:e8v39-hypercubic}.  A dimension-one or
dimension-three scalar cannot be made from a pure connection while
preserving background gauge invariance.  At dimension four,
the quadratic classification gives \(F^2\); reflection excludes
\operatorname{tr}(F\wedge F)\).
\end{proof}

\begin{corollary}[LOMH power-subtraction row]
\label{cor:e8v39-LOMH1}
Under EBC and the analyticity/hypercubic hypotheses of
Theorem~\ref{thm:e8v39-hypercubic}, LOMH item 1 holds: the complete
gauge, ghost, Haar, gauge-fixing, and chart-owner sum is transverse,
and all field-dependent power-divergent Taylor coefficients vanish.
\end{corollary}

\begin{proof}
Exact transversality is (39.7).  The absence of allowed
field-dependent local terms below dimension four is
Corollary~\ref{cor:e8v39-no-power}.  Since the statement concerns the
complete effective action (39.4), the cancellation includes every
row in its defining density.
\end{proof}

\subsection{Finite retained Ward rows}

\begin{theorem}[Ward closure of the forbidden retained coordinates]
\label{thm:e8v39-retained-Ward}
Assume EBC, full hypercubic symmetry, and reflection.  In the pure
Yang--Mills theta-zero sector:
\begin{enumerate}
\item the gauge-boson source and mass coordinates remain zero;
\item no dimension-four quadratic anisotropy independent of \(F^2\)
      is generated;
\item the parity-odd theta coordinate remains zero; and
\item finite gauge-orbit conservation identities are transported
      exactly by (39.6).
\end{enumerate}
Thus the Ward-invariance part of YMRTC item 4 is closed at finite
regulator.
\end{theorem}

\begin{proof}
The source vanishes by Lemma~\ref{lem:e8v39-one-point}; the mass and
anisotropy conclusions are
Theorem~\ref{thm:e8v39-hypercubic}.  Reflection preserves the
theta-zero subspace.  Equation (39.6) is the exact infinitesimal
conservation identity at every background.
\end{proof}

\begin{firewall}[Hypercubic symmetry is required at the blocking
level]
\label{fire:e8v39-block-symmetry}
If a block singles out the Euclidean time direction, separate
electric and magnetic \(F^2\) coefficients are allowed until
rotation restoration is proved.  Theorem
\ref{thm:e8v39-retained-Ward} applies to a fully hypercubic block or
to a symmetrized block cycle whose exact effective action has that
symmetry.
\end{firewall}

\subsection{Gauge parameter and Gribov boundary are separate}

Background covariance holds for every covariant gauge parameter, but
does not by itself show that the retained \(F^2\) coefficient is
parameter independent.

\begin{proposition}[Local BRST gauge-parameter implication]
\label{prop:e8v39-gauge-parameter}
Suppose on the complete regulated fibre:
\begin{enumerate}
\item there is a nilpotent BRST differential \(s\) preserving the
      measure and all owner boundaries;
\item the gauge-parameter derivative of the density is
      \(s\)-exact; and
\item the retained \(F^2\) insertion is \(s\)-closed.
\end{enumerate}
Then the retained \(F^2\) response is gauge-parameter independent.
\end{proposition}

\begin{proof}
If \(\xi\) is the gauge parameter and
\(\partial_\xi d\mu=s(\Psi_\xi d\mu)\), then
\[
 \partial_\xi\langle{\cal O}_{F^2}\rangle
 =
 \int s(\Psi_\xi{\cal O}_{F^2}\,d\mu)=0.
\]
The last equality is finite-dimensional BRST integration by parts;
there is no boundary term by item 1, and
\(s{\cal O}_{F^2}=0\) by item 3.
\end{proof}

\begin{firewall}[The BRST boundary hypothesis is nontrivial]
\label{fire:e8v39-Gribov}
A compact nonabelian gauge fibre can have Gribov copies, chart
boundaries, and cancellations in a global gauge-fixing formula.
EBC proves background covariance without solving those problems.
Proposition~\ref{prop:e8v39-gauge-parameter} is therefore a
conditional implication; its boundary hypothesis has not been
verified for the complete compact-link disintegration.
\end{firewall}

\subsection{Evenness of the complete retained response}

\begin{definition}[Complete inversion involution]
\label{def:e8v39-involution}
CII is an involution \({\cal C}\) of the full regulated fibre and
owner labels such that
\[
 {\cal C}:\quad(g,q,\overline A)\longmapsto
 (-g,-q,{\cal C}\overline A),                                   \tag{39.13}
\]
preserves the full Wilson/gauge-fixing/ghost/Haar/block density,
and leaves the retained \(F^2\) projection invariant.
\end{definition}

\begin{proposition}[CII proves the V38 evenness row]
\label{prop:e8v39-evenness}
Under CII,
\[
                         {\cal R}_j(-g,0)={\cal R}_j(g,0),        \tag{39.14}
\]
so the symmetry hypothesis (38.16) holds for the complete response.
\end{proposition}

\begin{proof}
In the response integral at \(-g\), make the change of variables and
owner labels supplied by \({\cal C}\).  The density and integration
domain are preserved, while the retained \(F^2\) observable is fixed.
The resulting integral is the response at \(g\).
\end{proof}

\begin{firewall}[Plaquette inversion alone is insufficient]
\label{fire:e8v39-involution-complete}
The identity
\(\operatorname{ReTr}U_p=\operatorname{ReTr}U_p^{-1}\)
checks only the Wilson factor.  CII must also pair the gauge-fixing,
ghost, Haar, block-kernel, chart-boundary, and bad-component owners.
Until that pairing is written, (39.14) remains conditional.
\end{firewall}

\subsection{Status after V39}

\begin{assessment}[Progress and V40 direction]
\label{ass:e8v39-status}
V39 proves, from exact finite-regulator background covariance:
\begin{enumerate}
\item the complete Ward identity and transverse lattice Hessian;
\item absence of a gauge-boson source, mass, and field-dependent
      power divergence;
\item uniqueness of the parity-even hypercubic \(F^2\) coefficient;
      and
\item closure of the corresponding YMRTC Ward rows.
\end{enumerate}
This resolves LOMH item 1, conditional only on EBC and the stated
analyticity/symmetry of the finite effective action.  Gauge-parameter
independence and complete \(g\)-evenness are separated into explicit
BRST-boundary and involution cards rather than inferred from the
Ward identity.  V40 is a compulsory companion audit; afterward it
should write the exact low-region Wilson/ghost background vertices
needed for LOMH item 2.
\end{assessment}

\begin{firewall}[Claim boundary after V39]
\label{fire:e8v39-claim}
V39 closes the finite Ward tensor classification but not the
background-symbol convergence, BRST boundary card, complete
inversion involution, beta remainder, global trajectory, continuum
limit, or mass gap.
\end{firewall}

\subsection{Parent record}

This V39 note was formed by appending the preceding material to V38,
whose validated record was
\[
\begin{array}{c|c}
\text{lines}&17570\\
\text{bytes}&610408\\
\text{SHA-256}&
\texttt{52D2F797C796220EFD1E7231966512B7ECF87E13E23691DB161702E2523C0262}.
\end{array}
\]

% =============================================================================
% END APPENDED EIGHTH-SERIES V39 MATERIAL
% =============================================================================

% =============================================================================
% BEGIN APPENDED EIGHTH-SERIES V40 MATERIAL
% =============================================================================

\section{Eighth-series V40: companion audit and complex packet-response
jets}

\subsection{Compulsory V40 companion audit}

The active companion files inspected at this checkpoint were
\[
\begin{array}{c|r|r|l}
\text{file}&\text{lines}&\text{bytes}&\text{SHA--256}\\ \hline
\texttt{Renewal\_SM\_Einstein\_Continuum\_Research\_Notes\_V25.tex}
&10052&382021&
\texttt{C1AB2B98B067D7720058F63EC252C656DBFD1C745876CE92DDF691029AE9C341}
\\
\texttt{Renewal\_NS\_Stability\_Research\_Notes\_V26.tex}
&5320&278283&
\texttt{82C887F8C2C69E4F28FAD7C3DC8DE5B9099CB5A4BF431EBA1FFF3E91935978AD}.
\end{array}                                                       \tag{40.1}
\]
The NS note is unchanged.  The SM note has advanced from V18 to V25.
Its relevant new results are: double Gevrey-one parameter/spatial
jets for a uniformly gapped analytic covariance; analytic Cauchy
control after an all-size forest/tree sum; large-field
good/transition/bad coercivity criteria; and exact finite interface
determinant factorization.

\begin{assessment}[V40 audit transfer]
\label{ass:e8v40-audit}
The following transfer is valid.
\begin{enumerate}
\item A complex-uniform branch inequality permits Cauchy estimates
      after the connected packet-tree sum.  This supplies the
      parameter derivatives required by the V38 beta remainder
      without differentiating every tree separately.
\item Analytic covariance-jet estimates apply only to the gapped
      shell/detail inverse.  Harmonic, zero-shell, and closing modes
      must remain retained, exactly as in the YM packet ledger.
\item Extensive log determinants must be assigned to local
      vacuum/block owners before volume-uniform Cauchy estimates are
      stated.
\end{enumerate}
The SM constraint-interface factorization may become useful at the
later OS gluing stage, but it does not change the present beta-symbol
calculation.  No SM or NS numerical constant is imported.
\end{assessment}

\begin{decision}[Audit-adjusted V40 direction]
\label{dec:e8v40-direction}
Before returning to the explicit Wilson background vertices, V40
closes the analytic implication
\[
 \text{complex packet branch card}
 \Longrightarrow
 \text{uniform response jets}
 \Longrightarrow
 \text{BMRC remainder}.                                          \tag{40.2}
\]
The remaining acquisition will then be the verification of that card
and LOMH for the actual compact-link block.
\end{decision}

\subsection{Analytic shell inverse}

Let \(z=(z_1,\ldots,z_p)\in{\mathbb C}^p\) collect the complexified
coupling, retained background, and finitely many marked source
parameters.  Let \(H(z)\) be a physical shell/detail Hessian, and put
\[
 E(z):=
 H(0)^{-1/2}\bigl(H(z)-H(0)\bigr)H(0)^{-1/2}.                    \tag{40.3}
\]

\begin{lemma}[Weighted analytic inverse]
\label{lem:e8v40-analytic-inverse}
Suppose \(E(z)\) is analytic on the polydisc
\({\mathbb D}_\rho^p\) and
\[
 \sup_{z\in{\mathbb D}_\rho^p}
 \|E(z)\|_{\rm row}\leq\eta<1.                                  \tag{40.4}
\]
Then \(H(z)^{-1}\) is analytic there and
\[
 \|H(0)^{1/2}H(z)^{-1}H(0)^{1/2}\|_{\rm row}
 \leq\frac1{1-\eta}.                                             \tag{40.5}
\]
For every smaller polydisc with radii
\(0<\rho_{0,a}<\rho_a\),
\[
 \|\partial_z^\alpha
 [H(0)^{1/2}H^{-1}H(0)^{1/2}](0)\|_{\rm row}
 \leq
 \frac{\alpha!}{1-\eta}
 \prod_{a=1}^p\rho_{0,a}^{-\alpha_a}.                           \tag{40.6}
\]
\end{lemma}

\begin{proof}
The normalized inverse is the row-algebra Neumann series
\[
 (I+E(z))^{-1}=\sum_{n\geq0}(-E(z))^n,
\]
which converges uniformly on the polydisc by (40.4).  This proves
analyticity and (40.5).  Apply the Banach-valued multivariate Cauchy
formula on the smaller distinguished boundary to obtain (40.6).
\end{proof}

\begin{firewall}[Only gapped shell/detail modes use this lemma]
\label{fire:e8v40-gapped-only}
If \(H(0)\) has a closing eigenvalue, the normalization in (40.3)
already diverges.  Lemma~\ref{lem:e8v40-analytic-inverse} is used
only after harmonic, zero-shell, and block-boundary modes have been
placed in the retained sector.
\end{firewall}

\subsection{Complex packet branch card}

Let \(z_{\gamma,j}^{\rm WP}(\zeta)\) be the complexified activity of
type \(\gamma\), and define
\({\mathbb K}_j^{\rm WP}(\zeta)\) by (33.6), now with absolute
complex moduli.

\begin{definition}[Complex packet branch card]
\label{def:e8v40-CPBC}
CPBC holds on a polydisc \({\mathbb D}_\rho^p\) when:
\begin{enumerate}
\item every good, transition, determinant, ghost, Haar, and bad
      activity is analytic in \(\zeta\);
\item its absolute packet row is locally uniformly summable there;
\item for one fixed positive type weight \(v\),
\[
 \sup_{j,\zeta\in{\mathbb D}_\rho^p}
 \max_\alpha
 \frac{({\mathbb K}_j^{\rm WP}(\zeta)v)_\alpha}{v_\alpha}
 \leq q_{\mathbb C}<1;                                          \tag{40.7}
\]
\item marked trunks are analytic with
      \(\sup_{\zeta}A_{r,j}^{\rm WP}(\zeta)\leq A_{r,\mathbb C}\);
      and
\item on complex chart tubes the real part of the bad-component
      action retains an extensive lower bound strong enough to sum
      its lattice-animal entropy.
\end{enumerate}
\end{definition}

\begin{firewall}[Real Peierls smallness does not imply CPBC]
\label{fire:e8v40-complex-bad}
A real positive bad-component weight may grow after complexification.
CPBC item 5 requires a coercive complex tube, possibly narrower than
the local Taylor disc.  Cauchy estimates cannot use a complex domain
on which the bad-polymer sum diverges.
\end{firewall}

\begin{theorem}[Analyticity after the complete packet-tree sum]
\label{thm:e8v40-connected-analytic}
Under CPBC, every \(r\)-marked connected response
\({\cal C}_{r,j}(\zeta)\), for fixed \(r\), is Banach-valued analytic
on \({\mathbb D}_\rho^p\), and
\[
 \sup_{\zeta\in{\mathbb D}_\rho^p}
 \|{\cal C}_{r,j}(\zeta)\|
 \leq
 \frac{A_{r,\mathbb C}}{(1-q_{\mathbb C})^r}.                    \tag{40.8}
\]
For every strictly smaller polydisc,
\[
\boxed{\qquad
 \|\partial_\zeta^\alpha{\cal C}_{r,j}(0)\|
 \leq
 \alpha!\left[
 \prod_{a=1}^p\rho_{0,a}^{-\alpha_a}
 \right]
 \frac{A_{r,\mathbb C}}{(1-q_{\mathbb C})^r}.
\qquad}                                                         \tag{40.9}
\]
The bounds are uniform in volume and \(j\).
\end{theorem}

\begin{proof}
For each finite selected tree, its integrand is analytic by CPBC
items 1 and 4.  The absolute selected-parent estimate used in
Theorem~\ref{thm:e8v33-marked-resolvent} is uniform on the complex
polydisc by (40.7), and gives the summable majorant (40.8).  Uniform
absolute convergence of analytic Banach-valued functions implies
analyticity of the sum.  Apply the multivariate Cauchy formula on a
smaller distinguished boundary to obtain (40.9).  Packet incidence,
owner, hull, and reblocking constants are volume independent by
V27--V31.
\end{proof}

\begin{remark}[No second factorial from tree size]
\label{rem:e8v40-one-factorial}
The factorial in (40.9) is the Cauchy factorial in parameter order.
The arbitrary tree-size sum has already been absorbed by
\((1-q_{\mathbb C})^{-r}\); it does not contribute another
factorial.
\end{remark}

\subsection{Analytic determinant density}

Let the extensive field-independent determinant be assigned to the
vacuum/block-density coordinate, and let
\[
 {\cal D}_j(\zeta):=
 \tau_\Lambda\left[
 \log H_j(\zeta)-\log H_j(0)
 \right]                                                        \tag{40.10}
\]
be the normalized response density.

\begin{proposition}[Volume-uniform determinant Cauchy bound]
\label{prop:e8v40-det-Cauchy}
Under the hypotheses of
Lemma~\ref{lem:e8v40-analytic-inverse}, and uniform local derivative
rows for \(H(\zeta)\), \({\cal D}_j\) is analytic on the same
polydisc.  On every smaller polydisc,
\[
 |\partial_\zeta^\alpha{\cal D}_j(0)|
 \leq
 \alpha!
 \prod_{a=1}^p\rho_{0,a}^{-\alpha_a}\,M_{\cal D},                \tag{40.11}
\]
where \(M_{\cal D}<\infty\) is independent of total volume.
\end{proposition}

\begin{proof}
Along the radial path \(t\zeta\),
\[
 {\cal D}_j(\zeta)
 =
 \int_0^1
 \tau_\Lambda\left[
 H_j(t\zeta)^{-1}D H_j(t\zeta)[\zeta]
 \right]dt.                                                     \tag{40.12}
\]
The inverse is bounded by (40.5), the derivative is a finite local
row, and normalized trace is bounded by
Lemma~\ref{lem:e8v38-trace-row}.  Thus (40.12) is uniformly bounded
and analytic.  Cauchy's formula gives (40.11).  The removed
field-independent determinant is extensive but belongs to the
vacuum density and never enters \(M_{\cal D}\).
\end{proof}

\subsection{The beta remainder from complex analyticity}

Let \({\cal R}_j(g,e)\) be the complete inverse-coupling response
after subtracting the coefficient \(b_{\rm YM}\).  Use one complex
coordinate for \(g\) and a complex line in each tested stable
direction \(e\).

\begin{theorem}[CPBC analytic beta remainder]
\label{thm:e8v40-beta-remainder}
Assume CPBC, the determinant hypotheses of
Proposition~\ref{prop:e8v40-det-Cauchy}, and the complete inversion
involution CII of V39.  Then there are finite volume-independent
constants \(M_g,M_e\) such that
\[
 \|\partial_g^2{\cal R}_j(g,0)\|\leq M_g,\qquad
 \|D_e{\cal R}_j(g,e)\|\leq M_e.                                \tag{40.13}
\]
Consequently
\[
 \boxed{\qquad
 \|{\cal R}_j(g,e)\|
 \leq\frac{M_g}{2}u+M_e\|e\|,
\qquad}                                                         \tag{40.14}
\]
and the BMRC remainder (37.14) holds.
\end{theorem}

\begin{proof}
The complete response is the finite retained projection of the sum
of its determinant density and connected conditional owners.
Proposition~\ref{prop:e8v40-det-Cauchy} and (40.9) bound its second
\(g\)-derivative and first stable derivative uniformly, giving
(40.13).  CII gives the evenness (39.14), hence
\(\partial_g{\cal R}_j(0,0)=0\).  Apply
Lemma~\ref{lem:e8v38-even-remainder} to obtain (40.14).
\end{proof}

\begin{corollary}[Exact status of BMRC item 5]
\label{cor:e8v40-BMRC5}
BMRC item 5 follows from CPBC, the weighted analytic shell inverse,
and CII.  No separate term-by-term all-order determinant expansion is
needed.
\end{corollary}

\begin{proof}
This is Theorem~\ref{thm:e8v40-beta-remainder}, followed by
Corollary~\ref{cor:e8v37-u-normal}.
\end{proof}

\subsection{What remains after the audit transfer}

\begin{assessment}[Progress and next acquisition after V40]
\label{ass:e8v40-status}
The audit closes the logical route from a complex-uniform packet
branch bound to the matched beta remainder.  The live regulator
tasks are now:
\begin{enumerate}
\item verify EBC/LOMH low-symbol convergence for the explicit
      Wilson, background-gauge, ghost, Haar, and alias vertices;
\item construct CII on the complete owner dictionary;
\item prove a complex bad-component tube giving CPBC item 5; and
\item verify the finite small-gain and normalization blocks in
      YMRTC.
\end{enumerate}
The first task determines the coefficient matching; the second and
third determine whether the analytic remainder theorem applies to
the complete compact-link measure.  V41 should return to the explicit
low-symbol expansion.
\end{assessment}

\begin{firewall}[Claim boundary after V40]
\label{fire:e8v40-claim}
V40 proves the analytic packet-response implication under CPBC and
records the compulsory companion audit.  It does not verify CPBC,
CII, the exact Wilson background symbols, the retained trajectory,
regulator removal, OS reconstruction, or the full mass gap.
\end{firewall}

\subsection{Parent record}

This V40 note was formed by appending the preceding material to V39,
whose validated record was
\[
\begin{array}{c|c}
\text{lines}&17965\\
\text{bytes}&624930\\
\text{SHA-256}&
\texttt{248E1F5A0B0CA3EC2B059C25837DBDDDF3A1FA34589AAC3DBFFCDA7D07CB3CC7}.
\end{array}
\]

% =============================================================================
% END APPENDED EIGHTH-SERIES V40 MATERIAL
% =============================================================================

% =============================================================================
% BEGIN APPENDED EIGHTH-SERIES V41 MATERIAL
% =============================================================================

\section{Eighth-series V41: the parity-tree quotient determinant and the
correction of the unmatched ghost row}
\label{sec:e8v41}

\subsection{Why the coefficient ledger must be corrected here}

The explicit reduced-fibre programme does not introduce a
Faddeev--Popov field.  Its adapted parity forest first puts every
internal tree link equal to the identity and then writes the exact
conditional law as
\[
 Z_{\beta,\Psi,M}(W)
 =\int_{{\rm SU}(3)^{45M^4}}
       \exp\{-\beta V_M(z,W)-\Psi(z,W)\}\,dm_{\rm red}(z).       \tag{41.1}
\]
The measure in (41.1) is literal product Haar measure.  Ordered
tree elimination uses only left and right Haar translations, so its
Jacobian is one.  Thus the background-gauge vector-minus-ghost
calculation of V37 is the correct continuum comparison target, but
it was not yet the one-loop determinant of (41.1).

This distinction is structural.  A ghost determinant can be a
coordinate representation of a quotient determinant, but it cannot
be inserted into a regulator that was defined without it unless the
coordinate-change identity, including its orbit Jacobian, is proved.
V41 proves that finite-dimensional identity and identifies the
additional comparison still needed for the compact-link shell.

\begin{firewall}[Correction of the V37--V40 coefficient claim]
\label{fire:e8v41-ghost-correction}
Equations (37.2)--(37.8) compute the universal continuum
background-gauge target
\[
 \frac12\log\det L_1-\log\det L_0.                              \tag{41.2}
\]
They do not, by themselves, compute the background response of
(41.1).  Every occurrence of the EBC/BMRC ghost row in
V37--V40 is henceforth read as a \emph{quotient-matching row}:
one must derive (41.2) from the actual parity-tree quotient, or
compute its scale difference directly.  None of the analytic
remainder results is invalidated; their regulator input was simply
not yet verified.
\end{firewall}

\subsection{Linearized rooted-tree coordinates}

All spaces in the next theorem are finite-dimensional real Euclidean
spaces; one copy is understood for every colour component.  Let
\(G\) be the rooted internal vertex-gauge space, \(E\) the fine-link
tangent space, and
\[
 D:G\longrightarrow E                                             \tag{41.3}
\]
the infinitesimal internal gauge action.  Rooting at the coarse
vertices makes \(D\) injective.  Let \(F:E\to G\) read the oriented
tree-link coordinates, let \(T=\ker F\), and denote the inclusion by
\(I_T:T\hookrightarrow E\).  Define
\[
 Q:G\oplus T\longrightarrow E,
 \qquad Q(\phi,t)=D\phi+I_Tt,
 \qquad J_T:=|\det Q|.                                           \tag{41.4}
\]

\begin{lemma}[Rooted parity-tree unimodularity]
\label{lem:e8v41-tree-unimodular}
Order each nonroot vertex after its parent and order the corresponding
tree edge in the same way.  At the identity background,
\(FD\) is triangular with diagonal blocks \(\pm I_{\mathfrak{su}(3)}\).
Consequently
\[
 |\det(FD)|=J_T=1.                                                \tag{41.5}
\]
At an arbitrary compact-link background the diagonal blocks are
adjoint maps of \({\rm SU}(3)\), so their real determinants are also
one and (41.5) persists in left-trivialized Haar coordinates.
\end{lemma}

\begin{proof}
For a tree edge from the parent \(p(x)\) to the child \(x\), its
linear gauge variation is
\(\phi_x-\phi_{p(x)}\) at the identity, up to the fixed orientation
sign.  The parent precedes the child, so the child coefficient is the
diagonal identity block and all other vertex coefficients occur
strictly earlier.  The non-tree coordinate block of \(Q\) is the
identity: after tree coordinates are listed first, its matrix is
\[
 \begin{pmatrix}FD&0\\ *&I_T\end{pmatrix}.
\]
This proves the first assertion.  Away from the identity, left or
right trivialization replaces the diagonal identity blocks by
\({\rm Ad}_U\).  The compact connected group \({\rm SU}(3)\) acts
orthogonally and orientation-preservingly on its Lie algebra, hence
\(\det{\rm Ad}_U=1\).  The same triangular determinant argument
applies.  It is the tangent form of the exact ordered Haar
factorization in (41.1).
\end{proof}

\begin{remark}[No gauge-volume constant is hidden]
\label{rem:e8v41-no-volume}
The gauge transformations in \(G\) equal the identity at every
coarse root.  Thus \(D\) has no constant stabilizer kernel and the
normalized internal Haar integral is exactly one.  The residual
coarse gauge group is retained rather than divided out.
\end{remark}

\subsection{Exact quotient-to-ghost determinant identity}

Let \(K:E\to E\) be self-adjoint, suppose
\[
 KD=0,
 \qquad \ker K=\operatorname{ran}D                                \tag{41.6}
\]
on the tested subspace, and put
\[
 L_0:=D^*D,
 \qquad K_{\rm gf}:=K+DD^*,
 \qquad H_T:=I_T^*KI_T.                                          \tag{41.7}
\]
If physical harmonic modes are present, first declare a common
coordinate complement and transport its measure through every
quotient map; all three operators below are then restricted to that
measured complement and all determinants are primed.  An arbitrary
Euclidean pseudodeterminant is not interchangeable with this
pushforward determinant.  This convention is necessary on a periodic
flat lattice.

\begin{theorem}[Finite-dimensional quotient/FP bridge with the orbit
Jacobian displayed]
\label{thm:e8v41-quotient-FP}
Under (41.3)--(41.7),
\[
 \boxed{
 (\det H_T)^{-1/2}
 =J_T^{-1}(\det K_{\rm gf})^{-1/2}\det L_0.}                     \tag{41.8}
\]
Equivalently,
\[
 \boxed{
 \frac12\log\det H_T
 =\log J_T+\frac12\log\det K_{\rm gf}-\log\det L_0.}             \tag{41.9}
\]
For the parity-tree Haar coordinates, \(J_T=1\).  Hence the
tree-slice Gaussian determinant is exactly the standard
vector-minus-ghost determinant, although no ghost variable occurs in
the defining quotient measure.
\end{theorem}

\begin{proof}
Let \(P\) be the orthogonal projection onto
\(P_E=(\operatorname{ran}D)^\perp\), and define the isometry
\[
 R:=DL_0^{-1/2}:G\longrightarrow\operatorname{ran}D.             \tag{41.10}
\]
The map \(S:=PI_T:T\to P_E\) is invertible.  Indeed, if \(PI_Tt=0\),
then \(I_Tt=D\phi\); applying \(F\) gives \(0=FD\phi\), hence
\(\phi=0\) and \(t=0\).  Domain and range have the same dimension.

Write the coordinate map \(Q\) in the orthogonal decomposition
\(E=\operatorname{ran}D\oplus P_E\).  Its diagonal blocks are
\(L_0^{1/2}\), from \(D=RL_0^{1/2}\), and \(S\).  Therefore
\[
 J_T=(\det L_0)^{1/2}|\det S|.                                  \tag{41.11}
\]
Self-adjointness and \(KD=0\) imply \(K=PKP\).  If
\(K_P=K|_{P_E}\), then
\[
 H_T=S^*K_PS,
 \quad
 \det H_T=(\det S)^2\det K_P
 =J_T^2\frac{\det K_P}{\det L_0}.                               \tag{41.12}
\]
In the same orthogonal decomposition, \(K_{\rm gf}\) has blocks
\(L_0\) and \(K_P\), because \(R^*DD^*R=L_0\).  Thus
\[
 \det K_{\rm gf}=\det L_0\det K_P.                              \tag{41.13}
\]
Substitution of (41.13) into (41.12), followed by the positive
square root, proves (41.8) and (41.9).  The primed version follows
after restricting every map to the stated common complement.
\end{proof}

\begin{firewall}[The missing square root in informal FP arguments]
\label{fire:e8v41-orbit-root}
The coordinate change \(\phi\mapsto D\phi\) has metric Jacobian
\((\det L_0)^{1/2}\).  A delta-function gauge condition supplies
one inverse determinant and the FP factor supplies one determinant;
the remaining square root is precisely the orbit-coordinate
Jacobian in (41.11).  Dropping it gives the wrong quotient measure.
Equation (41.8), rather than a mnemonic cancellation, is the ledger
used below.
\end{firewall}

\begin{corollary}[Pointwise background version]
\label{cor:e8v41-background-bridge}
Let \(U_*(W)\) be a smooth stationary background and let \(D(W)\)
be the rooted internal gauge tangent at that background.  If the
physical Hessian \(K(W)\) kills \(\operatorname{ran}D(W)\), is
positive on the tested quotient, and its retained kernel has
constant dimension, then (41.8)--(41.9) hold pointwise in \(W\).
For the parity-tree coordinates \(D_W\log J_T(W)=0\) to every
order.
\end{corollary}

\begin{proof}
Apply Theorem~\ref{thm:e8v41-quotient-FP} at each \(W\).  The
background version of Lemma~\ref{lem:e8v41-tree-unimodular} gives
\(J_T(W)=1\).  Constant kernel dimension makes the tested quotient
bundle smooth, so differentiation of its primed determinant is
legitimate.
\end{proof}

\subsection{The reduced/coarse Schur split}

The parity-tree slice is not integrated in one step.  Its coordinates
split as
\[
 T=R_{\rm red}\oplus C,                                         \tag{41.14}
\]
where \(R_{\rm red}\) consists of the \(45M^4\) reduced links and
\(C\) contains the designated second links carrying the coarse
field.  Write the tree-slice Hessian as
\[
 H_T=
 \begin{pmatrix}
 H_{rr}&H_{rc}\\ H_{cr}&H_{cc}
 \end{pmatrix},
 \qquad
 S_c:=H_{cc}-H_{cr}H_{rr}^{-1}H_{rc}.                           \tag{41.15}
\]

\begin{theorem}[Exact Gaussian fibre elimination]
\label{thm:e8v41-gaussian-fibre}
If \(H_{rr}\succ0\), then for every retained coordinate \(c\),
\[
 \begin{split}
 &\int_{R_{\rm red}}\exp\left\{-\frac12
 \left\langle(r,c),H_T(r,c)\right\rangle\right\}\,dr\\
 &\hspace{2em}=(2\pi)^{\dim R_{\rm red}/2}(\det H_{rr})^{-1/2}
 \exp\{-\tfrac12\langle c,S_cc\rangle\},                        \tag{41.16}
 \end{split}
\]
and
\[
 \boxed{\det H_T=\det H_{rr}\det S_c.}                          \tag{41.17}
\]
If \(S_c\) has retained harmonic kernel, define the matched
Schur-coordinate determinant by
\(\det_{\rm Sch}'H_T:=\det H_{rr}\det{}'S_c\), where the complement
and its measure are declared in the \(c\)-coordinates.
\end{theorem}

\begin{proof}
Set \(r'=r+H_{rr}^{-1}H_{rc}c\).  Direct multiplication gives
\[
 \langle(r,c),H_T(r,c)\rangle
 =\langle r',H_{rr}r'\rangle+\langle c,S_cc\rangle.
\]
Translation preserves Lebesgue measure, and the standard positive
Gaussian integral proves (41.16).  Taking a block triangular
factorization,
\[
 H_T=
 \begin{pmatrix}I&0\\H_{cr}H_{rr}^{-1}&I\end{pmatrix}
 \begin{pmatrix}H_{rr}&0\\0&S_c\end{pmatrix}
 \begin{pmatrix}I&H_{rr}^{-1}H_{rc}\\0&I\end{pmatrix},
\]
proves (41.17).  The same triangular change has determinant one on
\(R_{\rm red}\oplus C_\perp\), which proves the asserted
Schur-coordinate identity.  It need not equal the Euclidean
pseudodeterminant of \(H_T\), because the lifted retained kernel can
be oblique.
\end{proof}

\begin{corollary}[Fine-minus-coarse form of the shell determinant]
\label{cor:e8v41-scale-difference}
Suppose the Schur form \(S_c\) is the quadratic coarse quotient form
with its own rooted-tree data \((K_c,D_c,J_c)\), while \(H_T\) is
the fine quotient form with data \((K_f,D_f,J_f)\).  On the common
nonretained complement, using the transported Schur-coordinate
measures just declared,
\[
 \begin{split}
 \frac12\log\det H_{rr}
 &=\frac12\log\det H_T-\frac12\log\det{}'S_c\\
 &=\log\frac{J_f}{J_c}
 +\left[\frac12\log\det{}'K_{f,{\rm gf}}-\log\det L_{0,f}\right]\\
 &\quad-
 \left[\frac12\log\det{}'K_{c,{\rm gf}}-\log\det L_{0,c}\right].
                                                                    \tag{41.18}
 \end{split}
\]
For the exact parity-tree Haar quotient, \(J_f=J_c=1\).
Thus its vertical Gaussian normalization is exactly a scale
difference of vector-minus-ghost determinant combinations.
\end{corollary}

\begin{proof}
Take logarithms in (41.17) in the matched Schur coordinates and apply
(41.9) separately to the fine and coarse quotient forms with the
same transported complement measures.  The common retained kernel is
omitted before either determinant is taken.
Lemma~\ref{lem:e8v41-tree-unimodular} removes the final Jacobian ratio
for the parity trees.
\end{proof}

\begin{assessment}[What (41.18) repairs and what it does not]
\label{ass:e8v41-repair-scope}
Equation (41.18) is the missing algebraic bridge between the actual
tree quotient and V37.  It proves that ghosts are a valid
\emph{comparison representation} of the Gaussian quotient response.
It does not yet prove that the background-dependent fine and coarse
operators in the compact Wilson shell converge to the continuum
symbols used in (37.2), nor that their nonlinear remainder is
uniform.  Those are analytic and symbol-level tasks, not determinant
bookkeeping identities.
\end{assessment}

\subsection{The exact vertical gap and normalization}

For the pure Wilson action at the flat background, the inherited
normalization gives
\[
 D^2(\beta V_M)(0)[a,b]
 =\frac{\beta}{3}\langle d_1a,d_1b\rangle.                      \tag{41.19}
\]
If \(\Sigma_{rr}\) is the reduced-link curl Gram matrix, set
\[
 A_{\rm vert}:=\frac13\Sigma_{rr},
 \qquad H_{rr}^{(0)}=\beta A_{\rm vert}.                         \tag{41.20}
\]
The attached reduced-fibre programme proves by exact incidence
elimination and a Sturm certificate that the number \(\lambda_*\)
is the unique root in \((1/9,3/20)\) of
\[
 \begin{split}
 q_8(\lambda)={}&\lambda^8-16\lambda^7+98\lambda^6-296\lambda^5
 +472\lambda^4\\
 &-396\lambda^3+168\lambda^2-32\lambda+2,
 \end{split}                                                     \tag{41.21}
\]
and that
\[
 0.11760418512646926265<\lambda_*
 <0.11760418512646926269,\qquad
 \Sigma_{rr}\succeq\lambda_*I.                                  \tag{41.22}
\]
Consequently
\[
 \boxed{
 A_{\rm vert}\succeq\frac{\lambda_*}{3}I,
 \qquad
 \|(H_{rr}^{(0)})^{-1}\|
 \leq\frac{3}{\beta\lambda_*}.}                                 \tag{41.23}
\]
This result is imported rather than reproved here; its exact replay
is already present in the attached fibre note.

\begin{proposition}[The coupling power is a vacuum owner]
\label{prop:e8v41-beta-vacuum}
Let \(n_r=45M^4\dim\mathfrak{su}(3)=360M^4\).  Then
\[
 \frac12\log\det H_{rr}^{(0)}
 =\frac{n_r}{2}\log\beta+\frac12\log\det A_{\rm vert}.           \tag{41.24}
\]
The first term is independent of the retained field and belongs
entirely to the local vacuum density.  Every retained background
derivative of it vanishes.
\end{proposition}

\begin{proof}
Equation (41.24) follows from
\(\det(\beta A)=\beta^{n_r}\det A\).  The first summand depends only
on \(\beta\) and on the locally owned reduced-link count, so no
retained derivative can see it.
\end{proof}

\begin{firewall}[The vertical gap is not the mass gap]
\label{fire:e8v41-vertical-not-mass}
The number \(\lambda_*/3\) gaps only the eliminated block-fibre
coordinates.  Coarse transverse modes still have a curvature symbol
vanishing like \(|q|^2\) at small momentum.  Equation (41.23) makes
the shell integration local; it does not supply a physical
four-dimensional Yang--Mills mass.
\end{firewall}

\subsection{Background derivatives of the actual quotient determinant}

Let \(z_*(W)\) be a selected stationary centre in the reduced fibre
and write
\[
 H_{rr}(W)
 :=D_z^2[\beta V_M+\Psi](z_*(W),W).                             \tag{41.25}
\]
Every derivative of \(H_{rr}(W)\) below is a total derivative, so it
includes the motion of \(z_*(W)\).  The Gaussian quotient
contribution to the effective action is
\[
 \Gamma_{1,{\rm tree}}(W):=\frac12\log\det H_{rr}(W).            \tag{41.26}
\]

\begin{lemma}[Exact second response of the vertical determinant]
\label{lem:e8v41-logdet-response}
On a branch on which \(H_{rr}(W)\succ0\), for retained directions
\(a,b\),
\[
 \boxed{
 D^2\Gamma_{1,{\rm tree}}(W)[a,b]
 =\frac12\Tr\left(
 H_{rr}^{-1}D^2H_{rr}[a,b]
 -H_{rr}^{-1}DH_{rr}[a]H_{rr}^{-1}DH_{rr}[b]
 \right).}                                                      \tag{41.27}
\]
If \(H_{rr}=\beta A_{rr}\), all explicit powers of \(\beta\)
cancel in (41.27).
\end{lemma}

\begin{proof}
Jacobi's formula gives
\(D\log\det H[a]=\Tr(H^{-1}DH[a])\).  Differentiate once more and
use
\(D(H^{-1})[b]=-H^{-1}DH[b]H^{-1}\); cyclicity of the trace yields
(41.27).  Substituting \(H=\beta A\) gives one \(\beta^{-1}\) for
each inverse and one \(\beta\) for every derivative of \(H\), so
the factors cancel term by term.
\end{proof}

\begin{corollary}[Volume-free response after owner normalization]
\label{cor:e8v41-volume-free-logdet}
Assume a background tube on which
\(H_{rr}(W)\succeq c\beta I\), and assume the first and second
background derivative kernels have uniformly summable local rows.
Then the normalized trace of (41.27) is bounded independently of
the total periodic volume.
\end{corollary}

\begin{proof}
The inverse row is bounded by \((c\beta)^{-1}\), with exponential
off-diagonal decay supplied by the inherited finite-propagation
vertical inverse theorem.  Each term of (41.27) is a product of
finitely many uniformly summable rows.  The normalized trace estimate
of V38 bounds it by the corresponding row norm, with no factor of the
total volume.  The vacuum term was already removed in
Proposition~\ref{prop:e8v41-beta-vacuum}.
\end{proof}

\begin{assessment}[Available versus missing background control]
\label{ass:e8v41-background-control}
The exact flat floor (41.23) and the inherited small-field curvature
window provide a nonempty local tube in which \(H_{rr}^{-1}\) is
uniform and exponentially localized.  What is still missing is a
single common tube that also carries the stationary centre, every
generated interaction in \(\Psi\), and the complete marked complex
owner bounds required by CPBC.
\end{assessment}

\subsection{Why the deterministic Wilson symbol is not enough}

For the exact conditional law (41.1), put
\[
 S(z,W):=\beta V_M(z,W)+\Psi(z,W),
 \qquad
 A(W):=-\log\int e^{-S(z,W)}dm_{\rm red}(z).                    \tag{41.28}
\]
The product Haar measure in these coordinates is independent of
\(W\).

\begin{theorem}[Exact direct-minus-score-covariance Hessian]
\label{thm:e8v41-partition-hessian}
For any retained directions \(a,b\) for which differentiation under
the compact integral is valid,
\[
 \boxed{
 D^2A(W)[a,b]
 =\mathbb E_W[D_W^2S[a,b]]
 -\operatorname{Cov}_W(D_WS[a],D_WS[b]).}                       \tag{41.29}
\]
No ghost term and no derivative of the reduced Haar measure occurs.
\end{theorem}

\begin{proof}
Compactness makes all smooth derivatives integrable.  First,
\(DA[a]=\mathbb E_W[DS[a]]\).  Differentiating a normalized
expectation gives
\[
 D\mathbb E_W[f][b]
 =\mathbb E_W[Df[b]]
 -\mathbb E_W[f\,DS[b]]
 +\mathbb E_W[f]\mathbb E_W[DS[b]].
\]
Take \(f=DS[a]\) to obtain (41.29).  Independence of
\(dm_{\rm red}\) from \(W\) follows from the exact Haar quotient,
so there is no measure derivative.
\end{proof}

At a flat translation-invariant law, the inherited deterministic
coarse second-jet calculation has the form
\[
 \mathbb E K^W(q)
 =\beta w_{\beta,\Psi,M}\,C(q)\otimes I_{\mathfrak g},
 \qquad
 C(q)=\frac13\bigl(7I_4-u(q)u(q)^*\bigr),                       \tag{41.30}
\]
where
\[
 u(q)=(e^{iq_1},\ldots,e^{iq_4})^{\mathsf T}.                   \tag{41.31}
\]
Since \(\|u(q)\|^2=4\), the direction-space eigenvalues of \(C(q)\)
are \(1\) on \(\mathbb Cu(q)\) and \(7/3\) on its orthogonal
complement.  This is an exact direct-Hessian calculation, not the
full physical Hessian.

\begin{corollary}[The missing score symbol is a positive Gram]
\label{cor:e8v41-score-Gram}
Let \(J_a(z,W):=D_WS(z,W)[a]\) and centre it by
\(\widetilde J_a=J_a-\mathbb E_WJ_a\).  Then
\[
 \operatorname{Cov}_W(J_a,J_b)
 =\langle\widetilde J_a,\widetilde J_b\rangle_{L^2(\mu_W)},      \tag{41.32}
\]
so the second term of (41.29) is positive semidefinite as a form in
the retained direction.  It can reduce any eigenvalue of the direct
symbol (41.30).
\end{corollary}

\begin{proof}
Expanding the \(L^2(\mu_W)\) inner product gives
\(\mathbb E_W[J_aJ_b]-\mathbb E_WJ_a\,\mathbb E_WJ_b\), which is
the covariance.  Taking \(a=b\) proves positivity.
\end{proof}

\begin{firewall}[Raw positivity cannot determine the beta coefficient]
\label{fire:e8v41-raw-symbol}
Neither the positivity of (41.30) nor the vertical floor (41.23)
computes the signed shell response.  The physical retained Hessian
is the difference in (41.29).  Any argument that reads
\(b_{\rm YM}\) from \(C(q)\) alone omits the connected score
covariance and is incomplete.
\end{firewall}

\begin{proposition}[Gauge covariance of the complete conditional Hessian]
\label{prop:e8v41-coarse-Ward}
Let \(g\) be a residual coarse gauge transformation and
\(a^g,b^g\) the transported retained directions.  Gauge invariance
of \(S\) and Haar invariance imply
\[
 D^2A(W^g)[a^g,b^g]=D^2A(W)[a,b].                               \tag{41.33}
\]
At a flat background the Fourier symbol of (41.29) therefore kills
coarse gauge gradients and is transverse.  This statement applies
to the sum of the two terms, not necessarily to either term after an
incomplete owner split.
\end{proposition}

\begin{proof}
The residual coarse transformation maps \(z\) by products of left
and right translations and maps \(W\) to \(W^g\).  Change variables
in (41.28); product Haar measure and \(S\) are invariant, hence
\(A(W^g)=A(W)\).  Differentiate twice along the transported
directions.  At the flat background, differentiating this identity
along an infinitesimal coarse gauge orbit gives the Ward null
direction.  Translation invariance diagonalizes it in momentum.
\end{proof}

\begin{assessment}[The actual coefficient acquisition]
\label{ass:e8v41-score-acquisition}
The coefficient problem has now been reduced without changing
regulators: one must compute the transverse \(q^2\) coefficient of
the direct symbol (41.30) minus the Fourier transform of the Gram
(41.32), including the one-loop vertical determinant (41.27) and
the fine-minus-coarse subtraction (41.18).  The Ward identity
(41.33) removes longitudinal ambiguity but does not determine the
remaining transverse scalar.
\end{assessment}

\subsection{The quotient matching card}

\begin{definition}[Quotient Matching Card (QMC)]
\label{def:e8v41-QMC}
A compact-link shell satisfies QMC on a background tube when the
following six rows hold uniformly in scale and volume.
\begin{enumerate}
\item \emph{Exact quotient row.}  The relative parity-tree map is a
      global Haar factorization, with \(J_T=1\) and the residual
      coarse gauge group retained.
\item \emph{Vertical branch row.}  A stationary centre exists on
      every good chart; its reduced Hessian has a uniform positive
      floor, and transition/bad charts have assigned owners.
\item \emph{Determinant bridge row.}  The fine tree form, its
      retained Schur form, and their common harmonic complement obey
      (41.18), including every zero-mode and measure-normalization
      factor.
\item \emph{Low-symbol row.}  The scale difference in (41.18), or an
      equivalent direct evaluation of (41.27), converges in the
      transverse quadratic channel to
      \[
       b_{\rm YM}=\frac{22C_A}{3(4\pi)^2}\log L,                 \tag{41.34}
      \]
      with the convention fixed in V37.
\item \emph{Compact nonlinear row.}  Haar-chart vertices,
      stationary-centre motion, the covariance in (41.29), aliases,
      and chart transitions are included and satisfy the local
      owner/majorant bounds.
\item \emph{Analytic remainder row.}  The complete quotient owners
      satisfy CII and CPBC, so V40 applies after subtraction of
      (41.34).  The word ghost in the older CPBC dictionary denotes
      only the comparison determinant paired by (41.18), not an
      additional field in (41.1).
\end{enumerate}
\end{definition}

\begin{theorem}[QMC implies the actual-regulator beta row]
\label{thm:e8v41-QMC-beta}
If QMC holds and the trajectory normalization of V37 is used, then
the exact parity-tree shell map, not merely its continuum comparison,
satisfies
\[
 u_{j+1}=u_j-b_{\rm YM}u_j^2
 +O(u_j^3+u_j^2\|e_j\|),                                       \tag{41.35}
\]
with volume- and scale-uniform constants on the QMC tube.
\end{theorem}

\begin{proof}
Rows 1--3 identify the Gaussian quotient response with the
fine-minus-coarse expression (41.18), without changing the defining
measure.  Row 4 supplies its transverse coefficient.  Row 5 puts
every remaining compact-link contribution into the complete owner
dictionary.  Row 6 invokes
Theorem~\ref{thm:e8v40-beta-remainder} for the subtracted response.
The inverse-coupling-to-\(u\) conversion is
Corollary~\ref{cor:e8v37-u-normal}, yielding (41.35).
\end{proof}

\begin{assessment}[QMC rows currently closed]
\label{ass:e8v41-QMC-status}
For the attached regulator, row 1 is exact.  At the flat pure-Wilson
point, the reduced part of row 2 has the explicit floor (41.23).
The finite-dimensional algebra in row 3 is now exact whenever the
background branch and matched retained complement exist.  Rows
4--6 are not yet established for the actual compact-link law.  In
particular, V37 computes the target in row 4 but not the required
convergence to it.
\end{assessment}

\begin{proposition}[No circular use of the desired coefficient]
\label{prop:e8v41-no-coefficient-circle}
Rows 1--3 of QMC are independent of the value in (41.34).  Row 4 can
therefore be tested by subtracting the two finite quotient
determinants in (41.18) before any continuum coefficient is inserted.
A failure of the resulting low-symbol limit returns a finite
momentum--colour witness rather than an assumption of asymptotic
freedom.
\end{proposition}

\begin{proof}
The first three rows use only Haar factorization, finite Hessian
algebra, a chosen complement, and block Gaussian elimination.  Their
statements contain no \(b_{\rm YM}\).  On a finite torus the two
determinants are finite matrices and their differentiated difference
is explicit.  Translation and residual colour symmetry decompose
that difference into finitely many momentum--centralizer blocks.
If the uniform limit in row 4 fails, the maximum block loss is
attained; selecting a maximizing subsequence gives the asserted
witness.
\end{proof}

\subsection{Upstream acquisition forced by V41}

The next calculation must remain inside (41.1).  At a small retained
background it must construct the vertical Hessian and score kernels
in the same tree coordinates, then evaluate either
\[
 \frac12D_W^2\log\det H_{rr}(W)                                  \tag{41.36}
\]
together with the non-Gaussian cumulant owners, or the exactly
equivalent fine-minus-coarse quotient expression (41.18).  The
Brillouin-edge scale-difference cancellation of V38 can be used only
after that exact symbol has been derived.

\begin{decision}[V42 target]
\label{dec:e8v41-next}
V42 should derive the background Taylor tensors of the parity-tree
vertical curl Gram matrix and the associated score covariance in
(41.29), with explicit colour and alias routing.  If that expansion
does not close locally, the programme should move one step upstream
to a finite-volume quotient-torsion identity for the dyadic block
map, rather than reintroduce an unsupported ghost field.
\end{decision}

\begin{firewall}[Claim boundary after V41]
\label{fire:e8v41-claim}
V41 repairs the regulator identity and proves the exact
tree-quotient/FP and fine-minus-coarse determinant bridges.  It does
not verify QMC rows 2--6 on a full background tube, compute the
compact-link score covariance, prove the one-loop coefficient for
the actual shell, construct the global RG trajectory, remove the
regulators, establish Osterwalder--Schrader reconstruction, or prove
the Yang--Mills mass gap.
\end{firewall}

\subsection{Parent record}

This V41 note was formed by appending the preceding material to V40,
whose validated record was
\[
\begin{array}{c|c}
\text{lines}&18304\\
\text{bytes}&637111\\
\text{SHA-256}&
\texttt{4B432B88D176E10430E14F200FF480AD815793F476EB3B4AAAF6E49B2DA342EF}.
\end{array}
\]

% =============================================================================
% END APPENDED EIGHTH-SERIES V41 MATERIAL
% =============================================================================

% =============================================================================
% BEGIN APPENDED EIGHTH-SERIES V42 MATERIAL
% =============================================================================

\section{Eighth-series V42: the compensated tree one-loop symbol and a
cofinal coefficient estimator}
\label{sec:e8v42}

\subsection{Nonduplication audit and the actual next frontier}

The reduced-fibre attachment already contains more of the V41 target
than the V41 summary used.  In particular it proves, for the pure
Wilson law:
\begin{enumerate}
\item the harmonic tree lift \(b[a]\) of a retained coarse direction;
\item an exact Haar translation which cancels the linear Gaussian
      score without changing the partition function;
\item explicit local \(45\times45\) vertical vertices \(R(b)\) and
      \(Q(b)\);
\item the exact one-loop trace
      \[
       \gamma_M[a]
       =\frac12\Tr(GQ(b[a]))
        +\frac13\Tr(GR(b[a])GR(b[a]));                           \tag{42.1}
      \]
\item an exponentially summable infinite-volume kernel
      \(\mathcal J_\infty\) and an exponentially small finite-torus
      winding error.
\end{enumerate}
Those results are imported and not reproved.  V42 instead performs
three missing tasks: it identifies (42.1) with the quotient
determinant of V41, corrects the incomplete Haar wording in V41, and
reduces the pure-Wilson matching coefficient to an absolutely
convergent real-space moment and a cofinal finite-card estimator.

\subsection{The complete one-loop density and the Haar correction}

Use product exponential coordinates \(z_e=\exp x_e\) near the flat
minimum.  For a retained logarithmic field \(v\), let \(x_*(v)\) be
the reduced stationary centre, let
\[
 A(v):=D_x^2V_M(x_*(v),v),
 \qquad dm_{\rm red}=J(x)\,dx,                                  \tag{42.2}
\]
and define
\[
 \Gamma_M(v):=\frac12\log\det A(v)-\log J(x_*(v)).              \tag{42.3}
\]

\begin{firewall}[Completion of the V41 determinant wording]
\label{fire:e8v42-Haar-completion}
The quantity \(\Gamma_{1,{\rm tree}}\) in (41.26) is only the
vertical determinant part of the Laplace coefficient.  In fixed
exponential coordinates the actual one-loop density is (42.3).
The extra Haar term is not optional.  QMC row 4 must therefore use
\(D^2\Gamma_M\), or use compensated left-Haar coordinates in which
the same term is absorbed into the translated Hessian vertices.
Equation (41.18), being a Gaussian quotient identity, remains
correct.
\end{firewall}

For a retained direction \(a\), let
\[
 A_0:=A(0),\qquad
 B:=D_xD_vV_M(0,0),\qquad
 L_a:=D x_*(0)[a]=-A_0^{-1}Ba.                                 \tag{42.4}
\]
For \({\rm SU}(3)\) in the Lie metric of the fibre attachment,
\[
 D\log J(0)=0,\qquad -D^2\log J(0)=\frac12I.                   \tag{42.5}
\]

\begin{proposition}[Coordinate-complete vertical determinant response]
\label{prop:e8v42-coordinate-complete}
If \(A_a\) and \(A_{ab}\) denote total derivatives of \(A(v)\),
including the motion of \(x_*(v)\), then
\[
 \boxed{
 D^2\Gamma_M(0)[a,b]
 =\frac12\Tr\!\left(A_0^{-1}A_{ab}
 -A_0^{-1}A_bA_0^{-1}A_a\right)
 +\frac12\langle L_a,L_b\rangle.}                              \tag{42.6}
\]
\end{proposition}

\begin{proof}
The determinant terms follow from
Lemma~\ref{lem:e8v41-logdet-response}.  For the Haar term, the chain
rule and (42.5) give
\[
 -D^2[\log J(x_*(v))]_{v=0}[a,b]
 =\frac12\langle L_a,L_b\rangle;
\]
the term containing \(D^2x_*(0)\) vanishes because
\(D\log J(0)=0\).  Adding the two pieces proves (42.6).
\end{proof}

\begin{remark}[Why the Haar translation has no displayed Jacobian]
\label{rem:e8v42-Haar-translation}
The compensated path multiplies every reduced link by a prescribed
group element.  It is a product of left Haar translations and hence
preserves \(dm_{\rm red}\) exactly.  In those moving coordinates the
last term of (42.6) is carried by the derivatives of the translated
vertical Hessian.  Fixed exponential coordinates and compensated
Haar coordinates give the same \(D^2\Gamma_M\).
\end{remark}

\subsection{Crosswalk from the compensated vertices to the quotient trace}

Use one colour component and the parity-cell incidence blocks
\[
 A=\frac13D_R^*D_R,\qquad
 B=\frac13D_R^*D_C,\qquad
 G=A^{-1},\qquad L=-GB.                                        \tag{42.7}
\]
The harmonic lift of a coarse field \(a\) is
\[
 b[a]|_R=La,\qquad b[a]|_C=a,\qquad b[a]|_{\mathscr T}=0,
 \qquad D_R^*D\,b[a]=0.                                        \tag{42.8}
\]
Fix a unit \(X\in\mathfrak{su}(3)\), take
\(W_{y,\mu}(t)=\exp(ta_{y,\mu}X)\), and left-translate the reduced
links by
\[
 (T_tz)_e=\exp(tb_eX)z_e.                                      \tag{42.9}
\]
Let \(M_b(t)\) be the vertical logarithmic-coordinate Hessian of the
unit Wilson action at the translated flat configuration.  The
explicit plaquette calculation in the attachment gives a real skew
matrix \(R(b)\), a real symmetric matrix \(Q(b)\), and
\[
 M_b(0)=A\otimes I_8,\qquad
 M_b'(0)=\frac13R(b)\otimes\operatorname{ad}X,\qquad
 \Tr_{\rm colour}M_b''(0)=Q(b),                                \tag{42.10}
\]
with
\[
 \Tr_{\rm colour}(\operatorname{ad}X)^2=-6.                    \tag{42.11}
\]
It also proves that the translated vertical gradient is \(O(t^3)\);
therefore the translated stationary centre has zero first and
second derivatives at \(t=0\).

\begin{theorem}[The attached tadpole--bubble formula is the V41
quotient determinant]
\label{thm:e8v42-trace-crosswalk}
For every fixed finite \(M\) and every real one-colour retained
direction \(aX\),
\[
 \boxed{
 D^2\Gamma_M(0)[aX,aX]
 =\frac12\Tr(GQ(b[a]))
  +\frac13\Tr(GR(b[a])GR(b[a])).}                              \tag{42.12}
\]
Moreover
\[
 \Tr(GRGR)
 =-\|G^{1/2}R\,G^{1/2}\|_{\rm HS}^2\le0,                       \tag{42.13}
\]
so (42.12) is the signed tadpole minus the positive bubble.  It is
the one-loop response of the literal parity-tree quotient and,
through (41.18), the fine-minus-coarse quotient determinant response.
\end{theorem}

\begin{proof}
Equation (42.9) preserves product Haar measure and the compensated
partition function equals the original partition function along
\(W(t)\).  Because the new stationary centre moves only at order
\(t^3\), two derivatives of the Laplace amplitude are
\[
 \frac12\Tr\left[
 (G\otimes I_8)M_b''(0)
 -(G\otimes I_8)M_b'(0)
  (G\otimes I_8)M_b'(0)\right].                                \tag{42.14}
\]
The first identity in (42.10) has the same \(G=A^{-1}\) as the
vertical quotient block in (41.15).  Taking the colour trace in the
first term of (42.14) gives \(\frac12\Tr(GQ)\).  The second term,
using (42.10)--(42.11), is
\[
 -\frac12\cdot\frac19\cdot(-6)\Tr(GRGR)
 =\frac13\Tr(GRGR).
\]
This proves (42.12).  Since \(R^*=-R\),
\[
 \Tr(GRGR)
 =\Tr\!\left[(G^{1/2}RG^{1/2})^2\right]
 =-\|G^{1/2}RG^{1/2}\|_{\rm HS}^2,
\]
which proves (42.13).  Finally, the Hessian \(A\) is precisely the
reduced block of the parity-tree quadratic form; Theorem
\ref{thm:e8v41-gaussian-fibre} and Corollary
\ref{cor:e8v41-scale-difference} identify its determinant response
with the fine-minus-coarse quotient response.  Haar compensation
ensures that (42.12), unlike (41.27) alone, includes the coordinate
density contribution.
\end{proof}

\begin{corollary}[The leading score covariance has not been omitted]
\label{cor:e8v42-score-included}
In the weak-coupling limit at fixed \(M\), the second term in
(42.12) is exactly minus the leading variance of the compensated
Wilson score.  Thus the quotient trace already contains the
direct-minus-score-covariance combination (41.29); no additional
score Gram may be subtracted from (42.12).
\end{corollary}

\begin{proof}
Under \(x=\beta^{-1/2}\xi\), the leading compensated score is the
centred quadratic Gaussian form
\(\frac12\xi^TM_b'(0)\xi\).  For covariance \(G\otimes I_8\), its
variance is
\(\frac12\Tr[(G\otimes I_8)M_b'(0)
(G\otimes I_8)M_b'(0)]\), which is the negative of the second term
in (42.14).  Equation (42.12) already subtracts it.
\end{proof}

\begin{firewall}[Pure Wilson versus the selected generated family]
\label{fire:e8v42-pure-generated}
Equation (42.12) is the pure-Wilson one-loop coefficient.  A fixed
analytic generated action contributes its own second derivative
along the Wilson minimizer.  A scale-dependent response-selected
\(\Psi_{\beta,M}\) additionally requires uniform control of all
mixed scores; it cannot be discarded merely because the pure-Wilson
trace is explicit.
\end{firewall}

\subsection{An absolutely convergent coefficient moment}

Let \(\mathcal J_\infty(x)\) be the infinite-coarse-lattice kernel
of (42.12), with one colour component retained.  The attachment proves
that for explicit constants \(\mu,K_\mu,E_\mu>0\),
\[
 \sum_{x\in\mathbb Z^4}\sum_{\alpha,\nu}
 e^{\mu|x|_1}|(\mathcal J_\infty(x))_{\alpha\nu}|
 \le K_\mu,                                                     \tag{42.15}
\]
and
\[
 \|\mathcal J_M-\operatorname{Per}_M\mathcal J_\infty\|_{\rm op}
 \le E_\mu e^{-\mu M}.                                         \tag{42.16}
\]
For momentum in direction \(1\) and transverse polarization \(2\),
put
\[
 j(x):=(\mathcal J_\infty(x))_{22},\qquad
 \gamma_\infty(q):=\sum_{x\in\mathbb Z^4}j(x)e^{iqx_1}.         \tag{42.17}
\]
Reflection and self-adjointness make \(j(x)\) real and even in
\(x_1\).  The exact coarse Ward identity gives
\[
 \sum_xj(x)=\gamma_\infty(0)=0.                                \tag{42.18}
\]

\begin{theorem}[Exact low-angle moment of the tree quotient]
\label{thm:e8v42-moment}
The series
\[
 \boxed{
 \kappa_{\rm tree}:=-\frac12\sum_{x\in\mathbb Z^4}x_1^2j(x)}    \tag{42.19}
\]
is absolutely convergent, and for every real \(q\),
\[
 \gamma_\infty(q)
 =\kappa_{\rm tree}q^2+r_\gamma(q),\qquad
 |r_\gamma(q)|\le\frac{K_\mu}{\mu^4}|q|^4.                     \tag{42.20}
\]
In addition,
\[
 |\kappa_{\rm tree}|\le\frac{K_\mu}{\mu^2}.                    \tag{42.21}
\]
\end{theorem}

\begin{proof}
Absolute convergence follows from (42.15) and
\[
 t^n\le n!\,\mu^{-n}e^{\mu t}\qquad(t\ge0).
\]
By evenness and (42.18),
\[
 \gamma_\infty(q)
 =\sum_xj(x)(\cos(qx_1)-1).
\]
Taylor's formula and
\(|\cos y-1+y^2/2|\le |y|^4/24\) give
\[
 \left|\gamma_\infty(q)
 +\frac{q^2}{2}\sum_xx_1^2j(x)\right|
 \le\frac{|q|^4}{24}\sum_x|x_1|^4|j(x)|
 \le\frac{K_\mu}{\mu^4}|q|^4.
\]
The \(n=2\) bound similarly gives
\(\frac12\sum_xx_1^2|j(x)|\le K_\mu/\mu^2\), proving
(42.21).
\end{proof}

The exact flat coarse Schur reference imported from the attachment is
\[
 s_0(q)
 =\frac{t(t+8)(t+24)}{9(t^2+18t+16)},\qquad
 t=\sin^2(q/2),                                                 \tag{42.22}
\]
and
\[
 s_0(q)=\frac13q^2-\frac{31}{288}q^4+O(q^6).                   \tag{42.23}
\]

\begin{corollary}[The actual pure-Wilson inverse-coupling shift]
\label{cor:e8v42-tree-shift}
The downward one-block, pure-Wilson quotient shift exists and is
\[
 \boxed{
 \delta\beta_{\rm tree}^{\downarrow}
 :=\lim_{q\to0}\frac{\gamma_\infty(q)}{s_0(q)}
 =3\kappa_{\rm tree}
 =-\frac32\sum_xx_1^2j(x).}                                    \tag{42.24}
\]
QMC row 4 for the Gaussian pure-Wilson shell is therefore equivalent
to the single scalar identity
\[
 \boxed{
 \delta\beta_{\rm tree}^{\downarrow}=-b_{\rm YM}
 =-\frac{22C_A}{3(4\pi)^2}\log2.}                              \tag{42.25}
\]
The reverse ultraviolet step has the opposite sign.
\end{corollary}

\begin{proof}
Divide (42.20) by (42.23) and let \(q\to0\), obtaining the first two
equalities in (42.24); (42.19) gives the last.  The continuum
orientation and normalization are (37.10)--(37.11) with \(L=2\).
Thus the downward lattice shell matches precisely when (42.25)
holds.
\end{proof}

\begin{firewall}[A bounded moment is not its evaluated value]
\label{fire:e8v42-moment-not-value}
Theorem~\ref{thm:e8v42-moment} proves existence, absolute
convergence, and a constructive bound for the lattice coefficient.
It does not prove the identity (42.25).  That equality is now an
explicit scalar theorem to be established, rather than an implicit
appeal to universality.
\end{firewall}

\subsection{Equivalent compact Brillouin integral}

Let \(G(k)=[D_R(k)^*D_R(k)/3]^{-1}\).  Let
\(\mathcal R_q(k)\) and \(\mathcal Q_q(k)\) be the exactly normalized
linear and quadratic compensated vertices of the attachment,
including all \(96\) plaquette types and the harmonic lift (42.8).
Thus \(\mathcal Q_q\) already contains the factor \(1/2\) from the
zero-frequency part of a normalized real mode; no second factor
\(1/2\) is placed before its Fourier trace.

\begin{proposition}[One integral contains the complete pure-Wilson
coefficient]
\label{prop:e8v42-BZ-integral}
For \(|q|\) sufficiently small,
\[
 \boxed{
 \begin{split}
 \gamma_\infty(q)
 =\int_{\mathbb T^4}\frac{d^4k}{(2\pi)^4}
 \bigg\{&
 \Tr[G(k)\mathcal Q_q(k)]\\
 &-\frac13\Tr[
 G(k+q\widehat e_1)\mathcal R_q(k)
 G(k)\mathcal R_q(k)^*]\bigg\}.
 \end{split}}                                                   \tag{42.26}
\]
Every matrix is \(45\times45\).  The integrand is smooth in
\((k,q)\), and
\[
 \boxed{
 \delta\beta_{\rm tree}^{\downarrow}
 =\frac32\int_{\mathbb T^4}\frac{d^4k}{(2\pi)^4}
 \left.\partial_q^2\right|_{q=0}
 \bigg\{
 \Tr[G(k)\mathcal Q_q(k)]
 -\frac13\Tr[
 G(k+q\widehat e_1)\mathcal R_q(k)
 G(k)\mathcal R_q(k)^*]\bigg\}.}                               \tag{42.27}
\]
\end{proposition}

\begin{proof}
Fourier transformation of the tadpole trace in (42.12) produces the
first term of (42.26).  The two Green kernels in the bubble carry
momenta \(k\) and \(k+q\widehat e_1\), producing the second.  The
finite-torus formula in the attachment is the corresponding Riemann
sum; absolute summability (42.15) permits passage to the integral.

The vertical spectral floor (41.22) makes \(G(k)\) analytic and
uniformly bounded on the compact Brillouin torus.  The incidence
symbols and local vertices are trigonometric polynomials, while
\(L(q)=-G(q)B(q)\) is analytic by the same floor.  Hence the
integrand and its first two \(q\)-derivatives are uniformly bounded.
Differentiate under the integral and use
\(\delta\beta_{\rm tree}^{\downarrow}
=\tfrac32\gamma_\infty''(0)\), which follows from (42.20) and
(42.24).
\end{proof}

\subsection{A certified cofinal finite-card estimator}

Set
\[
 q_M:=\frac{2\pi}{M},\qquad
 \delta\beta_M:=\frac{\gamma_M(q_M)}{s_0(q_M)}.                 \tag{42.28}
\]
For \(0<|q|\le1\), define the finite explicit reference constant
\[
 C_s:=\sup_{0<|q|\le1}
 \frac{|s_0(q)-q^2/3|}{q^4}<\infty,\qquad
 m_s:=\frac5{7\pi^2}.                                          \tag{42.29}
\]
The exact reference bounds give
\[
 s_0(q)\ge\frac5{28}\,4\sin^2(q/2)
 \ge m_sq^2\qquad(|q|\le\pi).                                  \tag{42.30}
\]

\begin{theorem}[Finite matrices converge to the quotient beta
coefficient with an explicit rate]
\label{thm:e8v42-cofinal-estimator}
For every \(M\ge7\),
\[
 \boxed{
 |\delta\beta_M-\delta\beta_{\rm tree}^{\downarrow}|
 \le C_1q_M^2+
 \frac{E_\mu}{m_sq_M^2}e^{-\mu M},}                            \tag{42.31}
\]
where
\[
 C_1:=\frac{K_\mu}{m_s}
 \left(\frac{3C_s}{\mu^2}+\frac1{\mu^4}\right).                \tag{42.32}
\]
In particular, the exact finite \(45M^4\)-variable quotient cards
satisfy
\[
 \boxed{\delta\beta_M\longrightarrow
 \delta\beta_{\rm tree}^{\downarrow}.}                         \tag{42.33}
\]
\end{theorem}

\begin{proof}
By (42.20),
\[
 |\gamma_\infty(q)-\kappa_{\rm tree}q^2|
 \le K_\mu\mu^{-4}q^4.
\]
Using (42.29)--(42.30), (42.21), and
\(\delta\beta_{\rm tree}^{\downarrow}=3\kappa_{\rm tree}\),
\[
 \begin{split}
 \left|\frac{\gamma_\infty(q)}{s_0(q)}
       -3\kappa_{\rm tree}\right|
 &\le
 \frac{|\kappa_{\rm tree}|\,|q^2-3s_0(q)|
       +K_\mu\mu^{-4}q^4}{s_0(q)}\\
 &\le C_1q^2.                                                   \tag{42.34}
 \end{split}
\]
The winding estimate (42.16), compressed to the chosen unit
polarization, gives
\[
 |\gamma_M(q_M)-\gamma_\infty(q_M)|
 \le E_\mu e^{-\mu M}.                                         \tag{42.35}
\]
Divide by (42.30), set \(q=q_M\), and add (42.34).  Since
\(q_M^2=O(M^{-2})\) and \(q_M^{-2}e^{-\mu M}=O(M^2e^{-\mu M})\),
both terms vanish.
\end{proof}

\begin{corollary}[Three equivalent pure-Gaussian matching tests]
\label{cor:e8v42-three-tests}
The following are equivalent:
\begin{enumerate}
\item the downward pure-Wilson Gaussian quotient satisfies QMC row 4;
\item the absolutely convergent moment obeys
      \[
       -\frac32\sum_xx_1^2j(x)=-b_{\rm YM};                    \tag{42.36}
      \]
\item the finite exact-card sequence obeys
      \[
       \lim_{M\to\infty}\delta\beta_M=-b_{\rm YM}.              \tag{42.37}
      \]
\end{enumerate}
\end{corollary}

\begin{proof}
The equivalence of the first two statements is
Corollary~\ref{cor:e8v42-tree-shift}.  The equivalence of the last
two follows from (42.33).
\end{proof}

\begin{assessment}[What the estimator changes]
\label{ass:e8v42-estimator-scope}
The matching problem no longer requires a simultaneous limit of a
large finite-volume covariance divided by a closing \(q^2\) symbol.
The exponentially local infinite kernel first supplies the moment
(42.19); alternatively, the finite quotient cards converge with the
explicit error (42.31).  The certified \(M=4,q=\pi/2\) value in the
attachment is a high-angle calibration and is not a test of
(42.37).
\end{assessment}

\subsection{The remaining coefficient theorem and the upstream route}

The scalar integral (42.27) is finite and regulator specific in its
integrand.  Its expected value is universal because the tree
quotient is exactly another coordinate representation of the
fine-minus-coarse gauge quotient.  That expectation is not yet a
proof of (42.25).  A noncircular proof can now proceed by an operator
homotopy:
\begin{enumerate}
\item interpolate the parity-tree slice to an orthogonal
      background gauge slice while retaining the exact orbit
      Jacobian of Theorem~\ref{thm:e8v41-quotient-FP};
\item differentiate the scale-difference determinant along the
      interpolation;
\item prove that its transverse \(q^2\) coefficient is a
      scale-independent local counterterm, hence cancels between the
      fine and coarse resolutions;
\item evaluate the orthogonal endpoint by the continuum shell
      calculation of V37, with the low-symbol domination of V38.
\end{enumerate}

\begin{decision}[V43 target]
\label{dec:e8v42-next}
V43 should construct this measured-slice homotopy in finite
dimensions and prove invariance of the logarithmic scale-difference
coefficient.  If the homotopy acquires an oblique zero-mode
Jacobian, it must be retained explicitly; the calculation should
move upstream to the determinant line rather than assume that the
Jacobian is constant.
\end{decision}

\begin{firewall}[Claim boundary after V42]
\label{fire:e8v42-claim}
V42 identifies the existing compensated tadpole--bubble formula
with the actual parity-tree quotient determinant, restores the Haar
term omitted from the shorthand in V41, proves the absolutely
convergent low-angle coefficient formula, and gives the cofinal
finite-card estimate (42.31).  It does not evaluate the moment as
\(-b_{\rm YM}\), verify the generated/nonlinear QMC rows, construct
the full RG trajectory, remove the regulators, establish
Osterwalder--Schrader reconstruction, or prove the Yang--Mills mass
gap.
\end{firewall}

\subsection{Parent record}

This V42 note was formed by appending the preceding material to V41,
whose validated record was
\[
\begin{array}{c|c}
\text{lines}&19025\\
\text{bytes}&665089\\
\text{SHA-256}&
\texttt{0CC159AE5BA160C64AC1B57F5EBC88B2DC7C045429AF99A9232D6987BE3C1E4A}.
\end{array}
\]

% =============================================================================
% END APPENDED EIGHTH-SERIES V42 MATERIAL
% =============================================================================

% =============================================================================
% BEGIN APPENDED EIGHTH-SERIES V43 MATERIAL
% =============================================================================

\section{Eighth-series V43: measured slice invariance and the
Gaussian-reference incidence correction}
\label{sec:e8v43}

\subsection{The V42 moment is a calibration until the reference is matched}

V42 correctly identifies the scalar
\[
 \delta\beta_{\rm W\to Sch}^{\downarrow}
 :=\lim_{q\to0}\frac{\gamma_{\rm W\to Sch,\infty}(q)}{s_0(q)}
 =-\frac32\sum_xx_1^2j_{\rm W\to Sch}(x).                       \tag{43.1}
\]
It is the one-loop shift of the declared pure-Wilson-to-Schur block.
The abstract RG constructed earlier, however, carries an
\(O(1)\) quadratic pullback residue in a moving stable reference
graph \(h_j(m)\).  Its Gaussian endpoint need not equal the bare
Wilson quadratic form after only one block.

\begin{firewall}[Correction of the V42 QMC equivalence]
\label{fire:e8v43-reference-correction}
The equivalences (42.25) and (42.36)--(42.37) are valid only after a
reference-incidence statement identifies the cofinal Gaussian
reference shape with the pure-Wilson-to-Schur shape used in
\(\gamma_{\rm W\to Sch}\).  Without that statement, (43.1) is a
rigorous calibration coefficient, not yet
\(-b_{\rm YM}\).  The moment and finite-volume convergence theorems
of V42 remain valid.
\end{firewall}

\subsection{Exact invariance under a change of gauge slice}

Let \(D:G\to E\) and \(K:E\to E\) obey the hypotheses of
Theorem~\ref{thm:e8v41-quotient-FP}.  Let
\(I_i:T_i\hookrightarrow E\), \(i=0,1\), be two linear complements
of \(\operatorname{ran}D\), with fixed coordinate measures.  Put
\[
 H_i:=I_i^*KI_i,\qquad
 J_i:=|\det[D,I_i]|.                                            \tag{43.2}
\]

\begin{lemma}[Transition between two complements]
\label{lem:e8v43-complement-transition}
There are unique maps \(R:T_1\to T_0\) and \(C:T_1\to G\) such that
\[
 I_1=I_0R+DC,                                                   \tag{43.3}
\]
and \(R\) is invertible.
\end{lemma}

\begin{proof}
The map \([D,I_0]:G\oplus T_0\to E\) is an isomorphism, so every
column of \(I_1\) has a unique decomposition of the displayed form.
If \(Rt=0\), then \(I_1t=DCt\).  Since
\(\operatorname{ran}I_1\cap\operatorname{ran}D=\{0\}\), both sides
vanish and \(t=0\).  The two slice spaces have the same dimension,
so \(R\) is invertible.
\end{proof}

\begin{theorem}[Measured-slice determinant invariance]
\label{thm:e8v43-slice-invariance}
With (43.2)--(43.3),
\[
 H_1=R^*H_0R,\qquad
 J_1=J_0|\det R|,                                               \tag{43.4}
\]
and hence
\[
 \boxed{
 \frac12\log\det H_1-\log J_1
 =\frac12\log\det H_0-\log J_0.}                               \tag{43.5}
\]
The same identity holds on a matched measured complement of retained
physical zero modes.
\end{theorem}

\begin{proof}
Since \(KD=0\) and \(K\) is self-adjoint, substitution of (43.3)
gives
\[
 I_1^*KI_1=R^*I_0^*KI_0R.
\]
Moreover
\[
 [D,I_1](\phi,t)
 =[D,I_0](\phi+Ct,Rt).
\]
The second coordinate transformation is block triangular with
determinant \(\det R\), proving (43.4).  Taking determinants in the
first equality gives
\(\det H_1=(\det R)^2\det H_0\); substitution of the second equality
cancels \(\log|\det R|\) and proves (43.5).  If zero modes are
retained first, the identical argument applies to the transported
complement and its pushforward measure.
\end{proof}

\begin{corollary}[Tree gauge equals the orthogonal background-gauge
determinant]
\label{cor:e8v43-tree-orthogonal}
Let \(T_0\) be the parity-tree slice and let
\(T_1=(\operatorname{ran}D)^\perp\) with its Euclidean measure.
Then
\[
 J_0=1,\qquad J_1=(\det L_0)^{1/2},
\]
and
\[
 \boxed{
 \frac12\log\det H_{\rm tree}
 =\frac12\log\det K_P-\frac12\log\det L_0
 =\frac12\log\det K_{\rm gf}-\log\det L_0.}                     \tag{43.6}
\]
\end{corollary}

\begin{proof}
Tree unimodularity is
Lemma~\ref{lem:e8v41-tree-unimodular}.  For the orthogonal slice,
\(D=RL_0^{1/2}\) with \(R\) an isometry, so the orbit-coordinate
Jacobian is \((\det L_0)^{1/2}\).  Its restricted Hessian is \(K_P\).
Apply (43.5), then use
\(\det K_{\rm gf}=\det K_P\det L_0\) from (41.13).
\end{proof}

\begin{corollary}[All background jets are slice independent]
\label{cor:e8v43-background-jets}
Suppose \(D(W),K(W),I_i(W)\) are analytic on a background tube,
remain complementary, and have a common constant-dimensional
retained kernel.  Then every background derivative of
\[
 \frac12\log\det{}'H_i(W)-\log J_i(W)                           \tag{43.7}
\]
is independent of \(i\).  In particular the transverse \(F^2\)
coefficient and its scale difference cannot depend on whether the
parity-tree or orthogonal background slice is used.
\end{corollary}

\begin{proof}
Equation (43.5) holds pointwise as an identity of analytic functions
on the tube.  Differentiate it to any finite order.
\end{proof}

\begin{assessment}[The gauge-slice part of QMC is closed]
\label{ass:e8v43-slice-closed}
There is no remaining homotopy anomaly at the Gaussian level.
Whenever the slice stays complementary, its determinant change is
exactly its measure Jacobian and cancels in (43.5).  Thus the
tree-versus-background-gauge issue does not block the coefficient.
The remaining mismatch is between quadratic \emph{reference shapes},
not between gauge slices of one shape.
\end{assessment}

\subsection{The coefficient is a functional of the Gaussian reference shape}

Let \(\mathfrak Q\) be a class of translation-invariant, gauge
quotiented quadratic reference forms with a common retained
normalization.  For \(Q\in\mathfrak Q\), assume its transverse
classical symbol and its one-loop quotient kernel obey
\[
 s_Q(q)=\sigma(Q)q^2+O(q^4),\qquad \sigma(Q)>0,                 \tag{43.8}
\]
\[
 \|{\cal J}_Q\|_\mu
 :=\sum_xe^{\mu|x|_1}\sum_{\alpha,\nu}
 |({\cal J}_Q(x))_{\alpha\nu}|\le K.                           \tag{43.9}
\]
Put
\[
 \kappa(Q):=-\frac12\sum_xx_1^2({\cal J}_Q(x))_{22},
 \qquad
 \mathfrak b^\downarrow(Q):=\frac{\kappa(Q)}{\sigma(Q)}.        \tag{43.10}
\]
For the Wilson-to-Schur reference,
\(\sigma(Q_{\rm W\to Sch})=1/3\), so
\[
 \mathfrak b^\downarrow(Q_{\rm W\to Sch})
 =\delta\beta_{\rm W\to Sch}^{\downarrow}.                     \tag{43.11}
\]

\begin{lemma}[Lipschitz continuity of the one-loop moment]
\label{lem:e8v43-moment-Lipschitz}
Suppose, for \(Q,Q'\in\mathfrak Q\),
\[
 \|{\cal J}_Q-{\cal J}_{Q'}\|_\mu
 \le L_J\|Q-Q'\|_{\mathfrak Q},\qquad
 |\sigma(Q)-\sigma(Q')|
 \le L_\sigma\|Q-Q'\|_{\mathfrak Q},                           \tag{43.12}
\]
and \(\sigma(Q),\sigma(Q')\ge\sigma_0>0\).  Then
\[
 \boxed{
 |\mathfrak b^\downarrow(Q)-\mathfrak b^\downarrow(Q')|
 \le L_b\|Q-Q'\|_{\mathfrak Q},}                               \tag{43.13}
\]
where
\[
 L_b:=\frac{L_J}{\mu^2\sigma_0}
      +\frac{KL_\sigma}{\mu^2\sigma_0^2}.                       \tag{43.14}
\]
\end{lemma}

\begin{proof}
The moment estimate used in V42 gives
\[
 |\kappa(Q)-\kappa(Q')|
 \le\frac{L_J}{\mu^2}\|Q-Q'\|_{\mathfrak Q},
 \qquad |\kappa(Q')|\le\frac K{\mu^2}.                          \tag{43.15}
\]
Now write
\[
 \frac{\kappa(Q)}{\sigma(Q)}
 -\frac{\kappa(Q')}{\sigma(Q')}
 =
 \frac{\kappa(Q)-\kappa(Q')}{\sigma(Q)}
 +\kappa(Q')
 \frac{\sigma(Q')-\sigma(Q)}{\sigma(Q)\sigma(Q')}.
\]
Use (43.12), (43.15), and the lower floor \(\sigma_0\).
\end{proof}

\begin{theorem}[Cofinal Gaussian-shape coefficient]
\label{thm:e8v43-shape-limit}
Let the rescaled Gaussian reference map have an invariant shape
\(Q_*\in\mathfrak Q\), and suppose a reference sequence satisfies
\[
 \|Q_n-Q_*\|_{\mathfrak Q}\le C_Q\rho^n,\qquad 0<\rho<1.        \tag{43.16}
\]
Under the hypotheses of
Lemma~\ref{lem:e8v43-moment-Lipschitz},
\[
 \boxed{
 |\mathfrak b^\downarrow(Q_n)-\mathfrak b^\downarrow(Q_*)|
 \le L_bC_Q\rho^n.}                                            \tag{43.17}
\]
Consequently
\[
 \left|
 \frac1N\sum_{n=0}^{N-1}\mathfrak b^\downarrow(Q_n)
 -\mathfrak b^\downarrow(Q_*)\right|
 \le\frac{L_bC_Q}{N(1-\rho)}.                                  \tag{43.18}
\]
\end{theorem}

\begin{proof}
Equation (43.17) is (43.13) applied to (43.16).  Sum the geometric
bound for \(0\le n<N\) and divide by \(N\) to obtain (43.18).
\end{proof}

\begin{corollary}[Only the endpoint of the moving stable graph enters
the leading beta coefficient]
\label{cor:e8v43-moving-leading}
Suppose the moving reference shape satisfies
\[
 Q(h(u))=Q_*+O(u)                                               \tag{43.19}
\]
in \(\mathfrak Q\), and suppose the exact downward inverse-coupling
increment has the uniform expansion
\[
 \Delta(u,h)=\mathfrak b^\downarrow(Q(h))+O(u).                 \tag{43.20}
\]
Then
\[
 \boxed{\Delta(u,h(u))
 =\mathfrak b^\downarrow(Q_*)+O(u).}                            \tag{43.21}
\]
The desired ultraviolet coefficient is therefore
\(-\mathfrak b^\downarrow(Q_*)\), not automatically
\(-\mathfrak b^\downarrow(Q_{\rm W\to Sch})\).
\end{corollary}

\begin{proof}
Apply (43.13) to (43.19) and insert the result in (43.20).
\end{proof}

\begin{firewall}[Stable contraction does not identify the endpoint]
\label{fire:e8v43-contraction-not-incidence}
The abstract contraction of the moving stable graph proves uniqueness
and tracking after its map and norm are supplied.  It does not prove
that \(Q_*=Q_{\rm W\to Sch}\), nor does it evaluate
\(\mathfrak b^\downarrow(Q_*)\).  Those are finite Gaussian
reference-incidence data.
\end{firewall}

\subsection{Gaussian Reference Incidence Card}

\begin{definition}[Gaussian Reference Incidence Card (GRIC)]
\label{def:e8v43-GRIC}
GRIC consists of the following six regulator-specific rows.
\begin{enumerate}
\item \emph{Quadratic map.}  The exact Gaussian block, pullback
      rechart, canonical dilation, and marginal normalization define
      one map \({\cal F}_{\rm G}\) on \(\mathfrak Q\).
\item \emph{Invariant shape.}  There is a declared
      \(Q_*={\cal F}_{\rm G}(Q_*)\), and the physical reference
      sequence converges to it with a proved rate such as (43.16).
\item \emph{Quotient locality.}  On a neighborhood of \(Q_*\), the
      parity-tree vertical block has a uniform gap and its inverse,
      harmonic router, and one-loop kernel have common exponential
      weighted bounds.
\item \emph{Marginal normalization.}  The retained transverse symbol
      has (43.8) with \(\sigma\ge\sigma_0>0\), and the coefficient
      \(\kappa/\sigma\) is the inverse-coupling increment in the
      convention of (37.10).
\item \emph{Shape continuity.}  The kernel and normalization maps
      satisfy (43.12), including the Haar-density and quotient
      Jacobian terms exactly once.
\item \emph{Endpoint evaluation.}  Either the moment
      \[
       \mathfrak b^\downarrow(Q_*)
       =-\frac{1}{2\sigma(Q_*)}
         \sum_xx_1^2({\cal J}_{Q_*}(x))_{22}                    \tag{43.22}
      \]
      is evaluated directly, or a lattice-to-continuum symbol theorem
      proves
      \[
       \mathfrak b^\downarrow(Q_*)=-b_{\rm YM}.                 \tag{43.23}
      \]
\end{enumerate}
\end{definition}

\begin{theorem}[GRIC supplies the noncircular marginal coefficient]
\label{thm:e8v43-GRIC}
Assume GRIC and QMC rows 5--6.  Then the cofinal inverse-coupling
recursion has
\[
 \left(\frac1{u_{j+1}}\right)
 =\left(\frac1{u_j}\right)+b_{\rm YM}
 +O(u_j+\|e_j\|),                                               \tag{43.24}
\]
and hence
\[
 u_{j+1}=u_j-b_{\rm YM}u_j^2
 +O(u_j^3+u_j^2\|e_j\|).                                       \tag{43.25}
\]
\end{theorem}

\begin{proof}
GRIC rows 1--5 and
Corollary~\ref{cor:e8v43-moving-leading} identify the leading
downward increment with
\(\mathfrak b^\downarrow(Q_*)\), independently of gauge slice by
Corollary~\ref{cor:e8v43-background-jets}.  Row 6 gives
\(-b_{\rm YM}\) in the downward orientation, hence \(+b_{\rm YM}\)
in the ultraviolet orientation.  QMC rows 5--6 supply the
uniform \(O(u+\|e\|)\) remainder.  Apply
Corollary~\ref{cor:e8v37-u-normal}.
\end{proof}

\begin{assessment}[Status after the reference-incidence correction]
\label{ass:e8v43-status}
The following facts are now separated.
\begin{enumerate}
\item Gauge-slice matching is exact: (43.5) closes it to every
      background derivative on a constant-rank branch.
\item The pure-Wilson-to-Schur coefficient has the exact moment and
      finite-card estimators of V42.
\item The abstract moving reference graph gives a place for the
      \(O(1)\) quadratic pullback residue, but the supplied active
      notes do not construct its Gaussian endpoint \(Q_*\) as an
      explicit symbol or identify it with the V42 Wilson calibration.
\item Therefore GRIC rows 1--2 and 6, not the gauge quotient, are the
      immediate coefficient bottleneck.
\end{enumerate}
\end{assessment}

\subsection{Upstream construction selected for V44}

On a translation-invariant Gaussian quotient, the coarse compliance
is the blocked fine compliance.  Thus, before inversion and
normalization, the quadratic map has the form
\[
 C_{n+1}(q)
 =\sum_{\sigma\in\{0,1\}^4}
 B_\sigma(q)\,
 K_n(p_\sigma)^\dagger\,
 B_\sigma(q)^*,\qquad
 p_\sigma=\frac q2+\pi\sigma,                                  \tag{43.26}
\]
on the transverse quotient, followed by the declared dilation and
coupling normalization.  The exact Wilson reference formula
(42.22) is the first iterate of this mechanism, not yet its fixed
point.

\begin{decision}[V44 target]
\label{dec:e8v43-next}
V44 should define the complete normalized map following (43.26),
including the transverse projector, toron retention, dilation, and
the marginal scalar \(\sigma\).  It should then construct its
invariant Gaussian compliance, preferably by an absolutely
convergent all-alias formula, and compare the first Wilson iterate
with that endpoint.  If no strict shape contraction is available,
the all-alias fixed-point formula should replace an assumed
contraction.
\end{decision}

\begin{firewall}[Claim boundary after V43]
\label{fire:e8v43-claim}
V43 proves exact measured gauge-slice invariance and corrects the
reference-shape incidence needed to interpret V42.  It does not
construct \(Q_*\), evaluate (43.22), verify the nonlinear/complex
QMC rows, construct the full RG trajectory, remove the regulators,
establish Osterwalder--Schrader reconstruction, or prove the
Yang--Mills mass gap.
\end{firewall}

\subsection{Parent record}

This V43 note was formed by appending the preceding material to V42,
whose validated record was
\[
\begin{array}{c|c}
\text{lines}&19593\\
\text{bytes}&684932\\
\text{SHA-256}&
\texttt{417A6053DE6E9A12997ADCD84E61922CF56C419FD2E756BB03A0228C95DC43D3}.
\end{array}
\]

% =============================================================================
% END APPENDED EIGHTH-SERIES V43 MATERIAL
% =============================================================================

% =============================================================================
% BEGIN APPENDED EIGHTH-SERIES V44 MATERIAL
% =============================================================================

\section{Eighth-series V44: straight-link perimeter divergence and the
Gaussian reference germ}
\label{sec:e8v44}

\subsection{Why the naive all-alias fixed point must be tested first}

The fixed-shape GRIC of V43 is a sufficient card, not an established
property of the declared Wilson block.  The block retains a straight
two-link path in the link direction but performs no transverse
averaging.  Repeating it exposes more transverse ultraviolet aliases.
The resulting link-potential compliance need not converge as a
bounded function even though its Ward-normalized inverse germ does.
This section computes that effect exactly in the simplest axial
transverse sector.

\subsection{The repeated axial transverse compliance}

Keep component \(2\) and restrict fine momenta to \(p_2=0\).  Write
\(k=(p_1,p_3,p_4)\in\mathbb T^3\).  The unit-Wilson transverse
compliance is
\[
 c_0(k)
 :=\frac{3}{\omega_3(k)},\qquad
 \omega_3(k):=
 4\sum_{\nu\in\{1,3,4\}}\sin^2(k_\nu/2).                       \tag{44.1}
\]
It is defined away from \(k=0\).  The dyadic path block has Fourier
normalization \(1/4\), and its component-\(2\) path factor is two
when the second alias bit is zero and is zero when that bit is one.
Consequently the scalar compliance map on this reflection-invariant
sector is
\[
 ({\cal T}c)(k)
 :=\frac14\sum_{\sigma\in\{0,1\}^3}
 c\left(\frac k2+\pi\sigma\right).                             \tag{44.2}
\]
All arguments are understood modulo \(2\pi\).

\begin{lemma}[Exact \(n\)-fold alias formula]
\label{lem:e8v44-nfold}
For \(n\ge1\),
\[
 \boxed{
 c_n(k):=({\cal T}^nc_0)(k)
 =4^{-n}\!
 \sum_{m\in\{0,\ldots,2^n-1\}^3}
 c_0\left(\frac{k+2\pi m}{2^n}\right).}                         \tag{44.3}
\]
\end{lemma}

\begin{proof}
For \(n=1\), identify \(m=\sigma\) in (44.2).  Suppose (44.3)
holds at \(n\).  Applying (44.2) replaces \(k\) by
\(k/2+\pi\sigma\), and then
\[
 \frac{k/2+\pi\sigma+2\pi m}{2^n}
 =\frac{k+2\pi(\sigma+2m)}{2^{n+1}}.
\]
As \(\sigma\in\{0,1\}^3\) and
\(m\in\{0,\ldots,2^n-1\}^3\), the vector
\(\sigma+2m\) ranges bijectively over
\(\{0,\ldots,2^{n+1}-1\}^3\).  The prefactor gains one factor
\(1/4\), proving the induction.
\end{proof}

\begin{corollary}[The attached Schur reference is the first iterate]
\label{cor:e8v44-first-iterate}
For \(k=q\widehat e_1\ne0\),
\[
 c_1(q\widehat e_1)=s_0(q)^{-1},                               \tag{44.4}
\]
where \(s_0\) is (42.22).
\end{corollary}

\begin{proof}
Equation (44.2) is exactly the component-\(2\) compression of the
sixteen-alias formula
\[
 \frac3{16}\sum_{\sigma\in\{0,1\}^4}
 B_\sigma\frac{I-dd^*/\omega}{\omega}B_\sigma^*.
\]
The eight terms with \(\sigma_2=1\) vanish and the remaining
component-\(2\) terms give (44.2).  The explicit inversion carried
out in the attachment is (42.22).
\end{proof}

\subsection{The unsmeared compliance has no bounded fixed shape}

\begin{lemma}[Three-dimensional grid-shell estimate]
\label{lem:e8v44-grid-shell}
For every fixed \(k\in\mathbb T^3\setminus\{0\}\), there is a finite
\(C(k)\) such that
\[
 \sum_{m\in\{0,\ldots,2^n-1\}^3}
 c_0\left(\frac{k+2\pi m}{2^n}\right)
 \le C(k)\,2^{3n}.                                              \tag{44.5}
\]
\end{lemma}

\begin{proof}
Choose centered representatives in \([-\pi,\pi]^3\).  On that cube,
\[
 \omega_3(y)\ge\frac4{\pi^2}|y|^2,
 \qquad
 c_0(y)\le\frac{3\pi^2}{4|y|^2}.                               \tag{44.6}
\]
The grid spacing is \(h=2\pi2^{-n}\).  Apart from the point nearest
zero, the number of grid points whose centered distance lies between
\(rh\) and \((r+1)h\) is at most \(C_0(r+1)^2\).  Their contribution
to the inverse-square sum is at most
\[
 C_1h^{-2}\sum_{r=1}^{O(h^{-1})}1=O(h^{-3})=O(2^{3n}).
\]
The nearest point contributes at most \(C_2(k)4^n\), which is also
bounded by \(C_2(k)2^{3n}\) for \(n\ge1\).  Insert (44.6).
\end{proof}

\begin{theorem}[Straight-link contact/perimeter divergence]
\label{thm:e8v44-perimeter-divergence}
For every fixed nonzero \(k\in\mathbb T^3\),
\[
 \boxed{
 2^{n-2}\le c_n(k)\le C(k)2^n.}                                \tag{44.7}
\]
In particular \({\cal T}^nc_0\) has no finite pointwise limit and no
bounded compliance fixed shape in any norm controlling one nonzero
momentum value.
\end{theorem}

\begin{proof}
Since \(\omega_3(y)\le12\), every term in (44.3) obeys
\(c_0(y)\ge1/4\).  There are \(2^{3n}\) terms, hence
\[
 c_n(k)\ge4^{-n}2^{3n}\frac14=2^{n-2}.
\]
The upper bound is (44.5) multiplied by \(4^{-n}\).  Divergence at
one fixed \(k\ne0\) excludes convergence in any norm dominating
point evaluation there.
\end{proof}

\begin{assessment}[Meaning of the divergence]
\label{ass:e8v44-perimeter-meaning}
The growth in (44.7) is the Gaussian form of the ultraviolet contact
or perimeter renormalization of an unsmeared straight link.  It is
not a gauge-boson mass: it occurs in the link-potential compliance,
while the inverse curvature symbol retains a Ward zero.  It also
does not contradict the one-step vertical gap, because that gap
controls the variables eliminated at a fixed scale rather than the
all-refinement limit of a retained line observable.
\end{assessment}

\subsection{The marginal germ survives exactly}

For an axial argument set
\[
 s_n(q):=c_n(q\widehat e_1)^{-1}.                               \tag{44.8}
\]
The all-zero alias in (44.3) is
\[
 z_n(q)
 :=4^{-n}\frac{3}{4\sin^2(q/2^{n+1})}.                         \tag{44.9}
\]

\begin{theorem}[Exact Ward germ and divergent quartic reference
coefficient]
\label{thm:e8v44-Ward-germ}
For every fixed \(n\),
\[
 \boxed{\lim_{q\to0}q^2c_n(q\widehat e_1)=3,\qquad
 \lim_{q\to0}\frac{s_n(q)}{q^2}=\frac13.}                       \tag{44.10}
\]
More precisely, there is a finite \(p_n\ge0\) such that
\[
 c_n(q\widehat e_1)=\frac3{q^2}+p_n+O_n(q^2),                  \tag{44.11}
\]
\[
 s_n(q)=\frac{q^2}{3}-\frac{p_n}{9}q^4+O_n(q^6),               \tag{44.12}
\]
and constants \(c,C>0\), independent of \(n\), such that
\[
 \boxed{c\,2^n\le1+p_n\le C\,2^n.}                             \tag{44.13}
\]
\end{theorem}

\begin{proof}
Taylor expansion of (44.9) gives
\[
 z_n(q)=\frac3{q^2}+\frac{4^{-n}}4+O(4^{-2n}q^2).               \tag{44.14}
\]
Every other summand in (44.3) is analytic and even in \(q\) near
zero; the reflected alias pairs cancel its linear term.  Their
constant terms, together with the second term of (44.14), define
\(p_n\ge0\).  This proves (44.11), and inversion gives
(44.10) and (44.12).

For the lower bound, each of the \(2^{3n}-1\) nonzero aliases has
value at least \(1/4\), so its total constant contribution is at
least
\[
 4^{-n}\frac{2^{3n}-1}{4}
 =2^{n-2}-4^{-n-1}.                                             \tag{44.15}
\]
For the upper bound, repeat the grid-shell estimate of
Lemma~\ref{lem:e8v44-grid-shell} at \(k=0\) after deleting the
singular all-zero point.  The resulting sum is \(O(2^{3n})\);
multiplication by \(4^{-n}\) gives \(O(2^n)\).
\end{proof}

\begin{corollary}[The first Schur coefficient is not a stationary
bounded shape]
\label{cor:e8v44-no-first-fixed}
Although \(s_1=s_0\) has the correct marginal coefficient \(1/3\),
its further exact Gaussian iterates develop a quartic coefficient of
order \(2^n\).  Therefore
\[
 Q_{\rm W\to Sch}=Q_*
\]
cannot be justified by an unrenormalized bounded-compliance fixed
point.
\end{corollary}

\begin{proof}
Use (44.12)--(44.13).  A bounded analytic fixed shape would have a
uniformly bounded quartic Taylor coefficient, whereas the exact
iterates do not.
\end{proof}

\begin{firewall}[Order of limits]
\label{fire:e8v44-order-limits}
For every fixed \(n\), taking \(q\to0\) first gives the exact
marginal germ (44.10).  For every fixed \(q\ne0\), taking
\(n\to\infty\) first gives \(s_n(q)\to0\) by (44.7).  These limits
do not commute.  A cofinal beta argument must state its physical
momentum scaling and reference subtraction before interchanging
them.
\end{firewall}

\subsection{The physically scaled diagonal is well behaved}

\begin{lemma}[Uniform contact bound near the Ward point]
\label{lem:e8v44-contact-bound}
There is a numerical \(C_{\rm ct}<\infty\) such that, for
\(0<|q|\le1\) and every \(n\ge1\),
\[
 0\le
 c_n(q\widehat e_1)-\frac3{q^2}
 \le C_{\rm ct}2^n.                                             \tag{44.16}
\]
\end{lemma}

\begin{proof}
The all-zero term (44.9) is at least \(3/q^2\), because
\(|\sin x|\le|x|\), and its excess is uniformly \(O(4^{-n})\)
for \(|q|\le1\).  Delete that term from (44.3).  The centered grid
now stays at least a fixed fraction of one grid spacing away from
zero.  The shell count in
Lemma~\ref{lem:e8v44-grid-shell}, uniformly for the shift
\(|q|\le1\), bounds its unscaled sum by \(C2^{3n}\).
Multiplication by \(4^{-n}\) proves the upper bound.
\end{proof}

\begin{theorem}[Diagonal Ward normalization]
\label{thm:e8v44-diagonal-Ward}
Let \(q_n\ne0\) satisfy
\[
 q_n\longrightarrow0,\qquad 2^nq_n^2\longrightarrow0.          \tag{44.17}
\]
Then
\[
 \boxed{
 \frac{s_n(q_n)}{q_n^2/3}\longrightarrow1.}                    \tag{44.18}
\]
In particular (44.18) holds for fixed physical momentum under
dyadic refinement, \(q_n=2^{-n}p\).
\end{theorem}

\begin{proof}
Write
\[
 c_n(q\widehat e_1)=\frac3{q^2}+R_n(q),
 \qquad 0\le R_n(q)\le C_{\rm ct}2^n.
\]
Then
\[
 \frac{s_n(q)}{q^2/3}
 =\frac1{1+q^2R_n(q)/3},
\]
and hence
\[
 0\le1-\frac{s_n(q)}{q^2/3}
 \le\frac{C_{\rm ct}}3\,2^nq^2.                                \tag{44.19}
\]
Use (44.17).  For \(q_n=2^{-n}p\), the right side is
\((C_{\rm ct}|p|^2/3)2^{-n}\).
\end{proof}

\begin{assessment}[The appropriate reference object is a germ
sequence]
\label{ass:e8v44-germ-sequence}
The exact block has two simultaneous properties:
\[
 c_n(q)\asymp2^n\quad(q\ne0\ \hbox{fixed}),\qquad
 s_n(q_n)\sim q_n^2/3\quad(q_n=2^{-n}p).                        \tag{44.20}
\]
Thus a bounded global compliance \(Q_*\) is the wrong mandatory
object for this unsmeared link coordinate.  The cofinal programme
needs a scale-dependent reference sequence with a controlled Ward
germ and an explicitly owned contact/perimeter part.
\end{assessment}

\begin{definition}[Gaussian Germ Incidence Card (GGIC)]
\label{def:e8v44-GGIC}
GGIC replaces the fixed-shape portion of GRIC by the following
scale-dependent rows.
\begin{enumerate}
\item The exact Gaussian block/rechart produces reference symbols
      \(s_j(q)\) and contact owners \(P_j\).
\item On a declared physical diagonal \(q_j\), the condition
      \(2^jq_j^2\to0\) and a bound of the form (44.19) hold.
\item The marginal coefficient
      \(\lim_{q\to0}s_j(q)/q^2\) is normalized uniformly away from
      zero, while the contact owner is excluded from the marginal
      projection.
\item The measured quotient determinant is slice independent as in
      (43.5), scale by scale.
\item After the contact owner is removed, the one-loop kernels have
      uniform second moments and their normalized coefficients form
      a Cauchy sequence.
\item The limiting downward coefficient is evaluated as
      \(-b_{\rm YM}\), and the nonlinear QMC remainder is uniform on
      the same diagonal.
\end{enumerate}
\end{definition}

\begin{theorem}[GGIC is sufficient for the cofinal beta row]
\label{thm:e8v44-GGIC}
If GGIC and QMC rows 5--6 hold, then the ultraviolet recurrence
(43.24)--(43.25) follows.
\end{theorem}

\begin{proof}
Rows 1--3 define one unambiguous cofinal marginal projection despite
the contact divergence.  Row 4 permits evaluation in the orthogonal
background slice without changing the tree quotient.  Rows 5--6
give the leading coefficient and the
\(O(u_j+\|e_j\|)\) inverse-coupling remainder.  The algebraic
conversion to \(u_j\) is
Corollary~\ref{cor:e8v37-u-normal}.
\end{proof}

\begin{firewall}[No subtraction without an owner]
\label{fire:e8v44-no-free-subtraction}
The divergent quantity in (44.7) cannot simply be deleted from the
measure.  GGIC requires a declared contact/perimeter owner and proves
that the marginal projection is insensitive to it on the physical
diagonal.  This is a reference rechart, not a change of the exact
Wilson fibre integral.
\end{firewall}

\subsection{Status and the compulsory V45 audit}

\begin{assessment}[Progress after V44]
\label{ass:e8v44-status}
V44 proves that the naive bounded Gaussian endpoint requested in V43
is generally unavailable for the unsmeared straight-link block.
It simultaneously proves that the required \(q^2/3\) Ward germ
survives on the physical dyadic diagonal.  The immediate open row is
GGIC item 5: construct the contact-subtracted one-loop kernel sequence
and prove its second-moment coefficient is Cauchy.
\end{assessment}

\begin{decision}[V45 audit and next calculation]
\label{dec:e8v44-next}
V45 is the mandatory five-version companion audit.  It will inspect
the active SM and Navier--Stokes notes for a scale-dependent
reference/contact subtraction or a Cauchy moment mechanism.  After
that audit it should formulate the one-loop contact-difference
telescope for the Yang--Mills reference sequence, importing no
companion constant.
\end{decision}

\begin{firewall}[Claim boundary after V44]
\label{fire:e8v44-claim}
V44 proves the repeated-block contact divergence, the exact marginal
germ, and the diagonal normalization theorem.  It does not construct
the contact-subtracted one-loop Cauchy sequence, evaluate its limit,
verify the nonlinear/complex QMC rows, construct the full RG
trajectory, remove the regulators, establish
Osterwalder--Schrader reconstruction, or prove the Yang--Mills mass
gap.
\end{firewall}

\subsection{Parent record}

This V44 note was formed by appending the preceding material to V43,
whose validated record was
\[
\begin{array}{c|c}
\text{lines}&20017\\
\text{bytes}&699797\\
\text{SHA-256}&
\texttt{0430FCA45B706C2DB0B0A871024B19B9FE7DFFF8B59F8471548C06EB80D66CE1}.
\end{array}
\]

% =============================================================================
% END APPENDED EIGHTH-SERIES V44 MATERIAL
% =============================================================================

% =============================================================================
% BEGIN APPENDED EIGHTH-SERIES V45 MATERIAL
% =============================================================================

\section{Eighth-series V45: compulsory companion audit and a
weighted one-loop Cauchy path}
\label{sec:e8v45}

\subsection{Compulsory five-version companion audit}

The active companion notes inspected before choosing the V45
direction were
\[
\begin{array}{c|r|r|l}
\text{source}&\text{lines}&\text{bytes}&\text{SHA--256}\\ \hline
\text{SM--Einstein V35}&13369&511768&
\texttt{4A509CE0DD1609EF8BC5E4A6B6A6B065ADC3B2EEEA1482D0A99B906378993E85}\\
\text{Navier--Stokes V26}&5319&278283&
\texttt{82C887F8C2C69E4F28FAD7C3DC8DE5B9099CB5A4BF431EBA1FFF3E91935978AD}.
\end{array}                                                     \tag{45.1}
\]
Their filenames are
\[
\begin{split}
&\texttt{Renewal\_SM\_Einstein\_Continuum\_Research\_Notes\_V35.tex},\\
&\texttt{Renewal\_NS\_Stability\_Research\_Notes\_V26.tex}.
\end{split}
\]

\begin{assessment}[V45 audit transfer]
\label{ass:e8v45-audit}
The following structural transfers are valid.
\begin{enumerate}
\item SM V31 proves that a common spectral gap, a weighted resolvent
      bound, and finite weighted parameter-path length control the
      full quasi-local transport.  The same Banach-algebra argument
      applies to derivatives of the YM vertical inverse and its
      one-loop contractions.
\item SM V35 independently uses
      \(H_1=R^*H_0R\), \(J_1=J_0|\det R|\): it confirms the
      measured-slice magnitude identity of V43.  Its residual
      \(U(1)\) phase belongs to a complex chiral determinant line;
      it is not an extra factor in the real positive YM Gaussian
      determinant.
\item SM V34--V35 also separates small curvature from small
      potential by retaining flat holonomy.  This agrees with the YM
      toron/harmonic ledger and supplies no permission to put those
      modes into the vertical inverse.
\item NS V26 sharpens pointwise kernel derivative norms on rank-one
      inputs.  No NS numerical constant transfers.  For YM, one
      axial colour polarization may represent the coefficient only
      after the complete Ward and little-group symmetry reduction;
      it cannot replace the full operator norm used in the path
      estimate below.
\end{enumerate}
\end{assessment}

\begin{decision}[Audit-adjusted direction]
\label{dec:e8v45-direction}
The first unclosed GGIC row is the Cauchy property of the
contact-subtracted one-loop moments.  V45 therefore proves a weighted
path derivative estimate for the complete tadpole--bubble kernel.
The next version must populate the path length for the actual
scale-dependent Gaussian reference sequence.
\end{decision}

\subsection{Weighted inverse and harmonic-router derivatives}

Use the weighted convolution row norm
\[
 \|K\|_\mu:=\sup_x\sum_y e^{\mu|x-y|_1}\|K(x,y)\|,              \tag{45.2}
\]
which is a Banach-algebra norm.  Let \(t\mapsto(A_t,B_t)\) be a
\(C^1\) path of contact-subtracted vertical and mixed Gaussian
blocks.  Put
\[
 G_t:=A_t^{-1},\qquad
 {\cal H}_t:=(-G_tB_t,I),                                      \tag{45.3}
\]
where \({\cal H}_t\) is the harmonic lift from a retained source to
the reduced-plus-retained tree slice.

\begin{lemma}[Weighted path derivative of the vertical inverse]
\label{lem:e8v45-inverse-path}
Assume
\[
 \sup_t\|G_t\|_\mu\le g,\qquad
 \sup_t\|B_t\|_\mu\le b,                                       \tag{45.4}
\]
and write
\[
 a_1(t):=\|\dot A_t\|_\mu,\qquad
 b_1(t):=\|\dot B_t\|_\mu.
\]
Then
\[
 \|\dot G_t\|_\mu\le g^2a_1(t),                                \tag{45.5}
\]
\[
 \|{\cal H}_t\|_\mu\le h:=1+gb,\qquad
 \|\dot{\cal H}_t\|_\mu
 \le d_H(t):=g^2b\,a_1(t)+g\,b_1(t).                           \tag{45.6}
\]
\end{lemma}

\begin{proof}
Differentiate \(A_tG_t=I\) to obtain
\(\dot G_t=-G_t\dot A_tG_t\).  The Banach-algebra inequality gives
(45.5).  The first estimate in (45.6) follows from (45.3), and
\[
 \frac d{dt}(-G_tB_t)
 =G_t\dot A_tG_tB_t-G_t\dot B_t
\]
gives the second.
\end{proof}

\subsection{A complete tadpole--bubble path bound}

Let \({\mathbb Q}_t\) denote the local quadratic background vertex
and \({\mathbb R}_t\) the local linear background vertex after the
declared contact owner has been separated.  Absorb the finite
plaquette-incidence norms into their stocks and assume
\[
 \|{\mathbb Q}_t\|_\mu\le q,\quad
 \|{\mathbb R}_t\|_\mu\le r,\quad
 \|\dot{\mathbb Q}_t\|_\mu\le q_1(t),\quad
 \|\dot{\mathbb R}_t\|_\mu\le r_1(t).                           \tag{45.7}
\]
Write the contact-subtracted one-loop kernel schematically, but
with its exact multilinearity, as
\[
 {\cal J}_t
 =\frac12{\mathfrak T}({\cal H}_t,{\cal H}_t,G_t,{\mathbb Q}_t)
 -\frac13{\mathfrak B}({\cal H}_t,{\cal H}_t,
 G_t,G_t,{\mathbb R}_t,{\mathbb R}_t).                          \tag{45.8}
\]
Here \(\mathfrak T\) and \(\mathfrak B\) are the fixed owner,
convolution, trace, and index contractions of the parity-cell
grammar.  Their operator norms are normalized to at most one by the
stocks in (45.7).  In the pure Wilson specialization, (45.8) is
exactly (42.12), with the second sign representing the positive
bubble.

\begin{theorem}[Weighted derivative of the complete one-loop kernel]
\label{thm:e8v45-loop-path}
Under (45.4)--(45.8),
\[
 \boxed{\|\dot{\cal J}_t\|_\mu\le{\cal L}(t),}                  \tag{45.9}
\]
where
\[
 \begin{split}
 {\cal L}(t):={}&
 \frac12\left[
 2h\,d_H(t)\,gq+h^2g^2a_1(t)q+h^2gq_1(t)\right]\\
 &+\frac13\left[
 2h\,d_H(t)\,g^2r^2
 +2h^2g^3a_1(t)r^2
 +2h^2g^2r\,r_1(t)\right].
                                                                    \tag{45.10}
 \end{split}
\]
Consequently, for \(0\le s<t\le1\),
\[
 \|{\cal J}_t-{\cal J}_s\|_\mu
 \le\int_s^t{\cal L}(\tau)\,d\tau.                             \tag{45.11}
\]
\end{theorem}

\begin{proof}
Differentiate the four-linear tadpole contraction in (45.8).
The two lift derivatives contribute
\(2h\,d_Hgq\), the inverse derivative contributes
\(h^2g^2a_1q\) by (45.5), and the vertex derivative contributes
\(h^2gq_1\).  Multiply by \(1/2\).

The bubble is six-linear.  Its two lift slots contribute
\(2h\,d_Hg^2r^2\), its two inverse slots contribute
\(2h^2g^3a_1r^2\), and its two vertex slots contribute
\(2h^2g^2rr_1\).  Multiply by \(1/3\).  Submultiplicativity of
(45.2) bounds every differentiated contraction and gives
(45.9)--(45.10).  Integrate the Banach-valued derivative to obtain
(45.11).
\end{proof}

\begin{remark}[Why the full kernel norm is used]
\label{rem:e8v45-full-kernel}
The coefficient may later be read on one axial polarization because
Ward and hypercubic symmetry scalarize that block.  The path estimate
(45.9) precedes that symmetry compression and therefore uses the
full weighted operator kernel.  This is the safe analogue of the
rank-one/full-norm distinction seen in the NS audit.
\end{remark}

\subsection{The moment-null contact quotient}

For a real even one-loop kernel \(K\), define its axial transverse
second-moment functional by
\[
 {\mathfrak m}(K)
 :=-\frac12\sum_xx_1^2K_{22}(x).                               \tag{45.12}
\]
The weighted estimate used in V42 gives
\[
 |{\mathfrak m}(K)|\le\frac1{\mu^2}\|K\|_\mu.                  \tag{45.13}
\]
Let
\[
 {\cal N}_{\rm ct}:=\ker{\mathfrak m}.                          \tag{45.14}
\]
Every zero-displacement contact kernel belongs to
\({\cal N}_{\rm ct}\), irrespective of its size.

\begin{theorem}[Contact-difference Cauchy telescope]
\label{thm:e8v45-contact-Cauchy}
Let \({\cal J}_n\) be the complete scale-\(n\) one-loop kernels.
Suppose their adjacent differences admit the exact owner split
\[
 {\cal J}_{n+1}-{\cal J}_n=C_n+R_n,\qquad
 C_n\in{\cal N}_{\rm ct},                                      \tag{45.15}
\]
and
\[
 \sum_{n=0}^\infty\|R_n\|_\mu<\infty.                          \tag{45.16}
\]
Then
\[
 \kappa_n:={\mathfrak m}({\cal J}_n)
\]
is Cauchy and
\[
 \boxed{
 |\kappa_\infty-\kappa_N|
 \le\frac1{\mu^2}\sum_{n=N}^\infty\|R_n\|_\mu.}                \tag{45.17}
\]
No bound on \(\|C_n\|_\mu\) is required.
\end{theorem}

\begin{proof}
Apply \({\mathfrak m}\) to (45.15).  The contact term vanishes and
(45.13) gives
\[
 |\kappa_{n+1}-\kappa_n|
 \le\mu^{-2}\|R_n\|_\mu.
\]
Sum from \(N\) to \(M-1\), use (45.16), and let \(M\to\infty\).
\end{proof}

\begin{corollary}[Path-length criterion for GGIC item 5]
\label{cor:e8v45-GGIC5}
Join the contact-subtracted reference data at scales \(n\) and
\(n+1\) by a path satisfying Theorem
\ref{thm:e8v45-loop-path}, and put
\[
 \Lambda_n:=\int_0^1{\cal L}_n(t)\,dt.                          \tag{45.18}
\]
If the exact scale difference has a moment-null contact owner and
\[
 \sum_n\Lambda_n<\infty,                                       \tag{45.19}
\]
then GGIC item 5 holds and
\[
 |\kappa_\infty-\kappa_N|
 \le\frac1{\mu^2}\sum_{n=N}^\infty\Lambda_n.                   \tag{45.20}
\]
\end{corollary}

\begin{proof}
Equation (45.11) bounds the renormalized adjacent difference by
\(\Lambda_n\).  Apply
Theorem~\ref{thm:e8v45-contact-Cauchy}.
\end{proof}

\begin{proposition}[Normalized coefficient telescope]
\label{prop:e8v45-normalized-Cauchy}
Let \(\sigma_n\ge\sigma_0>0\) be the classical marginal
normalizations and
\[
 b_n^\downarrow:=\frac{\kappa_n}{\sigma_n}.                    \tag{45.21}
\]
If (45.16) holds and
\(\sum_n|\sigma_{n+1}-\sigma_n|<\infty\), then
\(b_n^\downarrow\) is Cauchy.  If, as in (44.10),
\(\sigma_n=1/3\) exactly, then
\[
 \boxed{
 |b_\infty^\downarrow-b_N^\downarrow|
 \le\frac3{\mu^2}\sum_{n=N}^\infty\|R_n\|_\mu.}                \tag{45.22}
\]
\end{proposition}

\begin{proof}
The sequence \(\kappa_n\) is bounded and Cauchy by
Theorem~\ref{thm:e8v45-contact-Cauchy}.  The summable increments of
\(\sigma_n\), together with its positive floor, make \(1/\sigma_n\)
Cauchy.  Their product is Cauchy.  For the constant value \(1/3\),
multiply (45.17) by three.
\end{proof}

\begin{firewall}[A large contact is harmless only after exact
projection]
\label{fire:e8v45-contact-exact}
Theorem~\ref{thm:e8v45-contact-Cauchy} does not say that any large
local term may be ignored.  The owner \(C_n\) must be shown exactly
to lie in \(\ker{\mathfrak m}\).  A local \(q^2\) counterterm has a
nonzero second moment and belongs to the beta coefficient, not to
\({\cal N}_{\rm ct}\).
\end{firewall}

\subsection{Status and the V46 acquisition}

\begin{assessment}[Progress after V45]
\label{ass:e8v45-status}
The companion audit converts GGIC item 5 into two concrete
scale-by-scale obligations:
\begin{enumerate}
\item identify the exact moment-null contact owner in
      \({\cal J}_{n+1}-{\cal J}_n\); and
\item prove summability of the remaining parameter-path lengths
      \(\Lambda_n\).
\end{enumerate}
Once those hold, the actual one-loop coefficient exists with the
tail bound (45.20).  Its value still has to be matched to
\(-b_{\rm YM}\).
\end{assessment}

\begin{decision}[V46 target]
\label{dec:e8v45-next}
V46 should differentiate the repeated Gaussian block map (44.2),
separate its zero-displacement contact increment, and test whether
the Ward-relative remainder has a geometric weighted path length.
If the raw difference is not summable, the calculation must move to
the inverse precision or scale-difference determinant before another
estimate is attempted.
\end{decision}

\begin{firewall}[Claim boundary after V45]
\label{fire:e8v45-claim}
V45 performs the mandatory audit, proves the weighted one-loop path
derivative estimate, and proves the contact-null Cauchy telescope.
It does not populate the actual adjacent reference path, evaluate
the limiting coefficient, verify the nonlinear/complex QMC rows,
construct the full RG trajectory, remove the regulators, establish
Osterwalder--Schrader reconstruction, or prove the Yang--Mills mass
gap.
\end{firewall}

\subsection{Parent record}

This V45 note was formed by appending the preceding material to V44,
whose validated record was
\[
\begin{array}{c|c}
\text{lines}&20430\\
\text{bytes}&713995\\
\text{SHA-256}&
\texttt{6BD947C0BEBA400D648BCA77FF2800F44364775C37087E548A50D040A0C7353B}.
\end{array}
\]

% =============================================================================
% END APPENDED EIGHTH-SERIES V45 MATERIAL
% =============================================================================

% =============================================================================
% BEGIN APPENDED EIGHTH-SERIES V46 MATERIAL
% =============================================================================

\section{Eighth-series V46: canonical rescaling turns the perimeter
contact into a summable Ward-germ correction}
\label{sec:e8v46}

\subsection{The correct rescaled inverse symbol}

Let \(c_n,s_n\) be the repeated raw compliance and inverse symbol of
V44.  For a fixed physical momentum variable \(p\), define
\[
 S_n(p):=4^n s_n(2^{-n}p).                                     \tag{46.1}
\]
The factor \(4^n\) is the canonical two-derivative rescaling.  It is
essential: comparing the unrescaled \(s_n(q)\) at fixed \(q\) gives
the noncommuting limit in V44.

\begin{theorem}[Geometric convergence of the physical Ward germ]
\label{thm:e8v46-germ-convergence}
Let \(P<\infty\).  For every \(n\) with \(2^{-n}P\le1\) and every
real \(|p|\le P\),
\[
 \boxed{
 0\le\frac{p^2}{3}-S_n(p)
 \le\frac{C_{\rm ct}}9\,2^{-n}p^4.}                            \tag{46.2}
\]
Consequently
\[
 \boxed{
 |S_{n+1}(p)-S_n(p)|
 \le\frac{C_{\rm ct}}6\,2^{-n}p^4,}                            \tag{46.3}
\]
and \(S_n(p)\to p^2/3\) uniformly on every bounded physical
momentum interval.
\end{theorem}

\begin{proof}
Lemma~\ref{lem:e8v44-contact-bound} gives
\[
 c_n(q\widehat e_1)=\frac3{q^2}+R_n(q),\qquad
 0\le R_n(q)\le C_{\rm ct}2^n.                                \tag{46.4}
\]
Set \(q=2^{-n}p\).  Direct inversion yields
\[
 S_n(p)
 =\frac{p^2}{3}
 \left(1+\frac{2^{-2n}p^2R_n(2^{-n}p)}3\right)^{-1}.            \tag{46.5}
\]
The quantity in parentheses exceeds one and its excess is at most
\((C_{\rm ct}/3)2^{-n}p^2\).  Since
\(0\le1-(1+x)^{-1}\le x\) for \(x\ge0\), (46.2) follows.
Apply (46.2) at \(n\) and \(n+1\) and use the triangle inequality;
\((1+1/2)/9=1/6\), proving (46.3).
\end{proof}

\begin{definition}[Quartic Ward-germ norm]
\label{def:e8v46-germ-norm}
For an even scalar symbol \(F\) with
\(F(0)=F''(0)=0\), set
\[
 \|F\|_{{\cal G},P}
 :=\sup_{0<|p|\le P}\frac{|F(p)|}{p^4}.                         \tag{46.6}
\]
\end{definition}

\begin{corollary}[Summable Gaussian reference path in germ norm]
\label{cor:e8v46-germ-Cauchy}
For all sufficiently large \(n\),
\[
 \left\|S_n-\frac{p^2}{3}\right\|_{{\cal G},P}
 \le\frac{C_{\rm ct}}9\,2^{-n},                                \tag{46.7}
\]
\[
 \|S_{n+1}-S_n\|_{{\cal G},P}
 \le\frac{C_{\rm ct}}6\,2^{-n}.                                \tag{46.8}
\]
In particular,
\[
 \sum_{n=N}^\infty
 \|S_{n+1}-S_n\|_{{\cal G},P}
 \le\frac{C_{\rm ct}}3\,2^{-N}.                                \tag{46.9}
\]
\end{corollary}

\begin{proof}
Divide (46.2)--(46.3) by \(p^4\), take the supremum, and sum the
geometric series.
\end{proof}

\begin{corollary}[The divergent compliance contact is a contracting
quartic precision owner]
\label{cor:e8v46-contact-to-quartic}
Let \(p_n\) be (44.11).  The quartic Taylor coefficient of \(S_n\)
is
\[
 -\frac{4^{-n}p_n}{9},                                        \tag{46.10}
\]
and
\[
 |4^{-n}p_n|=O(2^{-n}).                                       \tag{46.11}
\]
Thus the \(O(2^n)\) compliance contact belongs to a canonically
stable inverse-precision coordinate after physical rescaling.
\end{corollary}

\begin{proof}
Substitute \(q=2^{-n}p\) in (44.12) and multiply by \(4^n\).
Then use (44.13).
\end{proof}

\begin{firewall}[Germ summability is not weighted-kernel summability]
\label{fire:e8v46-two-norm}
Equations (46.7)--(46.9) control the low external-momentum germ.
The loop in (42.26) samples every internal Brillouin momentum, so
these estimates alone do not imply the weighted spatial path length
(45.19).  A low/edge decomposition must connect the germ norm to the
complete one-loop kernel without discarding a marginal \(q^2\)
counterterm.
\end{firewall}

\subsection{The exact nested Gaussian determinant cocycle}

The scale comparison is cleaner at determinant level than after
expanding every tadpole and bubble separately.  The following
finite-dimensional identity is the algebraic owner of that
comparison.

\begin{theorem}[Nested Schur determinant cocycle]
\label{thm:e8v46-Schur-cocycle}
Let
\[
 {\cal V}=R_0\oplus R_1\oplus\cdots\oplus R_{N-1}\oplus C
\]
and let \(H^{(0)}=H>0\) be a positive Hermitian operator on
\({\cal V}\).  Inductively, after \(R_0,\ldots,R_{j-1}\) have been
eliminated, write
\[
 H^{(j)}
 =\begin{pmatrix}A_j&B_j^*\\ B_j&D_j\end{pmatrix}
 \quad\hbox{on}\quad
 R_j\oplus(R_{j+1}\oplus\cdots\oplus C)
\]
and set
\[
 H^{(j+1)}:=D_j-B_jA_j^{-1}B_j^*.                             \tag{46.12}
\]
Then
\[
 \boxed{\det H=\left(\prod_{j=0}^{N-1}\det A_j\right)
                 \det H^{(N)}.}                              \tag{46.13}
\]
If all blocks depend \(C^2\)-smoothly on a background \(v\), then
for any two background directions \(u,w\),
\[
 D^2_{u,w}\log\det H
 =\sum_{j=0}^{N-1}D^2_{u,w}\log\det A_j
  +D^2_{u,w}\log\det H^{(N)}.                                \tag{46.14}
\]
\end{theorem}

\begin{proof}
For a positive block matrix,
\[
 \begin{pmatrix}I&0\\-B_jA_j^{-1}&I\end{pmatrix}
 H^{(j)}
 \begin{pmatrix}I&-A_j^{-1}B_j^*\\0&I\end{pmatrix}
 =\begin{pmatrix}A_j&0\\0&H^{(j+1)}\end{pmatrix}.
\]
The triangular factors have determinant one, hence
\(\det H^{(j)}=\det A_j\det H^{(j+1)}\).  Iteration proves
(46.13).  Positivity keeps every determinant positive, so taking
logarithms and differentiating twice proves (46.14).
\end{proof}

\begin{corollary}[No double counting of Gaussian shells]
\label{cor:e8v46-shell-no-double-count}
Let
\[
 \Gamma^{[0,N)}_{\rm G}(v)
 :=\frac12\log\det H(v)-\frac12\log\det H^{(N)}(v).
\]
Then
\[
 \boxed{\Gamma^{[0,N)}_{\rm G}(v)
 =\frac12\sum_{j=0}^{N-1}\log\det A_j(v).}                    \tag{46.15}
\]
Consequently the background Hessian of the left side is exactly
the sum of the shell tadpole--bubble Hessians on the right.  This is
an identity before estimates and not an asymptotic RG assertion.
\end{corollary}

\begin{proof}
Subtract \(\log\det H^{(N)}\) from the logarithm of (46.13), and
then use (46.14).
\end{proof}

\begin{remark}[Measured quotient version]
\label{rem:e8v46-measured-cocycle}
If a nontrivial slice is used, (46.15) applies to the Hessian
determinants while the measured action also contains the logarithm
of the orbit Jacobian.  The same cocycle is valid only after that
Jacobian has been factored shell by shell.  For the parity-tree
coordinates used by the attached exact regulator its finite
Jacobian is one, as recorded in V41.  This remark prevents an
unlicensed replacement of the quotient measure by a ghost
determinant.
\end{remark}

\subsection{A quantitative low-sector compiler}

We now state exactly what is needed to turn (46.8) into the
weighted one-loop path length of V45.  This isolates a finite
locality problem instead of hiding it in an informal continuity
claim for the vertices.

\begin{definition}[Uniform one-loop path multiplier]
\label{def:e8v46-loop-multiplier}
With the stocks \(g,b,h,q,r\) of (45.4)--(45.7), put
\[
\begin{split}
 \alpha_A&:=hg^3bq+\frac12h^2g^2q
       +\frac23hg^4br^2+\frac23h^2g^3r^2,\\
 \alpha_B&:=hg^2q+\frac23hg^3r^2,\\
 \alpha_Q&:=\frac12h^2g,\qquad
 \alpha_R:=\frac23h^2g^2r,\\
 C_{\rm loop}&:=\max\{\alpha_A,\alpha_B,\alpha_Q,\alpha_R\}.
                                                                    \tag{46.16}
\end{split}
\]
\end{definition}

\begin{lemma}[Linear form of the V45 path estimate]
\label{lem:e8v46-linear-loop}
Under the hypotheses of
Theorem~\ref{thm:e8v45-loop-path},
\[
 {\cal L}(t)
 =\alpha_Aa_1(t)+\alpha_Bb_1(t)
  +\alpha_Qq_1(t)+\alpha_Rr_1(t),                              \tag{46.17}
\]
and hence
\[
 \int_0^1{\cal L}(t)\,dt
 \le C_{\rm loop}\int_0^1
 (a_1+b_1+q_1+r_1)(t)\,dt.                                    \tag{46.18}
\]
\end{lemma}

\begin{proof}
Insert
\(d_H=g^2b\,a_1+g\,b_1\) from (45.6) into (45.10), collect the
four derivative stocks, and compare each coefficient with their
maximum.
\end{proof}

\begin{proposition}[Low-germ-to-loop implication]
\label{prop:e8v46-low-compiler}
Fix \(P\).  Suppose the actual contact-subtracted low-sector data
\[
 \Theta(S)=\bigl(A(S),B(S),{\mathbb Q}(S),{\mathbb R}(S)\bigr)
\]
admit, between \(S_n\) and \(S_{n+1}\), a path with the uniform
stocks of V45 and
\[
 \int_0^1(a_1+b_1+q_1+r_1)(t)\,dt
 \le C_\Theta\|S_{n+1}-S_n\|_{{\cal G},P}.                    \tag{46.19}
\]
Then the associated low-sector one-loop kernels satisfy
\[
 \|{\cal J}^{\rm low}_{n+1}-{\cal J}^{\rm low}_n\|_\mu
 \le C_{\rm low}2^{-n},\qquad
 C_{\rm low}:=\frac{C_{\rm loop}C_\Theta C_{\rm ct}}6.         \tag{46.20}
\]
\end{proposition}

\begin{proof}
Use (46.18), the path estimate (45.11), and (46.8).
\end{proof}

\begin{firewall}[What is conditional in (46.20)]
\label{fire:e8v46-low-conditional}
The contraction (46.8) is proved for the exact repeated Gaussian
symbol.  Estimate (46.19) is not yet proved for the actual
background-dependent \(45\)-component vertical block, harmonic
router, and cubic and quartic Wilson vertices.  In particular,
(46.20) is an implication with a sharply named acquisition, not a
completed estimate for the Yang--Mills shell.
\end{firewall}

\subsection{The edge cannot be renamed a contact}

\begin{lemma}[Constant and two-derivative local owners]
\label{lem:e8v46-local-moments}
Let \(\delta_0\) denote the origin kernel and let
\[
 (-\Delta_1\delta)_x:=2\delta_{x,0}
                  -\delta_{x,e_1}-\delta_{x,-e_1}.
\]
For the moment functional (45.12),
\[
 {\mathfrak m}(a\delta_0)=0,\qquad
 {\mathfrak m}\bigl(a(-\Delta_1\delta)\bigr)=a.                \tag{46.21}
\]
\end{lemma}

\begin{proof}
The first sum contains the factor \(x_1^2\) at \(x=0\).  In the
second sum only \(x=\pm e_1\) contribute, each with kernel
coefficient \(-a\).  Therefore
\(-\frac12((-a)+(-a))=a\).
\end{proof}

\begin{corollary}[Marginal edge terms must be estimated]
\label{cor:e8v46-edge-marginal}
A zero-displacement edge contribution may be placed in
\(\ker{\mathfrak m}\), but a local \(q_1^2\) contribution may not.
The latter changes the one-loop coefficient by its exact
coefficient and must either cancel by a determinant cocycle or have
a summable adjacent-scale increment.
\end{corollary}

\begin{definition}[Low/edge reference bridge card]
\label{def:e8v46-LERB}
At adjacent scales \(n,n+1\), a completed low/edge reference bridge
card (LERB) consists of:
\begin{enumerate}
\item an exact partition of the same regulated determinant Hessian,
      not two independently normalized approximations;
\item a low-sector path satisfying (46.19);
\item an identity
      \[
       {\cal J}^{\rm edge}_{n+1}-{\cal J}^{\rm edge}_n
       =C_n^{\rm edge}+R_n^{\rm edge},\quad
       {\mathfrak m}(C_n^{\rm edge})=0;                        \tag{46.22}
      \]
\item a verified constant \(C_{\rm edge}\), independent of \(n\),
      such that
      \[
       \|R_n^{\rm edge}\|_\mu\le C_{\rm edge}2^{-n};           \tag{46.23}
      \]
\item a common Ward restoration showing that the sum of the two
      sectors, including their overlap terms, is transverse.
\end{enumerate}
\end{definition}

\begin{theorem}[Two-scale Cauchy compiler]
\label{thm:e8v46-two-scale-compiler}
Assume a LERB is completed at every \(n\ge N\), with uniform
constants and with the exact total adjacent difference
\[
 {\cal J}_{n+1}-{\cal J}_n
 =C_n+R_n^{\rm low}+R_n^{\rm edge},\qquad
 {\mathfrak m}(C_n)=0,                                       \tag{46.24}
\]
where (46.20) and (46.23) hold.  Then the normalized coefficient
\(b_n^\downarrow=3{\mathfrak m}({\cal J}_n)\) converges and
\[
 \boxed{
 |b_\infty^\downarrow-b_N^\downarrow|
 \le\frac{6(C_{\rm low}+C_{\rm edge})}{\mu^2}\,2^{-N}.}        \tag{46.25}
\]
\end{theorem}

\begin{proof}
The residual in (46.24) has weighted norm at most
\((C_{\rm low}+C_{\rm edge})2^{-n}\).  Apply (45.22) and sum
\(\sum_{n=N}^\infty2^{-n}=2^{1-N}\).
\end{proof}

\begin{firewall}[A cutoff partition is not automatically a Ward
partition]
\label{fire:e8v46-cutoff-Ward}
Multiplying an integrand by low and high momentum cutoffs can create
boundary terms when the Ward identity shifts loop momentum.
LERB item 5 therefore requires either a periodic partition stable
under the exact finite-lattice Ward algebra or an explicit overlap
owner.  Transversality may be asserted only after the two sectors
are recombined.
\end{firewall}

\subsection{What V46 proves and where the calculation moves}

\begin{assessment}[V46 ledger]
\label{ass:e8v46-ledger}
The exact achievements are:
\begin{enumerate}
\item canonical rescaling converts the divergent raw compliance
      contact into the geometric Ward-germ estimates
      (46.2)--(46.11);
\item successive Gaussian eliminations obey the exact determinant
      cocycle (46.13)--(46.15);
\item the V45 tadpole--bubble bound has the explicit linear
      multiplier (46.16)--(46.18);
\item a completed LERB would give the quantitative coefficient tail
      (46.25); and
\item (46.21) forbids hiding a marginal \(q^2\) edge term among
      moment-null contacts.
\end{enumerate}
The unproved rows are precisely the actual low data lift (46.19)
and the edge identity and estimate (46.22)--(46.23).
\end{assessment}

\begin{decision}[Upstream V47 target]
\label{dec:e8v46-next}
Rather than estimate the edge-expanded vertices separately, V47
should apply the exact Schur cocycle to two consecutive dyadic
eliminations.  The target is an adjacent-scale determinant
commutator: identify the common two-step determinant, cancel the
shared shell before differentiation, and isolate the sole
associativity defect caused by the chosen background embedding.
If the embeddings commute exactly, the edge owner is algebraic; if
they do not, their commutator is the object that must satisfy the
geometric bound.
\end{decision}

\begin{firewall}[Claim boundary after V46]
\label{fire:e8v46-claim}
V46 proves a rescaled Gaussian germ contraction, an exact nested
determinant identity, and conditional compilers for the one-loop
coefficient.  It does not complete either missing LERB row,
evaluate the limiting coefficient, prove that it equals the
universal negative Yang--Mills beta coefficient, close the
nonlinear/complex quotient matching card, construct the full RG
trajectory, remove the regulators, establish
Osterwalder--Schrader reconstruction, or prove the Yang--Mills mass
gap.
\end{firewall}

\subsection{Parent record}

This V46 note was formed by appending the preceding material to V45,
whose validated record was
\[
\begin{array}{c|c}
\text{lines}&20794\\
\text{bytes}&726346\\
\text{SHA-256}&
\texttt{27552637526409080C99F352EE1E0ED28A320C9CE644B0B2B5FE0095E3A4DA8B}.
\end{array}
\]

% =============================================================================
% END APPENDED EIGHTH-SERIES V46 MATERIAL
% =============================================================================

% =============================================================================
% BEGIN APPENDED EIGHTH-SERIES V47 MATERIAL
% =============================================================================

\section{Eighth-series V47: exact blocking associativity and the
reference-reset commutator}
\label{sec:e8v47}

\subsection{Straight holonomy has no associativity defect}

Let \(G\) be the compact gauge group.  If
\(\gamma=(e_1,\ldots,e_m)\) is an oriented lattice path, use the
standard conventions
\[
 U_{\bar e}=U_e^{-1},\qquad
 U(\gamma):=U_{e_1}\cdots U_{e_m}.                             \tag{47.1}
\]

\begin{lemma}[Path concatenation]
\label{lem:e8v47-path-concat}
For composable oriented paths \(\gamma_1,\gamma_2\),
\[
 U(\gamma_1\star\gamma_2)=U(\gamma_1)U(\gamma_2).              \tag{47.2}
\]
Consequently, if \(Q_2\) is dyadic straight-link blocking and
\(Q_4\) is direct blocking along the same oriented length-four
path, then
\[
 \boxed{Q_2\circ Q_2=Q_4}                                    \tag{47.3}
\]
on every fine link configuration, not merely to linear order.
\end{lemma}

\begin{proof}
Equation (47.2) is the definition (47.1) followed by associativity
of multiplication in \(G\).  On a length-four path the first
application of \(Q_2\) produces
\(U_{e_1}U_{e_2}\) and \(U_{e_3}U_{e_4}\).  The second produces
\((U_{e_1}U_{e_2})(U_{e_3}U_{e_4})\), which equals the direct
ordered product.  Reversed orientations obey the same statement
because \((ab)^{-1}=b^{-1}a^{-1}\).
\end{proof}

\begin{corollary}[No logarithmic branch defect near the vacuum]
\label{cor:e8v47-no-log-defect}
Suppose the blocked background coordinate is defined by the
principal logarithm of the coarse holonomy on a neighbourhood
\({\cal U}\subset G\) where that logarithm is single valued.  If
both sides of (47.3) lie in \({\cal U}\), their Lie-algebra
coordinates agree exactly.
\end{corollary}

\begin{proof}
The group elements agree by (47.3), and a single-valued function
takes equal values on equal arguments.
\end{proof}

\begin{firewall}[The domain of the logarithm matters]
\label{fire:e8v47-log-domain}
Corollary~\ref{cor:e8v47-no-log-defect} is a local-vacuum statement.
It does not define a global logarithm on \(G\) and cannot be used to
discard large-field sectors.  The exact group identity (47.3),
which needs no logarithm, is the global statement.
\end{firewall}

\subsection{Exact finite-volume RG semigroup}

Let \(X_j=G^{E_j}\) be the link-configuration space on the
\(j\)-th lattice and let \(d\mu_j\) be normalized product Haar
measure.  Denote the group delta distribution by \(\delta_G\), with
\[
 \int_G f(W)\delta_G(W^{-1}V)\,dW=f(V).                        \tag{47.4}
\]
For a nonnegative integrable density \(\rho\) on \(X_j\), define
the exact pushforward density
\[
 ({\cal T}_{j}\rho)(V)
 :=\int_{X_j}\rho(U)
 \prod_{\ell\in E_{j+1}}
 \delta_G\bigl(V_\ell^{-1}(Q_jU)_\ell\bigr)\,d\mu_j(U).        \tag{47.5}
\]
Unused fine variables are integrated with normalized Haar measure.

\begin{theorem}[Exact two-step blocking law]
\label{thm:e8v47-RG-semigroup}
Let \(Q_{j:j+2}:=Q_{j+1}\circ Q_j\), and define
\({\cal T}_{j:j+2}\) by (47.5) using that composite block map.
Then, as densities on \(X_{j+2}\),
\[
 \boxed{{\cal T}_{j+1}{\cal T}_j
       ={\cal T}_{j:j+2}.}                                   \tag{47.6}
\]
For nested straight dyadic links, (47.3) identifies the right side
with direct length-four blocking.
\end{theorem}

\begin{proof}
For \(Z\in X_{j+2}\), insert (47.5) twice:
\[
\begin{split}
 ({\cal T}_{j+1}{\cal T}_j\rho)(Z)
 ={}&\int_{X_{j+1}}\int_{X_j}\rho(U)
 \prod_{\ell\in E_{j+1}}
 \delta_G\bigl(V_\ell^{-1}(Q_jU)_\ell\bigr)\\
 &\quad\cdot
 \prod_{L\in E_{j+2}}
 \delta_G\bigl(Z_L^{-1}(Q_{j+1}V)_L\bigr)
 \,d\mu_j(U)d\mu_{j+1}(V).                                  \tag{47.7}
\end{split}
\]
The integrand is nonnegative, so Tonelli's theorem permits either
order of integration.  Repeated use of (47.4) sets \(V=Q_jU\) in
the second delta product and gives
\[
 \int_{X_j}\rho(U)
 \prod_{L\in E_{j+2}}
 \delta_G\bigl(
 Z_L^{-1}(Q_{j+1}Q_jU)_L\bigr)\,d\mu_j(U),
\]
which is \({\cal T}_{j:j+2}\rho\).
\end{proof}

\begin{corollary}[Conservation of total mass]
\label{cor:e8v47-mass-conservation}
For every nonnegative integrable \(\rho\),
\[
 \int_{X_{j+1}}{\cal T}_j\rho\,d\mu_{j+1}
 =\int_{X_j}\rho\,d\mu_j.                                    \tag{47.8}
\]
\end{corollary}

\begin{proof}
Integrate (47.5) in \(V\), apply Tonelli, and use (47.4) once for
each coarse link.
\end{proof}

\begin{remark}[Gauge covariance]
\label{rem:e8v47-gauge-covariance}
The straight holonomy transforms by the gauge elements at its two
endpoints.  Hence \({\cal T}_j\) sends gauge-invariant densities to
gauge-invariant densities.  Equations (47.6)--(47.8) are statements
on the unreduced compact integral.  Any slice calculation must
reproduce them together with its orbit Jacobian.
\end{remark}

\subsection{Consequences for the one-loop determinant}

For a strictly positive density write
\[
 {\cal R}_j(A):=-\log{\cal T}_j(e^{-A}),                       \tag{47.9}
\]
up to a background-independent normalization.

\begin{corollary}[Exact action semigroup]
\label{cor:e8v47-action-semigroup}
On its finite-volume domain,
\[
 {\cal R}_{j+1}({\cal R}_j(A))
 =-\log{\cal T}_{j:j+2}(e^{-A})                               \tag{47.10}
\]
up to one additive constant independent of the retained field.
\end{corollary}

\begin{proof}
Exponentiating the left side and using (47.9) gives
\({\cal T}_{j+1}{\cal T}_j(e^{-A})\).  Apply (47.6) and take the
negative logarithm.
\end{proof}

\begin{proposition}[Associativity survives coefficient extraction]
\label{prop:e8v47-loop-associativity}
Assume the two sides of (47.10), in a fixed local vacuum chart and
with the same regulators, have asymptotic expansions in a parameter
\(\hbar\) whose coefficients are unique.  Then their coefficients
agree at every order.  In particular, the Gaussian determinant and
its background Hessian agree at one loop for direct and iterated
blocking.
\end{proposition}

\begin{proof}
The functions are equal by (47.10).  Subtract their asymptotic
series.  If the first nonzero coefficient occurred at order
\(\hbar^k\), division by \(\hbar^k\) and passage to
\(\hbar\downarrow0\) would give a nonzero limit for the zero
function, a contradiction.  Background differentiation preserves
the equality wherever the assumed expansion is differentiable.
At Gaussian order this is also the explicit Schur identity
(46.13)--(46.15).
\end{proof}

\begin{firewall}[What the semigroup theorem does not compare]
\label{fire:e8v47-semigroup-scope}
Theorem~\ref{thm:e8v47-RG-semigroup} compares two ways of performing
the same exact integral.  It does not say that an exact blocked
action lies again in the one-parameter Wilson family, nor that
replacing it by a selected reference action commutes with the next
block.
\end{firewall}

\subsection{The actual defect is the reference reset}

Let \({\cal P}_{j+1}\) be the declared retraction from the exact
effective action to the reference family used to define the next
Gaussian shell.  No linearity of \({\cal P}_{j+1}\) is assumed.
Define the pre-final-reset commutator
\[
 {\cal D}_j(A)
 :={\cal R}_{j+1}{\cal P}_{j+1}{\cal R}_j(A)
   -{\cal R}_{j+1}{\cal R}_j(A).                              \tag{47.11}
\]

\begin{proposition}[Exact localization of the associativity defect]
\label{prop:e8v47-reset-owner}
All failure of a reset reference evolution to agree with exact
two-step blocking is contained in \({\cal D}_j\).  In particular,
\[
 {\cal P}_{j+1}{\cal R}_j(A)={\cal R}_j(A)
 \quad\Longrightarrow\quad {\cal D}_j(A)=0.                   \tag{47.12}
\]
There is no additional straight-holonomy or Fubini associator.
\end{proposition}

\begin{proof}
The second term of (47.11) is exact direct two-step blocking by
(47.10); the first differs only by insertion of
\({\cal P}_{j+1}\).  If that insertion fixes the intermediate
action, the two arguments of \({\cal R}_{j+1}\) coincide.
Equations (47.3) and (47.6) exclude the two asserted additional
owners.
\end{proof}

\begin{definition}[One-loop reset kernel]
\label{def:e8v47-reset-kernel}
Let \(v\) be a retained background coordinate and let
\({\cal D}^{(1)}_j(v)\) denote the one-loop coefficient of
\({\cal D}_j(A)(v)\).  Its transverse quadratic kernel is
\[
 {\cal K}^{\rm reset}_j
 :=D_v^2{\cal D}^{(1)}_j(0).                                  \tag{47.13}
\]
\end{definition}

\begin{theorem}[Reset criterion for the missing LERB edge row]
\label{thm:e8v47-reset-LERB}
Suppose, after canonical rescaling and common Ward restoration,
\[
 {\cal K}^{\rm reset}_j=C_j^{\rm reset}+R_j^{\rm reset},
 \qquad {\mathfrak m}(C_j^{\rm reset})=0,\qquad
 \|R_j^{\rm reset}\|_\mu\le C_{\rm reset}2^{-j}.               \tag{47.14}
\]
Suppose also that all non-reset low-sector terms obey (46.20).
Then the LERB edge row can be populated with
\(C_{\rm edge}=C_{\rm reset}\), and
\[
 |b_\infty^\downarrow-b_N^\downarrow|
 \le
 \frac{6(C_{\rm low}+C_{\rm reset})}{\mu^2}\,2^{-N}.           \tag{47.15}
\]
\end{theorem}

\begin{proof}
Exact direct and iterated integration cancel by
Proposition~\ref{prop:e8v47-loop-associativity}.  Therefore the
only edge discrepancy introduced by the reference evolution is
(47.13).  Decomposition (47.14) is exactly (46.22)--(46.23).
Apply Theorem~\ref{thm:e8v46-two-scale-compiler}.
\end{proof}

\begin{corollary}[Second-moment matching is the exact marginal
reset condition]
\label{cor:e8v47-moment-matching}
If the reset is chosen so that
\[
 {\mathfrak m}({\cal K}^{\rm reset}_j)=0                      \tag{47.16}
\]
at every scale, then it produces no change in the normalized
one-loop coefficient at that scale.  Conversely, a reset with
nonzero value in (47.16) changes that coefficient by
\(3{\mathfrak m}({\cal K}^{\rm reset}_j)\).
\end{corollary}

\begin{proof}
The coefficient is
\(b_j^\downarrow=3{\mathfrak m}({\cal J}_j)\) by V45.
Linearity of \(\mathfrak m\) gives both claims.
\end{proof}

\begin{firewall}[Matching the classical kinetic term is
insufficient]
\label{fire:e8v47-classical-vs-loop}
A retraction may preserve the classical \(q^2/3\) Ward germ while
changing the one-loop second moment (47.16).  The reference reset
must therefore be matched at the determinant-Hessian level, not
only at the classical inverse symbol.
\end{firewall}

\subsection{Status and the V48 acquisition}

\begin{assessment}[The V47 upstream resolution]
\label{ass:e8v47-status}
The possible blocking-order bottleneck is resolved exactly:
straight group holonomy is associative, the compact finite-volume
pushforward has the semigroup property, and its one-loop expansion
inherits that identity.  The unknown edge term has consequently
been localized to a named object, the reference-reset kernel
\({\cal K}^{\rm reset}_j\), together with canonical rescaling and
any quotient-measure mismatch.
\end{assessment}

\begin{decision}[V48 target]
\label{dec:e8v47-next}
V48 should write the reference retraction explicitly as
second-moment extraction followed by Wilson-kernel reconstruction.
The first question is algebraic: whether this retraction makes
(47.16) exact.  The second is analytic: whether the complementary
irrelevant reset kernel contracts in \(\|\cdot\|_\mu\) at rate
\(2^{-j}\).  If weighted contraction fails, the calculation must
retain a finite irrelevant basis rather than repeatedly project to
one Wilson coupling.
\end{decision}

\begin{firewall}[Claim boundary after V47]
\label{fire:e8v47-claim}
V47 proves exact associativity of the compact blocking integral and
identifies the reference reset as the sole blocking-order defect.
It does not establish (47.14) or (47.16) for the actual retraction,
complete the LERB, evaluate the universal one-loop coefficient,
close the nonlinear/complex quotient matching card, construct the
full RG trajectory, remove the regulators, establish
Osterwalder--Schrader reconstruction, or prove the Yang--Mills mass
gap.
\end{firewall}

\subsection{Parent record}

This V47 note was formed by appending the preceding material to V46,
whose validated record was
\[
\begin{array}{c|c}
\text{lines}&21241\\
\text{bytes}&741227\\
\text{SHA-256}&
\texttt{9FE733C6324D19F6B6836E4F6AFAF41A151B64A4C1BC5A7968FC942B90C0F39E}.
\end{array}
\]

% =============================================================================
% END APPENDED EIGHTH-SERIES V47 MATERIAL
% =============================================================================

% =============================================================================
% BEGIN APPENDED EIGHTH-SERIES V48 MATERIAL
% =============================================================================

\section{Eighth-series V48: canonical moment retraction and scalar
reset mixing}
\label{sec:e8v48}

\subsection{The canonical projection onto the marginal direction}

Let \({\cal X}_\mu\) be the real Banach space of translation
invariant transverse quadratic kernels with the weighted norm used
in V45.  Retain the continuous functional
\[
 {\mathfrak m}(K)=-\frac12\sum_xx_1^2K_{22}(x),\qquad
 |{\mathfrak m}(K)|\le\mu^{-2}\|K\|_\mu.                      \tag{48.1}
\]
Choose once and for all the transverse Wilson kinetic Hessian
\({\cal W}\in{\cal X}_\mu\), normalized by
\[
 {\mathfrak m}({\cal W})=1,\qquad
 w_\mu:=\|{\cal W}\|_\mu.                                    \tag{48.2}
\]

\begin{definition}[Moment retraction]
\label{def:e8v48-moment-retraction}
Define bounded linear maps
\[
 {\cal P}K:={\mathfrak m}(K){\cal W},\qquad
 {\cal Q}:=I-{\cal P}.                                       \tag{48.3}
\]
\end{definition}

\begin{proposition}[Exact marginal--irrelevant splitting]
\label{prop:e8v48-splitting}
The maps in (48.3) obey
\[
 {\cal P}^2={\cal P},\quad {\cal Q}^2={\cal Q},\quad
 {\cal P}{\cal Q}={\cal Q}{\cal P}=0,\quad
 {\mathfrak m}({\cal Q}K)=0,                                 \tag{48.4}
\]
and
\[
 \|{\cal P}K\|_\mu\le w_\mu\mu^{-2}\|K\|_\mu,\qquad
 \|{\cal Q}K\|_\mu
 \le(1+w_\mu\mu^{-2})\|K\|_\mu.                              \tag{48.5}
\]
Moreover
\[
 {\cal X}_\mu
 =\operatorname{span}\{{\cal W}\}
  \mathbin{\dotplus}\ker{\mathfrak m}.                        \tag{48.6}
\]
\end{proposition}

\begin{proof}
Using (48.2),
\({\cal P}^2K={\mathfrak m}(K){\mathfrak m}({\cal W}){\cal W}
={\cal P}K\).  The remaining algebra in (48.4) follows from
\({\cal Q}=I-{\cal P}\).  In particular,
\({\mathfrak m}({\cal Q}K)
={\mathfrak m}(K)-{\mathfrak m}(K){\mathfrak m}({\cal W})=0\).
The first estimate in (48.5) follows from (48.1), and the second
from the triangle inequality.  Finally
\(K={\cal P}K+{\cal Q}K\), and the intersection of the two spaces
in (48.6) is zero by (48.2).
\end{proof}

\begin{corollary}[Same-scale Wilson reset is exactly moment
preserving]
\label{cor:e8v48-same-scale}
Replacing a quadratic kernel \(K\) by \({\cal P}K\) changes it by
\[
 {\cal P}K-K=-{\cal Q}K\in\ker{\mathfrak m}.                  \tag{48.7}
\]
Thus this reset neither changes the one-loop second moment nor its
normalized coefficient at the scale where the reset is made.
\end{corollary}

\begin{proof}
This is (48.4), followed by the definition
\(b^\downarrow=3{\mathfrak m}(K)\).
\end{proof}

\begin{firewall}[A nonlinear reference retraction needs its tangent]
\label{fire:e8v48-nonlinear-reset}
Corollary~\ref{cor:e8v48-same-scale} applies to the quadratic
one-loop sector if the differential of the full reference
retraction on that sector is precisely \({\cal P}\).  A nonlinear
retraction can mix lower-order couplings into the quadratic
one-loop kernel.  Those chain-rule terms must be included before
calling the reset moment preserving.
\end{firewall}

\subsection{One-step memory of the discarded irrelevant kernel}

Let \({\cal F}_j:{\cal X}_\mu\to{\cal X}_\mu\) denote the
canonically rescaled one-loop kernel map induced by the next exact
shell, with all other reference data held on their declared paths.
For an intermediate kernel \(K\), define
\[
 \Delta_j(K):={\cal F}_j({\cal P}K)-{\cal F}_j(K).             \tag{48.8}
\]

\begin{theorem}[Exact reset-memory formula]
\label{thm:e8v48-reset-memory}
Suppose \({\cal F}_j\) is continuously Fréchet differentiable on
the segment
\[
 Z_t:=K-t{\cal Q}K,\qquad0\le t\le1.                          \tag{48.9}
\]
Then
\[
 \boxed{\Delta_j(K)
 =-\int_0^1D{\cal F}_j(Z_t)[{\cal Q}K]\,dt.}                  \tag{48.10}
\]
Writing
\[
 \Delta_j(K)=C_j(K)+R_j(K),\qquad
 C_j:={\cal Q}\Delta_j,\quad R_j:={\cal P}\Delta_j,           \tag{48.11}
\]
gives the exact properties
\[
 {\mathfrak m}(C_j)=0,\qquad
 R_j={\mathfrak m}(\Delta_j){\cal W}.                         \tag{48.12}
\]
\end{theorem}

\begin{proof}
Since \(Z_0=K\), \(Z_1={\cal P}K\), and
\(\dot Z_t=-{\cal Q}K\), the Banach-space fundamental theorem of
calculus gives (48.10).  Apply the complementary projections
(48.3) and use Proposition~\ref{prop:e8v48-splitting}.
\end{proof}

\begin{definition}[Irrelevant-to-marginal mixing modulus]
\label{def:e8v48-mixing-modulus}
On the segment (48.9), set
\[
 \chi_j(K)
 :=\sup_{0\le t\le1}
 \bigl\|{\mathfrak m}\circ D{\cal F}_j(Z_t)\circ{\cal Q}
 \bigr\|_{{\cal X}_\mu^*}.                                   \tag{48.13}
\]
\end{definition}

\begin{corollary}[Scalar control replaces full reset control]
\label{cor:e8v48-scalar-control}
Under Theorem~\ref{thm:e8v48-reset-memory},
\[
 |{\mathfrak m}(\Delta_j(K))|
 \le\chi_j(K)\|{\cal Q}K\|_\mu,                               \tag{48.14}
\]
\[
 \|R_j(K)\|_\mu
 \le w_\mu\chi_j(K)\|{\cal Q}K\|_\mu.                         \tag{48.15}
\]
No norm estimate on the moment-null owner \(C_j\) is needed for
convergence of the scalar beta coefficient.
\end{corollary}

\begin{proof}
Apply \(\mathfrak m\) to (48.10), use (48.13), and then use
(48.12).
\end{proof}

\begin{theorem}[Summable reset-mixing criterion]
\label{thm:e8v48-summable-mixing}
Let \(K_j\) be the exact intermediate one-loop kernels.  If
\[
 \sum_{j=N}^\infty
 \chi_j(K_j)\|{\cal Q}K_j\|_\mu<\infty,                       \tag{48.16}
\]
then the cumulative normalized coefficient drift caused by all
subsequent Wilson resets is absolutely convergent and
\[
 \boxed{
 |\delta b^\downarrow_{[N,\infty)}|
 \le3\sum_{j=N}^\infty
 \chi_j(K_j)\|{\cal Q}K_j\|_\mu.}                             \tag{48.17}
\]
In particular, if
\[
 \chi_j(K_j)\|{\cal Q}K_j\|_\mu\le C_{\rm mix}2^{-j},          \tag{48.18}
\]
then
\[
 |\delta b^\downarrow_{[N,\infty)}|
 \le6C_{\rm mix}2^{-N}.                                      \tag{48.19}
\]
\end{theorem}

\begin{proof}
At step \(j\), the coefficient change is exactly
\(3{\mathfrak m}(\Delta_j)\).  Bound it by (48.14), sum, and use
the geometric series for (48.19).
\end{proof}

\begin{remark}[Coefficient convergence versus full locality]
\label{rem:e8v48-coeff-vs-locality}
The decomposition (48.11) is enough for the scalar coefficient
because \(C_j\in\ker{\mathfrak m}\).  It does not make a large
nonlocal \(C_j\) harmless for the full RG construction.  Uniform
locality and polymer bounds still require control of
\({\cal Q}\Delta_j\); V48 removes that stronger norm only from the
specific coefficient-convergence obligation.
\end{remark}

\subsection{Exact triangularity of the classical Gaussian germ}

The scalar reference map already exhibits a strict version of the
mixing condition.

\begin{proposition}[Regular compliance does not feed the quadratic
precision germ]
\label{prop:e8v48-classical-triangularity}
Let \(c\) be an even periodic compliance with
\[
 c(q\widehat e_1)=\frac{a}{q^2}+r(q),\qquad a>0,              \tag{48.20}
\]
where \(r\) is bounded near zero and \(c\) is bounded near every
nonzero alias \(\pi\sigma\), \(\sigma\in\{0,1\}^3\setminus\{0\}\).
For the dyadic alias map
\[
 ({\cal T}c)(q)
 =\frac14\sum_{\sigma\in\{0,1\}^3}
 c\left(\frac q2\widehat e_1+\pi\sigma\right),                \tag{48.21}
\]
one has
\[
 ({\cal T}c)(q)=\frac a{q^2}+O(1),\qquad
 ({\cal T}c)(q)^{-1}=\frac{q^2}{a}+O(q^4).                    \tag{48.22}
\]
Hence the quadratic precision coefficient \(1/a\) is independent
of all regular and alias data.  The classical
irrelevant-to-marginal derivative is exactly zero.
\end{proposition}

\begin{proof}
The \(\sigma=0\) term of (48.21) is
\[
 \frac14c(q/2)=\frac14\left(\frac{4a}{q^2}+O(1)\right)
 =\frac a{q^2}+O(1).
\]
All seven remaining terms are bounded by hypothesis, proving the
first relation in (48.22).  Write the result as
\((a/q^2)(1+O(q^2))\) and invert to obtain the second.  Since the
coefficient \(1/a\) contains none of the regular data, its
Fréchet derivative in those directions vanishes.
\end{proof}

\begin{corollary}[The Ward residue \(1/3\) is reset stable]
\label{cor:e8v48-one-third-stable}
For \(a=3\), every finite dyadic blocking and every moment
retraction preserving the pole residue has inverse-symbol germ
\[
 s(q)=\frac{q^2}{3}+O(q^4).                                  \tag{48.23}
\]
\end{corollary}

\begin{firewall}[Classical triangularity is not the one-loop Ward
identity]
\label{fire:e8v48-classical-not-loop}
Proposition~\ref{prop:e8v48-classical-triangularity} concerns
inversion of a scalar Gaussian compliance.  The one-loop
determinant integrates over internal edge momenta, so it can mix an
irrelevant deformation into the external \(q^2\) coefficient.
The required quantity is the actual \(\chi_j\) in (48.13), not its
classical analogue.
\end{firewall}

\subsection{The revised finite acquisition}

\begin{definition}[Reset mixing card]
\label{def:e8v48-RMC}
A reset mixing card (RMC) at scale \(j\) contains:
\begin{enumerate}
\item the precise differential of the nonlinear reference
      retraction, including chain-rule coupling terms;
\item the exact intermediate irrelevant kernel
      \({\cal Q}K_j\);
\item the differentiated one-loop shell map
      \(D{\cal F}_j(Z_t)\) on the segment (48.9);
\item a Ward-compatible evaluation or bound for
      \({\mathfrak m}D{\cal F}_j(Z_t){\cal Q}\); and
\item the verified product estimate (48.18), or a summable
      substitute.
\end{enumerate}
\end{definition}

\begin{assessment}[Progress after V48]
\label{ass:e8v48-status}
The same-scale marginal reset is now closed algebraically by
(48.3)--(48.7).  The propagated reset defect has the exact
moment-null plus marginal splitting (48.10)--(48.15).  Therefore
the coefficient problem no longer requires a weighted bound on the
entire edge reset kernel: it requires the scalar mixing product in
(48.16).  Full action locality remains a separate, stronger
obligation.
\end{assessment}

\begin{decision}[V49 target]
\label{dec:e8v48-next}
V49 should differentiate the exact tadpole--bubble formula
(42.12) in an irrelevant reference direction and apply
\(\mathfrak m\) before estimating.  Periodic loop-momentum
summation and the finite-lattice Ward identity must be used first,
because taking absolute values in (45.10) destroys the
cancellations that may make \(\chi_j\) small.  The output should be
an explicit finite-card formula for the mixing scalar, with the
toron mode retained.
\end{decision}

\begin{firewall}[Claim boundary after V48]
\label{fire:e8v48-claim}
V48 proves the canonical marginal projection, the exact
reset-memory formula, and a sharp scalar summability criterion.  It
does not verify the RMC for the actual \(45\)-component shell,
evaluate the limiting beta coefficient, prove its universal value,
close the nonlinear/complex quotient matching card, construct the
full RG trajectory, remove the regulators, establish
Osterwalder--Schrader reconstruction, or prove the Yang--Mills mass
gap.
\end{firewall}

\subsection{Parent record}

This V48 note was formed by appending the preceding material to V47,
whose validated record was
\[
\begin{array}{c|c}
\text{lines}&21591\\
\text{bytes}&753748\\
\text{SHA-256}&
\texttt{492AD8A91F8DD5293B4BD4FE8C0FC29A71998C895F887BCA7BB899A4A24AB11C}.
\end{array}
\]

% =============================================================================
% END APPENDED EIGHTH-SERIES V48 MATERIAL
% =============================================================================

% =============================================================================
% BEGIN APPENDED EIGHTH-SERIES V49 MATERIAL
% =============================================================================

\section{Eighth-series V49: differentiated finite-card formula for
irrelevant-to-marginal mixing}
\label{sec:e8v49}

\subsection{A reference-direction derivative of the quotient card}

Fix a finite \(M^4\) torus and write
\[
 {\mathbb T}_M^{4,*}
 :=\left\{\frac{2\pi}{M}m:
 m\in\{0,\ldots,M-1\}^4\right\}.                              \tag{49.1}
\]
The point \(k=0\) is included.  Let \(\theta\mapsto A_\theta(k)\)
be a \(C^1\) path of positive \(45\times45\) vertical symbols and
put
\[
 G_\theta(k):=A_\theta(k)^{-1},\qquad
 \dot G_\theta(k)
 =-G_\theta(k)\dot A_\theta(k)G_\theta(k).                    \tag{49.2}
\]
All dots in V49 denote the total derivative in the declared
irrelevant reference direction \(\eta=\dot\theta\).

Let \({\cal Q}_{\theta,q}(k)\) and
\({\cal R}_{\theta,q}(k)\) be the compensated quadratic and linear
vertices in the conventions of (42.26), including the harmonic
lift and Haar-density contribution.  Define
\[
\begin{split}
 \Phi_\theta(k,q):={}&
 \Tr\!\left[G_\theta(k){\cal Q}_{\theta,q}(k)\right]\\
 &-\frac13\Tr\!\left[
 G_\theta(k+q\widehat e_1){\cal R}_{\theta,q}(k)
 G_\theta(k){\cal R}_{\theta,q}(k)^*\right],                 \tag{49.3}\\
 \gamma_{M,\theta}(q):={}&
 \frac1{M^4}\sum_{k\in{\mathbb T}_M^{4,*}}\Phi_\theta(k,q).
                                                                    \tag{49.4}
\end{split}
\]

\begin{theorem}[Exact differentiated tadpole--bubble card]
\label{thm:e8v49-differentiated-card}
For every \(k,q\), let
\[
 G:=G_\theta(k),\quad G_+:=G_\theta(k+q\widehat e_1),\quad
 Q:={\cal Q}_{\theta,q}(k),\quad
 R:={\cal R}_{\theta,q}(k).
\]
Then
\[
\boxed{\begin{split}
 \dot\Phi_\theta(k,q)
 ={}&\Tr[\dot G\,Q+G\,\dot Q]\\
 &-\frac13\Tr[
 \dot G_+RGR^*
 +G_+\dot RGR^*
 +G_+R\dot G R^*
 +G_+RG\dot R^*],
\end{split}}                                                  \tag{49.5}
\]
where both inverse derivatives are given by (49.2).  Consequently
\[
 \boxed{\dot\gamma_{M,\theta}(q)
 =\frac1{M^4}\sum_{k\in{\mathbb T}_M^{4,*}}
 \dot\Phi_\theta(k,q).}                                      \tag{49.6}
\]
\end{theorem}

\begin{proof}
Differentiate (49.3) with the product rule.  There are two factors
in the tadpole and four in the bubble.  The finite sum (49.4)
commutes with differentiation.  Formula (49.2) follows by
differentiating \(A_\theta G_\theta=I\).
\end{proof}

\begin{remark}[The vertex dots are total dots]
\label{rem:e8v49-total-dots}
Writing \(L_\theta=-G_\theta B_\theta\) for the harmonic router,
\[
 \dot L_\theta
 =-G_\theta(\dot A_\theta L_\theta+\dot B_\theta).            \tag{49.7}
\]
Thus, if \(b_\theta[a]\) is the lifted retained direction,
\(\dot Q\) and \(\dot R\) in (49.5) include both the explicit
reference-action derivative and the terms
\(D_bQ[\dot b]\), \(D_bR[\dot b]\).  Omitting (49.7) would not
differentiate the quotient determinant.
\end{remark}

\begin{firewall}[The toron summand is not deleted]
\label{fire:e8v49-toron}
The sum (49.6) includes \(k=0\).  The vertical gap makes
\(A_\theta(0)\) invertible; it does not authorize deletion of the
constant retained background sector.  The external \(q=0\)
cancellation is a Ward identity of the complete sum, not a reason
to remove the zero internal momentum by hand.
\end{firewall}

\subsection{Ward extraction of the scalar mixing coefficient}

\begin{lemma}[Differentiated Ward zero]
\label{lem:e8v49-differentiated-Ward}
Suppose the reference path preserves the exact finite-lattice gauge
identity and reflection symmetry.  Then
\[
 \gamma_{M,\theta}(0)=0,\qquad
 \dot\gamma_{M,\theta}(0)=0,\qquad
 \partial_q\dot\gamma_{M,\theta}(0)=0.                        \tag{49.8}
\]
\end{lemma}

\begin{proof}
The first equality is the retained Ward zero (42.18) on the finite
card.  Differentiate it in \(\theta\) to obtain the second.
Reflection sends \(q\) to \(-q\), so
\(\dot\gamma_{M,\theta}\) is even; its first derivative at zero
vanishes.
\end{proof}

\begin{definition}[Finite-card reset-mixing scalar]
\label{def:e8v49-mixing-scalar}
The analytic symbols in (49.5) define, even at finite \(M\),
\[
 \boxed{
 \Xi_{M,\theta}[\eta]
 :=\frac32\left.\partial_q^2\right|_{q=0}
 \frac1{M^4}\sum_{k\in{\mathbb T}_M^{4,*}}
 \dot\Phi_\theta(k,q).}                                      \tag{49.9}
\]
The derivative acts on the trigonometric-polynomial continuation
in external \(q\); it is not a finite difference that drops the
toron.
\end{definition}

\begin{proposition}[Identification with
\(\mathfrak mD{\cal F}{\cal Q}\)]
\label{prop:e8v49-identification}
Let \(\eta={\cal Q}K\) be a moment-null reference direction and let
\({\cal F}_{M,\theta}\) be the finite-card one-loop shell map.
With the normalization of V48,
\[
 \boxed{
 \Xi_{M,\theta}[\eta]
 =3{\mathfrak m}\!\left(
 D{\cal F}_{M,\theta}(K)[\eta]\right).}                        \tag{49.10}
\]
\end{proposition}

\begin{proof}
By Lemma~\ref{lem:e8v49-differentiated-Ward}, the Taylor expansion
of (49.6) begins with
\(\dot\gamma_{M,\theta}(q)
=\frac12\dot\gamma_{M,\theta}''(0)q^2+O(q^4)\).
The moment functional extracts its \(q^2\) coefficient, while the
normalized beta coefficient is three times that coefficient.
Equation (49.9) is therefore exactly (49.10).
\end{proof}

\begin{corollary}[A completed finite-card RMC is a single scalar
identity]
\label{cor:e8v49-RMC-scalar}
For any declared irrelevant basis
\(\eta_1,\ldots,\eta_d\), the marginal mixing row of the reset
mixing card is the finite list
\[
 \Xi_{M,\theta}[\eta_\ell],\qquad1\le\ell\le d,               \tag{49.11}
\]
computed by (49.2), (49.5), and (49.9).  If every value is zero,
the linearized finite-card shell map is block triangular from that
irrelevant space into the marginal Wilson direction.
\end{corollary}

\begin{proof}
Linearity of the derivative and (49.10) reduce the operator on the
span of the basis to its values on the basis vectors.
\end{proof}

\subsection{A cofinal estimator for the mixing scalar}

Let \(q_M=2\pi/M\).  Suppose the irrelevant direction preserves the
quadratic classical moment, so its inverse-symbol derivative
satisfies
\[
 |\dot s_\theta(q)|\le C_{\dot s}|q|^4\qquad(|q|\le1).         \tag{49.12}
\]
Define the entirely finite ratio derivative
\[
\begin{split}
 \widehat\Xi_{M,\theta}[\eta]
 :={}&
 \frac{\dot\gamma_{M,\theta}(q_M)}{s_\theta(q_M)}
 -\frac{\gamma_{M,\theta}(q_M)\dot s_\theta(q_M)}
       {s_\theta(q_M)^2}.                                    \tag{49.13}
\end{split}
\]

\begin{theorem}[Cofinal convergence of differentiated cards]
\label{thm:e8v49-cofinal-mixing}
Assume uniformly on the reference segment:
\begin{enumerate}
\item the differentiated infinite-volume response has an
      exponentially summable kernel \(\dot j_\theta\) with
      \[
       \|\dot j_\theta\|_\mu\le K_{\dot\mu};                   \tag{49.14}
      \]
\item its finite-torus winding error is at most
      \(E_{\dot\mu}e^{-\mu M}\);
\item the undifferentiated bounds (42.15)--(42.16) hold; and
\item (42.29)--(42.30) and (49.12) hold.
\end{enumerate}
Then the absolutely convergent infinite-volume mixing scalar is
\[
 \boxed{
 \Xi_{\infty,\theta}[\eta]
 =-\frac32\sum_{x\in\mathbb Z^4}x_1^2\dot j_\theta(x),}        \tag{49.15}
\]
and there are explicit constants \(C_{\rm Tay},C_{\rm wind}\),
formed only from the displayed stocks, such that
\[
 \boxed{
 |\widehat\Xi_{M,\theta}[\eta]
       -\Xi_{\infty,\theta}[\eta]|
 \le C_{\rm Tay}q_M^2
 +C_{\rm wind}(1+q_M^{-2})e^{-\mu M}.}                        \tag{49.16}
\]
In particular,
\[
 \widehat\Xi_{M,\theta}[\eta]
 \longrightarrow\Xi_{\infty,\theta}[\eta].                   \tag{49.17}
\]
\end{theorem}

\begin{proof}
Exponential summability and the differentiated Ward zero give,
by the proof of Theorem~\ref{thm:e8v42-moment},
\[
 \dot\gamma_{\infty,\theta}(q)
 =\dot\kappa_\theta q^2+O(K_{\dot\mu}\mu^{-4}q^4),\qquad
 \dot\kappa_\theta
 =-\frac12\sum_xx_1^2\dot j_\theta(x),                        \tag{49.18}
\]
with
\(|\dot\kappa_\theta|\le K_{\dot\mu}\mu^{-2}\).
The first quotient in (49.13) differs from
\(3\dot\kappa_\theta\) by \(O(q_M^2)\), exactly as in
(42.34), with \(K_\mu\) replaced by \(K_{\dot\mu}\).

For the second quotient, the undifferentiated Taylor bound gives
\(|\gamma_{\infty,\theta}(q)|\le C_\gamma q^2\) for
\(|q|\le1\).  Equations (49.12) and
\(s_\theta(q)\ge m_sq^2\) imply
\[
 \left|
 \frac{\gamma_{\infty,\theta}(q)\dot s_\theta(q)}
      {s_\theta(q)^2}\right|
 \le\frac{C_\gamma C_{\dot s}}{m_s^2}q^2.                    \tag{49.19}
\]
The two winding errors contribute at worst a constant multiple of
\((1+q_M^{-2})e^{-\mu M}\), again using the same lower bound on
\(s_\theta\).  Combining these estimates proves (49.16).
Finally \(3\dot\kappa_\theta\) is (49.15), and both terms on the
right of (49.16) vanish.
\end{proof}

\begin{firewall}[Absolute values come after the Ward sum]
\label{fire:e8v49-sum-first}
The individual six terms in (49.5) need not have vanishing
constant or linear external-momentum parts.  Those pieces cancel
only after the trace, the complete periodic \(k\)-sum, the
Haar-compensated vertex terms, and the toron summand are combined.
Termwise absolute-value bounds such as (45.10) are valid for
locality, but they cannot decide whether the scalar (49.9)
vanishes or contracts.
\end{firewall}

\subsection{What remains to evaluate}

\begin{definition}[Differentiated quotient card]
\label{def:e8v49-DQC}
For each irrelevant direction, the differentiated quotient card
(DQC) records:
\begin{enumerate}
\item \(\dot A,\dot B\) and the total router derivative (49.7);
\item the total compensated vertices \(\dot Q,\dot R\);
\item the six trace owners in (49.5);
\item the exact Ward checks (49.8), including \(k=0\);
\item the scalar (49.9) and, on cofinal tori, the estimator
      (49.13) with the error certificate (49.16).
\end{enumerate}
\end{definition}

\begin{assessment}[Progress after V49]
\label{ass:e8v49-status}
The previously abstract mixing operator
\({\mathfrak m}D{\cal F}{\cal Q}\) is now an explicit finite
matrix card.  No infinite-volume inversion is needed to test a
finite irrelevant basis, and Theorem
\ref{thm:e8v49-cofinal-mixing} gives the route from those cards to
the infinite-volume scalar.  The actual dotted \(45\times45\)
symbols for the chosen reset complement have not yet been entered,
so the size or vanishing of \(\Xi\) remains open.
\end{assessment}

\begin{decision}[Mandatory V50 audit and calculation target]
\label{dec:e8v49-next}
V50 must first inspect the current Standard Model and
Navier--Stokes companion notes, as required every five versions.
It should then populate the DQC as far as the existing exact
incidence formulas permit.  The preferred order is: prove the
summed Ward cancellations in (49.8), isolate any total periodic
\(k\)-derivatives, and only then bound the residual scalar.  If the
one-dimensional Wilson reset is not stable, the reference family
must be enlarged by the smallest irrelevant directions detected by
(49.11).
\end{decision}

\begin{firewall}[Claim boundary after V49]
\label{fire:e8v49-claim}
V49 proves the exact differentiated tadpole--bubble formula,
identifies its moment with the reset-mixing scalar, and supplies a
cofinal estimator theorem.  It does not evaluate the actual DQC,
prove reset mixing is summable, complete the LERB, establish the
universal beta coefficient, close the nonlinear/complex quotient
matching card, construct the full RG trajectory, remove the
regulators, establish Osterwalder--Schrader reconstruction, or
prove the Yang--Mills mass gap.
\end{firewall}

\subsection{Parent record}

This V49 note was formed by appending the preceding material to V48,
whose validated record was
\[
\begin{array}{c|c}
\text{lines}&21933\\
\text{bytes}&765316\\
\text{SHA-256}&
\texttt{DA35DF2537EEB3FEDFD3554A048828E10F36537ABE5E80E9FA3B4BA6558891A6}.
\end{array}
\]

% =============================================================================
% END APPENDED EIGHTH-SERIES V49 MATERIAL
% =============================================================================

% =============================================================================
% BEGIN APPENDED EIGHTH-SERIES V50 MATERIAL
% =============================================================================

\section{Eighth-series V50: compulsory companion audit,
Ward--hypercubic quartic reduction, and the actual reset subspace}
\label{sec:e8v50}

\subsection{Compulsory companion audit}

The two current companion records in \texttt{latex\_notes} were
read before choosing the V50 direction:
\[
\begin{array}{c|r|r|l}
\text{record}&\text{lines}&\text{bytes}&\text{SHA-256}\\ \hline
\texttt{Renewal\_SM\_Einstein\_Continuum\_Research\_Notes\_V41.tex}
&15398&588848&
\texttt{C04DC27A93EF8E78FF9106285334F4F4C442A4D271896089AF1F3E596CD4401E}\\
\texttt{Renewal\_NS\_Stability\_Research\_Notes\_V26.tex}
&5320&278283&
\texttt{82C887F8C2C69E4F28FAD7C3DC8DE5B9099CB5A4BF431EBA1FFF3E91935978AD}.
\end{array}                                                    \tag{50.1}
\]

\begin{proposition}[Audit transfer decision]
\label{prop:e8v50-audit-transfer}
The SM record transfers the following structural conclusions:
\begin{enumerate}
\item raw gauge insertions have the wrong scale power; the complete
      gauge-invariant trace and its Ward cancellation must be formed
      before estimating;
\item contractible local owners and winding toron owners require
      separate ledgers;
\item every noncontracting local invariant must be projected before
      the complement is called irrelevant; and
\item Ward symmetry plus locality gives an \(O(p^2)\) response, but
      a quantitative covariant local expansion is still needed for
      higher-order contraction.
\end{enumerate}
The NS record supplies no equation or numerical constant for the
Yang--Mills problem.  Its rank-one refinement does, however, expose
a valid norm principle: an operator should be estimated on the
actual structured input set, not automatically on its whole ambient
space.  Applied here, this replaces the full
\(\ker{\mathfrak m}\) norm by the actual reset-image norm below.
\end{proposition}

\begin{proof}
The four SM conclusions are respectively the content of its
V41 insertion no-go and Ward-relative theorem, its local/toron
separation, its V39 scale-step projection theorem, and its stated
covariant-Poincare acquisition.  They concern the same lattice
gauge response structure and therefore transfer at the level of
their stated hypotheses.  The NS V26 estimates act on rank-one
Oseen-kernel inputs in a fluid equation and have no gauge,
determinant, or quotient content, so their constants do not
transfer.  Restricting an operator norm to a subset of its domain is
a general functional-analytic operation and needs no fluid
hypothesis.
\end{proof}

\begin{corollary}[Pure-Yang--Mills marginal census]
\label{cor:e8v50-YM-census}
For a reflection-even, gauge-invariant pure Yang--Mills local
action in four dimensions, the quadratic dimension-four
contractible row is the kinetic \(F_{\mu\nu}^2\) row.  The
parity-odd \(F\wedge F\) row is separately projected if reflection
is not imposed, and winding holonomies remain in the toron ledger.
Thus the projection (48.3) is the complete reflection-even
quadratic marginal projection, subject to the symmetry and locality
hypotheses made explicit next.
\end{corollary}

\begin{firewall}[The audit does not import an unproved lemma]
\label{fire:e8v50-audit-firewall}
The SM companion explicitly leaves its quantitative lattice
covariant Poincare lemma open.  V50 therefore does not cite that
missing lemma as a proof of Yang--Mills irrelevant contraction.
It proves only the quadratic Taylor classification obtainable from
Ward, reflection, hypercubic symmetry, and analyticity.
\end{firewall}

\subsection{Classification of the quadratic Ward tensor}

Let \(K_{\mu\nu}(p)\) be the Fourier symbol of a contractible,
translation-invariant quadratic owner on \(\mathbb Z^4\).  Assume
that it is real analytic near zero, real symmetric and even, is
covariant under the hypercubic group, and obeys the exact lattice
Ward identity
\[
 \sum_\mu\widehat p_\mu^*K_{\mu\nu}(p)=0,\qquad
 \widehat p_\mu=e^{ip_\mu}-1.                                \tag{50.2}
\]
Contractibility excludes a winding owner and gives \(K(0)=0\).

\begin{theorem}[Unique hypercubic transverse quadratic germ]
\label{thm:e8v50-quadratic-classification}
Under the preceding hypotheses there is a unique scalar \(a(K)\)
such that
\[
 K_{\mu\nu}(p)
 =a(K)\bigl(|p|^2\delta_{\mu\nu}-p_\mu p_\nu\bigr)
 +O(|p|^4).                                                   \tag{50.3}
\]
For the polarization used in (48.1),
\[
 a(K)={\mathfrak m}(K).                                      \tag{50.4}
\]
\end{theorem}

\begin{proof}
Evenness and \(K(0)=0\) make the first possible Taylor term
quadratic.  Reflections and coordinate permutations imply that its
diagonal and off-diagonal entries have the form
\[
 K_{\mu\mu}^{(2)}(p)=a_0|p|^2+b_0p_\mu^2,\qquad
 K_{\mu\nu}^{(2)}(p)=c_0p_\mu p_\nu\quad(\mu\ne\nu).           \tag{50.5}
\]
Since
\(\widehat p_\mu=ip_\mu+O(|p|^2)\), the lowest nonzero part of
(50.2), for fixed \(\nu\), is
\[
 p_\nu\bigl[(a_0+c_0)|p|^2+(b_0-c_0)p_\nu^2\bigr]=0.          \tag{50.6}
\]
As a polynomial identity this gives
\(c_0=-a_0\) and \(b_0=c_0=-a_0\).  Substitution in (50.5)
proves (50.3) with \(a(K)=a_0\).  On
\(p=q\widehat e_1\), its \(22\) entry is
\(a_0q^2+O(q^4)\).  The real-space Taylor identity used in
(42.19) says that the coefficient of \(q^2\) is
\(-\frac12\sum_xx_1^2K_{22}(x)={\mathfrak m}(K)\), proving
(50.4).  The next term is fourth order because the symbol is even.
\end{proof}

\begin{corollary}[Moment subtraction forces a quartic germ]
\label{cor:e8v50-quartic-after-projection}
For every \(K\) satisfying the theorem,
\[
 {\cal Q}K=O(|p|^4).                                         \tag{50.7}
\]
More precisely, for some \(P,C_4>0\),
\[
 \|{\cal Q}K\|_{{\cal G},P}
 :=\sup_{0<|p|\le P}
 \frac{\|({\cal Q}K)(p)\|}{|p|^4}
 \le C_4.                                                     \tag{50.8}
\]
\end{corollary}

\begin{proof}
The normalized Wilson Hessian \({\cal W}\) has the same quadratic
germ as (50.3) with coefficient one.  By (48.3)--(48.4) and
(50.4), subtraction of \({\cal P}K\) cancels the entire quadratic
Taylor tensor.  Analyticity and evenness give the fourth-order
Taylor remainder and its bounded quotient on a sufficiently small
closed ball.
\end{proof}

\begin{firewall}[Hypercubic symmetry is doing real work]
\label{fire:e8v50-hypercubic}
Ward transversality without the full coordinate-reflection and
permutation symmetries can leave several independent transverse
quadratic tensors.  A one-dimensional moment projection would then
be incomplete.  The symmetry rows must be verified for the actual
regulator and restored before applying
Corollary~\ref{cor:e8v50-quartic-after-projection}.
\end{firewall}

\subsection{Canonical contraction of the quartic complement}

For a two-derivative quadratic kernel define the canonical
rescaling
\[
 ({\cal S}_\ell K)(p):=\ell^2K(p/\ell),\qquad\ell\ge1.         \tag{50.9}
\]

\begin{lemma}[Quartic germ contracts relative to the kinetic row]
\label{lem:e8v50-quartic-rescaling}
If \(K(0)=D K(0)=D^2K(0)=D^3K(0)=0\), then
\[
 \|{\cal S}_\ell K\|_{{\cal G},P}
 \le\ell^{-2}\|K\|_{{\cal G},P}                              \tag{50.10}
\]
whenever \(P/\ell\le P\).
\end{lemma}

\begin{proof}
For \(0<|p|\le P\), set \(r=p/\ell\).  Then
\[
 \frac{\|\ell^2K(p/\ell)\|}{|p|^4}
 =\ell^{-2}\frac{\|K(r)\|}{|r|^4}
 \le\ell^{-2}\|K\|_{{\cal G},P}.
\]
\end{proof}

\begin{proposition}[Forced quartic recurrence]
\label{prop:e8v50-quartic-recurrence}
Suppose the reset complement stocks \(u_j\ge0\) satisfy
\[
 u_{j+1}\le\frac14u_j+c\,2^{-j}.                              \tag{50.11}
\]
Then, for \(j\ge N\),
\[
 u_j\le4^{-(j-N)}u_N+4c\,2^{-j},                             \tag{50.12}
\]
and \(\sum_{j=N}^\infty u_j<\infty\).
\end{proposition}

\begin{proof}
Iterating (50.11) gives
\[
 u_j\le4^{-(j-N)}u_N
 +c\sum_{r=N}^{j-1}4^{-(j-1-r)}2^{-r}.
\]
With \(s=j-1-r\), the last sum is at most
\(2^{-(j-1)}\sum_{s\ge0}2^{-s}=4\,2^{-j}\).
Both terms are summable in \(j\).
\end{proof}

\begin{remark}[Relation to V46]
\label{rem:e8v50-relation-V46}
The exact scalar reference in V46 has a generated quartic stock of
order \(2^{-j}\), rather than a fixed stock transported only by
(50.10).  It therefore fits the forcing term in (50.11), and its
proved estimate (46.7) is consistent with (50.12).
\end{remark}

\subsection{Restricting the mixing norm to the actual reset image}

Let \({\cal Y}_j\subset\ker{\mathfrak m}\) be a closed subspace
known to contain every discarded tangent produced by the declared
reference retraction at scale \(j\).  Define
\[
 \chi_j^{\rm act}
 :=\sup_{0\le t\le1}
 \left\|
 {\mathfrak m}\circ D{\cal F}_j(Z_t)
 \big|_{{\cal Y}_j}
 \right\|_{{\cal Y}_j^*}.                                    \tag{50.13}
\]

\begin{proposition}[Actual-image reset bound]
\label{prop:e8v50-actual-image}
If \({\cal Q}K_j\in{\cal Y}_j\), then
\[
 |{\mathfrak m}(\Delta_j(K_j))|
 \le\chi_j^{\rm act}\|{\cal Q}K_j\|_{{\cal Y}_j}.             \tag{50.14}
\]
Moreover \(\chi_j^{\rm act}\le\chi_j\) whenever the
\({\cal Y}_j\)-norm is the restriction of \(\|\cdot\|_\mu\).
\end{proposition}

\begin{proof}
Repeat the proof of (48.14), observing that the vector inserted in
(48.10) belongs to \({\cal Y}_j\).  A supremum over a subspace is
no larger than the supremum over the ambient unit ball.
\end{proof}

\begin{theorem}[Ward-compiled reset criterion]
\label{thm:e8v50-Ward-compiled-reset}
Assume at every \(j\ge N\):
\begin{enumerate}
\item the discarded owner is contractible, reflection even,
      hypercubic, exponentially local, and satisfies the exact Ward
      identity;
\item its moment is removed by \({\cal Q}\);
\item its quartic stock \(u_j=\|{\cal Q}K_j\|_{{\cal G},P}\)
      satisfies (50.11); and
\item the complete differentiated quotient card, after its Ward
      sum, satisfies the response-level bound
      \[
       |{\mathfrak m}D{\cal F}_j(Z)[\eta]|
       \le C_{\rm resp}\|\eta\|_{{\cal G},P}
       \quad(\eta\in{\cal Y}_j)                               \tag{50.15}
      \]
      uniformly on the reset segment.
\end{enumerate}
Then reset-induced marginal drift is absolutely summable:
\[
 \sum_{j=N}^\infty
 |{\mathfrak m}(\Delta_j(K_j))|
 \le C_{\rm resp}
 \left(\frac43u_N+8c\,2^{-N}\right).                          \tag{50.16}
\]
\end{theorem}

\begin{proof}
Items 1--2 and
Corollary~\ref{cor:e8v50-quartic-after-projection} place the actual
discarded tangent in the quartic germ space.  Equations
(48.10) and (50.15) give
\(|{\mathfrak m}(\Delta_j)|\le C_{\rm resp}u_j\).
Sum (50.12).  The first geometric sum is
\(\sum_{j=N}^\infty4^{-(j-N)}=4/3\); the second is
\(\sum_{j=N}^\infty4c2^{-j}=8c2^{-N}\).
\end{proof}

\begin{firewall}[The remaining edge theorem is (50.15)]
\label{fire:e8v50-response-open}
Theorem~\ref{thm:e8v50-Ward-compiled-reset} does not prove
(50.15).  That bound must be established on the full summed
tadpole--bubble expression (49.5), not on its six insertions
separately.  In particular, a quartic external germ does not by
itself control how an internal Brillouin-edge deformation enters
the loop.
\end{firewall}

\subsection{Status and the post-audit direction}

\begin{assessment}[V50 audit outcome]
\label{ass:e8v50-status}
The audit confirms the response-level route and sharpens the reset
problem in two ways.  First, the symmetry classification
(50.3)--(50.8) proves that exact moment subtraction removes the
entire quadratic local germ, leaving a canonically contracting
quartic complement.  Second, the DQC needs to bound the derivative
only on the actual reset space \({\cal Y}_j\).  The unresolved
analytic row is the uniform Ward-summed response bound (50.15).
\end{assessment}

\begin{decision}[V51 target]
\label{dec:e8v50-next}
V51 should construct the smallest explicit \({\cal Y}_j\) generated
by one dyadic Wilson reset.  Its basis should be obtained from the
fourth-order Taylor tensors of the exact blocked quadratic symbol,
with hypercubic and Ward relations imposed before enumeration.
The differentiated cards (49.5) are then required only on that
finite basis.  If nonlocal edge components remain outside this
basis, they must be retained as a separate analytic tail rather
than silently included in a Taylor polynomial.
\end{decision}

\begin{firewall}[Claim boundary after V50]
\label{fire:e8v50-claim}
V50 performs the compulsory companion audit, proves the unique
hypercubic transverse quadratic germ, proves quartic contraction
after moment subtraction, and restricts the reset norm to its
actual input space.  It does not prove (50.15), evaluate the DQC,
complete the LERB, establish the universal beta coefficient, close
the nonlinear/complex quotient matching card, construct the full
RG trajectory, remove the regulators, establish
Osterwalder--Schrader reconstruction, or prove the Yang--Mills mass
gap.
\end{firewall}

\subsection{Parent record}

This V50 note was formed by appending the preceding material to V49,
whose validated record was
\[
\begin{array}{c|c}
\text{lines}&22282\\
\text{bytes}&777537\\
\text{SHA-256}&
\texttt{75534557D106F85F6223381047680A1160BAD1CF412F00F9AD56521CCB7A0E9E}.
\end{array}
\]

% =============================================================================
% END APPENDED EIGHTH-SERIES V50 MATERIAL
% =============================================================================

% =============================================================================
% BEGIN APPENDED EIGHTH-SERIES V51 MATERIAL
% =============================================================================

\section{Eighth-series V51: the two-dimensional quartic Ward space
and a minimal local reset basis}
\label{sec:e8v51}

\subsection{Centered-link Ward coordinates}

Put
\[
 \bar p_\mu:=2\sin(p_\mu/2),\qquad
 x_\mu:=\bar p_\mu^2,\qquad
 S:=\sum_{\rho=1}^4x_\rho,\qquad
 T:=\sum_{\rho=1}^4x_\rho^2.                                 \tag{51.1}
\]
Since
\(\widehat p_\mu=e^{ip_\mu}-1
=i e^{ip_\mu/2}\bar p_\mu\), conjugation by the diagonal unitary
\[
 U(p)_{\mu\nu}:=e^{ip_\mu/2}\delta_{\mu\nu}                   \tag{51.2}
\]
turns the lattice Ward vector into \(i\bar p\).  Thus for
\[
 K^c(p):=U(p)^*K(p)U(p),                                     \tag{51.3}
\]
the two Ward identities read
\[
 K^c(p)\bar p=0,\qquad \bar p^{\,T}K^c(p)=0.                 \tag{51.4}
\]
The superscript \(c\) denotes centered-link coordinates, not a new
regulator.

\subsection{Two exact lattice-transverse quartic tensors}

Define symmetric matrix-valued trigonometric polynomials
\[
 {\cal W}^{(1)}_{\mu\nu}(p)
 :=S\bigl(S\delta_{\mu\nu}-\bar p_\mu\bar p_\nu\bigr),         \tag{51.5}
\]
and
\[
 {\cal W}^{(2)}_{\mu\nu}(p)
 :=\begin{cases}
 T+Sx_\mu-2x_\mu^2,&\mu=\nu,\\
 -\bar p_\mu\bar p_\nu(x_\mu+x_\nu),&\mu\ne\nu.
 \end{cases}                                                  \tag{51.6}
\]

\begin{lemma}[Exact symmetry and Ward properties]
\label{lem:e8v51-two-tensors}
Each \({\cal W}^{(r)}\), \(r=1,2\), is real, symmetric, even,
hypercubic covariant, finite range in centered-link coordinates,
and satisfies
\[
 {\cal W}^{(r)}(p)\bar p=0.                                  \tag{51.7}
\]
Their homogeneous Taylor terms of degree four are linearly
independent.
\end{lemma}

\begin{proof}
All symmetry statements are immediate from (51.1), (51.5), and
(51.6); trigonometric polynomial entries have finite Fourier
support.  For (51.5),
\[
 \sum_\mu\bar p_\mu{\cal W}^{(1)}_{\mu\nu}
 =S(S\bar p_\nu-S\bar p_\nu)=0.
\]
For (51.6), fix \(\nu\) and factor \(\bar p_\nu\).  The remaining
polynomial is
\[
 T+Sx_\nu-2x_\nu^2
 -\sum_{\mu\ne\nu}x_\mu(x_\mu+x_\nu)
 =T+Sx_\nu-2x_\nu^2
 -(T-x_\nu^2)-x_\nu(S-x_\nu)=0.
\]
At small \(p\), (51.5) has off-diagonal entry
\(-p_\mu p_\nu|p|^2\), whereas (51.6) has
\(-p_\mu p_\nu(p_\mu^2+p_\nu^2)\).  These are not scalar
multiples in four variables, proving independence.
\end{proof}

\begin{theorem}[Classification of quartic hypercubic Ward germs]
\label{thm:e8v51-quartic-classification}
Let \(H_{\mu\nu}(p)\) be a real symmetric homogeneous polynomial of
degree four, even under \(p\mapsto-p\), hypercubic covariant, and
transverse:
\[
 \sum_\mu p_\mu H_{\mu\nu}(p)=0.                              \tag{51.8}
\]
Then there are unique \(c_1,c_2\in\mathbb R\) such that
\[
 \boxed{
 H=c_1{\cal W}^{(1)}_{[4]}+c_2{\cal W}^{(2)}_{[4]},}           \tag{51.9}
\]
where the subscript denotes the homogeneous degree-four Taylor
part.  Hence the quartic hypercubic Ward space has dimension two.
\end{theorem}

\begin{proof}
Sign reflections and permutations preserving a chosen index force
the most general diagonal and off-diagonal entries to be
\[
\begin{split}
 H_{\mu\mu}
 &=A|p|^4+B\sum_\rho p_\rho^4
   +C|p|^2p_\mu^2+Dp_\mu^4,\\
 H_{\mu\nu}
 &=p_\mu p_\nu\bigl(E|p|^2+F(p_\mu^2+p_\nu^2)\bigr),
 \qquad\mu\ne\nu.                                            \tag{51.10}
\end{split}
\]
Set \(y_\rho=p_\rho^2\), \(S_0=\sum_\rho y_\rho\), and
\(T_0=\sum_\rho y_\rho^2\).  For fixed \(\nu\), division of
\(\sum_\mu p_\mu H_{\mu\nu}\) by \(p_\nu\) gives
\[
 (A+E)S_0^2+(B+F)T_0
 +(C-E+F)S_0y_\nu+(D-2F)y_\nu^2.                             \tag{51.11}
\]
The polynomial vanishes identically, so
\[
 E=-A,\qquad F=-B,\qquad C=-A+B,\qquad D=-2B.                \tag{51.12}
\]
Thus only \(A,B\) are free.  Taking \((A,B)=(1,0)\) gives the
degree-four part of (51.5), and taking \((A,B)=(0,1)\) gives that
of (51.6).  Their independence was proved in
Lemma~\ref{lem:e8v51-two-tensors}, so the representation is unique.
\end{proof}

\subsection{Exact extraction of the two local coefficients}

For an analytic centered symbol \(K^c=O(|p|^4)\), define
\[
 \lambda_{\rm ax}(K)
 :=\lim_{q\to0}
 \frac{K^c_{22}(q\widehat e_1)}{(2\sin(q/2))^4},              \tag{51.13}
\]
\[
 \lambda_{\rm dg}(K)
 :=\lim_{q\to0}
 \frac{K^c_{33}(q(\widehat e_1+\widehat e_2))}
      {(2\sin(q/2))^4}.                                      \tag{51.14}
\]

\begin{lemma}[Coefficient table]
\label{lem:e8v51-coefficient-table}
The two basis tensors have
\[
\begin{array}{c|cc}
 &\lambda_{\rm ax}&\lambda_{\rm dg}\\ \hline
 {\cal W}^{(1)}&1&4\\
 {\cal W}^{(2)}&1&2.
\end{array}                                                    \tag{51.15}
\]
Consequently the coefficients of the quartic Taylor germ of \(K\)
are
\[
 \boxed{
 c_1(K)=\frac{\lambda_{\rm dg}(K)-2\lambda_{\rm ax}(K)}2,
 \qquad
 c_2(K)=\frac{4\lambda_{\rm ax}(K)-\lambda_{\rm dg}(K)}2.}     \tag{51.16}
\]
\end{lemma}

\begin{proof}
On \(p=q\widehat e_1\), one has \(S=x_1\), \(T=x_1^2\), and both
\((22)\) entries equal \(x_1^2\).  On
\(p=q(\widehat e_1+\widehat e_2)\), one has
\(S=2x_1\), \(T=2x_1^2\), and the \((33)\) entries of (51.5) and
(51.6) are respectively \(4x_1^2\) and \(2x_1^2\).
This proves (51.15).  Solve
\(\lambda_{\rm ax}=c_1+c_2\) and
\(\lambda_{\rm dg}=4c_1+2c_2\) to obtain (51.16).
\end{proof}

\begin{definition}[Quartic local retraction]
\label{def:e8v51-quartic-retraction}
On the moment-null analytic Ward space define
\[
 {\cal P}_4K
 :=c_1(K){\cal W}^{(1)}+c_2(K){\cal W}^{(2)},\qquad
 {\cal Q}_6:=I-{\cal P}_4.                                   \tag{51.17}
\]
\end{definition}

\begin{proposition}[Projector and sixth-order remainder]
\label{prop:e8v51-sixth-remainder}
The map \({\cal P}_4\) is a projection onto
\[
 {\cal Y}_{\rm loc}^{(4)}
 :=\operatorname{span}\{{\cal W}^{(1)},{\cal W}^{(2)}\},       \tag{51.18}
\]
and for every analytic, even, hypercubic, transverse
\(K=O(|p|^4)\),
\[
 ({\cal Q}_6K)(p)=O(|p|^6).                                  \tag{51.19}
\]
\end{proposition}

\begin{proof}
The coefficient table (51.15) shows
\({\cal P}_4{\cal W}^{(r)}={\cal W}^{(r)}\); hence
\({\cal P}_4^2={\cal P}_4\).  Theorem
\ref{thm:e8v51-quartic-classification} says that subtracting
\({\cal P}_4K\) cancels the complete degree-four Taylor tensor.
Evenness excludes degree five, proving (51.19).
\end{proof}

\begin{firewall}[The exact basis is a germ-matching choice]
\label{fire:e8v51-germ-basis}
The tensors (51.5)--(51.6) are convenient exact
lattice-transverse completions of the two quartic Taylor tensors.
Other finite-range completions with the same fourth-order germs
differ by \(O(|p|^6)\).  Such a scheme change moves the
\({\cal Q}_6\) tail but cannot change the two-dimensional
classification.
\end{firewall}

\subsection{Improved scaling of the tail}

For \(F=O(|p|^6)\), set
\[
 \|F\|_{{\cal H},P}
 :=\sup_{0<|p|\le P}\frac{\|F(p)\|}{|p|^6}.                   \tag{51.20}
\]

\begin{lemma}[Sixth-order canonical contraction]
\label{lem:e8v51-sixth-contraction}
For the rescaling (50.9),
\[
 \|{\cal S}_\ell F\|_{{\cal H},P}
 \le\ell^{-4}\|F\|_{{\cal H},P}.                             \tag{51.21}
\]
\end{lemma}

\begin{proof}
With \(r=p/\ell\),
\[
 \frac{\|\ell^2F(p/\ell)\|}{|p|^6}
 =\ell^{-4}\frac{\|F(r)\|}{|r|^6}.
\]
Take the supremum.
\end{proof}

\begin{corollary}[Dyadic local reset ledger]
\label{cor:e8v51-dyadic-ledger}
After the marginal projection \({\cal P}\), one dyadic reset splits
the local quadratic complement exactly as
\[
 {\cal Q}K
 =c_1{\cal W}^{(1)}+c_2{\cal W}^{(2)}
  +{\cal Q}_6{\cal Q}K.                                     \tag{51.22}
\]
The first two coordinates have relative canonical factor \(1/4\)
per scale, while the sixth-order tail has factor \(1/16\), before
new shell forcing is added.
\end{corollary}

\subsection{Reduction of the differentiated quotient card}

\begin{theorem}[Two local cards plus one tail]
\label{thm:e8v51-two-cards-tail}
Let \(\Xi_j\) be the linear mixing functional (49.9) at scale \(j\).
For every local analytic reset complement,
\[
\boxed{
 \Xi_j[{\cal Q}K]
 =c_1(K)\Xi_j[{\cal W}^{(1)}]
  +c_2(K)\Xi_j[{\cal W}^{(2)}]
  +\Xi_j[{\cal Q}_6{\cal Q}K].}                              \tag{51.23}
\]
Thus the finite local fourth-order part of the DQC requires exactly
two matrix evaluations.  The last term is a separately declared
sixth-order or nonlocal-tail obligation.
\end{theorem}

\begin{proof}
Apply the linear functional \(\Xi_j\) to (51.22).
\end{proof}

\begin{corollary}[Sufficient restricted response stocks]
\label{cor:e8v51-response-stocks}
If
\[
 |\Xi_j[{\cal W}^{(r)}]|\le X_r\quad(r=1,2),\qquad
 |\Xi_j[F]|\le X_6\|F\|_{{\cal H},P},                         \tag{51.24}
\]
then
\[
 |\Xi_j[{\cal Q}K]|
 \le X_1|c_1(K)|+X_2|c_2(K)|
     +X_6\|{\cal Q}_6{\cal Q}K\|_{{\cal H},P}.                \tag{51.25}
\]
\end{corollary}

\begin{proof}
Use (51.23) and the triangle inequality after, not before, the
complete Ward-summed values \(\Xi_j\) have been formed.
\end{proof}

\begin{firewall}[Quadratic closure is not full action closure]
\label{fire:e8v51-quadratic-only}
The two tensors classify the quadratic fourth-order gauge response.
Dimension-six gauge invariants such as cubic curvature terms have
zero quadratic Hessian at the vacuum but contribute to higher
vertices.  They must be included in the nonlinear RG/polymer
ledger and are not eliminated by (51.22).
\end{firewall}

\subsection{Status and the V52 acquisition}

\begin{assessment}[Progress after V51]
\label{ass:e8v51-status}
The vague finite irrelevant basis requested in V50 is now explicit
and minimal in the local quadratic sector.  Hypercubic Ward
symmetry leaves exactly the tensors (51.5)--(51.6); their
coefficients are extracted by two low-momentum polarizations, and
the remaining analytic tail begins at sixth order.  The uncomputed
objects are the two complete finite-card responses
\(\Xi_j[{\cal W}^{(1)}]\), \(\Xi_j[{\cal W}^{(2)}]\), and a
uniform response bound on the tail.
\end{assessment}

\begin{decision}[V52 target]
\label{dec:e8v51-next}
V52 should derive a basis-independent formula for the two local
mixing numbers by differentiating the vertical Hessian with respect
to the added quadratic operators (51.5)--(51.6).  Because these
operators are quadratic, their direct cubic and quartic background
vertices must be separated from the induced router derivative.
The toron and Haar rows remain in the common card.  If the exact
local incidence data are insufficient to evaluate a number, V52
must state the missing matrix entry rather than substitute a
continuum value.
\end{decision}

\begin{firewall}[Claim boundary after V51]
\label{fire:e8v51-claim}
V51 proves the two-dimensional quartic Ward classification, an
explicit lattice-transverse basis, its coefficient extraction, and
the sixth-order tail split.  It does not evaluate either local DQC
number, bound the tail response, complete the LERB, establish the
universal beta coefficient, close the nonlinear/complex quotient
matching card, construct the full RG trajectory, remove the
regulators, establish Osterwalder--Schrader reconstruction, or
prove the Yang--Mills mass gap.
\end{firewall}

\subsection{Parent record}

This V51 note was formed by appending the preceding material to V50,
whose validated record was
\[
\begin{array}{c|c}
\text{lines}&22644\\
\text{bytes}&791044\\
\text{SHA-256}&
\texttt{ED4B9D43BBAED5C8CC4D7DAC34E30AC4E25974C115D67E7D9E7BFA779986DA60}.
\end{array}
\]

% =============================================================================
% END APPENDED EIGHTH-SERIES V51 MATERIAL
% =============================================================================

% =============================================================================
% BEGIN APPENDED EIGHTH-SERIES V52 MATERIAL
% =============================================================================

\section{Eighth-series V52: nonlinear completion of the quartic
reset basis and the stationary mixed-jet card}
\label{sec:e8v52}

\subsection{Why a quadratic basis does not determine a DQC}

The quantities \(\dot Q,\dot R\) in (49.5) are background jets of
the vertical Hessian.  They are not determined by the flat
quadratic Hessian alone.

\begin{proposition}[Nonuniqueness of nonlinear completion]
\label{prop:e8v52-nonlinear-nonuniqueness}
There exist reflection-even, gauge-invariant, finite-range lattice
actions \(V\) and \(\widetilde V\) with identical value, first
derivative, and Hessian at the flat configuration, but with
different fourth derivatives there.  Consequently they have the
same quadratic kernel and can have different one-loop tadpole
vertices.
\end{proposition}

\begin{proof}
Let \(P_{\mu\nu}(x)\) be an oriented plaquette holonomy and set
\[
 Z(U):=\sum_{x,\mu<\nu}
 \left(\Re\Tr\bigl(I-P_{\mu\nu}(x)\bigr)\right)^2.             \tag{52.1}
\]
The trace is invariant under conjugation at the plaquette base
point, so \(Z\) is gauge invariant; summing unoriented plaquettes
makes it reflection even and finite range.  In exponential
coordinates \(P=e^F\),
\[
 \Re\Tr(I-e^F)=-\frac12\Re\Tr(F^2)+O(\|F\|^4),                \tag{52.2}
\]
because \(F\) is anti-Hermitian and traceless.  Hence \(Z\) begins
at fourth order, its derivatives through order three at the flat
configuration vanish, and its fourth derivative is nonzero.
Take \(\widetilde V=V+\lambda Z\), \(\lambda\ne0\).
An \(x^2v^2\) component of the fourth derivative enters the
background second derivative of the vertical Hessian, which is the
tadpole vertex.  Thus equality of the flat Hessian does not imply
equality of the DQC.
\end{proof}

\begin{firewall}[No purely quadratic shortcut]
\label{fire:e8v52-no-quadratic-shortcut}
Declaring that the added operators (51.5)--(51.6) are purely
quadratic in a logarithmic gauge potential would set their higher
vertices by a coordinate convention, not by a compact
gauge-invariant action.  Such a declaration cannot populate the
quotient card.  A nonlinear compact-link completion must be fixed.
\end{firewall}

\subsection{Two explicit compact gauge-invariant completions}

Write \(P_{\mu\nu}(x)\) for the plaquette based at \(x\).  Parallel
transport its neighbour back to \(x\) by
\[
 (\tau_\rho P_{\mu\nu})(x)
 :=U_\rho(x)P_{\mu\nu}(x+\widehat e_\rho)U_\rho(x)^{-1},       \tag{52.3}
\]
and define the covariant plaquette difference
\[
 {\cal D}_\rho P_{\mu\nu}(x)
 :=(\tau_\rho P_{\mu\nu})(x)-P_{\mu\nu}(x).                   \tag{52.4}
\]
Both summands transform by conjugation with the gauge element at
\(x\).  With \(\|X\|_{\rm HS}^2=\Tr(X^*X)\), define
\[
 {\cal O}_1(U)
 :=\frac12\sum_{x,\mu<\nu}\sum_{\rho=1}^4
 \|{\cal D}_\rho P_{\mu\nu}(x)\|_{\rm HS}^2,                  \tag{52.5}
\]
\[
 {\cal O}_2(U)
 :=\frac12\sum_{x,\mu<\nu}
 \sum_{\rho\in\{\mu,\nu\}}
 \|{\cal D}_\rho P_{\mu\nu}(x)\|_{\rm HS}^2.                  \tag{52.6}
\]
If necessary, average each expression with its reflected image;
this leaves its quadratic Hessian unchanged.

\begin{theorem}[The completions realize the two quartic Ward
tensors]
\label{thm:e8v52-completions}
The actions (52.5)--(52.6) are nonnegative, finite range, and gauge
invariant.  After fixed positive scalar normalizations, their
flat quadratic Hessians in centered-link coordinates are exactly
\({\cal W}^{(1)}\) and \({\cal W}^{(2)}\), respectively.
\end{theorem}

\begin{proof}
Gauge invariance and nonnegativity follow from conjugation
covariance of (52.4) and unitary invariance of the Hilbert--Schmidt
norm.  Only links in two neighbouring plaquettes and one connector
occur, proving finite range.

Put \(U_e=e^{\varepsilon A_e}\).  To first order,
\[
 P_{\mu\nu}(x)=I+\varepsilon
 F_{\mu\nu}^{\rm lin}(x)+O(\varepsilon^2),\qquad
 {\cal D}_\rho P_{\mu\nu}
 =\varepsilon\nabla_\rho F_{\mu\nu}^{\rm lin}
 +O(\varepsilon^2).                                          \tag{52.7}
\]
The connector in (52.3) has no linear commutator with the identity.
In centered Fourier variables,
\[
 F_{\mu\nu}^{\rm lin}(p)
 =i(\bar p_\mu A_\nu-\bar p_\nu A_\mu)
 \quad\hbox{up to the unitary link phases}.                   \tag{52.8}
\]
Therefore (52.5) has quadratic form
\[
 S\sum_{\mu<\nu}
 |\bar p_\mu A_\nu-\bar p_\nu A_\mu|^2,
\]
whose Hessian is (51.5).  In (52.6), the weight for the
\((\mu,\nu)\) curvature is \(x_\mu+x_\nu\).  Its diagonal entry at
\(\nu\) is
\[
 \sum_{\mu\ne\nu}x_\mu(x_\mu+x_\nu)
 =T+Sx_\nu-2x_\nu^2,
\]
and its off-diagonal entry is
\(-\bar p_\mu\bar p_\nu(x_\mu+x_\nu)\), which is (51.6).
The harmless overall Hessian factors are removed by the stated
positive normalizations.
\end{proof}

\begin{remark}[A declared scheme, not a uniqueness assertion]
\label{rem:e8v52-completion-scheme}
Equations (52.5)--(52.6) give a concrete compact-link scheme for
the two local irrelevant couplings.  Proposition
\ref{prop:e8v52-nonlinear-nonuniqueness} shows that this scheme is
not forced by their quadratic germs.  Universality must eventually
show that changing the completion affects only summable irrelevant
terms.
\end{remark}

\subsection{Coordinate-complete mixed jets}

Let
\[
 V_\theta:=V_{\rm W}+\theta_1{\cal O}_1+\theta_2{\cal O}_2     \tag{52.9}
\]
on the parity-tree reduced fibre, with the same Haar measure as the
attached exact regulator.  For a retained one-parameter background
\(v(t)\), let \(x_\theta(t)\) be the true reduced stationary centre:
\[
 D_xV_\theta(x_\theta(t),v(t))=0.                             \tag{52.10}
\]
Define the total stationary vertical Hessian and density term
\[
 H_\theta(t)
 :=D_x^2V_\theta(x_\theta(t),v(t)),\qquad
 h_\theta(t):=-\log J(x_\theta(t)).                           \tag{52.11}
\]
At \(t=0\), put
\[
 A_\theta:=H_\theta(0),\quad G_\theta:=A_\theta^{-1},\quad
 P_\theta:=\partial_tH_\theta(0),\quad
 Q_\theta:=\partial_t^2H_\theta(0).                           \tag{52.12}
\]
All these are full colour--vertical matrices.

\begin{theorem}[Coordinate-complete one-loop response]
\label{thm:e8v52-coordinate-complete-response}
At every point where \(A_\theta>0\),
\[
\boxed{
 \gamma_\theta[v]
 =\frac12\Tr\!\left(
 G_\theta Q_\theta-G_\theta P_\theta G_\theta P_\theta
 \right)+h_\theta''(0).}                                     \tag{52.13}
\]
For \(r=1,2\), differentiation at \(\theta=0\) gives
\[
\boxed{\begin{split}
 \partial_{\theta_r}\gamma_\theta[v]\big|_0
 =\frac12\Tr\bigl(&
 \dot G_r Q+G\dot Q_r
 -\dot G_rPGP-G\dot P_rGP\\
 &-GP\dot G_rP-GPG\dot P_r\bigr)
 +\dot h_r''(0),
\end{split}}                                                  \tag{52.14}
\]
where
\[
 \dot G_r=-G\dot A_rG.                                       \tag{52.15}
\]
\end{theorem}

\begin{proof}
The measured one-loop density is
\[
 \Gamma_\theta(t)
 =\frac12\log\det H_\theta(t)+h_\theta(t).
\]
Twice differentiating the log determinant gives
\(\frac12\Tr(GQ-GPGP)\), proving (52.13).
Differentiate that finite-dimensional identity in \(\theta_r\),
use the product rule and (52.15), and commute the smooth
\(\theta,t\) derivatives to obtain (52.14).
\end{proof}

\begin{corollary}[Reduction to the attached Wilson card]
\label{cor:e8v52-Wilson-reduction}
For \(\theta=0\), if the compensated Wilson path makes the
stationary centre \(O(t^3)\), then the matrices in (52.12) reduce
to the attached \(A,R(b),Q(b)\) data and (52.13) reduces to
(42.12).  Formula (52.14) is the basis-independent version of
(49.5) for the two completion directions.
\end{corollary}

\begin{firewall}[The moving centre is part of the answer]
\label{fire:e8v52-moving-centre}
The attached Wilson compensation has the special property that its
stationary centre moves only at order \(t^3\).  That statement has
not been proved for (52.5)--(52.6).  Their mixed card must use the
total objects (52.11)--(52.12); setting the centre derivatives to
zero would omit legitimate Schur/score terms.
\end{firewall}

\subsection{Finite linear systems for the stationary jets}

Let \(F_\theta(x,t):=D_xV_\theta(x,v(t))\).  Suppress evaluation at
\((x,t,\theta)=(0,0,0)\).  The implicit equation (52.10) gives
\[
 x_t=-A^{-1}F_t,                                              \tag{52.16}
\]
\[
 x_{tt}
 =-A^{-1}\bigl(F_{tt}+2F_{xt}x_t
                    +F_{xx}[x_t,x_t]\bigr).                  \tag{52.17}
\]

\begin{proposition}[Mixed stationary-jet recursion]
\label{prop:e8v52-stationary-recursion}
For each completion direction \(r\), every entry of
\(\dot A_r,\dot P_r,\dot Q_r,\dot h_r''\) in (52.14) is obtained
by:
\begin{enumerate}
\item differentiating the finite loop words in
      \(V_{\rm W},{\cal O}_1,{\cal O}_2\) through total order four;
\item solving (52.16)--(52.17) with the gapped matrix \(A\); and
\item differentiating those two linear systems once in
      \(\theta_r\).
\end{enumerate}
No inverse other than the certified vertical inverse \(A^{-1}\) is
introduced.
\end{proposition}

\begin{proof}
Equations (52.16)--(52.17) follow by differentiating
\(F_\theta(x_\theta(t),t)=0\) once and twice in \(t\).
Their \(\theta_r\) derivatives are again linear equations with
left side \(A\).  The quantities in (52.12) are derivatives of
\(D_x^2V_\theta\) composed with \(x_\theta(t)\), so the chain rule
uses action derivatives of total order at most four and the
stationary jets just listed.  The density term is a fixed analytic
function of \(x_\theta(t)\) and uses the same jets.  Positivity of
\(A\) supplies its unique inverse.
\end{proof}

\begin{definition}[Nonlinear-completion DQC]
\label{def:e8v52-NCDQC}
For \(r=1,2\), the nonlinear-completion differentiated quotient
card consists of:
\begin{enumerate}
\item the normalized loop word (52.5) or (52.6);
\item its flat blocks \(\dot A_r,\dot B_r,\dot C_r\);
\item the stationary jets from (52.16)--(52.17) and their mixed
      derivatives;
\item the total matrices \(\dot P_r,\dot Q_r\) and density jet
      \(\dot h_r''\);
\item the toron-inclusive Fourier trace (52.14); and
\item the external-momentum extraction (49.9).
\end{enumerate}
\end{definition}

\subsection{Attachment check and the next finite task}

\begin{assessment}[What the attached V35 source supplies]
\label{ass:e8v52-attachment-check}
The reduced-fibre attachment supplies the \(96\times64\) plaquette
incidence, the blocks \(A,B,C\), the harmonic router, the exact
Wilson \(R(b),Q(b)\) entries, all \(96\) plaquette types, and the
toron-inclusive finite Fourier trace.  It does not supply the
two-plaquette-with-connector word derivatives required for
\({\cal O}_1,{\cal O}_2\).  Therefore a numerical value for
\(\Xi[{\cal W}^{(r)}]\) cannot rigorously be inferred from its
existing Wilson matrices.
\end{assessment}

\begin{decision}[V53 target]
\label{dec:e8v52-next}
V53 should enumerate the finite word types in (52.5)--(52.6) on
one parity cell and construct their flat \(RR,RC,CC\) Hessian
blocks.  Before third and fourth jets are attempted, it should
prove that these blocks Fourier-transform to (51.5)--(51.6) after
Schur restriction and verify their exact Ward zeros.  If the
parity-cell support crosses more than one coarse cell, the owner
offset bank must be enlarged explicitly.
\end{decision}

\begin{firewall}[Claim boundary after V52]
\label{fire:e8v52-claim}
V52 proves that a quadratic reset basis needs nonlinear
completion, gives two explicit compact gauge-invariant
completions, and derives the coordinate-complete mixed-jet card.
It does not enumerate the new parity-cell matrices, evaluate either
local mixing number, bound the analytic tail, complete the LERB,
establish the universal beta coefficient, close the
nonlinear/complex quotient matching card, construct the full RG
trajectory, remove the regulators, establish
Osterwalder--Schrader reconstruction, or prove the Yang--Mills mass
gap.
\end{firewall}

\subsection{Parent record}

This V52 note was formed by appending the preceding material to V51,
whose validated record was
\[
\begin{array}{c|c}
\text{lines}&23011\\
\text{bytes}&802864\\
\text{SHA-256}&
\texttt{5440A205887CA699582F57288F0820C764C98C7D0489758EE3AE27248496E96B}.
\end{array}
\]

% =============================================================================
% END APPENDED EIGHTH-SERIES V52 MATERIAL
% =============================================================================

% =============================================================================
% BEGIN APPENDED EIGHTH-SERIES V53 MATERIAL
% =============================================================================

\section{Eighth-series V53: exact parity-cell factorization of the
two quartic Hessians}
\label{sec:e8v53}

\subsection{The plaquette parity-shift matrices}

Index fine plaquette types by
\[
 {\cal P}:=\{(s,\mu\nu):
 s\in\{0,1\}^4,\ 1\le\mu<\nu\le4\},\qquad |{\cal P}|=96.      \tag{53.1}
\]
For a plaquette field \(f=(f_{s,\mu\nu})\) and coarse cell
momentum \(k\in{\mathbb T}^4\), define the \(96\times96\) forward
fine-shift difference \(E_\rho(k)\), preserving \(\mu\nu\), by
\[
 (E_\rho(k)f)_{s,\mu\nu}
 :=\begin{cases}
 f_{s+\widehat e_\rho,\mu\nu}-f_{s,\mu\nu},
 &s_\rho=0,\\
 e^{ik_\rho}f_{s-\widehat e_\rho,\mu\nu}-f_{s,\mu\nu},
 &s_\rho=1.
 \end{cases}                                                   \tag{53.2}
\]
Let \(\Pi_{\mu\nu}\) be the orthogonal projection onto the sixteen
parity types with orientation \(\mu\nu\).  Set
\[
 \Omega_1(k):=\sum_{\rho=1}^4E_\rho(k)^*E_\rho(k),             \tag{53.3}
\]
\[
 \Omega_2(k):=\sum_{\mu<\nu}\Pi_{\mu\nu}
 \bigl(E_\mu(k)^*E_\mu(k)+E_\nu(k)^*E_\nu(k)\bigr)
 \Pi_{\mu\nu}.                                                \tag{53.4}
\]

\begin{lemma}[Exact parity representation of fine differences]
\label{lem:e8v53-parity-shift}
Under coarse Fourier synthesis, (53.2) is exactly the forward
difference \(f(z+\widehat e_\rho)-f(z)\) on the fine lattice.
Moreover
\[
 \Omega_r(k)\succeq0,\qquad
 \|\Omega_1(k)\|\le16,\qquad
 \|\Omega_2(k)\|\le8.                                        \tag{53.5}
\]
\end{lemma}

\begin{proof}
Write the fine point as \(z=2y+s\).  If \(s_\rho=0\), its neighbour
has the same coarse cell \(y\) and parity
\(s+\widehat e_\rho\).  If \(s_\rho=1\), the neighbour is
\(2(y+\widehat e_\rho)+(s-\widehat e_\rho)\), producing the phase
\(e^{ik_\rho}\).  This proves the first assertion.
Each fine shift is unitary, so \(\|E_\rho\|\le2\) and
\(\|E_\rho^*E_\rho\|\le4\).  Positivity and the two norm bounds
follow from (53.3)--(53.4).
\end{proof}

\subsection{Factorization through the attached curl incidence}

Let \(D(k)\) be the attached \(96\times64\) fine-plaquette
incidence symbol (YMA3), and let \(D_R(k)\) and \(D_C(k)\) be its
forty-five reduced and four second-constituent column blocks.  The
fifteen parity-tree columns are fixed to zero on the section.

\begin{theorem}[Exact flat Hessian factorization]
\label{thm:e8v53-Hessian-factorization}
With the positive normalizations fixed in
Theorem~\ref{thm:e8v52-completions}, the scalar colour factors of
the flat Hessians of \({\cal O}_r\), \(r=1,2\), are
\[
 \boxed{{\mathsf H}_r(k)=D(k)^*\Omega_r(k)D(k).}               \tag{53.6}
\]
Their parity-tree-section blocks are therefore
\[
\boxed{\begin{aligned}
 A_r(k)&=D_R(k)^*\Omega_r(k)D_R(k),\\
 B_r(k)&=D_R(k)^*\Omega_r(k)D_C(k),\\
 C_r(k)&=D_C(k)^*\Omega_r(k)D_C(k).
\end{aligned}}                                                \tag{53.7}
\]
The full colour Hessian is
\({\mathsf H}_r\otimes I_8\).
\end{theorem}

\begin{proof}
At the flat field, (52.7) identifies the linearized plaquette
curvature with \(Da\).  The quadratic part of (52.5) is the squared
norm of all four forward differences of \(Da\), hence
\[
 \langle Da,\Omega_1Da\rangle
 =\langle a,D^*\Omega_1Da\rangle.
\]
For (52.6), only the \(\mu,\nu\) differences of a
\(\mu\nu\)-plaquette enter; the same calculation with (53.4) gives
\(D^*\Omega_2D\).  Restricting the columns to \(R,C\) gives
(53.7).  Adjoint invariance of the Hilbert--Schmidt metric makes
the flat colour factor the identity.
\end{proof}

\begin{corollary}[Positivity and finite support]
\label{cor:e8v53-positive-local}
For every \(k\),
\[
 \begin{pmatrix}A_r&B_r\\B_r^*&C_r\end{pmatrix}\succeq0.      \tag{53.8}
\]
Every entry is a finite Laurent polynomial in
\(e^{\pm ik_\rho}\).  Its offset bank is obtained by adding one
fine-neighbour shift from (53.2) to each of the two plaquette
incidences in (53.6); no infinite-range owner is introduced.
\end{corollary}

\begin{proof}
Equation (53.8) is the compression of
\((\Omega_r^{1/2}D)^*(\Omega_r^{1/2}D)\).
Both \(D\) and \(E_\rho\) are finite Laurent matrices, so their
finite products are finite Laurent matrices.
\end{proof}

\begin{corollary}[A signed deformation retains a vertical gap]
\label{cor:e8v53-gap-chart}
Let \(A_0=\frac13D_R^*D_R\) be the Wilson vertical block, with
\(A_0\succeq I/60\).  If
\[
 |\theta_1|\sup_k\|A_1(k)\|
 +|\theta_2|\sup_k\|A_2(k)\|<\frac1{60},                      \tag{53.9}
\]
then
\[
 A_\theta=A_0+\theta_1A_1+\theta_2A_2
 \succeq
 \left(\frac1{60}
 -|\theta_1|\|A_1\|_\infty
 -|\theta_2|\|A_2\|_\infty\right)I>0.                         \tag{53.10}
\]
For \(\theta_1,\theta_2\ge0\), the original floor \(1/60\) is
preserved.
\end{corollary}

\begin{proof}
Use the Wilson floor, the operator-norm perturbation inequality,
and \(A_r\succeq0\) for the last statement.
\end{proof}

\subsection{Exact curl--gradient Ward zero}

Let \(d(k)\) be the fine link symbol of an infinitesimal site gauge
gradient, including its parity components.

\begin{lemma}[Incidence complex]
\label{lem:e8v53-incidence-complex}
The oriented plaquette incidence satisfies
\[
 D(k)d(k)=0.                                                  \tag{53.11}
\]
\end{lemma}

\begin{proof}
For every plaquette, the four oriented edge differences of a site
potential telescope: each corner appears once with each sign.
Fourier transformation preserves that exact incidence identity.
\end{proof}

\begin{corollary}[Exact Ward zero of both deformations]
\label{cor:e8v53-Ward-zero}
For \(r=1,2\),
\[
 {\mathsf H}_r(k)d(k)=0,\qquad
 d(k)^*{\mathsf H}_r(k)=0.                                   \tag{53.12}
\]
Thus their unreduced flat Hessians pass the exact Ward row before
any numerical inversion.
\end{corollary}

\begin{proof}
Insert (53.11) into (53.6) and take adjoints.
\end{proof}

\begin{firewall}[Tree restriction versus gauge annihilation]
\label{fire:e8v53-tree-Ward}
The \(45\times45\) vertical block \(A_r\) is taken after fixing the
parity-tree directions and is not supposed to retain a gauge
kernel.  Equation (53.12) belongs to the complete link Hessian.
The retained Schur response inherits the coarse Ward identity from
the complete block; one must not demand that the invertible
vertical block annihilate a removed gauge direction.
\end{firewall}

\subsection{Router and Schur derivatives without entrywise algebra}

Retain the Wilson quantities
\[
 G=A_0^{-1},\qquad L=-GB_0,\qquad
 {\mathsf S}_0=C_0-B_0^*GB_0.                                \tag{53.13}
\]

\begin{theorem}[Exact first deformation of the harmonic router]
\label{thm:e8v53-router-derivative}
For \(r=1,2\),
\[
 \boxed{\dot L_r=-G(A_rL+B_r).}                              \tag{53.14}
\]
\end{theorem}

\begin{proof}
Differentiate
\((A_0+\theta A_r)L_\theta+(B_0+\theta B_r)=0\) at
\(\theta=0\) and multiply by \(G\).
\end{proof}

\begin{theorem}[Exact first deformation of the retained Schur
kernel]
\label{thm:e8v53-Schur-derivative}
The derivative of the retained flat Schur kernel is
\[
\boxed{\begin{split}
 \dot{\mathsf S}_r
 ={}&C_r-B_r^*GB_0-B_0^*GB_r+B_0^*GA_rGB_0\\
 ={}&C_r+L^*B_r+B_r^*L+L^*A_rL.
\end{split}}                                                  \tag{53.15}
\]
Equivalently, for every retained field \(a\),
\[
 \langle a,\dot{\mathsf S}_ra\rangle
 =\left\langle
 \binom{La}{a},
 \begin{pmatrix}A_r&B_r\\B_r^*&C_r\end{pmatrix}
 \binom{La}{a}\right\rangle\ge0.                             \tag{53.16}
\]
\end{theorem}

\begin{proof}
Differentiate
\({\mathsf S}_\theta
=C_\theta-B_\theta^*A_\theta^{-1}B_\theta\) and use
\(\partial_\theta A_\theta^{-1}|_0=-GA_rG\).  This gives the first
line.  Substitute \(L=-GB_0\) for the second.  Equation (53.16) is
the resulting block quadratic form, nonnegative by (53.8).
\end{proof}

\begin{corollary}[Finite exact construction of the flat DQC rows]
\label{cor:e8v53-flat-DQC}
At each discrete momentum, the flat rows
\[
 \dot A_r,\quad\dot B_r,\quad\dot C_r,\quad
 \dot G_r,\quad\dot L_r,\quad\dot{\mathsf S}_r                \tag{53.17}
\]
are determined by the attached integer-Laurent incidence \(D\),
the explicit two-point parity shifts (53.2), and one certified
Wilson inverse \(G\).  No two-plaquette word expansion is required
for these quadratic rows.
\end{corollary}

\subsection{What remains genuinely nonlinear}

\begin{assessment}[Progress after V53]
\label{ass:e8v53-status}
The V52 flat enumeration target is closed by the factorizations
(53.6)--(53.7).  They give positivity, locality, exact Ward zeros,
the router derivative, and the retained Schur derivative without a
large entry table.  The remaining NCDQC data are the third and
fourth mixed group-word jets entering
\(\dot P_r,\dot Q_r,\dot h_r''\) in (52.14).  Those jets are not
recoverable from (53.6), by Proposition
\ref{prop:e8v52-nonlinear-nonuniqueness}.
\end{assessment}

\begin{decision}[V54 target]
\label{dec:e8v53-next}
V54 should differentiate the covariant plaquette difference
\({\cal D}_\rho P\) itself.  A product-rule factorization should
express the required third and fourth jets through: the first two
plaquette-word derivatives, the first two connector-adjoint
derivatives, and the already known harmonic lift \(b[a]\) plus
(53.14).  This is preferable to enumerating complete link words
independently for every matrix entry.
\end{decision}

\begin{firewall}[Claim boundary after V53]
\label{fire:e8v53-claim}
V53 exactly populates the flat Hessian, gap, Ward, router, and
Schur rows for both nonlinear completions.  It does not populate
their third and fourth mixed jets, evaluate either local mixing
number, bound the analytic tail, complete the LERB, establish the
universal beta coefficient, close the nonlinear/complex quotient
matching card, construct the full RG trajectory, remove the
regulators, establish Osterwalder--Schrader reconstruction, or
prove the Yang--Mills mass gap.
\end{firewall}

\subsection{Parent record}

This V53 note was formed by appending the preceding material to V52,
whose validated record was
\[
\begin{array}{c|c}
\text{lines}&23351\\
\text{bytes}&815366\\
\text{SHA-256}&
\texttt{ED3B842F57E8BCBB584B8650CBC08F12BA1B5B7C357C92B76ECA90CA4D89D1A3}.
\end{array}
\]

% =============================================================================
% END APPENDED EIGHTH-SERIES V53 MATERIAL
% =============================================================================

% =============================================================================
% BEGIN APPENDED EIGHTH-SERIES V54 MATERIAL
% =============================================================================

\section{Eighth-series V54: factorized third and fourth jets of the
covariant plaquette-difference owners}
\label{sec:e8v54}

\subsection{Labeled mixed derivatives of a finite group word}

Let \(I=\{1,\ldots,n\}\) label \(n\) independent link-field
directions.  If \(F\) is a smooth matrix-valued function, write
\[
 F_I:=D^nF(0)[h_1,\ldots,h_n],\qquad F_\varnothing:=F(0).      \tag{54.1}
\]
For an ordered word \(W=g_1\cdots g_m\), an ordered partition
\(I_1\sqcup\cdots\sqcup I_m=I\) allows empty parts and preserves
the factor order.

\begin{lemma}[Finite-word derivative]
\label{lem:e8v54-word-derivative}
At every order \(n\),
\[
 \boxed{
 W_I
 =\sum_{I_1\sqcup\cdots\sqcup I_m=I}
 (g_1)_{I_1}\cdots(g_m)_{I_m}.}                              \tag{54.2}
\]
If \(\bar C=C^{-1}\) and \(C(0)=I\), its derivatives are determined
recursively by
\[
 \boxed{
 \bar C_I
 =-\sum_{\varnothing\ne J\subseteq I}
 C_J\bar C_{I\setminus J}\qquad(I\ne\varnothing).}            \tag{54.3}
\]
\end{lemma}

\begin{proof}
Repeated use of the product rule assigns each labeled derivative
to exactly one factor, giving (54.2).  Differentiate
\(C\bar C=I\) in all directions in \(I\):
\[
 0=\sum_{J\subseteq I}C_J\bar C_{I\setminus J}.
\]
The \(J=\varnothing\) term is \(\bar C_I\); moving the other terms
to the opposite side gives (54.3).
\end{proof}

\subsection{Derivatives of the covariant difference}

For one owner in (52.4), abbreviate
\[
 C:=U_\rho(x),\qquad P_+:=P_{\mu\nu}(x+\widehat e_\rho),
 \qquad P_0:=P_{\mu\nu}(x),
\]
\[
 Y:=CP_+C^{-1}-P_0.                                          \tag{54.4}
\]
At the flat configuration,
\[
 C=P_+=P_0=I,\qquad Y_\varnothing=0.                          \tag{54.5}
\]

\begin{proposition}[Three-jet factorization of one owner]
\label{prop:e8v54-Y-jets}
For every nonempty \(I\) with \(|I|\le3\),
\[
\boxed{
 Y_I
 =\sum_{I_1\sqcup I_2\sqcup I_3=I}
 C_{I_1}(P_+)_{I_2}\bar C_{I_3}
 -(P_0)_I,}                                                   \tag{54.6}
\]
where \(\bar C_I\) is generated by (54.3).  Thus the required
derivatives use only the first three derivatives of the connector
link and of the two plaquette words.
\end{proposition}

\begin{proof}
Apply (54.2) to the three-factor transported plaquette and subtract
the derivative of \(P_0\).  Formula (54.3) eliminates independent
inverse-word data.
\end{proof}

\begin{corollary}[The linear jet is the incidence difference]
\label{cor:e8v54-linear-jet}
For one direction \(h\),
\[
 Y_h=(P_+)_h-(P_0)_h.                                        \tag{54.7}
\]
The connector does not occur at first order.
\end{corollary}

\begin{proof}
In (54.6), the two connector terms are
\(C_h+\bar C_h=0\) by (54.3); the remaining terms give (54.7).
This is the word-level reason for the factorization (53.6).
\end{proof}

\subsection{The squared-norm jet compiler}

Use the real Hilbert--Schmidt pairing
\[
 \langle X,Z\rangle_{\mathbb R}:=\Re\Tr(X^*Z),\qquad
 o(Y):=\frac12\langle Y,Y\rangle_{\mathbb R}.                 \tag{54.8}
\]

\begin{theorem}[All action jets through order four]
\label{thm:e8v54-norm-jets}
At a point where \(Y=0\), the derivatives of the owner (54.8) are
\[
 D^2o[a,b]=\langle Y_a,Y_b\rangle_{\mathbb R},                \tag{54.9}
\]
\[
\begin{split}
 D^3o[a,b,c]={}&
 \langle Y_{ab},Y_c\rangle_{\mathbb R}
 +\langle Y_{ac},Y_b\rangle_{\mathbb R}
 +\langle Y_{bc},Y_a\rangle_{\mathbb R},                     \tag{54.10}
\end{split}
\]
and
\[
\begin{split}
 D^4o[a,b,c,d]={}&
 \langle Y_{ab},Y_{cd}\rangle_{\mathbb R}
 +\langle Y_{ac},Y_{bd}\rangle_{\mathbb R}
 +\langle Y_{ad},Y_{bc}\rangle_{\mathbb R}\\
 &+\langle Y_{abc},Y_d\rangle_{\mathbb R}
 +\langle Y_{abd},Y_c\rangle_{\mathbb R}\\
 &+\langle Y_{acd},Y_b\rangle_{\mathbb R}
 +\langle Y_{bcd},Y_a\rangle_{\mathbb R}.                    \tag{54.11}
\end{split}
\]
In particular, no fourth derivative \(Y_{abcd}\) is needed.
\end{theorem}

\begin{proof}
Differentiate \(\frac12\langle Y,Y\rangle_{\mathbb R}\) and use
the symmetry of the real pairing.  At order \(n\), terms are indexed
by unordered two-block partitions of the \(n\) labeled directions.
For \(n=2\) there is one \(1+1\) partition; for \(n=3\), three
\(2+1\) partitions; and for \(n=4\), three \(2+2\) and four
\(3+1\) partitions.  Terms containing \(Y_\varnothing\) vanish by
(54.5), giving (54.9)--(54.11).
\end{proof}

\begin{corollary}[The NCDQC uses only three word jets]
\label{cor:e8v54-three-word-jets}
Every action derivative required in the stationary recursion
(52.16)--(52.17) and in
\(\dot A_r,\dot P_r,\dot Q_r\) is a finite sum of pairings of the
three tensors (54.6).  Hence fourth-order automatic
differentiation of a long completed action is unnecessary.
\end{corollary}

\subsection{Explicit word-length majorants}

Suppose each labeled link direction satisfies
\(\|h_i\|\le d_i\).  For a link exponential or inverse exponential,
the standard ordered-integral derivative formula gives
\[
 \|(g_e)_I\|\le\prod_{i\in I}d_i.                             \tag{54.12}
\]

\begin{lemma}[Uniform three-jet owner bound]
\label{lem:e8v54-owner-bound}
The transported plaquette word \(CP_+C^{-1}\) has six link factors
and \(P_0\) has four.  Therefore, for \(1\le n\le3\),
\[
 \boxed{
 \|Y_I\|\le M_n\prod_{i\in I}d_i,\qquad
 M_n:=6^n+4^n.}                                               \tag{54.13}
\]
Explicitly,
\[
 M_1=10,\qquad M_2=52,\qquad M_3=280.                        \tag{54.14}
\]
\end{lemma}

\begin{proof}
In (54.2), each of the \(n\) labeled derivatives chooses one of
\(m\) factors, so there are \(m^n\) ordered assignments.  Bound
each term by (54.12), first for the six-factor transported word and
then for the four-factor plaquette, and use the triangle
inequality.
\end{proof}

\begin{proposition}[Explicit local action-jet stocks]
\label{prop:e8v54-action-stocks}
For one unnormalized owner \(o(Y)\),
\[
\begin{aligned}
 |D^2o[h_1,h_2]|
 &\le100\,d_1d_2,\\
 |D^3o[h_1,h_2,h_3]|
 &\le1560\,d_1d_2d_3,\\
 |D^4o[h_1,h_2,h_3,h_4]|
 &\le19312\,d_1d_2d_3d_4.                                  \tag{54.15}
\end{aligned}
\]
Per parity cell, \({\cal O}_1\) has \(16\cdot6\cdot4=384\)
translated owner terms and \({\cal O}_2\) has
\(16\cdot6\cdot2=192\).  Multiplication by these counts and by the
fixed normalizations of (52.5)--(52.6) gives volume-independent
local jet stocks.
\end{proposition}

\begin{proof}
Use (54.9)--(54.11), Cauchy--Schwarz, and (54.14):
\[
 M_1^2=100,\qquad
 3M_2M_1=1560,\qquad
 3M_2^2+4M_3M_1=19312.
\]
There are sixteen fine base parities and six plaquette
orientations.  The first action uses four difference directions
and the second uses two, proving the counts.  Each owner has a
fixed finite support, so a local row/column norm sees only a
volume-independent overlap multiple.
\end{proof}

\begin{firewall}[Crude stocks are not the physical mixing number]
\label{fire:e8v54-crude-stocks}
The constants (54.15) certify finiteness and volume independence.
They take absolute values before the colour trace, periodic
momentum sum, and Ward cancellation, so they must not be used as a
sign or sharp estimate for \(\Xi\).  The scalar DQC still uses
(52.14) followed by (49.9).
\end{firewall}

\subsection{Exact assembly of the mixed matrices}

For \(r=1,2\), sum (54.9)--(54.11) over the owner bank declared in
(52.5)--(52.6).  Insert reduced fluctuation directions in the first
two slots and background/harmonic directions in the remaining
slots.

\begin{proposition}[Finite mixed-jet assembly]
\label{prop:e8v54-mixed-assembly}
The direct completion contributions are:
\[
\begin{array}{c|c}
\text{matrix}&\text{owner derivative}\\ \hline
A_r&D^2{\cal O}_r[x,x]\\
P_r(b)&D^3{\cal O}_r[x,x,b]\\
Q_r(b_1,b_2)&D^4{\cal O}_r[x,x,b_1,b_2].
\end{array}                                                    \tag{54.16}
\]
The total dotted matrices in (52.14) are obtained by adding:
\begin{enumerate}
\item the Wilson vertex derivatives caused by the router change
      \(\dot b_r\), with \(\dot L_r\) given by (53.14);
\item the stationary-centre chain-rule terms generated by
      (52.16)--(52.17); and
\item the Haar-density stationary jet \(\dot h_r''\).
\end{enumerate}
Every item is a finite contraction of (54.6), the attached Wilson
\(A,B,R,Q\) matrices, and the single inverse \(G=A^{-1}\).
\end{proposition}

\begin{proof}
The vertical Hessian uses two fluctuation derivatives.  One or two
background derivatives give respectively the third and fourth
rows of (54.16), exactly as in (52.12).  The background inserted
in the reduced coordinates is the harmonic lift; differentiating
it gives (53.14).  Composition with the true stationary centre and
with the fixed Haar density gives the last two chain-rule families,
whose finite recursion was proved in
Proposition~\ref{prop:e8v52-stationary-recursion}.  No other source
of a \(\theta_r\) derivative occurs.
\end{proof}

\begin{definition}[Factorized mixed-jet certificate]
\label{def:e8v54-FMJC}
A factorized mixed-jet certificate (FMJC) stores, for each owner
type and each required direction tuple:
\begin{enumerate}
\item the link words \(C,P_+,P_0\);
\item the three derivative arrays \(Y_I\), \(|I|\le3\), from
      (54.2)--(54.6);
\item the pairings (54.9)--(54.11);
\item the parity-cell offsets and colour trace; and
\item the assembly residual against (53.6) in the quadratic row.
\end{enumerate}
\end{definition}

\subsection{Status and the V55 audit target}

\begin{assessment}[Progress after V54]
\label{ass:e8v54-status}
The higher-jet acquisition is now finite and factorized.  The
quartic action derivative needs only third derivatives of a
six-factor word, with explicit volume-free majorants.  The
quadratic row (53.6) provides an exact independent residual check
for the assembly.  What remains is to execute the finite FMJC and
then perform the Ward-summed trace; no analytic existence issue
remains at fixed volume for these two local completion directions.
\end{assessment}

\begin{decision}[Mandatory V55 audit and next calculation]
\label{dec:e8v54-next}
V55 must inspect the current SM and NS companion notes before
continuing.  It should then test whether either programme provides
a reusable certified finite-word or Ward-summed trace mechanism.
Absent such a transfer, the next Yang--Mills task is to instantiate
the FMJC on one parity cell, first verifying (53.6), then producing
\(\dot P_r,\dot Q_r\) for insertion into (52.14).
\end{decision}

\begin{firewall}[Claim boundary after V54]
\label{fire:e8v54-claim}
V54 proves the factorized word-derivative compiler and explicit
finite jet stocks.  It does not execute the FMJC, evaluate either
local mixing number, bound the analytic tail, complete the LERB,
establish the universal beta coefficient, close the
nonlinear/complex quotient matching card, construct the full RG
trajectory, remove the regulators, establish
Osterwalder--Schrader reconstruction, or prove the Yang--Mills mass
gap.
\end{firewall}

\subsection{Parent record}

This V54 note was formed by appending the preceding material to V53,
whose validated record was
\[
\begin{array}{c|c}
\text{lines}&23661\\
\text{bytes}&825903\\
\text{SHA-256}&
\texttt{22772C5A7342B35B4FCF776B9C0BDF31975C37CAD85AA755602162AA86A789E1}.
\end{array}
\]

% =============================================================================
% END APPENDED EIGHTH-SERIES V54 MATERIAL
% =============================================================================

% =============================================================================
% BEGIN APPENDED EIGHTH-SERIES V55 MATERIAL
% =============================================================================

\section{Eighth-series V55: compulsory audit and transfer of the
bounded-box nonabelian curvature compiler}
\label{sec:e8v55}

\subsection{Compulsory companion audit}

The current companion records at this checkpoint are
\[
\begin{array}{c|r|r|l}
\text{record}&\text{lines}&\text{bytes}&\text{SHA-256}\\ \hline
\texttt{Renewal\_SM\_Einstein\_Continuum\_Research\_Notes\_V44.tex}
&16545&633070&
\texttt{0E31F05ED80B6BBB42CCD7C972E68DC1C9517A74AF53E0A487738A6416D4CF20}\\
\texttt{Renewal\_NS\_Stability\_Research\_Notes\_V26.tex}
&5320&278283&
\texttt{82C887F8C2C69E4F28FAD7C3DC8DE5B9099CB5A4BF431EBA1FFF3E91935978AD}.
\end{array}                                                    \tag{55.1}
\]
The SM file advanced from V41 at the V50 audit to V44.  Its V42
closes, on bounded contractible small-curvature patches, the
quantitative covariant Poincare acquisition that V50 correctly
left open.  The NS file is unchanged.

\begin{proposition}[V55 audit decision]
\label{prop:e8v55-audit-decision}
The SM V42 cubical homotopy, nonabelian curvature fixed point, and
polymer box-reserve estimate transfer to the present
contractible Yang--Mills owners with their stated hypotheses.  Its
metric V43 and DeTurck--CMC V44 results do not transfer to the
pure-gauge reset card.  No NS formula or numerical constant
transfers.  The NS actual-input norm lesson was already incorporated
as \({\cal Y}_j\) in V50 and needs no repetition.
\end{proposition}

\begin{proof}
The SM V42 statements are theorems about compact lattice gauge
links, cubical cochains, local gauge-invariant analytic owners, and
polymer support.  None of their proofs uses a Standard Model
representation or a metric field, so the same \(SU(3)\) link
variables meet their hypotheses.  The later SM results concern
metric gauge fixing and constraint determinants, absent from the
present local pure-gauge card.  The NS record concerns Oseen
kernels and rank-one fluid inputs and contains no compact-link
curvature or determinant statement.
\end{proof}

\subsection{Imported bounded-box theorem}

The following theorem is imported from SM V42 rather than
reproved; this avoids duplicating its coordinate-by-coordinate
cubical homotopy and Banach fixed-point proof.

\begin{theorem}[Bounded-box nonabelian curvature reconstruction]
\label{thm:e8v55-imported-curvature}
Let \(B\) be a rectangular contractible cubical complex of diameter
\(r\).  There is a cochain homotopy
\[
 h_k:C^k(B)\longrightarrow C^{k-1}(B),\qquad
 dh_k+h_{k+1}d=I,\qquad
 \|h_k\|_{\infty\leftarrow\infty}\le C_4r.                   \tag{55.2}
\]
On a small logarithm chart, write plaquette curvature as
\[
 {\cal K}(a)=da+Q(a),\qquad
 \|Q(a)-Q(b)\|_\infty
 \le C_Q(\|a\|_\infty+\|b\|_\infty)\|a-b\|_\infty.            \tag{55.3}
\]
Put \(H=\|h_2\|\).  If
\[
 4H^2C_Q\|f\|_\infty<1,\qquad
 2H\|f\|_\infty<R_0,                                         \tag{55.4}
\]
then on the slice \(h_1a=0\) the equation
\[
 a=h_2(f-Q(a))                                               \tag{55.5}
\]
has a unique analytic solution \(a(f)\), supported in the same
box, with
\[
 \|a(f)\|_\infty\le2H\|f\|_\infty.                           \tag{55.6}
\]
Every analytic gauge-invariant owner \({\cal F}\) on that patch
therefore factors as
\[
 {\cal F}(a)=\widehat{\cal F}({\cal K}(a)).                   \tag{55.7}
\]
If an owner support \(X\) is assigned a box with
\[
 |B_X|\le c_B|X|,\qquad r_X\le c_r|X|,                       \tag{55.8}
\]
then its homotopy length and box enlargement are absorbed by
reducing a polymer exponent \(\kappa\) to any
\(\kappa'<\kappa/c_B\).
\end{theorem}

\begin{remark}[Exact scope of the import]
\label{rem:e8v55-import-scope}
Because \(H=O(r)\), (55.4) requires
\(\|f\|_\infty=O(r^{-2})\).  This is the block-smooth curvature
regime, not a global chart for arbitrary plaquette fields.
Nonzero flux, torons, multiply connected owners, and overlaps of
different local slices remain separate.
\end{remark}

\subsection{Application to the Yang--Mills reset tail}

\begin{lemma}[Dimension census after the Yang--Mills projections]
\label{lem:e8v55-dimension-census}
On a reflection-even contractible pure-gauge patch, project the
vacuum row and the marginal \(F^2\) row.  Then every remaining
nonconstant gauge-invariant local monomial in the curvature
representation (55.7) has engineering dimension at least six.
At dimension six the possible leading types include \(F^3\) and
\((D F)^2\); the latter contains the two quadratic tensors of V51.
\end{lemma}

\begin{proof}
Curvature has dimension two and each covariant difference has
dimension one.  Gauge invariance and index contraction exclude a
nonconstant scalar of dimension below four.  At dimension four the
reflection-even pure-gauge scalar is the kinetic contraction
\(\Tr F_{\mu\nu}F_{\mu\nu}\); the parity-odd topological contraction
is absent under reflection and otherwise belongs to its own
projected row.  Once these are removed, the next possible
gauge-invariant contractions have either three curvatures or two
curvatures and two covariant differences, hence dimension six.
The quadratic part of the second class is precisely a hypercubic
transverse quartic-momentum tensor, whose two-dimensional
classification was proved in V51.
\end{proof}

\begin{theorem}[Canonical local-tail contraction]
\label{thm:e8v55-local-tail-contraction}
Let \({\cal A}_j^{\rm loc}\) be a contractible analytic reset owner
at dyadic scale \(j\).  Assume:
\begin{enumerate}
\item its curvature patch satisfies (55.4);
\item its polymer boxes satisfy (55.8) with a fixed reserve;
\item the vacuum, \(F^2\), and any parity-odd topological rows are
      exactly projected; and
\item the curvature expansion (55.7) is absolutely summable in the
      reduced polymer norm.
\end{enumerate}
Then every homogeneous component of the remaining local owner has
relative scale factor at most
\[
 2^{4-d}\le\frac14,\qquad d\ge6,                              \tag{55.9}
\]
under one dyadic block.  Consequently, if new shell forcing has
norm \(f_j\),
\[
 \|{\cal A}_{j+1}^{\rm loc}\|_{\rm irr}
 \le\frac14\|{\cal A}_j^{\rm loc}\|_{\rm irr}+f_j.             \tag{55.10}
\]
\end{theorem}

\begin{proof}
Theorem~\ref{thm:e8v55-imported-curvature} writes every owner as an
absolutely convergent curvature series on its bounded box, while
the reserve in (55.8) absorbs the homotopy support and one
polynomial diameter loss.  Under block-smooth rescaling, curvature
has size \(2^{-2}\) and a covariant difference contributes an
additional \(2^{-1}\).  The block volume contributes \(2^4\).
Thus a dimension-\(d\) monomial has relative factor \(2^{4-d}\).
Lemma~\ref{lem:e8v55-dimension-census} gives \(d\ge6\), proving
(55.9).  Absolute summability permits termwise blocking and
summation.  Adding the newly generated owner gives (55.10).
\end{proof}

\begin{corollary}[The V51 completions are the leading local
irrelevant rows]
\label{cor:e8v55-leading-rows}
The actions \({\cal O}_1,{\cal O}_2\) of (52.5)--(52.6) have
leading type \((D F)^2\), dimension six, and relative dyadic factor
\(1/4\).  After they are retained as explicit coordinates, the
quadratic analytic tail begins at dimension eight and has relative
factor at most \(1/16\).
\end{corollary}

\begin{proof}
Their linearized curvatures and derivatives were computed in
(52.7)--(52.8), and their quadratic symbols begin at order
\(p^4\).  The remainder after \({\cal P}_4\) begins at order
\(p^6\), corresponding to two curvatures and four differences,
dimension eight.  Apply (55.9).
\end{proof}

\subsection{From local activities to the response bound}

Let \(C_{\rm jet}\) denote the fixed conversion from the reduced
polymer activity norm to the four derivative stocks
\((a_1,b_1,q_1,r_1)\) of V45.  For the explicit completions, it is
finite by (54.15) and the finite overlap number.

\begin{theorem}[Local analytic tail populates the response row]
\label{thm:e8v55-local-response}
Assume the uniform vertical stocks \(g,b,q,r\) of V45 and suppose
the compiled contractible reset tail satisfies
\[
 \int_0^1(a_1+b_1+q_1+r_1)(t)\,dt
 \le C_{\rm jet}\|{\cal A}_j^{\rm loc}\|_{\rm irr}.            \tag{55.11}
\]
Then its complete one-loop response obeys
\[
 \left|
 {\mathfrak m}\bigl(
 {\cal F}_j({\cal A}_j^{\rm loc})-{\cal F}_j(0)
 \bigr)\right|
 \le
 \mu^{-2}C_{\rm loop}C_{\rm jet}
 \|{\cal A}_j^{\rm loc}\|_{\rm irr}.                          \tag{55.12}
\]
If \(f_j\le c\,2^{-j}\) in (55.10), the local-tail contribution to
the marginal drift is absolutely summable.
\end{theorem}

\begin{proof}
First form the complete compensated tadpole--bubble kernel.  The
path bound (46.18) and (55.11) give its weighted norm difference at
most
\(C_{\rm loop}C_{\rm jet}\|{\cal A}_j^{\rm loc}\|_{\rm irr}\).
Apply the moment bound (48.1), proving (55.12).  Recurrence
(55.10) is (50.11), so Proposition
\ref{prop:e8v50-quartic-recurrence} makes the activity norms
summable when the forcing is \(O(2^{-j})\).
\end{proof}

\begin{firewall}[What remains outside the local-tail theorem]
\label{fire:e8v55-tail-scope}
Theorem~\ref{thm:e8v55-local-response} controls contractible
analytic owners satisfying the polymer gate.  It does not control
winding toron owners, flux sectors, transition functions between
different local curvature charts, or an edge contribution for
which no uniform polymer decomposition has been proved.  These
owners cannot be relabeled as part of the local tail.
\end{firewall}

\subsection{Audit outcome and next acquisition}

\begin{assessment}[Progress after V55]
\label{ass:e8v55-status}
The new companion theorem closes the local curvature-compilation
step that was open at V50.  Combined with V45 and the explicit V54
jet stocks, it yields the response bound and summability theorem
(55.12) for the contractible analytic reset tail.  The two leading
dimension-six coordinates still require their exact DQC values,
while the remaining global obstruction has been narrowed to toron,
flux, chart-overlap, and any as-yet-unlocalized edge owners.
\end{assessment}

\begin{decision}[V56 target]
\label{dec:e8v55-next}
V56 should return to the two leading coordinates and insert the
factorized jets (54.16) into the coordinate-complete trace (52.14).
Before numerical evaluation, it should prove a finite-volume
symmetry reduction of the owner bank under translations,
hypercubic transformations, reflection, and orientation reversal.
The resulting orbit representatives must retain the \(k=0\)
summand and exact multiplicities.
\end{decision}

\begin{firewall}[Claim boundary after V55]
\label{fire:e8v55-claim}
V55 performs the compulsory audit and, using the newly available
bounded-box curvature compiler, proves conditional contraction and
one-loop response control for the contractible analytic local
tail.  It does not evaluate the two dimension-six DQC numbers,
control toron/flux/overlap owners, complete the LERB, establish the
universal beta coefficient, close the nonlinear/complex quotient
matching card, construct the full RG trajectory, remove the
regulators, establish Osterwalder--Schrader reconstruction, or
prove the Yang--Mills mass gap.
\end{firewall}

\subsection{Parent record}

This V55 note was formed by appending the preceding material to V54,
whose validated record was
\[
\begin{array}{c|c}
\text{lines}&24001\\
\text{bytes}&837524\\
\text{SHA-256}&
\texttt{388B846BD35A5B19F29870BEBBEF40A5DC0F2188892048B82E143E566FDE30C9}.
\end{array}
\]

% =============================================================================
% END APPENDED EIGHTH-SERIES V55 MATERIAL
% =============================================================================

% =============================================================================
% BEGIN APPENDED EIGHTH-SERIES V56 MATERIAL
% =============================================================================

\section{Eighth-series V56: exact symmetry reduction of the
dimension-six owner bank}
\label{sec:e8v56}

\subsection{A symmetry-invariant marginal extractor}

The single polarization in (48.1) is sufficient after hypercubic
symmetry has been proved.  For orbit reduction it is safer to
average that polarization.  In dimension \(d=4\), define
\[
 {\mathfrak m}_{\rm sym}(K)
 :=-\frac1{2d(d-1)}
 \sum_{\substack{\alpha,\beta=1\\\alpha\ne\beta}}^d
 \sum_x x_\alpha^2K_{\beta\beta}(x).                          \tag{56.1}
\]

\begin{lemma}[Agreement with the axial moment]
\label{lem:e8v56-symmetric-moment}
If \(K\) is hypercubic covariant, then
\[
 {\mathfrak m}_{\rm sym}(K)={\mathfrak m}(K).                 \tag{56.2}
\]
For an arbitrary kernel, \({\mathfrak m}_{\rm sym}\) is invariant
under coordinate permutations and reflections.
\end{lemma}

\begin{proof}
For a hypercubic kernel, every ordered transverse pair
\((\alpha,\beta)\), \(\alpha\ne\beta\), is carried to
\((1,2)\) by a signed coordinate permutation, so all \(d(d-1)\)
summands in (56.1) agree.  Their average is (48.1).
For a general kernel, a coordinate symmetry merely permutes the
ordered pairs and changes \(x_\alpha\) by a sign whose square is
unchanged.
\end{proof}

For a finite analytic card, the corresponding normalized mixing
extractor is
\[
 \Xi_{\rm sym}
 :=\frac{3}{2d(d-1)}
 \sum_{\alpha\ne\beta}
 \left.\partial_q^2\right|_{q=0}
 \dot\gamma_{\beta\mid\alpha}(q),                            \tag{56.3}
\]
where the external momentum is \(q\widehat e_\alpha\) and the
retained polarization is \(\widehat e_\beta\).

\subsection{Equivariant owner sums}

Let \({\cal I}\) be a finite bank of local action owners and let a
finite lattice-symmetry group \({\mathscr G}\) act on owners,
fields, background polarizations, and momenta.  Denote the
coordinate-complete first response to an owner \(o\) by
\(\Xi[o]\); it includes the stationary centre and Haar term of
(52.14).

\begin{theorem}[Exact orbit-sum reduction]
\label{thm:e8v56-orbit-sum}
Assume:
\begin{enumerate}
\item the undeformed compact action, block constraint, Haar
      measure, and retained zero-mode ledger are
      \({\mathscr G}\)-equivariant;
\item the momentum grid is \({\mathscr G}\)-stable; and
\item the response is extracted by (56.3).
\end{enumerate}
Then
\[
 \Xi[g o]=\Xi[o]\qquad(g\in{\mathscr G}),                     \tag{56.4}
\]
and
\[
 \boxed{
 \Xi\!\left[\sum_{o\in{\cal I}}o\right]
 =\sum_{{\cal O}\in{\cal I}/{\mathscr G}}
 |{\cal O}|\,\Xi[o_{\cal O}],}                               \tag{56.5}
\]
where \(o_{\cal O}\) is any representative.  The \(k=0\) loop
summand is retained in every representative card.
\end{theorem}

\begin{proof}
Change variables by \(g\) in the finite compact integral.  Haar
measure and the undeformed action are unchanged, the owner becomes
\(go\), and background directions and momenta are permuted.
The average (56.3) is invariant by
Lemma~\ref{lem:e8v56-symmetric-moment}.  Thus (56.4) holds.
The first derivative with respect to the deformation action is
linear, including the implicitly differentiated stationary centre,
so partitioning \({\cal I}\) into orbits gives (56.5).
Every lattice automorphism fixes the zero momentum, so the change
of variables permutes the complete momentum grid without deleting
that summand.
\end{proof}

\begin{corollary}[Transported trees are allowed, unmeasured trees
are not]
\label{cor:e8v56-transported-tree}
If \(g\) maps the fixed parity tree \(T\) to a different tree
\(gT\), (56.4) remains valid provided the slice measure is
transported.  The measured determinant is unchanged by
Theorem~\ref{thm:e8v43-slice-invariance} and
Corollary~\ref{cor:e8v43-background-jets}.
\end{corollary}

\begin{proof}
The two tree spaces are complements of the same gauge-orbit
directions after applying \(g\).  The V43 identity cancels the
coordinate determinant against the orbit Jacobian, and its
background-jet corollary applies to the differentiated response.
\end{proof}

\begin{firewall}[A fixed ordered tree is not hypercubic]
\label{fire:e8v56-fixed-tree}
The attached parity forest orders directions by repeatedly
removing the least odd coordinate.  A coordinate permutation
therefore changes the tree.  Applying a permutation only to an
owner while leaving the \(45\times45\) reduced coordinates fixed
does not justify (56.4).  One must transform the tree as in
Corollary~\ref{cor:e8v56-transported-tree}, or retain all affected
fixed-tree representatives.
\end{firewall}

\subsection{Exact permutation-orbit census inside one parity cell}

Index a raw completion owner by
\[
 (s;P;\rho),\qquad
 s\in\{0,1\}^4,\quad P=\{\mu,\nu\},\quad |P|=2,               \tag{56.6}
\]
where \(\rho\) is its plaquette-difference direction.  Put
\[
 c=\begin{cases}\parallel,&\rho\in P,\\
                  \perp,&\rho\notin P,\end{cases}\quad
 w=|S(s)|,\quad a=s_\rho,\quad b=|P\cap S(s)|.                \tag{56.7}
\]

\begin{theorem}[Complete \(S_4\)-orbit classification]
\label{thm:e8v56-orbit-classification}
Two triples (56.6) lie in the same coordinate-permutation orbit if
and only if their signatures \((c,w,a,b)\) agree.  There are
twelve parallel and twelve transverse orbits.  Their sizes are
\[
 |{\cal O}_{\parallel}(w,a,b)|
 =12\binom2{w-b},                                             \tag{56.8}
\]
\[
 |{\cal O}_{\perp}(w,a,b)|
 =12\binom2b,                                                 \tag{56.9}
\]
when the indicated bits are admissible, and zero otherwise.
The parallel sizes sum to \(192\), the transverse sizes sum to
\(192\), and the total is \(384\).
\end{theorem}

\begin{proof}
The signature is invariant under a simultaneous permutation of
the four directions.  Conversely, after matching \(\rho\), a
permutation can match the other direction in \(P\) in the parallel
case, or the two-element set \(P\) in the transverse case.  The
remaining directions can then be permuted to match their occupied
bits because \(w,a,b\) fix the number of ones in each distinguished
class.  Thus the signature is complete.

In the parallel case, choose the distinguished \(\rho\) in four
ways and the other element of \(P\) in three ways.  Its bit is
\(b-a\); among the two directions outside \(P\), choose \(w-b\)
ones.  This gives (56.8).  In the transverse case, choose \(\rho\)
in four ways and \(P\) from the other three directions in three
ways; choose its \(b\) occupied elements in \(\binom2b\) ways.
The last outside direction has the forced bit \(w-a-b\), proving
(56.9).  Summing the admissible cases gives the stated totals,
also equal to
\(16\cdot6\cdot2\) in each class.
\end{proof}

For reproducibility, the complete representatives and
multiplicities are:
\[
\begin{array}{c|c|c|c|c|c|c|c|c}
\#&c&s&P&\rho&w&a&b&|{\cal O}|\\ \hline
1&\parallel&0000&12&1&0&0&0&12\\
2&\parallel&0001&12&1&1&0&0&24\\
3&\parallel&0001&14&1&1&0&1&12\\
4&\parallel&0001&14&4&1&1&1&12\\
5&\parallel&0011&12&1&2&0&0&12\\
6&\parallel&0011&13&1&2&0&1&24\\
7&\parallel&0011&13&3&2&1&1&24\\
8&\parallel&0011&34&3&2&1&2&12\\
9&\parallel&0111&12&1&3&0&1&12\\
10&\parallel&0111&12&2&3&1&1&12\\
11&\parallel&0111&23&2&3&1&2&24\\
12&\parallel&1111&12&1&4&1&2&12\\ \hline
13&\perp&0000&12&3&0&0&0&12\\
14&\perp&0001&12&3&1&0&0&12\\
15&\perp&0001&12&4&1&1&0&12\\
16&\perp&0001&14&2&1&0&1&24\\
17&\perp&0011&12&3&2&1&0&12\\
18&\perp&0011&13&2&2&0&1&24\\
19&\perp&0011&13&4&2&1&1&24\\
20&\perp&0011&34&1&2&0&2&12\\
21&\perp&0111&12&3&3&1&1&24\\
22&\perp&0111&23&1&3&0&2&12\\
23&\perp&0111&23&4&3&1&2&12\\
24&\perp&1111&12&3&4&1&2&12
\end{array}                                                    \tag{56.10}
\]
Here the four bits of \(s\) and the direction labels are ordered
\(1,2,3,4\).

\begin{corollary}[Exact reduction of the two completion cards]
\label{cor:e8v56-completion-reduction}
Let \(\xi_\lambda\) be the coordinate-complete, symmetrically
extracted DQC value for representative row \(\lambda\) of
(56.10), computed with a consistently transported tree.  If
\(\nu_1,\nu_2\) are the fixed normalizations in
(52.5)--(52.6), then
\[
 \boxed{
 \Xi[{\cal O}_1]
 =\nu_1\sum_{\lambda=1}^{24}
 |{\cal O}_\lambda|\xi_\lambda,}                              \tag{56.11}
\]
\[
 \boxed{
 \Xi[{\cal O}_2]
 =\nu_2\sum_{\lambda=1}^{12}
 |{\cal O}_\lambda|\xi_\lambda.}                              \tag{56.12}
\]
Thus \(384\) and \(192\) placements reduce to twenty-four and
twelve exact quotient cards, respectively.  The same twelve
parallel cards occur in both completions.
\end{corollary}

\begin{proof}
The first completion contains all parallel and transverse owners;
the second contains only the parallel ones.  Apply
Theorems~\ref{thm:e8v56-orbit-sum} and
\ref{thm:e8v56-orbit-classification}, including the action
normalizations.
\end{proof}

\begin{remark}[Two word shapes, twenty-four quotient placements]
\label{rem:e8v56-shapes-vs-cards}
At raw word level there are only the parallel and transverse
geometric shapes.  Their position relative to the dyadic block
root and the retained fibre is encoded by \((w,a,b)\), producing
the twenty-four quotient placements in (56.10).  Conflating these
two levels would erase genuine block-boundary data.
\end{remark}

\subsection{Momentum and reflection checks}

\begin{proposition}[Fourier covariance of a representative]
\label{prop:e8v56-Fourier-covariance}
Let \(g\in S_4\), let \(U_g(k)\) be its unitary permutation-phase
map on link types, and transport the tree.  Every quadratic or
mixed-jet representative satisfies
\[
 {\cal M}_{go}(k;v_1,\ldots,v_n)
 =U_g(k)^*
 {\cal M}_o(g^{-1}k;g^{-1}v_1,\ldots,g^{-1}v_n)
 U_g(k).                                                     \tag{56.13}
\]
Orientation reversal gives the corresponding Hermitian conjugate.
Therefore the complete trace is real after pairing reversed
owners, and the momentum sum is unchanged.
\end{proposition}

\begin{proof}
Equation (54.2) is natural under relabeling every link in a word.
The parity-cell Fourier transform turns the relabeling into
permutation phases \(U_g(k)\), giving (56.13).  Cyclicity of trace
cancels the outer unitaries; \(k\mapsto g^{-1}k\) permutes the
finite grid.  Reversing a word replaces every holonomy by its
inverse in reverse order, which conjugates the Hessian and mixed
jets.  The paired trace is therefore real.
\end{proof}

\begin{firewall}[No toron deletion in the orbit compiler]
\label{fire:e8v56-toron}
The fixed point \(k=0\) of every coordinate permutation occurs in
each sum used to define \(\xi_\lambda\).  Orbit reduction acts on
owner placements and polarizations, not by deleting a momentum
orbit of size one.  Winding background owners remain separately in
the toron ledger.
\end{firewall}

\subsection{Status and the V57 acquisition}

\begin{assessment}[Progress after V56]
\label{ass:e8v56-status}
The symmetry reduction is now exact and compatible with the
ordered parity tree.  The raw owner bank has two geometric word
shapes but twenty-four inequivalent placements relative to a
dyadic block; their complete multiplicities are (56.10).
Equations (56.11)--(56.12) reduce the two local mixing numbers to
twenty-four representative finite cards, with the parallel twelve
shared.  No momentum or toron mode has been discarded.
\end{assessment}

\begin{decision}[V57 target]
\label{dec:e8v56-next}
V57 should instantiate the twelve parallel FMJCs first, since they
populate both \({\cal O}_1\) and \({\cal O}_2\).  For each row of
(56.10), the card must verify the quadratic residual against
(53.6), the mixed-derivative identities (54.9)--(54.11), and the
Ward zero before its trace contribution is accepted.  The
transverse twelve follow only after the shared bank passes.
\end{decision}

\begin{firewall}[Claim boundary after V56]
\label{fire:e8v56-claim}
V56 proves the exact symmetry and multiplicity reduction of the
dimension-six owner bank.  It does not execute the twenty-four
representative cards, evaluate either local mixing number, control
toron/flux/overlap owners, complete the LERB, establish the
universal beta coefficient, close the nonlinear/complex quotient
matching card, construct the full RG trajectory, remove the
regulators, establish Osterwalder--Schrader reconstruction, or
prove the Yang--Mills mass gap.
\end{firewall}

\subsection{Parent record}

This V56 note was formed by appending the preceding material to V55,
whose validated record was
\[
\begin{array}{c|c}
\text{lines}&24296\\
\text{bytes}&849360\\
\text{SHA-256}&
\texttt{A969A6E53DCD6554DF1049B84FA4239EF1A033513C08C8B09BD4B4837CEC0B38}.
\end{array}
\]

% =============================================================================
% END APPENDED EIGHTH-SERIES V56 MATERIAL
% =============================================================================

% =============================================================================
% BEGIN APPENDED EIGHTH-SERIES V57 MATERIAL
% =============================================================================

\section{Eighth-series V57: redundant equation-of-motion direction
and the hypercubic anisotropy direction}
\label{sec:e8v57}

\subsection{A more informative quartic basis}

In centered momentum variables let
\[
 {\cal P}_{\mu\nu}(p)
 :=S\delta_{\mu\nu}-\bar p_\mu\bar p_\nu.                    \tag{57.1}
\]
This is the exact lattice Maxwell precision in the transverse
channel.  Since \(\bar p^{\,T}\bar p=S\),
\[
 {\cal P}(p)^2=S{\cal P}(p)={\cal W}^{(1)}(p).                \tag{57.2}
\]
Define
\[
 {\cal W}^{\rm R}:={\cal W}^{(1)},\qquad
 {\cal W}^{\rm A}:={\cal W}^{(2)}
                   -\frac23{\cal W}^{(1)}.                   \tag{57.3}
\]
The superscripts denote redundant-candidate and anisotropic
directions.

\begin{proposition}[The isotropic quartic direction is an
equation-of-motion square]
\label{prop:e8v57-EOM-square}
For the quadratic kinetic action
\[
 S_0(a):=\frac12\langle a,{\cal P}a\rangle,
\]
the linear field redefinition
\[
 a\longmapsto a_\varepsilon
 :=a+\frac{\varepsilon}{2}{\cal P}a                           \tag{57.4}
\]
satisfies
\[
 S_0(a_\varepsilon)
 =S_0(a)+\frac{\varepsilon}{2}
 \langle a,{\cal W}^{\rm R}a\rangle+O(\varepsilon^2).         \tag{57.5}
\]
\end{proposition}

\begin{proof}
Self-adjointness of \({\cal P}\) gives
\[
 \left.\frac d{d\varepsilon}\right|_0S_0(a_\varepsilon)
 =\frac12\langle{\cal P}a,{\cal P}a\rangle
 =\frac12\langle a,{\cal P}^2a\rangle.
\]
Use (57.2).
\end{proof}

\subsection{Rotational projection of the second tensor}

Let \({\cal W}^{(r)}_{[4]}\) denote the homogeneous degree-four
Taylor tensor.  For such a tensor define its rotational average
\[
 ({\cal A}_{O(4)}H)(p)
 :=\int_{O(4)}R^TH(Rp)R\,dR,                                 \tag{57.6}
\]
with normalized Haar measure.

\begin{theorem}[Exact isotropic--anisotropic decomposition]
\label{thm:e8v57-rotational-decomposition}
The rotationally averaged quartic tensors obey
\[
 {\cal A}_{O(4)}{\cal W}^{(1)}_{[4]}
 ={\cal W}^{(1)}_{[4]},\qquad
 {\cal A}_{O(4)}{\cal W}^{(2)}_{[4]}
 =\frac23{\cal W}^{(1)}_{[4]},                               \tag{57.7}
\]
and hence
\[
 \boxed{
 {\cal A}_{O(4)}{\cal W}^{\rm A}_{[4]}=0.}                    \tag{57.8}
\]
\end{theorem}

\begin{proof}
An \(O(4)\)-covariant, symmetric, transverse, homogeneous
degree-four tensor is a scalar multiple of
\[
 |p|^2(|p|^2I-pp^T)={\cal W}^{(1)}_{[4]}.
\]
The first equality in (57.7) follows directly.  To find the second
coefficient, take traces.  In \(d=4\),
\[
 \Tr{\cal W}^{(1)}_{[4]}=3|p|^4,\qquad
 \Tr{\cal W}^{(2)}_{[4]}
 =|p|^4+2\sum_{\mu=1}^4p_\mu^4.                              \tag{57.9}
\]
On a rotation orbit,
\[
 {\cal A}_{O(4)}\sum_\mu p_\mu^4
 =\frac{3}{d+2}|p|^4=\frac12|p|^4,                           \tag{57.10}
\]
because
\(\mathbb E(\omega_\mu^4)=3/[d(d+2)]\) on the unit sphere.
Thus the averaged trace in the second formula of (57.9) is
\(2|p|^4\), which is \(2/3\) of the first trace.  This proves
(57.7), and (57.8) follows from (57.3).
\end{proof}

\begin{corollary}[Physical interpretation of the two coordinates]
\label{cor:e8v57-two-interpretations}
The local quartic Ward space splits as
\[
 {\cal Y}^{(4)}_{\rm loc}
 =\operatorname{span}\{{\cal W}^{\rm R}\}
  \mathbin{\dotplus}
  \operatorname{span}\{{\cal W}^{\rm A}\}.                   \tag{57.11}
\]
The first direction is tangent, at quadratic order, to a field
redefinition generated by the kinetic equation of motion.  The
second has zero rotational projection and measures a purely
hypercubic lattice anisotropy at fourth order.
\end{corollary}

\begin{proof}
The change of basis (57.3) is triangular with unit diagonal, so it
is invertible.  Apply Proposition~\ref{prop:e8v57-EOM-square} and
Theorem~\ref{thm:e8v57-rotational-decomposition}.
\end{proof}

\subsection{Exact measured redundancy identity}

Let \(X\) be a compact finite-dimensional manifold without
boundary, \(d\mu\) a smooth positive measure, and
\(\operatorname{div}_\mu V\) the divergence of a smooth vector
field \(V\) relative to \(d\mu\).

\begin{theorem}[Measured field-redefinition owner]
\label{thm:e8v57-measured-redundancy}
For every smooth action \(S\),
\[
 {\cal O}_V:=DS[V]-\operatorname{div}_\mu V                   \tag{57.12}
\]
obeys
\[
 \boxed{\int_X{\cal O}_V e^{-S}\,d\mu=0.}                    \tag{57.13}
\]
If \(S=S(\,\cdot\,;W)\) depends smoothly on a retained background
\(W\), and \(V(\,\cdot\,;W)\) is tangent to the integration fibre
for every \(W\), then every background derivative of (57.13)
also vanishes.
\end{theorem}

\begin{proof}
The divergence theorem on a compact manifold without boundary
gives
\[
 0=\int_X\operatorname{div}_\mu(e^{-S}V)\,d\mu
 =\int_Xe^{-S}
 \bigl(\operatorname{div}_\mu V-DS[V]\bigr)\,d\mu,
\]
which is (57.13).  If the identity holds for every \(W\) on one
smooth fibre bundle and all data depend smoothly on \(W\),
differentiate the zero function with respect to \(W\).
\end{proof}

\begin{corollary}[A true redundant DQC vanishes]
\label{cor:e8v57-redundant-DQC}
On the compact reduced-link fibre with product Haar measure, an
owner of the form (57.12), with \(V\) vertical for the exact block
constraint and with all retained zero modes held fixed, has
\[
 \Xi[{\cal O}_V]=0.                                          \tag{57.14}
\]
\end{corollary}

\begin{proof}
Theorem~\ref{thm:e8v57-measured-redundancy} says that its first
deformation derivative of the fibre integral is zero for every
retained background.  Take the coordinate-complete background
Hessian and then the marginal extractor.
\end{proof}

\begin{remark}[The divergence term is the Jacobian]
\label{rem:e8v57-divergence-Jacobian}
The two terms in (57.12) are inseparable.  \(DS[V]\) is the action
change and \(-\operatorname{div}_\mu V\) is the infinitesimal
Jacobian of the measured coordinate change.  Dropping the latter
can manufacture a nonzero mixing coefficient for a redundant
direction.
\end{remark}

\subsection{The verticality obstruction}

Let \({\cal B}\) denote the exact compact block map and let
\({\cal V}_{\rm EOM}\) be a gauge-covariant nonlinear completion
whose linearization at the vacuum is
\(\frac12{\cal P}a\).

\begin{definition}[Redundancy completion card]
\label{def:e8v57-RCC}
A redundancy completion card (RCC) consists of:
\begin{enumerate}
\item a local analytic vector field
      \({\cal V}_{\rm EOM}\) on fine links;
\item the exact verticality identity
      \[
       D{\cal B}(U)[{\cal V}_{\rm EOM}(U)]=0;                 \tag{57.15}
      \]
\item its Haar divergence;
\item equality of the quadratic Hessian of
      \(DS[{\cal V}_{\rm EOM}]-\operatorname{div}_{\rm Haar}
       {\cal V}_{\rm EOM}\)
      with \({\cal W}^{\rm R}\); and
\item a local comparison between that measured owner and the
      completion \({\cal O}_1\).
\end{enumerate}
\end{definition}

\begin{firewall}[Quadratic redundancy does not prove fibre
redundancy]
\label{fire:e8v57-verticality}
Equation (57.5) is a flat linear statement.  The natural covariant
equation-of-motion vector field need not be tangent to the
nonlinear block fibre, and projecting it vertically can introduce
a nonlocal block inverse.  Therefore (57.14) cannot yet be applied
to \({\cal O}_1\); the missing identity is the RCC, especially
(57.15).
\end{firewall}

\subsection{Change of completion basis}

Define the explicit anisotropy completion
\[
 {\cal O}_{\rm A}:={\cal O}_2-\frac23{\cal O}_1.              \tag{57.16}
\]
Its quadratic Hessian is \({\cal W}^{\rm A}\).  If an RCC owner
\({\cal O}_{\rm R}\) is constructed, define the completion defect
\[
 {\cal E}_{\rm comp}:={\cal O}_1-{\cal O}_{\rm R}.            \tag{57.17}
\]

\begin{proposition}[DQC reduction in the new basis]
\label{prop:e8v57-DQC-basis}
Whenever \({\cal O}_{\rm R}\) satisfies the RCC,
\[
 \Xi[{\cal O}_1]=\Xi[{\cal E}_{\rm comp}],                    \tag{57.18}
\]
\[
 \Xi[{\cal O}_2]
 =\frac23\Xi[{\cal E}_{\rm comp}]
  +\Xi[{\cal O}_{\rm A}].                                    \tag{57.19}
\]
\end{proposition}

\begin{proof}
Linearity of the first deformation response gives
\(\Xi[{\cal O}_1]
=\Xi[{\cal O}_{\rm R}]+\Xi[{\cal E}_{\rm comp}]\).
The first term vanishes by (57.14), proving (57.18).
Equation (57.16) gives
\({\cal O}_2=\frac23{\cal O}_1+{\cal O}_{\rm A}\); apply
(57.18).
\end{proof}

\begin{firewall}[Rotational averaging is not a finite-lattice
cancellation]
\label{fire:e8v57-anisotropy}
Equation (57.8) says that the leading continuum rotational
projection of \({\cal O}_{\rm A}\) vanishes.  It does not say that
its finite-lattice DQC is zero.  A quantitative symmetry-restoration
estimate is needed to show that its mixing into \(F^2\) contracts
cofinally.
\end{firewall}

\subsection{Status and the V58 acquisition}

\begin{assessment}[Progress after V57]
\label{ass:e8v57-status}
The two dimension-six coordinates now have distinct structural
roles.  One is an equation-of-motion square at quadratic order and
would have identically zero DQC if completed by a vertical measured
field redefinition.  The other has zero \(O(4)\) projection and is
a pure hypercubic anisotropy.  This explains which cancellations
the twenty-four cards must test: vertical Jacobian cancellation for
the redundant direction and cofinal rotational restoration for the
anisotropic direction.
\end{assessment}

\begin{decision}[V58 target]
\label{dec:e8v57-next}
V58 should test the natural covariant EOM vector field against the
exact straight-holonomy block map.  Its horizontal defect
\(D{\cal B}[{\cal V}_{\rm EOM}]\) must be computed first.  If
nonzero, split the vector field into a vertical component and an
explicit retained-field coordinate change, carrying both
Jacobians.  In parallel, formulate the anisotropy response as the
difference between a card and its rotational average; no claim of
vanishing is allowed without a quantitative bound.
\end{decision}

\begin{firewall}[Claim boundary after V57]
\label{fire:e8v57-claim}
V57 proves the redundant/anisotropic quartic basis, the exact
measured field-redefinition identity, and the conditional DQC
reduction.  It does not complete the RCC, prove rotational
restoration, evaluate the remaining cards, control
toron/flux/overlap owners, complete the LERB, establish the
universal beta coefficient, close the nonlinear/complex quotient
matching card, construct the full RG trajectory, remove the
regulators, establish Osterwalder--Schrader reconstruction, or
prove the Yang--Mills mass gap.
\end{firewall}

\subsection{Parent record}

This V57 note was formed by appending the preceding material to V56,
whose validated record was
\[
\begin{array}{c|c}
\text{lines}&24641\\
\text{bytes}&862304\\
\text{SHA-256}&
\texttt{31498B878E30871E4E6FFF2760F12B168CBDB518906AAAEAD2417072F9018301}.
\end{array}
\]

% =============================================================================
% END APPENDED EIGHTH-SERIES V57 MATERIAL
% =============================================================================

% =============================================================================
% BEGIN APPENDED EIGHTH-SERIES V58 MATERIAL
% =============================================================================

\section{Eighth-series V58: pushforward of measured redundant
owners and failure of strict fibre verticality}
\label{sec:e8v58}

\subsection{The natural EOM vector is not vertical}

At the linearized level, straight dyadic blocking of a fine link
field is
\[
 ({\cal B}_0a)_\mu(y)
 =a_\mu(2y)+a_\mu(2y+\widehat e_\mu).                         \tag{58.1}
\]
For a plane wave \(a_\mu(x)=A_\mu e^{ip\cdot x}\),
\[
 ({\cal B}_0a)_\mu(y)
 =(1+e^{ip_\mu})A_\mu e^{i(2p)\cdot y}.                       \tag{58.2}
\]

\begin{proposition}[Explicit nonvertical EOM direction]
\label{prop:e8v58-EOM-not-vertical}
Let \({\cal P}(p)\) be (57.1).  In general
\[
 {\cal B}_0{\cal P}\ne0.                                     \tag{58.3}
\]
For example, take
\[
 p=(\pi/2,0,0,0),\qquad A=\widehat e_2.
\]
Then \(A\) is transverse, \(S(p)=2\), and
\[
 {\cal B}_0({\cal P}A)_2=4e^{i(2p)\cdot y}\ne0.               \tag{58.4}
\]
\end{proposition}

\begin{proof}
The chosen polarization obeys
\(\bar p^{\,T}A=0\), so \({\cal P}(p)A=S(p)A=2A\).
Its second component has block factor
\(1+e^{ip_2}=2\) in (58.2), giving (58.4).
\end{proof}

\begin{corollary}[Correction to the V57 acquisition]
\label{cor:e8v58-verticality-correction}
The strict verticality row (57.15) cannot hold for a nonlinear EOM
vector whose linearization is \(\frac12{\cal P}a\).  Verticality is
therefore too strong a definition of redundancy for the exact
block pushforward.
\end{corollary}

\begin{proof}
Differentiating (57.15) at the vacuum would give
\({\cal B}_0{\cal P}=0\), contradicting (58.4).
\end{proof}

\subsection{A general pushforward continuity theorem}

Let \(X,Y\) be compact oriented smooth finite-dimensional
manifolds without boundary, with smooth positive measures
\(d\mu_X,d\mu_Y\), and let
\({\cal B}:X\to Y\) be a smooth submersion.  For a smooth action
\(S\), define the pushed density \(\rho\) by
\[
 \rho\,d\mu_Y:={\cal B}_*(e^{-S}d\mu_X).                      \tag{58.5}
\]
For a smooth vector field \(V\) on \(X\), define its pushed current
\(J\), a vector density on \(Y\), weakly by
\[
 \int_Y\alpha(J)\,d\mu_Y
 :=\int_X e^{-S(x)}
 \alpha_{{\cal B}(x)}\!\left(D{\cal B}(x)V(x)\right)
 d\mu_X(x)                                                    \tag{58.6}
\]
for every smooth one-form \(\alpha\) on \(Y\).

\begin{theorem}[Measured redundancy commutes with pushforward]
\label{thm:e8v58-pushforward-redundancy}
Let
\[
 {\cal O}_V:=DS[V]-\operatorname{div}_{\mu_X}V.               \tag{58.7}
\]
Then, as a distributional density on \(Y\),
\[
 \boxed{
 {\cal B}_*({\cal O}_Ve^{-S}d\mu_X)
 =-(\operatorname{div}_{\mu_Y}J)d\mu_Y.}                      \tag{58.8}
\]
If \(\rho>0\) and
\[
 \overline V:=J/\rho,                                        \tag{58.9}
\]
then the first variation of the effective action
\[
 A:=-\log\rho
\]
under the fine action deformation \(S+\varepsilon{\cal O}_V\) is
\[
 \boxed{
 \dot A
 =DA[\overline V]-\operatorname{div}_{\mu_Y}\overline V.}     \tag{58.10}
\]
Thus a measured redundant tangent pushes to a measured redundant
tangent even when \(V\) is not vertical.
\end{theorem}

\begin{proof}
For a smooth test function \(f\) on \(Y\), apply the divergence
theorem on \(X\) to \(e^{-S}(f\circ{\cal B})V\):
\[
\begin{split}
 0
 &=\int_X\operatorname{div}_{\mu_X}
 \bigl(e^{-S}(f\circ{\cal B})V\bigr)d\mu_X\\
 &=\int_Xe^{-S}\left[
 -(f\circ{\cal B}){\cal O}_V
 +Df_{{\cal B}(x)}(D{\cal B}(x)V(x))\right]d\mu_X.
                                                                    \tag{58.11}
\end{split}
\]
The second term is \(\int_YDf(J)d\mu_Y\), which equals
\(-\int_Yf\operatorname{div}_{\mu_Y}J\,d\mu_Y\).
This is the weak form of (58.8).

Under \(S+\varepsilon{\cal O}_V\), (58.8) gives
\[
 \dot\rho
 =-\rho\,\mathbb E({\cal O}_V\mid{\cal B}=y)
 =\operatorname{div}_{\mu_Y}J.                               \tag{58.12}
\]
Since \(J=\rho\overline V\) and \(D\rho=-\rho\,DA\),
\[
 -\frac{\dot\rho}{\rho}
 =-\frac{\operatorname{div}(\rho\overline V)}{\rho}
 =DA[\overline V]-\operatorname{div}\overline V,
\]
which proves (58.10).
\end{proof}

\begin{corollary}[Redundant tangents form an exact RG ideal]
\label{cor:e8v58-redundant-ideal}
At every finite regulator for which the block map is a smooth
submersion and the pushed density is positive, the exact RG
differential maps measured field-redefinition tangents into
measured field-redefinition tangents.  Consequently it induces a
well-defined differential on the quotient of action tangents by
the subspace of measured redundant tangents.
\end{corollary}

\begin{proof}
Equation (58.10) has exactly the form (58.7), now on the retained
space \(Y\).  Therefore equivalent fine tangents have equivalent
coarse tangents, which is precisely the condition for a map on the
quotient.
\end{proof}

\begin{remark}[Compact-group formulation]
\label{rem:e8v58-compact-group}
For the exact regulator, \(X\) is the product of fine-link copies
of \(SU(3)\), \(Y\) the retained-link product, and both measures are
product Haar.  Formula (58.6) is defined without choosing a
logarithm.  Local charts are needed only to write the vector
components and divergences, so no branch choice enters (58.8).
\end{remark}

\subsection{Gaussian conditional current}

The theorem becomes explicit for a linear block and a Gaussian
reference.  Let \(P>0\) on a finite real vector space \(E\), let
\(B:E\to C\) be onto, and put
\[
 S(x)=\frac12\langle x,Px\rangle,\qquad
 H:=P^{-1}B^*(BP^{-1}B^*)^{-1},\qquad
 S_C:=(BP^{-1}B^*)^{-1}.                                    \tag{58.13}
\]
Then \(BH=I\) and the conditional mean is
\(\mathbb E(x\mid Bx=c)=Hc\).

\begin{proposition}[Coarse velocity of the linear EOM flow]
\label{prop:e8v58-Gaussian-current}
For
\[
 V(x)=\frac12Px,                                              \tag{58.14}
\]
the coarse velocity (58.9) is
\[
 \boxed{
 \overline V(c)=L_Cc,\qquad
 L_C:=\frac12BPH=\frac12BB^*S_C.}                            \tag{58.15}
\]
The pushed redundant owner is
\[
 {\cal O}_{\overline V}(c)
 =\langle S_Cc,L_Cc\rangle-\Tr L_C.                           \tag{58.16}
\]
Its constant term is a vacuum contact.  If \(S_C(p)=O(|p|^2)\)
and \(BB^*(p)=O(1)\) near zero, its quadratic symbol is
\(O(|p|^4)\), so its second moment vanishes.
\end{proposition}

\begin{proof}
The conditional expectation of (58.14) is
\(\frac12PHc\).  Applying \(B\) gives the first expression in
(58.15).  The definition of \(H\) gives
\(PH=B^*S_C\), proving the second.  The Gaussian effective action
is \(\frac12\langle c,S_Cc\rangle\) up to a constant.
Insert \(\overline V=L_Cc\) in (58.10); the divergence of a linear
field is \(\Tr L_C\), proving (58.16).
Finally \(L_C=\frac12BB^*S_C=O(|p|^2)\), so the symmetric
quadratic kernel in the first term of (58.16) is
\(\frac12(S_CL_C+L_C^*S_C)=O(|p|^4)\).
\end{proof}

\begin{corollary}[No Gaussian marginal feed from the EOM
direction]
\label{cor:e8v58-no-Gaussian-feed}
For the exact Gaussian block reference, the equation-of-motion
field redefinition pushes to a quartic-momentum redundant owner and
has zero \(F^2\) moment.  Its only zeroth-order contribution is the
vacuum Jacobian contact \(-\Tr L_C\).
\end{corollary}

\subsection{Replacement of the strict RCC}

V57 required a vertical vector field.  The pushforward theorem
shows that the correct object is a current card.

\begin{definition}[Pushforward redundancy card]
\label{def:e8v58-PRC}
A pushforward redundancy card (PRC) for the isotropic
dimension-six direction consists of:
\begin{enumerate}
\item a local gauge-covariant fine vector field \(V_{\rm EOM}\)
      whose linearization is \(\frac12{\cal P}a\);
\item its complete Haar divergence and measured owner
      \({\cal O}_{V_{\rm EOM}}\);
\item the conditional current \(J\) in (58.6), without deleting
      toron sectors;
\item smoothness and locality bounds for
      \(\overline V=J/\rho\);
\item proof that the retained redundant owner (58.10) begins at
      fourth external-momentum order after its vacuum contact is
      removed; and
\item a local comparison of
      \({\cal O}_1-{\cal O}_{V_{\rm EOM}}\).
\end{enumerate}
\end{definition}

\begin{theorem}[Corrected redundant-direction reduction]
\label{thm:e8v58-corrected-reduction}
If the PRC is complete, then the exact RG image of
\({\cal O}_{V_{\rm EOM}}\) has zero marginal \(F^2\) projection
and is retained in the redundant-coordinate ledger.  The only
physical local DQC left from \({\cal O}_1\) is its completion
difference
\[
 {\cal E}_{\rm comp}
 :={\cal O}_1-{\cal O}_{V_{\rm EOM}}.                         \tag{58.17}
\]
\end{theorem}

\begin{proof}
Theorem~\ref{thm:e8v58-pushforward-redundancy} identifies the RG
image as the measured redundant owner (58.10).  PRC item 5 makes
its second moment zero, and item 4 places it in the local
coordinate-change ledger.  Linearity of the differentiated RG
response splits \({\cal O}_1\) into that owner and (58.17).
\end{proof}

\begin{firewall}[Redundancy is not numerical zero in a fixed
coordinate section]
\label{fire:e8v58-redundancy-section}
When \(V\) is not vertical, its fine insertion need not give a
zero fixed-coordinate DQC; it gives the coarse redundant owner
(58.10).  It disappears only after passing to the quotient by
measured field redefinitions, or after transporting the retained
coordinate and its Jacobian.  Setting its raw card to zero would
repeat the measure error corrected in V41--V43.
\end{firewall}

\subsection{Status and the V59 acquisition}

\begin{assessment}[Progress after V58]
\label{ass:e8v58-status}
The strict vertical RCC is disproved at linear order by (58.4) and
replaced by the exact pushforward continuity theorem.  Measured
redundant tangents form an RG ideal, and the Gaussian conditional
current proves that the EOM direction has no marginal feed at the
reference level.  The nonlinear remaining acquisition is no longer
verticality: it is locality and derivative order of the conditional
coarse current, plus the explicit completion difference (58.17).
\end{assessment}

\begin{decision}[V59 target]
\label{dec:e8v58-next}
V59 should derive a conditional-current locality estimate from the
uniform vertical gap and the weighted covariance machinery of V45.
The numerator and denominator defining \(J/\rho\) must be expanded
on the same fibre so vacuum-volume factors cancel before division.
If a uniform positive lower bound for the conditional density is
unavailable, use connected cumulants of the current rather than a
pointwise ratio.
\end{decision}

\begin{firewall}[Claim boundary after V58]
\label{fire:e8v58-claim}
V58 proves the general pushforward redundancy theorem and its
Gaussian Yang--Mills consequence.  It does not complete the
nonlinear PRC, bound the conditional current, prove anisotropy
contraction, evaluate the remaining cards, control
toron/flux/overlap owners, complete the LERB, establish the
universal beta coefficient, close the nonlinear/complex quotient
matching card, construct the full RG trajectory, remove the
regulators, establish Osterwalder--Schrader reconstruction, or
prove the Yang--Mills mass gap.
\end{firewall}

\subsection{Parent record}

This V58 note was formed by appending the preceding material to V57,
whose validated record was
\[
\begin{array}{c|c}
\text{lines}&24973\\
\text{bytes}&873450\\
\text{SHA-256}&
\texttt{26321C107687AF284575E1BDD6B7DB35DD01CB45B3378DCB3E3001292A464522}.
\end{array}
\]

% =============================================================================
% END APPENDED EIGHTH-SERIES V58 MATERIAL
% =============================================================================

% =============================================================================
% BEGIN APPENDED EIGHTH-SERIES V59 MATERIAL
% =============================================================================

\section{Eighth-series V59: connected-current locality without a
volume denominator}
\label{sec:e8v59}

\subsection{Normalized current and exact score identities}

On the reduced fibre over a retained background \(W\), write
\[
 d\nu_W(z):=Z(W)^{-1}e^{-S(z,W)}\,dm_{\rm red}(z),\qquad
 \langle F\rangle_W:=\int F\,d\nu_W.                         \tag{59.1}
\]
For a fine vector field \(V\), put
\[
 j(z,W):=D{\cal B}(z,W)[V(z,W)].                              \tag{59.2}
\]
Then the coarse velocity in (58.9) is simply
\[
 \overline V(W)=\langle j\rangle_W.                          \tag{59.3}
\]
No lower bound on the unnormalized partition function appears in
(59.3).

\begin{lemma}[First and second score derivatives]
\label{lem:e8v59-score-derivatives}
For retained background directions \(a,b\), let
\(S_a=D_aS\), \(S_{ab}=D_aD_bS\).  Then
\[
 D_a\langle F\rangle
 =\langle F_a\rangle-\kappa(F,S_a),                           \tag{59.4}
\]
and
\[
\boxed{\begin{split}
 D_bD_a\langle F\rangle
 ={}&\langle F_{ab}\rangle
 -\kappa(F_a,S_b)-\kappa(F_b,S_a)-\kappa(F,S_{ab})\\
 &+\kappa(F,S_a,S_b),
\end{split}}                                                  \tag{59.5}
\]
where \(\kappa\) denotes the connected cumulant under \(\nu_W\).
\end{lemma}

\begin{proof}
Differentiate the normalized quotient in (59.1):
\[
 D_a\langle F\rangle
 =\langle F_a-FS_a\rangle+\langle F\rangle\langle S_a\rangle,
\]
which is (59.4).  For two observables,
\[
 D_b\kappa(F,G)
 =\kappa(F_b,G)+\kappa(F,G_b)-\kappa(F,G,S_b).                \tag{59.6}
\]
This follows by differentiating
\(\langle FG\rangle-\langle F\rangle\langle G\rangle\) and using
(59.4).  Differentiate (59.4), apply (59.6) with \(G=S_a\), and
collect terms to obtain (59.5).
\end{proof}

\begin{corollary}[Current jets use connected quantities]
\label{cor:e8v59-current-jets}
The first two background jets of the conditional current
\(\overline V=\langle j\rangle\) are given by (59.4)--(59.5) with
\(F=j\).  Vacuum-volume factors cancel inside connected cumulants
before any spatial estimate is taken.
\end{corollary}

\begin{firewall}[Do not bound numerator and denominator
separately]
\label{fire:e8v59-no-ratio-bound}
A lower bound on \(Z(W)\) typically decays exponentially with
volume and is useless for locality of \(J/\rho\).  Equations
(59.3)--(59.5) are identities for normalized expectations and
contain only connected responses.  Estimating the two
unnormalized quantities first would destroy this cancellation.
\end{firewall}

\subsection{A connected-cumulant locality compiler}

For finite supports \(X_0,\ldots,X_n\), let
\(\tau(X_0,\ldots,X_n)\) be the length of a shortest lattice tree
meeting all supports.

\begin{theorem}[Connected-current locality]
\label{thm:e8v59-current-locality}
Assume that every local component \(j_X\) of \(j\) and every local
background score component \(S_{a,Y}\) obeys, for \(0\le n\le2\),
\[
 \left|
 \kappa(F_{0,X_0},F_{1,X_1},\ldots,F_{n,X_n})
 \right|
 \le C_n\prod_{\ell=0}^n\|F_{\ell,X_\ell}\|_{\rm loc}
 e^{-\mu\tau(X_0,\ldots,X_n)},                               \tag{59.7}
\]
including the derivative observables occurring in (59.4)--(59.5).
Assume also a bounded number of owner supports through every site.
Then the kernels of
\[
 \overline V,\qquad D\overline V,\qquad D^2\overline V        \tag{59.8}
\]
are exponentially summable in every weight
\(0<\mu'<\mu\), with constants independent of volume.
\end{theorem}

\begin{proof}
Insert the local decompositions of \(j,S_a,S_{ab}\) into
(59.3)--(59.5).  Each term is a connected cumulant with at most
three marked supports and therefore has the tree bound (59.7).
Fix one output support and one or two input sites.  Summing over
intermediate owner supports costs a lattice convolution of
exponentials.  For \(\mu'<\mu\),
\[
 \sum_z e^{-\mu|x-z|_1}e^{-\mu|z-y|_1}
 \le C_{\mu-\mu'}e^{-\mu'|x-y|_1}.                           \tag{59.9}
\]
The same tree-pruning argument handles three supports.  The
bounded incidence number absorbs choices of owners through a
site.  Thus each kernel in (59.8) has a volume-independent
\(\mu'\)-weighted sum.
\end{proof}

\begin{corollary}[Polymer-gate version]
\label{cor:e8v59-polymer-current}
The V55 bounded-box curvature compiler and any polymer measure
satisfying the connected clustering bound (59.7) yield a local
coarse velocity \(\overline V\) with two controlled background
jets.  Toron and winding components must first be removed from the
contractible owner bank.
\end{corollary}

\subsection{Verification at the Gaussian one-loop reference}

\begin{lemma}[Connected Wick graph bound]
\label{lem:e8v59-Wick-bound}
Let \(\xi\) be a centered Gaussian field with covariance \(G\)
satisfying \(\|G\|_\mu\le g\).  Let
\(F_0,\ldots,F_n\) be local polynomials of degrees at most \(D\),
with fixed coefficient \(\ell^1\) norms and \(n\le2\).  Then their
connected cumulant is the sum of Wick contractions whose graph on
the marked supports is connected.  Consequently it satisfies
(59.7), with a volume-independent constant depending only on
\(D,g,\mu\), the coefficient norms, and the support sizes.
\end{lemma}

\begin{proof}
Introduce sources \(t_0,\ldots,t_n\) and take
\(\partial_{t_0}\cdots\partial_{t_n}\) of the logarithm of the
Gaussian moment-generating function for the polynomial
insertions.  The logarithm selects connected Wick graphs.  Every
connected graph contains a spanning tree between the marked
supports.  Bound propagators on that tree by their
\(\mu\)-weighted kernel and all remaining propagators by \(g\).
There are finitely many pairings because the degrees are bounded.
Successive weighted convolutions along the tree give (59.7).
\end{proof}

\begin{theorem}[One-loop conditional-current locality]
\label{thm:e8v59-one-loop-current}
For the finite-degree current and action jets entering the
dimension-six PRC, the Gaussian vertical reference with the V41
spectral floor satisfies the hypotheses of
Theorem~\ref{thm:e8v59-current-locality}.  Hence the zeroth and
one-loop conditional-current jets have exponentially summable
kernels, uniformly in volume.
\end{theorem}

\begin{proof}
The block differential, EOM vector, Wilson score vertices, and
the completion words have fixed finite supports and finite degrees
at every required perturbative order.  V45 supplies the weighted
covariance \(G\), while Lemma~\ref{lem:e8v59-Wick-bound} supplies
the connected clustering bound for each finite list of
insertions.  Apply Theorem~\ref{thm:e8v59-current-locality}.
\end{proof}

\begin{firewall}[A stationary Hessian gap is not a nonlinear
clustering theorem]
\label{fire:e8v59-gap-not-clustering}
Theorem~\ref{thm:e8v59-one-loop-current} is a Gaussian/one-loop
statement.  A positive Hessian at the vacuum alone does not imply
(59.7) for the full compact interacting fibre.  The nonlinear PRC
requires either a uniform log-Sobolev/cluster theorem or the
small-activity polymer gate; large-field suppression must be
included.
\end{firewall}

\subsection{Ward structure of the pushed current}

\begin{proposition}[Equivariance of the conditional current]
\label{prop:e8v59-current-equivariance}
If the fine action, block map, measure, and vector field \(V\) are
gauge covariant, then the current \(J\), density \(\rho\), and
coarse velocity \(\overline V=J/\rho\) are gauge covariant.
The same statement holds for translations, reflections, and
hypercubic transformations preserved by the regulator.
\end{proposition}

\begin{proof}
Apply a symmetry as a change of variables in (58.5)--(58.6).
Covariance of \(D{\cal B}[V]\) transports the vector index of
\(J\), while invariance of the action and measure leaves its
weight unchanged.  Division by the scalar invariant density
\(\rho\) preserves covariance.
\end{proof}

\begin{proposition}[Quadratic derivative order of a local EOM
current]
\label{prop:e8v59-current-order}
Assume the contractible part of \(\overline V\) is exponentially
local, reflection even, and has no constant or one-derivative
gauge-covariant vector term.  Then its linearization \(L(p)\) at
the vacuum satisfies
\[
 L(p)d(p)=0,\qquad L(p)=O(|p|^2).                             \tag{59.10}
\]
If the effective-action Hessian is \(K(p)=O(|p|^2)\), the
quadratic part of \(DA[\overline V]\) is \(O(|p|^4)\).
\end{proposition}

\begin{proof}
Gauge covariance differentiated at the vacuum annihilates a pure
gauge gradient, giving the first identity.  Exponential locality
makes \(L(p)\) analytic.  Contractibility removes a constant
toron owner, reflection removes the odd Taylor term, and the stated
absence hypothesis leaves order two.  The quadratic Hessian of
\(DA[\overline V]\) is the symmetrization
\(K L+L^*K\), plus terms containing the vanished first derivative
of \(A\); it is therefore \(O(|p|^4)\).
\end{proof}

\subsection{The Jacobian marginal is a separate row}

The full pushed redundant owner is
\[
 {\cal O}_{\overline V}
 =DA[\overline V]-\operatorname{div}_{\rm Haar}\overline V.   \tag{59.11}
\]
Even when the first term starts at \(p^4\), a nonlinear Haar
divergence may have a local \(F^2\) Hessian.

\begin{definition}[Conditional-current Jacobian card]
\label{def:e8v59-CCJC}
The conditional-current Jacobian card (CCJC) records
\[
 K_{\rm grad}:=D_W^2(DA[\overline V])\big|_{W=I},\qquad
 K_{\rm Jac}:=-D_W^2\operatorname{div}_{\rm Haar}
 \overline V\big|_{W=I},                                    \tag{59.12}
\]
\[
 \zeta_{\rm grad}:={\mathfrak m}(K_{\rm grad}),\qquad
 \zeta_{\rm Jac}:={\mathfrak m}(K_{\rm Jac}).                 \tag{59.13}
\]
The redundant marginal coordinate is their sum.
\end{definition}

\begin{theorem}[Local decomposition of the pushed redundant
owner]
\label{thm:e8v59-redundant-decomposition}
Under Theorem~\ref{thm:e8v59-current-locality} and the exact
symmetry hypotheses of V50, the contractible Hessian
\[
 K_{\rm red}:=K_{\rm grad}+K_{\rm Jac}
\]
has the unique decomposition
\[
 \boxed{
 K_{\rm red}
 =(\zeta_{\rm grad}+\zeta_{\rm Jac}){\cal W}
  +{\cal Q}K_{\rm red},\qquad
 {\mathfrak m}({\cal Q}K_{\rm red})=0.}                       \tag{59.14}
\]
If Proposition~\ref{prop:e8v59-current-order} applies, then
\(\zeta_{\rm grad}=0\).  At the Gaussian reference,
\(\zeta_{\rm Jac}=0\) as well.
\end{theorem}

\begin{proof}
Current locality makes both kernels exponentially summable.
Gauge covariance and reflection give the Ward and hypercubic
hypotheses of Theorem
\ref{thm:e8v50-quadratic-classification}.  The moment projection
(48.3) therefore gives (59.14).  The \(O(p^4)\) conclusion of
Proposition~\ref{prop:e8v59-current-order} makes the first moment
zero.  In the Gaussian formula (58.16), the divergence is the
background-independent constant \(-\Tr L_C\), whose Hessian
vanishes.
\end{proof}

\begin{firewall}[A redundant Jacobian coefficient is not silently
dropped]
\label{fire:e8v59-Jacobian-marginal}
The coefficient \(\zeta_{\rm Jac}\) is a coordinate-Jacobian row.
It may be irrelevant to observables after quotienting by measured
field redefinitions, but it contributes in a fixed action
coordinate.  Universal beta matching must either include it or
prove invariance under the declared coordinate change.
\end{firewall}

\subsection{Status and the V60 audit target}

\begin{assessment}[Progress after V59]
\label{ass:e8v59-status}
The conditional current no longer requires a volume-uniform lower
bound on a partition density.  Its first two jets are exact
connected cumulants, and the weighted Gaussian covariance proves
one-loop exponential locality.  The pushed EOM direction has
quartic action-gradient order; the only possible fixed-coordinate
marginal feed is isolated as the Haar-Jacobian coefficient
\(\zeta_{\rm Jac}\).  Full nonlinear current locality remains tied
to a genuine clustering or polymer theorem.
\end{assessment}

\begin{decision}[Mandatory V60 audit and next target]
\label{dec:e8v59-next}
V60 must inspect the current SM and NS companions.  It should look
specifically for a normalized-current, logarithmic-derivative, or
connected-polymer theorem that can strengthen the nonlinear version
of (59.7).  The Yang--Mills calculation should then derive the Haar
divergence of a concrete covariant EOM vector field and determine
whether \(\zeta_{\rm Jac}\) vanishes or is a removable marginal
coordinate term.
\end{decision}

\begin{firewall}[Claim boundary after V59]
\label{fire:e8v59-claim}
V59 proves the connected-current identities and one-loop locality,
and isolates the nonlinear Haar-Jacobian marginal.  It does not
prove full interacting clustering, evaluate
\(\zeta_{\rm Jac}\), prove anisotropy contraction, evaluate the
remaining cards, control toron/flux/overlap owners, complete the
LERB, establish the universal beta coefficient, close the
nonlinear/complex quotient matching card, construct the full RG
trajectory, remove the regulators, establish
Osterwalder--Schrader reconstruction, or prove the Yang--Mills mass
gap.
\end{firewall}

\subsection{Parent record}

This V59 note was formed by appending the preceding material to V58,
whose validated record was
\[
\begin{array}{c|c}
\text{lines}&25309\\
\text{bytes}&885221\\
\text{SHA-256}&
\texttt{9AEFB8DF212D760E636B0883FAD3817B5519C480E0649206FB5AFE79E50E7D31}.
\end{array}
\]

% =============================================================================
% END APPENDED EIGHTH-SERIES V59 MATERIAL
% =============================================================================

% =============================================================================
% BEGIN APPENDED EIGHTH-SERIES V60 MATERIAL
% =============================================================================

\section{Eighth-series V60: compulsory companion audit, the third-jet
debit, and a Haar-compensated EOM coordinate}
\label{sec:e8v60}

\subsection{Compulsory audit of the current companion notes}

This audit was made against the live companion files present on
8 September 2026.  The Standard Model programme advanced while the
audit was being performed, so the final record uses the later file
and not the superseded intermediate version:
\[
\begin{array}{c|r|r|l}
\text{file}&\text{lines}&\text{bytes}&\text{SHA--256}\\ \hline
\texttt{Renewal\_SM\_Einstein\_Continuum\_Research\_Notes\_V53.tex}
 &19253&731432&
 \texttt{F85E61D99E971DC4ADC6A9819C76C267621D899D7CBDEA2DA7F95D6187E13D3D}
 \\
\texttt{Renewal\_NS\_Stability\_Research\_Notes\_V26.tex}
 &5320&278283&
 \texttt{82C887F8C2C69E4F28FAD7C3DC8DE5B9099CB5A4BF431EBA1FFF3E91935978AD}.
\end{array}                                                     \tag{60.1}
\]

SM V51 constructs the exact vertical projection
\(\Pi_{\rm v}=I-JD{\cal B}\), proves the measured vertical and
projectable-horizontal identities, and writes the divergence of
\(\Pi_{\rm v}V\) with the derivatives of the inverse chart included.
It also states explicitly that a nonprojectable horizontal component
requires a conditional, measure-dependent transport law.  The
pushforward current of V58 is precisely that missing law.  SM V52
extends the bookkeeping to superdivergence, while SM V53 proves a
volume-free anchored closed-word bound for a small relative
superdeterminant.  The latter is useful for determinant sectors but is
not a clustering theorem for normalized Wilson expectations.

The NS file is unchanged from the V55 audit.  Its rank-one first- and
second-derivative bounds do not supply a normalized-current or
connected-polymer theorem.  Thus neither companion proves the full
interacting version of (59.7).  The SM divergence ledger does,
however, expose the derivative count corrected next.

\begin{assessment}[Transfer from the V60 audit]
\label{ass:e8v60-audit-transfer}
The useful transfer is twofold.  First, a Jacobian row must retain the
derivatives of every nonlinear projection or conditional transport.
Second, a volume-free connected determinant expansion is not by
itself a connected-correlation estimate for the normalized compact
measure.  The latter still requires a small-activity polymer or a
uniform functional-inequality argument.
\end{assessment}

\subsection{The Haar-Jacobian Hessian spends a third current jet}

Let \(q=(q^1,\ldots,q^m)\) be exponential coordinates at the identity
of a retained product group and write Haar measure as
\[
 d\mu_{\rm H}=h(q)\,dq,\qquad
 \gamma_i(q):=\partial_i\log h(q).
\]
The density is even at the identity, so \(\gamma_i(0)=0\).  For a
vector field \(U=U^i\partial_i\),
\[
 \operatorname{div}_{\rm H}U
 =\partial_iU^i+\gamma_iU^i.                                  \tag{60.2}
\]

\begin{proposition}[Third-jet debit]
\label{prop:e8v60-third-jet-debit}
Assume \(U(0)=0\).  Then
\[
\begin{split}
 \partial_a\partial_b\operatorname{div}_{\rm H}U\big|_0
 ={}&\partial_a\partial_b\partial_iU^i\big|_0\\
 &+(\partial_a\gamma_i)(0)\,\partial_bU^i(0)
  +(\partial_b\gamma_i)(0)\,\partial_aU^i(0).                \tag{60.3}
\end{split}
\]
Consequently the kernel \(K_{\rm Jac}\) in (59.12) requires the
third background jet of \(\overline V\), not merely the two jets
listed in (59.8).
\end{proposition}

\begin{proof}
Differentiate (60.2) twice.  The four product-rule terms from
\(\gamma_iU^i\) are
\[
 (\partial_{ab}\gamma_i)U^i
 +(\partial_a\gamma_i)(\partial_bU^i)
 +(\partial_b\gamma_i)(\partial_aU^i)
 +\gamma_i\partial_{ab}U^i.
\]
The first and last vanish at the identity by \(U(0)=0\) and
\(\gamma(0)=0\).  The derivative of \(\partial_iU^i\) gives the
first line of (60.3).
\end{proof}

This does not invalidate the exact decomposition (59.14), but it
corrects its stated locality input: two current jets establish
locality of the action-gradient Hessian, whereas the Jacobian
Hessian needs one additional jet.

\subsection{All-order connected score formula}

For a finite labelled set \(I\) of background directions, write
\(D_I=\prod_{i\in I}D_i\).  If \(J\subset I\), put
\(F_J=D_JF\).  For a nonempty block \(B\subset I\), put
\(S_B=D_BS\).  Let \(\Pi(A)\) be the set of set partitions of
\(A\), including the empty partition of the empty set, and use the
one-entry convention \(\kappa(G)=\langle G\rangle\).

\begin{theorem}[Connected score formula to arbitrary order]
\label{thm:e8v60-all-score}
For the normalized measure (59.1),
\[
 \boxed{
 D_I\langle F\rangle
 =\sum_{J\subset I}\ 
   \sum_{\pi\in\Pi(I\setminus J)}
   (-1)^{|\pi|}
   \kappa\bigl(F_J,(S_B)_{B\in\pi}\bigr).}                 \tag{60.4}
\]
In particular, the third derivative is a sum of connected
cumulants with at most four entries.  No unnormalized partition
function occurs.
\end{theorem}

\begin{proof}
For normalized expectations, differentiation of a cumulant gives
\[
 D_a\kappa(G_1,\ldots,G_k)
 =\sum_{j=1}^k
   \kappa(G_1,\ldots,D_aG_j,\ldots,G_k)
  -\kappa(G_1,\ldots,G_k,S_a).                               \tag{60.5}
\]
This follows by differentiating the logarithm of the joint moment
generating function, exactly as in (59.6).  Induct on \(|I|\).
When a new derivative \(a\) is applied to a term of (60.4), it
either joins \(J\), joins one existing score block \(B\), or forms
a new singleton score block \(\{a\}\).  These are, respectively,
the first sum in (60.5), differentiation of an existing \(S_B\),
and the final negative cumulant.  They enumerate each pair
\((J',\pi')\) for \(I\cup\{a\}\) exactly once and give the sign
\((-1)^{|\pi'|}\).  The empty-set case is the expectation itself.
\end{proof}

\begin{corollary}[Corrected connected-current locality input]
\label{cor:e8v60-third-current-locality}
To control \(D^r\overline V\) it suffices to have the local score
derivatives through order \(r\) and connected tree decay for at
most \(r+1\) marked observables.  In particular, locality of
\(K_{\rm Jac}\) follows from a four-point version of (59.7),
together with (60.3).  The Gaussian proof extends immediately:
connected Wick graphs with at most four marked supports contain a
spanning tree and obey the same weighted convolution estimate.
\end{corollary}

\begin{proof}
Apply (60.4) to \(F=j\).  Every term for \(|I|=r\) has at most
\(r+1\) cumulant entries.  Sum the associated tree bound by the
tree-pruning argument of Theorem~\ref{thm:e8v59-current-locality}.
For a Gaussian measure, the proof of
Lemma~\ref{lem:e8v59-Wick-bound} is independent of the number of
marked supports as long as that number and the polynomial degrees
are fixed; take four supports.
\end{proof}

\subsection{Exact Casimir calculation on product Haar space}

We now calculate a concrete nonlinear completion.  Let
\(G=SU(N)\), and choose an anti-Hermitian orthonormal basis
\(T^A\) with
\[
 -\operatorname{Tr}(T^AT^B)=\delta^{AB},\qquad
 \sum_AT^AT^A=-C_FI,\qquad C_F=\frac{N^2-1}{N}.             \tag{60.6}
\]
Let \(X=G^E\) have the product bi-invariant metric and normalized
product Haar measure.  Its Laplacian is
\(\Delta_X=\sum_{e,A}(L_e^A)^2\), so
\(\operatorname{div}_{\rm H}\nabla f=\Delta_Xf\).
For a non-self-intersecting elementary plaquette \(p\), define
\[
 r_p(U):=1-\frac1N\operatorname{ReTr}U_p,\qquad
 V_{\rm W}:=\sum_pr_p,\qquad
 \Phi:=\sum_pr_p^2.                                         \tag{60.7}
\]

\begin{lemma}[One-plaquette Haar identities]
\label{lem:e8v60-plaquette-Haar}
For every elementary four-link plaquette,
\[
 \boxed{\Delta_Xr_p=4C_F(1-r_p),}                            \tag{60.8}
\]
and, in exponential link coordinates \(U_e=e^{a_e}\),
\[
 r_p=r_p^{(2)}+O(\|a\|^3),\qquad
 r_p^{(2)}=\frac1{2N}\|q_pa\|^2,                            \tag{60.9}
\]
\[
 \|\nabla r_p\|^2
 =\frac8N r_p^{(2)}+O(\|a\|^3),                             \tag{60.10}
\]
where \(q_p\) is the signed linear plaquette curl.
\end{lemma}

\begin{proof}
For any one link occurring once in \(U_p\), the Casimir identity
(60.6) gives
\[
 \sum_A(L_e^A)^2\operatorname{ReTr}U_p
 =-C_F\operatorname{ReTr}U_p.
\]
The four link Laplacians add, proving (60.8).  At the identity,
the linearized plaquette logarithm is \(q_pa\), and
\(\operatorname{ReTr}e^A=N-\|A\|^2/2+O(\|A\|^3)\), which proves
(60.9).  The quadratic gradient is
\(N^{-1}q_p^*q_pa\).  Since the four signed link incidences obey
\(q_pq_p^*=4I_{\mathfrak{su}(N)}\),
\[
 \|N^{-1}q_p^*q_pa\|^2
 =\frac4{N^2}\|q_pa\|^2
 =\frac8N r_p^{(2)}.
\]
This proves (60.10).
\end{proof}

\begin{corollary}[Square-plaquette Laplacian]
\label{cor:e8v60-square-Laplacian}
The exact identity
\[
 \Delta_X\Phi
 =\sum_p\left\{8C_Fr_p(1-r_p)
                    +2\|\nabla r_p\|^2\right\}              \tag{60.11}
\]
has quadratic part
\[
 (\Delta_X\Phi)^{(2)}
 =8\left(C_F+\frac2N\right)\sum_pr_p^{(2)}.                 \tag{60.12}
\]
\end{corollary}

\begin{proof}
Use
\(\Delta(r_p^2)=2r_p\Delta r_p+2\|\nabla r_p\|^2\), then
insert (60.8)--(60.10).
\end{proof}

\subsection{A local Haar-compensated EOM vector}

For a real parameter \(\lambda\), set
\[
 {\cal V}_\lambda
 :=\frac12\nabla V_{\rm W}+\lambda\nabla\Phi.               \tag{60.13}
\]
This is a local analytic gauge-covariant vector field.  It also
respects translations, reflections, and hypercubic transformations.
Because \(\nabla\Phi=O(\|a\|^3)\), its linearization is the same
as that of \(\frac12\nabla V_{\rm W}\), namely
\(\frac12{\cal P}a\) in the normalization of (57.1).

\begin{theorem}[Exact Haar divergence and canonical compensation]
\label{thm:e8v60-compensated-divergence}
The product-Haar divergence is
\[
\boxed{\begin{split}
 \operatorname{div}_{\rm H}{\cal V}_\lambda
 ={}&2C_F\sum_p(1-r_p)\\
 &+\lambda\sum_p\left\{8C_Fr_p(1-r_p)
                         +2\|\nabla r_p\|^2\right\}.
\end{split}}                                                  \tag{60.14}
\]
Its quadratic part, apart from the vacuum constant, is
\[
 \left[-2C_F+8\lambda\left(C_F+\frac2N\right)\right]
 \sum_pr_p^{(2)}.                                            \tag{60.15}
\]
Therefore the canonical value
\[
 \boxed{
 \lambda_*:=\frac{C_F}{4(C_F+2/N)}
 =\frac{N^2-1}{4(N^2+1)}}                                   \tag{60.16}
\]
gives
\[
 D^2\operatorname{div}_{\rm H}{\cal V}_{\lambda_*}(I)=0.
                                                                    \tag{60.17}
\]
For \(SU(3)\), \(\lambda_*=1/5\).
\end{theorem}

\begin{proof}
The divergence of a gradient is its Laplacian.  Apply (60.8) to
\(V_{\rm W}\) and (60.11) to \(\Phi\), obtaining (60.14).
Equations (60.9)--(60.12) give (60.15).  Solving its scalar
coefficient for zero gives (60.16), hence (60.17).
\end{proof}

\begin{remark}[Why the uncompensated gradient is not enough]
\label{rem:e8v60-uncompensated-feed}
For \(\lambda=0\), the Jacobian part
\(-\operatorname{div}_{\rm H}{\cal V}_0\) contains
\(+2C_FV_{\rm W}^{(2)}\).  Thus the most immediate compact-group
completion of the flat EOM vector has a nonzero marginal
Haar-Jacobian term.  Its omission would make the quadratic
field-redefinition argument incorrect.  The cubic vector
\(\lambda_*\nabla\Phi\) cancels exactly this term without changing
the prescribed linearization.
\end{remark}

\begin{theorem}[No fine Wilson marginal for the compensated owner]
\label{thm:e8v60-no-fine-marginal}
Let the pure Wilson action be \(S=\beta V_{\rm W}\), and define
\[
 {\cal O}_*:=DS[{\cal V}_{\lambda_*}]
             -\operatorname{div}_{\rm H}{\cal V}_{\lambda_*}.
                                                                    \tag{60.18}
\]
After removal of its vacuum constant, the quadratic Hessian of
\({\cal O}_*\) is \(\beta{\cal P}^2\).  In particular its
\(F^2\) moment vanishes and its quadratic owner begins at fourth
external-momentum order.
\end{theorem}

\begin{proof}
At the vacuum,
\(V_{\rm W}=\frac12\langle a,{\cal P}a\rangle+O(\|a\|^3)\)
and
\({\cal V}_{\lambda_*}=\frac12{\cal P}a+O(\|a\|^2)\).
The compensation is cubic and hence contributes no quadratic term
to \(DS[{\cal V}_{\lambda_*}]\).  Therefore that quadratic term is
\(\frac\beta2\langle a,{\cal P}^2a\rangle\).  Equation (60.17)
removes the quadratic divergence.  Finally
\({\cal P}^2=S{\cal P}\) by (57.2), which is fourth order in
centered momentum and has zero second moment.
\end{proof}

\subsection{What has and has not been decided about
\texorpdfstring{\(\zeta_{\rm Jac}\)}{zetaJac}}

The preceding construction also proves that a retained
Haar-Jacobian marginal is a coordinate-section coefficient rather
than a universal physical coefficient.  Let \(V_{{\rm W},Y}\) and
\(\Phi_Y\) be (60.7) on the retained lattice, and put
\[
 K_{{\rm W},Y}:=D^2V_{{\rm W},Y}(I),\qquad
 \omega_Y:={\mathfrak m}(K_{{\rm W},Y})\ne0.                 \tag{60.19}
\]

\begin{proposition}[Local removability of the Jacobian marginal]
\label{prop:e8v60-Jacobian-removable}
For any local equivariant retained velocity \(\overline V\), replace it by
\[
 \overline V^{\,\prime}
 =\overline V+\lambda\nabla\Phi_Y.                           \tag{60.20}
\]
This does not change its linearization at the vacuum.  It changes
the Jacobian marginal by
\[
 \zeta_{\rm Jac}^{\,\prime}
 =\zeta_{\rm Jac}
 -8\lambda\left(C_F+\frac2N\right)\omega_Y.                 \tag{60.21}
\]
Hence the unique choice
\[
 \lambda=\frac{\zeta_{\rm Jac}}
 {8(C_F+2/N)\omega_Y}
\]
sets \(\zeta_{\rm Jac}^{\,\prime}=0\).  The added action-gradient
term has no quadratic Hessian, so this operation changes only the
redundant coordinate section at marginal order.
\end{proposition}

\begin{proof}
The zero linearization follows from
\(\nabla\Phi_Y=O(\|a\|^3)\).  Equation (60.12) says that
\[
 D^2\operatorname{div}_{\rm H}(\lambda\nabla\Phi_Y)(I)
 =8\lambda(C_F+2/N)K_{{\rm W},Y}.
\]
Remembering the minus sign in the definition of \(K_{\rm Jac}\)
and applying \({\mathfrak m}\) proves (60.21).  Since the retained
effective action is stationary at the vacuum, its differential
applied to a cubic vector has degree at least four and contributes
no quadratic Hessian.
\end{proof}

\begin{firewall}[Vanishing in a chosen section is not a raw-current
calculation]
\label{fire:e8v60-section-vs-raw}
Proposition~\ref{prop:e8v60-Jacobian-removable} proves that the
marginal Jacobian coefficient is removable in the quotient by local
measured field redefinitions.  It does not prove that the raw
pushforward current (59.3) has \(\zeta_{\rm Jac}=0\) before the
coordinate change.  Computing that raw coefficient requires the
third current jet and hence the four-point connected bounds identified
in Corollary~\ref{cor:e8v60-third-current-locality}.
\end{firewall}

\subsection{Status and the V61 acquisition}

\begin{assessment}[Progress after V60]
\label{ass:e8v60-status}
V60 completes the compulsory companion audit, corrects the
conditional-current derivative count, and proves the all-order
connected score formula.  It evaluates the exact product-Haar
divergence of a concrete compact-group EOM vector and constructs the
local cubic compensation \(\lambda_*=1/5\) for \(SU(3)\).  The
resulting measured Wilson owner has no quadratic \(F^2\) component.
At the retained level, any remaining \(\zeta_{\rm Jac}\) is proved
to be a removable local coordinate coefficient, although its value
in the unmodified raw pushforward current is not yet evaluated.
\end{assessment}

\begin{decision}[V61 target]
\label{dec:e8v60-next}
V61 should formulate and prove a four-marked-observable cluster
gate for the normalized compact Wilson fibre in a declared
small-activity polymer chamber.  It must give an explicit
entropy-versus-activity inequality, include the large-field
activities, and return the third conditional-current jet in a
volume-independent weighted norm.  If that gate cannot be verified
from the present Wilson constants, the programme should move
upstream to the single-block large-field suppression estimate rather
than assume (59.7).
\end{decision}

\begin{firewall}[Claim boundary after V60]
\label{fire:e8v60-claim}
V60 proves a compensated fine EOM completion and the coordinate
removability of the Haar-Jacobian marginal.  It does not prove the
four-point interacting cluster gate, evaluate the raw pushed
\(\zeta_{\rm Jac}\), compare the compensated owner with
\({\cal O}_1\), prove anisotropy contraction, evaluate the remaining
cards, control toron/flux/overlap owners, complete the LERB,
establish the universal beta coefficient, close the nonlinear/complex
quotient matching card, construct the full RG trajectory, remove the
regulators, establish Osterwalder--Schrader reconstruction, or prove
the Yang--Mills mass gap.
\end{firewall}

\subsection{Parent record}

This V60 note was formed by appending the preceding material to V59,
whose validated record was
\[
\begin{array}{c|c}
\text{lines}&25667\\
\text{bytes}&898830\\
\text{SHA-256}&
\texttt{E2CD25BCE7163BBA632A5E5861873E89D2786897415067832503DF7E755D6497}.
\end{array}
\]

% =============================================================================
% END APPENDED EIGHTH-SERIES V60 MATERIAL
% =============================================================================

% =============================================================================
% BEGIN APPENDED EIGHTH-SERIES V61 MATERIAL
% =============================================================================

\section{Eighth-series V61: an explicit four-mark polymer gate for
the third conditional-current jet}
\label{sec:e8v61}

\subsection{Purpose and exact scope}

Corollary~\ref{cor:e8v60-third-current-locality} reduced locality of
the Haar-Jacobian Hessian to connected clustering of at most four
local observables.  This note proves a volume-independent theorem
which supplies that clustering from an exact small-activity polymer
representation.  It also separates the good-field perturbative
activity from the bad-field Peierls activity.  The theorem is an
implication with an explicit numerical gate; the final subsection
records which Wilson constants are not yet loaded.

Let \({\cal G}\) be the block adjacency graph, of maximum degree
\(\Delta\).  A polymer is a nonempty finite connected vertex set.
Two polymers are incompatible when they intersect.  For a family
\(\Gamma\) of pairwise compatible polymers, write
\(z_t(\Gamma)=\prod_{Y\in\Gamma}z_t(Y)\), where
\(t=(t_0,\ldots,t_r)\) is a vector of local complex sources.

\subsection{A self-contained Kotecky--Preiss majorant}

We first use a deliberately coarse animal count whose proof requires
no asymptotics.

\begin{lemma}[Connected-set entropy]
\label{lem:e8v61-animal-count}
For every vertex \(x\) and every \(n\ge1\), the number of connected
vertex sets \(Y\ni x\) with \(|Y|=n\) is at most
\[
                         \Delta^{2(n-1)}.                    \tag{61.1}
\]
\end{lemma}

\begin{proof}
Fix an ordering of the neighbours at every vertex.  Give each
connected set its deterministic depth-first spanning tree rooted at
\(x\), using the first unused neighbour whenever a choice occurs.
The depth-first traversal has exactly \(2(n-1)\) oriented steps.
The traversal word determines every visited vertex and hence the
set.  At each step there are at most \(\Delta\) choices, giving
(61.1).  Some words do not encode a tree, which only enlarges the
upper bound.
\end{proof}

For \(a,\mu>0\) and \(u\ge0\), define
\[
 Q_{a,\mu}(u)
 :=\frac{u e^{a+\mu}}
 {1-\Delta^2u e^{a+\mu}},                                   \tag{61.2}
\]
whenever the denominator is positive.

\begin{lemma}[Explicit KP gate]
\label{lem:e8v61-KP-gate}
Suppose, uniformly on a source polydisc \({\cal D}_\rho\),
\[
                         |z_t(Y)|\le u^{|Y|}.                 \tag{61.3}
\]
If
\[
 \Delta^2u e^{a+\mu}<1,\qquad Q_{a,\mu}(u)\le a,            \tag{61.4}
\]
then, for every polymer \(Y\),
\[
 \sum_{Y':Y'\cap Y\ne\varnothing}
 |z_t(Y')|e^{(a+\mu)|Y'|}
 \le a|Y|.                                                   \tag{61.5}
\]
For the four-dimensional nearest-neighbour block graph,
\(\Delta=8\).  The simple sufficient choice \(a=1\) is
\[
                  \boxed{65u e^{1+\mu}\le1.}                \tag{61.6}
\]
\end{lemma}

\begin{proof}
For a fixed vertex, Lemma~\ref{lem:e8v61-animal-count} and (61.3)
give
\[
 \sum_{Y'\ni x}|z_t(Y')|e^{(a+\mu)|Y'|}
 \le\sum_{n\ge1}\Delta^{2(n-1)}
       (ue^{a+\mu})^n=Q_{a,\mu}(u).                          \tag{61.7}
\]
An incompatible \(Y'\) contains at least one vertex of \(Y\).
Sum (61.7) over those vertices and use (61.4), proving (61.5).
When \(a=1\), the inequality
\(Q_{1,\mu}(u)\le1\) is equivalent to
\((\Delta^2+1)ue^{1+\mu}\le1\), which gives (61.6).
\end{proof}

\begin{lemma}[Rooted cluster majorant]
\label{lem:e8v61-rooted-cluster}
Under (61.5), the hard-core partition function
\[
 {\cal Z}(t)=
 \sum_{\Gamma\ {\rm pairwise\ compatible}}z_t(\Gamma)       \tag{61.8}
\]
is nonzero on \({\cal D}_\rho\), and
\[
 \log{\cal Z}(t)
 =\sum_{n\ge1}\frac1{n!}
   \sum_{Y_1,\ldots,Y_n}
   \phi^T(Y_1,\ldots,Y_n)\prod_{j=1}^nz_t(Y_j),              \tag{61.9}
\]
where \(\phi^T\) is the hard-core Ursell coefficient.  Moreover,
for a fixed root polymer \(Y_0\),
\[
\begin{split}
 &\sum_{n\ge0}\frac1{n!}
  \sum_{Y_1,\ldots,Y_n}
  |\phi^T(Y_0,Y_1,\ldots,Y_n)|\\
 &\hspace{35mm}\times
  \prod_{j=1}^n\{|z_t(Y_j)|e^{\mu|Y_j|}\}
 \le e^{a|Y_0|}.                                             \tag{61.10}
\end{split}
\]
All sums are uniform in the volume.
\end{lemma}

\begin{proof}
For polymers \(Y_i,Y_j\), put \(f_{ij}=-1\) if they intersect and
zero otherwise.  Expanding the compatibility product
\(\prod_{i<j}(1+f_{ij})\) and sorting graphs by connected
components gives (61.9) formally, with
\[
 \phi^T(Y_1,\ldots,Y_n)
 =\sum_{G\ {\rm connected}}\prod_{\{i,j\}\in E(G)}f_{ij}.
\]
The Penrose deletion map assigns each connected graph to a spanning
tree and cancels alternating supergraphs in every fibre; hence
\[
 |\phi^T(Y_1,\ldots,Y_n)|
 \le\sum_{T\ {\rm tree}}
      \prod_{\{i,j\}\in E(T)}{\bf1}_{Y_i\cap Y_j\ne\varnothing}.
                                                                    \tag{61.11}
\]
Root a tree at \(Y_0\).  Successively sum its terminal branches.
For a parent polymer \(Y\), (61.5), applied to the activity
\(|z_t(Y')|e^{\mu|Y'|}\), supplies at most \(a|Y|\) after the
standard exponential generating sum over unordered children.
Induction from the leaves yields \(e^{a|Y_0|}\), proving (61.10).
The absolute majorant just obtained validates the formal component
sorting, normal convergence of (61.9), and nonvanishing of
\({\cal Z}\) on the branch connected to \({\cal Z}=1\) at zero
activities.
\end{proof}

\subsection{Marked sources and tree decay}

Let \(X_0,\ldots,X_r\) be fixed local supports, with \(r\le3\).
Assume \(z_t(Y)\) is analytic for \(|t_i|<\rho_i\), satisfies
(61.3) on the closed smaller polydisc, and has the locality property
\[
 D_{t_i}z_t(Y)=0\quad\hbox{whenever }Y\cap X_i=\varnothing.   \tag{61.12}
\]
Let \(\tau(X_0,\ldots,X_r)\) be the length of a shortest block
tree meeting all the marked supports.

\begin{theorem}[Four-mark polymer clustering]
\label{thm:e8v61-four-mark-cluster}
Under (61.3)--(61.5), the connected source response satisfies
\[
\boxed{
 \left|
 D_{t_0}\cdots D_{t_r}\log{\cal Z}(0)
 \right|
 \le
 2^{r+1}|X_0|Q_{a,\mu}(u)
 \left(\prod_{i=0}^r\rho_i^{-1}\right)
 e^{-\mu\tau(X_0,\ldots,X_r)}.}                             \tag{61.13}
\]
The bound is independent of the total number of blocks.
\end{theorem}

\begin{proof}
Use (61.9).  Inclusion--exclusion over the replacements
\(t_i\mapsto0\) selects exactly those cluster terms which depend on
every source; it costs at most \(2^{r+1}\).  By (61.12), the union
of the polymers in any selected connected cluster meets every
\(X_i\).  The incompatibility graph is connected, so the union is a
connected block set and
\[
 \sum_j|Y_j|\ge \tau(X_0,\ldots,X_r).                       \tag{61.14}
\]
Extract \(e^{-\mu\tau}\) from the factor
\(e^{\mu\sum_j|Y_j|}\).  Choose as root the first polymer meeting
\(X_0\).  Lemma~\ref{lem:e8v61-rooted-cluster}, followed by
(61.7) and a sum over the at most \(|X_0|\) anchor vertices, bounds
the remaining absolute series by \(|X_0|Q_{a,\mu}(u)\).
Multivariate Cauchy on the source polydisc contributes
\(\prod_i\rho_i^{-1}\).  This proves (61.13).
\end{proof}

\begin{corollary}[Third conditional-current jet]
\label{cor:e8v61-third-current}
Suppose the source-coupled reduced Wilson fibre integral admits the
exact polymer representation (61.8), with one source for the local
current \(j\) and up to three sources for the score insertions in
(60.4).  If (61.3)--(61.6) hold uniformly in the retained
background chart, then all connected cumulants needed for
\(D^3\overline V\) obey (59.7).  Consequently
\(D^3\overline V\) and the Haar-Jacobian Hessian are exponentially
summable in every smaller weight.
\end{corollary}

\begin{proof}
Connected cumulants are mixed derivatives of the logarithm of the
source-coupled normalized partition function.  The normalization
subtracts only the source-free value and hence does not change a
mixed derivative containing the current source.  Apply
Theorem~\ref{thm:e8v61-four-mark-cluster}, then (60.4), (60.3), and
the weighted tree-pruning argument of V59.
\end{proof}

\subsection{The large-field activity is a separate input}

The compactness of the link group makes a one-block suppression
statement elementary once all zero-action sectors have been assigned
to good charts.  It does not, by itself, prove the multi-block
Peierls estimate.

\begin{lemma}[Uniform one-block bad-field suppression]
\label{lem:e8v61-one-block-bad}
Let \(F\) be a compact \(n\)-dimensional reduced one-block fibre
with probability Haar measure \(dm\), and let \(v_W\ge0\) depend
continuously on a retained background \(W\) in a compact chart
\({\cal C}\).  Suppose:
\begin{enumerate}
\item every zero of \(v_W\) lies in a declared good set \(G_W\);
\item on \(F\setminus G_W\),
      \(v_W\ge\delta_b>0\), uniformly in \(W\); and
\item each \(G_W\) contains a coordinate ball on which
      \(v_W(x)\le C_0|x|^2\) and the Haar density is at least
      \(h_0>0\), with uniform radius \(r_0\).
\end{enumerate}
Then, for \(\beta\ge r_0^{-2}\),
\[
 \nu_{\beta,W}(F\setminus G_W)
 \le C_b\beta^{n/2}e^{-\beta\delta_b},\qquad
 C_b:=\frac1{h_0|B_1|e^{-C_0}},                              \tag{61.15}
\]
uniformly in \(W\).
\end{lemma}

\begin{proof}
The bad numerator is at most \(e^{-\beta\delta_b}\), because
\(dm\) is a probability measure.  Integrate the denominator over
the coordinate ball of radius \(\beta^{-1/2}\).  There
\(\beta v_W\le C_0\), while the measure is at least
\(h_0|B_1|\beta^{-n/2}\).  Dividing gives (61.15).
\end{proof}

For a connected bad component \(Y\), one additionally needs a
Peierls additivity estimate after the boundary links and retained
torons are fixed.  If that estimate and the good-field expansion give
\[
 \|z^{\rm bad}_t(Y)\|
 \le u_b(\beta)^{|Y|},\qquad
 u_b(\beta):=C_b\beta^{n/2}e^{-\beta\delta_b},                \tag{61.16}
\]
\[
 \|z^{\rm good}_t(Y)\|
 \le u_g(\beta)^{|Y|},                                      \tag{61.17}
\]
then the triangle inequality gives (61.3) with
\[
                         u(\beta)=u_g(\beta)+u_b(\beta).     \tag{61.18}
\]

\begin{theorem}[Explicit nonlinear current gate]
\label{thm:e8v61-Wilson-gate}
Assume the exact good/bad polymerization of the reduced Wilson fibre
satisfies (61.16)--(61.18), including the four source derivatives,
and that
\[
 \boxed{
 65e^{1+\mu}\left{
 u_g(\beta)+C_b\beta^{n/2}e^{-\beta\delta_b}
 \right}\le1.}                                             \tag{61.19}
\]
Then the nonlinear conditional current has three
volume-uniform exponentially local background jets.  The raw
\(K_{\rm Jac}\) in (59.12) is therefore a well-defined local
kernel, and its marginal coefficient can be computed before applying
the removable-coordinate result of V60.
\end{theorem}

\begin{proof}
Equation (61.19) is (61.6) with (61.18).  Apply
Corollary~\ref{cor:e8v61-third-current}.
\end{proof}

\begin{firewall}[Compactness is not Peierls additivity]
\label{fire:e8v61-compact-not-Peierls}
Lemma~\ref{lem:e8v61-one-block-bad} controls one isolated block.
It does not prove (61.16) for a connected set of bad blocks, because
plaquettes cross block boundaries and flat centre or toron sectors
can correlate neighbouring blocks.  Those sectors must be retained
or charted first, and an energy cost proportional to \(|Y|\) must be
proved for what remains.
\end{firewall}

\subsection{Loaded versus missing Wilson inputs}

The present notes already contain three ingredients for the good
part: the exact vertical Hessian floor (41.22), the weighted inverse
bound of V45, and the finite analytic Wilson-jet bounds of V54--V55.
They do not yet combine those data into the source-uniform activity
number \(u_g(\beta)\) in (61.17).  For the bad part, compactness
proves existence of \(\delta_b\) only after the complete zero-action
set has been charted, but no quantitative multi-block Peierls value
has been established.  Thus (61.19) is an explicit gate, not a
verified claim about every sufficiently large \(\beta\).

\begin{assessment}[Progress after V61]
\label{ass:e8v61-status}
V61 proves a volume-free four-mark cluster theorem with the explicit
four-dimensional entropy condition (61.6), derives the third
conditional-current jet from it, and proves a uniform one-block
large-field suppression lemma.  It reduces full nonlinear current
locality to two concrete numbers, \(u_g(\beta)\) and the connected
bad-field cost \(\delta_b\), combined in (61.19).  Neither number is
silently inferred from the vacuum Hessian.
\end{assessment}

\begin{decision}[V62 upstream target]
\label{dec:e8v61-next}
V62 should classify the exact zero set of the unit Wilson action on
one straight-block reduced fibre with retained toron and centre-flux
data fixed.  On the contractible sector it should use the parity-tree
coordinates and the incidence floor (41.22) to derive a quantitative
energy bound outside a declared good ball.  If an additional flat
sector survives, it must be promoted to the retained ledger before a
positive \(\delta_b\) is asserted.
\end{decision}

\begin{firewall}[Claim boundary after V61]
\label{fire:e8v61-claim}
V61 proves the abstract four-point polymer implication and the
one-block compact suppression estimate.  It does not verify the
Wilson activity gate (61.19), prove multi-block Peierls additivity,
evaluate the raw pushed \(\zeta_{\rm Jac}\), compare the compensated
owner with \({\cal O}_1\), prove anisotropy contraction, evaluate the
remaining cards, control toron/flux/overlap owners, complete the LERB,
establish the universal beta coefficient, close the nonlinear/complex
quotient matching card, construct the full RG trajectory, remove the
regulators, establish Osterwalder--Schrader reconstruction, or prove
the Yang--Mills mass gap.
\end{firewall}

\subsection{Parent record}

This V61 note was formed by appending the preceding material to V60,
whose validated record was
\[
\begin{array}{c|c}
\text{lines}&26127\\
\text{bytes}&916133\\
\text{SHA-256}&
\texttt{CB265825C8CDF7D8359E23530A107B6A791CBD551429D75179806CCB072C13DB}.
\end{array}
\]

% =============================================================================
% END APPENDED EIGHTH-SERIES V61 MATERIAL
% =============================================================================

% =============================================================================
% BEGIN APPENDED EIGHTH-SERIES V62 MATERIAL
% =============================================================================

\section{Eighth-series V62: retained-background fibre coercivity
and a localized bad-block Peierls cost}
\label{sec:e8v62}

\subsection{Imported result and the nonduplication boundary}

The attached file
\[
\begin{split}
&\texttt{Renewal\_Yang\_Mills\_Reduced\_Fibre\_Wilson\_Shell\_Symbol\_and\_}\[-0.2em]
&\hspace{20mm}\texttt{Local\_Vacuum\_Programme\_V35.tex}
\end{split}
\]
has \(39839\) lines, \(1560166\) bytes, and SHA--256
\[
\texttt{7C6056D8149D3015D469B57474A75B23C88A6FE9F64FC81C7DCA85C361C647A3}.
                                                                    \tag{62.1}
\]
Its Version V2 already proves, on the complete pinned reduced fibre
at \(W=I\),
\[
 \sum_{e\in{\mathscr E}^{\rm red}}
 \|z_e-I\|_F^2
 \le20\sum_p\|U_p-I\|_F^2=120V_M(z,I),                      \tag{62.2}
\]
and hence proves that \(z=I\) is the unique fibre minimum.  It also
proves the exact forty-five-link filling table and its maximal
weighted plaquette occupancy twenty.  None of those results is
reproved here.  The new work is to retain \(W\), identify the moving
comparison section, classify the zero set, and localize (62.2) to a
set of bad coarse blocks.

\subsection{The moving comparison section}

Recall the exact parity-tree section.  A fine vertex is
\(2y+s\), with \(s\in\{0,1\}^4\).  The tree path \(P_s\) lists the
occupied directions in decreasing order.  For a reduced positive
edge \(e=(2y+s,\mu)\), the attached filling compares
\[
 A_{s,\mu}=P_s\mu
 \quad\hbox{with}\quad
 B_{s,\mu}=\begin{cases}
 P_{s+\widehat e_\mu},&s_\mu=0,\\
 \mu\mu P_{s-\widehat e_\mu},&s_\mu=1.
 \end{cases}                                                  \tag{62.3}
\]
In the section, every tree link is \(I\), and the second link of
the straight two-link coarse path is \(W_{y,\mu}\).  Define
\[
 H_e(W):=\operatorname{Hol}(B_{s,\mu})
 =\begin{cases}
 I,&s_\mu=0,\\
 W_{y,\mu},&s_\mu=1,
 \end{cases}                                                  \tag{62.4}
\]
and
\[
 D_W(z):=\sum_{e\in{\mathscr E}^{\rm red}}
          \|z_e-H_e(W)\|_F^2.                                \tag{62.5}
\]

\begin{theorem}[Uniform retained-background fibre coercivity]
\label{thm:e8v62-moving-coercivity}
For every retained configuration \(W\), every \(M\ge4\), and every
point \(z\) of the exact reduced fibre,
\[
 \boxed{D_W(z)\le20\sum_p\|U_p(z,W)-I\|_F^2
                   =120V_M(z,W).}                            \tag{62.6}
\]
The constant is independent of \(W\) and of the torus side.
\end{theorem}

\begin{proof}
The holonomy of \(A_{s,\mu}\) is \(z_e\), because its prefix is a
tree path.  Equation (62.4) is the exact holonomy of the comparison
path \(B_{s,\mu}\).  The adjacent-swap sequence imported with
(62.2) changes one path into the other.  A swap changes its holonomy
by a unitary left and right factor times \(U_p-I\); hence
\[
 \|z_e-H_e(W)\|_F
 \le\sum_{p\in\Sigma_e}\|U_p(z,W)-I\|_F.                    \tag{62.7}
\]
The filling lengths and the incidence of every plaquette are
combinatorial and do not depend on the values of \(W\).  The
imported Cauchy--Schwarz and occupancy-twenty summation therefore
applies verbatim to (62.7), giving the first inequality in (62.6).
For \(SU(3)\),
\(1-\operatorname{ReTr}U/3=\|U-I\|_F^2/6\), which gives the
equality with \(120V_M\).
\end{proof}

\begin{remark}[No logarithm and no small-background hypothesis]
\label{rem:e8v62-no-chart}
The distance in (62.5) is a chordal product-group distance to a
moving section.  The proof uses only unitary invariance of the
Frobenius norm.  It is therefore global on the compact fibre and
does not assume that either \(z\) or \(W\) lies in a logarithm
chart.
\end{remark}

\subsection{Exact zero-set classification}

Call \(W\) coarse-flat when every contractible elementary coarse
plaquette has holonomy \(I\).  Nontrivial commuting toron
holonomies are allowed; they are retained data.  A nontrivial
centre-flux plaquette is not coarse-flat in this definition.

\begin{theorem}[Zero set on a fixed retained fibre]
\label{thm:e8v62-zero-set}
The zero set
\[
 {\cal Z}_W:=\{z:V_M(z,W)=0\}
\]
is empty when \(W\) is not coarse-flat.  When \(W\) is
coarse-flat and belongs to the declared retained toron sector,
\[
                         \boxed{{\cal Z}_W=\{H(W)\}.}         \tag{62.8}
\]
Thus centre flux and toron data do not create a second reduced-fibre
minimum: centre flux has positive action, while a flat toron moves
the unique minimum through the retained section.
\end{theorem}

\begin{proof}
Every plaquette defect is nonnegative, and it vanishes only when its
plaquette holonomy is \(I\).  If \(V_M(z,W)=0\), all fine
plaquettes are therefore trivial.  The straight block identity
writes a coarse plaquette as the ordered product of the fine
plaquettes in its \(2\times2\) tiling, conjugated to one basepoint;
hence every coarse plaquette is \(I\).  This proves emptiness for a
non-flat \(W\).

If the action is zero, (62.6) also gives \(z=H(W)\), proving
uniqueness.  Conversely, suppose \(W\) is coarse-flat.  Start at a
coarse root, use the identity tree links inside each parity cell,
and use the designated second links to transport by \(W\) between
coarse roots.  Coarse flatness makes this transport independent of
every contractible coarse path.  Formula (62.4) assigns to every
remaining fine edge exactly the endpoint transport determined in
this way.  Consequently every fine elementary plaquette has
holonomy \(I\), so \(V_M(H(W),W)=0\).  Global toron holonomies are
unchanged because the construction never contracts a winding coarse
cycle.
\end{proof}

\begin{firewall}[Flat torons are retained, not integrated copies]
\label{fire:e8v62-toron-zero}
Equation (62.8) is a statement at fixed \(W\).  Varying a commuting
toron produces a family of minima in the total fine configuration
space, but exactly one point in each reduced fibre.  Integrating that
family as if it were a vertical degeneracy would destroy the block
map and the toron ledger.
\end{firewall}

\subsection{Localization of the filling inequality}

Let \({\cal C}_y\) be the forty-five reduced links based in the
coarse parity cell rooted at \(y\), and put
\[
 D_y(z;W):=\sum_{e\in{\cal C}_y}
             \|z_e-H_e(W)\|_F^2.                             \tag{62.9}
\]
Let \({\cal P}_y\) be the plaquettes occurring in the imported
fillings \(\Sigma_e\) for \(e\in{\cal C}_y\), and for a set of
coarse blocks \(Y\), put
\[
 {\cal P}(Y):=\bigcup_{y\in Y}{\cal P}_y,\qquad
 Y^+:=Y+\{0,1\}^4.                                          \tag{62.10}
\]

\begin{lemma}[Finite support of the filling owner]
\label{lem:e8v62-filling-support}
Every plaquette of \({\cal P}_y\) is based in a fine cube belonging
to one of the coarse cells \(y+\epsilon\),
\(\epsilon\in\{0,1\}^4\).  Consequently
\[
 {\cal P}(Y)\subset Y^+,\qquad |Y^+|\le16|Y|.               \tag{62.11}
\]
\end{lemma}

\begin{proof}
The words in (62.3) use positive coordinate steps only.  Each
coordinate is used at most once, except that \(\mu\) may occur
twice.  Every intermediate path in an adjacent swap therefore stays
inside \(2y+[0,2]^4\).  The plaquette swept by the swap is in one of
the sixteen coarse cells listed in (62.11).  Taking unions proves
the claim.
\end{proof}

\begin{theorem}[Localized coercivity for an arbitrary block set]
\label{thm:e8v62-local-coercivity}
For every set \(Y\) of coarse blocks,
\[
 \boxed{
 \sum_{y\in Y}D_y(z;W)
 \le20\sum_{p\in{\cal P}(Y)}\|U_p(z,W)-I\|_F^2.}             \tag{62.12}
\]
If a block is called \(r\)-bad when \(D_y(z;W)\ge r^2\), then
every set \(Y\) of \(r\)-bad blocks obeys the pointwise Peierls
bound
\[
 \boxed{
 \sum_{p\in{\cal P}(Y)}
 \left(1-\frac13\operatorname{ReTr}U_p\right)
 \ge\frac{r^2}{120}|Y|.}                                    \tag{62.13}
\]
In particular the missing energy density of V61 may be taken as
\[
                         \delta_b(r)=\frac{r^2}{120}.         \tag{62.14}
\]
\end{theorem}

\begin{proof}
Sum the squared form of (62.7) only over links based in \(Y\).
The coefficient of any plaquette in this restricted sum is no
larger than its coefficient in the complete filling table.  The
imported maximal weighted occupancy is twenty, proving (62.12).
If all blocks of \(Y\) are bad, its left side is at least
\(r^2|Y|\).  Divide (62.12) by \(120=20\cdot6\) and use the exact
\(SU(3)\) chord identity to obtain (62.13)--(62.14).
\end{proof}

\begin{corollary}[Unnormalized contour suppression]
\label{cor:e8v62-contour-suppression}
For the pure Wilson weight and any prescribed \(r\)-bad set \(Y\),
\[
 {\bf1}_{\{Y\ {\rm bad}\}}e^{-\beta V_M(z,W)}
 \le
 e^{-\beta r^2|Y|/120}
 \exp\left\{-\beta\!\sum_{p\notin{\cal P}(Y)}r_p(U)\right\}.
                                                                    \tag{62.15}
\]
The same estimate holds with a generated action \(\Psi\) after
multiplication by \(e^{-\Psi}\); no sign of \(\Psi\) is inferred.
\end{corollary}

\begin{proof}
Split the nonnegative Wilson action into plaquettes in
\({\cal P}(Y)\) and its complement, and apply (62.13) to the first
sum.
\end{proof}

\subsection{Insertion into the V61 activity gate}

The attachment also proves, for normalized \(SU(3)\) Haar measure,
\[
 I(t):=\int e^{-t\|U-I\|_F^2}\,dU
 \ge C_-t^{-4},\qquad
 C_-:=\frac{3e^{-43/32}}{1024\pi^8},\quad t\ge1.             \tag{62.16}
\]
Define the inverse forty-five-link reference normalizer
\[
 L_\beta:=I(4\beta)^{-45}
 \le\left\{C_-^{-1}(4\beta)^4\right\}^{45}.                 \tag{62.17}
\]

\begin{proposition}[Bad-activity normalization ledger]
\label{prop:e8v62-bad-ledger}
Suppose an exact good/bad contour extraction assigns to a bad
component \(Y\) only variables in \(Y^+\), divides by at most one
reference factor \(I(4\beta)^{45}\) for every block of \(Y^+\),
and assigns every remaining interaction to the good activity.
Then (62.15)--(62.17) give
\[
 |z_t^{\rm bad}(Y)|
 \le u_b(\beta,r)^{|Y|},                                    \tag{62.18}
\]
where one completely explicit, deliberately coarse choice is
\[
 \boxed{
 u_b(\beta,r):=
 \left\{C_-^{-1}(4\beta)^4\right\}^{720}
 e^{-\beta r^2/120}.}                                       \tag{62.19}
\]
For every fixed \(r>0\), this bound tends to zero exponentially as
\(\beta\to\infty\).
\end{proposition}

\begin{proof}
Lemma~\ref{lem:e8v62-filling-support} gives at most
\(16|Y|\) block normalizers.  Each costs at most \(L_\beta\), and
Corollary~\ref{cor:e8v62-contour-suppression} gives the exponential
factor.  Combining them yields
\(L_\beta^{16|Y|}e^{-\beta r^2|Y|/120}\).  Insert (62.17) and use
\(45\cdot16=720\).  Exponential decay dominates the fixed power
\(\beta^{2880}\).
\end{proof}

The hypothesis of Proposition~\ref{prop:e8v62-bad-ledger} is a
normalization and ownership statement, not an energy estimate.  It
must be proved by an exact contour extraction.  What V62 removes is
the more serious possible obstruction: a zero-energy bad component
or a fibre-flat centre copy.  Neither exists after the retained
section is handled correctly.

\begin{assessment}[Progress after V62]
\label{ass:e8v62-status}
V62 imports rather than repeats the attachment's pinned coercivity.
It proves the global moving-section inequality (62.6), classifies
the exact zero set over every retained background, and proves the
localized Peierls cost \(r^2/120\) for an arbitrary collection of
bad blocks.  It also turns that cost into the explicit provisional
bad-activity majorant (62.19) under a stated local-normalization
ledger.  The retained toron sector is preserved and centre flux is
not mistaken for a second fibre vacuum.
\end{assessment}

\begin{decision}[V63 target]
\label{dec:e8v62-next}
V63 should construct the exact source-dependent good/bad partition
of unity and its connected-component contour extraction on the
block graph.  It must prove the ownership and normalizer hypotheses
of Proposition~\ref{prop:e8v62-bad-ledger}, including plaquettes
shared by two components, and place every leftover small-field
interaction in a specific \(u_g(\beta)\) activity.  If a volume
factor appears, the extraction must be rewritten before applying
the V61 gate.
\end{decision}

\begin{firewall}[Claim boundary after V62]
\label{fire:e8v62-claim}
V62 proves an exact geometric multi-block energy cost, not yet the
complete normalized polymer expansion.  It does not verify the
activity-extraction hypothesis of (62.18), bound
\(u_g(\beta)\), close (61.19), evaluate the raw pushed
\(\zeta_{\rm Jac}\), compare the compensated owner with
\({\cal O}_1\), prove anisotropy contraction, evaluate the remaining
cards, control overlap owners, complete the LERB, establish the
universal beta coefficient, close the nonlinear/complex quotient
matching card, construct the full RG trajectory, remove the
regulators, establish Osterwalder--Schrader reconstruction, or prove
the Yang--Mills mass gap.
\end{firewall}

\subsection{Parent record}

This V62 note was formed by appending the preceding material to V61,
whose validated record was
\[
\begin{array}{c|c}
\text{lines}&26501\\
\text{bytes}&930400\\
\text{SHA-256}&
\texttt{BA70D472BFA0A1D39BFD13A819EC7BE9A7654747FBF1EB4A5C561ED1E1229F3D}.
\end{array}
\]

% =============================================================================
% END APPENDED EIGHTH-SERIES V62 MATERIAL
% =============================================================================

% =============================================================================
% BEGIN APPENDED EIGHTH-SERIES V63 MATERIAL
% =============================================================================

\section{Eighth-series V63: a dominated product-Haar reference and
volume-free bad-block normalization}
\label{sec:e8v63}

\subsection{Why the V62 provisional ledger is replaced}

Proposition~\ref{prop:e8v62-bad-ledger} gave a valid conditional
bound if a contour extractor spent as many as sixteen forty-five-link
normalizers per bad block.  Its prefactor grows like
\(\beta^{2880}\).  More importantly, it did not construct the
reference measure whose global normalizer cancels before the
polymer estimate.  The moving-section coercivity (62.6) permits a
cleaner exact choice.

Let \(\Lambda_c\) be the coarse block set and put
\[
                         t:=\frac\beta{120}.                  \tag{63.1}
\]
For a reduced link \(e\), define the translated chord measure
\[
 d\gamma_{t,H_e(W)}(z_e)
 :=I(t)^{-1}e^{-t\|z_e-H_e(W)\|_F^2}\,dz_e,                  \tag{63.2}
\]
where \(dz_e\) is normalized Haar measure and \(I(t)\) is (62.16).
Left or right Haar invariance makes its normalizer independent of
\(H_e(W)\).  Let \(\gamma_{t,W}\) be the product of (63.2) over
the \(45|\Lambda_c|\) reduced links.

\begin{theorem}[Exact dominated-reference factorization]
\label{thm:e8v63-dominated-reference}
Define
\[
 R_{\beta,W}(z):=\beta V_M(z,W)-tD_W(z).                     \tag{63.3}
\]
Then
\[
                         R_{\beta,W}(z)\ge0                  \tag{63.4}
\]
pointwise, and the exact reduced-fibre density factors as
\[
\boxed{
 e^{-\beta V_M(z,W)}dm_{\rm red}(z)
 =I(t)^{45|\Lambda_c|}
  e^{-R_{\beta,W}(z)}d\gamma_{t,W}(z).}                      \tag{63.5}
\]
Consequently, for every observable \(F\),
\[
 \langle F\rangle_{\beta,W}
 =\frac{\mathbb E_{\gamma_{t,W}}
       [F e^{-R_{\beta,W}}]}
       {\mathbb E_{\gamma_{t,W}}[e^{-R_{\beta,W}}]}.        \tag{63.6}
\]
The extensive factor in (63.5) cancels exactly before any
correlation or current estimate.
\end{theorem}

\begin{proof}
Theorem~\ref{thm:e8v62-moving-coercivity} gives
\(V_M\ge D_W/120\), which is (63.4).  From (63.2),
\[
 dm_{\rm red}
 =I(t)^{45|\Lambda_c|}e^{tD_W}d\gamma_{t,W}.
\]
Multiply by \(e^{-\beta V_M}\) and use (63.3), proving (63.5).
Division of the two instances of (63.5), with and without \(F\),
gives (63.6).
\end{proof}

\begin{firewall}[The denominator is handled by connected expansion]
\label{fire:e8v63-no-denominator-bound}
Equation (63.6) is not an invitation to lower-bound its denominator
by a number uniform in the volume.  Source derivatives are to be
taken from the logarithm of the same \(\gamma_{t,W}\) expectation,
as in (60.4) and V61.  The useful new fact is that the reference is a
normalized product probability and \(0\le e^{-R}\le1\).
\end{firewall}

\subsection{An exact smooth good/bad decomposition}

Fix \(r>0\).  Choose smooth functions
\(\chi_{\rm g},\chi_{\rm b}:[0,\infty)\to[0,1]\) such that
\[
 \chi_{\rm g}+\chi_{\rm b}=1,
 \quad \chi_{\rm g}=1\ \hbox{on }[0,r^2],
 \quad \chi_{\rm g}=0\ \hbox{on }[4r^2,\infty).            \tag{63.7}
\]
In particular
\(\operatorname{supp}\chi_{\rm b}\subset[r^2,\infty)\).
For \(B\subset\Lambda_c\), set
\[
 \chi_B(z;W):=
 \prod_{y\in B}\chi_{\rm b}(D_y(z;W))
 \prod_{y\notin B}\chi_{\rm g}(D_y(z;W)).                  \tag{63.8}
\]

\begin{proposition}[Exact selector identity]
\label{prop:e8v63-selector}
Pointwise on the reduced fibre,
\[
                         1=\sum_{B\subset\Lambda_c}\chi_B.  \tag{63.9}
\]
The identity remains valid after any retained-background or local
source derivative; the derivative acts on the displayed factors and
does not create an unassigned remainder.
\end{proposition}

\begin{proof}
Expand
\(\prod_y\{\chi_{\rm g}(D_y)+\chi_{\rm b}(D_y)\}\) and use
(63.7).  Since the identity holds as a smooth identity in all of its
arguments, differentiation preserves it.
\end{proof}

The block variables in distinct \({\cal C}_y\) are disjoint, so
they are independent under \(\gamma_{t,W}\).  Translation by
\(H_e(W)\) also shows that the law of \(D_y\) is independent of
\(y\) and \(W\).

\subsection{A volume-free bad-set bound}

Define
\[
 q_b(\beta,r):=
 C_-^{-45}t^{180}e^{-tr^2},
 \qquad t=\beta/120,                                        \tag{63.10}
\]
with the explicit \(C_-\) from (62.16).

\begin{lemma}[One-block chord tail]
\label{lem:e8v63-chord-tail}
For \(\beta\ge120\),
\[
 \gamma_{t,W}\{D_y\ge r^2\}
 \le q_b(\beta,r).                                          \tag{63.11}
\]
\end{lemma}

\begin{proof}
On \(D_y\ge r^2\),
\(e^{-tD_y}\le e^{-tr^2}\).  The unnormalized Haar measure of the
event is at most one.  Since the block contains forty-five
independent links, its reference normalizer is \(I(t)^{45}\).
Therefore
\[
 \gamma_{t,W}\{D_y\ge r^2\}
 \le I(t)^{-45}e^{-tr^2}.
\]
For \(t\ge1\), insert \(I(t)\ge C_-t^{-4}\), obtaining (63.11).
\end{proof}

\begin{theorem}[Joint bad-block domination]
\label{thm:e8v63-joint-bad}
For every block set \(Y\subset\Lambda_c\),
\[
\boxed{
 \mathbb E_{\gamma_{t,W}}
 \left[e^{-R_{\beta,W}}
       \prod_{y\in Y}\chi_{\rm b}(D_y)\right]
 \le q_b(\beta,r)^{|Y|}.}                                   \tag{63.12}
\]
The bound is uniform in the volume and in every real retained
background \(W\).  If bounded local sources obey
\(|F_i|\le M_i\) and \(|s_i|\le\rho_i\), insertion of
\(e^{\sum_i s_iF_i}\) multiplies the right side by at most
\(e^{\sum_i\rho_iM_i}\).
\end{theorem}

\begin{proof}
By (63.4), \(e^{-R}\le1\).  By (63.7), each bad selector is
bounded by the indicator of \(D_y\ge r^2\).  The sets of forty-five
link variables belonging to distinct blocks are independent under
the product reference.  Apply Lemma~\ref{lem:e8v63-chord-tail} in
each marked block.  The source statement follows from the pointwise
bound on its exponential.
\end{proof}

This is the desired normalization repair.  It improves the
provisional V62 power \(\beta^{2880}\) to \(\beta^{180}\), and it
does so before division by a partition function.

\subsection{The correct contour adjacency and entropy constant}

The filling owner of a bad block lies in
\(y+\{0,1\}^4\) by Lemma~\ref{lem:e8v62-filling-support}.
Therefore define the contour graph \({\cal G}_*\) by
\[
 y\sim_*y'\quad\Longleftrightarrow\quad
 0<\|y-y'\|_\infty\le1.                                    \tag{63.13}
\]
It has maximum degree
\[
                         \Delta_*=3^4-1=80.                  \tag{63.14}
\]
Connected components of a bad selector set \(B\) are henceforth
taken in \({\cal G}_*\).  Distinct components then have disjoint
filling neighbourhoods.

\begin{proposition}[Correct entropy gate for filling contours]
\label{prop:e8v63-contour-entropy}
The nearest-neighbour numerical constant \(65\) in (61.6) applies
only to primitive activities supported on the degree-eight graph.
For contours carrying the filling neighbourhood (62.10), the same
proved animal bound gives the sufficient condition
\[
 \boxed{6401e^{1+\mu_*}u\le1,}                               \tag{63.15}
\]
where \(\mu_*\) is measured in the \({\cal G}_*\) distance.
\end{proposition}

\begin{proof}
Use Lemma~\ref{lem:e8v61-KP-gate} with \(a=1\) and
\(\Delta=\Delta_*=80\).  Then \(\Delta_*^2+1=6401\).
\end{proof}

Thus V61's theorem was correct for its declared graph, but an
application to the actual filling contours must spend the larger
finite-range entropy.  Since one \({\cal G}_*\) step moves at most
one coarse lattice unit in each coordinate, exponential decay in
that metric still implies physical lattice exponential decay after
a fixed change of rate.

\subsection{Exact component ownership and the remaining good factor}

For a selector set \(B\), write
\(B=Y_1\mathbin{\dot\cup}\cdots\mathbin{\dot\cup}Y_k\) for its
\({\cal G}_*\)-connected components.  By construction,
\({\cal P}(Y_i)\cap{\cal P}(Y_j)=\varnothing\) for \(i\ne j\).
Assign a plaquette lying in one of these sets to that unique bad
component; assign every other plaquette to the good sea.  Assign a
local source to the lexicographically first block of its support,
and enlarge the component if that support meets its filling
neighbourhood.  This is a deterministic, nonduplicating ownership
rule.

The exact selector sum (63.9) and the ownership rule do not yet make
the weight a product over the \(Y_i\), because the residual
\(R_{\beta,W}\) contains local interactions in the good sea which
can join two components.  Those interactions must be expanded into
connected good polymers.  Once an exact expansion supplies source-
analytic activities with
\[
                         \|z^{\rm g}_t(X)\|
                         \le u_g(\beta,r)^{|X|},              \tag{63.16}
\]
the combined incidence graph has two kinds of primitive objects:
bad marked blocks, bounded by (63.12), and good interaction
polymers, bounded by (63.16).  Connected-component collection then
gives the V61 majorant with
\[
                         u(\beta,r)
                         =q_b(\beta,r)+u_g(\beta,r).          \tag{63.17}
\]

\begin{theorem}[Revised explicit four-mark gate]
\label{thm:e8v63-revised-gate}
If the source-dependent good residual admits the exact activity
bound (63.16) on the contour graph and
\[
\boxed{
 6401e^{1+\mu_*}
 \left\{
 C_-^{-45}\left(\frac\beta{120}\right)^{180}
 e^{-\beta r^2/120}
 +u_g(\beta,r)\right\}\le1,}                                \tag{63.18}
\]
then the conclusions of
Corollary~\ref{cor:e8v61-third-current} hold for the full
good/bad reduced-fibre representation.
\end{theorem}

\begin{proof}
The first primitive activity is bounded by
Theorem~\ref{thm:e8v63-joint-bad}; the second is (63.16).
The deterministic ownership rule prevents duplicate plaquettes or
sources.  Summing the two majorants gives (63.17).  Condition
(63.18) is precisely (63.15), so apply
Theorem~\ref{thm:e8v61-four-mark-cluster} and its current
corollary.
\end{proof}

\begin{firewall}[Domination does not make the good residual small]
\label{fire:e8v63-good-not-small}
The pointwise inequality \(0\le e^{-R}\le1\) is enough for the bad
tail (63.12), but not for (63.16).  The crude reference precision
\(1/60\) touches only the global lower Hessian edge and need not be a
small relative perturbation of the full vertical Wilson Hessian.
The good activity must use the sharper Hessian/covariance chart and
the analytic Wilson jets already developed in V41--V55; it cannot be
set to zero because the bad estimate is strong.
\end{firewall}

\begin{assessment}[Progress after V63]
\label{ass:e8v63-status}
V63 constructs an exact retained-background product-Haar reference
whose extensive normalizer cancels identically.  It proves the
volume-free joint bad-block bound (63.12), replaces the provisional
bad activity by the sharper explicit function (63.10), constructs a
smooth exact selector and a nonduplicating component ownership rule,
and corrects the actual contour entropy constant from \(65\) to
\(6401\).  The nonlinear four-point gate is now reduced to the single
good-field activity number \(u_g(\beta,r)\).
\end{assessment}

\begin{decision}[V64 target]
\label{dec:e8v63-next}
V64 should construct \(u_g(\beta,r)\) by rescaling the exact
stationary vertical variable by \(\beta^{-1/2}\), using the V41
precision, the V45 weighted covariance, and the V54 Wilson jet
bounds.  It must keep the smooth cutoff derivatives and background
scores, and it must decide whether a fixed \(r\) gives a genuinely
small activity or whether a nested small/intermediate-field
decomposition is required.
\end{decision}

\begin{firewall}[Claim boundary after V63]
\label{fire:e8v63-claim}
V63 closes bad-block normalization but not the good-field expansion.
It does not prove (63.16), close (63.18), evaluate the raw pushed
\(\zeta_{\rm Jac}\), compare the compensated owner with
\({\cal O}_1\), prove anisotropy contraction, evaluate the remaining
cards, complete the LERB, establish the universal beta coefficient,
close the nonlinear/complex quotient matching card, construct the
full RG trajectory, remove the regulators, establish
Osterwalder--Schrader reconstruction, or prove the Yang--Mills mass
gap.
\end{firewall}

\subsection{Parent record}

This V63 note was formed by appending the preceding material to V62,
whose validated record was
\[
\begin{array}{c|c}
\text{lines}&26856\\
\text{bytes}&943934\\
\text{SHA-256}&
\texttt{20E071116E134593C194291187968E80EC8E1B3ED911D7EAFF472BAF8A5C6568}.
\end{array}
\]

% =============================================================================
% END APPENDED EIGHTH-SERIES V63 MATERIAL
% =============================================================================

% =============================================================================
% BEGIN APPENDED EIGHTH-SERIES V64 MATERIAL
% =============================================================================

\section{Eighth-series V64: the running small-field window and
simultaneous good/bad decay}
\label{sec:e8v64}

\subsection{The fixed-radius tension}

V63 deliberately left the good activity \(u_g(\beta,r)\) open.  A
fixed unscaled radius \(r\) makes the bad factor exponentially
small, but a supremum bound for the cubic Taylor remainder of
\(\beta V_M\) is then of order \(\beta r^3\), not small.  Taking
\(r\) of Gaussian size \(\beta^{-1/2}\) makes the Taylor expansion
uniform but gives no decaying bad-field exponent.  The two demands
are reconciled by an intermediate running radius.

For \(\alpha\in(0,1/2)\), put
\[
                         r_\beta:=\beta^{-\alpha}.            \tag{64.1}
\]

\begin{lemma}[The scale window]
\label{lem:e8v64-scale-window}
Suppose a local stationary action owner has a uniform third
derivative bound on a fixed logarithm chart.  On
\(\|x\|\le C r_\beta\), its scaled cubic Taylor remainder obeys
\[
                  \beta|R_3(x)|\le C_3'\beta^{1-3\alpha}.    \tag{64.2}
\]
The V63 bad exponent is
\[
                         t r_\beta^2
                         =\frac1{120}\beta^{1-2\alpha}.      \tag{64.3}
\]
Thus the open interval
\[
                         \boxed{\frac13<\alpha<\frac12}      \tag{64.4}
\]
makes the good Taylor factor tend to zero and the bad suppression
tend to infinity simultaneously.
\end{lemma}

\begin{proof}
Taylor's formula at a stationary point gives
\(|R_3(x)|\le C_3\|x\|^3/6\), which is (64.2) after inserting
(64.1) and absorbing the fixed local dimension.  Equation (64.3)
is (63.1).  Its exponent grows exactly when \(\alpha<1/2\), while
(64.2) decays exactly when \(\alpha>1/3\).
\end{proof}

We henceforth choose the rational interior value
\[
                         \boxed{\alpha=\frac25}.             \tag{64.5}
\]
Then the good supremum debit is \(O(\beta^{-1/5})\), and the bad
exponential is \(e^{-\beta^{1/5}/120}\).

\subsection{Passage to the weak-field normalization}

Write a reduced link relative to the moving section as
\[
 z_e=e^{g a_e}H_e(W),\qquad g:=\beta^{-1/2}.                 \tag{64.6}
\]
The quantitative Wilson remainder estimate in this subsection is
made at a coarse-flat retained background (including a fixed
commuting toron), where \(H(W)\) is the unique stationary fibre
minimum by Theorem~\ref{thm:e8v62-zero-set}.  A non-flat complex
background polydisc requires transport to the analytic stationary
centre; that additional stock is assigned to the forest constant
below and is not inferred from (64.6).
On the support of \(\chi_{\rm g}(D_y)\), (63.7) gives
\(D_y\le4r_\beta^2\).  For \(\beta\) large enough that
\(2r_\beta<2\), every relative link lies in the principal
logarithm chart.  The scalar inequality
\(|\theta|\le(\pi/2)|e^{i\theta}-1|\), \(|\theta|\le\pi\),
applied to eigenangles gives
\[
 \|a_e\|_F
 \le\frac\pi{2g}\|z_eH_e(W)^{-1}-I\|_F
 \le\pi\beta^{1/10}.                                       \tag{64.7}
\]

\begin{lemma}[Explicit local rescaled-field stock]
\label{lem:e8v64-field-stock}
Every filling neighbourhood contains at most
\[
                         n_*=45\cdot16=720                  \tag{64.8}
\]
reduced links.  On its good support,
\[
 \|a\|_{2,{\rm loc}}^2
 \le720\pi^2\beta^{1/5},\qquad
 g\|a\|_{\infty,{\rm loc}}
 \le\pi\beta^{-2/5}.                                       \tag{64.9}
\]
\end{lemma}

\begin{proof}
The support enlargement has at most sixteen coarse blocks by
(62.11), each with forty-five reduced links.  Apply (64.7) to each
of those links and sum the squares.  The second estimate follows
directly from \(g=\beta^{-1/2}\).
\end{proof}

\subsection{The loaded Wilson remainder estimate}

V23 already proves for the normalized unit Wilson action
\({\cal V}_g(a)=g^{-2}V_M(e^{ga})\) that, on
\(g\|a\|_\infty\le1\),
\[
 |{\cal V}_g(a)-{\cal V}_g(0)
   -\tfrac12\langle a,Aa\rangle|
 \le\frac{32}{3}g\|a\|_\infty\|a\|_2^2.                  \tag{64.10}
\]
This is imported from the existing note, not derived again.

\begin{proposition}[Explicit good-neighbourhood remainder at the
flat retained base]
\label{prop:e8v64-good-remainder}
On a filling neighbourhood on which every block selector is good,
the local Wilson nonlinear remainder obeys the deliberately coarse
bound
\[
                         |R_{\rm nl}^{\rm loc}|
                         \le\eta_g(\beta),                   \tag{64.11}
\]
where
\[
 \boxed{
 \eta_g(\beta):=7680\pi^3\beta^{-1/5}.}                     \tag{64.12}
\]
Consequently its primitive exponential factor satisfies
\[
 |e^{-R_{\rm nl}^{\rm loc}}-1|
 \le e^{\eta_g(\beta)}-1.                                  \tag{64.13}
\]
\end{proposition}

\begin{proof}
For \(\beta\) large enough, (64.9) makes the chart condition in
(64.10) valid.  Insert both estimates of (64.9) into the right side
of (64.10):
\[
 \frac{32}{3}
 (\pi\beta^{-2/5})(720\pi^2\beta^{1/5})
 =7680\pi^3\beta^{-1/5}.
\]
The elementary inequality
\(|e^{-u}-1|\le e^{|u|}-1\) gives (64.13).
\end{proof}

The constant in (64.12) is intentionally large because it uses the
sixteen-block filling neighbourhood and takes absolute values before
all Ward cancellations.  Its importance is that it is explicit,
volume independent, and tends to zero.

\subsection{Cutoff derivatives do not consume powers of
\texorpdfstring{\(r_\beta^{-1}\)}{r-beta inverse}}

A naive background derivative of
\(\chi_{\rm g}(D_y(z;W))\) appears to cost a power of
\(r_\beta^{-1}\).  This is a coordinate artefact.  Make on each
reduced link the Haar translation
\[
                         \zeta_e=z_eH_e(W)^{-1}.              \tag{64.14}
\]
Then
\[
 D_y(z;W)=\sum_{e\in{\cal C}_y}\|\zeta_e-I\|_F^2,           \tag{64.15}
\]
and the reference measure (63.2) is independent of \(W\) in the
\(\zeta\) variables.

\begin{proposition}[Stationary selector convention]
\label{prop:e8v64-stationary-selector}
If all retained-background derivatives are taken after the exact
change (64.14), neither the selector (63.8) nor the product reference
measure contributes a background score.  All background scores lie
in the transported Wilson residual, block map, current insertion,
and any generated action.  Hence the first three current jets do not
acquire artificial factors \(r_\beta^{-1}\) from differentiating
the cutoff.
\end{proposition}

\begin{proof}
Equation (64.15) removes \(W\) from every selector argument.
Haar invariance removes it from the product differential, and the
normalizer \(I(t)\) is constant.  Differentiation at fixed
\(\zeta\) therefore annihilates both objects.  The chain rule moves
the derivative to every remaining occurrence of the translated
link, which are precisely the listed physical score terms.
\end{proof}

\begin{firewall}[Moving coordinates must carry their current]
\label{fire:e8v64-moving-current}
The translation (64.14) is used only to organize the conditional
integral and its derivatives.  Its induced retained-coordinate
current and Haar Jacobian are already governed by the pushforward
identity of V58.  Treating \(\zeta\) as fixed while also treating
\(z\) as fixed would omit that chain-rule current.
\end{firewall}

Proposition~\ref{prop:e8v64-stationary-selector} removes a cutoff
derivative, not the stationary-centre problem.  For a non-flat
background source, the action must be recentered at the analytic
critical point already used in the V44--V45 covariance chart.  Its
first three background jets and the Cauchy-radius loss remain part of
the marked forest extraction.

\subsection{Asymptotic form of the complete gate}

With (64.5), the bad primitive of V63 becomes
\[
 q_b(\beta):=
 C_-^{-45}\left(\frac\beta{120}\right)^{180}
 \exp\left\{-\frac1{120}\beta^{1/5}\right\}.               \tag{64.16}
\]
It tends to zero, although only beyond an enormous numerical scale
because the present constants are crude.

The remaining conversion is from local nonlinear factors such as
(64.13), integrated against the Gaussian precision \(A\), to
hard-core connected activities.  Denote by \(C_{\rm GF}\) the norm
of a source-marked Gaussian forest extraction in the V45 weighted
covariance algebra, including the finite overlap, analytic stationary-
centre transport, and Cauchy-radius stocks.  The required output is
\[
 u_g(\beta)
 \le C_{\rm GF}\{e^{\eta_g(\beta)}-1\}.                     \tag{64.17}
\]

\begin{proposition}[Eventual closure once the forest constant is loaded]
\label{prop:e8v64-eventual-gate}
If \(C_{\rm GF}<\infty\) is independent of the volume and \(\beta\),
and the exact source-marked forest extraction proves (64.17), then
for every fixed \(\mu_*>0\) there exists \(\beta_0<\infty\) such
that, for \(\beta\ge\beta_0\),
\[
 6401e^{1+\mu_*}
 \left[q_b(\beta)+
 C_{\rm GF}\{e^{7680\pi^3\beta^{-1/5}}-1\}\right]\le1.     \tag{64.18}
\]
Thus the four-mark gate (63.18) closes at sufficiently weak bare
coupling once this single forest constant is proved finite.
\end{proposition}

\begin{proof}
The exponential in (64.16) dominates its fixed polynomial, so
\(q_b(\beta)\to0\).  Equation (64.12) tends to zero, and hence so
does the second term in brackets in (64.18).  Multiplication by the
fixed entropy factor preserves convergence to zero.  The definition
of convergence supplies a finite \(\beta_0\).
\end{proof}

\begin{firewall}[A named forest constant is not yet a theorem]
\label{fire:e8v64-forest-debit}
V45 proves that the Gaussian covariance belongs to a weighted
convolution algebra, and V54 proves finite local derivative stocks.
Those facts strongly constrain \(C_{\rm GF}\), but they do not by
themselves prove the all-order, source-marked forest reorganization
needed for (64.17).  Proposition~\ref{prop:e8v64-eventual-gate} is
conditional on that extraction and is not a declaration that the
nonlinear cluster theorem has already been proved.
\end{firewall}

\begin{assessment}[Progress after V64]
\label{ass:e8v64-status}
V64 resolves the incompatible fixed-radius estimates by the exact
window \(1/3<\alpha<1/2\), chooses \(\alpha=2/5\), and proves both
an explicit \(O(\beta^{-1/5})\) local Wilson remainder at the flat
or commuting-toron retained base and an
explicit \(\beta^{180}e^{-\beta^{1/5}/120}\) bad primitive.  It
also removes spurious cutoff-derivative losses through the moving
Haar coordinate.  The nonlinear cluster gate is reduced to an
all-order Gaussian forest extraction with a volume-independent
constant \(C_{\rm GF}\).
\end{assessment}

\begin{decision}[Mandatory V65 audit and next target]
\label{dec:e8v64-next}
V65 must inspect the then-current SM and NS companion notes.  It
should look for a marked forest formula, a connected Gaussian
activity estimate, or an analytic determinant-word construction
that can be adapted to prove (64.17).  If none is present, the YM
programme should prove a finite-degree source-marked forest lemma
directly, first for the four cumulants required by the current
Jacobian rather than for an unnecessary all-order observable bank.
\end{decision}

\begin{firewall}[Claim boundary after V64]
\label{fire:e8v64-claim}
V64 proves simultaneous decay of the primitive good and bad
quantities but not the forest conversion (64.17).  It does not yet
close the nonlinear four-point cluster gate, evaluate the raw pushed
\(\zeta_{\rm Jac}\), compare the compensated owner with
\({\cal O}_1\), prove anisotropy contraction, evaluate the remaining
cards, complete the LERB, establish the universal beta coefficient,
close the nonlinear/complex quotient matching card, construct the
full RG trajectory, remove the regulators, establish
Osterwalder--Schrader reconstruction, or prove the Yang--Mills mass
gap.
\end{firewall}

\subsection{Parent record}

This V64 note was formed by appending the preceding material to V63,
whose validated record was
\[
\begin{array}{c|c}
\text{lines}&27204\\
\text{bytes}&956642\\
\text{SHA-256}&
\texttt{2E21B8FC1A2E0E0D18A25E5B7259F3C4E81F161808B4DB3B2BED7E90355D182F}.
\end{array}
\]

% =============================================================================
% END APPENDED EIGHTH-SERIES V64 MATERIAL
% =============================================================================

% =============================================================================
% BEGIN APPENDED EIGHTH-SERIES V65 MATERIAL
% =============================================================================

\section{Eighth-series V65: compulsory audit, transfer of the
Gaussian forest theorem, and an analytic selector repair}
\label{sec:e8v65}

\subsection{Companion files inspected}

The live companions at the end of this audit were
\[
\begin{array}{c|r|r|l}
\text{file}&\text{lines}&\text{bytes}&\text{SHA--256}\\ \hline
\texttt{Renewal\_SM\_Einstein\_Continuum\_Research\_Notes\_V60.tex}
 &21252&800272&
 \texttt{799F08670BEA0696125B2216A33BB291D1BECF2E00FF3FCC4028DAB1F2F1AB62}
 \\
\texttt{Renewal\_NS\_Stability\_Research\_Notes\_V26.tex}
 &5320&278283&
 \texttt{82C887F8C2C69E4F28FAD7C3DC8DE5B9099CB5A4BF431EBA1FFF3E91935978AD}.
\end{array}                                                     \tag{65.1}
\]
The SM file itself advanced during the audit; the final row above is
the later V60 record.  Its current V58--V60 material independently
confirms that a third conditional-current derivative can require a
fourth cumulant and that moving coordinates remove artificial
selector derivatives but not the physical chart and action scores.
The NS file remains unchanged and supplies no forest or compact-link
input.

The decisive transferable result is earlier in the same SM file.
Its V21 proves the Brydges--Kennedy-type positive forest
interpolation, the spanning-tree formula for Gaussian cumulants, the
Cayley bound, and a connected activity estimate.  Those statements
depend only on positivity of the reduced Gaussian covariance and on
local analytic vertex norms; they do not use the Standard Model field
content.

\begin{assessment}[V65 audit decision]
\label{ass:e8v65-audit-decision}
The SM V21 positive-Gaussian forest and tree-summation theorems
transfer to the reduced Yang--Mills covariance after the V41 gauge
quotient and V45 weighted inverse.  The SM gravitational chart-score
stocks and its running-interface constants do not transfer.  The NS
rank-one estimates have already been accounted for and add no new
row.
\end{assessment}

\subsection{The transferred forest estimate}

Let \(C\) be a positive Gaussian covariance on the reduced physical
space.  For \(n\) replicated fields and a spanning tree \(T\), write
\[
 \Delta_{ij,C}:=\sum_{a,b}C_{ab}
 \frac{\partial}{\partial\phi_{i,a}}
 \frac{\partial}{\partial\phi_{j,b}}.                        \tag{65.2}
\]
The imported interpolation gives, for local analytic vertices
\(F_1,\ldots,F_n\),
\[
 \kappa_C(F_1,\ldots,F_n)
 =\sum_{T\ {\rm tree}}\int_{[0,1]^T}
 \mathbb E_{C^T(w)}
 \left[
  \prod_{\{i,j\}\in T}\Delta_{ij,C}
  \prod_{i=1}^nF_i(\phi_i)
 \right]dw.                                                   \tag{65.3}
\]
Positivity of every weakened covariance \(C^T(w)\) is part of the
imported theorem.

\begin{theorem}[Transferred connected-tree bound]
\label{thm:e8v65-transferred-tree}
Assume a source-marked Yang--Mills vertex has strong analytic/polymer
norm at most \(z\), and every fixed labelled tree integrand in
(65.3) obeys
\[
 \|{\cal A}_T(F_1,\ldots,F_n)\|_{\rm weak}
 \le z^nK_T^{n-1},                                           \tag{65.4}
\]
where \(K_T\) includes the V45 weighted covariance edge, finite
owner types, and the one-time analytic-radius loss.  If
\[
                         q_T:=ezK_T<1,                        \tag{65.5}
\]
then the complete connected exponential satisfies
\[
\boxed{
 \left\|\sum_{n\ge1}\frac1{n!}
 \kappa_C(F,\ldots,F)\right\|_{\rm weak}
 \le z+\frac{e^2z^2K_T}{1-ezK_T}.}                           \tag{65.6}
\]
The same estimate is analytic in four additional marked sources on
every smaller source polydisc on which (65.4)--(65.5) are uniform.
\end{theorem}

\begin{proof}
This is the SM V21 theorem specialized to bosonic reduced
Yang--Mills vertices.  For completeness, the only counting step is
\[
 \frac{|{\cal T}_n|}{n!}
 =\frac{n^{n-2}}{n!}\le e^n.
\]
Combining it with (65.4) makes the \(n\)-th term, \(n\ge2\), at
most \(K_T^{-1}(ezK_T)^n\); summing the geometric series gives
(65.6).  Uniform normal convergence followed by multivariate Cauchy
gives the four marked-source statement.  The forest identity and
positivity needed to justify (65.3) are exactly the imported parts
and are not repeated here.
\end{proof}

\begin{corollary}[Explicit replacement for the abstract forest
constant]
\label{cor:e8v65-forest-replacement}
If the good Wilson vertices can be placed in (65.4) with
\[
 z_\beta\le C_{\rm an}
 \{e^{7680\pi^3\beta^{-1/5}}-1\},                            \tag{65.7}
\]
then the abstract expression (64.17) may be replaced by
\[
 u_g(\beta)\le C_{\rm own}\left\{
 z_\beta+\frac{e^2z_\beta^2K_T}{1-ez_\beta K_T}
 \right\}.                                                   \tag{65.8}
\]
For fixed finite \(C_{\rm an},C_{\rm own},K_T\), this tends to
zero and satisfies the branch condition for all sufficiently large
\(\beta\).
\end{corollary}

\begin{proof}
Apply Theorem~\ref{thm:e8v65-transferred-tree} and the deterministic
owner-to-activity conversion \(C_{\rm own}\).  Equation (64.12)
makes the right side of (65.7) tend to zero; the claims follow by
continuity at \(z=0\).
\end{proof}

\subsection{The compact-selector obstruction}

The smooth selector (63.7) is adequate for the exact bad-set
identity and for finitely many ordinary derivatives.  It cannot be
inserted without modification into the analytic all-tree norm used
in (65.4).

\begin{proposition}[No nontrivial analytic compact selector]
\label{prop:e8v65-no-analytic-cutoff}
There is no real-analytic function on a connected interval which is
one on a nonempty open subinterval and zero on another nonempty open
subinterval.  In particular the compact transition profile (63.7)
is not analytic.
\end{proposition}

\begin{proof}
If an analytic function equals one on an open interval, all Taylor
coefficients of its difference from one vanish at an interior point.
The identity theorem makes it identically one on the connected
domain, contradicting its vanishing on the other interval.
\end{proof}

\begin{firewall}[Finite differentiability does not pass an all-tree
card]
\label{fire:e8v65-selector-forest}
The current Jacobian needs only four marked observables, but the
interacting exponential contains arbitrarily many unmarked
vertices.  A star tree can differentiate one selector to arbitrarily
high order.  Bounds through the third or fourth derivative therefore
do not prove (65.4).  The selector must be treated by a separate
boundary/large-field argument or replaced by an analytic soft split.
\end{firewall}

\subsection{An exact analytic soft split}

The dominated product reference of V63 permits the second option.
Keep \(r_\beta=\beta^{-2/5}\) and define, for one block,
\[
 g_{\beta,y}:=\exp\{-D_y/r_\beta^2\},\qquad
 b_{\beta,y}:=1-g_{\beta,y}.                                 \tag{65.9}
\]
Both are real analytic on the compact group and holomorphic on the
fixed complexified \(SL(3,\mathbb C)\) chart used by the forest
norm, satisfy \(g_{\beta,y}+b_{\beta,y}=1\), and obey
\[
                         0\le b_{\beta,y}\le D_y/r_\beta^2. \tag{65.10}
\]
Thus the subset expansion (63.9) remains exact with
\(\chi_{\rm g},\chi_{\rm b}\) replaced by \(g,b\), but no
compactly supported derivative enters a forest.

Let the one-link constants from the attachment be
\[
 C_-:=\frac{3e^{-43/32}}{1024\pi^8},\qquad
 C_+:=\frac{25\pi^7}{512}.                                  \tag{65.11}
\]

\begin{lemma}[First chord moment of the reference]
\label{lem:e8v65-chord-moment}
For \(t\ge2\) and the one-link probability (63.2),
\[
 \mathbb E_{\gamma_t}\|U-I\|_F^2
 \le\frac{32C_+}{eC_-}\,t^{-1}.                             \tag{65.12}
\]
\end{lemma}

\begin{proof}
For \(d\ge0\),
\[
 d e^{-td}
 \le\frac2{et}e^{-(t/2)d},
\]
because the maximum of \(d e^{-(t/2)d}\) is \(2/(et)\).
Integrate and use the attachment's bounds
\(I(t/2)\le C_+(t/2)^{-4}\) and
\(I(t)\ge C_-t^{-4}\).  Their ratio gives (65.12).
\end{proof}

\begin{theorem}[Soft-bad primitive]
\label{thm:e8v65-soft-bad}
For \(t=\beta/120\), \(\beta\ge240\), and every block set \(Y\),
\[
\boxed{
 \mathbb E_{\gamma_{t,W}}
 \left[e^{-R_{\beta,W}}
       \prod_{y\in Y}b_{\beta,y}\right]
 \le q_s(\beta)^{|Y|},}                                     \tag{65.13}
\]
where
\[
 \boxed{
 q_s(\beta):=
 \frac{172800C_+}{eC_-}\,\beta^{-1/5}.}                     \tag{65.14}
\]
This is volume independent, uniform in real \(W\), and analytic in
the link variables.
\end{theorem}

\begin{proof}
The block contains forty-five independent links.  Equations
(65.10) and (65.12) give
\[
 \mathbb E_{\gamma_{t,W}}b_{\beta,y}
 \le\frac{45}{r_\beta^2}
       \frac{32C_+}{eC_-}t^{-1}.
\]
Insert \(r_\beta^2=\beta^{-4/5}\) and
\(t^{-1}=120\beta^{-1}\), obtaining (65.14).
For several blocks use independence of their disjoint link sets and
the pointwise domination \(e^{-R}\le1\), exactly as in (63.12).
\end{proof}

The soft primitive decays only polynomially, whereas (64.16) has a
stretched-exponential tail.  The gain is analytic compatibility.
One may retain the stretched-exponential estimate for genuinely far
blocks and use (65.14) only for the analytic transition, but they
must be entered as distinct primitive types in the activity ledger.

\begin{proposition}[Soft all-good factor is perturbative in rescaled
variables]
\label{prop:e8v65-soft-good}
In the weak-field coordinate (64.6),
\[
 \prod_y g_{\beta,y}
 =\exp\left\{-\beta^{-1/5}
 \sum_y Q_y(a)+O(\beta^{-3/5}\|a\|_{\rm loc}^3)\right\},     \tag{65.15}
\]
where \(Q_y\) is the positive quadratic chord form at the identity.
Thus the soft selector adds a local analytic quadratic vertex of
order \(\beta^{-1/5}\), the same order as the Wilson remainder in
(64.12).
\end{proposition}

\begin{proof}
With \(g=\beta^{-1/2}\), the chord expansion is
\[
 D_y(e^{ga}H;W)=g^2Q_y(a)+O(g^3\|a\|_{\rm loc}^3).
\]
Divide by \(r_\beta^2=\beta^{-4/5}\).  The quadratic coefficient is
\(\beta^{-1}/\beta^{-4/5}=\beta^{-1/5}\), and the cubic coefficient
is \(\beta^{-3/2}/\beta^{-4/5}=\beta^{-7/10}\), which is bounded by
the displayed, deliberately weaker \(O(\beta^{-3/5})\) on the
declared local analytic ball.  Summing the exponents of the product
gives (65.15).
\end{proof}

\subsection{Post-audit direction}

The transferred theorem now gives the exact shape of the missing
constant, and the soft split removes the compact-cutoff obstruction.
The remaining fixed-tree task is finite and Yang--Mills-specific:
put the Wilson remainder, the soft quadratic selector vertex, the
current insertion, and the three score marks into one strong
analytic norm; then bound every edge of (65.3) by the V45 weighted
covariance without spending the strong radius repeatedly.

\begin{assessment}[Progress after V65]
\label{ass:e8v65-status}
V65 performs the compulsory current-file audit and imports the
positive Gaussian forest theorem legitimately.  It replaces the
abstract \(C_{\rm GF}\) by the explicit tree expression (65.8),
identifies why the compact smooth selector fails the all-tree
analytic card, and constructs an exact entire soft split.  The
soft-bad primitive and soft-good perturbation are both explicitly
\(O(\beta^{-1/5})\).  What remains is the one-time strong-to-weak
analytic-radius estimate (65.4) for the concrete Yang--Mills vertex
bank.
\end{assessment}

\begin{decision}[V66 target]
\label{dec:e8v65-next}
V66 should define the concrete strong analytic norm for the four
marked Yang--Mills source types and prove the fixed-tree estimate
(65.4).  It must use a global tree convolution, not iterated binary
strong-to-weak estimates, and it must keep the soft-transition and
far-bad primitives as distinct owner types.  If the source radii
collapse with \(\beta\), their powers must be displayed rather than
absorbed into \(K_T\).
\end{decision}

\begin{firewall}[Claim boundary after V65]
\label{fire:e8v65-claim}
V65 transfers the forest theorem and repairs selector analyticity,
but it does not yet prove the concrete fixed-tree estimate (65.4),
close the numerical branch (63.18), evaluate the raw pushed
\(\zeta_{\rm Jac}\), compare the compensated owner with
\({\cal O}_1\), prove anisotropy contraction, evaluate the remaining
cards, complete the LERB, establish the universal beta coefficient,
close the nonlinear/complex quotient matching card, construct the
full RG trajectory, remove the regulators, establish
Osterwalder--Schrader reconstruction, or prove the Yang--Mills mass
gap.
\end{firewall}

\subsection{Parent record}

This V65 note was formed by appending the preceding material to V64,
whose validated record was
\[
\begin{array}{c|c}
\text{lines}&27526\\
\text{bytes}&968943\\
\text{SHA-256}&
\texttt{7B96EB8AFDCAA460856EE8C0E56513F946231D7E5C9EEB62EF32A3A027F7BE9A}.
\end{array}
\]

% =============================================================================
% END APPENDED EIGHTH-SERIES V65 MATERIAL
% =============================================================================

\end{document}
